\clearpage
\thispagestyle{plain}
\markboth{예비사항}{예비사항}
\oldpage[0\textsubscript{III}]{5}
\phantomsection
\label{chapter:0III-fr}

\begin{center}
{\large 제0장 (계속)\footnotemark\par}
\vspace{1.5em}
{\Large\bfseries 예비사항\par}
\end{center}

\footnotetext{참조를 쉽게 찾을 수 있도록, 이제부터 제I장과 함께 출판된
제0장의 절은 기호 $0_{\mathrm I}$로 인용한다.}

\setcounter{section}{7}
\section{표현 가능 함자}
\label{section:0.8-fr}

\setcounter{subsection}{0}
\subsection{표현 가능 함자}
\label{subsection:0.8.1-fr}

\begin{env}[8.1.1]
\label{0.8.1.1-fr}
집합의 범주를 $\mathrm{Ens}$로 나타내기로 한다. $C$를 하나의
범주라 하자~; $C$의 두 대상 $X$, $Y$에 대하여
$h_X(Y)=\operatorname{Hom}(Y,X)$로 놓고~; $C$의 임의의 사상
$u:Y\to Y'$에 대하여 $h_X(u)$를 $\operatorname{Hom}(Y',X)$에서
$\operatorname{Hom}(Y,X)$로 가는 함수 $v\mapsto vu$로 나타낸다. 이
정의들에 따르면 $h_X:C\to\mathrm{Ens}$가 \emph{반변 함자}, 즉 범주
$C$의 쌍대범주인 $C^0$에서 범주 $\mathrm{Ens}$로 가는 공변 함자들의
범주 $\operatorname{Hom}(C^0,\mathrm{Ens})$의 대상임은 자명하다
(T, 1.7, d) 및 [29]).
\end{env}

\begin{env}[8.1.2]
\label{0.8.1.2-fr}
이제 $w:X\to X'$을 $C$의 사상이라 하자~; 임의의 $Y\in C$와 임의의
$v\in\operatorname{Hom}(Y,X)=h_X(Y)$에 대하여
$wv\in\operatorname{Hom}(Y,X')=h_{X'}(Y)$이다~; $h_w(Y)$를
$h_X(Y)$에서 $h_{X'}(Y)$로 가는 함수 $v\mapsto wv$로 나타내자.
$C$의 임의의 사상 $u:Y\to Y'$에 대하여 다음 도식은
\[
  \xymatrix{
    h_X(Y')\ar[r]^{h_X(u)}\ar[d]_{h_w(Y')} &
    h_X(Y)\ar[d]^{h_w(Y)}\\
    h_{X'}(Y')\ar[r]_{h_{X'}(u)} &
    h_{X'}(Y)
  }
\]
가환한다~; 다시 말해 $h_w$는 \emph{자연변환}
$h_X\to h_{X'}$이다 (T, 1.2).
또는 범주 $\operatorname{Hom}(C^0,\mathrm{Ens})$의 사상이다
(T, 1.7, d)). 따라서 $h_X$와 $h_w$의 정의는 \emph{표준 공변 함자}
\[
  h:C\to\operatorname{Hom}(C^0,\mathrm{Ens}),\qquad X\mapsto h_X.
  \tag{8.1.2.1}
\]
를 정의한다.
\end{env}

\begin{env}[8.1.3]
\label{0.8.1.3-fr}
$X$를 $C$의 대상, $F$를 $C$에서 $\mathrm{Ens}$로 가는 반변 함자
($\operatorname{Hom}(C^0,\mathrm{Ens})$의 대상)라 하자.
$g:h_X\to F$를 \emph{자연변환}이라 하자~: 임의의 $Y\in C$에 대하여
\oldpage[0\textsubscript{III}]{6}
$g(Y)$는 $h_X(Y)$에서 $F(Y)$로 가는 함수이고, $C$의 임의의 사상
$u:Y\to Y'$에 대하여 다음 도식이 가환한다.
\[
  \xymatrix{
    h_X(Y')\ar[r]^{h_X(u)}\ar[d]_{g(Y')} &
    h_X(Y)\ar[d]^{g(Y)}\\
    F(Y')\ar[r]_{F(u)} & F(Y)
  }
  \tag{8.1.3.1}\label{0.8.1.3.1-fr}
\]
특히 함수 $g(X):h_X(X)=\operatorname{Hom}(X,X)\to F(X)$가 있으므로,
다음 원소를 얻는다.
\[
  \alpha(g)=(g(X))(1_X)\in F(X)
  \tag{8.1.3.2}
\]
따라서 표준 함수
\[
  \alpha:\operatorname{Hom}(h_X,F)\to F(X).
  \tag{8.1.3.3}
\]
가 정의된다.

반대로 원소 $\xi\in F(X)$를 생각하자~; $C$의 임의의 사상
$v:Y\to X$에 대하여 $F(v)$는 $F(X)$에서 $F(Y)$로 가는 함수이다~;
다음 함수를 생각하자.
\[
  v\mapsto(F(v))(\xi)
  \tag{8.1.3.4}
\]
이는 $h_X(Y)$에서 $F(Y)$로 가는 함수이다~; 이 함수를
$(\beta(\xi))(Y)$로 나타내면,
\[
  \beta(\xi):h_X\to F
  \tag{8.1.3.5}
\]
는 \emph{자연변환}이다. 실제로 $C$의 임의의 사상 $u:Y\to Y'$에
대하여 $(F(vu))(\xi)=(F(u)\circ F(v))(\xi)$이고, 이는
$g=\beta(\xi)$일 때 (\hyperref[0.8.1.3.1-fr]{8.1.3.1})의 가환성을
확인해 준다. 이로써 표준 함수
\[
  \beta:F(X)\to\operatorname{Hom}(h_X,F).
  \tag{8.1.3.6}
\]
를 정의하였다.
\end{env}

\begin{proposition}[8.1.4]
\label{0.8.1.4-fr}
함수 $\alpha$와 $\beta$는 서로 역인 전단사이다.
\end{proposition}

$\xi\in F(X)$에 대하여 $\alpha(\beta(\xi))$를 계산하자~; 임의의
$Y\in C$에 대하여 $(\beta(\xi))(Y)$는 $h_X(Y)$에서 $F(Y)$로 가는
함수 $g_1(Y):v\mapsto(F(v))(\xi)$이다. 따라서
\[
  \alpha(\beta(\xi))=(g_1(X))(1_X)=(F(1_X))(\xi)
  =1_{F(X)}(\xi)=\xi.
\]
이제 $g\in\operatorname{Hom}(h_X,F)$에 대하여
$\beta(\alpha(g))$를 계산하자~; 임의의 $Y\in C$에 대하여
$(\beta(\alpha(g)))(Y)$는 함수
$v\mapsto(F(v))((g(X))(1_X))$이다~; (\hyperref[0.8.1.3.1-fr]{8.1.3.1})의
가환성에 따라, 이 함수는 $h_X(v)$의 정의에 의해
$v\mapsto(g(Y))((h_X(v))(1_X))=(g(Y))(v)$와 다르지 않다. 다시 말해
$g(Y)$와 같으며, 이로써 명제가 증명된다.

\begin{env}[8.1.5]
\label{0.8.1.5-fr}
범주 $C$의 \emph{부분범주} $C'$은 다음 조건으로 정의됨을 상기하자.
$C'$의 대상들은 $C$의 대상들이고, $X'$, $Y'$이 $C'$의 두 대상이면
\emph{$C'$ 안에서} $X'\to Y'$인 사상들의 집합
$\operatorname{Hom}_{C'}(X',Y')$은 \emph{$C$ 안에서} $X'\to Y'$인
사상들의 집합 $\operatorname{Hom}_C(X',Y')$의 부분집합이며, 표준적인
<<~사상의 합성~>> 함수
\[
  \operatorname{Hom}_{C'}(X',Y')\times\operatorname{Hom}_{C'}(Y',Z')
  \to\operatorname{Hom}_{C'}(X',Z')
\]
\oldpage[0\textsubscript{III}]{7}
는 다음 표준 함수의 제한이다.
\[
  \operatorname{Hom}_C(X',Y')\times\operatorname{Hom}_C(Y',Z')
  \to\operatorname{Hom}_C(X',Z').
\]

$C'$의 임의의 두 대상에 대하여
$\operatorname{Hom}_{C'}(X',Y')=\operatorname{Hom}_C(X',Y')$이면,
$C'$을 $C$의 \emph{충만한 부분범주}라고 한다. 이때 $C'$의 대상들과
동형인 $C$의 대상들로 이루어진 $C$의 부분범주 $C''$도 $C$의 충만한
부분범주이며, 쉽게 확인할 수 있듯이 $C'$과 \emph{동치}이다 (T, 1.2).

공변 함자 $F:C_1\to C_2$에 대하여, $C_1$의 임의의 두 대상 $X_1$,
$Y_1$에 대해 $\operatorname{Hom}(X_1,Y_1)$에서
$\operatorname{Hom}(F(X_1),F(Y_1))$로 가는 함수 $u\mapsto F(u)$가
\emph{전단사}이면, $F$가 \emph{충만하고 충실하다}고 한다~; 이 조건은
$C_2$의 부분범주 $F(C_1)$이 \emph{충만함}을 함의한다. 더욱이 두 대상
$X_1$, $X'_1$의 상이 같은 대상 $X_2$이면, $F(u)=1_{X_2}$를 만족하는
동형사상 $u:X_1\to X'_1$가 유일하게 존재한다. 이제 $F(C_1)$의 임의의
대상 $X_2$에 대하여 $F(X_1)=X_2$를 만족하는 $C_1$의 대상 $X_1$ 중
하나를 $G(X_2)$로 택하자 ($G$는 선택공리를 사용하여 정의한다)~;
$F(C_1)$의 임의의 사상 $v:X_2\to Y_2$에 대하여 $G(v)$를
$F(u)=v$를 만족하는 유일한 사상 $u:G(X_2)\to G(Y_2)$로 정한다~;
그러면 $G$는 $F(C_1)$에서 $C_1$로 가는 \emph{함자}이다~;
$FG$는 $F(C_1)$의 항등함자이고, 앞의 논의에 따라 함자의 동형사상
$\varphi:1_{C_1}\to GF$가 존재하여 $F$, $G$, $\varphi$ 및 항등사상
$1_{F(C_1)}\to FG$가 범주 $C_1$과 $C_2$의 충만한 부분범주
$F(C_1)$ 사이의 \emph{동치}를 정의한다 (T, 1.2).
\end{env}

\begin{env}[8.1.6]
\label{0.8.1.6-fr}
함자 $F$가 $h_{X'}$이고 $X'$이 $C$의 임의의 대상인 경우에 명제
(\hyperref[0.8.1.4-fr]{8.1.4})를 적용하자~; 이때 함수
$\beta:\operatorname{Hom}(X,X')\to\operatorname{Hom}(h_X,h_{X'})$는
(\hyperref[0.8.1.2-fr]{8.1.2})에서 정의한 함수 $w\mapsto h_w$와
다르지 않다~; 이 함수가 \emph{전단사}이므로,
(\hyperref[0.8.1.5-fr]{8.1.5})의 용어를 사용하면 다음을 얻는다~:
\end{env}

\begin{proposition}[8.1.7]
\label{0.8.1.7-fr}
표준 함자 $h:C\to\operatorname{Hom}(C^0,\mathrm{Ens})$는 충만하고
충실하다.
\end{proposition}

\begin{env}[8.1.8]
\label{0.8.1.8-fr}
$F$를 $C$에서 $\mathrm{Ens}$로 가는 반변 함자라 하자~; 어떤 대상
$X\in C$가 존재하여 $F$가 $h_X$와 \emph{동형}이면, $F$가
\emph{표현 가능하다}고 한다~; (\hyperref[0.8.1.7-fr]{8.1.7})에 따라
$X\in C$와 함자의 동형사상 $g:h_X\to F$를 주면 $X$는 유일한
동형사상까지 결정된다. 명제 (\hyperref[0.8.1.7-fr]{8.1.7})는 또한
$h$가 $C$와, $\operatorname{Hom}(C^0,\mathrm{Ens})$에서 \emph{표현 가능한
반변 함자}들로 이루어진 충만한 부분범주 사이의 \emph{동치}를
정의한다는 뜻이다. 한편 (\hyperref[0.8.1.4-fr]{8.1.4})에 따라
자연변환 $g:h_X\to F$를 주는 것은 원소 $\xi\in F(X)$를 주는 것과
동치이다~; $g$가 \emph{동형사상}이라는 것은 $\xi$에 대하여 다음
조건과 동치이다~: \emph{$C$의 임의의 대상 $Y$에 대하여
$\operatorname{Hom}(Y,X)$에서 $F(Y)$로 가는 함수
$v\mapsto(F(v))(\xi)$가 전단사이다}. $\xi$가 이 조건을 만족하면,
쌍 $(X,\xi)$가 표현 가능 함자 $F$를 \emph{표현한다}고 한다. 용어를
다소 남용하여, 어떤 $\xi\in F(X)$에 대하여 $(X,\xi)$가 $F$를
표현할 때, 즉 $h_X$가 $F$와 동형일 때, 대상 $X\in C$가 $F$를
표현한다고도 한다.

$F$, $F'$을 $C$에서 $\mathrm{Ens}$로 가는 표현 가능한 두 반변
함자라 하고, $h_X\to F$와 $h_{X'}\to F'$을 함자의 두 동형사상이라
하자. 그러면 (\hyperref[0.8.1.6-fr]{8.1.6})에 따라
$\operatorname{Hom}(X,X')$과, 자연변환 $F\to F'$들의 집합
$\operatorname{Hom}(F,F')$ 사이에 표준적인 일대일 대응이 있다.
\end{env}

\begin{env}[8.1.9]
\label{0.8.1.9-fr}
\textit{예. I: 사영극한.} 표현 가능한 반변 함자의 개념은 특히
통상적인 <<~보편 문제의
\oldpage[0\textsubscript{III}]{8}
해~>>라는 개념의 <<~쌍대~>> 개념을 포괄한다. 더 일반적으로,
\emph{사영극한}의 개념이 표현 가능 함자의 개념의 특수한 경우임을
살펴보겠다. 범주 $C$에서 \emph{사영계}는 준순서 집합 $I$, $C$의
대상들의 족 $(A_\alpha)_{\alpha\in I}$, 그리고
$\alpha\leq\beta$인 임의의 두 지표 $(\alpha,\beta)$에 대한 사상
$u_{\alpha\beta}:A_\beta\to A_\alpha$를 주어 정의함을 상기하자.
$C$에서 이 사영계의 \emph{사영극한}은 $C$의 대상 $B$
($\varprojlim A_\alpha$로 표기한다)와 각 $\alpha\in I$에 대한 사상
$u_\alpha:B\to A_\alpha$로 이루어지며, 다음 조건들을 만족한다~:
1\textsuperscript{o} $\alpha\leq\beta$일 때
$u_\alpha=u_{\alpha\beta}u_\beta$이다~;
2\textsuperscript{o} $C$의 임의의 대상 $X$와, $\alpha\leq\beta$일 때
$v_\alpha=u_{\alpha\beta}v_\beta$를 만족하는 사상
$v_\alpha:X\to A_\alpha$들의 임의의 족 $(v_\alpha)_{\alpha\in I}$에
대하여, 모든 $\alpha\in I$에 대해 $v_\alpha=u_\alpha v$를 만족하는
유일한 사상 $v:X\to B$가 존재한다 ($\varprojlim v_\alpha$로 표기한다)
(T, 1.8). 이는 다음과 같이 해석된다~: $u_{\alpha\beta}$들은 표준적으로
함수
\[
  \overline{u}_{\alpha\beta}:\operatorname{Hom}(X,A_\beta)
  \to\operatorname{Hom}(X,A_\alpha)
\]
들을 정의하며, 이 함수들은 집합들의 \emph{사영계}
$(\operatorname{Hom}(X,A_\alpha),\overline{u}_{\alpha\beta})$를
정의한다. 그리고 $(v_\alpha)$는 정의에 따라 집합
$\varprojlim_\alpha\operatorname{Hom}(X,A_\alpha)$의 원소이다~;
$X\mapsto\varprojlim_\alpha\operatorname{Hom}(X,A_\alpha)$가 $C$에서
$\mathrm{Ens}$로 가는 \emph{반변 함자}임은 명백하며, 사영극한 $B$가
존재한다는 것은 $(v_\alpha)\mapsto\varprojlim v_\alpha$가 $X$에 관한
함자들의 \emph{동형사상}
\[
  \varprojlim_\alpha\operatorname{Hom}(X,A_\alpha)
  \xrightarrow{\sim}\operatorname{Hom}(X,B)
  \tag{8.1.9.1}
\]
이라는 것, 다시 말해 함자
$X\mapsto\varprojlim_\alpha\operatorname{Hom}(X,A_\alpha)$가
\emph{표현 가능하다}는 것과 동치이다.
\end{env}

\begin{env}[8.1.10]
\label{0.8.1.10-fr}
\textit{예. II: 끝 대상.} $C$를 범주라 하고, $\{a\}$를 한 원소로만
이루어진 집합이라 하자. 임의의 대상 $X$에 집합 $\{a\}$를 대응시키고,
$C$의 임의의 사상 $X\to X'$에 유일한 함수 $\{a\}\to\{a\}$를
대응시키는 반변 함자 $F:C\to\mathrm{Ens}$를 생각하자. 이 함자가
\emph{표현 가능하다}는 것은, 임의의 $Y\in C$에 대하여
$\operatorname{Hom}(Y,e)=h_e(Y)$가 \emph{한 원소로만 이루어지는} 대상
$e\in C$가 존재한다는 뜻이다~; 이러한 $e$를 $C$의 \emph{끝 대상}이라
한다. $C$의 두 끝 대상이 동형임은 명백하다 (이에 따라, 일반적으로
선택공리를 사용하여, $C$의 끝 대상을 정하고 이를 $e_C$로 나타낼 수
있다). 예를 들어 범주 $\mathrm{Ens}$에서 끝 대상은 한 원소로만
이루어진 집합이다~; 체 $K$ 위의 \emph{증대 대수}의 범주(사상은
증대와 양립하는 대수 준동형이다)에서는 $K$가 끝 대상이다~;
\emph{$S$-준스킴}의 범주 (I, 2.5.1)에서는 $S$가 끝 대상이다.
\end{env}

\begin{env}[8.1.11]
\label{0.8.1.11-fr}
범주 $C$의 두 대상 $X$, $Y$에 대하여
$h'_X(Y)=\operatorname{Hom}(X,Y)$로 놓고, 임의의 사상 $u:Y\to Y'$에
대하여 $h'_X(u)$를 $\operatorname{Hom}(X,Y)$에서
$\operatorname{Hom}(X,Y')$로 가는 함수 $v\mapsto uv$로 두자~;
그러면 $h'_X$는 \emph{공변 함자} $C\to\mathrm{Ens}$이다. 따라서
(\hyperref[0.8.1.2-fr]{8.1.2})에서와 같이 표준 공변 함자
$h':C^0\to\operatorname{Hom}(C,\mathrm{Ens})$의 정의를 얻는다~;
$C$에서 $\mathrm{Ens}$로 가는 \emph{공변 함자} $F$, 즉
$\operatorname{Hom}(C,\mathrm{Ens})$의 대상은, 어떤 대상 $X\in C$
(필연적으로 유일한 동형사상까지 유일하다)가 존재하여 $F$가
\emph{$h'_X$와 동형}이면 \emph{표현 가능하다}고 한다~; 이 개념에
대하여 앞의 논의와 <<~쌍대인~>> 논의를 전개하는 일은 독자에게
맡긴다. 이 개념은 \emph{귀납극한}의 개념, 특히 통상적인
<<~보편 문제의 해~>>라는 개념을 포괄한다.
\end{env}

\subsection{범주에서의 대수적 구조}
\label{subsection:0.8.2-ko}

\begin{env}[8.2.1]
\label{0.8.2.1-ko}
\oldpage[0\textsubscript{III}]{9}
$C$에서 $\mathrm{Ens}$로 가는 두 반변 함자 $F$, $F'$이 주어졌다고 하자.
모든 대상 $Y\in C$에 대하여
$(F\times F')(Y)=F(Y)\times F'(Y)$로 놓고, $C$의 모든 사상
$u:Y\to Y'$에 대하여 $(F\times F')(u)=F(u)\times F'(u)$로 놓음을
상기하자. 뒤의 사상은 $F(Y')\times F'(Y')$에서 $F(Y)\times F'(Y)$로
가는 사상
$(t,t')\mapsto((F(u))(t),(F'(u))(t'))$이다. 따라서
$F\times F':C\to\mathrm{Ens}$는 
\emph{반변 함자}이다(이는 사실 범주
$\operatorname{Hom}(C^0,\mathrm{Ens})$에서 대상 $F$, $F'$의
\emph{곱}과 다르지 않다). 대상 $X\in C$가 주어졌을 때, 함자적 사상
\[
  \gamma_X:h_X\times h_X\to h_X
  \tag{8.2.1.1}
\]
을 
\emph{$X$ 위의 내부 합성법칙}이라 하겠다.

다시 말해 (T, 1.2), 모든 대상 $Y\in C$에 대하여 $\gamma_X(Y)$는
사상 $h_X(Y)\times h_X(Y)\to h_X(Y)$이고(따라서 정의상 집합
$h_X(Y)$ 위의 \emph{내부 합성법칙}이다), $C$의 모든 사상
$u:Y\to Y'$에 대하여 도식
\[
  \xymatrix{
    h_X(Y')\times h_X(Y')
      \ar[rr]^{h_X(u)\times h_X(u)}\ar[dd]_{\gamma_X(Y')} &&
    h_X(Y)\times h_X(Y)\ar[dd]^{\gamma_X(Y)}\\\\
    h_X(Y')\ar[rr]_{h_X(u)} && h_X(Y)
  }
\]
이 가환이라는 조건을 만족한다. 이는 합성법칙 $\gamma_X(Y)$와
$\gamma_X(Y')$에 대하여 $h_X(u)$가 $h_X(Y')$에서 $h_X(Y)$로 가는
\emph{준동형}임을 뜻한다.

마찬가지로 $C$의 두 대상 $Z$, $X$가 주어졌을 때, 함자적 사상
\[
  \omega_{X,Z}:h_Z\times h_X\to h_X
  \tag{8.2.1.2}
\]
을 
\emph{$Z$를 작용자 영역으로 하는 $X$ 위의 외부 합성법칙}이라 한다.

위와 같이 모든 $Y\in C$에 대하여 $\omega_{X,Z}(Y)$는
$h_Z(Y)$를 작용자 영역으로 하는 $h_X(Y)$ 위의 외부 합성법칙이고,
모든 사상 $u:Y\to Y'$에 대하여 $h_X(u)$와 $h_Z(u)$는
$(h_Z(Y'),h_X(Y'))$에서 $(h_Z(Y),h_X(Y))$로 가는
\emph{디-준동형}을 이룸을 알 수 있다.
\end{env}

\begin{env}[8.2.2]
\label{0.8.2.2-ko}
$C$의 두 번째 대상 $X'$에 내부 합성법칙 $\gamma_{X'}$가 주어졌다고
하자. $C$의 사상 $w:X\to X'$가 모든 $Y\in C$에 대하여
$h_w(Y):h_X(Y)\to h_{X'}(Y)$가 합성법칙 $\gamma_X(Y)$와
$\gamma_{X'}(Y)$에 관한 준동형이면, $w$를 이 합성법칙들에 관한
\emph{준동형}이라 하겠다. $X''$이 내부 합성법칙 $\gamma_{X''}$를 갖춘
$C$의 세 번째
\oldpage[0\textsubscript{III}]{10}
대상이고, $w':X'\to X''$이 $\gamma_{X'}$와 $\gamma_{X''}$에 관한
준동형인 $C$의 사상이면, 사상 $w'w:X\to X''$이 합성법칙
$\gamma_X$와 $\gamma_{X''}$에 관한 준동형임은 명백하다. $C$의 동형
$w:X\xrightarrow{\sim}X'$가 이 합성법칙들에 관한 준동형이고 그 역 사상
$w^{-1}$가 합성법칙 $\gamma_{X'}$와 $\gamma_X$에 관한 준동형이면,
$w$를 \emph{합성법칙 $\gamma_X$와 $\gamma_{X'}$에 관한 동형}이라 한다.

외부 합성법칙을 갖춘 $C$의 대상쌍들 사이의 \emph{디-준동형}도 같은
방식으로 정의한다.
\end{env}

\begin{env}[8.2.3]
\label{0.8.2.3-ko}
대상 $X\in C$ 위의 내부 합성법칙 $\gamma_X$가 \emph{모든}
$Y\in C$에 대하여 $h_X(Y)$ 위의 \emph{군법칙} $\gamma_X(Y)$를 주면,
이 법칙을 갖춘 $X$를 \emph{$C$-군} 또는 \emph{범주 $C$의 군 대상}이라
한다. \emph{$C$-환}, \emph{$C$-가군} 등도 마찬가지로 정의한다.
\end{env}

\begin{env}[8.2.4]
\label{0.8.2.4-ko}
대상 $X\in C$와 자기 자신의 \emph{곱} $X\times X$가 $C$에서
존재한다고 하자. 이는 사영극한의 특수한 경우이므로
(\hyperref[0.8.1.9-fr]{8.1.9}), 정의에 따라 표준 동형을 제외하고
$h_{X\times X}=h_X\times h_X$이다. 따라서 $X$ 위의 내부 합성법칙은
함자적 사상 $\gamma_X:h_{X\times X}\to h_X$로 볼 수 있고, 이에 따라
(\hyperref[0.8.1.6-fr]{8.1.6}) $h_{c_X}=\gamma_X$인 원소
$c_X\in\operatorname{Hom}(X\times X,X)$를 표준적으로 정한다. 그러므로
이 경우 $X$ 위의 내부 합성법칙을 주는 것은 사상 $X\times X\to X$를
주는 것과 동치이다. $C$가 범주 $\mathrm{Ens}$일 때에는 집합 위의 내부
합성법칙이라는 고전적 개념을 다시 얻는다. 곱 $Z\times X$가 $C$에서
존재하는 경우의 외부 합성법칙에 대해서도 같은 결과가 성립한다.
\end{env}

\begin{env}[8.2.5]
\label{0.8.2.5-ko}
앞의 표기 아래에서 $X\times X\times X$도 $C$에서 존재한다고 하자.
함자를 표현하는 대상으로서 곱의 특징짓기
(\hyperref[0.8.1.9-fr]{8.1.9})로부터 표준 동형
\[
  (X\times X)\times X\xrightarrow{\sim}X\times X\times X
  \xrightarrow{\sim}X\times(X\times X)
\]
들이 존재한다. $X\times X\times X$를 $(X\times X)\times X$와
표준적으로 동일시하면, 모든 $Y\in C$에 대하여 사상
$\gamma_X(Y)\times 1_{h_X(Y)}$는 $h_{c_X\times 1_X}(Y)$와 동일시된다.
따라서 모든 $Y\in C$에 대하여 내부 합성법칙 $\gamma_X(Y)$가
\emph{결합적}이라고 하는 것과 사상들의 도식
\[
  \xymatrix{
    h_X(Y)\times h_X(Y)\times h_X(Y)
      \ar[rr]^{\gamma_X(Y)\times 1}\ar[dd]_{1\times\gamma_X(Y)} &&
    h_X(Y)\times h_X(Y)\ar[dd]^{\gamma_X(Y)}\\\\
    h_X(Y)\times h_X(Y)\ar[rr]_{\gamma_X(Y)} && h_X(Y)
  }
\]
이 가환이라고 하는 것, 그리고 사상들의 도식
\oldpage[0\textsubscript{III}]{11}
\[
  \xymatrix{
    X\times X\times X
      \ar[rr]^{c_X\times 1_X}\ar[dd]_{1_X\times c_X} &&
    X\times X\ar[dd]^{c_X}\\\\
    X\times X\ar[rr]_{c_X} && X
  }
\]
이 가환이라고 하는 것은 서로 동치이다.
\end{env}

\begin{env}[8.2.6]
\label{0.8.2.6-ko}
(\hyperref[0.8.2.5-ko]{8.2.5})의 가정 아래에서 모든 $Y\in C$에 대하여
내부 합성법칙 $\gamma_X(Y)$가 \emph{군법칙}임을 나타내려면, 한편으로
그것이 결합적임을 나타내고 다른 한편으로 군에서의 \emph{역원}의 성질을
갖는 사상 $\alpha_X(Y):h_X(Y)\to h_X(Y)$가 존재함을 나타내야 한다.
$C$의 모든 사상 $u:Y\to Y'$에 대하여 $h_X(u)$가 군 $h_X(Y')$에서
$h_X(Y)$로 가는 준동형이어야 함을 앞에서 보았으므로, 먼저
$\alpha_X:h_X\to h_X$는 \emph{함자적 사상}이어야 한다. 한편 군 $G$에서
역원을 취하는 사상 $s\mapsto s^{-1}$의 특징적인 성질은 항등원을 쓰지
않고도 나타낼 수 있다. 즉 두 합성사상
\[
  (s,t)\mapsto(s,s^{-1},t)\mapsto(s,s^{-1}t)\mapsto s(s^{-1}t),
\]
\[
  (s,t)\mapsto(s,s^{-1},t)\mapsto(s,ts^{-1})\mapsto(ts^{-1})s
\]
이 $G\times G$에서 $G$로 가는 두 번째 사영 $(s,t)\mapsto t$와 같다고
쓰면 충분하다. (\hyperref[0.8.1.3-fr]{8.1.3})에 따라
$a_X\in\operatorname{Hom}(X,X)$인 $\alpha_X=h_{a_X}$이다. 그러면 앞의
첫째 조건은 합성사상
\[
  X\times X\xrightarrow{(1_X,a_X)\times 1_X}X\times X\times X
  \xrightarrow{1_X\times c_X}X\times X\xrightarrow{c_X}X
\]
이 $C$에서의 두 번째 사영 $X\times X\to X$라는 것을 나타내며,
둘째 조건도 같은 방식으로 옮겨진다.
\end{env}

\begin{env}[8.2.7]
\label{0.8.2.7-ko}
이제 $C$에 \emph{끝 대상} $e$가 존재한다고 하자
(\hyperref[0.8.1.10-fr]{8.1.10}). 여전히 모든 $Y\in C$에 대하여
$\gamma_X(Y)$가 $h_X(Y)$ 위의 군법칙이라고 하고, $\eta_X(Y)$로
$\gamma_X(Y)$의 항등원을 나타내자. $C$의 모든 사상 $u:Y\to Y'$에
대하여 $h_X(u)$가 군 준동형이므로
$\eta_X(Y)=(h_X(u))(\eta_X(Y'))$이다. 특히 $Y'=e$로 취하면 $u$는
$\operatorname{Hom}(Y,e)$의 유일한 원소 $\varepsilon$이므로,
$\eta_X(e)$가 모든 $Y\in C$에 대하여 $\eta_X(Y)$를 완전히 정함을 알 수
있다. $h_X(X)=\operatorname{Hom}(X,X)$인 군의 항등원
$e_X=\eta_X(X)$로 놓자. 도식
\[
  \xymatrix{
    h_X(e)\ar[r]^{h_X(\varepsilon)}\ar[d]_{h_{e_X}(e)} &
    h_X(Y)\ar[d]^{h_{e_X}(Y)}\\
    h_X(e)\ar[r]_{h_X(\varepsilon)} & h_X(Y)
  }
\]
의 가환성((\hyperref[0.8.1.2-fr]{8.1.2}) 참조)은 집합 $h_X(Y)$에서
사상 $h_{e_X}(Y)$가 모든 원소를 항등원으로 보내는
\oldpage[0\textsubscript{III}]{12}
$s\mapsto\eta_X(Y)$와 다르지 않음을 보인다. 그러면 모든 $Y\in C$에
대하여 $\eta_X(Y)$가 $\gamma_X(Y)$의 항등원이라는 것은 합성사상
\[
  X\xrightarrow{(1_X,1_X)}X\times X
  \xrightarrow{1_X\times e_X}X\times X\xrightarrow{c_X}X,
\]
그리고 $1_X$와 $e_X$의 자리를 바꾼 유사한 사상이 둘 다 $1_X$와
\emph{같다}고 하는 것과 동치임을 확인할 수 있다.
\end{env}

\begin{env}[8.2.8]
\label{0.8.2.8-ko}
물론 범주에서의 대수적 구조에 관한 예는 어렵지 않게 얼마든지 더 들 수
있다. 군의 예는 충분히 상세히 다루었지만, 앞으로 만나게 될 대수적 구조의
예에서는 이와 비슷한 논의를 전개하는 일을 대체로 독자에게 맡기겠다.
\end{env}

\section{구성가능 집합}
\label{section:0.9-ko}

\subsection{구성가능 집합}
\label{subsection:0.9.1-ko}

\begin{definition}[9.1.1]
\label{0.9.1.1-ko}
연속사상 $f:X\to Y$가 $Y$의 모든 준콤팩트 열린집합 $U$에 대하여
$f^{-1}(U)$가 준콤팩트이면, $f$를 \emph{준콤팩트}라 한다. 위상공간
$X$의 부분집합 $Z$에 대한 표준 단사사상 $Z\to X$가 준콤팩트이면,
$Z$를 \emph{$X$에서 역콤팩트}라 한다. 다시 말해 $X$의 모든 준콤팩트
열린집합 $U$에 대하여 $U\cap Z$가 준콤팩트인 경우이다.
\end{definition}

$X$의 \emph{닫힌} 부분집합은 $X$에서 역콤팩트이지만, $X$의 준콤팩트
부분집합은 반드시 $X$에서 역콤팩트인 것은 아니다. $X$가 준콤팩트이면
$X$에서 역콤팩트인 모든 열린 부분집합은 준콤팩트이다. $X$에서
역콤팩트인 집합들의 모든 \emph{유한} 합집합은 $X$에서 역콤팩트임이
명백하다. 이는 준콤팩트 집합들의 모든 유한 합집합이 준콤팩트이기
때문이다. $X$에서 역콤팩트인 열린집합들의 모든 유한 교집합은
$X$에서 역콤팩트인 열린집합이다. \emph{국소 뇌터} 공간 $X$에서는
모든 준콤팩트 집합이 뇌터 부분공간이고, 따라서 $X$의 \emph{모든
부분집합}은 $X$에서 역콤팩트이다.

\begin{definition}[9.1.2]
\label{0.9.1.2-ko}
위상공간 $X$가 주어졌다고 하자. $X$에서 역콤팩트인 모든 열린
부분집합을 포함하고 유한 교집합과 여집합을 취하는 것에 관하여 닫혀
있는 $X$의 부분집합들의 최소 집합 $\mathfrak{F}$에 속하는 $X$의
부분집합을 \emph{구성가능}이라 한다(따라서 $\mathfrak{F}$는 유한
합집합에 관해서도 닫혀 있다).
\end{definition}

\begin{proposition}[9.1.3]
\label{0.9.1.3-ko}
$X$의 부분집합이 구성가능이기 위한 필요충분조건은 $X$에서 역콤팩트인
열린집합 $U$, $V$에 대한 $U\cap\complement{V}$ 꼴의 집합들의 유한
합집합인 것이다.
\end{proposition}

이 조건이 충분함은 명백하다. 필요함을 보이기 위하여, $X$에서
역콤팩트인 열린집합 $U$, $V$에 대한 $U\cap\complement{V}$ 꼴의
집합들의 유한 합집합으로 이루어진 집합 $\mathfrak{G}$를 생각하자.
$\mathfrak{G}$의 모든 원소의 여집합이 $\mathfrak{G}$에 속함을 보이면
충분하다. 따라서 $I$가 유한하고 $U_i$, $V_i$가 $X$에서 역콤팩트인
열린집합인
$Z=\bigcup_{i\in I}(U_i\cap\complement{V_i})$를 취하자. 이때
$\complement{Z}=\bigcap_{i\in I}(V_i\cup\complement{U_i})$이므로,
$\complement{Z}$는 몇 개의 $V_i$와 몇 개의 $\complement{U_i}$를
교차하여 얻는 집합들의 유한 합집합이다. 따라서 각 집합은
\oldpage[0\textsubscript{III}]{13}
$V\cap\complement{U}$ 꼴이며, 여기서 $U$는 몇 개의 $U_i$의
합집합이고 $V$는 몇 개의 $V_i$의 교집합이다. 그런데 $X$에서
역콤팩트인 열린집합들의 유한 합집합과 유한 교집합은 역콤팩트인
열린집합임을 앞에서 보았으므로 결론이 따른다.

\begin{corollary}[9.1.4]
\label{0.9.1.4-ko}
$X$의 모든 구성가능 부분집합은 $X$에서 역콤팩트이다.
\end{corollary}

$U$, $V$가 $X$에서 역콤팩트인 열린집합일 때
$U\cap\complement{V}$가 $X$에서 역콤팩트임을 보이면 충분하다.
그런데 $W$가 $X$의 준콤팩트 열린집합이면
$W\cap U\cap\complement{V}$는 준콤팩트 공간 $W\cap U$의 닫힌
부분집합이므로 준콤팩트이다.

특히 다음을 얻는다.

\begin{corollary}[9.1.5]
\label{0.9.1.5-ko}
$X$의 열린 부분집합 $U$가 구성가능이기 위한 필요충분조건은 $X$에서
역콤팩트인 것이다. $X$의 닫힌 부분집합 $F$가 구성가능이기 위한
필요충분조건은 열린집합 $\complement{F}$가 역콤팩트인 것이다.
\end{corollary}

\begin{env}[9.1.6]
\label{0.9.1.6-ko}
$X$의 모든 준콤팩트 열린 부분집합이 역콤팩트인 경우, 다시 말해 $X$의
두 준콤팩트 열린 부분집합의 교집합이 준콤팩트인 경우가 중요하다
((\hyperref[I.5.5.6-ko]{I, 5.5.6}) 참조). $X$ 자체가 준콤팩트이면,
이는 $X$에서 역콤팩트인 열린 부분집합들이 $X$의 준콤팩트 열린
부분집합들과 같고, $X$의 구성가능 부분집합들이 준콤팩트 열린집합
$U$, $V$에 대한 $U\cap\complement{V}$ 꼴의 집합들의 유한 합집합과
같다는 뜻이다.
\end{env}

\begin{corollary}[9.1.7]
\label{0.9.1.7-ko}
뇌터 공간 $X$의 부분집합이 구성가능이기 위한 필요충분조건은 $X$의
국소 닫힌 부분집합들의 유한 합집합인 것이다.
\end{corollary}

\begin{proposition}[9.1.8]
\label{0.9.1.8-ko}
$X$를 위상공간, $U$를 $X$의 열린 부분집합이라 하자.
\begin{enumerate}
  \item[(i)] $T$가 $X$의 구성가능 부분집합이면 $T\cap U$는 $U$의
    구성가능 부분집합이다.
  \item[(ii)] 더 나아가 $U$가 $X$에서 역콤팩트라고 하자. $U$의
    부분집합 $Z$가 $X$에서 구성가능이기 위한 필요충분조건은 $U$에서
    구성가능인 것이다.
\end{enumerate}
\end{proposition}

\begin{enumerate}
  \item[(i)] (\hyperref[0.9.1.3-ko]{9.1.3})을 사용하면, $T$가 $X$에서
    역콤팩트인 열린집합일 때 $T\cap U$가 $U$에서 역콤팩트인
    열린집합임을 보이는 것으로 환원된다. 다시 말해 모든 준콤팩트
    열린집합 $W\subset U$에 대하여 $T\cap U\cap W=T\cap W$가
    준콤팩트임을 보이면 되는데, 이는 가정에서 곧바로 따른다.
  \item[(ii)] (i)에 따라 조건은 필요하므로, 충분함만 보이면 된다.
    (\hyperref[0.9.1.3-ko]{9.1.3})을 고려하면 $Z$가 \emph{$U$에서}
    역콤팩트인 열린집합인 경우만 생각하면 충분하다. 실제로 그러면
    $U\setminus Z$가 $X$에서 구성가능이고, $Z$, $Z'$이 \emph{$U$에서}
    역콤팩트인 두 열린집합이면 $Z\cap(U\setminus Z')$가 참으로
    $X$에서 구성가능임이 따른다. 이제 $W$가 $X$의 준콤팩트
    열린집합이고 $Z$가 $U$에서 역콤팩트인 열린집합이면
    $Z\cap W=Z\cap(W\cap U)$이고, 가정에 따라 $W\cap U$는 $U$의
    준콤팩트 열린집합이다. 따라서 $W\cap Z$는 준콤팩트이고, 이에 따라
    $Z$는 $X$에서 역콤팩트인 열린집합이며 \emph{더구나} $X$에서
    구성가능이다.
\end{enumerate}

\begin{corollary}[9.1.9]
\label{0.9.1.9-ko}
$X$를 위상공간, $(U_i)_{i\in I}$를 $X$에서 역콤팩트인 열린집합들로
이루어진 $X$의 유한 덮개라 하자. $X$의 부분집합 $Z$가 $X$에서
구성가능이기 위한 필요충분조건은 모든 $i\in I$에 대하여 $Z\cap U_i$가
$U_i$에서 구성가능인 것이다.
\end{corollary}

\begin{env}[9.1.10]
\label{0.9.1.10-ko}
특히 $X$가 준콤팩트이고 $X$의 모든 점이
\oldpage[0\textsubscript{III}]{14}
$X$에서 역콤팩트인 열린
근방들(따라서 \emph{더구나} 준콤팩트 열린 근방들)로 이루어진 기본계를
갖는다고 하자. 그러면 $X$의 부분집합 $Z$가 $X$에서 구성가능이라는
조건은 \emph{국소적}이다. 다시 말해 모든 $x\in X$에 대하여 $x$의
열린 근방 $V$가 존재하여 $V\cap Z$가 $V$에서 구성가능인 것이
필요충분하다. 실제로 이 조건이 성립하면, 모든 $x\in X$에 대하여
$X$에서 \emph{역콤팩트}이고 $V\cap Z$가 $V$에서 구성가능인 열린 근방
$V$가 존재한다. 이는 $X$에 관한 가정과
(\hyperref[0.9.1.8-ko]{9.1.8, (i)})에서 따른다. 이제 이 근방들의 유한
개로 $X$를 덮고 (\hyperref[0.9.1.9-ko]{9.1.9})를 적용하면 충분하다.
\end{env}

\begin{definition}[9.1.11]
\label{0.9.1.11-ko}
$X$를 위상공간이라 하자. 모든 $x\in X$에 대하여 $x$의 열린 근방 $V$가
존재하여 $T\cap V$가 $V$에서 구성가능이면, $X$의 부분집합 $T$를
\emph{$X$에서 국소 구성가능}이라 한다.
\end{definition}

$V\cap T$가 $V$에서 구성가능이면 모든 열린집합 $W\subset V$에 대하여
$W\cap T$가 $W$에서 구성가능임이
(\hyperref[0.9.1.8-ko]{9.1.8, (i)})에서 곧바로 따른다. $T$가 $X$에서
국소 구성가능이면, 앞의 주석에 따라 $X$의 모든 열린집합 $U$에 대하여
$T\cap U$는 $U$에서 국소 구성가능이다. 같은 주석으로부터 $X$에서
국소 구성가능인 부분집합들의 집합은 유한 합집합과 유한 교집합에
관하여 닫혀 있음을 알 수 있다. 한편 이 집합이 여집합을 취하는 것에
관해서도 닫혀 있음은 명백하다.

\begin{proposition}[9.1.12]
\label{0.9.1.12-ko}
$X$를 위상공간이라 하자. $X$에서 구성가능인 모든 집합은 $X$에서
국소 구성가능이다. $X$가 준콤팩트이고 $X$의 위상이 $X$에서 역콤팩트인
집합들로 이루어진 기저를 가지면 그 역도 성립한다.
\end{proposition}

첫째 주장은 정의 (\hyperref[0.9.1.11-ko]{9.1.11})에서 따르고, 둘째
주장은 (\hyperref[0.9.1.10-ko]{9.1.10})에서 따른다.

\begin{corollary}[9.1.13]
\label{0.9.1.13-ko}
$X$를 그 위상이 $X$에서 역콤팩트인 집합들로 이루어진 기저를 갖는
위상공간이라 하자. $X$에서 국소 구성가능인 모든 부분집합 $T$는
$X$에서 역콤팩트이다.
\end{corollary}

실제로 $U$를 $X$의 준콤팩트 열린집합이라 하자. $T\cap U$는 $U$에서
국소 구성가능이므로 (\hyperref[0.9.1.12-ko]{9.1.12})에 따라 $U$에서
구성가능이고, 따라서 (\hyperref[0.9.1.4-ko]{9.1.4})에 따라
준콤팩트이다.

\subsection{뇌터 공간에서의 구성가능 집합}
\label{subsection:0.9.2-ko}

\begin{env}[9.2.1]
\label{0.9.2.1-ko}
뇌터 공간 $X$에서 $X$의 구성가능 부분집합들은 \emph{$X$의 국소 닫힌
부분집합들의 유한 합집합}임을
(\hyperref[0.9.1.7-ko]{9.1.7})에서 보았다.

뇌터 공간 $X'$에서 $X$로 가는 연속사상에 의한 $X$의 구성가능 집합의
역상은 $X'$에서 구성가능이다. $Y$가 뇌터 공간 $X$의 구성가능
부분집합이면, $Y$의 부분공간으로서 구성가능인 $Y$의 부분집합들과
$X$의 부분공간으로서 구성가능인 것들은 서로 같다.
\end{env}

\begin{proposition}[9.2.2]
\label{0.9.2.2-ko}
$X$를 뇌터 기약공간, $E$를 $X$의 구성가능 부분집합이라 하자. $E$가
$X$에서 모든 곳에서 조밀하기 위한 필요충분조건은 $E$가 $X$의 공집합이
아닌 열린 부분집합 하나를 포함하는 것이다.
\end{proposition}

\oldpage[0\textsubscript{III}]{15}
공집합이 아닌 모든 열린집합은 $X$에서 조밀하므로 이 조건이 충분함은
명백하다. 역으로 $E=\bigcup_{i=1}^n(U_i\cap F_i)$를 $X$의 구성가능
부분집합이라 하자. 여기서 $U_i$들은 공집합이 아닌 열린집합이고
$F_i$들은 $X$의 닫힌집합이다. 그러면
$\overline{E}\subset\bigcup_i F_i$이다. 따라서 $\overline{E}=X$이면
$X$는 $F_i$ 중 하나와 같고, 이에 따라 $E\supset U_i$이다. 이로써
증명이 끝난다.

$X$가 일반점 $x$를 가지면
(\hyperref[0.2.1.2-ko]{$0_{\mathrm I}$, 2.1.2}),
(\hyperref[0.9.2.2-ko]{9.2.2})의 조건은 관계 $x\in E$와 동치이다.

\begin{proposition}[9.2.3]
\label{0.9.2.3-ko}
$X$를 뇌터 공간이라 하자. $X$의 부분집합 $E$가 구성가능이기 위한
필요충분조건은 $X$의 모든 닫힌 기약 부분집합 $Y$에 대하여 $E\cap Y$가
$Y$에서 어디에서도 조밀하지 않거나 $Y$의 공집합이 아닌 열린
부분집합 하나를 포함하는 것이다.
\end{proposition}

조건의 필요성은 $E\cap Y$가 $Y$의 구성가능 부분집합이어야 한다는
사실과 (\hyperref[0.9.2.2-ko]{9.2.2})에서 따른다. 실제로 $Y$에서 조밀하지
않은 부분집합은 기약공간 $Y$에서 반드시 어디에서도 조밀하지 않다
(\hyperref[0.2.1.1-ko]{$0_{\mathrm I}$, 2.1.1}). 조건이 충분함을 보이기
위하여, $Y\cap E$가 구성가능인(여기서 $Y$에 관하여 구성가능이든
$X$에 관하여 구성가능이든 같은 뜻이다) $X$의 닫힌 부분집합 $Y$들의
집합 $\mathfrak{F}$에 뇌터 귀납법의 원리
(\hyperref[0.2.2.2-ko]{$0_{\mathrm I}$, 2.2.2})를 적용하자. 따라서
$X$의 모든 닫힌 부분집합 $Y\ne X$에 대하여 $E\cap Y$가 구성가능이라고
가정할 수 있다. 먼저 $X$가 기약이 아니라고 하고, 그 기약 성분들을
$X_i$ ($1\leq i\leq m$)라 하자. 이들은 반드시 유한 개이다
(\hyperref[0.2.2.5-ko]{$0_{\mathrm I}$, 2.2.5}). 가정에 따라
$E\cap X_i$들은 구성가능이고, 따라서 그 합집합 $E$도 구성가능이다.
이제 $X$가 기약이라고 하자. 이때 가정에 따라 $E$가 어디에서도
조밀하지 않아서 $\overline{E}\ne X$이고
$E=E\cap\overline{E}$가 구성가능이거나, 또는 $E$가 $X$의 공집합이
아닌 열린집합 $U$를 포함하여 $U$와 $E\cap(X\setminus U)$의
합집합이다. 그런데 $X\setminus U$는 $X$와 다른 닫힌집합이므로
$E\cap(X\setminus U)$는 구성가능이다. 따라서 $E$ 자체도 구성가능이고,
이로써 증명이 끝난다.

\begin{corollary}[9.2.4]
\label{0.9.2.4-ko}
$X$를 뇌터 공간, $(E_\alpha)$를 다음 조건을 만족하는 $X$의 구성가능
부분집합들의 증가 여과족이라 하자.
\begin{enumerate}
  \item[$1^\circ$] $X$는 족 $(E_\alpha)$의 합집합이다.
  \item[$2^\circ$] $X$의 모든 닫힌 기약 부분집합은 어떤 $E_\alpha$의
    폐포에 포함된다.
\end{enumerate}
그러면 $X=E_\alpha$인 지표 $\alpha$가 존재한다.

$X$의 모든 닫힌 기약 부분집합이 일반점을 가지면 가정 $2^\circ$를
생략할 수 있다.
\end{corollary}

$X$의 닫힌 부분집합들 가운데 적어도 하나의 $E_\alpha$에 포함되는
것들의 집합 $\mathfrak{M}$에 뇌터 귀납법의 원리
(\hyperref[0.2.2.2-ko]{$0_{\mathrm I}$, 2.2.2})를 적용하자. 따라서
$X$의 모든 닫힌 부분집합 $Y\ne X$가 어떤 $E_\alpha$에 포함된다고
가정할 수 있다. $X$가 기약이 아니면 명제는 명백하다. 실제로 $X$의
각 기약 성분 $X_i$ ($1\leq i\leq m$)는 어떤 $E_{\alpha_i}$에 포함되고,
모든 $E_{\alpha_i}$를 포함하는 $E_\alpha$가 존재한다. 따라서 $X$가
기약인 경우를 생각하자. 가정에 따라 $X=\overline{E_\beta}$인 $\beta$가
존재하고, 따라서 (\hyperref[0.9.2.2-ko]{9.2.2})에 따라 $E_\beta$는
$X$의 공집합이 아닌 열린집합 $U$를 포함한다. 그런데 닫힌집합
$X\setminus U$는 어떤 $E_\gamma$에 포함되므로, $E_\beta$와
$E_\gamma$를 포함하는 $E_\alpha$를 취하면 충분하다. $X$의 모든 닫힌
기약 부분집합 $Y$가
\oldpage[0\textsubscript{III}]{16}
일반점 $y$를 가지면 $y\in E_\alpha$인 $\alpha$가 존재한다. 따라서
$Y=\overline{\{y\}}\subset\overline{E_\alpha}$이고, 조건 $2^\circ$는
$1^\circ$의 귀결이다.

\begin{proposition}[9.2.5]
\label{0.9.2.5-ko}
$X$를 뇌터 공간, $x$를 $X$의 점, $E$를 $X$의 구성가능 부분집합이라
하자. $E$가 $x$의 근방이기 위한 필요충분조건은 $x$를 포함하는 $X$의
모든 닫힌 기약 부분집합 $Y$에 대하여 $E\cap Y$가 $Y$에서 조밀한
것이다($Y$에 일반점 $y$가 존재하면, 이는
(\hyperref[0.9.2.2-ko]{9.2.2})에 따라 $y\in E$라는 뜻이기도 하다).
\end{proposition}

조건이 필요함은 명백하다. 충분함을 보이자. $x$를 포함하며 $E\cap Y$가
$Y$에서 $x$의 근방인 $X$의 닫힌 부분집합 $Y$들의 집합
$\mathfrak{M}$에 뇌터 귀납법의 원리를 적용하면, $x$를 포함하는
$X$의 모든 닫힌 부분집합 $Y\ne X$가 $\mathfrak{M}$에 속한다고 가정할
수 있다. $X$가 기약이 아니면 $x$를 포함하는 $X$의 각 기약 성분
$X_i$는 $X$와 다르므로 $E\cap X_i$는 $X_i$에서 $x$의 근방이다.
따라서 $E$는 $x$를 포함하는 $X$의 기약 성분들의 합집합에서 $x$의
근방이고, 이 합집합 자체가 $X$에서 $x$의 근방이므로 $E$도 그러하다.
$X$가 기약이면 가정에 따라 $E$는 $X$에서 조밀하고, 따라서
(\hyperref[0.9.2.2-ko]{9.2.2})에 따라 $X$의 공집합이 아닌 열린
부분집합 $U$를 포함한다. $x\in U$이면 명제는 명백하다. 그렇지 않으면
가정에 따라 $x$는 $X\setminus U$에 관한 $E\cap(X\setminus U)$의
내점이다. 따라서 $X$에서 $X\setminus E$의 폐포는 $x$를 포함하지 않고,
이 폐포의 여집합은 $E$에 포함되는 $x$의 근방이다. 이로써 증명이
끝난다.

\begin{corollary}[9.2.6]
\label{0.9.2.6-ko}
$X$를 뇌터 공간, $E$를 $X$의 부분집합이라 하자. $E$가 $X$의
열린집합이기 위한 필요충분조건은 $E$와 만나는 $X$의 모든 닫힌 기약
부분집합 $Y$에 대하여 $E\cap Y$가 $Y$의 공집합이 아닌 열린 부분집합
하나를 포함하는 것이다.
\end{corollary}

조건이 필요함은 명백하다. 역으로 조건이 성립하면
(\hyperref[0.9.2.3-ko]{9.2.3})에 따라 $E$는 구성가능이다. 더 나아가
(\hyperref[0.9.2.5-ko]{9.2.5})에 따라 $E$는 자신의 모든 점의 근방이고,
따라서 결론이 따른다.

\subsection{구성가능 함수}
\label{subsection:0.9.3-ko}

\begin{definition}[9.3.1]
\label{0.9.3.1-ko}
$h$를 위상공간 $X$에서 집합 $T$로 가는 사상이라 하자. 모든 $t\in T$에
대하여 $h^{-1}(t)$가 구성가능이고 유한 개의 $t$를 제외하면 공집합이면,
$h$를 \emph{구성가능}이라 한다. 이때 $T$의 모든 부분집합 $S$에 대하여
$h^{-1}(S)$는 구성가능이다. 모든 $x\in X$가 $h|V$가 구성가능인 열린
근방 $V$를 가지면 $h$를 \emph{국소 구성가능}이라 한다.
\end{definition}

모든 구성가능 함수는 국소 구성가능이다. $X$가 준콤팩트이고 $X$에서
역콤팩트인 열린집합들로 이루어진 기저를 가지면(특히 $X$가 뇌터이면)
그 역도 성립한다.

\begin{proposition}[9.3.2]
\label{0.9.3.2-ko}
$h$를 뇌터 공간 $X$에서 집합 $T$로 가는 사상이라 하자. $h$가
구성가능이기 위한 필요충분조건은 $X$의 모든 닫힌 기약 부분집합 $Y$에
대하여 $Y$에서 상대적으로 열려 있는 공집합이 아닌 부분집합 $U$가
존재하여 $U$에서 $h$가 상수인 것이다.
\end{proposition}

조건은 필요하다. 실제로 가정에 따라 $h$가 $Y$에서 취하는 값들은
유한 개의 $t_i$뿐이고, 각 집합 $h^{-1}(t_i)\cap Y$는 $Y$에서
구성가능이다(\hyperref[0.9.2.1-ko]{9.2.1}). 이 집합들이 모두 $Y$에서
어디에서도 조밀하지 않을 수는 없으므로, 그중 적어도 하나는 공집합이
아닌 열린집합을 포함한다
(\hyperref[0.9.2.3-ko]{9.2.3}).

\oldpage[0\textsubscript{III}]{17}
조건이 충분함을 보이기 위하여, $h|Y$가 구성가능인 $X$의 닫힌 부분집합
$Y$들의 집합 $\mathfrak{M}$에 뇌터 귀납법의 원리를 적용하자. 따라서
$X$의 모든 닫힌 부분집합 $Y\ne X$에 대하여 $h|Y$가 구성가능이라고
가정할 수 있다. $X$가 기약이 아니면 $X$의 각 기약 성분 $X_i$로 제한한
$h$는 구성가능이다. 이 성분들은 유한 개이므로 정의
(\hyperref[0.9.3.1-ko]{9.3.1})에서 $h$가 구성가능임이 곧바로 따른다.
$X$가 기약이면 가정에 따라 $h$가 상수인 $X$의 공집합이 아닌 열린
부분집합 $U$가 존재한다. 한편 가정에 따라 $X\setminus U$로 제한한
$h$는 구성가능이므로, 이로부터 $h$가 구성가능임이 곧바로 따른다.

\begin{corollary}[9.3.3]
\label{0.9.3.3-ko}
$X$를 모든 닫힌 기약 부분집합이 일반점을 갖는 뇌터 공간이라 하자.
$h$가 $X$에서 집합 $T$로 가는 사상이고 모든 $t\in T$에 대하여
$h^{-1}(t)$가 구성가능이면, $h$는 구성가능이다.
\end{corollary}

실제로 $Y$가 $X$의 닫힌 기약 부분집합이고 $y$가 그 일반점이면,
$Y\cap h^{-1}(h(y))$는 구성가능이고 $y$를 포함한다. 따라서
(\hyperref[0.9.2.2-ko]{9.2.2})에 따라 이 집합은 $Y$의 공집합이 아닌
열린 부분집합 하나를 포함하며,
(\hyperref[0.9.3.2-ko]{9.3.2})를 적용하면 충분하다.

\begin{proposition}[9.3.4]
\label{0.9.3.4-ko}
$X$를 모든 닫힌 기약 부분집합이 일반점을 갖는 뇌터 공간, $h$를
$X$에서 순서집합으로 가는 구성가능 사상이라 하자. $h$가 $X$에서
상반연속이기 위한 필요충분조건은 모든 $x\in X$와 $x$의 모든 일반화
$x'$에 대하여
(\hyperref[0.2.1.2-ko]{$0_{\mathrm I}$, 2.1.2})
$h(x')\leq h(x)$인 것이다.
\end{proposition}

함수 $h$는 유한 개의 값만 취한다. 따라서 $h$가 상반연속이라는 것은
모든 $x\in X$에 대하여 $h(y)\leq h(x)$인 $y\in X$들의 집합 $E$가
$x$의 근방이라는 뜻이다. 가정에 따라 $E$는 $X$의 구성가능
부분집합이다. 한편 $X$의 닫힌 기약 부분집합 $Y$가 $x$를 포함한다는
것은 그 일반점 $y$가 $x$의 일반화라는 뜻이다. 따라서 결론은
(\hyperref[0.9.2.5-ko]{9.2.5})에서 따른다.

\section{평탄 가군에 관한 보충}
\label{section:0.10-fr}

(\hyperref[subsection:0.10.1-fr]{10.1})과
(\hyperref[subsection:0.10.2-fr]{10.2})에서 증명 없이 서술한 성질들의
증명은 부르바키, \emph{Alg. comm.}, 제II장과 제III장을 참조한다.

\subsection{평탄 가군과 자유 가군 사이의 관계}
\label{subsection:0.10.1-fr}

\begin{env}[10.1.1]
\label{0.10.1.1-fr}
$A$를 환, $\mathfrak{J}$를 $A$의 아이디얼, $M$을 $A$-가군이라 하자. 모든
정수 $p\geq0$에 대하여 다음과 같은 표준 $(A/\mathfrak{J})$-가군
준동형이 있다.
\[
\label{0.10.1.1.1-fr}
  \varphi_p:(M/\mathfrak{J}M)\otimes_{A/\mathfrak{J}}
  (\mathfrak{J}^p/\mathfrak{J}^{p+1})
  \longrightarrow \mathfrak{J}^pM/\mathfrak{J}^{p+1}M
  \tag{10.1.1.1}
\]
이 준동형은 명백히 \emph{전사}이다. $\mathfrak{J}^p$들로 여과된 $A$에
연관된 등급환을
$\operatorname{gr}(A)=\bigoplus_{p\geq0}\mathfrak{J}^p/\mathfrak{J}^{p+1}$로,
$\mathfrak{J}^pM$들로 여과된 $M$에 연관된 등급 $\operatorname{gr}(A)$-가군을
$\operatorname{gr}(M)=\bigoplus_{p\geq0}\mathfrak{J}^pM/\mathfrak{J}^{p+1}M$로
나타내자. 그러면
$\operatorname{gr}_p(A)=\mathfrak{J}^p/\mathfrak{J}^{p+1}$,
$\operatorname{gr}_p(M)=\mathfrak{J}^pM/\mathfrak{J}^{p+1}M$이고,
$\varphi_p$들은 다음의 \emph{전사}인 등급 $\operatorname{gr}(A)$-가군
준동형을 정의한다.
\[
\label{0.10.1.1.2-fr}
  \varphi:\operatorname{gr}_0(M)\otimes_{\operatorname{gr}_0(A)}
  \operatorname{gr}(A)\longrightarrow\operatorname{gr}(M).
  \tag{10.1.1.2}
\]
\end{env}

\oldpage[0\textsubscript{III}]{18}
\begin{env}[10.1.2]
\label{0.10.1.2-fr}
다음 가정 가운데 하나가 성립한다고 하자.
\begin{enumerate}
  \item[(i)] $\mathfrak{J}$는 멱영이다.
  \item[(ii)] $A$는 뇌터이고, $\mathfrak{J}$는 $A$의 근기에 포함되며,
    $M$은 유한 생성이다.
\end{enumerate}
그러면 다음 성질들은 서로 동치이다.
\begin{enumerate}
  \item[$a)$] $M$은 자유 $A$-가군이다.
  \item[$b)$] $M/\mathfrak{J}M=M\otimes_A(A/\mathfrak{J})$는 자유
    $(A/\mathfrak{J})$-가군이고,
    $\operatorname{Tor}_1^A(M,A/\mathfrak{J})=0$이다.
  \item[$c)$] $M/\mathfrak{J}M$은 자유 $(A/\mathfrak{J})$-가군이고
    표준 준동형 (\hyperref[0.10.1.1.2-fr]{10.1.1.2})는 단사이다
    (따라서 전단사이다).
\end{enumerate}
\end{env}

\begin{env}[10.1.3]
\label{0.10.1.3-fr}
$A/\mathfrak{J}$가 \emph{체}라고(다시 말해 $\mathfrak{J}$가 극대
아이디얼이라고) 하고, (\hyperref[0.10.1.2-fr]{10.1.2})의 가정 (i),
(ii) 가운데 하나가 성립한다고 하자. 그러면 다음 성질들은 서로 동치이다.
\begin{enumerate}
  \item[$a)$] $M$은 자유 $A$-가군이다.
  \item[$b)$] $M$은 사영 $A$-가군이다.
  \item[$c)$] $M$은 평탄 $A$-가군이다.
  \item[$d)$] $\operatorname{Tor}_1^A(M,A/\mathfrak{J})=0$이다.
  \item[$e)$] 표준 준동형
    (\hyperref[0.10.1.1.2-fr]{10.1.1.2})는 전단사이다.
\end{enumerate}
이 결과는 특히 다음 두 경우에 적용된다.
\begin{enumerate}
  \item[(i)] 극대 아이디얼 $\mathfrak{J}$가 \emph{멱영}인 국소환 $A$ 위의
    \emph{임의의} 가군 $M$(예를 들어 $A$가 아르틴 국소환인 경우).
  \item[(ii)] \emph{뇌터} 국소환 위의 \emph{유한 생성} 가군 $M$.
\end{enumerate}
\end{env}

\subsection{평탄성의 국소 판정법}
\label{subsection:0.10.2-fr}

\begin{env}[10.2.1]
\label{0.10.2.1-fr}
가정과 표기는 (\hyperref[0.10.1.1-fr]{10.1.1})에서와 같다고 하고, 다음
조건들을 생각하자.
\begin{enumerate}
  \item[$a)$] $M$은 평탄 $A$-가군이다.
  \item[$b)$] $M/\mathfrak{J}M$은 평탄 $(A/\mathfrak{J})$-가군이고
    $\operatorname{Tor}_1^A(M,A/\mathfrak{J})=0$이다.
  \item[$c)$] $M/\mathfrak{J}M$은 평탄 $(A/\mathfrak{J})$-가군이고
    표준 준동형 (\hyperref[0.10.1.1.2-fr]{10.1.1.2})는 전단사이다.
  \item[$d)$] 모든 $n\geq1$에 대하여 $M/\mathfrak{J}^nM$은 평탄
    $(A/\mathfrak{J}^n)$-가군이다.
\end{enumerate}
그러면 다음 함의들이 성립한다.
\[
  a\Rightarrow b\Rightarrow c\Rightarrow d
\]
또한 $\mathfrak{J}$가 \emph{멱영}이면 네 조건 $a)$, $b)$, $c)$, $d)$는
\emph{서로 동치}이다. $A$가 뇌터이고, 더 나아가 $M$이
\emph{아이디얼적으로 분리}되어 있을 때에도 마찬가지이다. 여기서
아이디얼적으로 분리되어 있다는 것은 $A$의 모든 아이디얼
$\mathfrak{a}$에 대하여 $A$-가군 $\mathfrak{a}\otimes_A M$이
$\mathfrak{J}$-준아딕 위상에 관해 \emph{분리}되어 있다는 뜻이다.
\end{env}

\begin{env}[10.2.2]
\label{0.10.2.2-fr}
$A$를 뇌터환, $B$를 가환 뇌터 $A$-대수, $\mathfrak{J}$를
$\mathfrak{J}B$가 $B$의 근기에 포함되도록 하는 $A$의 아이디얼,
$M$을 유한 생성 $B$-가군이라 하자. 그러면 $M$을 $A$-가군으로 볼 때
(\hyperref[0.10.2.1-fr]{10.2.1})의 조건 $a)$, $b)$, $c)$, $d)$는 서로
동치이다.

이 결과는 특히 $A$와 $B$가 뇌터 \emph{국소환}이고 준동형 $A\to B$가
\emph{국소 준동형}일 때 적용된다. 더 구체적으로 이때
$\mathfrak{J}$가 $A$의 \emph{극대} 아이디얼이면, 조건 $b)$와 $c)$에서
$M/\mathfrak{J}M$이 평탄이라는 자동으로 성립하는 가정을 생략할 수
있고, 조건 $d)$는 가군 $M/\mathfrak{J}^nM$이
$A/\mathfrak{J}^n$ 위에서 \emph{자유}라는 뜻이다.
\end{env}

\oldpage[0\textsubscript{III}]{19}
\begin{env}[10.2.3]
\label{0.10.2.3-fr}
$A$, $B$, $\mathfrak{J}$, $M$에 대한 가정은
(\hyperref[0.10.2.2-fr]{10.2.2})의 처음에 서술한 것과 같다고 하자.
$\widehat A$를 $A$의 $\mathfrak{J}$-준아딕 위상에 대한 분리 완비화,
$\widehat M$을 $M$의 $\mathfrak{J}B$-준아딕 위상에 대한 분리
완비화라 하자. 그러면 $M$이 평탄 $A$-가군이기 위한 필요충분조건은
$\widehat M$이 평탄 $\widehat A$-가군인 것이다.
\end{env}

\begin{env}[10.2.4]
\label{0.10.2.4-fr}
$\rho:A\to B$를 뇌터 국소환들의 국소 준동형, $k$를 $A$의 잉여체,
$M$, $N$을 두 유한 생성 $B$-가군이라 하고, $N$은 $A$-\emph{평탄}이라고
가정하자. $u:M\to N$을 $B$-준동형이라 하자. 그러면 다음 조건들은 서로
동치이다.
\begin{enumerate}
  \item[$a)$] $u$는 단사이고 $\operatorname{Coker}(u)$는 평탄
    $A$-가군이다.
  \item[$b)$] $u\otimes1:M\otimes_A k\to N\otimes_A k$는 단사이다.
\end{enumerate}
\end{env}

\begin{env}[10.2.5]
\label{0.10.2.5-fr}
$\rho:A\to B$, $\sigma:B\to C$를 뇌터 국소환들의 국소 준동형,
$k$를 $A$의 잉여체, $M$을 유한 생성 $C$-가군이라 하자. $B$가 평탄
$A$-가군이라고 가정한다. 그러면 다음 조건들은 서로 동치이다.
\begin{enumerate}
  \item[$a)$] $M$은 평탄 $B$-가군이다.
  \item[$b)$] $M$은 평탄 $A$-가군이고, $M\otimes_A k$는 평탄
    $(B\otimes_A k)$-가군이다.
\end{enumerate}
\end{env}

\begin{proposition}[10.2.6]
\label{0.10.2.6-fr}
$A$, $B$를 두 뇌터 국소환, $\rho:A\to B$를 국소 준동형,
$\mathfrak{J}$를 극대 아이디얼에 포함되는 $B$의 아이디얼, $M$을
유한 생성 $B$-가군이라 하자. 모든 $n\geq0$에 대하여
$M_n=M/\mathfrak{J}^{n+1}M$이 평탄 $A$-가군이라고 가정한다. 그러면
$M$은 평탄 $A$-가군이다.
\end{proposition}

유한 생성 $A$-가군들의 모든 단사 준동형 $u:N'\to N$에 대하여
$v=1\otimes u:M\otimes_A N'\to M\otimes_A N$이 단사임을 증명해야 한다.
그런데 $M\otimes_A N'$과 $M\otimes_A N$은 유한 생성 $B$-가군이므로
$\mathfrak{J}$-준아딕 위상에 관해 분리되어 있다
(\hyperref[0.7.3.5-ko]{$0_{\mathrm I}$, 7.3.5}). 따라서 분리 완비화들
사이의 준동형
$\widehat v:\widehat{(M\otimes_A N')}\to\widehat{(M\otimes_A N)}$이
단사임을 증명하면 충분하다. 한편 $\widehat v=\varprojlim v_n$이고,
여기서 $v_n$은 준동형
$1\otimes u:M_n\otimes_A N'\to M_n\otimes_A N$이다. 가정에 따라
$M_n$은 $A$-평탄이므로 모든 $n$에 대해 $v_n$은 단사이고, 따라서
$\widehat v$도 단사이다. 실제로 함자 $\varprojlim$은 왼쪽 완전이다.

\begin{corollary}[10.2.7]
\label{0.10.2.7-fr}
$A$를 뇌터환, $B$를 뇌터 국소환, $\rho:A\to B$를 준동형,
$f$를 $B$의 극대 아이디얼의 원소, $M$을 유한 생성 $B$-가군이라 하자.
$M$의 곱셈 준동형 $f_M:x\mapsto fx$가 단사이고 $M/fM$이 평탄
$A$-가군이라고 가정한다. 그러면 $M$은 평탄 $A$-가군이다.
\end{corollary}

$i\geq0$에 대하여 $M_i=f^iM$이라 두자. $f_M$이 단사이므로
$M_i/M_{i+1}$은 $M/fM$과 동형이고, 따라서 모든 $i\geq0$에 대하여
$A$-평탄이다. 완전열
\[
  0\longrightarrow M_i/M_{i+1}\longrightarrow M/M_{i+1}
  \longrightarrow M/M_i\longrightarrow0
\]
에서 $i$에 대한 귀납법으로 모든 $i\geq0$에 대해 $M/M_i$가
$A$-\emph{평탄}임을 얻는다
(\hyperref[0.6.1.2-ko]{$0_{\mathrm I}$, 6.1.2}). 따라서
(\hyperref[0.10.2.6-fr]{10.2.6})을 적용할 수 있다. 다음과 같이 직접
논증할 수도 있다. 모든 유한 생성 $A$-가군 $N$에 대하여
$M\otimes_A N$은 유한 생성 $B$-가군이다. $f$가 $B$의 근기
$\mathfrak n$에 속하므로 $M\otimes_A N$ 위의 $(f)$-아딕 위상은
\oldpage[0\textsubscript{III}]{20}
$\mathfrak n$-아딕 위상보다 더 섬세하며, 후자는 \emph{분리}되어 있음을
알고 있다 (\hyperref[0.7.3.5-ko]{$0_{\mathrm I}$, 7.3.5}). 또한
$M/M_i$가 $A$-평탄이므로
\[
  f^i(M\otimes_A N)=\operatorname{Im}(M_i\otimes_A N\to M\otimes_A N)
  =\operatorname{Ker}(M\otimes_A N\to(M/M_i)\otimes_A N)
\]
이다 (\hyperref[0.6.1.2-ko]{$0_{\mathrm I}$, 6.1.2}). 이제 $N$을 유한
생성 $A$-가군, $N'$을 $N$의 부분가군, $j:N'\to N$을 표준 포함
준동형이라 하자. 가환 도식
\[
\xymatrix{
  M\otimes_A N'\ar[r]\ar[d]_{1_M\otimes j} &
  (M/M_i)\otimes_A N'\ar[d]^{1_{M/M_i}\otimes j}\\
  M\otimes_A N\ar[r] & (M/M_i)\otimes_A N
}
\]
에서 $M/M_i$가 $A$-평탄이므로 $1_{M/M_i}\otimes j$는 단사이다.
따라서 모든 $i$에 대하여
\[
  \operatorname{Ker}(M\otimes_A N'\to M\otimes_A N)
  \subset\operatorname{Ker}(M\otimes_A N'\to(M/M_i)\otimes_A N')
\]
이다. 위에서 본 대로 우변들의 교집합은 $0$만으로 이루어져 있으므로
좌변도 마찬가지이고, 따라서 $M$은 $A$-평탄이다.

\begin{proposition}[10.2.8]
\label{0.10.2.8-fr}
$A$를 축소 뇌터환, $M$을 유한 생성 $A$-가군이라 하자. 이산
값매김환인 모든 $A$-대수 $B$에 대하여 $M\otimes_A B$가 평탄
$B$-가군이라고(따라서 자유 가군이라고
(\hyperref[0.10.1.3-fr]{10.1.3})) 가정한다. 그러면 $M$은 평탄
$A$-가군이다.
\end{proposition}

$M$이 평탄이기 위한 필요충분조건은 $A$의 모든 극대 아이디얼
$\mathfrak m$에 대하여 $M_{\mathfrak m}$이 평탄
$A_{\mathfrak m}$-가군인 것임을 알고 있다
(\hyperref[0.6.3.3-ko]{$0_{\mathrm I}$, 6.3.3}). 따라서 $A$가
\emph{국소}인 경우로 한정해도 된다
(\hyperref[0.1.2.8-ko]{$0_{\mathrm I}$, 1.2.8}). 이제
$\mathfrak m$을 $A$의 극대 아이디얼, $\mathfrak p_i$
$(1\leq i\leq r)$를 그 극소 소아이디얼들, $k$를 잉여체
$A/\mathfrak m$이라 하자. 각 $i$에 대하여 정역
$A/\mathfrak p_i$와 같은 분수체 $K_i$를 가지며 이 정역을 지배하는
이산 값매김환 $B_i$가 존재함을 알고 있다
(\hyperref[II.7.1.7-ko]{II, 7.1.7}). $M_i=M\otimes_A B_i$라 두자.
가정에 따라 $M_i$는 $B_i$ 위에서 자유이므로, $k_i$를 $B_i$의
잉여체라 하면
\[
\label{0.10.2.8.1-fr}
  \operatorname{rg}_{k_i}(M_i\otimes_{B_i}k_i)
  =\operatorname{rg}_{K_i}(M_i\otimes_{B_i}K_i).
  \tag{10.2.8.1}
\]
이다. 그런데 합성 준동형 $A\to A/\mathfrak p_i\to B_i$는 명백히
국소 준동형이므로 $k_i$는 $k$의 확대체이고,
$M_i\otimes_{B_i}k_i=M\otimes_A k_i=(M\otimes_A k)\otimes_k k_i$이며,
한편 $M_i\otimes_{B_i}K_i=M\otimes_A K_i$이다. 따라서 등식
(\hyperref[0.10.2.8.1-fr]{10.2.8.1})로부터
\[
  \operatorname{rg}_k(M\otimes_A k)=\operatorname{rg}_{K_i}(M\otimes_A K_i)
  \qquad(1\leq i\leq r\text{에 대하여})
\]
를 얻는다. $A$는 축소환이므로 이 조건으로부터 $M$이 자유
$A$-가군임을 알고 있다(부르바키, \emph{Alg. comm.}, 제II장,
\S~3, n\textsuperscript{o}~2, 명제 7).

\subsection{국소환의 평탄 확대의 존재}
\label{subsection:0.10.3-fr}

\begin{proposition}[10.3.1]
\label{0.10.3.1-fr}
$A$를 뇌터 국소환, $\mathfrak J$를 그 극대 아이디얼,
$k=A/\mathfrak J$를 그 잉여체라 하자. $K$를 $k$의 확대체라 하자.
$B/\mathfrak J B$가 $K$와 $k$-동형이고 $B$가 평탄 $A$-가군이 되도록
하는, $A$에서 뇌터 국소환 $B$로의 국소 준동형이 존재한다.
\end{proposition}

이 명제를 여러 단계로 증명하겠다.

\begin{env}[10.3.1.1]
\label{0.10.3.1.1-fr}
먼저 $T$가 부정원이고 $K=k(T)$라고 하자. 다항식환 $A'=A[T]$에서
계수들이 아이디얼 $\mathfrak J$에 속하는 다항식들로 이루어진
소아이디얼 $\mathfrak J'=\mathfrak J A'$을 생각하자.
\oldpage[0\textsubscript{III}]{21}
$A'/\mathfrak J'$은 $k[T]$와 표준적으로 동형임이 명백하다. 분수환
$B=A'_{\mathfrak J'}$이 조건을 충족함을 보이자. 이는 극대 아이디얼이
$\mathfrak L=\mathfrak J B$인 뇌터 국소환임이 명백하다. 또한
$B/\mathfrak L=(A'/\mathfrak J')_{\mathfrak J'}=(k[T])_{\mathfrak J'}$은
$k[T]$의 분수체 $K$에 다름 아니다. 마지막으로 $B$는 평탄
$A'$-가군이고 $A'$은 자유 $A$-가군이므로, $B$는 평탄 $A$-가군이다
(\hyperref[0.6.2.1-ko]{$0_{\mathrm I}$, 6.2.1}).
\end{env}

\begin{env}[10.3.1.2]
\label{0.10.3.1.2-fr}
다음으로 $t$가 $k$ 위에서 대수적이고 $K=k(t)=k[t]$라고 하자.
$f\in k[T]$를 $t$의 최소다항식이라 하자. $k[T]$에서의 표준 상이
$f$인 일계수 다항식 $F\in A[T]$가 존재한다. 다시 $A'=A[T]$라 두고,
$\mathfrak J'$을 $A'$의 아이디얼 $\mathfrak J A'+(F)$라 하자. 이번에는
몫환 $B=A'/(F)$가 조건을 충족함을 보이겠다. 우선 $B$가 자유
$A$-가군이고, 따라서 평탄임이 명백하다. 환 $A'/\mathfrak J'$은
\[
  (A'/\mathfrak J A')/((\mathfrak J A'+(F))/\mathfrak J A')
  =k[T]/(f)=K
\]
와 동형이다. 따라서 $\mathfrak J'$의 $B$에서의 상 $\mathfrak L$은
극대이고, 명백히 $\mathfrak L=\mathfrak J B$이다. 마지막으로 $B$는
유한 생성 $A$-가군이므로 반국소환이다(부르바키,
\emph{Alg. comm.}, 제IV장, \S~2, n\textsuperscript{o}~5, 명제 9의
따름정리 3). 그리고 $B$의 극대 아이디얼들은 $B/\mathfrak J B$의
극대 아이디얼들과 일대일 대응한다([13], 제I권, 259쪽). 따라서 앞의
결과로부터 $B$가 국소환임을 알 수 있다.
\end{env}

\begin{lemma}[10.3.1.3]
\label{0.10.3.1.3-fr}
$(A_\lambda,f_{\mu\lambda})$를 국소환들의 여과 귀납계라 하고,
$f_{\mu\lambda}$들은 국소 준동형이라고 하자. $\mathfrak m_\lambda$를
$A_\lambda$의 극대 아이디얼, $K_\lambda=A_\lambda/\mathfrak m_\lambda$라
하자. 그러면 $A'=\varinjlim A_\lambda$는 국소환이고,
$\mathfrak m'=\varinjlim\mathfrak m_\lambda$는 그 극대 아이디얼이며,
$K=\varinjlim K_\lambda$는 그 잉여체이다. 또한 $\lambda<\mu$일 때
$\mathfrak m_\mu=\mathfrak m_\lambda A_\mu$이면, 모든 $\lambda$에 대해
$\mathfrak m'=\mathfrak m_\lambda A'$이다. 더 나아가
$\lambda<\mu$일 때 $A_\mu$가 평탄 $A_\lambda$-가군이고 모든
$A_\lambda$가 뇌터이면, $A'$은 뇌터이며 모든 $\lambda$에 대하여
평탄 $A_\lambda$-가군이다.
\end{lemma}

가정에 따라 $\lambda<\mu$이면
$f_{\mu\lambda}(\mathfrak m_\lambda)\subset\mathfrak m_\mu$이므로,
$\mathfrak m_\lambda$들은 귀납계를 이루고 그 극한 $\mathfrak m'$은
명백히 $A'$의 아이디얼이다. 또한 $x'\notin\mathfrak m'$이면 어떤
$\lambda$와 $x_\lambda\in A_\lambda$에 대하여
$x'=f_\lambda(x_\lambda)$이다. 여기서
$f_\lambda:A_\lambda\to A'$은 표준 준동형이다. $x'$가
$\mathfrak m'$에 속하지 않으므로 필연적으로
$x_\lambda\notin\mathfrak m_\lambda$이고, 따라서 $x_\lambda$는
$A_\lambda$에서 역원 $y_\lambda$를 가진다. 그러면
$y'=f_\lambda(y_\lambda)$는 $A'$에서 $x'$의 역원이다. 이로써 $A'$이
극대 아이디얼 $\mathfrak m'$을 갖는 국소환임이 증명된다. $K$에 관한
주장은 $\varinjlim$이 완전 함자라는 사실에서 곧바로 따른다. 가정
$\mathfrak m_\mu=\mathfrak m_\lambda A_\mu$는 표준 사상
$\mathfrak m_\lambda\otimes_{A_\lambda}A_\mu\to\mathfrak m_\mu$가
전사라는 뜻이다. 따라서 관계
$\mathfrak m'=\mathfrak m_\lambda A'$도 함자 $\varinjlim$의 완전성과
이 함자가 텐서곱과 가환한다는 사실에서 따른다.

이제 $\lambda<\mu$일 때
$\mathfrak m_\mu=\mathfrak m_\lambda A_\mu$이고 $A_\mu$가 평탄
$A_\lambda$-가군이라고 하자. 그러면
(\hyperref[0.6.2.3-ko]{$0_{\mathrm I}$, 6.2.3})에 따라 모든
$\lambda$에 대하여 $A'$은 평탄 $A_\lambda$-가군이다. $A'$과
$A_\lambda$가 국소환이고 $\mathfrak m'=\mathfrak m_\lambda A'$이므로,
$A'$은 실제로 충실 평탄 $A_\lambda$-가군이다
(\hyperref[0.6.6.2-ko]{$0_{\mathrm I}$, 6.6.2}). 끝으로 모든
$A_\lambda$가 \emph{뇌터}라고 하자. 그러면
$\mathfrak m_\lambda$-준아딕 위상들은 분리되어 있다
(\hyperref[0.7.3.5-ko]{$0_{\mathrm I}$, 7.3.5}). 이로부터 먼저 $A'$의
$\mathfrak m'$-아딕 위상이 \emph{분리}되어 있음을 보이자. 실제로
$x'\in A'$이 모든 $\mathfrak m'^n$ ($n>0$)에 속하면, 어떤 지표
$\mu$와 $x_\mu\in A_\mu$에 대하여 $x'$은 $x_\mu$의 상이다. 그런데
$A_\mu$에서 $\mathfrak m'^n=\mathfrak m_\mu^nA'$의 역상은
$\mathfrak m_\mu^n$이므로
(\hyperref[0.6.6.1-ko]{$0_{\mathrm I}$, 6.6.1}), $x_\mu$는 모든
$\mathfrak m_\mu^n$에 속한다. 따라서 가정에 의해 $x_\mu=0$이고,
결국 $x'=0$이다. $\widehat A'$을 $A'$의 $\mathfrak m'$-아딕 위상에
대한 완비화라 하자. 앞의 결과로부터 $A'\subset\widehat A'$이다.
$\widehat A'$이 \emph{뇌터}이고 모든 $\lambda$에 대하여
$A_\lambda$-\emph{평탄}임을 보이겠다. 그러면
\oldpage[0\textsubscript{III}]{22}
$\widehat A'$은 $A'$-평탄이고
(\hyperref[0.6.2.3-ko]{$0_{\mathrm I}$, 6.2.3}),
$\mathfrak m'\widehat A'\ne\widehat A'$이므로 $\widehat A'$은 충실 평탄
$A'$-가군이다 (\hyperref[0.6.6.2-ko]{$0_{\mathrm I}$, 6.6.2}). 따라서
마침내 $A'$이 \emph{뇌터}임을 얻으며
(\hyperref[0.6.5.2-ko]{$0_{\mathrm I}$, 6.5.2}), 이로써 보조정리의
증명이 끝난다.

$\widehat A'=\varprojlim_n A'/\mathfrak m'^n$이다. $A'$이
$A_\lambda$-평탄이므로
\[
  \mathfrak m'^{n}/\mathfrak m'^{n+1}
  =(\mathfrak m_\lambda^n/\mathfrak m_\lambda^{n+1})
    \otimes_{A_\lambda}A'
  =(\mathfrak m_\lambda^n/\mathfrak m_\lambda^{n+1})
    \otimes_{K_\lambda}(K_\lambda\otimes_{A_\lambda}A')
  =(\mathfrak m_\lambda^n/\mathfrak m_\lambda^{n+1})
    \otimes_{K_\lambda}K
\]
이다. $\mathfrak m_\lambda^n/\mathfrak m_\lambda^{n+1}$은 유한 차원
$K_\lambda$-벡터공간이므로, 모든 $n\geq0$에 대하여
$\mathfrak m'^n/\mathfrak m'^{n+1}$은 유한 차원 $K$-벡터공간이다.
따라서 (\hyperref[0.7.2.12-ko]{$0_{\mathrm I}$, 7.2.12})와
(\hyperref[0.7.2.8-ko]{$0_{\mathrm I}$, 7.2.8})에 의해
$\widehat A'$은 뇌터이다. 또한 $\widehat A'$의 극대 아이디얼은
$\mathfrak m'\widehat A'$이고,
$\widehat A'/\mathfrak m'^n\widehat A'$은 $A'/\mathfrak m'^n$과
동형임을 알고 있다. 그리고
$A'/\mathfrak m'^n=(A_\lambda/\mathfrak m_\lambda^n)
\otimes_{A_\lambda}A'$이므로, $A'/\mathfrak m'^n$은 평탄
$(A_\lambda/\mathfrak m_\lambda^n)$-가군이다
(\hyperref[0.6.2.1-ko]{$0_{\mathrm I}$, 6.2.1}). 따라서 판정법
(\hyperref[0.10.2.2-fr]{10.2.2})을 뇌터 $A_\lambda$-대수
$\widehat A'$에 적용할 수 있고, 이로부터 $\widehat A'$이
$A_\lambda$-평탄임을 얻는다. C.Q.F.D.

\begin{env}[10.3.1.4]
\label{0.10.3.1.4-fr}
이제 일반적인 경우를 다루자. 순서수 $\gamma$와, 모든 순서수
$\lambda\leq\gamma$에 대하여 $k$를 포함하는 $K$의 부분체
$k_\lambda$가 다음 조건들을 만족하도록 존재한다. $1^\circ$ 모든
$\lambda<\gamma$에 대하여 $k_{\lambda+1}$은 단 하나의 원소로 생성되는
$k_\lambda$의 확대체이다. $2^\circ$ 바로 앞 원소가 없는 모든 순서수
$\mu$에 대하여
$k_\mu=\bigcup_{\lambda<\mu}k_\lambda$이다. $3^\circ$ $K=k_\gamma$이다.
실제로 적당한 $\beta$에 대하여 순서수 $\xi\leq\beta$들의 집합에서
$K$로 가는 전단사 $\xi\mapsto t_\xi$를 택하면 충분하다.
$\lambda\leq\beta$에 대하여 $k_\lambda$를 초한 귀납법으로 다음과 같이
정의한다. $\lambda$에 바로 앞 원소가 없으면 $\mu<\lambda$인
$k_\mu$들의 합집합으로 두고, $\lambda=\nu+1$이면
$t_\xi\notin k_\nu$가 되는 가장 작은 순서수 $\xi$를 택하여
$k_\nu(t_\xi)$로 둔다. 그러면 $\gamma$는 정의상
$k_\gamma=K$가 되는 $\beta$ 이하의 가장 작은 순서수이다.

이제 초한 귀납법으로 $\lambda\leq\gamma$에 대한 뇌터 국소환
$A_\lambda$의 족과, $\lambda\leq\mu$에 대한 국소 준동형
$f_{\mu\lambda}:A_\lambda\to A_\mu$를 다음 조건들을 만족하도록
정의하겠다.
\begin{enumerate}
  \item[(i)] $(A_\lambda,f_{\mu\lambda})$는 귀납계이고 $A_0=A$이다.
  \item[(ii)] 모든 $\lambda$에 대하여 $k$-동형
    $A_\lambda/\mathfrak J A_\lambda\xrightarrow{\sim}k_\lambda$가 있다.
  \item[(iii)] $\lambda\leq\mu$이면 $A_\mu$는 평탄
    $A_\lambda$-가군이다.
\end{enumerate}

$\lambda<\mu<\xi$에 대해 $A_\lambda$와 $f_{\mu\lambda}$가 정의되어
있다고 하자. 먼저 $\xi=\zeta+1$이어서
$k_\xi=k_\zeta(t)$인 경우를 생각하자. $t$가 $k_\zeta$ 위에서
초월적이면, (\hyperref[0.10.3.1.1-fr]{10.3.1.1})의 절차에 따라
$A_\xi=(A_\zeta[t])_{\mathfrak J A_\zeta[t]}$로 정의한다.
$f_{\xi\zeta}$는 표준 사상이고, $\lambda<\zeta$이면
$f_{\xi\lambda}=f_{\xi\zeta}\circ f_{\zeta\lambda}$로 둔다.
(\hyperref[0.10.3.1.1-fr]{10.3.1.1})에서 증명한 것에 비추어 조건
(i)--(iii)의 확인은 즉시 이루어진다. 다음으로 $t$가 대수적이라고 하고,
$h$를 $k_\zeta[T]$에서의 그 최소다항식, $H$를 $k_\zeta[T]$에서의 상이
$h$인 $A_\zeta[T]$의 일계수 다항식이라 하자. 이때
$A_\xi=A_\zeta[T]/(H)$로 두고, $f_{\xi\lambda}$들은 앞에서와 같이
정의한다. 그러면 조건 (i)--(iii)은
(\hyperref[0.10.3.1.2-fr]{10.3.1.2})에서 본 결과로부터 따른다.

이제 $\xi$에 바로 앞 원소가 없다고 하자. $A_\xi$를
$\lambda<\xi$에 대한 국소환들의 귀납계
$(A_\lambda,f_{\mu\lambda})$의 귀납극한으로 두고,
$\lambda<\xi$에 대한 $f_{\xi\lambda}$는 표준 사상으로 정의한다.
$A_\xi$가 뇌터 국소환이라는 사실, $f_{\xi\lambda}$들이 국소
준동형이라는 사실, 그리고 $\lambda\leq\xi$에 대한 조건 (i)--(iii)은
\oldpage[0\textsubscript{III}]{23}
귀납가정과 보조정리
(\hyperref[0.10.3.1.3-fr]{10.3.1.3})에서 따른다. 이 구성을 마치면
$B=A_\gamma$가 (\hyperref[0.10.3.1-fr]{10.3.1})의 명제를 만족함이
명백하다.
\end{env}

(\hyperref[0.10.2.1-fr]{10.2.1, $c)$})에 따라 표준 동형
\[
\label{0.10.3.1.5-fr}
  \operatorname{gr}(A)\otimes_k K\xrightarrow{\sim}\operatorname{gr}(B).
  \tag{10.3.1.5}
\]
이 있음을 유의하자.

한편 (\hyperref[0.10.3.1-fr]{10.3.1})의 결론을 바꾸지 않고 $B$를 그
$\mathfrak J B$-아딕 완비화 $\widehat B$로 바꿀 수 있다. 실제로
$\widehat B$는 평탄 $B$-가군이고
(\hyperref[0.7.3.3-ko]{$0_{\mathrm I}$, 7.3.3}), 따라서 평탄
$A$-가군이다 (\hyperref[0.6.2.1-ko]{$0_{\mathrm I}$, 6.2.1}).

또한 다음 따름정리가 증명되었다.

\begin{corollary}[10.3.2]
\label{0.10.3.2-fr}
$K$가 유한 차수 확대체이면, $B$가 유한 $A$-대수라고 가정할 수 있다.
\end{corollary}

\section{호몰로지 대수학에 관한 보충}
\label{section:0.11-fr}

\subsection{스펙트럼 열에 관한 복습}
\label{subsection:0.11.1-fr}

\begin{env}[11.1.1]
\label{0.11.1.1-fr}
이후에는 (T, 2.4)에서 정의한 것보다 더 일반적인 스펙트럼 열의 개념을
사용한다. (T, 2.4)의 표기를 유지하면서, 아벨 범주 $\mathcal C$에서
다음 자료들로 이루어진 계 E를 \emph{스펙트럼 열}이라 부르겠다.
\begin{enumerate}
  \item[(a)] $p\in\mathbb Z$, $q\in\mathbb Z$, $r\geq2$에 대하여
    정의되는 $\mathcal C$의 대상들의 족 $(E_r^{pq})$.
  \item[(b)]
    $d_r^{p+r,q-r+1}d_r^{pq}=0$을 만족하는 사상들의 족
    $d_r^{pq}:E_r^{pq}\to E_r^{p+r,q-r+1}$. 다음과 같이 둔다.
    $Z_{r+1}(E_r^{pq})=\operatorname{Ker}(d_r^{pq})$,
    $B_{r+1}(E_r^{pq})=\operatorname{Im}(d_r^{p-r,q+r-1})$. 그러면
    \[
      B_{r+1}(E_r^{pq})\subset Z_{r+1}(E_r^{pq})\subset E_r^{pq}.
    \]
  \item[(c)] 동형사상들의 족
    $\alpha_r^{pq}:Z_{r+1}(E_r^{pq})/B_{r+1}(E_r^{pq})
    \xrightarrow{\sim}E_{r+1}^{pq}$.
\end{enumerate}

이제 $k\geq r+1$에 대하여 $k$에 대한 귀납법으로 부분대상
$B_k(E_r^{pq})$와 $Z_k(E_r^{pq})$를 다음과 같이 정의한다. 이들은
표준 사상 $E_r^{pq}\to E_r^{pq}/B_{r+1}(E_r^{pq})$에 의한, 각각
$\alpha_r^{pq}$를 통해 부분대상 $B_k(E_{r+1}^{pq})$와
$Z_k(E_{r+1}^{pq})$로 식별되는 이 몫의 부분대상들의 역상이다. 그러면
동형을 제외하면 명백히
\[
\label{0.11.1.1.1-fr}
  Z_k(E_r^{pq})/B_k(E_r^{pq})=E_k^{pq}
  \qquad(k\geq r+1\text{에 대하여})
  \tag{11.1.1.1}
\]
이고, 다시 $B_r(E_r^{pq})=0$, $Z_r(E_r^{pq})=E_r^{pq}$로 두면 다음
포함관계들이 성립한다.
\[
\label{0.11.1.1.2-fr}
  0=B_r(E_r^{pq})\subset B_{r+1}(E_r^{pq})\subset B_{r+2}(E_r^{pq})
  \subset\cdots\subset Z_{r+2}(E_r^{pq})\subset Z_{r+1}(E_r^{pq})
  \subset Z_r(E_r^{pq})=E_r^{pq}.
  \tag{11.1.1.2}
\]

E의 나머지 자료는 다음과 같다.
\begin{enumerate}
  \item[(d)] 다음을 만족하는 $E_2^{pq}$의 두 부분대상
    $B_\infty(E_2^{pq})$와 $Z_\infty(E_2^{pq})$:
    $B_\infty(E_2^{pq})\subset Z_\infty(E_2^{pq})$이고, 모든
    $k\geq2$에 대하여
    \[
      B_k(E_2^{pq})\subset B_\infty(E_2^{pq})
      \qquad\text{및}\qquad
      Z_\infty(E_2^{pq})\subset Z_k(E_2^{pq}).
    \]
    다음과 같이 둔다.
    \[
    \label{0.11.1.1.3-fr}
      E_\infty^{pq}=Z_\infty(E_2^{pq})/B_\infty(E_2^{pq}).
      \tag{11.1.1.3}
    \]
  \item[(e)]
    \oldpage[0\textsubscript{III}]{24}
    $\mathcal C$의 대상들의 족 $(E^n)$. 각 대상에는 \emph{감소 여과}
    $(F^p(E^n))_{p\in\mathbb Z}$가 주어져 있다. 통상과 같이
    $\operatorname{gr}(E^n)$으로 여과 대상 $E^n$에 연관된 등급 대상을
    나타낸다. 이는
    $\operatorname{gr}_p(E^n)=F^p(E^n)/F^{p+1}(E^n)$들의 직합이다.
  \item[(f)] 모든 쌍 $(p,q)\in\mathbb Z\times\mathbb Z$에 대한 동형사상
    $\beta^{pq}:E_\infty^{pq}\xrightarrow{\sim}
    \operatorname{gr}_p(E^{p+q})$.
\end{enumerate}

여과를 제외한 족 $(E^n)$을 스펙트럼 열 E의 \emph{수렴 대상}이라
부른다.

범주 $\mathcal C$가 무한 직합을 가지거나, 모든 $r\geq2$와 모든
$n\in\mathbb Z$에 대하여 $p+q=n$이고 $E_r^{pq}\ne0$인 쌍 $(p,q)$이
유한 개라고 하자($r=2$에 대해서만 그러하다고 해도 충분하다). 그러면
\[
  E_r^{(n)}=\sum_{p+q=n}E_r^{pq}
\]
이 정의된다. $p+q=n$인 각 쌍 $(p,q)$에 대하여 $E_r^{pq}$로의 제한이
$d_r^{pq}$인 사상 $E_r^{(n)}\to E_r^{(n+1)}$을 $d_r^{(n)}$으로
나타내면 $d_r^{(n+1)}\circ d_r^{(n)}=0$이다. 다시 말해
$(E_r^{(n)})_{n\in\mathbb Z}$는 $\mathcal C$에서 차수 $+1$인 미분
연산자를 갖는 \emph{복합체 $E_r^{(\bullet)}$}이다. 그리고 c)로부터
모든 $r\geq2$에 대하여 $\mathrm H^n(E_r^{(\bullet)})$가
$E_{r+1}^{(n)}$과 동형임을 얻는다.
\end{env}

\begin{env}[11.1.2]
\label{0.11.1.2-fr}
스펙트럼 열 E에서 스펙트럼 열
$E'=(E_r^{\prime pq},E^{\prime n})$로 가는 \emph{사상 $u:E\to E'$}은
사상들의 계
$u_r^{pq}:E_r^{pq}\to E_r^{\prime pq}$,
$u^n:E^n\to E^{\prime n}$으로 이루어진다. 여기서 $u^n$들은
$E^n$과 $E^{\prime n}$의 여과와 양립하고, 도식
\[
  \xymatrix{
    E_r^{pq}\ar[r]^-{d_r^{pq}}\ar[d]_{u_r^{pq}} &
    E_r^{p+r,q-r+1}\ar[d]^{u_r^{p+r,q-r+1}}\\
    E_r^{\prime pq}\ar[r]_{d_r^{\prime pq}} &
    E_r^{\prime p+r,q-r+1}
  }
\]
은 가환한다. 또한 몫을 취하면 $u_r^{pq}$는 사상
\[
  \overline u_r^{pq}:
  Z_{r+1}(E_r^{pq})/B_{r+1}(E_r^{pq})\longrightarrow
  Z_{r+1}(E_r^{\prime pq})/B_{r+1}(E_r^{\prime pq})
\]
을 주며,
$\alpha_r^{\prime pq}\circ\overline u_r^{pq}
=u_{r+1}^{pq}\circ\alpha_r^{pq}$이어야 한다. 끝으로
$u_2^{pq}(B_\infty(E_2^{pq}))\subset B_\infty(E_2^{\prime pq})$,
$u_2^{pq}(Z_\infty(E_2^{pq}))\subset Z_\infty(E_2^{\prime pq})$이어야
한다. 몫을 취하면 $u_2^{pq}$는 사상
$u_\infty^{pq}:E_\infty^{pq}\to E_\infty^{\prime pq}$을 주며, 도식
\[
  \xymatrix{
    E_\infty^{pq}\ar[r]^-{u_\infty^{pq}}\ar[d]_{\beta^{pq}} &
    E_\infty^{\prime pq}\ar[d]^{\beta^{\prime pq}}\\
    \operatorname{gr}_p(E^{p+q})
      \ar[r]_{\operatorname{gr}_p(u^{p+q})} &
    \operatorname{gr}_p(E^{\prime p+q})
  }
\]
은 가환해야 한다.

앞의 정의들로부터 $r$에 대한 귀납법에 의해 다음을 알 수 있다.
$u_2^{pq}$들이 \emph{동형사상}이면 $r\geq2$에 대한
$u_r^{pq}$들도 동형사상이다. \emph{더 나아가}
$u_2^{pq}(B_\infty(E_2^{pq}))=B_\infty(E_2^{\prime pq})$,
$u_2^{pq}(Z_\infty(E_2^{pq}))=Z_\infty(E_2^{\prime pq})$이고
$u^n$들이 동형사상임을 \emph{알고 있다면}, $u$가
\emph{동형사상}이라고 결론지을 수 있다.
\end{env}

\begin{env}[11.1.3]
\label{0.11.1.3-fr}
\oldpage[0\textsubscript{III}]{25}
$(F^p(X))_{p\in\mathbb Z}$가 대상 $X\in\mathcal C$의 (감소) 여과라고
하자. $\inf(F^p(X))=0$이면 이 여과를 \emph{분리}, 어떤 $p$에 대해
$F^p(X)=0$이면 \emph{이산}, $\sup(F^p(X))=X$이면 \emph{소진적}(또는
\emph{공분리}), 어떤 $p$에 대해 $F^p(X)=X$이면 \emph{공이산}이라고
함을 상기하자.

스펙트럼 열 $E=(E_r^{pq},E^n)$가
\[
  B_\infty(E_2^{pq})=\sup_k(B_k(E_2^{pq})),\qquad
  Z_\infty(E_2^{pq})=\inf_k(Z_k(E_2^{pq}))
\]
를 만족하면 \emph{약수렴}한다고 한다. 다시 말해 대상
$B_\infty(E_2^{pq})$와 $Z_\infty(E_2^{pq})$가 스펙트럼 열 E의 자료
a)--c)에 의해 결정된다는 뜻이다. 스펙트럼 열 E가 약수렴하고, 더
나아가 다음을 만족하면 \emph{정칙}이라고 한다.
\begin{enumerate}
  \item[$1^\circ$] 모든 쌍 $(p,q)$에 대하여 감소열
    $(Z_k(E_2^{pq}))_{k\geq2}$는 \emph{어떤 단계 이후로 일정하다}.
    그러면 E가 약수렴한다는 가정에 따라, $p$와 $q$에 의존하는 충분히
    큰 $k$에 대하여
    $Z_\infty(E_2^{pq})=Z_k(E_2^{pq})$이다.
  \item[$2^\circ$] 모든 $n$에 대하여 $E^n$의 여과
    $(F^p(E^n))_{p\in\mathbb Z}$는 \emph{이산}이고 \emph{소진적}이다.
\end{enumerate}

스펙트럼 열 E가 약수렴하고, 더 나아가 다음을 만족하면
\emph{공정칙}이라고 한다.
\begin{enumerate}
  \item[$3^\circ$] 모든 쌍 $(p,q)$에 대하여 증가열
    $(B_k(E_2^{pq}))_{k\geq2}$는 \emph{어떤 단계 이후로 일정하다}.
    따라서 $B_\infty(E_2^{pq})=B_k(E_2^{pq})$이고, 이어서
    $E_\infty^{pq}=\inf E_k^{pq}$이다.
  \item[$4^\circ$] 모든 $n$에 대하여 $E^n$의 여과는
    \emph{공이산}이다.
\end{enumerate}

끝으로 E가 정칙인 동시에 공정칙이면 \emph{쌍정칙}이라고 한다. 다시
말해 다음 조건들이 성립한다.
\begin{enumerate}
  \item[(a)] 모든 쌍 $(p,q)$에 대하여 열
    $(B_k(E_2^{pq}))_{k\geq2}$과 $(Z_k(E_2^{pq}))_{k\geq2}$은
    \emph{어떤 단계 이후로 일정}하고, 충분히 큰 $k$에 대하여
    $B_\infty(E_2^{pq})=B_k(E_2^{pq})$ 및
    $Z_\infty(E_2^{pq})=Z_k(E_2^{pq})$이다. 따라서
    $E_\infty^{pq}=E_k^{pq}$이다.
  \item[(b)] 모든 $n$에 대하여 여과
    $(F^p(E^n))_{p\in\mathbb Z}$는 \emph{이산}이고 \emph{공이산}이다.
    이는 이 여과가 \emph{유한}이라고도 표현한다.
\end{enumerate}

따라서 (T, 2.4)에서 정의한 스펙트럼 열들은 쌍정칙 스펙트럼 열이다.
\end{env}

\begin{env}[11.1.4]
\label{0.11.1.4-fr}
범주 $\mathcal C$에서 여과 귀납극한들이 존재하고 함자 $\varinjlim$이
\emph{완전}하다고 하자. 이는 (T, 1.5)의 공리 AB 5)가 성립한다는 것과
동치이다(T, 1.8 참조). 그러면 대상 $X\in\mathcal C$의 여과
$(F^p(X))_{p\in\mathbb Z}$가 소진적이라는 조건은
$\varinjlim_{p\to-\infty}F^p(X)=X$로도 쓸 수 있다. 스펙트럼 열 E가
약수렴하면
$B_\infty(E_2^{pq})=\varinjlim_{k\to\infty}B_k(E_2^{pq})$이다. 더
나아가 $u:E\to E'$가 E에서 $\mathcal C$의 약수렴 스펙트럼 열 $E'$로
가는 사상이면, $\varinjlim$의 완전성에 의해
$u_2^{pq}(B_\infty(E_2^{pq}))=B_\infty(E_2^{\prime pq})$이다. 또한
다음이 성립한다.
\end{env}

\begin{proposition}[11.1.5]
\label{0.11.1.5-fr}
$\mathcal C$를 여과 귀납극한들이 완전한 아벨 범주, E와 $E'$을
$\mathcal C$의 두 정칙 스펙트럼 열, $u:E\to E'$을 스펙트럼 열들의
사상이라 하자. $u_2^{pq}$들이 동형사상이면 $u$도 동형사상이다.
\end{proposition}

\begin{proof}
$u_r^{pq}$들이 동형사상이고
\[
  u_2^{pq}(B_\infty(E_2^{pq}))=B_\infty(E_2^{\prime pq})
\]
임을 이미 알고 있다 (\hyperref[0.11.1.2-fr]{11.1.2}).
\oldpage[0\textsubscript{III}]{26}
E와 $E'$이 정칙이라는 가정으로부터
$u_2^{pq}(Z_\infty(E_2^{pq}))=Z_\infty(E_2^{\prime pq})$도 따른다.
$u_2^{pq}$가 동형사상이므로 $u_\infty^{pq}$도 동형사상이고, 따라서
$\operatorname{gr}_p(u^{p+q})$도 동형사상이라고 결론지을 수 있다.
그런데 $E^n$과 $E^{\prime n}$의 여과들은 이산이고 소진적이므로,
이로부터 $u^n$들도 동형사상임이 따른다(부르바키,
\emph{Alg. comm.}, 제III장, \S~2, n\textsuperscript{o} 8, 정리 1).
\end{proof}

\begin{env}[11.1.6]
\label{0.11.1.6-fr}
(\hyperref[0.11.1.1.2-fr]{11.1.1.2})와 정의
(\hyperref[0.11.1.1.3-fr]{11.1.1.3})로부터, 스펙트럼 열 E에서
$E_r^{pq}=0$이면 $k\geq r$에 대하여 $E_k^{pq}=0$이고
$E_\infty^{pq}=0$임이 따른다. 어떤 정수 $r\geq2$와, 모든 정수
$n\in\mathbb Z$에 대한 정수 $q(n)$이 존재하여 $q\ne q(n)$이면
$E_r^{n-q,q}=0$일 때, 스펙트럼 열을 \emph{퇴화}한다고 한다. 앞의
관찰로부터 우선 $k\geq r$($k=\infty$ 포함)이고 $q\ne q(n)$이면
$E_k^{n-q,q}=0$임을 얻는다. 또한 $E_{r+1}^{pq}$의 정의로부터
$E_{r+1}^{n-q(n),q(n)}=E_r^{n-q(n),q(n)}$이다. E가 약수렴하면 따라서
$E_\infty^{n-q(n),q(n)}=E_r^{n-q(n),q(n)}$도 성립한다. 다시 말해 모든
$n\in\mathbb Z$에 대하여 $p\ne n-q(n)$이면
$\operatorname{gr}_p(E^n)=0$이고,
$\operatorname{gr}_{n-q(n)}(E^n)=E_r^{n-q(n),q(n)}$이다. 더 나아가
$E^n$의 여과가 이산이고 소진적이면, 열 E는 정칙이고 동형을 제외하면
$E^n=E_r^{n-q(n),q(n)}$이다.
\end{env}

\begin{env}[11.1.7]
\label{0.11.1.7-fr}
범주 $\mathcal C$에서 여과 귀납극한들이 존재하고 완전하다고 하자.
$(E_\lambda,u_{\mu\lambda})$를 여과 지표집합에 따른 $\mathcal C$의
스펙트럼 열들의 귀납계라 하자. 그러면 이 귀납계의 \emph{귀납극한}은
$\mathcal C$의 대상들로 이루어진 스펙트럼 열들의 가법 범주에서
존재한다. 이를 보려면 $E_r^{pq}$, $d_r^{pq}$, $\alpha_r^{pq}$,
$B_\infty(E_2^{pq})$, $Z_\infty(E_2^{pq})$, $E^n$, $F^p(E^n)$,
$\beta^{pq}$를 각각 $E_{r,\lambda}^{pq}$, $d_{r,\lambda}^{pq}$,
$\alpha_{r,\lambda}^{pq}$, $B_\infty(E_{2,\lambda}^{pq})$,
$Z_\infty(E_{2,\lambda}^{pq})$, $E_\lambda^n$,
$F^p(E_\lambda^n)$, $\beta_\lambda^{pq}$의 귀납극한으로 정의하면
충분하다. (\hyperref[0.11.1.1-fr]{11.1.1})의 조건들은 $\mathcal C$에서
함자 $\varinjlim$이 완전하다는 사실로부터 확인된다.
\end{env}

\begin{remark}[11.1.8]
\label{0.11.1.8-fr}
범주 $\mathcal C$가 \emph{뇌터}환 $A$ 위의 $A$-가군들의 범주(각각,
환 $A$ 위의 $A$-가군들의 범주)라고 하자. 그러면
(\hyperref[0.11.1.1-fr]{11.1.1})의 정의들로부터, 주어진 $r$에 대하여
$E_r^{pq}$들이 \emph{유한 생성} $A$-가군(각각, \emph{유한 길이}
$A$-가군)이면 모든 $s\geq r$에 대한 $E_s^{pq}$와
$E_\infty^{pq}$도 그러함을 알 수 있다. 더 나아가 모든 $n$에 대해
수렴 대상 $(E^n)$의 여과가 이산이고 공이산이면, 각 $E^n$도
\emph{유한 생성} $A$-가군(각각, \emph{유한 길이} $A$-가군)이라고
결론지을 수 있다.
\end{remark}

\begin{env}[11.1.9]
\label{0.11.1.9-fr}
스펙트럼 열 E가 $p+q=n$에 관하여 ``균일하게'' 쌍정칙임을 보장하는
조건들을 살펴볼 것이다. 이때 다음 보조정리를 사용한다.
\end{env}

\begin{lemma}[11.1.10]
\label{0.11.1.10-fr}
$(E_r^{pq})$를 (\hyperref[0.11.1.1-fr]{11.1.1})의 자료 a), b), c)로
연결된 $\mathcal C$의 대상들의 족이라 하자. 고정된 정수 $n$에 대하여
다음 성질들은 서로 동치이다.
\begin{enumerate}
  \item[(a)] 정수 $r(n)$이 존재하여, $r\geq r(n)$이고 $p+q=n$ 또는
    $p+q=n-1$이면 모든 사상 $d_r^{pq}$가 0이다.
  \item[(b)] 정수 $r(n)$이 존재하여, $p+q=n$ 또는 $p+q=n+1$이면
    $s\geq r\geq r(n)$에 대하여
    $B_r(E_2^{pq})=B_s(E_2^{pq})$이다.
  \item[(c)] 정수 $r(n)$이 존재하여, $p+q=n$ 또는 $p+q=n-1$이면
    $s\geq r\geq r(n)$에 대하여
    $Z_r(E_2^{pq})=Z_s(E_2^{pq})$이다.
  \item[(d)] 정수 $r(n)$이 존재하여, $p+q=n$이면
    $s\geq r\geq r(n)$에 대하여
    $B_r(E_2^{pq})=B_s(E_2^{pq})$ 및
    $Z_r(E_2^{pq})=Z_s(E_2^{pq})$이다.
\end{enumerate}
\end{lemma}

\begin{proof}
\oldpage[0\textsubscript{III}]{27}
실제로 (\hyperref[0.11.1.1-fr]{11.1.1})의 조건 a), b), c)에 따라
$Z_{r+1}(E_2^{pq})=Z_r(E_2^{pq})$이라는 것은 $d_r^{pq}=0$이라는
것과 동치이고,
$B_r(E_2^{p+r,q-r+1})=B_{r+1}(E_2^{p+r,q-r+1})$이라는 것도
$d_r^{pq}=0$이라는 것과 동치이다. 보조정리는 이 관찰에서 즉시 따른다.
\end{proof}

\subsection{여과 복합체의 스펙트럼 열}
\label{subsection:0.11.2-fr}

\begin{env}[11.2.1]
\label{0.11.2.1-fr}
아벨 범주 $\mathcal C$가 주어졌다고 하자. 미분 연산자의 차수가 $+1$인
$\mathcal C$의 대상들의 \emph{복합체} $(K^i)_{i\in\mathbb Z}$는
$K^\bullet$과 같은 표기로 나타내고, 미분 연산자의 차수가 $-1$인
$\mathcal C$의 대상들의 복합체 $(K_i)_{i\in\mathbb Z}$는 $K_\bullet$과 같은
표기로 나타내기로 한다. 미분 연산자 $d$의 차수가 $+1$인 모든 복합체
$K^\bullet=(K^i)$에 대하여 $K'_i=K^{-i}$로 놓으면 복합체
$K'_\bullet=(K'_i)$를 대응시킬 수 있으며, 이때 미분 연산자
$K'_i\to K'_{i-1}$는 연산자 $d:K^{-i}\to K^{-i+1}$이다. 그 역도
마찬가지이다. 따라서 상황에 따라 어느 한 종류의 복합체를 고찰하고, 한
종류에 대한 모든 결과를 다른 종류에 대한 결과로 옮길 수 있다. 마찬가지로
$K^{\bullet\bullet}=(K^{ij})$ (각각 $K_{\bullet\bullet}=(K_{ij})$)와 같은
표기로 두 미분 연산자의 차수가 모두 $+1$ (각각 $-1$)인 $\mathcal C$의
대상들의 \emph{이중복합체}를 나타낸다. 첨자의 부호를 바꾸면 다시 한
종류에서 다른 종류로 넘어가며, 임의의 다중복합체에도 이와 유사한 표기와
고찰을 적용한다. $K^\bullet$ 또는 $K_\bullet$이라는 표기는 반드시 복합체일
필요는 없는 $\mathbb Z$형의 $\mathcal C$의 \emph{등급 대상}에도 사용한다
(또는 영인 미분 연산자를 주어 이를 복합체로 볼 수 있다). 예를 들어 미분
연산자의 차수가 $+1$인 복합체 $K^\bullet$의 \emph{코호몰로지}를
$\mathrm H^\bullet(K^\bullet)=(\mathrm H^i(K^\bullet))_{i\in\mathbb Z}$로
쓰고, 미분 연산자의 차수가 $-1$인 복합체 $K_\bullet$의
\emph{호몰로지}를
$\mathrm H_\bullet(K_\bullet)=(\mathrm H_i(K_\bullet))_{i\in\mathbb Z}$로
쓴다. 위에서 설명한 연산으로 $K^\bullet$에서 $K'_\bullet$로 넘어가면
$\mathrm H_i(K'_\bullet)=\mathrm H^{-i}(K^\bullet)$이다.

이와 관련하여, 복합체 $K^\bullet$ (각각 $K_\bullet$)에 대하여 일반적으로
$Z^i(K^\bullet)=\operatorname{Ker}(K^i\to K^{i+1})$ (``공순환 대상'')와
$B^i(K^\bullet)=\operatorname{Im}(K^{i-1}\to K^i)$ (``공경계 대상'')
(각각 $Z_i(K_\bullet)=\operatorname{Ker}(K_i\to K_{i-1})$ (``순환 대상'')와
$B_i(K_\bullet)=\operatorname{Im}(K_{i+1}\to K_i)$ (``경계 대상''))로
쓰기로 한다. 따라서
$\mathrm H^i(K^\bullet)=Z^i(K^\bullet)/B^i(K^\bullet)$ (각각
$\mathrm H_i(K_\bullet)=Z_i(K_\bullet)/B_i(K_\bullet)$)이다.

$K^\bullet=(K^i)$ (각각 $K_\bullet=(K_i)$)가 $\mathcal C$의 복합체이고
$T:\mathcal C\to\mathcal C'$가 $\mathcal C$에서 아벨 범주 $\mathcal C'$로
가는 함자이면, $\mathcal C'$의 복합체 $(T(K^i))$ (각각 $(T(K_i))$)를
$T(K^\bullet)$ (각각 $T(K_\bullet)$)로 나타낸다.

$\partial$-함자의 정의(T, 2.1)는 되풀이하지 않는다. 다만 사상
$\partial$가 차수를 하나 낮출 때에는 $\partial^*$-함자 대신 역시
$\partial$-함자라고 부르며, 혼동의 여지가 있을 때에는 매번 문맥으로 이를
명시한다.

마지막으로, $\mathcal C$의 \emph{등급 대상} $(A_i)_{i\in\mathbb Z}$에
대하여 어떤 $i_0$가 존재하여 $i<i_0$ (각각 $i>i_0$)이면 $A_i=0$일 때,
이를 \emph{아래로 유계} (각각 \emph{위로 유계})라고 한다.
\end{env}

\begin{env}[11.2.2]
\label{0.11.2.2-fr}
$K^\bullet$을 미분 연산자 $d$의 차수가 $+1$인 $\mathcal C$의 복합체라
하자. 또한 $K^\bullet$에 그 등급 부분대상들로 이루어진 \emph{여과}
$F(K^\bullet)=(F^p(K^\bullet))_{p\in\mathbb Z}$가 주어졌다고 하자. 다시
말해
\oldpage[0\textsubscript{III}]{28}
$F^p(K^\bullet)=(K^i\cap F^p(K^\bullet))_{i\in\mathbb Z}$이다. 더 나아가
모든 $p\in\mathbb Z$에 대하여
$d(F^p(K^\bullet))\subset F^p(K^\bullet)$라고 가정한다. 이제
$K^\bullet$으로부터 \emph{함자적으로} 스펙트럼 열 $E(K^\bullet)$을
정의하는 방법을 간단히 상기하자(M, XV, 4 및 G, I, 4.3). $r\geq2$일 때,
표준 사상
$F^p(K^\bullet)/F^{p+r}(K^\bullet)\to
F^p(K^\bullet)/F^{p+1}(K^\bullet)$은 코호몰로지에 대한 사상
\[
  \mathrm H^{p+q}(F^p(K^\bullet)/F^{p+r}(K^\bullet))
  \longrightarrow
  \mathrm H^{p+q}(F^p(K^\bullet)/F^{p+1}(K^\bullet))
\]
을 정의한다. 이 사상의 상을 $Z_r^{pq}(K^\bullet)$로 나타낸다. 마찬가지로
완전열
\[
  0\to F^p(K^\bullet)/F^{p+1}(K^\bullet)
  \to F^{p-r+1}(K^\bullet)/F^{p+1}(K^\bullet)
  \to F^{p-r+1}(K^\bullet)/F^p(K^\bullet)\to0
\]
로부터 코호몰로지 완전열에 의해 사상
\[
  \mathrm H^{p+q-1}(F^{p-r+1}(K^\bullet)/F^p(K^\bullet))
  \longrightarrow
  \mathrm H^{p+q}(F^p(K^\bullet)/F^{p+1}(K^\bullet))
\]
을 얻으며, 이 사상의 상을 $B_r^{pq}(K^\bullet)$로 나타낸다. 이때
$B_r^{pq}(K^\bullet)\subset Z_r^{pq}(K^\bullet)$임을 보일 수 있고,
$E_r^{pq}(K^\bullet)=Z_r^{pq}(K^\bullet)/B_r^{pq}(K^\bullet)$로 놓는다.
여기서는 $d_r^{pq}$와 $\alpha_r^{pq}$의 정의를 명시하지 않는다.

여기서 $p$, $q$를 고정하면 모든 $Z_r^{pq}(K^\bullet)$와
$B_r^{pq}(K^\bullet)$는 같은 대상
$\mathrm H^{p+q}(F^p(K^\bullet)/F^{p+1}(K^\bullet))$의 부분대상이며, 이
대상을 $Z_1^{pq}(K^\bullet)$로 쓴다는 점에 유의하자.
$B_1^{pq}(K^\bullet)=0$으로 놓으면 위의 $Z_r^{pq}(K^\bullet)$와
$B_r^{pq}(K^\bullet)$의 정의는 $r=1$에도 적용된다. 또한
$E_1^{pq}(K^\bullet)=Z_1^{pq}(K^\bullet)$로 놓는다.
(\hyperref[0.11.1.1-fr]{11.1.1})의 조건들이 $r=1$에 대해서도 성립하도록
$d_1^{pq}$와 $\alpha_1^{pq}$를 정의한다. 한편 부분대상
$Z_\infty^{pq}(K^\bullet)$는 사상
\[
  \mathrm H^{p+q}(F^p(K^\bullet))\longrightarrow
  \mathrm H^{p+q}(F^p(K^\bullet)/F^{p+1}(K^\bullet))
  =E_1^{pq}(K^\bullet)
\]
의 상으로 정의하고, $B_\infty^{pq}(K^\bullet)$는 사상
\[
  \mathrm H^{p+q-1}(K^\bullet/F^p(K^\bullet))\longrightarrow
  \mathrm H^{p+q}(F^p(K^\bullet)/F^{p+1}(K^\bullet))
  =E_1^{pq}(K^\bullet)
\]
의 상으로 정의한다. 이 사상은 위와 같이 코호몰로지 완전열에서 얻는다.
$Z_\infty(E_2^{pq}(K^\bullet))$와 $B_\infty(E_2^{pq}(K^\bullet))$로는
$Z_\infty^{pq}(K^\bullet)$와 $B_\infty^{pq}(K^\bullet)$의
$E_2^{pq}(K^\bullet)$ 안에서의 표준 상을 취한다.

마지막으로 $F^p(\mathrm H^n(K^\bullet))$는 표준 포함 사상
$F^p(K^\bullet)\to K^\bullet$에서 오는 사상
$\mathrm H^n(F^p(K^\bullet))\to\mathrm H^n(K^\bullet)$의
$\mathrm H^n(K^\bullet)$ 안에서의 상으로 나타낸다. 코호몰로지 완전열에
의하면 이는 또한 사상
$\mathrm H^n(K^\bullet)\to\mathrm H^n(K^\bullet/F^p(K^\bullet))$의
핵이다. 이로써 $E^n(K^\bullet)=\mathrm H^n(K^\bullet)$ 위의 여과를
정의한다. 여기서는 동형사상 $\beta^{pq}$의 정의도 제시하지 않는다.
\end{env}

\begin{env}[11.2.3]
\label{0.11.2.3-fr}
$E(K^\bullet)$의 \emph{함자성}은 다음과 같은 뜻으로 이해해야 한다.
$\mathcal C$의 두 \emph{여과 복합체} $K^\bullet$,
$K^{\prime\bullet}$과 \emph{여과와 양립하는} 복합체 사상
$u:K^\bullet\to K^{\prime\bullet}$이 주어지면, 이로부터 자명한 방식으로
사상 $u_r^{pq}$ ($r\geq1$)와 $u^n$을 얻는다. 이 사상들이
(\hyperref[0.11.1.2-fr]{11.1.2})의 의미에서 $d_r^{pq}$,
$\alpha_r^{pq}$ 및 $\beta^{pq}$와 양립함을 보일 수 있으므로, 이들은
실제로 스펙트럼 열의 사상
$E(u):E(K^\bullet)\to E(K^{\prime\bullet})$을 정의한다. 더 나아가 $u$와
$v$가 위와 같은 종류의 $K^\bullet\to K^{\prime\bullet}$ 사상들이고
\emph{차수 $\leq k$로 호모토픽}이면, $r>k$에 대하여
$u_r^{pq}=v_r^{pq}$이고 모든 $n$에 대하여 $u^n=v^n$이다(M, XV, 3.1).
\end{env}

\begin{env}[11.2.4]
\label{0.11.2.4-fr}
\oldpage[0\textsubscript{III}]{29}
$\mathcal C$에서 여과 귀납극한들이 완전하다고 가정하자. 그러면
$K^\bullet$의 여과 $(F^p(K^\bullet))$가 \emph{소진적}이면 모든 $n$에
대하여 여과 $(F^p(\mathrm H^n(K^\bullet)))$도 소진적이다. 실제로 가정에
의해 $K^\bullet=\varinjlim_{p\to-\infty}F^p(K^\bullet)$이고,
$\mathcal C$에 대한 가정에 의해 코호몰로지는 귀납극한과 교환한다. 같은
이유로
$B_\infty(E_2^{pq}(K^\bullet))=\sup_k B_k(E_2^{pq}(K^\bullet))$이다.
모든 $n$에 대하여 어떤 정수 $u(n)$가 존재하여 $p>u(n)$이면
$\mathrm H^n(F^p(K^\bullet))=0$일 때, $K^\bullet$의 여과
$(F^p(K^\bullet))$를 \emph{정칙}이라고 한다. 특히 $K^\bullet$의 여과가
\emph{이산}이면 그러하다. $K^\bullet$의 여과가 정칙이고 소진적이며
$\mathcal C$에서 여과 귀납극한들이 완전하면, 스펙트럼 열
$E(K^\bullet)$이 \emph{정칙}임을 보일 수 있다(M, XV, 4).
\end{env}

\subsection{이중복합체의 스펙트럼 열}
\label{subsection:0.11.3-fr}

\begin{env}[11.3.1]
\label{0.11.3.1-fr}
이중복합체에 관한 규약으로는 (M)의 규약보다 (T, 2.4)의 규약을 따른다.
따라서 그러한 이중복합체 $K^{\bullet\bullet}=(K^{ij})$의 두 미분 연산자
$d'$, $d''$ (차수 $+1$)는 \emph{서로 교환한다}고 가정한다. 다음 두 조건
중 하나가 성립한다고 하자.
\begin{enumerate}
  \item[$1^\circ$] $\mathcal C$에서 무한 직합들이 존재한다.
  \item[$2^\circ$] 모든 $n\in\mathbb Z$에 대하여 $p+q=n$이고
    $K^{pq}\ne0$인 쌍 $(p,q)$는 \emph{유한}개뿐이다.
\end{enumerate}
그러면 이중복합체 $K^{\bullet\bullet}$는
$K^{\prime n}=\sum_{i+j=n}K^{ij}$인 (전체) 복합체
$(K^{\prime n})_{n\in\mathbb Z}$를 정의한다. 이 복합체의 미분 연산자
$d$ (차수 $+1$)는 $x\in K^{ij}$에 대하여
$dx=d'x+(-1)^i d''x$로 주어진다. 이하에서 \emph{이중복합체
$K^{\bullet\bullet}$가 정의하는 (전체) 복합체}를 말할 때에는 언제나 위
조건들 중 하나가 성립한다고 암묵적으로 가정한다. 다중복합체에 대해서도
이와 유사한 규약을 채택한다.

$K^{i,\bullet}$ (각각 $K^{\bullet,j}$)로 전체 복합체
$(K^{ij})_{j\in\mathbb Z}$ (각각 $(K^{ij})_{i\in\mathbb Z}$)를 나타내고,
$Z_{\mathrm{II}}^p(K^{i,\bullet})$, $B_{\mathrm{II}}^p(K^{i,\bullet})$,
$\mathrm H_{\mathrm{II}}^p(K^{i,\bullet})$ (각각
$Z_{\mathrm I}^p(K^{\bullet,j})$, $B_{\mathrm I}^p(K^{\bullet,j})$,
$\mathrm H_{\mathrm I}^p(K^{\bullet,j})$)로 각각 그 $p$번째 공순환 대상,
공경계 대상 및 코호몰로지 대상을 나타낸다. 미분
$d':K^{i,\bullet}\to K^{i+1,\bullet}$은 복합체의 사상이므로 공순환,
공경계 및 코호몰로지 위에 다음 연산자들을 준다.
\[
\begin{aligned}
  d'&:Z_{\mathrm{II}}^p(K^{i,\bullet})
      \longrightarrow Z_{\mathrm{II}}^p(K^{i+1,\bullet}),\\
  d'&:B_{\mathrm{II}}^p(K^{i,\bullet})
      \longrightarrow B_{\mathrm{II}}^p(K^{i+1,\bullet}),\\
  d'&:\mathrm H_{\mathrm{II}}^p(K^{i,\bullet})
      \longrightarrow \mathrm H_{\mathrm{II}}^p(K^{i+1,\bullet}).
\end{aligned}
\]
이 연산자들에 대하여
$(Z_{\mathrm{II}}^p(K^{i,\bullet}))_{i\in\mathbb Z}$,
$(B_{\mathrm{II}}^p(K^{i,\bullet}))_{i\in\mathbb Z}$ 및
$(\mathrm H_{\mathrm{II}}^p(K^{i,\bullet}))_{i\in\mathbb Z}$가
복합체임은 자명하다. 복합체
$(\mathrm H_{\mathrm{II}}^p(K^{i,\bullet}))_{i\in\mathbb Z}$를
$\mathrm H_{\mathrm{II}}^p(K^{\bullet\bullet})$로 나타내고, 그 $q$번째
공순환 대상, 공경계 대상 및 코호몰로지 대상을 각각
$Z_{\mathrm I}^q(\mathrm H_{\mathrm{II}}^p(K^{\bullet\bullet}))$,
$B_{\mathrm I}^q(\mathrm H_{\mathrm{II}}^p(K^{\bullet\bullet}))$ 및
$\mathrm H_{\mathrm I}^q(\mathrm H_{\mathrm{II}}^p(K^{\bullet\bullet}))$로
나타낸다. 마찬가지로 복합체
$\mathrm H_{\mathrm I}^p(K^{\bullet\bullet})$와 그 코호몰로지 대상
$\mathrm H_{\mathrm{II}}^q(\mathrm H_{\mathrm I}^p(K^{\bullet\bullet}))$를
정의한다. 한편 $\mathrm H^n(K^{\bullet\bullet})$는
$K^{\bullet\bullet}$가 정의하는 (전체) 복합체의 $n$번째 코호몰로지
대상을 나타냄을 상기하자.
\end{env}

\begin{env}[11.3.2]
\label{0.11.3.2-fr}
이중복합체 $K^{\bullet\bullet}$가 정의하는 복합체 위에는 다음과 같이
주어지는 두 표준 여과 $(F_{\mathrm I}^p(K^{\bullet\bullet}))$와
$(F_{\mathrm{II}}^p(K^{\bullet\bullet}))$를 고찰할 수 있다.
\[
\label{0.11.3.2.1-fr}
  F_{\mathrm I}^p(K^{\bullet\bullet})
  =\left(\sum_{\substack{i+j=n\\i\geq p}}K^{ij}\right)_{n\in\mathbb Z},
  \qquad
  F_{\mathrm{II}}^p(K^{\bullet\bullet})
  =\left(\sum_{\substack{i+j=n\\j\geq p}}K^{ij}\right)_{n\in\mathbb Z}.
  \tag{11.3.2.1}
\]
\oldpage[0\textsubscript{III}]{30}
정의에 따라 이들은 실제로 $K^{\bullet\bullet}$가 정의하는 (전체)
복합체의 등급 부분대상들이므로 이 복합체를 여과 복합체로 만든다. 또한 이
여과들이 소진적이고 분리되어 있음은 자명하다.

각 여과에는 스펙트럼 열
(\hyperref[0.11.2.2-fr]{11.2.2})이 대응한다.
$(F_{\mathrm I}^p(K^{\bullet\bullet}))$와
$(F_{\mathrm{II}}^p(K^{\bullet\bullet}))$에 각각 대응하는 스펙트럼 열을
$'E(K^{\bullet\bullet})$와 $''E(K^{\bullet\bullet})$로 나타내며, 이들을
$K^{\bullet\bullet}$의 \emph{스펙트럼 열}이라 한다. 두 열은 모두
코호몰로지 $(\mathrm H^n(K^{\bullet\bullet}))$를 수렴 대상으로 갖는다.
더 나아가 다음이 성립함을 보일 수 있다(M, XV, 6).
\[
\label{0.11.3.2.2-fr}
  'E_2^{pq}(K^{\bullet\bullet})
  =\mathrm H_{\mathrm I}^p(\mathrm H_{\mathrm{II}}^q(K^{\bullet\bullet})),
  \qquad
  ''E_2^{pq}(K^{\bullet\bullet})
  =\mathrm H_{\mathrm{II}}^q(\mathrm H_{\mathrm I}^p(K^{\bullet\bullet})).
  \tag{11.3.2.2}
\]

이중복합체의 모든 사상
$u:K^{\bullet\bullet}\to K^{\prime\bullet\bullet}$은 그 자체로
$K^{\bullet\bullet}$와 $K^{\prime\bullet\bullet}$의 같은 종류의 여과와
양립하므로, 두 스펙트럼 열 각각에 대한 사상을 정의한다. 더 나아가 호모토픽인
두 사상은 대응하는 여과 (전체) 복합체들의 \emph{차수 $\leq1$인
호모토피}를 정의하므로, 두 스펙트럼 열 각각에 대하여 \emph{같은} 사상을
정의한다(M, XV, 6.1).
\end{env}

\begin{proposition}[11.3.3]
\label{0.11.3.3-fr}
$K^{\bullet\bullet}=(K^{ij})$를 아벨 범주 $\mathcal C$의 이중복합체라
하자.
\begin{enumerate}
  \item[(i)] 어떤 $i_0$와 $j_0$가 존재하여 $i<i_0$ 또는 $j<j_0$이면
    $K^{ij}=0$일 때(각각 $i>i_0$ 또는 $j>j_0$이면 $K^{ij}=0$일 때), 두
    스펙트럼 열 $'E(K^{\bullet\bullet})$와
    $''E(K^{\bullet\bullet})$는 쌍정칙이다.
  \item[(ii)] 어떤 $i_0$와 $i_1$이 존재하여 $i<i_0$ 또는 $i>i_1$이면
    $K^{ij}=0$일 때(각각 어떤 $j_0$와 $j_1$이 존재하여 $j<j_0$ 또는
    $j>j_1$이면 $K^{ij}=0$일 때), 두 스펙트럼 열
    $'E(K^{\bullet\bullet})$와 $''E(K^{\bullet\bullet})$는 쌍정칙이다.
\end{enumerate}
더 나아가 $\mathcal C$에서 여과 귀납극한들이 존재하고 완전하다고 하자.
그러면 다음이 성립한다.
\begin{enumerate}
  \item[(iii)] 어떤 $i_0$가 존재하여 $i>i_0$이면 $K^{ij}=0$일
    때(각각 어떤 $j_0$가 존재하여 $j<j_0$이면 $K^{ij}=0$일 때), 열
    $'E(K^{\bullet\bullet})$은 정칙이다.
  \item[(iv)] 어떤 $i_0$가 존재하여 $i<i_0$이면 $K^{ij}=0$일
    때(각각 어떤 $j_0$가 존재하여 $j>j_0$이면 $K^{ij}=0$일 때), 열
    $''E(K^{\bullet\bullet})$은 정칙이다.
\end{enumerate}
\end{proposition}

\begin{proof}
이 명제는 정의 (\hyperref[0.11.1.3-fr]{11.1.3})과
(\hyperref[0.11.2.4-fr]{11.2.4}) 및 여과 $F_{\mathrm I}$에 관한 다음
관찰들에서 곧바로 따른다($K^{\bullet\bullet}$의 두 첨자의 역할을 바꾸면
$F_{\mathrm{II}}$에 관한 유사한 관찰들을 얻는다).
\begin{enumerate}
  \item[$1^\circ$] 어떤 $i_0$가 존재하여 $i>i_0$이면 $K^{ij}=0$일 때,
    여과 $F_{\mathrm I}(K^{\bullet\bullet})$는 이산이다.
  \item[$2^\circ$] 어떤 $i_0$가 존재하여 $i<i_0$이면 $K^{ij}=0$일 때,
    여과 $F_{\mathrm I}(K^{\bullet\bullet})$는 쌍대 이산이다. 이로부터 모든
    $n$에 대하여 대응하는 여과
    $F_{\mathrm I}(\mathrm H^n(K^{\bullet\bullet}))$도 쌍대 이산임을
    곧바로 얻는다. 더 나아가 여과
    $F_{\mathrm I}(K^{\bullet\bullet})$에 대응하는 $B_r^{pq}$의 정의
    (\hyperref[0.11.2.2-fr]{11.2.2})에 의하면 모든 쌍 $(p,q)$에 대하여
    열 $(B_r^{pq})_{r\geq2}$는 정상적이다.
  \item[$3^\circ$] 어떤 $j_0$가 존재하여 $j<j_0$이면 $K^{ij}=0$일 때,
    $p+r+j_0>n$이면
    \[
      F_{\mathrm I}^{p+r}(K^{\bullet\bullet})
      \cap\left(\sum_{i+j=n}K^{ij}\right)=0
    \]
    이다. 따라서 $r>q-j_0+1$이면
    $Z_r^{pq}=Z_\infty(E_2^{pq})$이다. 한편 $p>n-j_0+1$이면
    $\mathrm H^n(F_{\mathrm I}^p(K^{\bullet\bullet}))=0$이다.
  \item[$4^\circ$] 어떤 $j_0$가 존재하여 $j>j_0$이면 $K^{ij}=0$일 때,
    $p-r+1+j_0<n$이면
    \[
      F_{\mathrm I}^{p-r+1}(K^{\bullet\bullet})
      \cap\left(\sum_{i+j=n}K^{ij}\right)=\sum_{i+j=n}K^{ij}.
    \]
    \oldpage[0\textsubscript{III}]{31}
    따라서 $r>j_0-q+1$이면
    $B_r^{pq}=B_\infty(E_2^{pq})$이다. 한편 $p+j_0<n-1$이면
    $\mathrm H^n(F_{\mathrm I}^p(K^{\bullet\bullet}))
    =\mathrm H^n(K^{\bullet\bullet})$이다.
\end{enumerate}
\end{proof}

\begin{env}[11.3.4]
\label{0.11.3.4-fr}
이중복합체 $K^{\bullet\bullet}=(K^{ij})$가 $i<0$ 또는 $j<0$이면
$K^{ij}=0$이라고 하자. 그러면 모든 $p\in\mathbb Z$에 대하여 표준
«~가장자리 준동형~»
\[
\label{0.11.3.4.1-fr}
  'E_2^{p0}(K^{\bullet\bullet})\longrightarrow
  \mathrm H^p(K^{\bullet\bullet})
  \tag{11.3.4.1}
\]
을 정의할 수 있음이 알려져 있다(M, XV, 6). 그 이유를 간단히 상기하자.
한편으로 스펙트럼 열 $'E(K^{\bullet\bullet})$에서
$2\leq r\leq+\infty$일 때
$Z_r^{p0}=Z_{\mathrm I}^p(Z_{\mathrm{II}}^0(K^{\bullet\bullet}))$이고,
다른 한편으로
$\mathrm H^p(F_{\mathrm I}^{p+1}(K^{\bullet\bullet}))=0$이다. 따라서
동형사상
$\beta^{p0}:'E_\infty^{p0}\xrightarrow{\sim}
\mathrm H^p(F_{\mathrm I}^p)/\mathrm H^p(F_{\mathrm I}^{p+1})$은
준동형
$'E_\infty^{p0}\to\mathrm H^p(F_{\mathrm I}^p(K^{\bullet\bullet}))
\to\mathrm H^p(K^{\bullet\bullet})$을 준다. 모든 $Z_r^{p0}$가 같으므로
$r\leq s$에 대하여 표준 준동형
$'E_r^{p0}\to'E_s^{p0}$을 정의할 수 있고, 특히 준동형
$'E_2^{p0}\to'E_\infty^{p0}$을 정의할 수 있다. 이를 합성하여 가장자리
준동형 $'E_2^{p0}\to\mathrm H^p(K^{\bullet\bullet})$을 얻는다. 더 나아가
$d'z=0$을 만족하는 원소
$z\in Z_{\mathrm{II}}^0(K^{\bullet\bullet})\subset K^{p0}$의
법 $B_2^{p0}$인 유에 위에서 정의한 가장자리 준동형을 적용하면,
$'E_\infty^{p0}$에서 $z$의 법 $B_\infty^{p0}$인 유를 얻고, 이어서
$\mathrm H^p(K^{\bullet\bullet})$에서 $z$의 코호몰로지류를 얻음을 곧바로
확인할 수 있다. 그러므로 끝으로 \emph{가장자리 준동형
(\hyperref[0.11.3.4.1-fr]{11.3.4.1})은 표준 포함 사상}
$Z_{\mathrm{II}}^0(K^{\bullet\bullet})\to K^{\bullet\bullet}$
\emph{에 코호몰로지를 취하여 얻어진다}($K^{\bullet\bullet}$는 전체
복합체로 본다). 마찬가지로 가장자리 준동형
\[
\label{0.11.3.4.2-fr}
  ''E_2^{p0}(K^{\bullet\bullet})\longrightarrow
  \mathrm H^p(K^{\bullet\bullet})
  \tag{11.3.4.2}
\]
은 자연스럽게 표준 포함 사상
$Z_{\mathrm I}^0(K^{\bullet\bullet})\to K^{\bullet\bullet}$에서
얻어지는 것으로 해석된다.
\end{env}

\begin{env}[11.3.5]
\label{0.11.3.5-fr}
이제 $K_{\bullet\bullet}=(K_{ij})$를 두 미분 연산자의 차수가 모두
$-1$인 $\mathcal C$의 이중복합체라 하자. 이때 전체 복합체
$(K_{ij})_{j\in\mathbb Z}$ (각각 $(K_{ij})_{i\in\mathbb Z}$)를
$K_{i,\bullet}$ (각각 $K_{\bullet,j}$)로 쓰고, 그 $p$번째 호몰로지 대상을
$\mathrm H_p^{\mathrm{II}}(K_{i,\bullet})$ (각각
$\mathrm H_p^{\mathrm I}(K_{\bullet,j})$)로 쓴다. 복합체
$(\mathrm H_p^{\mathrm{II}}(K_{i,\bullet}))_{i\in\mathbb Z}$ (각각
$(\mathrm H_p^{\mathrm I}(K_{\bullet,j}))_{j\in\mathbb Z}$)은
$\mathrm H_p^{\mathrm{II}}(K_{\bullet\bullet})$ (각각
$\mathrm H_p^{\mathrm I}(K_{\bullet\bullet})$)로 쓰고, 그 $q$번째
호몰로지 대상은
$\mathrm H_q^{\mathrm I}(\mathrm H_p^{\mathrm{II}}(K_{\bullet\bullet}))$
(각각
$\mathrm H_q^{\mathrm{II}}(\mathrm H_p^{\mathrm I}(K_{\bullet\bullet}))$)
로 쓴다. 순환 대상과 경계 대상에도 이와 유사한 표기를 쓴다. 마지막으로
$\mathrm H_n(K_{\bullet\bullet})$는 $K_{\bullet\bullet}$가 정의하는
(차수가 $-1$인 미분 연산자를 갖는) 전체 복합체의 $n$번째 호몰로지 대상이
존재할 때 이를 나타낸다.

$K^{\prime ij}=K_{-i,-j}$로 놓고, $K_{\bullet\bullet}$에 대응하는 차수
$+1$의 미분 연산자들을 갖는 이중복합체를
$K^{\prime\bullet\bullet}=(K^{\prime ij})$라 하자. 정의에 따라
$K_{\bullet\bullet}$의 \emph{스펙트럼 열}은
$K^{\prime\bullet\bullet}$의 스펙트럼 열이며, 이를
$'E(K_{\bullet\bullet})$와 $''E(K_{\bullet\bullet})$로 쓴다. 다만 표기는
다음과 같이 바꾼다.
\[
  'E_{pq}^r(K_{\bullet\bullet})
  ='E_r^{-p,-q}(K^{\prime\bullet\bullet}),\qquad
  ''E_{pq}^r(K_{\bullet\bullet})
  =''E_r^{-p,-q}(K^{\prime\bullet\bullet})
\]
$2\leq r\leq\infty$에 대하여 이 표기를 쓰면 다음을 얻는다.
\[
  'E_{pq}^2(K_{\bullet\bullet})
  =\mathrm H_p^{\mathrm I}(\mathrm H_q^{\mathrm{II}}(K_{\bullet\bullet})),
  \qquad
  ''E_{pq}^2(K_{\bullet\bullet})
  =\mathrm H_q^{\mathrm{II}}(\mathrm H_p^{\mathrm I}(K_{\bullet\bullet})).
\]
부호 오류를 피하려면 일반적으로 이 스펙트럼 열들과 그 수렴 대상 사이의
관계는 복합체 $K^{\prime\bullet\bullet}$로 돌아가 고찰하는 편이 낫다.
다만 (\hyperref[0.11.3.3-fr]{11.3.3})에 대응하는 판정법은 다음과 같다.
\end{env}

\begin{env}[11.3.6]
\label{0.11.3.6-fr}
스펙트럼 열 $'E(K_{\bullet\bullet})$와
$''E(K_{\bullet\bullet})$는 다음 경우에 \emph{쌍정칙}이다.
\begin{enumerate}
  \item[(a)] 어떤 $i_0$와 $j_0$가 존재하여 $i>i_0$ 또는 $j>j_0$이면
    $K_{ij}=0$이다(각각 $i<i_0$
    \oldpage[0\textsubscript{III}]{32}
    또는 $j<j_0$이면 $K_{ij}=0$이다).
  \item[(b)] 어떤 $i_0$와 $i_1$이 존재하여 $i<i_0$이고 $i>i_1$이면
    $K_{ij}=0$이다.
  \item[(c)] 어떤 $j_0$와 $j_1$이 존재하여 $j<j_0$이고 $j>j_1$이면
    $K_{ij}=0$이다.
\end{enumerate}

어떤 $i_0$가 존재하여 $i<i_0$이면 $K_{ij}=0$이거나, 어떤 $j_0$가
존재하여 $j>j_0$이면 $K_{ij}=0$일 때, 열
$'E(K_{\bullet\bullet})$은 정칙이다.

어떤 $i_0$가 존재하여 $i>i_0$이면 $K_{ij}=0$이거나, 어떤 $j_0$가
존재하여 $j<j_0$이면 $K_{ij}=0$일 때, 열
$''E(K_{\bullet\bullet})$은 정칙이다.
\end{env}

\subsection{복합체 $K^\bullet$에 관한 함자의 초코호몰로지}
\label{subsection:0.11.4-fr}

\begin{env}[11.4.1]
\label{0.11.4.1-fr}
$\mathcal C$를 아벨 범주라 하자. $\mathcal C$의 대상 $A$의
\emph{오른쪽 분해}(또는 \emph{코호몰로지 분해})란, 미분작용소의 차수가
$+1$인 $\mathcal C$의 대상들의 복합체
\[
  0\longrightarrow L^0\longrightarrow L^1\longrightarrow L^2
  \longrightarrow\cdots
\]
와, 분해의 \emph{증대}라 불리는 사상 $\varepsilon:A\to L^0$로
이루어지며, 이 증대는 복합체의 사상
\[
\xymatrix{
  0\ar[r] & A\ar[r]\ar[d]_{\varepsilon} & 0\ar[r]\ar[d]
    & 0\ar[r]\ar[d] & \cdots\\
  0\ar[r] & L^0\ar[r] & L^1\ar[r] & L^2\ar[r] & \cdots
})
\]
로 볼 수 있고, 열
\[
  0\longrightarrow A\xrightarrow{\varepsilon}L^0\longrightarrow L^1
  \longrightarrow\cdots
\]
이 완전한 것을 말한다. 마찬가지로 $A$의 \emph{왼쪽 분해}(또는
\emph{호몰로지 분해})란, 미분작용소의 차수가 $-1$인 $\mathcal C$의
대상들의 복합체 $0\leftarrow L_0\leftarrow L_1\leftarrow\cdots$와
증대 $\varepsilon:L_0\to A$로 이루어지며, 열
\[
  0\longleftarrow A\xleftarrow{\varepsilon}L_0\longleftarrow L_1
  \longleftarrow\cdots
\]
이 완전한 것을 말한다.

대상 $A$의 오른쪽 분해 $(L^i)_{i\geq0}$가 $i\geq n+1$일 때
$L^i=0$을 만족하면, 이 분해의 \emph{길이는 $\leq n$}이라고 한다.
길이가 $\leq n$인 왼쪽 분해도 같은 방식으로 정의한다. 어떤 정수
$n$에 대하여 길이가 $\leq n$인 분해를 \emph{유한} 분해라 한다.

$A$의 분해를 이루는 $A$ 이외의 $\mathcal C$의 대상들이
\emph{사영 대상}(각각 \emph{단사 대상})이면 그 분해를
\emph{사영 분해}(각각 \emph{단사 분해})라 한다. $\mathcal C$가 한
환 위의 가군들(예를 들어 왼쪽 가군들)의 범주인 경우, 분해를 이루는
$A$ 이외의 가군들이 \emph{평탄}(각각 \emph{자유})이면 마찬가지로
$A$의 분해를 \emph{평탄 분해}(각각 \emph{자유 분해})라 한다.
\end{env}

\begin{env}[11.4.2]
\label{0.11.4.2-fr}
$K^\bullet=(K^i)_{i\in\mathbb Z}$를 미분작용소의 차수가 $+1$인
$\mathcal C$의 대상들의 복합체라 하자.

$K^\bullet$의 \emph{오른쪽 카르탕--아일렌베르크 분해}란,
미분작용소들의 차수가 $+1$이고 $j<0$일 때 $L^{ij}=0$인 이중복합체
$L^{\bullet\bullet}=(L^{ij})$와 단순복합체의 사상
$\varepsilon:K^\bullet\to L^{\bullet,0}$의 쌍으로서, 다음 조건들을
만족하는 것을 말한다.
\oldpage[0\textsubscript{III}]{33}
\begin{enumerate}
  \item[(i)] 각 지표 $i$에 대하여 열들
    \[
    \begin{aligned}
      0\longrightarrow{}&K^i\xrightarrow{\varepsilon}L^{i0}
        \longrightarrow L^{i1}\longrightarrow\cdots\\
      0\longrightarrow{}&\mathrm B^i(K^\bullet)\xrightarrow{\varepsilon}
        \mathrm B_{\mathrm I}^i(L^{\bullet,0})
        \longrightarrow\mathrm B_{\mathrm I}^i(L^{\bullet,1})
        \longrightarrow\cdots\\
      0\longrightarrow{}&\mathrm Z^i(K^\bullet)\xrightarrow{\varepsilon}
        \mathrm Z_{\mathrm I}^i(L^{\bullet,0})
        \longrightarrow\mathrm Z_{\mathrm I}^i(L^{\bullet,1})
        \longrightarrow\cdots\\
      0\longrightarrow{}&\mathrm H^i(K^\bullet)\xrightarrow{\varepsilon}
        \mathrm H_{\mathrm I}^i(L^{\bullet,0})
        \longrightarrow\mathrm H_{\mathrm I}^i(L^{\bullet,1})
        \longrightarrow\cdots
    \end{aligned}
    \]
    은 완전하다. 다시 말해 $(L^{i,\bullet})$,
    $(\mathrm B_{\mathrm I}^i(L^{\bullet,\bullet}))$,
    $(\mathrm Z_{\mathrm I}^i(L^{\bullet,\bullet}))$,
    $(\mathrm H_{\mathrm I}^i(L^{\bullet,\bullet}))$은 각각 $K^i$,
    $\mathrm B^i(K^\bullet)$, $\mathrm Z^i(K^\bullet)$,
    $\mathrm H^i(K^\bullet)$의 \emph{분해}이다.
  \item[(ii)] 각 $j$에 대하여 단순복합체 $L^{\bullet,j}$는
    \emph{분할}된다. 다시 말해 완전열들
    \[
      \label{0.11.4.2.1-fr}
      0\longrightarrow\mathrm B_{\mathrm I}^i(L^{\bullet,j})
      \longrightarrow\mathrm Z_{\mathrm I}^i(L^{\bullet,j})
      \longrightarrow\mathrm H_{\mathrm I}^i(L^{\bullet,j})
      \longrightarrow0
      \tag{11.4.2.1}
    \]
    \[
      \label{0.11.4.2.2-fr}
      0\longrightarrow\mathrm Z_{\mathrm I}^i(L^{\bullet,j})
      \longrightarrow L^{ij}
      \longrightarrow\mathrm B_{\mathrm I}^{i+1}(L^{\bullet,j})
      \longrightarrow0
      \tag{11.4.2.2}
    \]
    은 분할된다.
\end{enumerate}

$\mathcal C$의 모든 대상이 어떤 단사 대상의 부분대상이면,
$\mathcal C$의 모든 복합체 $K^\bullet$은 \emph{단사} 카르탕--아일렌베르크
분해, 즉 단사 대상들 $L^{ij}$로 이루어진 분해를 가짐이 증명된다
(M, XVII, 1.2). 이때 위의 조건 (ii)에 따라
$\mathrm B_{\mathrm I}^i(L^{\bullet,j})$,
$\mathrm Z_{\mathrm I}^i(L^{\bullet,j})$,
$\mathrm H_{\mathrm I}^i(L^{\bullet,j})$도 단사 대상이다. 더 나아가,
$\mathcal C$의 복합체의 임의의 사상 $f:K^\bullet\to K'^\bullet$,
$K^\bullet$의 임의의 카르탕--아일렌베르크 분해
$L^{\bullet\bullet}$, 그리고 $K'^\bullet$의 임의의 단사
카르탕--아일렌베르크 분해 $L'^{\bullet\bullet}$에 대하여, $f$ 및
증대들과 양립하는 이중복합체의 사상
$F:L^{\bullet\bullet}\to L'^{\bullet\bullet}$가 존재한다. 또한 $f$와
$g$가 $K^\bullet$에서 $K'^\bullet$로 가는 서로 호모토픽한 두 사상이면,
이에 대응하는 $L^{\bullet\bullet}$에서 $L'^{\bullet\bullet}$로 가는
사상들도 서로 호모토픽하다(\emph{같은 곳}).

$K^\bullet$이 \emph{아래로 유계}(각각 \emph{위로 유계})이고 $i<i_0$
(각각 $i>i_0$)일 때 $K^i=0$이면, $i<i_0$ (각각 $i>i_0$)일 때
$L^{ij}=0$이 되도록 $L^{\bullet\bullet}$을 택할 수 있다
(M, XVII, 1.3).

한편 $\mathcal C$의 모든 대상이 길이 $\leq n$인 \emph{단사 분해}를
갖게 하는 정수 $n$이 존재한다고 하자. 그러면 $j>n$일 때
$L^{ij}=0$이라고 가정할 수 있다 (M, XVII, 1.4).
\end{env}

\begin{env}[11.4.3]
\label{0.11.4.3-fr}
이제 $T$를 $\mathcal C$에서 아벨 범주 $\mathcal C'$으로 가는
\emph{가법 공변 함자}라 하자. $\mathcal C$의 복합체 $K^\bullet$과
그 단사 카르탕--아일렌베르크 분해 $L^{\bullet\bullet}$이 주어졌다고
하고, 이중복합체 $T(L^{\bullet\bullet})$이 정의하는 (단순)복합체가
존재한다고 가정하자
(\hyperref[0.11.3.1-fr]{11.3.1} 참조). 이 이중복합체의 두 스펙트럼 열
$'E(T(L^{\bullet\bullet}))$와 $''E(T(L^{\bullet\bullet}))$을
\emph{복합체 $K^\bullet$에 관한 $T$의 초코호몰로지 스펙트럼 열}이라
한다. (\hyperref[0.11.4.2-fr]{11.4.2})와
(\hyperref[0.11.3.2-fr]{11.3.2})에 따라 이들은 실제로 선택한 단사
카르탕--아일렌베르크 분해 $L^{\bullet\bullet}$이 아니라
$K^\bullet$에만 의존하며, 더 나아가 $K^\bullet$에 \emph{함자적으로}
의존한다. 두 열은 공통 수렴대상
$\mathrm H^\bullet(T(L^{\bullet\bullet}))$을 가지며, 이 또한
\emph{복합체 $K^\bullet$에 관한 $T$의 초코호몰로지}라 하고
$\mathrm R^\bullet T(K^\bullet)$로 나타낸다. 앞의 두 스펙트럼 열의
$E_2$항은 다음과 같음이 증명된다.
\[
  \label{0.11.4.3.1-fr}
  'E_2^{pq}=\mathrm H^p(\mathrm R^qT(K^\bullet))
  \tag{11.4.3.1}
\]
\[
  \label{0.11.4.3.2-fr}
  ''E_2^{pq}=\mathrm R^pT(\mathrm H^q(K^\bullet))
  \tag{11.4.3.2}
\]
\oldpage[0\textsubscript{III}]{34}
여기서 $p\in\mathbb Z$에 대하여 $\mathrm R^pT$는 통상과 같이 $T$의
$p$번째 \emph{유도함자}이고, $\mathrm R^qT(K^\bullet)$는 복합체
$(\mathrm R^qT(K^i))_{i\in\mathbb Z}$를 뜻한다.

\phantomsection\label{0.11.4.4-fr}
달리 명시하지 않는 한, 이제부터 $\mathcal C$의 모든 대상이
$\mathcal C$의 어떤 단사 대상의 부분대상이라고 가정한다. 따라서
$\mathcal C$의 모든 복합체에는 단사 카르탕--아일렌베르크 분해가
존재한다. $j<0$일 때 $L^{ij}=0$이므로
(\hyperref[0.11.3.3-fr]{11.3.3})의 판정법에 따라 다음 두 경우에는
$K^\bullet$에 관한 $T$의 초코호몰로지 스펙트럼 열 \emph{둘 다} 존재하며
\emph{쌍정칙}이다.
\begin{enumerate}
  \item[$1^\circ$] $K^\bullet$이 아래로 유계이다.
  \item[$2^\circ$] $\mathcal C$의 모든 대상이, 고려하는 대상과 무관한
    어떤 정수 $n$ 이하의 길이를 갖는 단사 분해를 가진다.
\end{enumerate}
실제로 첫째 경우에는 (\hyperref[0.11.4.2-fr]{11.4.2})에 따라
$i<i_0$일 때 $L^{ij}=0$이 되게 하는 $i_0$가 존재한다고 가정할 수
있고, 둘째 경우에는 $j>j_1$일 때 $L^{ij}=0$이 되게 하는 $j_1$이
존재한다고 가정할 수 있다. 두 경우 모두 주어진 $n$에 대하여
$L^{ij}\neq0$이고 $i+j=n$인 쌍 $(i,j)$는 유한 개뿐임도 명백하므로,
주장이 성립한다.

$\mathcal C'$에서 여과 귀납극한들이 존재하고 완전하다고 가정하면
(특히 $\mathcal C'$에서 무한 직합들이 존재한다), 이중복합체
$T(L^{\bullet\bullet})$이 정의하는 복합체가 존재하고,
(\hyperref[0.11.3.3-fr]{11.3.3})의 판정법에 따라 스펙트럼 열
$'E(T(L^{\bullet\bullet}))$은 항상 \emph{정칙}이다.

$K^\bullet$이 단 하나의 항 $K^{i_0}$을 제외하고 모든 항 $K^i$가
0인 복합체이면, $\mathrm R^nT(K^\bullet)$은
$\mathrm R^{n-i_0}T(K^{i_0})$와 동형이다. 실제로 $i\neq i_0$일 때
$L^{ij}=0$인 카르탕--아일렌베르크 분해 $L^{\bullet\bullet}$을 취하면
정의에서 곧바로 따른다.

$K^\bullet$, $K'^\bullet$이 $\mathcal C$의 두 복합체이고 $f$, $g$가
$K^\bullet$에서 $K'^\bullet$로 가는 서로 \emph{호모토픽한} 두
사상이면, $f$와 $g$로부터 얻는 사상
$\mathrm R^\bullet T(K^\bullet)\to
\mathrm R^\bullet T(K'^\bullet)$은 서로 같으며, 코호몰로지 스펙트럼
열의 사상들도 마찬가지이다.
\end{env}

\begin{proposition}[11.4.5]
\label{0.11.4.5-fr}
$\mathcal C'$에서 여과 귀납극한들이 존재하고 완전하다고 가정하자.
모든 $n>0$과 모든 $i\in\mathbb Z$에 대하여
$\mathrm R^nT(K^i)=0$이면, 함자적 동형
\[
  \label{0.11.4.5.1-fr}
  \mathrm R^iT(K^\bullet)\xrightarrow{\sim}
  \mathrm H^i(T(K^\bullet))\qquad(i\in\mathbb Z)
  \tag{11.4.5.1}
\]
이 성립한다.
\end{proposition}

\begin{proof}
실제로 첫째 스펙트럼 열 (\hyperref[0.11.4.3.1-fr]{11.4.3.1})에서
0이 아닌 $E_2$항은 $'E_2^{p0}=\mathrm H^p(T(K^\bullet))$뿐이다.
다시 말해 이 스펙트럼 열은 \emph{퇴화}한다. 이 열은 정칙이므로
(\hyperref[0.11.4.4-fr]{11.4.4}), 결론은
(\hyperref[0.11.1.6-fr]{11.1.6})에서 따른다.
\end{proof}

\begin{env}[11.4.6]
\label{0.11.4.6-fr}
이제 예를 들어 $\mathcal C\times\mathcal C'$에서 $\mathcal C''$으로
가는 \emph{공변 쌍함자} $(M,N)\mapsto T(M,N)$을 생각하자. 여기서
$\mathcal C$, $\mathcal C'$, $\mathcal C''$은 세 아벨 범주이다. 간단히
하기 위하여 $T$는 각 변수에 관하여 가법이고, $\mathcal C$와
$\mathcal C'$의 모든 대상은 각각 어떤 단사 대상의 부분대상이며,
$\mathcal C''$에서 여과 귀납극한들이 존재하고 완전하다고 가정한다.
그러면 미분작용소의 차수가 $+1$인 $\mathcal C$와 $\mathcal C'$의 두
복합체 $K^\bullet$, $K'^\bullet$에 관한 $T$의 \emph{초코호몰로지}를
다음과 같이 정의한다. $K^\bullet$ (각각 $K'^\bullet$)의 단사
카르탕--아일렌베르크 분해 $L^{\bullet\bullet}$ (각각
$L'^{\bullet\bullet}$)을 취한다. 그러면
$T(L^{\bullet\bullet},L'^{\bullet\bullet})$은 $\mathcal C''$의
사중복합체이며, $T(L^{ij},L'^{hk})$의 차수를 $i+h$와 $j+k$로 취하여
$\mathcal C''$의 \emph{이중복합체}로 본다. \emph{$K^\bullet$과
$K'^\bullet$에 관한 $T$의 초코호몰로지}는 정의에 따라 이
이중복합체의 코호몰로지(즉, 연관된 단순복합체의 코호몰로지)
$\mathrm H^\bullet(T(L^{\bullet\bullet},L'^{\bullet\bullet}))$이며,
$\mathrm R^\bullet T(K^\bullet,K'^\bullet)$로 나타낸다. 이는 두
스펙트럼 열의 수렴대상이고, 그 $E_2$항들은 다음과 같다.
\oldpage[0\textsubscript{III}]{35}
\[
  \label{0.11.4.6.1-fr}
  'E_2^{pq}=\mathrm H^p(\mathrm R^qT(K^\bullet,K'^\bullet))
  \tag{11.4.6.1}
\]
\[
  \label{0.11.4.6.2-fr}
  ''E_2^{pq}=\sum_{q'+q''=q}
  \mathrm R^pT(\mathrm H^{q'}(K^\bullet),\mathrm H^{q''}(K'^\bullet))
  \tag{11.4.6.2}
\]
(M, XVII, 2 참조).

여기서 $\mathrm R^qT(K^\bullet,K'^\bullet)$은 이중복합체
$(\mathrm R^qT(K^i,K'^j))_{(i,j)\in\mathbb Z\times\mathbb Z}$이고,
(\hyperref[0.11.4.6.1-fr]{11.4.6.1})의 우변은 이를 단순복합체로 볼 때의
코호몰로지이다.

더 나아가 첫째 스펙트럼 열은 항상 정칙이다. $\mathcal C$와
$\mathcal C'$의 모든 대상이 길이 $\leq n$인 단사 분해를 갖게 하는
$n$이 존재하거나, $K^\bullet$과 $K'^\bullet$이 아래로 유계이면 두
스펙트럼 열은 쌍정칙이다. 뒤의 경우에는 $\mathcal C''$에서
귀납극한들이 존재한다는 가정도 생략할 수 있다.

$K_1^\bullet$, $K_1'^\bullet$을 각각 $\mathcal C$와 $\mathcal C'$의
다른 두 복합체라 하고, $f$, $g$를 $K^\bullet$에서 $K_1^\bullet$로
가는 서로 호모토픽한 두 사상, $f'$, $g'$를 $K'^\bullet$에서
$K_1'^\bullet$로 가는 서로 호모토픽한 두 사상이라 하자. 그러면
$f$, $f'$에서 얻는 사상과 $g$, $g'$에서 얻는 사상
\[
\mathrm R^\bullet T(K^\bullet,K'^\bullet)\to
\mathrm R^\bullet T(K_1^\bullet,K_1'^\bullet)
\]
은 서로 같고, 초코호몰로지 스펙트럼 열의 사상들도 마찬가지이다.

이는 임의의 가법 공변 다중함자로 쉽게 일반화된다.
\end{env}

\begin{proposition}[11.4.7]
\label{0.11.4.7-fr}
$\mathcal C$의 모든 단사 대상 $I$ (각각 $\mathcal C'$의 단사 대상
$I'$)에 대하여 $A'\mapsto T(I,A')$ (각각 $A\mapsto T(A,I')$)가
완전함자라고 가정하자. 그러면 (\hyperref[0.11.4.6-fr]{11.4.6})의
표기 아래 표준 동형
\[
  \label{0.11.4.7.1-fr}
  \mathrm R^\bullet T(K^\bullet,K'^\bullet)
  \xrightarrow{\sim}\mathrm H^\bullet(T(L^{\bullet\bullet},K'^\bullet))
  \xrightarrow{\sim}\mathrm H^\bullet(T(K^\bullet,L'^{\bullet\bullet}))
  \tag{11.4.7.1}
\]
이 성립한다. 여기서 마지막 두 항은 각각 삼중복합체
$T(L^{\bullet\bullet},K'^\bullet)$과
$T(K^\bullet,L'^{\bullet\bullet})$이 정의하는 단순복합체의
코호몰로지이다.
\end{proposition}

\begin{proof}
예를 들어 이 동형들 중 첫째 것을 정의하자. 사중복합체
$T(L^{\bullet\bullet},L'^{\bullet\bullet})$은
$T(L^{ij},L'^{hk})$의 차수를 $i+j+h$와 $k$로 취하여
\emph{이중복합체}로 볼 수 있다. 각 $h$에 대하여
$L'^{h,\bullet}$은 $K'^h$의 \emph{분해}이므로, $T$에 관한 가정에 따라
이 이중복합체에 대하여
$\mathrm H_{\mathrm{II}}^q(T(L^{\bullet\bullet},L'^{\bullet\bullet}))=0$
($q\neq0$)이고
\[
\mathrm H_{\mathrm{II}}^0(T(L^{\bullet\bullet},L'^{\bullet\bullet}))
=T(L^{\bullet\bullet},K'^\bullet).
\]
따라서 이 이중복합체의 첫째 스펙트럼 열은 \emph{퇴화}한다. 또한
$k<0$일 때 $L'^{hk}=0$이므로 이 열은 정칙이고
(\hyperref[0.11.3.3-fr]{11.3.3}), 결론은
(\hyperref[0.11.1.6-fr]{11.1.6})에서 따른다.
\end{proof}

임의의 수 $n$개의 변수를 갖는 공변 다중함자에 대해서도 비슷한 결과가
성립한다. 초코호몰로지를 계산할 때 모든 복합체를 카르탕--아일렌베르크
분해로 바꿀 필요는 없고, 그중 $n-1$개만 바꾸면 된다. 다만 임의의
$n-1$개 변수를 단사 대상들로 고정했을 때 나머지 변수에 관한 공변
함자는 완전해야 한다.

\subsection{초코호몰로지에서 귀납극한으로의 이행}
\label{subsection:0.11.5-ko}

\begin{lemma}[11.5.1]
\label{0.11.5.1-ko}
$K^\bullet=(K^i)_{i\in\mathbb Z}$를 $\mathcal C$의 복합체라 하고, 각 정수
$r\in\mathbb Z$에 대하여 $K_{(r)}^i=0$ ($i<r$),
$K_{(r)}^i=K^i$ ($i\geq r$)인 복합체를 $K_{(r)}^\bullet$이라 하자.
$T$를 $\mathcal C$에서 $\mathcal C'$로 가는 공변 가법 함자로서
\oldpage[0\textsubscript{III}]{36}
귀납극한과 가환하는 함자라 하자(여과 귀납극한은 $\mathcal C$와
$\mathcal C'$에서 존재하고 완전하다고 가정한다). 그러면 $r$이
$-\infty$로 갈 때 $\mathrm R^\bullet T(K^\bullet)$는 귀납극한
\[
  \varinjlim_{r\to-\infty}\mathrm R^\bullet T(K_{(r)}^\bullet)
\]
과 동형이다.
\end{lemma}

$K^\bullet$의 단사 카르탕--아일렌베르크 분해는 각 $i$에 대하여
$\mathrm B^i(K^\bullet)$의 단사 분해
$(X_{\mathrm B}^{ij})_{j\geq0}$와 $\mathrm H^i(K^\bullet)$의 단사 분해
$(X_{\mathrm H}^{ij})_{j\geq0}$를 임의로 택하여 구성한다. 이렇게 하면
구성법으로부터 $K^i$의 단사 분해 $(L^{ij})_{j\geq0}$와 미분작용소
$L^{i,j}\to L^{i+1,j}$는 \emph{오직 분해들}
$(X_{\mathrm B}^{i,\bullet})$, $(X_{\mathrm H}^{i,\bullet})$ 및
$(X_{\mathrm B}^{i+1,\bullet})$에만 의존함을 알 수 있다
(M, XVII, 1.2). 한편 $i\geq r+1$이면 분명히
$\mathrm B^i(K_{(r)}^\bullet)=\mathrm B^i(K^\bullet)$이고
$\mathrm H^i(K_{(r)}^\bullet)=\mathrm H^i(K^\bullet)$이다. 또 각 $r$에
대하여 표준 단사사상
$\varphi_{r-1,r}^\bullet:K_{(r)}^\bullet\to K_{(r-1)}^\bullet$이 있으며,
$\varphi_{r-1,r}^i$는 $i\neq r-1$일 때 항등사상이다. 앞의 관찰에
따르면 $L^{\bullet\bullet}=(L^{ij})$가 $K^\bullet$의 단사
카르탕--아일렌베르크 분해일 때, 각 $r$에 대하여 $K_{(r)}^\bullet$의
단사 카르탕--아일렌베르크 분해
$L_{(r)}^{\bullet\bullet}=(L_{(r)}^{ij})$를
$i<r$이면 $L_{(r)}^{ij}=0$, $i\geq r+1$이면
$L_{(r)}^{ij}=L^{ij}$이 되도록 정의할 수 있다. 또한
$\varphi_{r-1,r}^\bullet$에 대응하는 이중복합체의 사상
$\Phi_{r-1,r}^{\bullet\bullet}:L_{(r)}^{\bullet\bullet}\to
L_{(r-1)}^{\bullet\bullet}$을 정의할 수 있다. 이 사상의 정의법
(\emph{같은 곳})을 다시 적용하면 $i\neq r$이고 $i\neq r-1$일 때
$\Phi_{r-1,r}^{ij}$가 항등사상이 되도록 구성할 수 있다. 이와 같이
$\mathcal C$의 이중복합체들로 이루어진 귀납계
$(L_{(r)}^{\bullet\bullet})$를 정의하였다. $r$이 $-\infty$로 갈 때
그 귀납극한은 분명히 $L^{\bullet\bullet}$이다. 직합과 귀납극한이
가환하므로, $L^{\bullet\bullet}$에 결부된 단순복합체도
$L_{(r)}^{\bullet\bullet}$에 결부된 단순복합체들의 귀납극한이다.
$T$는 가정에 따라 귀납극한과 가환하고, 코호몰로지도 그러하므로
(함자 $\varinjlim$의 완전성에 의하여), 동형을 제외하면
\[
  \mathrm H^\bullet(T(L^{\bullet\bullet}))=
  \varinjlim\mathrm H^\bullet(T(L_{(r)}^{\bullet\bullet}))
\]
이다.

보조정리 (\hyperref[0.11.5.1-ko]{11.5.1})을 쓰면 아래로 유계인
복합체들에 대하여 성립하는 성질을 귀납극한으로 이행시켜 임의의
복합체 $K^\bullet$로 확장할 수 있다. 첫 번째 예로 다음을 증명한다.

\begin{proposition}[11.5.2]
\label{0.11.5.2-ko}
$\mathcal C$, $\mathcal C'$ 및 $T$에 관한
(\hyperref[0.11.5.1-ko]{11.5.1})의 가정 아래에서
$\mathrm R^\bullet T(K^\bullet)$는 $\mathcal C$의 복합체들의 아벨
범주에서 코호몰로지 함자이다.
\end{proposition}

\begin{proof}
아래로 유계인 복합체들의 경우로 환원할 수 있음을 보이자. 복합체들의
완전열
\[
  0\longrightarrow K'^\bullet\longrightarrow K^\bullet
  \longrightarrow K''^\bullet\longrightarrow0
\]
이 주어지면, 각 $r$에 대하여 분명히 완전열
\[
  0\longrightarrow K_{(r)}'^\bullet\longrightarrow K_{(r)}^\bullet
  \longrightarrow K_{(r)}''^\bullet\longrightarrow0
\]
을 얻으며, 따라서 가정에 의하여 완전열
\[
  \cdots\longrightarrow\mathrm R^nT(K_{(r)}'^\bullet)
  \longrightarrow\mathrm R^nT(K_{(r)}^\bullet)
  \longrightarrow\mathrm R^nT(K_{(r)}''^\bullet)
  \xrightarrow{\partial}\mathrm R^{n+1}T(K_{(r)}'^\bullet)
  \longrightarrow\cdots
\]
을 얻는다. 이 완전열들은 귀납계를 이룬다. 보조정리
(\hyperref[0.11.5.1-ko]{11.5.1})과 함자 $\varinjlim$의 완전성으로부터
완전열
\[
  \cdots\longrightarrow\mathrm R^nT(K'^\bullet)
  \longrightarrow\mathrm R^nT(K^\bullet)
  \longrightarrow\mathrm R^nT(K''^\bullet)
  \xrightarrow{\partial}\mathrm R^{n+1}T(K'^\bullet)
  \longrightarrow\cdots
\]
을 얻는다.
\end{proof}

아래로 유계인 복합체들을 다룰 때에는 $i<0$이면 $K^i=0$인 것들만
고려해도 충분하다. 이들은 분명히 아벨 범주 $\mathcal K$를 이룬다.

\begin{lemma}[11.5.2.1]
\label{0.11.5.2.1-ko}
$\mathcal K$에서 다음 성질들을 갖는 복합체
$Q^\bullet=(Q^i)_{i\geq0}$의 집합을 $\mathfrak J$라 하자.
\oldpage[0\textsubscript{III}]{37}
\begin{enumerate}
  \item[$1^\circ$] 각 $Q^i$는 $\mathcal C$의 단사 대상이다.
  \item[$2^\circ$] 각 $i\geq0$에 대하여
    $\mathrm Z^i(Q^\bullet)=\mathrm B^i(Q^\bullet)$이고,
    $\mathrm Z^i(Q^\bullet)$는 $Q^i$의 직합인자이다.
\end{enumerate}
그러면 다음이 성립한다.
\begin{enumerate}
  \item[(i)] 모든 $Q^\bullet\in\mathfrak J$는 $\mathcal K$의 단사
    대상이다.
  \item[(ii)] $\mathcal K$의 모든 대상은 $\mathfrak J$에 속하는 어떤
    복합체의 부분복합체와 동형이다.
\end{enumerate}
\end{lemma}

\begin{proof}
\textup{(i)} $A^\bullet=(A^i)$를 $\mathcal K$의 대상,
$A'^\bullet=(A'^i)$를 $A^\bullet$의 부분대상,
$Q^\bullet=(Q^i)$를 $\mathfrak J$의 대상이라 하자. 사상
$f=(f^i):A'^\bullet\to Q^\bullet$이 주어졌다고 하고, 이를 사상
$g=(g^i):A^\bullet\to Q^\bullet$으로 연장하려 한다. 표기를
간단히 하기 위하여 가군 범주의 언어를 쓰겠다([27] 참조).

$Q^i$를
$\mathrm B^i(Q^\bullet)\oplus\mathrm B^{i+1}(Q^\bullet)$와
동일시하자. $i$에 관한 귀납법을 쓰며, 따라서 $j<i$에 대하여
$g^j$가 정의되어 있고 $j<i-1$에 대하여 미분작용소
$d^j:A^j\to A^{j+1}$ 및 $d'^j:Q^j\to Q^{j+1}$과 양립하며, 더 나아가
다음을 만족한다고 가정하자.
\begin{enumerate}
  \item[$1^\circ$] $g^{i-1}(\mathrm Z^{i-1}(A^\bullet))\subset
    \mathrm Z^{i-1}(Q^\bullet)$.
  \item[$2^\circ$] 모든 $j$에 대하여
    $C^j=(d^j)^{-1}(A'^{j+1})$라 놓으면,
    $d'^{i-1}\circ g^{i-1}$과 $f^i\circ d^{i-1}$은 $C^{i-1}$ 위에서
    일치한다.
\end{enumerate}
사상 $f^i:A'^i\to Q^i$를 두 사영과 합성하면 두 사상
$f'^i:A'^i\to\mathrm B^i(Q^\bullet)$ 및
$f''^i:A'^i\to\mathrm B^{i+1}(Q^\bullet)$을 얻는다.
$d'^{i-1}\circ g^{i-1}$은 $A^{i-1}$을 $\mathrm B^i(Q^\bullet)$로
보내고 $\mathrm Z^{i-1}(A^\bullet)$ 위에서 0이므로 사상
$h^i:\mathrm B^i(A^\bullet)\to\mathrm B^i(Q^\bullet)$을 정의한다.
또 $d'^{i-1}\circ g^{i-1}$과 $f^i\circ d^{i-1}$은 $C^{i-1}$ 위에서
일치하므로, $h^i$와 $f'^i$는
$\mathrm B^i(A^\bullet)\cap A'^i$ 위에서 일치한다.
$\mathrm B^i(Q^\bullet)$은 $Q^i$의 직합인자이므로 단사 대상이다.
따라서 $\mathrm B^i(A^\bullet)$ 위에서는 $h^i$와, $A'^i$ 위에서는
$f'^i$와 일치하는 사상
$g'^i:A^i\to\mathrm B^i(Q^\bullet)$가 존재한다. 한편 사상
$f'^{i+1}\circ d^i:C^i\to\mathrm B^{i+1}(Q^\bullet)$을 생각하자.
이는 $\mathrm Z^i(A^\bullet)$ 위에서 0이다. 그런데
$\mathrm B^{i+1}(Q^\bullet)$은 단사 대상이므로, $C^i$ 위에서는
$f'^{i+1}\circ d^i$와 일치하고 $\mathrm Z^i(A^\bullet)$ 위에서는
$0$과 일치하는 사상
$g''^i:A^i\to\mathrm B^{i+1}(Q^\bullet)$가 존재한다. 이제
$g^i=g'^i+g''^i$로 놓으면 귀납을 계속할 수 있다.

\textup{(ii)} $A^\bullet=(A^i)$를 $\mathfrak J$에 속하는 복합체에
매입하기 위하여, 각 $i\geq1$에 대하여 단사사상
$f'^i:A^i\to Q'^i$가 존재하도록 $\mathcal C$의 단사 대상 $Q'^i$를
택한다. 이제
\[
  Q^i=0\quad(i<0),\qquad Q^0=Q'^1,\qquad
  Q^i=Q'^i\oplus Q'^{i+1}\quad(i\geq1)
\]
로 놓고, 자명한 미분작용소를 준다. 모든 $i$에 대하여
$f^i=f'^i+(f'^{i+1}\circ d^i)$로 놓자($i\leq0$이면
$f'^i=0$으로 한다). 각 $f^i$가 단사이고 $f^i$들이 미분작용소와
양립함은 곧바로 알 수 있다.
\end{proof}

\begin{corollary}[11.5.2.2]
\label{0.11.5.2.2-ko}
$\mathcal K$의 모든 대상 $K^\bullet$은 $\mathfrak J$의 대상들로
이루어진 오른쪽 분해를 갖는다. $L^{\bullet\bullet}$이 그러한
분해이고, $L'^{\bullet\bullet}$이 $\mathcal K$의 대상들로 이루어진
$K^\bullet$의 임의의 분해이면, 확대사상들과 양립하는 이중복합체의
사상 $F:L'^{\bullet\bullet}\to L^{\bullet\bullet}$이 존재하고, 그러한
두 사상 $F$, $F'$은 서로 호모토픽이다.
\end{corollary}

\begin{proof}
이는 아벨 범주 $\mathcal K$에 (M, V, 1.1 a))를 적용한 것에 지나지
않는다.
\end{proof}

\begin{env}[11.5.2.3]
\label{0.11.5.2.3-ko}
이 준비들을 갖추었으므로, $K^\bullet$의 단사 카르탕--아일렌베르크
분해 $L'^{\bullet\bullet}$과 $\mathfrak J$의 대상들로 이루어진
$K^\bullet$의 분해 $L^{\bullet\bullet}$을 생각하고, 다음 동형이
성립함을 보이자.
\[
  \mathrm H^\bullet(T(L'^{\bullet\bullet}))\xrightarrow{\sim}
  \mathrm H^\bullet(T(L^{\bullet\bullet})).
\]
실제로 (\hyperref[0.11.5.2.2-ko]{11.5.2.2})로부터 이중복합체의 사상
$T(L'^{\bullet\bullet})\to T(L^{\bullet\bullet})$을 얻고, 따라서 이
이중복합체들의 첫째 분광수열 사이의 사상
$'E(T(L'^{\bullet\bullet}))\to{}'E(T(L^{\bullet\bullet}))$을 얻는다.
(\hyperref[0.11.3.3-fr]{11.3.3})에 의하여 이 분광수열들은 정칙이므로,
앞의 사상이 $E_2$ 항들에서 동형임을 보이면 충분하다
(\hyperref[0.11.1.5-fr]{11.1.5}). 다시 말해
$\mathrm H_{\mathrm{II}}^q(T(L^{i,\bullet}))$가
$\mathrm R^qT(K^i)$와 같음을 보이면 충분하다. 그런데
$L^{i,\bullet}$은 $K^i$의 오른쪽 분해이므로 다음 보조정리를
증명하는 문제로 환원된다.
\end{env}

\oldpage[0\textsubscript{III}]{38}
\begin{lemma}[11.5.2.4]
\label{0.11.5.2.4-ko}
$L^\bullet=(L^i)_{i\in\mathbb Z}$가 $\mathcal C$의 대상 $A$의 오른쪽
분해이고, 모든 $i\in\mathbb Z$와 모든 $n>0$에 대하여
$\mathrm R^nT(L^i)=0$이라 하자. 그러면
$\mathrm H^\bullet T(L^\bullet)=\mathrm R^\bullet T(A)$이다.
\end{lemma}

\begin{proof}
이는 (T, 2.5.2)의 특수한 경우이다.
\end{proof}

\begin{env}[11.5.2.5]
\label{0.11.5.2.5-ko}
이제 (\hyperref[0.11.5.2-ko]{11.5.2})의 증명은 곧바로 끝난다.
실제로 $L'^{\bullet\bullet}$은 아벨 범주 $\mathcal K$에서
$K^\bullet$의 단사 분해이다. 다시 말해
$K^\bullet\mapsto\mathrm H^\bullet(T(L'^{\bullet\bullet}))$은 바로
범주 $\mathcal K$에서 $T$의 오른쪽 유도 코호몰로지 함자이다
(T, 2.3).
\end{env}

\begin{proposition}[11.5.3]
\label{0.11.5.3-ko}
$\mathcal C$, $\mathcal C'$ 및 $T$에 관한
(\hyperref[0.11.5.1-ko]{11.5.1})의 가정 아래에서,
$L^{\bullet\bullet}=(L^{ij})$를 $j<0$이면 $L^{ij}=0$이고 각 $i$에
대하여 $L^{i,\bullet}$이 $K^i$의 분해인 이중복합체라 하자. 끝으로
모든 쌍 $(i,j)$와 모든 $n>0$에 대하여
$\mathrm R^nT(L^{ij})=0$이라고 가정하자. 그러면 함자적 동형
\[
  \label{0.11.5.3.1-ko}
  \mathrm R^\bullet T(K^\bullet)\xrightarrow{\sim}
  \mathrm H^\bullet(T(L^{\bullet\bullet})).
  \tag{11.5.3.1}
\]
이 존재한다.
\end{proposition}

\begin{proof}
$i<r$이면 $L_{(r)}^{ij}=0$, $i\geq r$이면
$L_{(r)}^{ij}=L^{ij}$인 이중복합체를
$L_{(r)}^{\bullet\bullet}=(L_{(r)}^{ij})$라 하자. $r$이
$-\infty$로 갈 때 $L^{\bullet\bullet}$이
$L_{(r)}^{\bullet\bullet}$의 귀납극한임은 곧바로 알 수 있다.
$T$에 대한 가정과 (\hyperref[0.11.5.1-ko]{11.5.1})에 의하여, 따라서
$K^\bullet$이 아래로 유계인 경우만 증명하면 충분하다. 예를 들어
$i<0$이면 $K^i=0$이고 $L^{ij}=0$인 경우만 다루면 된다. 이제
$L'^{\bullet\bullet}=(L'^{ij})$을 $\mathfrak J$의 대상들로 이루어진
$K^\bullet$의 오른쪽 분해라 하자
(\hyperref[0.11.5.2.2-ko]{11.5.2.2}). 확대사상들과 양립하는
이중복합체의 사상 $L^{\bullet\bullet}\to L'^{\bullet\bullet}$이
존재하므로, 첫째 분광수열 사이의 사상
$'E(T(L^{\bullet\bullet}))\to{}'E(T(L'^{\bullet\bullet}))$을 얻는다.
보조정리 (\hyperref[0.11.5.2.4-ko]{11.5.2.4})에 의하여,
(\hyperref[0.11.5.2.3-ko]{11.5.2.3})에서와 같이 이 사상은
\emph{동형}이고, 이로부터 결론이 따른다.
\end{proof}

\begin{remark}[11.5.4]
\label{0.11.5.4-ko}
앞의 논증은 아래로 유계인 복합체 $K^\bullet$의 범주에서는
\emph{$T$가 여과 귀납극한과 가환한다고 가정하지 않아도}
(\hyperref[0.11.5.2-ko]{11.5.2})와
(\hyperref[0.11.5.3-ko]{11.5.3})의 결론들이 성립함을 증명한다.
더욱이 $i<0$이면 $K^i=0$인 복합체 $K^\bullet$들의 범주
$\mathcal K$만을 생각할 때에는, $\mathrm R^\bullet T(K^\bullet)$가
$\mathcal K$에서 $T$의 \emph{오른쪽 유도 함자}들의 계라는 특징으로부터
이 코호몰로지 함자가 \emph{보편적}임을 알 수 있다(T, 2.3).

$T$에 관하여 추가 가정을 하지 않고도
(\hyperref[0.11.5.2-ko]{11.5.2})가 성립하는 또 하나의 경우는,
$\mathcal C$의 모든 대상이 길이 $\leq m$인 단사 분해를 갖게 하는
정수 $m>0$이 존재하는 경우이다. 실제로
(\hyperref[0.11.5.1-ko]{11.5.1})의 증명에 나타나는
$\mathcal C$의 대상들의 단사 분해를 모두 길이 $\leq m$인 것으로
택할 수 있다. 따라서 $-r$이 충분히 크면 이중복합체
$T(L_{(r)}^{\bullet\bullet})$의 전체차수 $n$인 항들은
$T(L^{\bullet\bullet})$의 항들과 같고 그 수가 유한하다. 이에 따라
각 $n$에 대하여 $-r$이 충분히 크면
$\mathrm H^n(T(L^{\bullet\bullet}))=
\mathrm H^n(T(L_{(r)}^{\bullet\bullet}))$이다.
(\hyperref[0.11.5.2-ko]{11.5.2})의 표기를 쓰면, 따라서 $-r$이 충분히
클 때($n$에 의존한다)
$\mathrm R^nT(K_{(r)}^\bullet)=\mathrm R^nT(K^\bullet)$이고,
$K'^\bullet$과 $K''^\bullet$에 대해서도 마찬가지이므로 결론이 따른다.
같은 방법으로, $\mathcal C$가 앞의 가정을 만족하고 분해들
$L^{i,\bullet}$의 길이가 $\leq m$이라고 가정하면
(\hyperref[0.11.5.3-ko]{11.5.3})도 $T$에 대한 추가 조건 없이
성립한다.
\end{remark}

\begin{env}[11.5.5]
\label{0.11.5.5-ko}
(\hyperref[0.11.5.2-ko]{11.5.2})의 결과는 공변 다변함자로
일반화된다. 예를 들어 (\hyperref[0.11.4.6-fr]{11.4.6})의 상황을
생각하자. 여기서는 $\mathcal C$, $\mathcal C'$ 및 $\mathcal C''$에서
\oldpage[0\textsubscript{III}]{39}
여과 귀납극한이 존재하고 완전하며, $T$가 귀납극한과 가환한다고
가정한다. 그러면 $\mathrm R^\bullet T(K^\bullet,K'^\bullet)$은
각 복합체 $K^\bullet$, $K'^\bullet$에 관하여 코호몰로지 쌍함자이다.
이를 보이기 위하여 (\hyperref[0.11.5.2-ko]{11.5.2})에서와 같이
$K^\bullet$과 $K'^\bullet$이 아래로 유계인 경우로 환원한다. 그런 다음
$K^\bullet$과 $K'^\bullet$에 대하여
(\hyperref[0.11.5.2.2-ko]{11.5.2.2})에 기술된 유형의 단사 분해를
택하면, (M, V, 4.1)에서 증명한 일반 성질로 환원된다.
\end{env}

\begin{env}[11.5.6]
\label{0.11.5.6-ko}
마찬가지로 (\hyperref[0.11.4.7-fr]{11.4.7})과
(\hyperref[0.11.5.3-ko]{11.5.3})의 결과들은 다음과 같이 일반화된다.
(\hyperref[0.11.5.5-ko]{11.5.5})의 가정 아래에서 두 이중복합체
$L^{\bullet\bullet}=(L^{ij})$,
$L'^{\bullet\bullet}=(L'^{ij})$가 주어졌다고 하자. $j<0$이면
$L^{ij}=0$이고 $L'^{ij}=0$이며, 각 $i$에 대하여
$L^{i,\bullet}$은 $K^i$의 분해이고 $L'^{i,\bullet}$은 $K'^i$의
분해라고 하자. 끝으로 모든 지표계 $(i,j,h,k)$와 $n>0$에 대하여
$\mathrm R^nT(L^{ij},L'^{hk})=0$이라고 하자. 그러면
$K^\bullet$과 $K'^\bullet$에 관한 함자적 동형
\[
  \label{0.11.5.6.1-ko}
  \mathrm R^\bullet T(K^\bullet,K'^\bullet)\xrightarrow{\sim}
  \mathrm H^\bullet(T(L^{\bullet\bullet},L'^{\bullet\bullet})).
  \tag{11.5.6.1}
\]
이 성립한다. 이는 (\hyperref[0.11.5.3-ko]{11.5.3})에서와 같이
$K^\bullet$과 $K'^\bullet$이 아래로 유계인 경우로 환원하여 증명한다.

더 나아가 각 쌍 $(i,j)$와 각 쌍 $(h,k)$에 대하여 함자
$A\mapsto T(A,L'^{hk})$와 $A'\mapsto T(L^{ij},A')$가 각각
$\mathcal C$와 $\mathcal C'$에서 완전하다고 가정하자. 그러면 또한
함자적 동형
\[
  \label{0.11.5.6.2-ko}
  \mathrm R^\bullet T(K^\bullet,K'^\bullet)
  \xrightarrow{\sim}\mathrm H^\bullet(T(L^{\bullet\bullet},K'^\bullet))
  \xrightarrow{\sim}\mathrm H^\bullet(T(K^\bullet,L'^{\bullet\bullet})).
  \tag{11.5.6.2}
\]
이 성립한다. 증명은 (\hyperref[0.11.4.7-fr]{11.4.7})의 증명과
비슷하다.
\end{env}

\begin{env}[11.5.7]
\label{0.11.5.7-ko}
끝으로 (\hyperref[0.11.5.5-ko]{11.5.5})와
(\hyperref[0.11.5.6-ko]{11.5.6})의 결과들은 다음 두 경우에 $T$가
귀납극한과 가환한다고 가정하지 않아도 성립함에 유의하자. 하나는
아래로 유계인 복합체 $K^\bullet$, $K'^\bullet$에 한정하는 경우이고,
다른 하나는 $\mathcal C$(각각 $\mathcal C'$)의 모든 대상이 길이가
유계인 단사 분해를 갖는다고 가정하고
(\hyperref[0.11.5.6-ko]{11.5.6})에서 이중복합체
$L^{\bullet\bullet}$과 $L'^{\bullet\bullet}$의 둘째 차수가 위로
유계라고 가정하는 경우이다.
\end{env}

\subsection{복합체 $K_\bullet$에 관한 함자의 초호몰로지}
\label{subsection:0.11.6-fr}

\begin{env}[11.6.1]
\label{0.11.6.1-fr}
$\mathcal C$를 아벨 범주,
$K_\bullet=(K_i)_{i\in\mathbb Z}$를 미분작용소의 차수가 $-1$인
$\mathcal C$의 대상들의 복합체라 하자. $K_\bullet$의
\emph{왼쪽 카르탕--아일렌베르크 분해}란, 미분작용소들의 차수가
$-1$이고 $j<0$이면 $L_{ij}=0$인 이중복합체
$L_{\bullet\bullet}=(L_{ij})$와 단순복합체들의 사상
\[
\varepsilon:L_{\bullet,0}\to K_\bullet
\]
로 이루어지며, (\hyperref[0.11.4.2-fr]{11.4.2})의 조건들에서
``화살표의 방향을 뒤집어서'' 얻은 조건들을 만족하는 것이다.
$\mathcal C$의 모든 대상이 어떤 사영 대상의 몫이면,
$\mathcal C$의 모든 복합체 $K_\bullet$은
\emph{사영 카르탕--아일렌베르크 분해}, 즉 사영 대상들 $L_{ij}$로
이루어진 분해를 가지며, 이는
(\hyperref[0.11.4.2-fr]{11.4.2})와 비슷한 함자적 성질들을 갖는다.
또한 $K_\bullet$이 아래로 유계이면 (각각 위로 유계이면),
$i<i_0$일 때 $K_i=0$인 경우 (각각 $i>i_0$일 때 $K_i=0$인 경우)
$i<i_0$이면 $L_{ij}=0$이라고 (각각 $i>i_0$이면
$L_{ij}=0$이라고) 가정할 수 있다. $\mathcal C$의 모든 대상이
길이 $\leq n$인 사영 분해를 가지면 $j>n$일 때 $L_{ij}=0$이라고
가정할 수 있다.
\end{env}

\begin{env}[11.6.2]
\label{0.11.6.2-fr}
\oldpage[0\textsubscript{III}]{40}
$T$를 $\mathcal C$에서 아벨 범주 $\mathcal C'$으로 가는 공변 가법
함자라고 가정하자. $\mathcal C$의 복합체 $K_\bullet$에 관한
$T$의 \emph{초호몰로지} $L_\bullet T(K_\bullet)$와
\emph{초호몰로지 스펙트럼 열}들이 존재할 때의 정의는 여전히
(\hyperref[0.11.4.3-fr]{11.4.3})에서 ``화살표의 방향을
뒤집어서'' 얻는다. 이렇게 얻은 두 스펙트럼 열의 $E_2$ 항은
\[
\label{0.11.6.2.1-fr}
\tag{11.6.2.1}
{}'E^2_{pq}=H_p(L_qT(K_\bullet))
\]
및
\[
\label{0.11.6.2.2-fr}
\tag{11.6.2.2}
{}''E^2_{pq}=L_pT(H_q(K_\bullet))
\]
이다. 여기서 $L_pT$는 $p\geq0$이면 $T$의 $p$번째 왼쪽 유도
함자이고 $p<0$이면 $0$이며, $L_qT(K_\bullet)$는 복합체
$(L_qT(K_i))_{i\in\mathbb Z}$이다.

초호몰로지의 성질들이 모두 초코호몰로지의 성질들에서 단순히
``화살표의 방향을 뒤집어서'' 따라오는 것은 아니다. 범주
$\mathcal C'$에 (T, 1.5)의 (AB~5*) 유형의 추가 가정을 두지 않으면,
앞의 두 스펙트럼 열의 정칙성 조건 때문에 이번에는
(\hyperref[0.11.3.4-fr]{11.3.4})의 판정법들을 적용해야 하기
때문이다. 이 판정법들에 의하면 $\mathcal C'$에서 여과 귀납극한들이
존재하고 완전하다고 가정할 때, 이중복합체
$T(L_{\bullet\bullet})$가 정하는 복합체가 존재하고 둘째 스펙트럼 열
${}''E(T(L_{\bullet\bullet}))$은 정칙이다. $K_\bullet$이 아래로
유계라고 가정하거나, $\mathcal C$의 모든 대상이 길이 $\leq n$인
사영 분해를 갖게 하는 어떤 정수 $n$이 존재한다고 가정하면, 두
초호몰로지 스펙트럼 열은 $\mathcal C'$에 대한 아무 가정 없이도
존재하며 쌍정칙이다.
\end{env}

\begin{proposition}[11.6.3]
\label{0.11.6.3-fr}
$\mathcal C$, $\mathcal C'$을 두 아벨 범주, $T$를 $\mathcal C$에서
$\mathcal C'$으로 가는 공변 가법 함자라 하자. 그러면 다음이
성립한다.
\begin{enumerate}
\item[\rm(i)]
\emph{초호몰로지 $L_\bullet T(K_\bullet)$은 $\mathcal C$의 아래로
유계인 복합체들의 아벨 범주에서 호몰로지 함자이다.}
\item[\rm(ii)]
\emph{$K_\bullet$을 $\mathcal C$의 아래로 유계인 복합체라 하자.
모든 $n>0$와 모든 $i\in\mathbb Z$에 대해 $L_nT(K_i)=0$이면
함자적 동형}
\[
\label{0.11.6.3.1-fr}
\tag{11.6.3.1}
L_iT(K_\bullet)\simeq H_i(T(K_\bullet))\qquad(i\in\mathbb Z)
\]
\emph{이 성립한다.}
\item[\rm(iii)]
\emph{$K_\bullet$을 $\mathcal C$의 아래로 유계인 복합체라 하자.
$L_{\bullet\bullet}=(L_{ij})$를 $j<0$이면 $L_{ij}=0$이고, 모든
$i$에 대해 $L_{i,\bullet}$이 $K_i$의 분해인 이중복합체라 하자.
끝으로 모든 쌍 $(i,j)$와 모든 $n>0$에 대해
$L_nT(L_{ij})=0$이라고 가정하자. 그러면 함자적 동형}
\[
\label{0.11.6.3.2-fr}
\tag{11.6.3.2}
L_\bullet T(K_\bullet)\simeq H_\bullet(T(L_{\bullet\bullet}))
\]
\emph{이 존재한다.}
\end{enumerate}
아래로 유계인 복합체의 경우 증명은 각각
(\hyperref[0.11.5.2-ko]{11.5.2}),
(\hyperref[0.11.4.5-fr]{11.4.5}),
(\hyperref[0.11.5.3-ko]{11.5.3})의 증명과 같은 방식으로 한다.
이 논증들의 세부사항은 독자에게 맡긴다.
\end{proposition}

\begin{env}[11.6.4]
\label{0.11.6.4-fr}
각 변수에 관해 공변이고 가법인 다변함자에 대해서도 완전히 비슷한
결과들이 성립한다. 예를 들어 쌍함자 $T$에 대한 두 초호몰로지
스펙트럼 열의 $E_2$ 항은 다음과 같다.
\oldpage[0\textsubscript{III}]{41}
\[
\label{0.11.6.4.1-fr}
\tag{11.6.4.1}
{}'E^2_{pq}=H_p(L_qT(K_\bullet,K'_\bullet))
\]
\[
\label{0.11.6.4.2-fr}
\tag{11.6.4.2}
{}''E^2_{pq}=
\sum_{q'+q''=q}L_pT(H_{q'}(K_\bullet),H_{q''}(K'_\bullet)).
\]
여기에서도 둘째 스펙트럼 열이 정칙이다. 복합체 $K_\bullet$과
$K'_\bullet$이 아래로 유계이거나, 고려하는 아벨 범주들의 대상들이
고정된 길이의 사영 분해를 가질 때에는 두 스펙트럼 열이 쌍정칙이다.

또한 다음이 성립한다.
\end{env}

\begin{proposition}[11.6.5]
\label{0.11.6.5-fr}
$\mathcal C$, $\mathcal C'$, $\mathcal C''$을 세 아벨 범주,
$T$를 $\mathcal C\times\mathcal C'$에서 $\mathcal C''$으로 가는
공변 쌍가법 쌍함자라 하자.
\begin{enumerate}
\item[\rm(i)]
\emph{$L_\bullet T(K_\bullet,K'_\bullet)$은 아래로 유계인 각 복합체
$K_\bullet$, $K'_\bullet$에 관해 호몰로지 쌍함자이다. 여기서 이
복합체들은 각각 $\mathcal C$의 대상들과 $\mathcal C'$의 대상들로
이루어진다.}
\item[\rm(ii)]
\emph{$K_\bullet$과 $K'_\bullet$이 아래로 유계라고 가정하자.
$L_{\bullet\bullet}=(L_{ij})$,
$L'_{\bullet\bullet}=(L'_{ij})$를 $j<0$이면
$L_{ij}=L'_{ij}=0$이고, 모든 $i$에 대해 $L_{i,\bullet}$은
$K_i$의 분해이며 $L'_{i,\bullet}$은 $K'_i$의 분해인 두
이중복합체라 하자. 끝으로 모든 $n>0$와 모든 지표계
$(i,j,h,k)$에 대해 $L_nT(L_{ij},L'_{hk})=0$이라고 하자. 그러면
함자적 동형}
\[
\label{0.11.6.5.1-fr}
\tag{11.6.5.1}
L_\bullet T(K_\bullet,K'_\bullet)
\simeq H_\bullet(T(L_{\bullet\bullet},L'_{\bullet\bullet}))
\]
\emph{이 성립한다.}
\item[\rm(iii)]
\emph{더 나아가 모든 쌍 $(i,j)$와 모든 쌍 $(h,k)$에 대해 함자}
\[
A\longmapsto T(A,L'_{hk}),\qquad
A'\longmapsto T(L_{ij},A')
\]
\emph{가 각각 $\mathcal C$와 $\mathcal C'$에서 완전하다고 하자.
그러면 함자적 동형}
\[
\label{0.11.6.5.2-fr}
\tag{11.6.5.2}
L_\bullet T(K_\bullet,K'_\bullet)
\simeq H_\bullet(T(L_{\bullet\bullet},K'_\bullet))
\simeq H_\bullet(T(K_\bullet,L'_{\bullet\bullet}))
\]
\emph{들이 성립한다.}
\end{enumerate}
증명은 (\hyperref[0.11.5.5-ko]{11.5.5})와
(\hyperref[0.11.5.6-ko]{11.5.6})의 증명과 비슷하다.
\end{proposition}

\subsection{이중복합체 $K_{\bullet\bullet}$에 관한 함자의 초호몰로지}
\label{subsection:0.11.7-fr}

\begin{env}[11.7.1]
\label{0.11.7.1-fr}
$\mathcal C$를 모든 대상이 어떤 사영 대상의 몫인 아벨 범주라 하자.
$\mathcal C$의 대상들로 이루어지고 두 차수가 모두 아래로 유계인
이중복합체 $K_{\bullet\bullet}=(K_{ij})$를 생각하자. 언제나
$i<0$ 또는 $j<0$이면 $K_{ij}=0$인 경우로 한정할 수 있으므로,
이제부터 그렇게 하겠다. $K_{\bullet\bullet}$를 $\mathcal C$의
대상들로 이루어진 양의 차수 복합체들의 아벨 범주 $\mathcal K$의 대상
\[
K_{i,\bullet}=(K_{ij})_{j\geq0}
\]
들로 이루어진 (단순)복합체로 볼 수 있다. 보조정리
(\hyperref[0.11.5.2.1-ko]{11.5.2.1})(보다 정확히는 ``화살표의
방향을 뒤집어서'' 얻은 ``쌍대'' 보조정리)와 (M, V, 2.2)에 의해
$K_{\bullet\bullet}$은 범주 $\mathcal K$에서
\emph{사영 카르탕--아일렌베르크 분해}를 갖는다. 그러한 분해는 모든
차수가 $\geq0$이고 사영 대상들로 이루어진 $\mathcal C$의 삼중복합체
\[
M_{\bullet\bullet\bullet}=(M_{ijk})
\]
이다. 모든 $i$에 대해
\[
M_{i,\bullet\bullet},\quad
\mathrm B_i^{\mathrm I}(M_{\bullet\bullet\bullet}),\quad
\mathrm Z_i^{\mathrm I}(M_{\bullet\bullet\bullet}),\quad
H_i^{\mathrm I}(M_{\bullet\bullet\bullet})
\]
는 각각 범주 $\mathcal K$에서
\[
K_{i,\bullet},\quad
\mathrm B_i^{\mathrm I}(K_{\bullet\bullet}),\quad
\mathrm Z_i^{\mathrm I}(K_{\bullet\bullet}),\quad
H_i^{\mathrm I}(K_{\bullet\bullet})
\]
의 사영 분해를 이룬다. 특히 모든 쌍 $(i,j)$에 대해
$M_{ij,\bullet}$은 $\mathcal C$에서 $K_{ij}$의 사영 분해이다.
\end{env}

\begin{proposition}[11.7.2]
\label{0.11.7.2-fr}
$T$를 $\mathcal C$에서 아벨 범주 $\mathcal C'$으로 가는 공변 가법
함자라 하자. (\hyperref[0.11.7.1-fr]{11.7.1})의 표기 아래에서,
\emph{삼중복합체 $T(M_{\bullet\bullet\bullet})$에 대응하는
단순복합체의 호몰로지
$H_\bullet(T(M_{\bullet\bullet\bullet}))$는
$K_{\bullet\bullet}$에 대응하는 단순복합체의 초호몰로지
$L_\bullet T(K_{\bullet\bullet})$와 표준적으로 동형이다}
(\hyperref[0.11.6.2-fr]{11.6.2}). \emph{또한 이는}

\oldpage[0\textsubscript{III}]{42}

\emph{$t=a,b,a',b',c,d$에 대해 $\mathop{}^{(t)}E$로 나타내는 여섯
쌍정칙 스펙트럼 열의 공통 수렴 대상이며, 그 $E_2$ 항들은 다음
공식들로 주어진다.}
\[
\begin{aligned}
{}^{(a)}E^2_{pq}&=L_pT(H_q(K_{\bullet\bullet})),&
{}^{(b)}E^2_{pq}&=H_p(L_q^{\mathrm{II}}T(K_{\bullet\bullet})),\\
{}^{(a')}E^2_{pq}&=L_pT(H_q^{\mathrm I}(K_{\bullet\bullet})),&
{}^{(b')}E^2_{pq}&=H_p(L_qT(K_{\bullet\bullet})),\\
{}^{(c)}E^2_{pq}&=L_pT(H_q^{\mathrm{II}}(K_{\bullet\bullet})),&
{}^{(d)}E^2_{pq}&=H_p(L_q^{\mathrm I}T(K_{\bullet\bullet})).
\end{aligned}
\]
\end{proposition}

복합체 $A_\bullet=(A_i)$마다 대상 $F(A_i)$들로 이루어진 복합체를
$F(A_\bullet)$로 나타낸다는 표기를 사용함을 상기하자. 예를 들어
$L_q^{\mathrm{II}}T(K_{\bullet\bullet})$는 복합체
\[
(L_q^{\mathrm{II}}T(K_{i,\bullet}))_{i\geq0}
\]
를 나타내며, 여기서 $L_q^{\mathrm{II}}T(K_{i,\bullet})$는
단순복합체 $K_{i,\bullet}$에 관한 함자 $T$의 지표 $q$인
초호몰로지이다.

\begin{proof}
$L_\bullet$을 $K_{\bullet\bullet}$에 대응하는 단순복합체라 하자.
그러면
\[
L_i=\bigoplus_{r+s=i}K_{rs}.
\]
또한
\[
N_{ij}=\bigoplus_{r+s=i}M_{rsj}
\]
로 놓자. 모든 $i$에 대해 $N_{i,\bullet}$은 $\mathcal C$에서
$L_i$의 사영 분해임이 분명하다. 따라서
(\hyperref[0.11.6.3-fr]{11.6.3})과
(\hyperref[0.11.6.4-fr]{11.6.4})에 의해 함자적 동형
\[
L_\bullet T(L_\bullet)\simeq H_\bullet(T(N_{\bullet\bullet}))
\]
이 성립한다. 이중복합체 $T(N_{\bullet\bullet})$에 대응하는
단순복합체는 삼중복합체 $T(M_{\bullet\bullet\bullet})$에도
대응하므로, 이는 명제의 첫째 주장을 증명한다.

또한 $L_\bullet T(L_\bullet)$은 단순복합체 $L_\bullet$에 관한
$T$의 두 초호몰로지 스펙트럼 열
(\hyperref[0.11.6.2-fr]{11.6.2})의 수렴 대상이고, 이 두 열은 각각
$\mathop{}^{(b')}E$와 $\mathop{}^{(a)}E$이다.

이제 $M_{\bullet\bullet\bullet}$을
\[
U_{ij}=\bigoplus_{r+s=j}M_{irs}
\]
인 이중복합체 $U_{\bullet\bullet}$로 보자.
$H_\bullet(T(M_{\bullet\bullet\bullet}))$은 여전히 이중복합체
$T(U_{\bullet\bullet})$의 두 스펙트럼 열의 수렴 대상이다. 그런데
모든 $i$에 대해 $M_{i,\bullet\bullet}$은 단순복합체
$K_{i,\bullet}$에 관하여
(\hyperref[0.11.6.3-fr]{11.6.3})의 조건을 만족하는 이중복합체이다.
따라서
\[
H_q^{\mathrm{II}}(T(U_{i,\bullet}))
=L_q^{\mathrm{II}}T(K_{i,\bullet}),
\]
이고 $T(U_{\bullet\bullet})$의 첫째 스펙트럼 열은 바로
$\mathop{}^{(b)}E$이다. 한편 모든 $r$에 대해
$M_{\bullet,r,\bullet}$은 단순복합체 $K_{\bullet,r}$의
카르탕--아일렌베르크 분해이다. (M, XV, 2)의 계산에 의해
\[
H_q^{\mathrm I}(T(M_{\bullet,rs}))
=T(H_q^{\mathrm I}(M_{\bullet,rs})),
\]
따라서
\[
H_q^{\mathrm I}(T(U_{\bullet,j}))
=\bigoplus_{r+s=j}T(H_q^{\mathrm I}(M_{\bullet,rs})).
\]
다시 말해 단순복합체
$H_q^{\mathrm I}(T(U_{\bullet\bullet}))$은 이중복합체
$T(H_q^{\mathrm I}(M_{\bullet\bullet\bullet}))$에 대응하는
단순복합체이다. 그런데 모든 $q$에 대해
$H_q^{\mathrm I}(M_{\bullet\bullet\bullet})$은 범주 $\mathcal K$에서
단순복합체 $H_q^{\mathrm I}(K_{\bullet\bullet})$의 사영 분해이다.
(\hyperref[0.11.6.3-fr]{11.6.3})을 적용하면
\[
H_p^{\mathrm{II}}
\bigl(T(H_q^{\mathrm I}(M_{\bullet\bullet\bullet}))\bigr)
=L_pT(H_q^{\mathrm I}(K_{\bullet\bullet}))
\]
을 얻고, 따라서 스펙트럼 열 $\mathop{}^{(a')}E$를 얻는다. 끝으로
삼중복합체 $M_{\bullet\bullet\bullet}$의 정의에서 지표 $i$와 $j$의
역할을 서로 바꾸고 앞의 논증을 이 새 삼중복합체에 적용하면
$\mathop{}^{(c)}E$와 $\mathop{}^{(d)}E$를 얻는다.
\end{proof}

$L_\bullet T(K_{\bullet\bullet})$를 이중복합체
$K_{\bullet\bullet}$에 관한 $T$의 \emph{초호몰로지}라 한다.

\begin{env}[11.7.3]
\label{0.11.7.3-fr}
\emph{주석.}
\begin{enumerate}
\item[\rm(i)]
(\hyperref[0.11.6.3-fr]{11.6.3})에 의해
$L_\bullet T(K_{\bullet\bullet})$은 두 차수가 각각 아래로 유계인
$\mathcal C$의 이중복합체 $K_{\bullet\bullet}$들의 범주에서
호몰로지 함자이다.

\oldpage[0\textsubscript{III}]{43}

\item[\rm(ii)]
$M_{\bullet\bullet\bullet}$을 모든 쌍 $(r,s)$에 대해
$M_{rs,\bullet}$이 $K_{rs}$의 분해이고, 모든 삼중지표 $(i,j,k)$와
모든 $n>0$에 대해 $L_nT(M_{ijk})=0$인 $\mathcal C$의
삼중복합체라 하자. 그러면 동형
\[
L_\bullet T(K_{\bullet\bullet})
\simeq H_\bullet(T(M_{\bullet\bullet\bullet}))
\]
이 성립한다. 실제로 (\hyperref[0.11.7.2-fr]{11.7.2})의 증명에서
사용한 표기 아래에서 $N_{i,\bullet}$은 $L_i$의 분해이고 모든
$n>0$와 모든 쌍 $(i,j)$에 대해 $L_nT(N_{ij})=0$이므로,
(\hyperref[0.11.6.3-fr]{11.6.3}, (iii))을 적용하면 충분하다.

\item[\rm(iii)]
(\hyperref[0.11.7.2-fr]{11.7.2})의 결과는 공변 다변함자로 곧바로
일반화된다. 예를 들어 $\mathcal C'$을 모든 대상이 어떤 사영 대상의
몫인 또 하나의 아벨 범주, $K'_{\bullet\bullet}$을 두 차수가 아래로
유계인 $\mathcal C'$의 이중복합체, $T$를
$\mathcal C\times\mathcal C'$에서 아벨 범주 $\mathcal C''$으로
가는 공변 가법 쌍함자라 하자. $L_\bullet$과 $L'_\bullet$을 각각
$K_{\bullet\bullet}$과 $K'_{\bullet\bullet}$에 대응하는
단순복합체라 하면, 두 이중복합체 $K_{\bullet\bullet}$,
$K'_{\bullet\bullet}$에 관한 $T$의 초호몰로지를 초호몰로지
\[
L_\bullet T(L_\bullet,L'_\bullet)
\]
로 정의한다. (\hyperref[0.11.6.4-fr]{11.6.4})와
(\hyperref[0.11.6.5-fr]{11.6.5})를 적용하면
(\hyperref[0.11.7.2-fr]{11.7.2})에서와 같이 이 초호몰로지를 공통
수렴 대상으로 갖는 여섯 스펙트럼 열을 얻는다. 이 열들을 써 보는
일은 독자에게 맡긴다.
\end{enumerate}
\end{env}

\subsection{단체 복합체의 코호몰로지에 관한 보충}
\label{subsection:0.11.8-fr}

\begin{env}[11.8.1]
\label{0.11.8.1-fr}
$A$를 유한집합이라 하고, $\Sigma(A)$를 $A$의 원소들로 이루어진 유한열
$\sigma=(\alpha_0,\ldots,\alpha_h)$들의 집합($A$의 ``단체''들의 집합)이라 하자.
$|\sigma|=\{\alpha_0,\ldots,\alpha_h\}$로 놓는다. 사슬 복합체
$\mathrm C_\bullet(A)$는 $\Sigma(A)$의 원소들이 생성하는 등급 자유
아벨군이며, $(\alpha_0,\ldots,\alpha_h)$의 차수는 $h$이고, 그 미분은
다음과 같이 정의됨을 상기하자.
\[
d(\alpha_0,\ldots,\alpha_h)=
\sum_{j=0}^{h}(-1)^j
(\alpha_0,\ldots,\widehat{\alpha_j},\ldots,\alpha_h).
\]
$\mathrm C_\bullet(A)$에서 두 $\alpha_i$가 같은 사슬
$\sigma=(\alpha_0,\ldots,\alpha_h)$들과 다음 사슬들이 생성하는 부분군을
$\mathrm D_\bullet(A)$라 하자.
$\pi(\sigma)-\varepsilon_\pi\sigma$. 여기서 각 순열 $\pi$에 대하여
\[
\pi(\sigma)=(\alpha_{\pi(0)},\ldots,\alpha_{\pi(h)})
\]
이고 $\varepsilon_\pi$는 $\pi$의 부호이다. 그러면
$\mathrm D_\bullet(A)$는 $\mathrm C_\bullet(A)$의 부분복합체이다. 그
원소들을 \emph{퇴화 사슬}이라 부른다. 이제
$\mathrm L_\bullet(A)=\mathrm C_\bullet(A)/\mathrm D_\bullet(A)$로 놓으면
복합체의 자연 준동형
\[
p:\mathrm C_\bullet(A)\longrightarrow\mathrm L_\bullet(A)
\]
이 있다. 한편 복합체의 준동형
\[
j:\mathrm L_\bullet(A)\longrightarrow\mathrm C_\bullet(A)
\]
을 다음과 같이 정의한다. $A$에 전순서를 하나 정한다. 단체
$\sigma=(\alpha_0,\ldots,\alpha_h)$의 mod.~$\mathrm D_\bullet(A)$ 동치류에
대하여, 두 $\alpha_i$가 같으면 $0$을 대응시키고, 그렇지 않으면
$\alpha_i$들을 오름차순으로 배열한 열 $(\beta_0,\ldots,\beta_h)$를
대응시킨다. $p\circ j$가 $\mathrm L_\bullet(A)$의 항등사상임은 분명하다.
\end{env}

\begin{env}[11.8.2]
\label{0.11.8.2-fr}
$B$를 또 하나의 유한집합이라 하자. $d'$, $d''$가 각각
$\mathrm C_\bullet(A)$와 $\mathrm C_\bullet(B)$의 미분이면, 텐서곱 복합체
$\mathrm C_\bullet(A)\otimes\mathrm C_\bullet(B)$는
$\Sigma(A)\times\Sigma(B)$의 원소들이 생성하는 자유 아벨군으로 볼 수
있으며, 그 미분은 다음과 같다.
\[
d(\sigma,\tau)=(d'\sigma,\tau)+(-1)^{h+1}(\sigma,d''\tau)
\quad\text{$\operatorname{Card}|\sigma|=h+1$이면.}
\]
자연 준동형들
$\mathrm C_\bullet(A)\to\mathrm L_\bullet(A)$,
$\mathrm C_\bullet(B)\to\mathrm L_\bullet(B)$는 준동형
\[
p:\mathrm C_\bullet(A)\otimes\mathrm C_\bullet(B)
\longrightarrow\mathrm L_\bullet(A)\otimes\mathrm L_\bullet(B)
\]
을 정의한다. 뒤의 텐서곱은 다음과 동형이다.
\[
\frac{\mathrm C_\bullet(A)\otimes\mathrm C_\bullet(B)}
{\mathrm C_\bullet(A)\otimes\mathrm D_\bullet(B)
+\mathrm D_\bullet(A)\otimes\mathrm C_\bullet(B)}.
\]
마찬가지로, (\hyperref[0.11.8.1-fr]{11.8.1})에서 정의한 준동형
$\mathrm L_\bullet(A)\to\mathrm C_\bullet(A)$와
$\mathrm L_\bullet(B)\to\mathrm C_\bullet(B)$를($A$와 $B$에 전순서를
정하여) 이용하면, $p\circ j$가 항등사상이 되도록 하는 준동형
\[
j:\mathrm L_\bullet(A)\otimes\mathrm L_\bullet(B)
\longrightarrow\mathrm C_\bullet(A)\otimes\mathrm C_\bullet(B)
\]
을 정의할 수 있다.
\end{env}

\oldpage[0\textsubscript{III}]{44}
이 표기 아래에서 다음이 성립한다.
\begin{proposition}[11.8.3]
\label{0.11.8.3-fr}
\emph{호모토피
\[
h:\mathrm C_\bullet(A)\otimes\mathrm C_\bullet(B)
\longrightarrow\mathrm C_\bullet(A)\otimes\mathrm C_\bullet(B)
\]
가 존재하여, $h(\sigma,\tau)$는 $|\sigma_i|\subset|\sigma|$,
$|\tau_i|\subset|\tau|$를 만족하는 단체쌍 $(\sigma_i,\tau_i)$들의
선형결합이고, $f=j\circ p$에 대하여 다음이 성립한다.}
\[
\label{0.11.8.3.1-fr}
\tag{11.8.3.1}
f-1=h\circ d+d\circ h.
\]
\end{proposition}

\begin{proof}
$\sigma$와 $\tau$의 차수의 합에 관한 귀납법으로 각 단체쌍
$(\sigma,\tau)$ 위에서 $h$를 정의하면 충분하다. 그 합이 $0$이면
$h=0$으로 취할 수 있기 때문이다. 다음과 같이 놓자.
\[
\omega=f(\sigma,\tau)-(\sigma,\tau)-h(d(\sigma,\tau)).
\]
귀납 가정과 $d$의 정의에 따라
\[
\omega\in\mathrm C_\bullet(|\sigma|)\otimes\mathrm C_\bullet(|\tau|)
\]
이다. 또한
\[
d\omega=f(d(\sigma,\tau))-d(\sigma,\tau)-d(h(d(\sigma,\tau)))
=h(d(d(\sigma,\tau)))=0
\]
이다. 이는 (\hyperref[0.11.8.3.1-fr]{11.8.3.1})과 귀납 가정에서
따른다. 그런데 $q>0$이면 $H_q(\mathrm C_\bullet(A))=0$이고
(G, I, 3.7.4), 따라서 퀸네트 공식(G, I, 5.5.2)에 의해
\[
H_q(\mathrm C_\bullet(A)\otimes\mathrm C_\bullet(B))=0\quad(q>0)
\]
이기도 하다. 이 사실에서 $A$를 $|\sigma|$로, $B$를 $|\tau|$로
바꾸어 적용하면
\[
\omega'\in\mathrm C_\bullet(|\sigma|)\otimes\mathrm C_\bullet(|\tau|)
\]
이며 $\omega=d\omega'$인 원소가 존재함을 알 수 있다.
$h(\sigma,\tau)=\omega'$으로 취하면 단체쌍 $(\sigma,\tau)$에 대하여
(\hyperref[0.11.8.3.1-fr]{11.8.3.1})을 확인할 수 있으므로 귀납을 계속할
수 있다.
\end{proof}

\begin{env}[11.8.4]
\label{0.11.8.4-fr}
표기는 (\hyperref[0.11.8.2-fr]{11.8.2})에서와 같이 하자.
$|\sigma|\subset|\sigma'|$이고 $|\tau|\subset|\tau'|$일 때
$(\sigma,\tau)\leq(\sigma',\tau')$로 놓는다.
$\Sigma(A)\times\Sigma(B)$ 위의 \emph{계수계} $\Gamma$는 다음 데이터로
이루어진다. 첫째, $\Gamma_{\sigma,\tau}$가 오직 집합 $|\sigma|$와
$|\tau|$에만 의존하는 아벨군들의 족 $(\Gamma_{\sigma,\tau})$가 있고,
둘째,
\[
\Gamma_{\sigma,\tau}\longrightarrow\Gamma_{\sigma',\tau'}
\quad\text{$(\sigma,\tau)\leq(\sigma',\tau')$이면}
\]
인 준동형들의 족이 있어서 이 전순서 관계에 관한 귀납계를 이룬다. 이제
공사슬 복합체 $\mathrm C^\bullet(A,B;\Gamma)$를 다음과 같이 정의한다.
그 원소는 $\Sigma(A)\times\Sigma(B)$를 움직이는
$(\sigma,\tau)$에 대한 족 $\lambda=(\lambda(\sigma,\tau))$이며, 각 쌍에
대하여 $\lambda(\sigma,\tau)\in\Gamma_{\sigma,\tau}$이다. 미분은 다음과
같이 준다. 만일
\[
d(\sigma,\tau)=\sum_i\pm(\sigma_i,\tau_i)
\]
이면, 모든 $i$에 대하여 $|\sigma_i|\subset|\sigma|$이고
$|\tau_i|\subset|\tau|$이며,
\[
d\lambda(\sigma,\tau)=\sum_i\pm\lambda_i(\sigma_i,\tau_i)
\]
로 취한다. 여기서 $\lambda_i(\sigma_i,\tau_i)$는
$\lambda(\sigma_i,\tau_i)$의 $\Gamma_{\sigma,\tau}$ 안에서의 표준 상을
나타낸다.

공사슬 $\lambda\in\mathrm C^\bullet(A,B;\Gamma)$가 다음 조건을 만족하면
\emph{이중 교대적}이라고 하겠다. 두 단체 $\sigma$, $\tau$ 중 하나에
같은 항이 두 개 있으면 $\lambda(\sigma,\tau)=0$이고, 지표들의 임의의
순열 $\pi$, $\pi'$에 대하여
\[
\lambda(\pi(\sigma),\tau)=\varepsilon_\pi\lambda(\sigma,\tau),
\qquad
\lambda(\sigma,\pi'(\tau))=\varepsilon_{\pi'}\lambda(\sigma,\tau)
\]
이다. 이러한 공사슬들이 $\mathrm C^\bullet(A,B;\Gamma)$의 부분복합체
$\mathrm L^\bullet(A,B;\Gamma)$를 이룸은 분명하다.
\end{env}

\begin{proposition}[11.8.5]
\label{0.11.8.5-fr}
\emph{표준 단사 준동형
\[
\mathrm L^\bullet(A,B;\Gamma)\longrightarrow
\mathrm C^\bullet(A,B;\Gamma)
\]
은 이 두 복합체의 코호몰로지 사이의 동형을 정의한다.}
\end{proposition}

\begin{proof}
$p$와 $j$를 (\hyperref[0.11.8.2-fr]{11.8.2})에서 정의한 뜻으로 쓰면,
사상들
\[
{}'p:\lambda\longmapsto\lambda\circ p,\qquad
{}'j:\lambda\longmapsto\lambda\circ j
\]
은 각각 $\mathrm L^\bullet(A,B;\Gamma)$와
$\mathrm C^\bullet(A,B;\Gamma)$에서 정의되며, 첫째 것은 바로 표준 단사
준동형이다. $\mathop{}'j\circ{}'p$가 항등사상이므로
$\mathop{}'p\circ{}'j$가 항등사상과 호모토픽임을 보이면 충분하다. 그런데
(\hyperref[0.11.8.3-fr]{11.8.3})에 따라
$\mathop{}'h:\lambda\mapsto\lambda\circ h$가
$\mathrm C^\bullet(A,B;\Gamma)$에서 정의되므로, 항등식
(\hyperref[0.11.8.3.1-fr]{11.8.3.1})을 전치하여 원하는 결과를 얻는다.
\end{proof}

\begin{env}[11.8.6]
\label{0.11.8.6-fr}
\oldpage[0\textsubscript{III}]{45}
명제 (\hyperref[0.11.8.5-fr]{11.8.5})에 의해
$\mathrm L^\bullet(A,B;\Gamma)$의 코호몰로지 계산은
$\mathrm C^\bullet(A,B;\Gamma)$의 코호몰로지 계산으로 귀착된다. 한편
아일렌베르크--질버 정리(G, I, 3.10.2)에 따라 뒤의 코호몰로지는 다음과
같이 정의되는 공사슬 복합체의 코호몰로지와 표준적으로 동형임을 상기하자.
사슬 복합체 $\mathrm P_\bullet(A,B)$를 구성한다. 이는 $\sigma$와 $\tau$의
차수가 같은 $(\sigma,\tau)\in\Sigma(A)\times\Sigma(B)$들의 선형결합으로
이루어지며, 그 미분은
\[
d(\sigma,\tau)=\sum_{j,k}(-1)^{j+k}(\sigma_j,\tau_k)
\]
이다. 여기서
\[
d\sigma=\sum_j(-1)^j\sigma_j,\qquad
d\tau=\sum_k(-1)^k\tau_k
\]
이다. 이때 복합체의 표준 준동형 두 개
\[
f:\mathrm P_\bullet(A,B)\longrightarrow
\mathrm C_\bullet(A)\otimes\mathrm C_\bullet(B),\qquad
g:\mathrm C_\bullet(A)\otimes\mathrm C_\bullet(B)
\longrightarrow\mathrm P_\bullet(A,B)
\]
가 있으며, 호모토피 $h$, $h'$이 존재하여 다음을 만족함이
증명되어 있다(\emph{인용한 곳}).
\[
f\circ g-1=d\circ h+h\circ d,\qquad
g\circ f-1=d\circ h'+h'\circ d.
\]
더욱이
\[
f(\sigma,\tau)\in\mathrm C_\bullet(|\sigma|)
\otimes\mathrm C_\bullet(|\tau|),\qquad
g(\sigma,\tau)\in\mathrm P_\bullet(|\sigma|,|\tau|)
\]
이고, 호모토피 $h$, $h'$은 다음을 만족하도록 취할 수 있다.
\[
h(\sigma,\tau)\in\mathrm C_\bullet(|\sigma|)
\otimes\mathrm C_\bullet(|\tau|),\qquad
h'(\sigma,\tau)\in\mathrm P_\bullet(|\sigma|,|\tau|).
\]
이는 $h(\sigma,\tau)$와 $h'(\sigma,\tau)$를 $\sigma$와 $\tau$의 차수의
합에 관한 귀납법으로 정의할 수 있고,
\[
H_q(\mathrm C_\bullet(|\sigma|)\otimes\mathrm C_\bullet(|\tau|))
\quad\text{와}\quad
H_q(\mathrm P_\bullet(|\sigma|,|\tau|))
\]
가 $q>0$일 때 $0$이기 때문이다(\emph{인용한 곳}). 그러면
(\hyperref[0.11.8.3-fr]{11.8.3})에서와 같이 논증하여 결론을 얻는다.

이제 $\mathrm P^\bullet(A,B;\Gamma)$를 다음 족들의 집합으로 정의한다.
$(\sigma,\tau)$는 두 항의 차수가 같은 쌍들을 움직이고,
$\lambda=(\lambda(\sigma,\tau))$이며
$\lambda(\sigma,\tau)\in\Gamma_{\sigma,\tau}$이다. 또한
\[
d\sigma=\sum_j(-1)^j\sigma_j,\qquad
d\tau=\sum_k(-1)^k\tau_k
\]
이므로
\[
d\lambda(\sigma,\tau)=
\sum_{j,k}(-1)^{j+k}\lambda(\sigma_j,\tau_k)
\]
가 정의되며, 이것이 복합체 $\mathrm P^\bullet(A,B;\Gamma)$의 미분을
준다. 이때 앞의 고찰에 따라 사상들
\[
{}'f:\lambda\mapsto\lambda\circ f,\quad
{}'g:\lambda\mapsto\lambda\circ g,\quad
{}'h:\lambda\mapsto\lambda\circ h,\quad
{}'h':\lambda\mapsto\lambda\circ h'
\]
은 모두 정의된다. 따라서 $\mathop{}'f\circ{}'g$와
$\mathop{}'g\circ{}'f$는 항등사상과 호모토픽이고, 이로부터
$\mathrm C^\bullet(A,B;\Gamma)$의 코호몰로지와
$\mathrm P^\bullet(A,B;\Gamma)$의 코호몰로지 사이의 원하는 동형을
얻는다.
\end{env}

\begin{env}[11.8.7]
\label{0.11.8.7-fr}
\emph{주의}. (\hyperref[0.11.8.3-fr]{11.8.3})에서와 같은 논증을
$\mathrm C_\bullet(A)$와 $\mathrm L_\bullet(A)$에 적용하면 다음을 알 수
있다. $j$와 $p$를 (\hyperref[0.11.8.1-fr]{11.8.1})에서와 같이 정의할 때
$f=j\circ p$는 다시 (\hyperref[0.11.8.3.1-fr]{11.8.3.1})의 관계식을
만족하고, $|h(\sigma)|\subset|\sigma|$이다. 따라서
(\hyperref[0.11.8.5-fr]{11.8.5})에서와 같이, 자명한 방식으로 정의되는 두
복합체 $\mathrm L^\bullet(A;\Gamma)$와
$\mathrm C^\bullet(A;\Gamma)$의 코호몰로지 사이의 동형을 얻는다. 이것이
(G, I, 3.8.1)에서 증명이 개략적으로 제시된 결과이다.
\end{env}

\begin{env}[11.8.8]
\label{0.11.8.8-fr}
이제 (\hyperref[0.11.8.2-fr]{11.8.2})의 표기와 가정으로 돌아가서,
$\Sigma(A)\times\Sigma(B)$ 위의 계수계들의 복합체
$\Gamma^\bullet=(\Gamma^k)$를 생각하자. 따라서 각 $(\sigma,\tau)$에
대하여 $\Gamma^k_{\sigma,\tau}$들은 아벨군의 복합체를 이루며
$(k\in\mathbb Z)$, $(\sigma,\tau)\leq(\sigma',\tau')$일 때 다음 가환
도표들이 있다.
\[
\xymatrix{
\Gamma^k_{\sigma,\tau}\ar[r]\ar[d] &
\Gamma^{k+1}_{\sigma,\tau}\ar[d]\\
\Gamma^k_{\sigma',\tau'}\ar[r] &
\Gamma^{k+1}_{\sigma',\tau'}
}
\]
\oldpage[0\textsubscript{III}]{46}
그러면
\[
\mathrm C^\bullet(A,B;\Gamma^\bullet)
=\bigl(\mathrm C^h(A,B;\Gamma^k)\bigr)_{(h,k)\in\mathbb Z\times\mathbb Z}
\]
가 아벨군의 이중복합체이고,
\[
\mathrm L^\bullet(A,B;\Gamma^\bullet)
=\bigl(\mathrm L^h(A,B;\Gamma^k)\bigr)
\]
가 그 부분 이중복합체임을 곧바로 확인할 수 있다.
\end{env}

\begin{proposition}[11.8.9]
\label{0.11.8.9-fr}
\emph{표준 단사 준동형
\[
\mathrm L^\bullet(A,B;\Gamma^\bullet)\longrightarrow
\mathrm C^\bullet(A,B;\Gamma^\bullet)
\]
은 이 두 이중복합체의 코호몰로지 사이의 동형을 정의한다.}
\end{proposition}

\begin{proof}
간단히 쓰기 위해 $\mathrm C^{\bullet\bullet}=
\mathrm C^\bullet(A,B;\Gamma^\bullet)$,
$\mathrm L^{\bullet\bullet}=
\mathrm L^\bullet(A,B;\Gamma^\bullet)$로 놓자.
$h<0$이면 $\mathrm C^{hk}=\mathrm L^{hk}=0$이므로 이 이중복합체들의 둘째
스펙트럼 열들은 정칙이다
(\hyperref[0.11.3.3-fr]{11.3.3}). 따라서 준동형
$\mathrm L^{\bullet\bullet}\to\mathrm C^{\bullet\bullet}$은 스펙트럼 열의
사상
\[
{}''E(\mathrm L^{\bullet\bullet})\longrightarrow
{}''E(\mathrm C^{\bullet\bullet})
\]
을 주고, $E_2$ 항에서는 이것이 다음으로 귀착된다.
\[
\label{0.11.8.9.1-fr}
\tag{11.8.9.1}
H^p_{\mathrm{II}}(H^q_{\mathrm I}(\mathrm L^{\bullet\bullet}))
\longrightarrow
H^p_{\mathrm{II}}(H^q_{\mathrm I}(\mathrm C^{\bullet\bullet})).
\]
그런데 모든 $k\in\mathbb Z$에 대하여
(\hyperref[0.11.8.5-fr]{11.8.5})로부터 준동형
$H^q_{\mathrm I}(\mathrm L^{\bullet,k})\to
H^q_{\mathrm I}(\mathrm C^{\bullet,k})$가 전단사임을 알 수 있다. 따라서
(\hyperref[0.11.1.5-fr]{11.1.5})로부터 결론이 따른다.
\end{proof}

\begin{env}[11.8.10]
\label{0.11.8.10-fr}
마찬가지로 (\hyperref[0.11.8.6-fr]{11.8.6})의 표기 아래에서 이중복합체의
표준 준동형들
\[
\mathrm C^\bullet(A,B;\Gamma^\bullet)\longrightarrow
\mathrm P^\bullet(A,B;\Gamma^\bullet)
\]
이 있다(표기는 자명하다). 이번에는
(\hyperref[0.11.8.6-fr]{11.8.6})에 근거하여
(\hyperref[0.11.8.9-fr]{11.8.9})에서와 같이 논증하면, 이 준동형도
코호몰로지 사이의 동형을 줌을 알 수 있다.
\end{env}

\subsection{유한 생성 복합체에 관한 보조정리}
\label{subsection:0.11.9-fr}

\begin{proposition}[11.9.1]
\label{0.11.9.1-fr}
$\mathcal C$를 아벨 범주라 하고, $K'$과 $K''$을 $\mathcal C$의 대상들의
집합의 부분집합이라 하자. $K''\subset K'$이고 다음 조건들을 만족한다고
하자.
\begin{enumerate}
\item[{\rm(i)}] \emph{모든 대상 $A'\in K'$과 $\mathcal C$ 안의 모든
에피사상 $u:A\to A'$에 대하여, 어떤 대상 $B\in K''$과 사상
$v:B\to A$가 존재하여 $u\circ v$가 에피사상이다.}
\item[{\rm(ii)}] \emph{모든 대상쌍 $A\in K'$, $B\in K'$과 모든
에피사상 $u:A\to B$에 대하여 $\operatorname{Ker}(u)$는 $K'$에 속한다.}
\item[{\rm(iii)}] \emph{$K''$의 두 대상의 곱은 $K''$에 속한다.}
\end{enumerate}
\emph{$P_\bullet=(P_i)_{i\in\mathbb Z}$를 $\mathcal C$ 안의 복합체라 하자.
모든 $i$에 대하여 $H_i(P_\bullet)\in K'$이고, 어떤 $d$가 존재하여
$i<d$일 때 $H_i(P_\bullet)=0$이라고 하자. 그러면 $\mathcal C$ 안의
복합체 $Q_\bullet=(Q_i)_{i\in\mathbb Z}$와 복합체의 사상
$u:Q_\bullet\to P_\bullet$가 존재하여, 모든 $i$에 대하여 $Q_i\in K''$이고
$i<d$일 때 $Q_i=0$이며, 이에 대응하는 사상
$H_\bullet(Q_\bullet)\to H_\bullet(P_\bullet)$은 동형이다.}
\end{proposition}

\begin{proof}
먼저 성질 (i)에서 따르는 다음 결과를 증명하자.

\emph{(i bis) $u:C\to B$를 $\mathcal C$ 안의 에피사상, $A$를 $K'$의
대상, $v:A\to B$를 $\mathcal C$ 안의 사상이라 하자. 그러면 대상
$D\in K''$, 에피사상 $u':D\to A$, 사상 $v':D\to C$가 존재하여 도표}
\[
\xymatrix{
D\ar[r]^{u'}\ar[d]_{v'} & A\ar[d]^{v}\\
C\ar[r]_{u} & B
}
\tag{11.9.1.1}\label{0.11.9.1.1-fr}
\]
\emph{가 가환이다.}

실제로 $\mathcal C$ 안의 섬유곱 $C\times_B A$를 생각하고, 도표
\oldpage[0\textsubscript{III}]{47}
\[
\xymatrix{
C\times_B A\ar[r]^{q}\ar[d]_{p} & A\ar[d]^{v}\\
C\ar[r]_{u} & B
}
\tag{11.9.1.2}\label{0.11.9.1.2-fr}
\]
가 가환이 되게 하는 표준 사영 $p:C\times_B A\to C$,
$q:C\times_B A\to A$를 취하자. $q$의 여핵은
$v^{-1}(u(C))$에 의한 $A$의 몫임이 알려져 있다([27], p.~1-12).
$u$가 에피사상이므로 $u(C)=B$이고 $v^{-1}(u(C))=A$이다. 따라서 $q$는
에피사상이다. 이제 에피사상 $q:C\times_B A\to A$에 (i)를 적용하면,
어떤 대상 $D\in K''$과 사상 $w:D\to C\times_B A$가 존재하여
$q\circ w$가 에피사상이다. $u'=q\circ w$, $v'=p\circ w$로 취하면 된다.

이제 명제를 증명하기 위해 귀납법을 쓰자. 어떤 $i\geq d-1$에 대하여,
$j\leq i$인 대상 $Q_j$, 사상 $d_j:Q_j\to Q_{j-1}$, 사상
$u_j:Q_j\to P_j$가 이미 구성되어 다음을 만족한다고 하자. $j<d$이면
$Q_j=0$이고, $j\leq i$에 대하여 $d_{j-1}\circ d_j=0$이며
$d_j\circ u_j=u_{j-1}\circ d_j$이다. 더욱이 다음 조건들이 성립한다고
가정한다.
\begin{enumerate}
\item[$(\mathrm I_i)$] $j\leq i$이면 $Q_j\in K''$이고, $j<i$이면
$B_j(Q_\bullet)\in K'$이다.
\item[$(\mathrm{II}_i)$] $j<i$이면, 족 $(u_k)_{k\leq i}$에서 유도되는
준동형 $H_j(Q_\bullet)\to H_j(P_\bullet)$은 동형이다.
\item[$(\mathrm{III}_i)$] 합성사상
\[
v_i:Z_i(Q_\bullet)\longrightarrow Z_i(P_\bullet)
\longrightarrow H_i(P_\bullet)
\]
은 에피사상이다. 여기서 왼쪽 화살표는 $u_i$의 제한이고 오른쪽 화살표는
표준 사상이다.
\end{enumerate}
(ii)에 따라, 에피사상 $Q_i\to B_{i-1}(Q_\bullet)$의 핵
$Z_i(Q_\bullet)$은 가정 $(\mathrm I_i)$에 의해 $K'$에 속한다. 또한 (ii)와
가정 $(\mathrm{III}_i)$에 따라 $N_i=\operatorname{Ker}(v_i)$도 $K'$에
속한다. (i bis)에 의해 $Q'_{i+1}\in K''$, 에피사상
$d'_{i+1}:Q'_{i+1}\to N_i$, 사상 $u'_{i+1}:Q'_{i+1}\to P_{i+1}$이
존재하여 도표
\[
\xymatrix{
Q'_{i+1}\ar[r]^{d'_{i+1}}\ar[d]_{u'_{i+1}} &
N_i\ar[d]^{u_i}\\
P_{i+1}\ar[r]_{d_{i+1}} & B_i(P_\bullet)
}
\tag{11.9.1.3}\label{0.11.9.1.3-fr}
\]
가 가환이다.

표준 사상 $Z_{i+1}(P_\bullet)\to H_{i+1}(P_\bullet)$은 에피사상이고,
가정에 따라 $H_{i+1}(P_\bullet)\in K'$이다. 따라서 (i)에 의해 대상
$Q''_{i+1}\in K''$과 사상
$u''_{i+1}:Q''_{i+1}\to Z_{i+1}(P_\bullet)$이 존재하여 합성사상
$Q''_{i+1}\to Z_{i+1}(P_\bullet)\to H_{i+1}(P_\bullet)$이
에피사상이다. $d''_{i+1}:Q''_{i+1}\to N_i$를 $0$으로 취하면 도표
\[
\xymatrix{
Q''_{i+1}\ar[r]^{d''_{i+1}}\ar[d]_{u''_{i+1}} &
N_i\ar[d]^{u_i}\\
Z_{i+1}(P_\bullet)\ar[r]_{d_{i+1}} & P_i
}
\tag{11.9.1.4}\label{0.11.9.1.4-fr}
\]
는 가환이다. 아래쪽 가로 화살표가 $0$이기 때문이다.
\oldpage[0\textsubscript{III}]{48}
이제 다음과 같이 취하자.
\[
Q_{i+1}=Q'_{i+1}\times Q''_{i+1},\qquad
d_{i+1}=d'_{i+1}+d''_{i+1},\qquad
u_{i+1}=u'_{i+1}+u''_{i+1}.
\]
$d_{i+1}(Q_{i+1})=d'_{i+1}(Q'_{i+1})=N_i\subset Z_i(Q_\bullet)$이므로
$d_i\circ d_{i+1}=0$이고, 통상적인 표기 아래에서
$B_i(Q_\bullet)=N_i$이다. 따라서 $(\mathrm I_{i+1})$이 확인된다. 도표
(\hyperref[0.11.9.1.3-fr]{11.9.1.3})과
(\hyperref[0.11.9.1.4-fr]{11.9.1.4})의 가환성에 따라
$d_{i+1}\circ u_{i+1}=u_i\circ d_{i+1}$이다. $N_i$의 정의에 의해,
$u_k$들의 계($k\leq i+1$)에서 유도되는 사상
\[
H_i(Q_\bullet)=Z_i(Q_\bullet)/N_i\longrightarrow H_i(P_\bullet)
\]
은 몫으로 내려서 $v_i$에서 유도되는 사상이다. $v_i$가 에피사상이므로
이 사상은 동형이고, 따라서 $(\mathrm{II}_{i+1})$이 성립한다. 마지막으로
정의에 의해 $Q''_{i+1}\subset Z_{i+1}(Q_\bullet)$이다.
$u''_{i+1}$의 선택에 따라 사상
\[
v_{i+1}:Z_{i+1}(Q_\bullet)\longrightarrow Z_{i+1}(P_\bullet)
\longrightarrow H_{i+1}(P_\bullet)
\]
은 에피사상이다. 실제로 $Q''_{i+1}$에 대한 그 제한이 이미
에피사상이다. 따라서 $(\mathrm{III}_{i+1})$이 성립한다. 그러므로 귀납을
계속할 수 있고, 명제가 증명된다.
\end{proof}

\begin{corollary}[11.9.2]
\label{0.11.9.2-fr}
$A$를 뇌터 환이라 하자(가환일 필요는 없다). $P_\bullet=(P_i)_{i\in\mathbb Z}$를
오른쪽 $A$-가군들의 복합체라 하자. $H_i(P_\bullet)$들이 유한 생성
$A$-가군이고, 어떤 $d$가 존재하여 $i<d$일 때 $H_i(P_\bullet)=0$이라고
가정하자. 그러면 유한 계수의 자유 오른쪽 $A$-가군들로 이루어진 복합체
$Q_\bullet=(Q_i)_{i\in\mathbb Z}$가 존재하여 $i<d$일 때 $Q_i=0$이고,
복합체의 준동형 $u:Q_\bullet\to P_\bullet$가 존재하여 $u$에 대응하는
준동형 $H_\bullet(Q_\bullet)\to H_\bullet(P_\bullet)$은 전단사이다.
\end{corollary}

\begin{proof}
(\hyperref[0.11.9.1-fr]{11.9.1})에서 $\mathcal C$를 오른쪽
$A$-가군들의 범주로, $K'$을 유한 생성 $A$-가군들의 집합으로, $K''$을
유한 계수의 자유 $A$-가군들의 집합으로 취하여 적용한다.
$A$가 뇌터라는 가정에 유의하면
(\hyperref[0.11.9.1-fr]{11.9.1})의 조건 (i), (ii), (iii)은 곧바로
확인된다.
\end{proof}

\begin{env}[11.9.3]
\label{0.11.9.3-fr}
\emph{주의}.
\begin{enumerate}
\item[{\rm(i)}] (\hyperref[0.11.9.2-fr]{11.9.2})의 조건 아래에서
$P_i$들이 평탄 오른쪽 $A$-가군이라고 더 가정하자. 그러면 모든 왼쪽
$A$-가군 $M$에 대하여 복합체의 준동형
\[
u\otimes 1:Q_\bullet\otimes_A M\longrightarrow P_\bullet\otimes_A M
\]
은 여전히 호몰로지의 동형
$H_\bullet(Q_\bullet\otimes_A M)\simeq
H_\bullet(P_\bullet\otimes_A M)$을 정의한다. 이는 제III장에서 보게 될
것이다.
\item[{\rm(ii)}] $A$가 뇌터라고 가정하지 않으면
(\hyperref[0.11.9.2-fr]{11.9.2})의 결론은 더 이상 반드시 참인 것은
아니다. 실제로 한 항을 제외하고는 모두 $0$인 복합체에 이를 적용하면,
모든 유한 생성 오른쪽 $A$-가군이 유한 생성 자유 가군들에 의한 분해를
가진다는 결론을 얻게 되지만, 이는 일반적으로 참이 아니다
(cf. Bourbaki, \emph{Alg. comm.}, chap.~I, \S~2, exerc.~6).

그러나 $A$가 뇌터라고 가정하는 대신 $H_i(P_\bullet)$들이 유한
$\infty$-표시를 갖는다고만 가정할 수도 있다(cf. chap.~IV).
\end{enumerate}
\end{env}

\subsection{유한 길이 가군 복합체의 오일러--푸앵카레 특성}
\label{subsection:0.11.10-fr}

\begin{env}[11.10.1]
\label{0.11.10.1-fr}
$A$를 환이라 하자(가환일 필요는 없다). 유한 길이의 왼쪽
$A$-가군들로 이루어진 유한 복합체
\[
\label{0.11.10.1.1-fr}
\tag{11.10.1.1}
M^\bullet:0\longrightarrow M^0\longrightarrow M^1
\longrightarrow\cdots\longrightarrow M^n\longrightarrow0
\]
를 생각하자. 이 복합체의 \emph{오일러--푸앵카레 특성}은 다음 수로
정의된다.
\oldpage[0\textsubscript{III}]{49}
\[
\label{0.11.10.1.2-fr}
\tag{11.10.1.2}
\chi(M^\bullet)=\sum_{i=0}^{n}(-1)^i\operatorname{long}(M^i).
\]
\end{env}

\begin{proposition}[11.10.2]
\label{0.11.10.2-fr}
\emph{유한 길이의 왼쪽 $A$-가군들로 이루어진 모든 유한 복합체
$M^\bullet$에 대하여
\[
\chi(M^\bullet)=\chi(H^\bullet(M^\bullet))
\]
이다. 여기서 $H^\bullet(M^\bullet)$은 영 미분을 갖는 복합체로 본다.
특히 열 (\hyperref[0.11.10.1.1-fr]{11.10.1.1})이 완전하면
$\chi(M^\bullet)=0$이다.}
\end{proposition}

\begin{proof}
간단히 $B^i=B^i(M^\bullet)$, $Z^i=Z^i(M^\bullet)$,
$H^i=H^i(M^\bullet)=Z^i/B^i$로 놓자. $B^i$, $Z^i$, $H^i$는 유한
길이이다. 완전열
\[
0\longrightarrow B^i\longrightarrow Z^i\longrightarrow H^i
\longrightarrow0,
\qquad
0\longrightarrow Z^i\longrightarrow M^i\longrightarrow B^{i+1}
\longrightarrow0
\]
로부터 관계식
\[
\operatorname{long}(Z^i)=\operatorname{long}(H^i)+
\operatorname{long}(B^i),
\qquad
\operatorname{long}(M^i)=\operatorname{long}(Z^i)+
\operatorname{long}(B^{i+1})
\]
을 얻는다. 따라서
\[
\operatorname{long}(M^i)-\operatorname{long}(H^i)=
\operatorname{long}(B^{i+1})+\operatorname{long}(B^i).
\]
이 관계식에 $(-1)^i$를 곱하고 $0\leq i\leq n$에 대하여 합하자.
$B^0=B^{n+1}=0$임에 유의하면 원하는 등식을 얻는다.
\end{proof}

\begin{corollary}[11.10.3]
\label{0.11.10.3-fr}
$E=(E_r^{pq})$를 환 $A$ 위의 가군들의 범주에서의 스펙트럼 열이라 하자.
$E_2^{pq}$들이 유한 길이 $A$-가군이고, $E_2^{pq}\neq0$인 쌍 $(p,q)$는
유한 개뿐이라고 가정하자. 그러면 모든 복합체
$E_r^\bullet=(E_r^{(n)})_{n\in\mathbb Z}$
(\hyperref[0.11.1.1-fr]{11.1.1})의 오일러--푸앵카레 특성
$\chi(E_r^\bullet)$들은 모두 같다. 더 나아가 스펙트럼 열 $E$가
약수렴하고
\[
E_\infty^{(n)}=\bigoplus_{p+q=n}E_\infty^{pq}\quad(n\in\mathbb Z)
\]
로 놓으면 $\chi(E_\infty^\bullet)=\chi(E_2^\bullet)$이기도 하다. 여기서
$E_\infty^\bullet=(E_\infty^{(n)})_{n\in\mathbb Z}$는 영 미분을 갖는
복합체로 본다.
\end{corollary}

\begin{proof}
먼저 $E_2^{pq}=0$이면 $2\leq r<+\infty$에 대하여 $E_r^{pq}=0$임에
유의하자. 따라서 모든 복합체 $E_r^\bullet$은 유한하고 유한 길이
$A$-가군들로 이루어진다. 모든 유한한 $r$에 대한 관계식
$\chi(E_r^\bullet)=\chi(E_{r+1}^\bullet)$은
(\hyperref[0.11.10.2-fr]{11.10.2})와, 영 미분을 갖는 복합체로서의
$H^\bullet(E_r^\bullet)$과 $E_{r+1}^\bullet$ 사이의 동형으로부터
따른다. $E_r^{pq}$들이 유한 길이라는 가정에 따라 모든 쌍 $(p,q)$에
대하여 열들
$(B_k(E_2^{pq}))_{k\geq2}$와 $(Z_k(E_2^{pq}))_{k\geq2}$는 각각 어느
단계부터 일정하다. 따라서 $E$가 약수렴한다는 가정과 유한 개의 쌍
$(p,q)$를 제외하면 $E_2^{pq}=0$이라는 가정에 의해, 모든 쌍 $(p,q)$에
대하여 $E_\infty^{pq}=E_r^{pq}$인 어떤 정수 $r\geq2$가 존재한다. 이로부터
$\chi(E_\infty^\bullet)$에 관한 주장이 따른다.
\end{proof}

\section{층의 코호몰로지에 관한 보충}
\label{section:0.12-fr}

\subsection{환 달린 공간 위의 가군층의 코호몰로지}
\label{subsection:0.12.1-fr}

\begin{env}[12.1.1]
\label{0.12.1.1-fr}
$(X,\mathcal O_X)$를 환 달린 공간이라 하자. 임의의
$\mathcal O_X$-가군 $\mathcal F$에 대하여 코호몰로지
$H^\bullet(X,\mathcal F)$를 정의함을 상기하자. 이는
$\mathcal O_X$-가군들의 범주 $C(X)$에서 아벨군들의 범주로 가는
보편 코호몰로지 함자 (T, 2.2)이며, 왼쪽 완전 함자
$\mathcal F\to\Gamma(X,\mathcal F)$의 유도함자이다. 함자
$H^\bullet(X,\mathcal F)$는
\oldpage[0\textsubscript{III}]{50}
$X$ 위의 \emph{아벨군의 층들}의 범주에서 같은 방식으로 정의되는
코호몰로지 함자 (G, II, 7.2.1)를 범주 $C(X)$로 제한한 것과 동형이다.
\end{env}

\begin{env}[12.1.2]
\label{0.12.1.2-fr}
$A=\Gamma(X,\mathcal O_X)$로 놓자. $A$의 각 원소는 아벨군
$\Gamma(X,\mathcal F)$의 자기준동형을 정의하므로, 함자성에 의해
$H^\bullet(X,\mathcal F)$의 $\delta$-함자의 자기준동형을 정의한다.
이 자기준동형들은 각 $H^p(X,\mathcal F)$에 $A$-가군 구조를 정의하며,
연산자 $\partial$는 $A$-선형이다. 또한 임의의 양의 정수 $p$, $q$와
두 $\mathcal O_X$-가군 $\mathcal F$, $\mathcal G$에 대하여
\emph{컵곱}이라 하는 $A$-가군 준동형
\[
\label{0.12.1.2.1-fr}
\tag{12.1.2.1}
H^p(X,\mathcal F)\otimes_A H^q(X,\mathcal G)
\longrightarrow H^{p+q}(X,\mathcal F\otimes_{\mathcal O_X}\mathcal G)
\]
이 있다 (G, II, 6.6). 이 준동형들에 의해 $H^p(X,\mathcal O_X)$
($p\geq0$)들의 직합 $S$는 반가환 등급 $A$-대수가 되고,
$H^p(X,\mathcal F)$들의 직합은 등급 $S$-가군이 된다.

$X$의 임의의 \emph{열린 덮개} $\mathfrak U$에 대하여, 앞으로
$\check C^\bullet(\mathfrak U,\mathcal F)$는 (G, II, 5.1)과 달리 항상
$\Gamma(U_\alpha,\mathcal F)$를 계수계로 하는 $\mathfrak U$의 신경의
\emph{교대} 공사슬 복합체를 나타낸다. 분명
$\check C^\bullet(\mathfrak U,\mathcal F)$는 등급 $A$-가군이므로,
이 복합체의 코호몰로지군 $\check H^p(\mathfrak U,\mathcal F)$에는
$A$-가군 구조가 주어진다. 또한 표준 사상
\[
\check H^p(\mathfrak U,\mathcal F)\longrightarrow H^p(X,\mathcal F)
\]
(G, II, 5.4)는 $A$-선형이다.
\end{env}

\begin{env}[12.1.3]
\label{0.12.1.3-fr}
$(X',\mathcal O_{X'})$를 두 번째 환 달린 공간이라 하고,
$f=(\psi,\theta)$를 $X'$에서 $X$로 가는 사상이라 하자.

$A'=\Gamma(X',\mathcal O_{X'})$로 놓자. $\psi$-사상 $\theta$는 표준적으로
환 준동형 $A\to A'$을 정의한다. $\mathcal F$를 $\mathcal O_X$-가군,
$\mathcal F'$을 $\mathcal O_{X'}$-가군이라 하자. 임의의 $f$-사상
$u:\mathcal F\to\mathcal F'$
(\hyperref[0.4.4.1-ko]{0\textsubscript{I}, 4.4.1})에 대하여, 모든
$p\geq0$에 대해 \emph{디-준동형}
\[
\label{0.12.1.3.1-fr}
\tag{12.1.3.1}
u_p:H^p(X,\mathcal F)\longrightarrow H^p(X',\mathcal F')
\]
을 정의할 수 있음을 보이겠다. 실제로 $\psi^*$는 $X$ 위의 아벨군의
층들의 범주에서 완전하므로,
$\mathcal F\to H^\bullet(X',\psi^*(\mathcal F))$는 이 범주에서
$\delta$-함자이며, 표준 $\delta$-함자 준동형
\[
\label{0.12.1.3.2-fr}
\tag{12.1.3.2}
H^\bullet(X,\mathcal F)\longrightarrow H^\bullet(X',\psi^*(\mathcal F))
\]
이 있음을 안다. 이는 차수 0에서 표준 준동형
$\Gamma(X,\mathcal F)\to\Gamma(X',\psi^*(\mathcal F))$로 귀착한다는
조건에 의하여 유일하게 정해진다 (T, 3.2.2). 또한 $A$의 각 원소는
$\Gamma(X,\mathcal F)$의 자기준동형 $\mu$와
$\Gamma(X',\psi^*(\mathcal F))$의 자기준동형 $\mu'$을 정하며, 이때 도식
\[
\label{0.12.1.3.3-fr}
\tag{12.1.3.3}
\xymatrix{
\Gamma(X,\mathcal F)\ar[r]\ar[d]_\mu
&\Gamma(X',\psi^*(\mathcal F))\ar[d]^{\mu'}\\
\Gamma(X,\mathcal F)\ar[r]
&\Gamma(X',\psi^*(\mathcal F))
}
\]
은 가환한다. 보편 코호몰로지 함자에서 사상 연장의 유일성
성질 (T, 2.2)에 의해, $\mu$와 $\mu'$에는
(\hyperref[0.12.1.3.3-fr]{12.1.3.3})과 같은 도식들을 가환하게
하는 코호몰로지로의 유일한 연장이 있다. 이는
(\hyperref[0.12.1.3.2-fr]{12.1.3.2})가 $A$-가군 준동형임을 뜻한다.
이제
\[
f^*(\mathcal F)=\psi^*(\mathcal F)
\otimes_{\psi^*(\mathcal O_X)}\mathcal O_{X'}
\]
이고, 따라서 $\psi^*(\mathcal O_X)$-가군 $\psi^*(\mathcal F)$에서
$\mathcal O_{X'}$-가군 $f^*(\mathcal F)$로 가는 표준 디-준동형
$\psi^*(\mathcal F)\to f^*(\mathcal F)$가 있음을 유의하자. 함자성에
의해 이에 대응하는 환들이 $A$와 $A'$인 함자적 디-준동형
\[
\label{0.12.1.3.4-fr}
\tag{12.1.3.4}
H^p(X',\psi^*(\mathcal F))\longrightarrow H^p(X',f^*(\mathcal F))
\]
을 얻는다. 이 디-준동형과
(\hyperref[0.12.1.3.2-fr]{12.1.3.2})를 합성하면 $\mathcal F$에 대해
함자적인 표준 디-준동형
\[
\label{0.12.1.3.5-fr}
\tag{12.1.3.5}
\theta_p:H^p(X,\mathcal F)\longrightarrow H^p(X',f^*(\mathcal F))
\]
을 얻는다. 끝으로 함자성에 의해 준동형
$u^\sharp:f^*(\mathcal F)\to\mathcal F'$에서 $A'$-가군 준동형
$H^p(X',f^*(\mathcal F))\to H^p(X',\mathcal F')$을 얻고, 이를
(\hyperref[0.12.1.3.5-fr]{12.1.3.5})와 합성하면
(\hyperref[0.12.1.3.1-fr]{12.1.3.1})을 얻는다.

$f'=(\psi',\theta'):X''\to X'$을 환 달린 공간의 두 번째 사상이라 하고,
$f''=ff'$을 합성사상이라 하자. 함자 $\psi'^*$와 텐서곱의
교환가능성 (\hyperref[0.4.3.3-ko]{0\textsubscript{I}, 4.3.3})을 고려하면,
디-준동형
$H^p(X',f^*(\mathcal F))\to H^p(X'',f'^*(f^*(\mathcal F)))$와
(\hyperref[0.12.1.3.5-fr]{12.1.3.5})의 합성은 이에 대응하는 디-준동형
$H^p(X,\mathcal F)\to H^p(X'',f''^*(\mathcal F))$임을 곧바로 확인한다.
\end{env}

\begin{env}[12.1.4]
\label{0.12.1.4-fr}
준동형 (\hyperref[0.12.1.3.2-fr]{12.1.3.2})의 직접적인 정의는 다음과
같이 얻을 수 있다. $X$ 위의 아벨군의 층들로 이루어진 $\mathcal F$의
단사 분해 $\mathcal L^\bullet=(\mathcal L^i)$를 생각하자. 함자
$\psi^*$가 완전하므로, $\psi^*(\mathcal L^\bullet)$은 $X'$ 위의 층들로
이루어진 $\psi^*(\mathcal F)$의 분해이다. $\mathcal L'^\bullet=
(\mathcal L'^i)$가 $X'$ 위의 아벨군의 층들의 범주에서
$\psi^*(\mathcal F)$의 단사 분해이면, 증대와 양립하는 아벨군의
층들의 복합체 사상
$\psi^*(\mathcal L^\bullet)\to\mathcal L'^\bullet$이 있으며
(M, V, 1.1.a)), 이는 호모토피를 제외하면 잘 정해진다. 이로부터
아벨군의 복합체 준동형
\[
\Gamma(X,\mathcal L^\bullet)\longrightarrow
\Gamma(X',\psi^*(\mathcal L^\bullet))\longrightarrow
\Gamma(X',\mathcal L'^\bullet)
\]
을 얻는다. 그 합성은 코호몰로지를 취함으로써 $\delta$-함자의 사상
$H^\bullet(X,\mathcal F)\to H^\bullet(X',\psi^*(\mathcal F))$을 준다.
이는 차수 0에서 (\hyperref[0.12.1.3.2-fr]{12.1.3.2})와 일치하므로,
그 사상과 같다 (T, 2.2).

이제 $X$의 \emph{열린 덮개} $\mathfrak U=(U_\alpha)$를 생각하고,
$\mathfrak U$의 역상인 $X'$의 열린 덮개를
$\mathfrak U'=(f^{-1}(U_\alpha))$라 하자. $X$의 모든 열린집합 $V$에
대한 표준 준동형
$\Gamma(V,\mathcal F)\to\Gamma(f^{-1}(V),f^*(\mathcal F))$들은
곧바로 (G, II, 5.1 참조) 복합체 준동형
$\check C^\bullet(\mathfrak U,\mathcal F)\to
\check C^\bullet(\mathfrak U',f^*(\mathcal F))$을 정의하며, 따라서
표준 준동형
\[
\label{0.12.1.4.1-fr}
\tag{12.1.4.1}
\theta_p:\check H^p(\mathfrak U,\mathcal F)
\longrightarrow\check H^p(\mathfrak U',f^*(\mathcal F))
\]
을 얻는다.

\oldpage[0\textsubscript{III}]{52}
또한 다음 도식들은 가환한다.
\[
\label{0.12.1.4.2-fr}
\tag{12.1.4.2}
\xymatrix{
\check H^p(\mathfrak U,\mathcal F)\ar[r]^{\theta_p}\ar[d]
&\check H^p(\mathfrak U',f^*(\mathcal F))\ar[d]\\
H^p(X,\mathcal F)\ar[r]_{\theta_p}
&H^p(X',f^*(\mathcal F))
}
\]
여기서 수직 화살표들은 (G, II, 5.2)의 표준 준동형들이다.
(\hyperref[0.12.1.4.2-fr]{12.1.4.2})의 가환성을 보이기 위하여,
$\Gamma(X,\mathcal C^\bullet(\mathfrak U,\mathcal F))=
\check C^\bullet(\mathfrak U,\mathcal F)$인 (G, II, 5.2)
$\mathfrak U$에 대한 $\mathcal F$의 (교대) 공사슬층의 복합체
$\mathcal C^\bullet(\mathfrak U,\mathcal F)$를 생각하자. 표준 준동형
$\Gamma(V,\mathcal F)\to
\Gamma(f^{-1}(V),\psi^*(\mathcal F))$들은 이때 $\psi$-사상
$\mathcal C^\bullet(\mathfrak U,\mathcal F)\to
\mathcal C^\bullet(\mathfrak U',\psi^*(\mathcal F))$을 정의하며,
위의 표기법으로 다음 가환도식을 얻는다.
\[
\xymatrix{
\Gamma(X,\mathcal C^\bullet(\mathfrak U,\mathcal F))\ar[r]\ar[d]
&\Gamma(X',\mathcal C^\bullet(\mathfrak U',\psi^*(\mathcal F)))\ar[d]\\
\Gamma(X,\mathcal L^\bullet)\ar[r]
&\Gamma(X',\psi^*(\mathcal L^\bullet))\ar[r]
&\Gamma(X',\mathcal L'^\bullet)
}
\]
코호몰로지를 취하면 이는 가환도식
\[
\xymatrix{
\check H^p(\mathfrak U,\mathcal F)\ar[r]\ar[d]
&\check H^p(\mathfrak U',\psi^*(\mathcal F))\ar[d]\\
H^p(X,\mathcal F)\ar[r]
&H^p(X',\psi^*(\mathcal F))
}
\]
을 준다. 여기서 수직 화살표들은 (G, II, 5.2)의 표준 준동형들이다.
이 도식들과 다음 가환도식들을 결합하면 충분하다.
\[
\xymatrix{
\check H^p(\mathfrak U',\psi^*(\mathcal F))\ar[r]\ar[d]
&\check H^p(\mathfrak U',f^*(\mathcal F))\ar[d]\\
H^p(X',\psi^*(\mathcal F))\ar[r]
&H^p(X',f^*(\mathcal F))
}
\]
\oldpage[0\textsubscript{III}]{53}
이 도식들은 준동형 $\psi^*(\mathcal F)\to f^*(\mathcal F)$와
(G, II, 5.2)의 표준 준동형들의 함자성에서 나오며, 이들을 결합하면
(\hyperref[0.12.1.4.2-fr]{12.1.4.2})의 가환성을 얻는다.

$A=\Gamma(X,\mathcal O_X)$, $A'=\Gamma(X',\mathcal O_{X'})$일 때
준동형 (12.1.4.1)은 환 $A$와 $A'$에 대응하는 가군의 디-준동형임을
유의하자. 두 사상의 합성에 대하여
(\hyperref[0.12.1.4.1-fr]{12.1.4.1})은
(\hyperref[0.12.1.3.5-fr]{12.1.3.5})의 추이성과 비슷한
\emph{추이성}을 갖는다. 끝으로 앞의 정의들에서 $\mathcal F$의 단사
분해 $\mathcal L^\bullet$ 대신, 모든 $i$와 모든 $p>0$에 대하여
$H^p(X,\mathcal L^i)=0$인 분해로 시작해도 똑같았을 것임을 유의하자
(G, II, 4.7.1).
\end{env}

\begin{env}[12.1.5]
\label{0.12.1.5-fr}
$\mathcal F$, $\mathcal G$, $\mathcal H$를 세 $\mathcal O_X$-가군이라 하고,
코호몰로지 위에 컵곱 (\hyperref[0.12.1.2.1-fr]{12.1.2.1})에서 유도되는
준동형
\[
\label{0.12.1.5.1-fr}
\tag{12.1.5.1}
H^p(X,\mathcal F)\otimes_A H^q(X,\mathcal G)
\longrightarrow H^{p+q}(X,\mathcal H)
\]
을 주는 $\mathcal O_X$-준동형
$u:\mathcal F\otimes_{\mathcal O_X}\mathcal G\to\mathcal H$를 생각하자.
(\hyperref[0.12.1.3-fr]{12.1.3})의 가정과 표기법 아래 다음 도식들이
가환함을 보이자.
\[
\label{0.12.1.5.2-fr}
\tag{12.1.5.2}
\xymatrix@C=10pt{
H^p(X,\mathcal F)\otimes_A H^q(X,\mathcal G)\ar[r]\ar[d]
&H^{p+q}(X,\mathcal H)\ar[d]\\
H^p(X',f^*(\mathcal F))\otimes_{A'}H^q(X',f^*(\mathcal G))\ar[r]
&H^{p+q}(X',f^*(\mathcal H))
}
\]
여기서 수직 화살표들은 표준 준동형
(\hyperref[0.12.1.3.5-fr]{12.1.3.5})에서 나온다. 이를 위해
(\hyperref[0.12.1.5.1-fr]{12.1.5.1})을 얻는 방법을 상기하자.
$\mathcal F$, $\mathcal G$, $\mathcal H$ 각각의 \emph{표준} 분해
(G, II, 4.3) $\mathcal L^\bullet$, $\mathcal M^\bullet$,
$\mathcal N^\bullet$에서 시작한다. 이들은 $\mathcal O_X$-가군들로
이루어져 있다. $u$에 대응하는 $\mathcal O_X$-가군의 복합체의 선형 사상
$\mathcal L^\bullet\otimes_{\mathcal O_X}\mathcal M^\bullet
\to\mathcal N^\bullet$은 $A$-가군의 복합체 준동형
$\Gamma(X,\mathcal L^\bullet)\otimes_A\Gamma(X,\mathcal M^\bullet)
\to\Gamma(X,\mathcal N^\bullet)$을 주며, 코호몰로지를 취하면 준동형
\[
H^p(\Gamma(X,\mathcal L^\bullet))\otimes_A
H^q(\Gamma(X,\mathcal M^\bullet))
\longrightarrow H^{p+q}(\Gamma(X,\mathcal N^\bullet))
\]
을 얻는다 (G, II, 6.6). 그런데 다음 도식은 분명 가환한다.
\[
\label{0.12.1.5.3-fr}
\tag{12.1.5.3}
\xymatrix@C=8pt{
\Gamma(X,\mathcal L^\bullet)\otimes_A\Gamma(X,\mathcal M^\bullet)
\ar[r]\ar[d]
&\Gamma(X,\mathcal N^\bullet)\ar[d]\\
\Gamma(X',\psi^*(\mathcal L^\bullet))
\otimes_{\Gamma(X',\psi^*(\mathcal O_X))}
\Gamma(X',\psi^*(\mathcal M^\bullet))\ar[r]
&\Gamma(X',\psi^*(\mathcal N^\bullet))
}
\]
\oldpage[0\textsubscript{III}]{54}
코호몰로지를 취하면 이 도식은 다음 가환도식들을 준다.
\[
\label{0.12.1.5.4-fr}
\tag{12.1.5.4}
\xymatrix@C=7pt{
H^p(X,\mathcal F)\otimes_A H^q(X,\mathcal G)\ar[r]\ar[d]
&H^{p+q}(X,\mathcal H)\ar[d]\\
H^p(\Gamma(X',\psi^*(\mathcal L^\bullet)))
\otimes_{\Gamma(X',\psi^*(\mathcal O_X))}
H^q(\Gamma(X',\psi^*(\mathcal M^\bullet)))\ar[r]
&H^{p+q}(\Gamma(X',\psi^*(\mathcal N^\bullet)))
}
\]
한편 $\psi^*(\mathcal L^\bullet)$, $\psi^*(\mathcal M^\bullet)$,
$\psi^*(\mathcal N^\bullet)$은 각각 $\psi^*(\mathcal F)$,
$\psi^*(\mathcal G)$, $\psi^*(\mathcal H)$의 \emph{분해}이므로,
다음 가환도식이 있다 (G, II, 6.6.1).
\[
\label{0.12.1.5.5-fr}
\tag{12.1.5.5}
\xymatrix@C=7pt{
H^p(\Gamma(X',\psi^*(\mathcal L^\bullet)))
\otimes_{\Gamma(X',\psi^*(\mathcal O_X))}
H^q(\Gamma(X',\psi^*(\mathcal M^\bullet)))\ar[r]\ar[d]
&H^{p+q}(\Gamma(X',\psi^*(\mathcal N^\bullet)))\ar[d]\\
H^p(X',\psi^*(\mathcal F))
\otimes_{\Gamma(X',\psi^*(\mathcal O_X))}
H^q(X',\psi^*(\mathcal G))\ar[r]
&H^{p+q}(X',\psi^*(\mathcal H))
}
\]
끝으로 함자성에 의해 다음 가환도식을 얻는다.
\[
\label{0.12.1.5.6-fr}
\tag{12.1.5.6}
\xymatrix@C=7pt{
H^p(X',\psi^*(\mathcal F))
\otimes_{\Gamma(X',\psi^*(\mathcal O_X))}
H^q(X',\psi^*(\mathcal G))\ar[r]\ar[d]
&H^{p+q}(X',\psi^*(\mathcal H))\ar[d]\\
H^p(X',f^*(\mathcal F))\otimes_{A'}H^q(X',f^*(\mathcal G))\ar[r]
&H^{p+q}(X',f^*(\mathcal H))
}
\]
세 도식 (\hyperref[0.12.1.5.4-fr]{12.1.5.4}),
(\hyperref[0.12.1.5.5-fr]{12.1.5.5}),
(\hyperref[0.12.1.5.6-fr]{12.1.5.6})을 결합하면 구하는 가환도식
(\hyperref[0.12.1.5.2-fr]{12.1.5.2})를 얻는다.
\end{env}

\begin{remark}[12.1.6]
\label{0.12.1.6-fr}
\oldpage[0\textsubscript{III}]{55}
(\hyperref[0.12.1.3-fr]{12.1.3})의 표기법 아래 다음 가환도식이 있다고
가정하자.
\[
\label{0.12.1.6.1-fr}
\tag{12.1.6.1}
\xymatrix{
0\ar[r]&\mathcal F\ar[r]^r\ar[d]_u
&\mathcal G\ar[r]^s\ar[d]_v
&\mathcal H\ar[r]\ar[d]_w&0\\
0\ar[r]&\mathcal F'\ar[r]_{r'}
&\mathcal G'\ar[r]_{s'}
&\mathcal H'\ar[r]&0
}
\]
여기서 $r$, $s$는 $\mathcal O_X$-가군 준동형, $r'$, $s'$은
$\mathcal O_{X'}$-가군 준동형, $u$, $v$, $w$는 $f$-사상이고 두 행은
\emph{완전}하다. 그러면 가환도식
\[
\label{0.12.1.6.2-fr}
\tag{12.1.6.2}
\xymatrix@C=6pt{
\cdots\ar[r]&H^p(X,\mathcal F)\ar[r]\ar[d]_{u_p}
&H^p(X,\mathcal G)\ar[r]\ar[d]_{v_p}
&H^p(X,\mathcal H)\ar[r]^\partial\ar[d]_{w_p}
&H^{p+1}(X,\mathcal F)\ar[r]\ar[d]_{u_{p+1}}&\cdots\\
\cdots\ar[r]&H^p(X',\mathcal F')\ar[r]
&H^p(X',\mathcal G')\ar[r]
&H^p(X',\mathcal H')\ar[r]_\partial
&H^{p+1}(X',\mathcal F')\ar[r]&\cdots
}
\]
을 얻는다. 실제로 (\hyperref[0.12.1.6.1-fr]{12.1.6.1})은 다음과 같이
인수분해된다.
\[
\xymatrix@R=14pt{
0\ar[r]&\mathcal F\ar[r]\ar[d]
&\mathcal G\ar[r]\ar[d]
&\mathcal H\ar[r]\ar[d]&0\\
0\ar[r]&\psi^*(\mathcal F)\ar[r]\ar[d]
&\psi^*(\mathcal G)\ar[r]\ar[d]
&\psi^*(\mathcal H)\ar[r]\ar[d]&0\\
0\ar[r]&\mathcal F'\ar[r]
&\mathcal G'\ar[r]
&\mathcal H'\ar[r]&0
}
\]
여기서 가운데 행은 \emph{완전}하다
(\hyperref[0.3.7.2-ko]{0\textsubscript{I}, 3.7.2}). 이제
(\hyperref[0.12.1.3.2-fr]{12.1.3.2})가 $\delta$-함자의 준동형이라는
사실과 $H^p(X',\mathcal F')$들이 $\mathcal F'$에 대한
$\delta$-함자를 이룬다는 사실을 쓰면 충분하다.
\end{remark}

\begin{env}[12.1.7]
\label{0.12.1.7-fr}
가정과 표기법을 (\hyperref[0.12.1.3-fr]{12.1.3})과 같이 하고,
이제 $\mathcal F=f_*(\mathcal F')=\psi_*(\mathcal F')$인 경우를 생각하자.
(\hyperref[0.12.1.3-fr]{12.1.3})에서 정의한 디-준동형
\[
\label{0.12.1.7.1-fr}
\tag{12.1.7.1}
H^p(X,f_*(\mathcal F'))\longrightarrow H^p(X',\mathcal F')
\]
은 $H^p(X',\mathcal F')$의 어떤 자기동형의 차이를 제외하면 합성함자
$\mathcal F'\leadsto\Gamma(X,\psi_*(\mathcal F'))$의 스펙트럼 열의
\emph{가장자리 준동형}으로 얻을 수 있음을 보이겠다 (T, 2.4).
(\hyperref[0.12.1.4-fr]{12.1.4})에서 준동형
(\hyperref[0.12.1.7.1-fr]{12.1.7.1})에 대해 준 설명은 여기서 이 준동형을
다음과 같이 얻을 수 있음을 보인다. $\psi_*(\mathcal F')$와
$\mathcal F'$의 단사 분해 $\mathcal L^\bullet$과
$\mathcal L'^\bullet$을 각각 택한 뒤, 표준 준동형
$\psi^*(\psi_*(\mathcal F'))\to\mathcal F'$의 ``위에 놓인'' 복합체 준동형
$v:\psi^*(\mathcal L^\bullet)\to\mathcal L'^\bullet$을 택한다.
\oldpage[0\textsubscript{III}]{56}
이어서 $\Gamma(X',\mathcal L'^\bullet)=
\Gamma(X,\psi_*(\mathcal L'^\bullet))$임을 유의하면, 합성 준동형
\[
\Gamma(X,\mathcal L^\bullet)\longrightarrow
\Gamma(X',\psi^*(\mathcal L^\bullet))
\xrightarrow{\Gamma(v)}\Gamma(X',\mathcal L'^\bullet)
\]
은 다름 아닌
\[
\label{0.12.1.7.2-fr}
\tag{12.1.7.2}
\Gamma(v^\flat):\Gamma(X,\mathcal L^\bullet)
\longrightarrow\Gamma(X,\psi_*(\mathcal L'^\bullet))
\]
이다 (\hyperref[0.3.7.1-ko]{0\textsubscript{I}, 3.7.1}). 그리고
(\hyperref[0.12.1.7.2-fr]{12.1.7.2})에서 코호몰로지를 취하면
(\hyperref[0.12.1.7.1-fr]{12.1.7.1})을 얻는다. 한편 합성함자
$\mathcal F'\leadsto\Gamma(X,\psi_*(\mathcal F'))$의 스펙트럼 열들은
$X$ 위의 아벨군의 층들의 범주에서 복합체
$\psi_*(\mathcal L'^\bullet)$의 단사 카르탕--아일렌베르크 분해
$\mathcal M^{\bullet\bullet}=(\mathcal M^{ij})$를 생각하여 얻는다.
문제의 스펙트럼 열들은 이중복합체
$\Gamma(X,\mathcal M^{\bullet\bullet})$의 스펙트럼 열들이다
($i<0$ 또는 $j<0$이면 $\mathcal M^{ij}=0$이므로 이들은 쌍정칙이다).
그런데 이 이중복합체의 첫 번째 스펙트럼 열은 \emph{퇴화}한다.
실제로 층 $\psi_*(\mathcal L'^i)$들이 플라스크이므로 (G, II, 3.1.1),
$q>0$이면
$H_{\mathrm{II}}^q(\Gamma(X,\mathcal M^{i,\bullet}))
=H^q(\psi_*(\mathcal L'^i))=0$이다 (G, II, 4.4.3). 따라서 전단사인
가장자리 준동형 (\hyperref[0.11.1.6-fr]{11.1.6})
\[
\label{0.12.1.7.3-fr}
\tag{12.1.7.3}
'E_2^{i0}=H^i(H_{\mathrm{II}}^0(\Gamma(X,\mathcal M^{\bullet\bullet})))
\longrightarrow H^i(\Gamma(X,\mathcal M^{\bullet\bullet}))
\]
이 있다. 이 준동형은 증대
\[
\label{0.12.1.7.4-fr}
\tag{12.1.7.4}
\Gamma(X,\psi_*(\mathcal L'^\bullet))
\longrightarrow\Gamma(X,\mathcal M^{\bullet\bullet})
\]
에서 코호몰로지를 취하여 나온다는 것을 안다
(\hyperref[0.11.3.4-fr]{11.3.4}). 이 증대 자체는 증대
$\eta:\psi_*(\mathcal L'^\bullet)\to\mathcal M^{\bullet\bullet}$에서 나온다.
한편 두 번째 스펙트럼 열에는 가장자리 준동형
\[
\label{0.12.1.7.5-fr}
\tag{12.1.7.5}
''E_2^{i0}=H^i(H_{\mathrm I}^0(\Gamma(X,\mathcal M^{\bullet\bullet})))
\longrightarrow H^i(\Gamma(X,\mathcal M^{\bullet\bullet}))
\]
이 있다. 이는 (\hyperref[0.11.3.4-fr]{11.3.4}) 복합체 준동형
$Z_{\mathrm I}^0(\Gamma(X,\mathcal M^{\bullet\bullet}))
\to\Gamma(X,\mathcal M^{\bullet\bullet})$에서 코호몰로지를 취하여 나온다.
그런데 $\psi_*$가 왼쪽 완전하므로 열
\[
0\longrightarrow\psi_*(\mathcal F')\longrightarrow
\psi_*(\mathcal L'^0)\longrightarrow\psi_*(\mathcal L'^1)
\]
은 완전하다. 따라서 카르탕--아일렌베르크 분해의 정의
(\hyperref[0.11.4.2-fr]{11.4.2})에 의해
$B_{\mathrm I}^0(\mathcal M^{\bullet\bullet})=0$,
$Z_{\mathrm I}^0(\mathcal M^{\bullet\bullet})=\mathcal L^\bullet$로
택할 수 있다. 도식
\[
\xymatrix@C=50pt@R=24pt{
\mathcal L^0\ar[r]^{i^0}&\mathcal M^{00}\\
\psi_*(\mathcal F')\ar[u]^\varepsilon
\ar[r]_{\varepsilon'}\ar[ur]^{\varepsilon''}
&\psi_*(\mathcal L'^0)\ar[u]_{\eta^0}
}
\]
은 가환하므로, 복합체 단사사상
$i:\mathcal L^\bullet\to\mathcal M^{\bullet\bullet}$은 증대
$\varepsilon$ 및 $\varepsilon''$과 양립한다. 이로써
$\mathcal L^\bullet$에서 $\mathcal M^{\bullet\bullet}$로 가는 두 복합체
준동형
\[
\xymatrix@C=52pt@R=22pt{
\mathcal L^\bullet\ar[r]^i\ar[dr]_{v^\flat}
&\mathcal M^{\bullet\bullet}\\
&\psi_*(\mathcal L'^\bullet)\ar[u]_\eta
}
\]
\oldpage[0\textsubscript{III}]{57}
을 얻으며, 이들은 증대 $\varepsilon$ 및 $\varepsilon''$과 양립한다.
$\mathcal L^\bullet$은 \emph{단사 분해}이고
$\mathcal M^{\bullet\bullet}$은 단사층들로 이루어져 있으므로,
(M, V, 1.1a))에 의해 이 두 준동형은 \emph{서로 호모토픽}이다.
따라서 이에 대응하는 두 준동형
$\Gamma(X,\mathcal L^\bullet)\to
\Gamma(X,\mathcal M^{\bullet\bullet})$도 서로 호모토픽이고,
코호몰로지를 취하면 \emph{같은} 준동형을 얻는다. 다시 말해
(\hyperref[0.12.1.7.5-fr]{12.1.7.5})의 가장자리 준동형, 즉
$H^p(X,\psi_*(\mathcal F'))\to
H^p(\Gamma(X,\mathcal M^{\bullet\bullet}))$은
(\hyperref[0.12.1.7.1-fr]{12.1.7.1})과
(\hyperref[0.12.1.7.3-fr]{12.1.7.3})의 합성임을 보였다. 여기서 후자는
$H^p(X',\mathcal F')\to
H^p(\Gamma(X,\mathcal M^{\bullet\bullet}))$로 쓰이며 앞에서
\emph{동형}임을 보았다. 이로부터 주장이 따른다.
\end{env}

\subsection{고차 직상}
\label{subsection:0.12.2-fr}

\begin{env}[12.2.1]
\label{0.12.2.1-fr}
$(X,\mathcal O_X)$, $(Y,\mathcal O_Y)$를 두 환 달린 공간이라 하고,
$f=(\psi,\theta)$를 $X$에서 $Y$로 가는 사상이라 하자. 이 사상은
직상 함자
\[
f_*:C(X)\longrightarrow C(Y)
\]
를 정의한다. 이는 사실 $X$ 위의 아벨군층들의 범주에서 정의된 함자
$\psi_*$를 $C(X)$에 제한한 것과 같다. 뒤의 함자는 가법이고 왼쪽
완전하다. $X$ 위의 모든 아벨군층은 어떤 단사 아벨군층의 부분층과
동형이므로, 함자 $\psi_*$의 오른쪽 유도함자
$\mathcal F\leadsto R^p\psi_*(\mathcal F)$가 정의된다. 각
$R^p\psi_*(\mathcal F)$는 $Y$ 위의 아벨군층이고, $R^p\psi_*$들은
보편 코호몰로지 함자를 이룬다 (T, 2.3).

또한 층 $R^p\psi_*(\mathcal F)$는 전층
$V\leadsto H^p(f^{-1}(V),\mathcal F)$에 연관된 층이다 (T, 3.7.2).
이제 $\mathcal F$가 $\mathcal O_X$-가군이라고 가정하면,
$H^p(f^{-1}(V),\mathcal F)$는 자연스럽게
$\Gamma(f^{-1}(V),\mathcal O_X)$-가군, 따라서
$\Gamma(V,\psi_*(\mathcal O_X))$-가군의 구조를 갖는다. 준동형
$\theta:\mathcal O_Y\to\psi_*(\mathcal O_X)$로부터 다시
$\Gamma(V,\mathcal O_Y)$-가군의 구조를 얻는다. 이렇게 정의한 구조에
대하여 열린집합 $V$에서 열린집합 $V'\subset V$로의 제한은
디-준동형을 정의함이 명백하다. 따라서 각 $R^p\psi_*(\mathcal F)$에
$\mathcal O_Y$-가군 구조를 정의할 수 있다. 이 $\mathcal O_Y$-가군을
$R^p f_*(\mathcal F)$로 나타내며, 이로써 $R^p f_*$는 $C(X)$에서
$C(Y)$로 가는 가법 함자로 정의된다. 더 나아가 $R^p f_*$들은
$\partial$-함자를 이룬다. 실제로
\[
0\longrightarrow\mathcal F'\longrightarrow\mathcal F
\longrightarrow\mathcal F''\longrightarrow0
\]
이 $\mathcal O_X$-가군들의 완전열이면, 위에서 준
$R^p\psi_*$와 $R^p\psi_*(\mathcal F)$ 위의 $\mathcal O_Y$-가군
구조의 설명으로부터 준동형
\[
\partial:R^p\psi_*(\mathcal F'')\longrightarrow
R^{p+1}\psi_*(\mathcal F')
\]
이 $\mathcal O_Y$-가군의 준동형임이 곧바로 따른다. 끝으로
$R^p f_*$는 $f_*$의 오른쪽 유도함자와 동일시된다. 실제로 모든
$\mathcal O_X$-가군은 $\mathcal O_X$-가군들로 이루어진 단사 분해를
갖고, 이러한 분해는 아벨군의 플라스크층들로 이루어지므로
(G, II, 7.1), $R^p\psi_*(\mathcal F)$를 계산하는 데 쓸 수 있다.
모든 플라스크층 $\mathcal G$와 모든 $n\geq1$에 대하여
$R^n\psi_*(\mathcal G)=0$이기 때문이다 (T, 2.4.1, 주석 3 및 명제
3.3.2의 따름정리). 따라서 $R^p f_*$들은 $C(X)$에서 $C(Y)$로 가는
보편 코호몰로지 함자를 이룬다 (T, 2.3).
\end{env}

\begin{env}[12.2.2]
\label{0.12.2.2-fr}
$\mathcal F$, $\mathcal G$를 두 $\mathcal O_X$-가군이라 하자.
(\hyperref[0.12.2.1-fr]{12.2.1})의 표기 아래 $Y$의 모든 열린집합
$V$에 대하여 컵곱 준동형 (\hyperref[0.12.1.2.1-fr]{12.1.2.1})
\[
H^p(f^{-1}(V),\mathcal F)
\otimes_{\Gamma(f^{-1}(V),\mathcal O_X)}
H^q(f^{-1}(V),\mathcal G)
\longrightarrow
H^{p+q}(f^{-1}(V),\mathcal F\otimes_{\mathcal O_X}\mathcal G)
\]
이 있다.
\oldpage[0\textsubscript{III}]{58}
컵곱의 정의 (G, II, 6.6)에서 이 준동형들이 $V$에서 그 열린
부분공간 $V'$로의 이행과 가환함이 곧바로 따른다. 한편 $\theta$에서
오는 환 준동형
\[
\Gamma(V,\mathcal O_Y)\longrightarrow\Gamma(V,\psi_*(\mathcal O_X))
=\Gamma(f^{-1}(V),\mathcal O_X)
\]
이 있으므로, 텐서곱의 표준 준동형
\[
\begin{aligned}
&H^p(f^{-1}(V),\mathcal F)\otimes_{\Gamma(V,\mathcal O_Y)}
H^q(f^{-1}(V),\mathcal G)\\
&\qquad\longrightarrow
H^p(f^{-1}(V),\mathcal F)
\otimes_{\Gamma(f^{-1}(V),\mathcal O_X)}
H^q(f^{-1}(V),\mathcal G)
\end{aligned}
\]
을 얻으며, 이것도 $V$에서 $V'$로의 제한과 양립한다. 따라서 합성으로
$\Gamma(V,\mathcal O_Y)$-가군의 준동형을 얻고, 고려하는 전층들에
연관된 층에 대하여 $\mathcal F$, $\mathcal G$에 함자적인 표준 준동형
\[
\label{0.12.2.2.1-fr}
\tag{12.2.2.1}
R^p f_*(\mathcal F)\otimes_{\mathcal O_Y}R^q f_*(\mathcal G)
\longrightarrow
R^{p+q}f_*(\mathcal F\otimes_{\mathcal O_X}\mathcal G)
\]
을 정의한다. $p=q=0$일 때 이 준동형은
(\hyperref[0.4.2.2.1-ko]{0\textsubscript{I}, 4.2.2.1})로 환원된다.
\end{env}

\begin{proposition}[12.2.3]
\label{0.12.2.3-fr}
\emph{모든 $\mathcal O_X$-가군 $\mathcal F$와 유한 계수 국소 자유
$\mathcal O_Y$-가군 $\mathcal L$에 대하여 함자적인 표준 동형
\[
\label{0.12.2.3.1-fr}
\tag{12.2.3.1}
R^p f_*(\mathcal F)\otimes_{\mathcal O_Y}\mathcal L
\xrightarrow{\sim}
R^p f_*(\mathcal F\otimes_{\mathcal O_X}f^*(\mathcal L))
\]
이 성립한다.}
\end{proposition}

\begin{proof}
준동형 (\hyperref[0.12.2.3.1-fr]{12.2.3.1})은
(\hyperref[0.12.2.2.1-fr]{12.2.2.1})의 특수한 경우인 준동형
\[
\label{0.12.2.3.2-fr}
\tag{12.2.3.2}
R^p f_*(\mathcal F)\otimes_{\mathcal O_Y}f_*(f^*(\mathcal L))
\longrightarrow
R^p f_*(\mathcal F\otimes_{\mathcal O_X}f^*(\mathcal L))
\]
과, 표준 준동형
(\hyperref[0.4.4.3.2-ko]{0\textsubscript{I}, 4.4.3.2})에서 오는
(\hyperref[0.12.2.3.1-fr]{12.2.3.1})의 좌변에서
(\hyperref[0.12.2.3.2-fr]{12.2.3.2})의 좌변으로 가는 준동형을
합성하여 얻는다. $\mathcal L$이 국소 자유일 때
(\hyperref[0.12.2.3.1-fr]{12.2.3.1})이 동형임을 확인하려면 곧바로
$\mathcal L=\mathcal O_Y$인 경우로 환원할 수 있다. 문제는 $Y$ 위에서
국소적이고, 고려하는 함자들은 $\mathcal L$에 관하여 가법이기 때문이다.
그런데 이 경우에는 (\hyperref[0.12.2.2.1-fr]{12.2.2.1})의 정의에 따라,
연관된 전층의 준동형이 전단사임을 확인하면 충분하고, 이는 관계
$f^*(\mathcal O_Y)=\mathcal O_X$에서 즉시 따른다.
\end{proof}

\begin{env}[12.2.4]
\label{0.12.2.4-fr}
$(Z,\mathcal O_Z)$를 셋째 환 달린 공간, $g:Y\to Z$를 환 달린 공간의
사상이라 하자. 모든 단사 $\mathcal O_X$-가군 $\mathcal G$에 대하여
$f_*(\mathcal G)$가 아벨군의 플라스크층임을 알고 있다
(G, II, 7.1 및 3.1.1). 따라서
(\hyperref[0.12.2.1-fr]{12.2.1})에 의해 모든 $p>0$에 대하여
$R^p g_*(f_*(\mathcal G))=0$이다. 이에 따라 (T, 2.4.1) 합성함자의
르레 스펙트럼 열을 합성함자 $g_*f_*$에 적용할 수 있다. 즉
$h=g\circ f$일 때 수렴대상이 함자 $R^\bullet h_*$이고 $E_2$항이
\[
\label{0.12.2.4.1-fr}
\tag{12.2.4.1}
E_2^{pq}=R^p g_*(R^q f_*(\mathcal F))
\]
로 주어지는 쌍정칙 스펙트럼 열이 있다.
\end{env}

\begin{env}[12.2.5]
\label{0.12.2.5-fr}
(\hyperref[0.12.2.4-fr]{12.2.4})의 조건 아래 $\mathcal O_Z$-가군의
표준 준동형
\[
\label{0.12.2.5.1-fr}
\tag{12.2.5.1}
R^n g_*(f_*(\mathcal F))\longrightarrow R^n h_*(\mathcal F)
\]
\[
\label{0.12.2.5.2-fr}
\tag{12.2.5.2}
R^n h_*(\mathcal F)\longrightarrow g_*(R^n f_*(\mathcal F))
\]
을 직접 정의하자.
\oldpage[0\textsubscript{III}]{59}
이들은 르레 스펙트럼 열의 ``가장자리 준동형''과 동일시할 수 있다
(\hyperref[0.12.1.7-fr]{12.1.7} 참조). 고차 직상층들에 연관되는
전층들 위에서 작업하면 충분하다
(\hyperref[0.12.2.1-fr]{12.2.1}). 이를 위하여 $Z$의 임의의 열린집합
$W$와 $Y$에서의 그 역상 $g^{-1}(W)$를 생각하자. 대응하는 환들이
$\Gamma(g^{-1}(W),\mathcal O_Y)$와
$\Gamma(h^{-1}(W),\mathcal O_X)$일 때 표준 디-준동형
\[
\label{0.12.2.5.3-fr}
\tag{12.2.5.3}
H^n(g^{-1}(W),f_*(\mathcal F))\longrightarrow
H^n(f^{-1}(g^{-1}(W)),f^*(f_*(\mathcal F)))
\]
이 있다. 한편 표준 준동형
(\hyperref[0.4.4.3.3-ko]{0\textsubscript{I}, 4.4.3.3})은 함자성에 따라
표준 준동형
\[
\label{0.12.2.5.4-fr}
\tag{12.2.5.4}
H^n(h^{-1}(W),f^*(f_*(\mathcal F)))\longrightarrow
H^n(h^{-1}(W),\mathcal F)
\]
을 주며, 이는 $\Gamma(h^{-1}(W),\mathcal O_X)$-가군의 준동형이다.
환 준동형
$\Gamma(W,\mathcal O_Z)\to\Gamma(g^{-1}(W),\mathcal O_Y)$을 고려하면,
(\hyperref[0.12.2.5.4-fr]{12.2.5.4})와
(\hyperref[0.12.2.5.3-fr]{12.2.5.3})의 합성으로 전층의 준동형을 얻고,
이는 층의 준동형 (\hyperref[0.12.2.5.1-fr]{12.2.5.1})을 준다.

(\hyperref[0.12.2.5.2-fr]{12.2.5.2})의 정의는 더 간단하다. 정의에
따라 $R^n h_*(\mathcal F)$는 전층
$W\leadsto H^n(f^{-1}(g^{-1}(W)),\mathcal F)$에 연관되고,
$R^n f_*(\mathcal F)$는 전층
$V\leadsto H^n(f^{-1}(V),\mathcal F)$에 연관된다. 따라서 표준 준동형
\[
H^n(f^{-1}(g^{-1}(W)),\mathcal F)\longrightarrow
\Gamma(g^{-1}(W),R^n f_*(\mathcal F))
\]
이 있고, 이 준동형들이 전층의 준동형을 정의함은 즉시 알 수 있다.
이 전층의 준동형이 다시
(\hyperref[0.12.2.5.2-fr]{12.2.5.2})를 정의한다.
\end{env}

\begin{env}[12.2.6]
\label{0.12.2.6-fr}
(\hyperref[0.12.2.4-fr]{12.2.4})의 가정 아래 $\mathcal F$,
$\mathcal G$, $\mathcal K$를 세 $\mathcal O_X$-가군이라 하고,
$u:\mathcal F\otimes\mathcal G\to\mathcal K$를
$\mathcal O_X$-준동형이라 하자. 그러면 다음 두 도식은 가환한다.
\[
\xymatrix@C=10pt@R=24pt{
R^p g_*(f_*(\mathcal F))\otimes_{\mathcal O_Z}R^q g_*(f_*(\mathcal G))
\ar[r]\ar[d]&R^{p+q}g_*(f_*(\mathcal K))\ar[d]\\
R^p h_*(\mathcal F)\otimes_{\mathcal O_Z}R^q h_*(\mathcal G)
\ar[r]&R^{p+q}h_*(\mathcal K)
}
\tag{12.2.6.1}\label{0.12.2.6.1-fr}
\]
그리고
\[
\xymatrix@C=10pt@R=24pt{
R^p h_*(\mathcal F)\otimes_{\mathcal O_Z}R^q h_*(\mathcal G)
\ar[r]\ar[d]&R^{p+q}h_*(\mathcal K)\ar[d]\\
g_*(R^p f_*(\mathcal F))\otimes_{\mathcal O_Z}g_*(R^q f_*(\mathcal G))
\ar[r]&g_*(R^{p+q}f_*(\mathcal K))
}
\tag{12.2.6.2}\label{0.12.2.6.2-fr}
\]
\oldpage[0\textsubscript{III}]{60}
여기서 수평 화살표들은
(\hyperref[0.12.2.2.1-fr]{12.2.2.1})에서 오며, 마지막 수평 화살표에는
(\hyperref[0.4.2.2.1-ko]{0\textsubscript{I}, 4.2.2.1})도 합성하였다.
수직 화살표들은 각각 준동형
(\hyperref[0.12.2.5.1-fr]{12.2.5.1})과
(\hyperref[0.12.2.5.2-fr]{12.2.5.2})에서 온다.

실제로 대응하는 전층의 준동형들에 대하여 확인하면 충분하다. 이
준동형들에 관하여 (\hyperref[0.12.2.2-fr]{12.2.2})와
(\hyperref[0.12.2.5-fr]{12.2.5})에서 준 정의로 돌아가면,
(\hyperref[0.12.2.6.1-fr]{12.2.6.1})의 경우 곧바로 가환도식들
(\hyperref[0.12.1.5.2-fr]{12.1.5.2})로 환원된다.
(\hyperref[0.12.2.6.2-fr]{12.2.6.2})의 확인은 더욱 간단하다.
\end{env}

\subsection{층의 $\operatorname{Ext}$ 함자에 관한 보충}
\label{subsection:0.12.3-fr}

\begin{env}[12.3.1]
\label{0.12.3.1-fr}
환 달린 공간 $(X,\mathcal O_X)$를 생각하자. $\mathcal O_X$-가군의
범주에서 $\Gamma(X,\mathcal O_X)$-가군의 범주로 가는 쌍함자
$\operatorname{Ext}_{\mathcal O_X}^p(X;\mathcal F,\mathcal G)$와,
$\mathcal O_X$-가군의 범주에서 그 자신으로 가는 쌍함자
$\mathcal{E}xt_{\mathcal O_X}^p(\mathcal F,\mathcal G)$의 정의와 주요
성질들, 그리고 이들을 잇는 쌍정칙 스펙트럼 열
$E(\mathcal F,\mathcal G)$에 대해서는 다시 논하지 않는다
(T, 4.2 및 G, II, 7.3).
\end{env}

\begin{env}[12.3.2]
\label{0.12.3.2-fr}
(M, XIV, 1)에서와 같은 방식으로 $\mathcal O_X$-가군 $\mathcal F$의
$\mathcal O_X$-가군 $\mathcal G$에 의한 확장이라는 개념과 동치인
확장류들 사이의 합성법칙을 정의한다. 실제로 가군에 대하여 한 논증은
임의의 아벨 범주에 명백하게 적용된다. 단사 대상에의 매장이 존재한다는
사실만을 쓰는 (M, XIV, 1.1)의 둘째 증명은 $\mathcal O_X$-가군의
범주에서도 여전히 유효하다. 따라서
$\operatorname{Ext}_{\mathcal O_X}^1(X;\mathcal F,\mathcal G)$는
$\mathcal F$의 $\mathcal G$에 의한 확장류들의 아벨군과 표준적으로
동일시된다.
\end{env}

\begin{proposition}[12.3.3]
\label{0.12.3.3-fr}
\emph{$(X,\mathcal O_X)$를 환층 $\mathcal O_X$가 연접인 환 달린
공간이라 하자. 임의의 두 연접 $\mathcal O_X$-가군 $\mathcal F$,
$\mathcal G$와 모든 $p\geq0$에 대하여
$\mathcal{E}xt_{\mathcal O_X}^p(\mathcal F,\mathcal G)$는 연접
$\mathcal O_X$-가군이다.}
\end{proposition}

\begin{proof}
$\mathcal{E}xt_{\mathcal O_X}^p(\mathcal F,\mathcal G)$들은
$\mathcal F$에 관하여 반변 코호몰로지 함자를 이룸에 유의하자.
$\mathcal F$가 연접이므로, 모든 $p$와 모든 점 $x\in X$에 대하여
$x$의 열린 근방 $U$와 $(\mathcal O_X|U)$-가군들의 완전열
\[
0\longrightarrow\mathcal R\longrightarrow\mathcal L_{p-1}
\longrightarrow\cdots\longrightarrow\mathcal L_0
\longrightarrow\mathcal F|U\longrightarrow0
\]
이 존재하여, 각 $\mathcal L_i$ ($0\leq i\leq p-1$)는 어떤
$\mathcal O_X^{n_i}|U$와 동형이고 $\mathcal R$은 연접이다. 이는
$\mathcal O_X$가 연접이라는 가정 아래
(\hyperref[0.5.3.2-ko]{0\textsubscript{I}, 5.3.2})와
(\hyperref[0.5.3.4-ko]{0\textsubscript{I}, 5.3.4})를 $p$에 관하여
귀납적으로 적용하여 얻는다.

이제 $p\geq1$일 때 $\mathcal L|U$가 어떤 $\mathcal O_X^n|U$와
동형인 모든 $\mathcal O_X$-가군 $\mathcal L$에 대하여
\[
\mathcal{E}xt_{\mathcal O_X|U}^p(\mathcal L|U,\mathcal G|U)=0
\]
임을 유의하자 (T, 4.2.3). 따라서 (M, V, 7.2)의 논증을 반변
코호몰로지 함자
\[
\mathcal F\leadsto
\mathcal{H}om_{\mathcal O_X|U}(\mathcal F|U,\mathcal G|U)
\]
에 적용할 수 있고, 완전열
\[
\mathcal{H}om_{\mathcal O_X|U}(\mathcal L_{p-1},\mathcal G|U)
\longrightarrow
\mathcal{H}om_{\mathcal O_X|U}(\mathcal R,\mathcal G|U)
\longrightarrow
\mathcal{E}xt_{\mathcal O_X|U}^p(\mathcal F|U,\mathcal G|U)
\longrightarrow0
\]
을 얻는다. 이 열의 처음 두 항은 연접 $(\mathcal O_X|U)$-가군이므로
(\hyperref[0.5.3.5-ko]{0\textsubscript{I}, 5.3.5}), 셋째 항도 연접이다
(\hyperref[0.5.3.4-ko]{0\textsubscript{I}, 5.3.4}).
\end{proof}

\oldpage[0\textsubscript{III}]{61}
\begin{proposition}[12.3.4]
\label{0.12.3.4-fr}
\emph{$f:X\to Y$를 환 달린 공간의 평탄 사상이라 하고,
$\mathcal F$, $\mathcal G$를 두 $\mathcal O_Y$-가군이라 하자.}
\begin{enumerate}
\item[{\rm(i)}] \emph{코호몰로지 쌍함자들의 준동형
\[
\label{0.12.3.4.1-fr}
\tag{12.3.4.1}
f^*(\mathcal{E}xt_{\mathcal O_Y}^{\bullet}(\mathcal F,\mathcal G))
\longrightarrow
\mathcal{E}xt_{\mathcal O_X}^{\bullet}(f^*(\mathcal F),f^*(\mathcal G))
\]
이 존재하며, 차수 0에서는 표준 준동형
(\hyperref[0.4.4.6-ko]{0\textsubscript{I}, 4.4.6})으로 환원된다.}

\item[{\rm(ii)}] \emph{스펙트럼 열의 표준 사상
\[
\label{0.12.3.4.2-fr}
\tag{12.3.4.2}
E(\mathcal F,\mathcal G)\longrightarrow
E(f^*(\mathcal F),f^*(\mathcal G))
\]
이 존재하며, $E_2$항에서는
(\hyperref[0.12.3.4.1-fr]{12.3.4.1})과
(\hyperref[0.12.1.3.1-fr]{12.1.3.1})에서 얻는 준동형
\[
\label{0.12.3.4.3-fr}
\tag{12.3.4.3}
H^p(Y,\mathcal{E}xt_{\mathcal O_Y}^q(\mathcal F,\mathcal G))
\longrightarrow
H^p(X,\mathcal{E}xt_{\mathcal O_X}^q(f^*(\mathcal F),f^*(\mathcal G)))
\]
으로 환원된다.}
\end{enumerate}
\end{proposition}

\begin{proof}
\begin{enumerate}
\item[{\rm(i)}]
$f^*$는 $\mathcal O_Y$-가군의 범주에서 완전함자이므로
(\hyperref[0.6.7.2-ko]{0\textsubscript{I}, 6.7.2}), 함자들
\[
\mathcal G\leadsto
f^*(\mathcal{H}om_{\mathcal O_Y}(\mathcal F,\mathcal G))
\quad\hbox{및}\quad
\mathcal G\leadsto
\mathcal{H}om_{\mathcal O_X}(f^*(\mathcal F),f^*(\mathcal G))
\]
은 왼쪽 완전하다. 따라서
(\hyperref[0.4.4.6-ko]{0\textsubscript{I}, 4.4.6})에서 이들의 유도함자
사이를 잇는 준동형을 표준적으로 얻는다.

이를 계산하기 위하여 $\mathcal G$의 단사 분해
$\mathcal L^\bullet=(\mathcal L^i)$을 취하면, 층의 복합체들의
코호몰로지 사상
\[
\mathcal H^p(f^*(\mathcal{H}om_{\mathcal O_Y}
(\mathcal F,\mathcal L^\bullet)))
\longrightarrow
\mathcal H^p(\mathcal{H}om_{\mathcal O_X}
(f^*(\mathcal F),f^*(\mathcal L^\bullet)))
\]
을 얻는다. $f^*$가 완전하므로 정의에 따라
\[
\mathcal H^p(f^*(\mathcal{H}om_{\mathcal O_Y}
(\mathcal F,\mathcal L^\bullet)))
=f^*(\mathcal H^p(\mathcal{H}om_{\mathcal O_Y}
(\mathcal F,\mathcal L^\bullet)))
=f^*(\mathcal{E}xt_{\mathcal O_Y}^p(\mathcal F,\mathcal G)).
\]
한편 $f^*$의 완전성에 따라 $f^*(\mathcal L^\bullet)$은
$f^*(\mathcal G)$의 분해이다. $\mathcal L'^\bullet=(\mathcal L'^i)$을
$\mathcal O_X$-가군 범주에서 $f^*(\mathcal G)$의 단사 분해라 하면,
호모토피를 제외하고 정해지는 복합체의 준동형
$f^*(\mathcal L^\bullet)\to\mathcal L'^\bullet$이 존재하며, 따라서
코호몰로지 위에 잘 정의된 준동형을 준다. 이 준동형과 위에서 정의한
준동형을 합성하여 (\hyperref[0.12.3.4.1-fr]{12.3.4.1})을 얻는다.

\item[{\rm(ii)}]
앞의 표기 아래 $\mathcal O_X$-가군층 복합체의 준동형
\[
f^*(\mathcal{H}om_{\mathcal O_Y}(\mathcal F,\mathcal L^\bullet))
\longrightarrow
\mathcal{H}om_{\mathcal O_X}(f^*(\mathcal F),\mathcal L'^\bullet)
\]
이 있다. $\mathcal M^{\bullet\bullet}$을
$\mathcal O_Y$-가군 범주에서 복합체
$\mathcal{H}om_{\mathcal O_Y}(\mathcal F,\mathcal L^\bullet)$의 단사
카르탕--아일렌베르크 분해라 하자. $f^*$의 완전성에 따라
$f^*(\mathcal M^{\bullet\bullet})$은 복합체
$f^*(\mathcal{H}om_{\mathcal O_Y}(\mathcal F,\mathcal L^\bullet))$의
카르탕--아일렌베르크 분해이다. $\mathcal M'^{\bullet\bullet}$을
복합체
$\mathcal{H}om_{\mathcal O_X}(f^*(\mathcal F),\mathcal L'^\bullet)$의
단사 카르탕--아일렌베르크 분해라 하면,
(\hyperref[0.11.4.2-fr]{11.4.2})에 따라 위의 준동형과 양립하는,
호모토피를 제외하고 정해지는 준동형
\[
f^*(\mathcal M^{\bullet\bullet})\longrightarrow
\mathcal M'^{\bullet\bullet}
\]
이 있다. 다시 말해 이는 층의 이중복합체들의 $f$-사상
$\mathcal M^{\bullet\bullet}\to\mathcal M'^{\bullet\bullet}$이다.
여기서 호모토피를 제외하고 정해지는 가군 이중복합체의 디-준동형
$\Gamma(Y,\mathcal M^{\bullet\bullet})\to
\Gamma(X,\mathcal M'^{\bullet\bullet})$을 얻고, 이어서 잘 정의된
스펙트럼 열의 사상을 얻는다
(\hyperref[0.11.3.2-fr]{11.3.2}). 이것이 구하는 사상
(\hyperref[0.12.3.4.2-fr]{12.3.4.2})이며,
(\hyperref[0.12.3.4.3-fr]{12.3.4.3})의 특징은 정의에서 곧바로 따른다.
\end{enumerate}
\end{proof}

\begin{proposition}[12.3.5]
\label{0.12.3.5-fr}
\emph{(\hyperref[0.12.3.4-fr]{12.3.4})의 가정 아래 환층
$\mathcal O_Y$도 연접이라고 하자. 그러면 모든 연접
$\mathcal O_Y$-가군 $\mathcal F$에 대하여 표준 준동형
(\hyperref[0.12.3.4.1-fr]{12.3.4.1})은 전단사이다.}
\end{proposition}

\begin{proof}
\oldpage[0\textsubscript{III}]{62}
문제는 $Y$ 위에서 국소적이므로 완전열
\[
0\longrightarrow\mathcal R\longrightarrow\mathcal O_Y^n
\longrightarrow\mathcal F\longrightarrow0
\]
이 존재한다고 가정할 수 있다. 이때 $\mathcal R$도 연접
$\mathcal O_Y$-가군이다
(\hyperref[0.5.3.4-ko]{0\textsubscript{I}, 5.3.4}). 준동형
\[
f^*(\mathcal{E}xt_{\mathcal O_Y}^p(\mathcal F,\mathcal G))
\longrightarrow
\mathcal{E}xt_{\mathcal O_X}^p(f^*(\mathcal F),f^*(\mathcal G))
\]
이 전단사임을 보이기 위하여 $p$에 관한 귀납법을 쓴다. $p=0$일 때는
(\hyperref[0.6.7.6.1-ko]{0\textsubscript{I}, 6.7.6.1})에서 명제가
따른다. 이제 가환도식
\[
{\tiny
\xymatrix@C=2pt@R=24pt{
f^*(\mathcal{E}xt_{\mathcal O_Y}^{p-1}(\mathcal O_Y^n,\mathcal G))
\ar[r]\ar[d]&
f^*(\mathcal{E}xt_{\mathcal O_Y}^{p-1}(\mathcal R,\mathcal G))
\ar[r]^\partial\ar[d]&
f^*(\mathcal{E}xt_{\mathcal O_Y}^{p}(\mathcal F,\mathcal G))
\ar[r]\ar[d]&
f^*(\mathcal{E}xt_{\mathcal O_Y}^{p}(\mathcal O_Y^n,\mathcal G))
\ar[d]\\
\mathcal{E}xt_{\mathcal O_X}^{p-1}(\mathcal O_X^n,f^*(\mathcal G))
\ar[r]&
\mathcal{E}xt_{\mathcal O_X}^{p-1}(f^*(\mathcal R),f^*(\mathcal G))
\ar[r]^\partial&
\mathcal{E}xt_{\mathcal O_X}^{p}(f^*(\mathcal F),f^*(\mathcal G))
\ar[r]&
\mathcal{E}xt_{\mathcal O_X}^{p}(\mathcal O_X^n,f^*(\mathcal G)).
}}
\]
을 얻는다. 여기서 $f^*(\mathcal O_Y)=\mathcal O_X$이고, $f^*$가
완전하므로 두 행은 완전하다. 또한
\[
\mathcal{E}xt_{\mathcal O_Y}^p(\mathcal O_Y^n,\mathcal G)=0
\quad\hbox{모든 }p>0\hbox{에 대하여},
\quad\hbox{그리고 마찬가지로}\quad
\mathcal{E}xt_{\mathcal O_X}^p(\mathcal O_X^n,f^*(\mathcal G))=0
\quad\hbox{모든 }p>0\hbox{에 대하여}
\]
이다 (T, 4.2.3). 귀납가정에 따라 앞 도식의 처음 두 수직 화살표는
동형이고 오른쪽 항들은 0이므로,
\[
f^*(\mathcal{E}xt_{\mathcal O_Y}^p(\mathcal F,\mathcal G))
\longrightarrow
\mathcal{E}xt_{\mathcal O_X}^p(f^*(\mathcal F),f^*(\mathcal G))
\]
은 동형이다.
\end{proof}

\subsection{직상 함자의 초코호몰로지}
\label{subsection:0.12.4-fr}

\begin{env}[12.4.1]
\label{0.12.4.1-fr}
$(X,\mathcal O_X)$, $(Y,\mathcal O_Y)$를 두 환 달린 공간이라 하고,
$f:X\to Y$를 환 달린 공간의 사상이라 하자. 임의의
$\mathcal O_X$-가군 복합체
$\mathcal K^\bullet=(\mathcal K^i)_{i\in\mathbb Z}$에 대하여
$f_*$의 \emph{초코호몰로지}를 취할 수 있다
(\hyperref[0.11.4.4-fr]{11.4.4})~; 실제로 $\mathcal O_Y$-가군들의 아벨
범주에서는 여과 귀납극한이 존재하고 완전하다 (T, 3.1.1).
초코호몰로지 $\mathcal O_Y$-가군 $R^p f_*(\mathcal K^\bullet)$는
$\mathcal H^p(f,\mathcal K^\bullet)$ 또는
$\mathcal H_f^p(\mathcal K^\bullet)$로도 나타내기로 한다.
$\mathcal H^\bullet(f,\mathcal K^\bullet)$은
$\mathcal O_Y$-가군의 이중복합체 $f_*(\mathcal L^{\bullet\bullet})$의
코호몰로지임을 상기하자. 여기서 $\mathcal L^{\bullet\bullet}$은
$\mathcal O_X$-가군들의 범주에서 $\mathcal K^\bullet$의
카르탕--아일렌베르크 단사 분해이다~;
$\mathcal H^\bullet(f,\mathcal K^\bullet)$은 두 스펙트럼 열
${}'E(f,\mathcal K^\bullet)$과 ${}''E(f,\mathcal K^\bullet)$의
수렴대상이며, 그 $E_2$항들은 다음과 같다.
\begin{equation}
\label{0.12.4.1.1-fr}
\tag{12.4.1.1}
{}'E_2^{pq}=\mathcal H^p(\mathcal H^q(f,\mathcal K^\bullet))
\end{equation}
\begin{equation}
\label{0.12.4.1.2-fr}
\tag{12.4.1.2}
{}''E_2^{pq}=\mathcal H^p(f,\mathcal H^q(\mathcal K^\bullet))
\quad(=R^p f_*(\mathcal H^q(\mathcal K^\bullet))).
\end{equation}
이 공식들에서는 함자를 복합체에 적용한 것에 대한 일반 표기
$T(A^\bullet)$를 채택하였고
(\hyperref[0.11.2.1-fr]{11.2.1}), $\mathcal O_X$-가군
$\mathcal F$에 대하여 $R^p f_*(\mathcal F)$ 대신
$\mathcal H^p(f,\mathcal F)$를 쓴다. 또한 스펙트럼 열
${}'E(f,\mathcal K^\bullet)$은 언제나 \emph{정칙}임을 상기하자~;
$\mathcal K^\bullet$이 아래로 유계이면 두 스펙트럼 열
${}'E(f,\mathcal K^\bullet)$과 ${}''E(f,\mathcal K^\bullet)$은
\emph{쌍정칙}이고,
\oldpage[0\textsubscript{III}]{63}
또는 모든 $\mathcal O_X$-가군이 길이가 $\leq m$인 플라스크 분해를
갖도록 하는 정수 $m$이 존재할 때에도 그러하다
(\hyperref[0.11.4.4-fr]{11.4.4}).
\end{env}

\begin{env}[12.4.2]
\label{0.12.4.2-fr}
마찬가지로 $\mathbf H^\bullet(X,\mathcal K^\bullet)$으로
$\mathcal O_X$-가군 복합체 $\mathcal K^\bullet$에 대한 함자
$\Gamma$의 초코호몰로지를 나타내기로 한다~; 따라서
$\mathbf H^p(X,\mathcal K^\bullet)$은
$\Gamma(X,\mathcal O_X)$-가군이다. 또한
$\mathbf H^\bullet(X,\mathcal K^\bullet)$을
$\mathcal H^\bullet(f,\mathcal K^\bullet)$의 특수한 경우로 볼 수 있다.
여기서 $f$는 $(X,\mathcal O_X)$에서 환
$\Gamma(X,\mathcal O_X)$을 갖춘 한 점짜리 환 달린 공간으로 가는
사상이다.

$X$의 임의의 열린집합 $V$에 대하여
$\mathbf H^\bullet(V,\mathcal K^\bullet|V)$ 대신
$\mathbf H^\bullet(V,\mathcal K^\bullet)$을 쓰기로 한다.
\end{env}

\begin{proposition}[12.4.3]
\label{0.12.4.3-fr}
\emph{임의의 정수 $p\in\mathbb Z$에 대하여 $\mathcal O_Y$-가군
$\mathcal H^p(f,\mathcal K^\bullet)$은 $Y$ 위의 준층
$U\mapsto\mathbf H^p(f^{-1}(U),\mathcal K^\bullet)$에 대응하는 층과
표준적으로 동형이다.}
\end{proposition}

\begin{proof}
실제로 (\hyperref[0.12.4.1-fr]{12.4.1})의 표기 아래에서 코호몰로지층
$\mathcal H^p(f_*(\mathcal L^{\bullet\bullet}))$은 다음 준층에
대응하는 층이다.
\[
U\longmapsto H^p(\Gamma(U,f_*(\mathcal L^{\bullet\bullet})))
=H^p(\Gamma(f^{-1}(U),\mathcal L^{\bullet\bullet})).
\]
그런데 $\mathcal L^{\bullet\bullet}|f^{-1}(U)$가
$\mathcal K^\bullet|f^{-1}(U)$의 카르탕--아일렌베르크 단사 분해임은
분명하므로 (T, 3.1.3), 정의에 따라
\[
H^p(\Gamma(f^{-1}(U),\mathcal L^{\bullet\bullet}))
=\mathbf H^p(f^{-1}(U),\mathcal K^\bullet)
\]
이다.
\end{proof}

\begin{proposition}[12.4.4]
\label{0.12.4.4-fr}
\emph{초코호몰로지 $\mathcal H^\bullet(f,\mathcal K^\bullet)$은 다음
각 경우에 $\mathcal K^\bullet$에 관한 코호몰로지 함자이다~:}
\begin{enumerate}
\item[\emph{a)}]
\emph{$\mathcal K^\bullet$이 아래로 유계인 복합체들의 범주를 움직인다.}
\item[\emph{b)}]
\emph{모든 $\mathcal O_X$-가군이 길이가 $\leq m$인 플라스크 분해를
갖도록 하는 정수 $m$이 존재한다.}
\item[\emph{c)}]
\emph{$X$는 뇌터 공간이다.}
\end{enumerate}
\end{proposition}

\begin{proof}
경우 a)와 b)는 (\hyperref[0.11.5.4-ko]{11.5.4})의 특수한 경우이다.
한편 이 경우 함자 $f_*$가 귀납극한과 가환함이 알려져 있으므로
(G, II, 3.10.1), 경우 c)는
(\hyperref[0.11.5.2-ko]{11.5.2})에서 따른다.
\end{proof}

\begin{env}[12.4.5]
\label{0.12.4.5-fr}
이제 $X$의 열린 덮개 $\mathfrak U=(U_\alpha)$를 생각하자. $X$ 위의
임의의 \emph{준층} 복합체
$\mathcal K^\bullet=(\mathcal K^j)$에 대하여 이중복합체
$\check C^\bullet(\mathfrak U,\mathcal K^\bullet)$를 생각하자. 이
이중복합체의 $(i,j)$ 성분은
$\check C^i(\mathfrak U,\mathcal K^j)$이며, 이는
$\mathcal K^j$에 값을 갖는 $\mathfrak U$의 신경의 교대
$i$-공사슬군이다 (G, II, 5.1). 이 이중복합체의 코호몰로지를
\emph{$\mathcal K^\bullet$을 계수로 하는 덮개 $\mathfrak U$의
초코호몰로지}라 하고, 다음과 같이 나타내기로 한다.
\[
\mathbf H^\bullet(\mathfrak U,\mathcal K^\bullet)
=H^\bullet(\check C^\bullet(\mathfrak U,\mathcal K^\bullet)).
\]
덮개의 르레 스펙트럼 열 (T, 3.8.1 및 G, II, 5.9.1)은 다음과 같이
초코호몰로지로 일반화된다~:
\end{env}

\begin{proposition}[12.4.6]
\label{0.12.4.6-fr}
\emph{$\mathcal K^\bullet=(\mathcal K^i)$를 $\mathcal O_X$-가군의
복합체라 하자. $\mathcal K^\bullet$에 관하여 정칙인 스펙트럼 함자로서
초코호몰로지 $\mathbf H^\bullet(X,\mathcal K^\bullet)$로 수렴하고,
그 $E_2$항이 다음과 같은 것이 존재한다.}
\begin{equation}
\label{0.12.4.6.1-fr}
\tag{12.4.6.1}
E_2^{pq}=\check H^p(\mathfrak U,h^q(\mathcal K^\bullet)),
\end{equation}
\emph{여기서 $h^q(\mathcal K^\bullet)$은 $X$ 위의 준층 복합체
$V\mapsto H^q(V,\mathcal K^\bullet)$를 나타낸다. 앞의 스펙트럼 열은
$\mathcal K^\bullet$이 아래로 유계이면 쌍정칙이다.}
\end{proposition}

\begin{proof}
$\mathcal K^\bullet$의 카르탕--아일렌베르크 단사 분해
$\mathcal L^{\bullet\bullet}$와 삼중복합체
$\check C^\bullet(\mathfrak U,\mathcal L^{\bullet\bullet})
=(\check C^i(\mathfrak U,\mathcal L^{jk}))$를 생각하자~; 먼저 이
삼중복합체를 차수 $i$와 $j+k$에 관한 이중복합체로 보자. $i$는
$\geq0$인 값만 취하므로 이 이중복합체의 두 번째 스펙트럼 열은
정칙이다 (\hyperref[0.11.3.3-fr]{11.3.3}). 또한 모든 $q>0$에 대하여
$\check H^q(\mathfrak U,\mathcal L^{jk})=0$이므로 퇴화한다. 실제로
$\mathcal O_X$-가군 $\mathcal L^{jk}$들은 플라스크층이다
\oldpage[0\textsubscript{III}]{64}
(G, II, 5.2.3). 따라서 (\hyperref[0.11.1.6-fr]{11.1.6})에 의해 표준
동형
\[
H^n(\check C^\bullet(\mathfrak U,\mathcal L^{\bullet\bullet}))
\xrightarrow{\sim}
H^n(\Gamma(X,\mathcal L^{\bullet\bullet}))
\]
을 얻는다 ((G, II, 5.2.2)에 의함). 그러므로 정의
(\hyperref[0.12.4.2-fr]{12.4.2})에 따라 동형
\[
H^n(\check C^\bullet(\mathfrak U,\mathcal L^{\bullet\bullet}))
\xrightarrow{\sim}\mathbf H^n(X,\mathcal K^\bullet)
\]
을 얻는다. 한편 삼중복합체
$\check C^\bullet(\mathfrak U,\mathcal L^{\bullet\bullet})$를 차수
$i+j$와 $k$에 관한 이중복합체로 보자. $k$는 $\geq0$인 값만 취하므로
이 이중복합체의 첫 번째 스펙트럼 열은 언제나 정칙이다~;
$\mathcal L^{jk}=0$ $(j<j_0)$이면, 즉 $\mathcal K^\bullet$이 아래로
유계이면 쌍정칙이다 (\hyperref[0.11.3.3-fr]{11.3.3}). 이 스펙트럼 열이
구하는 것이다~; 실제로 모든 $j$에 대하여
$\mathcal L^{j,\bullet}$은 $\mathcal K^j$의 단사 분해이다~; 따라서
$H^q(\check C^i(\mathfrak U,\mathcal L^{j,\bullet}))$은 다름 아닌
공사슬 복합체 $\check C^i(\mathfrak U,h^q(\mathcal K^j))$이다. 이로써
증명이 끝난다.
\end{proof}

\begin{corollary}[12.4.7]
\label{0.12.4.7-fr}
\emph{$\mathfrak U$의 신경의 모든 단체 $\sigma$와 모든 정수 $i$에
대하여 $H^q(U_\sigma,\mathcal K^i)=0$ $(q>0)$이면 표준 동형}
\begin{equation}
\label{0.12.4.7.1-fr}
\tag{12.4.7.1}
\mathbf H^\bullet(\mathfrak U,\mathcal K^\bullet)
\xrightarrow{\sim}\mathbf H^\bullet(X,\mathcal K^\bullet).
\end{equation}
\emph{이 존재한다.}
\end{corollary}

\begin{proof}
실제로 가정에 의해 $q>0$이면
$\check C^\bullet(\mathfrak U,h^q(\mathcal K^i))=0$이고, 따라서
$E_2^{pq}=0$이다~; 열 (\hyperref[0.12.4.6.1-fr]{12.4.6.1})이 퇴화하고
정칙이므로, 결론은 $\mathbf H^\bullet(\mathfrak U,\mathcal K^\bullet)$의
정의 (\hyperref[0.12.4.5-fr]{12.4.5})와
(\hyperref[0.11.1.6-fr]{11.1.6})에서 따른다.
\end{proof}

\begin{env}[12.4.8]
\label{0.12.4.8-fr}
$(X',\mathcal O_{X'})$을 두 번째 환 달린 공간이라 하고,
$f=(\psi,\theta)$를 $X'$에서 $X$로 가는 사상이라 하자.
(\hyperref[0.12.1.3-fr]{12.1.3})과
(\hyperref[0.12.1.4-fr]{12.1.4})에서와 같은 방법으로
$\mathcal O_X$-가군 복합체 $\mathcal K^\bullet$의 초코호몰로지에 대한
디-준동형
\begin{equation}
\label{0.12.4.8.1-fr}
\tag{12.4.8.1}
\mathbf H^p(X,\mathcal K^\bullet)
\longrightarrow\mathbf H^p(X',f^*(\mathcal K^\bullet)).
\end{equation}
을 정의한다. $\mathcal K^\bullet$의 카르탕--아일렌베르크 단사 분해
$\mathcal L^{\bullet\bullet}$에서 출발하자. $\psi^*$가 완전하므로
$\psi^*(\mathcal L^{\bullet\bullet})$은 $\psi^*(\mathcal O_X)$-가군들의
범주에서 $\psi^*(\mathcal K^\bullet)$의 카르탕--아일렌베르크 분해이다~;
이때 사상
$\psi^*(\mathcal L^{\bullet\bullet})\to
\mathcal L'^{\bullet\bullet}$이 존재한다. 여기서
$\mathcal L'^{\bullet\bullet}$은 $\psi^*(\mathcal K^\bullet)$의
카르탕--아일렌베르크 단사 분해이다. 이로부터 코호몰로지에 대한 사상
\[
\mathbf H^\bullet(X,\mathcal K^\bullet)
\longrightarrow\mathbf H^\bullet(X',\psi^*(\mathcal K^\bullet))~;
\]
을 얻는다. 이를 함자성으로부터 얻는 사상
$\psi^*(\mathcal K^\bullet)\to f^*(\mathcal K^\bullet)$에서 유도되는
사상과 합성하면, 구하는 사상
(\hyperref[0.12.4.8.1-fr]{12.4.8.1})을 얻는다.

(\hyperref[0.12.4.8.1-fr]{12.4.8.1})과
(\hyperref[0.12.4.3-fr]{12.4.3})에서 출발하여
(\hyperref[0.12.2.5-fr]{12.2.5})에서와 같이 논하면, 환 달린 공간의
두 사상 $f:X\to Y$, $g:Y\to Z$에 대하여 $\mathcal O_X$-가군 복합체
$\mathcal K^\bullet$의 초코호몰로지에 대한 준동형
\begin{equation}
\label{0.12.4.8.2-fr}
\tag{12.4.8.2}
\mathcal H^n(g,f_*(\mathcal K^\bullet))
\longrightarrow\mathcal H^n(h,\mathcal K^\bullet)
\end{equation}
\begin{equation}
\label{0.12.4.8.3-fr}
\tag{12.4.8.3}
\mathcal H^n(h,\mathcal K^\bullet)
\longrightarrow g_*(\mathcal H^n(f,\mathcal K^\bullet)).
\end{equation}
들을 정의할 수 있다. 정의의 세부사항은 독자에게 맡긴다.
\end{env}

\section{호몰로지 대수학에서의 사영극한}
\label{section:0.13-fr}

\subsection{미타그--레플러 조건}
\label{subsection:0.13.1-fr}

\begin{env}[13.1.1]
\label{0.13.1.1-fr}
$\mathcal C$를 무한곱이 존재하는 아벨 범주라 하자 (T, 1.5의 공리
AB~3*). 그러면 $\mathcal C$의 한 대상의 부분대상들로 이루어진 임의의
족의 하한이 존재하고, $\mathcal C$의 대상들의 모든 사영계는
사영극한을 갖는다. 이 사영극한은 고려하는 사영계의 왼쪽 완전함자이다
(T, 1.8). $(A_\alpha,f_{\alpha\beta})$를 지표집합 $I$가 오른쪽
여과인 $\mathcal C$의 대상들의 사영계라 하자.
\oldpage[0\textsubscript{III}]{65}
$A=\varprojlim A_\alpha$라 하고, 모든 $\alpha\in I$에 대하여
$f_\alpha:A\to A_\alpha$를 표준 사상이라 하자. 각 $\alpha\in I$에
대하여 $\alpha\leq\beta$일 때의 $f_{\alpha\beta}(A_\beta)$들은
$A_\alpha$의 부분대상들로 이루어진 감소 여과족을 이룬다. 부분대상
\[
A'_\alpha=\inf_{\beta\geq\alpha}f_{\alpha\beta}(A_\beta)
\]
을 $A_\alpha$ 안의 ``보편상'' 부분대상이라 한다. 명백히
$f_\alpha(A)\subset A'_\alpha$이고, $\alpha\leq\beta$일 때
$f_{\alpha\beta}(A'_\beta)\subset A'_\alpha$이다. 따라서
$(A'_\alpha,f_{\alpha\beta}|A'_\beta)$는 사영계이고
$A=\varprojlim A'_\alpha$이다.
\end{env}

\begin{env}[13.1.2]
\label{0.13.1.2-fr}
$\mathcal C$ 안의 사영계 $(A_\alpha,f_{\alpha\beta})$에 대하여 다음
조건을 \emph{미타그--레플러 조건}이라 한다.
\begin{itemize}
\item[\textbf{(ML)}]
\emph{모든 지표 $\alpha$에 대하여 $\beta\geq\alpha$가 존재하여,
모든 $\gamma\geq\beta$에 대해
$f_{\alpha\gamma}(A_\gamma)=f_{\alpha\beta}(A_\beta)$이다.}
\end{itemize}
$f_{\alpha\beta}$들이 \emph{에피사상}이면 조건 (ML)이 성립함은
명백하다. 역으로 (ML)이 성립하고, 모든 $\alpha\in I$에 대하여
$A'_\alpha$가 $A_\alpha$ 안의 ``보편상'' 부분대상이면,
$\alpha\leq\beta$일 때 $f_{\alpha\beta}$의 $A'_\beta$에 대한 제한은
\emph{에피사상} $A'_\beta\to A'_\alpha$이다. 실제로
$\delta\geq\gamma$일 때
$f_{\beta\delta}(A_\delta)=f_{\beta\gamma}(A_\gamma)$가 되도록
$\gamma\geq\beta$를 택하면
$A'_\beta=f_{\beta\gamma}(A_\gamma)$이다. 한편 이는
$\delta\geq\gamma$일 때
$f_{\alpha\delta}(A_\delta)=f_{\alpha\gamma}(A_\gamma)$를 함의하므로,
$A'_\alpha=f_{\alpha\gamma}(A_\gamma)=f_{\alpha\beta}(A'_\beta)$이다.

또한 각 대상 $A_\alpha$가 $\mathcal C$에서 \emph{아르틴}, 즉
$A_\alpha$의 부분대상들로 이루어진 모든 족이 극소 원소를 가지면
조건 (ML)이 성립한다. 실제로 $A_\alpha$의 부분대상들로 이루어진
감소 여과족 $(f_{\alpha\beta}(A_\beta))$의 극소 원소는 반드시 이
부분대상들 중 가장 작은 것이다.
\end{env}

\begin{remark}[13.1.3]
\label{0.13.1.3-fr}
예를 들어 $\mathcal C$가 집합의 범주일 때에도 조건 (ML)을 같은 방식으로
서술할 수 있다. 이 경우에도 $A_\alpha$의 ``보편상'' 부분집합을 정의할
수 있고, (\hyperref[0.13.1.1-fr]{13.1.1})과
(\hyperref[0.13.1.2-fr]{13.1.2})에서 한 주석들은 여전히 성립한다.
\end{remark}

\subsection{아벨군에 대한 미타그--레플러 조건}
\label{subsection:0.13.2-fr}

\begin{proposition}[13.2.1]
\label{0.13.2.1-fr}
\emph{
\[
0\longrightarrow A_\alpha\xrightarrow{u_\alpha}B_\alpha
\xrightarrow{v_\alpha}C_\alpha\longrightarrow0
\]
을 아벨군들의 사영계들의 완전열이라 하자. 세 사영계는 같은 여과
지표집합 $I$를 갖는다.}
\begin{enumerate}
\item[\emph{(i)}] \emph{$(B_\alpha)$가 (ML)을 만족하면
$(C_\alpha)$도 (ML)을 만족한다.}
\item[\emph{(ii)}] \emph{$(A_\alpha)$와 $(C_\alpha)$가 (ML)을
만족하면 $(B_\alpha)$도 (ML)을 만족한다.}
\end{enumerate}
\end{proposition}

\begin{proof}
$(f_{\alpha\beta})$, $(g_{\alpha\beta})$,
$(h_{\alpha\beta})$를 각각 사영계 $(A_\alpha)$, $(B_\alpha)$,
$(C_\alpha)$를 정의하는 준동형계라 하자.

(i) $\lambda\geq\beta$일 때
$g_{\alpha\beta}(B_\beta)=g_{\alpha\lambda}(B_\lambda)$라고 하자.
$v_\beta$와 $v_\lambda$가 전사이므로,
\[
h_{\alpha\beta}(C_\beta)
=v_\alpha(g_{\alpha\beta}(B_\beta))
=v_\alpha(g_{\alpha\lambda}(B_\lambda))
=h_{\alpha\lambda}(C_\lambda)
\]
가 모든 $\lambda\geq\beta$에 대하여 성립한다.

(ii) $\alpha\in I$라 하고, $\lambda\geq\beta$일 때
$f_{\alpha\beta}(A_\beta)=f_{\alpha\lambda}(A_\lambda)$가 되도록
$\beta\geq\alpha$를 택하자. 또한 $\lambda\geq\gamma$일 때
$h_{\beta\gamma}(C_\gamma)=h_{\beta\lambda}(C_\lambda)$가 되도록
$\gamma\geq\beta$를 택하자. 이제
$y_\alpha\in g_{\alpha\gamma}(B_\gamma)$라 하면,
$y_\gamma\in B_\gamma$가 존재하여
$y_\alpha=g_{\alpha\gamma}(y_\gamma)$이다.
$y_\beta=g_{\beta\gamma}(y_\gamma)$로 놓으면
$v_\beta(y_\beta)=h_{\beta\gamma}(v_\gamma(y_\gamma))$이다. 모든
$\lambda\geq\gamma$에 대하여 가정에 의해 $y_\lambda\in B_\lambda$가
존재하여
\[
h_{\beta\gamma}(v_\gamma(y_\gamma))
=h_{\beta\lambda}(v_\lambda(y_\lambda))
=v_\beta(g_{\beta\lambda}(y_\lambda)).
\]
따라서 $v_\beta(y_\beta-g_{\beta\lambda}(y_\lambda))=0$이고,
$x_\beta\in A_\beta$가 존재하여
\oldpage[0\textsubscript{III}]{66}
\[
y_\beta=g_{\beta\lambda}(y_\lambda)+u_\beta(x_\beta).
\]
그러므로
\[
y_\alpha=g_{\alpha\lambda}(y_\lambda)
+u_\alpha(f_{\alpha\beta}(x_\beta)).
\]
그런데 $\lambda\geq\beta$이므로 $x_\lambda\in A_\lambda$가 존재하여
$f_{\alpha\beta}(x_\beta)=f_{\alpha\lambda}(x_\lambda)$이고, 끝으로
\[
y_\alpha
=g_{\alpha\lambda}(y_\lambda)+u_\alpha(f_{\alpha\lambda}(x_\lambda))
=g_{\alpha\lambda}(y_\lambda+u_\lambda(x_\lambda))
\in g_{\alpha\lambda}(B_\lambda).
\]
이로써 증명이 끝난다.
\end{proof}

\begin{proposition}[13.2.2]
\label{0.13.2.2-fr}
\emph{$I$를 가산 공종 부분집합을 갖는 여과 순서집합이라 하자.
\[
0\longrightarrow A_\alpha\xrightarrow{u_\alpha}B_\alpha
\xrightarrow{v_\alpha}C_\alpha\longrightarrow0
\]
을 지표집합이 $I$인 아벨군들의 사영계들의 완전열이라 하자.
$(A_\alpha)$가 조건 (ML)을 만족하면 열
\[
0\longrightarrow\varprojlim A_\alpha
\longrightarrow\varprojlim B_\alpha
\longrightarrow\varprojlim C_\alpha\longrightarrow0
\]
은 완전하다.}
\end{proposition}

\begin{proof}
준동형
\[
v=\varprojlim v_\alpha:\varprojlim B_\alpha\longrightarrow
\varprojlim C_\alpha
\]
가 전사임을 보이면 충분하다. $z=(z_\alpha)$를
$\varprojlim C_\alpha$의 원소라 하고
$E_\alpha=v_\alpha^{-1}(z_\alpha)$로 놓자. 준동형
$g_{\alpha\beta}:B_\beta\to B_\alpha$의 제한에 대하여
$E_\alpha$들이 공집합이 아닌 집합들의 사영계를 이룸은 명백하다.
이 사영계가 조건 (ML)을 만족함을 보이자. $u_\alpha$를 통하여
$A_\alpha$를 $B_\alpha$의 부분집합으로 동일시한다. 모든
$\alpha\in I$에 대하여 $\beta\geq\alpha$가 존재하여
$\lambda\geq\beta$이면
$g_{\alpha\beta}(A_\beta)=g_{\alpha\lambda}(A_\lambda)$이다. 이때
$\lambda\geq\beta$에 대하여
$g_{\alpha\beta}(E_\beta)=g_{\alpha\lambda}(E_\lambda)$도 성립함을
보이자. 실제로 $y_\lambda\in E_\lambda$를 택하고
$y_\beta=g_{\beta\lambda}(y_\lambda)$,
$y_\alpha=g_{\alpha\lambda}(y_\lambda)$로 놓자. 또한
$y'_\alpha\in g_{\alpha\beta}(E_\beta)$라 하여, 어떤
$y'_\beta\in E_\beta$에 대하여
$y'_\alpha=g_{\alpha\beta}(y'_\beta)$라 하자. 그러면
$y'_\beta-y_\beta=x_\beta\in A_\beta$이고, 가정에 따라
$x_\lambda\in A_\lambda$가 존재하여
$g_{\alpha\beta}(x_\beta)=g_{\alpha\lambda}(x_\lambda)$이다. 따라서
\[
y'_\alpha=g_{\alpha\beta}(y_\beta)+g_{\alpha\beta}(x_\beta)
=g_{\alpha\lambda}(y_\lambda)+g_{\alpha\lambda}(x_\lambda)
=g_{\alpha\lambda}(y_\lambda+x_\lambda)
\in g_{\alpha\lambda}(E_\lambda).
\]
주장이 증명되었다. 이제 $I$에 관한 가정 아래 (ML)을 만족하는 공집합이
아닌 집합들의 사영계는 공집합이 아닌 사영극한을 가짐이 알려져 있다
(Bourbaki, \emph{Top. gén.}, chap.~II, 3\textsuperscript{e} éd., §~3,
th.~1). 따라서 $y=(y_\alpha)\in\varprojlim E_\alpha$가 존재하고,
정의에 의해 모든 $\alpha$에 대하여 $v_\alpha(y_\alpha)=z_\alpha$이므로
$z=v(y)$이다. 증명 끝.
\end{proof}

\begin{proposition}[13.2.3]
\label{0.13.2.3-fr}
\emph{$I$에 관한 가정은 (\hyperref[0.13.2.2-fr]{13.2.2})와 같다고
하자. $(K_\alpha^\bullet)_{\alpha\in I}$를 미분작용소의 차수가 $+1$인
아벨군 복합체
$K_\alpha^\bullet=(K_\alpha^n)_{n\in\mathbb Z}$들의 사영계라 하자.
각 $n$에 대하여 함자적인 표준 준동형}
\begin{equation}
\label{0.13.2.3.1-fr}
\tag{13.2.3.1}
h_n:H^n(\varprojlim K_\alpha^\bullet)
\longrightarrow\varprojlim H^n(K_\alpha^\bullet)
\end{equation}
\emph{이 존재한다. 모든 차수 $n$에 대하여 아벨군들의 사영계
$(K_\alpha^n)_{\alpha\in I}$가 (ML)을 만족하면 모든 준동형 $h_n$은
전사이다. 더 나아가 한 차수 $n$에 대하여 사영계
$(H^{n-1}(K_\alpha^\bullet))_{\alpha\in I}$가 (ML)을 만족하면
$h_n$은 전단사이다.}
\end{proposition}

\begin{proof}
모든 $n$에 대하여 $K^n=\varprojlim K_\alpha^n$로 놓자. 준동형
$h_n$의 정의는 도식
\[
\xymatrix{
\cdots\ar[r]&K^{n-1}\ar[r]\ar[d]&K^n\ar[r]\ar[d]
&K^{n+1}\ar[r]\ar[d]&\cdots\\
\cdots\ar[r]&K_\alpha^{n-1}\ar[r]&K_\alpha^n\ar[r]
&K_\alpha^{n+1}\ar[r]&\cdots
}
\]
의 가환성에서 온다.
\oldpage[0\textsubscript{III}]{67}
여기서 $K^\bullet$의 미분작용소들은 $K_\alpha^\bullet$의 대응하는
미분작용소들의 사영극한이다.

완전열들
\[
(*_n)\qquad
0\longrightarrow \mathrm B^n(K_\alpha^\bullet)
\longrightarrow \mathrm Z^n(K_\alpha^\bullet)
\longrightarrow H^n(K_\alpha^\bullet)\longrightarrow0
\]
\[
(**_n)\qquad
0\longrightarrow \mathrm Z^{n-1}(K_\alpha^\bullet)
\longrightarrow K_\alpha^{n-1}
\longrightarrow \mathrm B^n(K_\alpha^\bullet)\longrightarrow0
\]
을 생각하자. 가정과 명제
(\hyperref[0.13.2.1-fr]{13.2.1}, (i))에 따라 모든 $n$에 대하여
사영계 $(\mathrm B^n(K_\alpha^\bullet))_{\alpha\in I}$가 (ML)을
만족한다. 그러므로 (\hyperref[0.13.2.2-fr]{13.2.2})에 따라 열
\[
(***_n)\qquad
0\longrightarrow\varprojlim_\alpha\mathrm B^n(K_\alpha^\bullet)
\longrightarrow\varprojlim_\alpha\mathrm Z^n(K_\alpha^\bullet)
\longrightarrow\varprojlim_\alpha H^n(K_\alpha^\bullet)
\longrightarrow0
\]
은 완전하다. 그런데
$\varprojlim_\alpha\mathrm B^n(K_\alpha^\bullet)$은
$\mathrm B^n(K^\bullet)$을 포함하는 $K^n$의 부분군과 동일시되고,
$\varprojlim_\alpha\mathrm Z^n(K_\alpha^\bullet)$은
$\mathrm Z^n(K^\bullet)$의 부분군과 동일시된다. 따라서 $h_n$은
전사이다. 이제 사영계
$(H^{n-1}(K_\alpha^\bullet))_{\alpha\in I}$도 (ML)을 만족한다고
가정하면, 완전열 $(*_{n-1})$과 명제
(\hyperref[0.13.2.1-fr]{13.2.1}, (ii))에 따라 사영계
$(\mathrm Z^{n-1}(K_\alpha^\bullet))_{\alpha\in I}$가 (ML)을
만족한다. 그러면 (\hyperref[0.13.2.2-fr]{13.2.2})를 완전열
$(**_n)$에 적용하여 완전열
\[
0\longrightarrow
\varprojlim_\alpha\mathrm Z^{n-1}(K_\alpha^\bullet)
\longrightarrow K^{n-1}\xrightarrow{u}
\varprojlim_\alpha\mathrm B^n(K_\alpha^\bullet)
\longrightarrow0
\]
을 얻는다. 이제
$\varprojlim_\alpha\mathrm B^n(K_\alpha^\bullet)
\supset\mathrm B^n(K^\bullet)$이고, $u$와 포함
$\varprojlim_\alpha\mathrm B^n(K_\alpha^\bullet)\to K^n$의 합성이
미분작용소 $K^{n-1}\to K^n$이다. 따라서 $u$가 전사이면
$\varprojlim_\alpha\mathrm B^n(K_\alpha^\bullet)
=\mathrm B^n(K^\bullet)$이고, 이에 따라 $h_n$은 단사이다. 증명 끝.
\end{proof}

\begin{env}[13.2.4]
\label{0.13.2.4-fr}
\emph{주석} (13.2.4).
\begin{enumerate}
\item[\emph{(i)}]
(\hyperref[0.13.2.2-fr]{13.2.2})의 논증(Bourbaki,
\emph{같은 곳} 참조)에 따르면, 다음 조건들만을 가정해도 그 명제의
결론은 여전히 성립한다. 각 $A_\alpha$에 \emph{완비 거리화 가능}
공간의 구조를 줄 수 있고, 그 구조에서 평행이동은 위상동형이며,
사영계 $(A_\alpha)$를 정의하는 사상
$f_{\alpha\beta}:A_\beta\to A_\alpha$는 고려하는 거리들에 대하여
\emph{균등연속}이고, 끝으로 사영계 $(A_\alpha)$는 다음 조건을
만족한다고 하자.
\begin{itemize}
\item[\textbf{(ML')}]
\emph{모든 지표 $\alpha$에 대하여 $\beta\geq\alpha$가 존재하여,
모든 $\gamma\geq\beta$에 대해 $f_{\alpha\gamma}(A_\gamma)$가
$f_{\alpha\beta}(A_\beta)$에서 조밀하다.}
\end{itemize}

이로부터 (\hyperref[0.13.2.3-fr]{13.2.3})과 비슷한 보충을 얻는다.
모든 $\alpha$와 $n<0$에 대하여 $K_\alpha^n=0$이라고 하자. 또한
$n\geq0$일 때 $(K_\alpha^n)_{\alpha\in I}$가 (ML)을 만족하고,
$A_\alpha=H^0(K_\alpha^\bullet)$에 위의 성질들을 만족하는 거리공간의
구조를 줄 수 있다고 하자. 그러면 $n\geq2$에 대해서는
(\hyperref[0.13.2.3-fr]{13.2.3})의 결론들이 바뀌지 않고, 더 나아가
$h_1$은 \emph{전단사}이다. 실제로
(\hyperref[0.13.2.2-fr]{13.2.2})의 논증은 사영계
$(\mathrm B^1(K_\alpha^\bullet))_{\alpha\in I}$가 (ML)을 만족하고,
열 $(***_1)$이 완전하며, 위의 결과에 따라
$\varprojlim\mathrm B^1(K_\alpha^\bullet)=\mathrm B^1(K^\bullet)$임을
여전히 보여 준다. 이로써 특히 (T, 3.10.2)의 주장들을 얻었다.

\item[\emph{(ii)}]
함자 $\varprojlim$의 오른쪽 유도함자 $\varprojlim^{(n)}$을 도입하여
앞의 명제들보다 더 완전한 명제들을 얻을 수 있다 [28].
\end{enumerate}
\end{env}

\subsection{응용: 층들의 사영극한의 코호몰로지}
\label{subsection:0.13.3-fr}

\begin{proposition}[13.3.1]
\label{0.13.3.1-fr}
\oldpage[0\textsubscript{III}]{68}
\emph{$X$를 위상공간,
$(\mathcal F_k)_{k\in\mathbb N}$를 $X$ 위의 아벨군의 층들로 이루어진
사영계라 하고,
$\mathcal F=\varprojlim_k\mathcal F_k$라 하자. 다음 조건들이 성립한다고
가정한다.}
\begin{enumerate}
\item[\emph{(i)}]
\emph{$X$의 위상의 기저 $\mathfrak B$가 존재하여, 모든
$U\in\mathfrak B$와 모든 $i\geq0$에 대하여 사영계
$(H^i(U,\mathcal F_k))_{k\in\mathbb N}$는 (ML)을 만족한다.}
\item[\emph{(ii)}]
\emph{모든 $x\in X$와 모든 $i>0$에 대하여}
\[
\varinjlim_U\bigl(\varprojlim_k H^i(U,\mathcal F_k)\bigr)=0
\]
\emph{이다. 여기서 $U$는 $\mathfrak B$에 속하는 $x$의 근방들의 집합을
움직인다.}
\item[\emph{(iii)}]
\emph{사영계 $(\mathcal F_k)$를 정의하는 준동형
$u_{hk}:\mathcal F_k\to\mathcal F_h$ $(h\leq k)$는 전사이다.}
\end{enumerate}
\emph{그러면 모든 $i>0$에 대하여 표준 준동형}
\[
h_i:H^i(X,\mathcal F)\longrightarrow
\varprojlim_k H^i(X,\mathcal F_k)
\]
\emph{은 전사이다~; 더 나아가 주어진 $i$에 대하여 사영계
$(H^{i-1}(X,\mathcal F_k))_{k\in\mathbb N}$가 (ML)을 만족하면 $h_i$는
전단사이다.}
\end{proposition}

\begin{proof}
\emph{a)} 먼저 $\mathcal F_k$들과 $u_{hk}$들의 핵 $\mathcal N_{hk}$들이
플라스크라고 가정하자~; 이때 명제의 조건 (iii)로부터
$i>0$이면 $H^i(X,\mathcal F)=0$임을 보이겠다. $X$의 모든 열린집합
$U$와 $U$의 열린집합들로 이루어진 $U$의 모든 덮개 $\mathfrak U$에
대하여 $i>0$이면 $\check H^i(\mathfrak U,\mathcal F)=0$임을 보이면 충분하다.
실제로 이로부터 먼저 체흐 코호몰로지에 대하여 모든 $i>0$에 대해
$\check H^i(U,\mathcal F)=0$을 얻고, 이어서 ($X$의 모든 열린집합의 집합에
(G, II, 5.9.2)를 적용하여) 모든 $i>0$에 대해
$H^i(X,\mathcal F)=0$을 얻는다. $\mathcal F_k$들은 플라스크이므로
$i>0$이면 $\check H^i(\mathfrak U,\mathcal F_k)=0$이다
(G, II, 5.2.3)~; 각 $k$에 대하여 덮개 $\mathfrak U$의 신경의 교대
공사슬 복합체 $\check C^\bullet(\mathfrak U,\mathcal F_k)$를 생각하자
(G, II, 5.1). 이들은 자명하게 아벨군의 복합체들로 이루어진 사영계를
이룬다. 모든 사상
$\check C^\bullet(\mathfrak U,\mathcal F_k)\to
\check C^\bullet(\mathfrak U,\mathcal F_h)$ $(h\leq k)$가 전사임을 보이자.
정의에 따라, $X$의 모든 열린집합 $V$에 대하여 사상
$\Gamma(V,\mathcal F_k)\to\Gamma(V,\mathcal F_h)$가 전사임을 보이면
충분하다~; 그런데 가정에 의해 열
\[
0\longrightarrow\mathcal N_{hk}\longrightarrow\mathcal F_k
\longrightarrow\mathcal F_h\longrightarrow0
\]
은 완전하므로 코호몰로지 완전열
\[
\Gamma(V,\mathcal F_k)\longrightarrow\Gamma(V,\mathcal F_h)
\longrightarrow H^1(V,\mathcal N_{hk})=0
\]
을 준다. 여기서 $\mathcal N_{hk}$가 플라스크임을 썼다. 따라서 사영계
$(\check C^\bullet(\mathfrak U,\mathcal F_k))_{k\in\mathbb N}$는
(ML)을 만족한다~; 모든 $i\geq0$에 대하여
$(\check H^i(\mathfrak U,\mathcal F_k))_{k\in\mathbb N}$도 그러하다.
실제로 $i>0$일 때에는 자명하고,
$\check H^0(\mathfrak U,\mathcal F_k)=\Gamma(U,\mathcal F_k)$이므로
(G, II, 5.2.2), 앞에서 본 바에 의해 $i=0$일 때에도 (ML)이 성립한다.
따라서 (\hyperref[0.13.2.3-fr]{13.2.3})을 적용하면
\[
\check H^i(\mathfrak U,\mathcal F)
=\varprojlim_k\check H^i(\mathfrak U,\mathcal F_k)=0
\quad\hbox{모든 }i>0\hbox{에 대하여}.
\]

\emph{b)} 이제 일반적인 경우로 돌아가자. 각 $k\in\mathbb N$에 대하여
$\mathcal F_k$의 플라스크층들에 의한 표준 분해
$\mathcal C^\bullet(X,\mathcal F_k)
=(\mathcal C^i(X,\mathcal F_k))_{i\geq0}$를 생각하자
(G, II, 4.3). 각 $i\geq0$에 대하여
$(\mathcal C^i(X,\mathcal F_k))_{k\in\mathbb N}$가 플라스크층들로
이루어진 사영계임은 분명하다~; 이것이 \emph{a)}의 조건들을
만족함을 보이자.
\oldpage[0\textsubscript{III}]{69}
실제로 $h\leq k$일 때 $\mathcal N_{hk}$를 $u_{hk}$의 핵이라 하면,
열
\[
0\longrightarrow\mathcal N_{hk}\longrightarrow\mathcal F_k
\longrightarrow\mathcal F_h\longrightarrow0
\]
은 (iii)에 의해 완전하고, 우리의 주장은 함자
$\mathcal A\mapsto\mathcal C^i(X,\mathcal A)$가 완전하다는 사실에서
따른다 (G, II, 4.3). 다음과 같이 놓자.
$\mathcal G^i=\varprojlim_k\mathcal C^i(X,\mathcal F_k)$.
그러면 \emph{a)}에 의해 $j>0$, $i\geq0$이면
$H^j(X,\mathcal G^i)=0$이다. 이제
$\mathcal G^\bullet=(\mathcal G^i)_{i\geq0}$가 층 $\mathcal F$의
분해임을 보이자~; $j>0$이면 $H^j(X,\mathcal G^i)=0$이므로 코호몰로지
$H^\bullet(X,\mathcal F)$는 $H^\bullet(\Gamma(X,\mathcal G^\bullet))$와
같다 (G, II, 4.7.1).

사영극한을 취하면 완전열들
\[
0\longrightarrow\mathcal F_k
\longrightarrow\mathcal C^0(X,\mathcal F_k)
\longrightarrow\mathcal C^1(X,\mathcal F_k)\longrightarrow\cdots
\]
로부터 아벨군의 층들의 복합체
\[
0\longrightarrow\mathcal F\longrightarrow\mathcal G^0
\longrightarrow\mathcal G^1\longrightarrow\cdots
\]
를 얻음은 분명하다. 우리의 주장을 증명하려면 $i>0$일 때
$\mathcal H^i(\mathcal G^\bullet)=0$임을 보여야 한다. 이 층은 준층
$U\mapsto H^i(\Gamma(U,\mathcal G^\bullet))$에 의해 생성된다
(G, II, 4.1)~; 한편 복합체 $\Gamma(U,\mathcal G^\bullet)$는 아벨군의
복합체들의 사영계
\[
(\Gamma(U,\mathcal C^\bullet(X,\mathcal F_k)))_{k\in\mathbb N}
\]
의 사영극한이다 (0\textsubscript{I}, 3.2.6). \emph{a)}에서 각
$i\geq0$에 대하여 사상
$\Gamma(U,\mathcal C^i(X,\mathcal F_k))\to
\Gamma(U,\mathcal C^i(X,\mathcal F_h))$ $(h\leq k)$가 전사임을 보았다~;
다른 한편 표준 분해 $\mathcal C^\bullet(U,\mathcal F_k|U)$는
$\mathcal C^\bullet(X,\mathcal F_k)$가 $U$ 위에 유도하는 것이므로
$H^i(U,\mathcal F_k)=
H^i(\Gamma(U,\mathcal C^\bullet(X,\mathcal F_k)))$이다~; 따라서 가정
(i)에 의해 모든 $U\in\mathfrak B$에 대하여 복합체의 사영계
$(\Gamma(U,\mathcal C^\bullet(X,\mathcal F_k)))_{k\in\mathbb N}$에
(\hyperref[0.13.2.3-fr]{13.2.3})을 적용할 수 있고, 이에 따라
\[
H^i(\Gamma(U,\mathcal G^\bullet))
=\varprojlim_k H^i(U,\mathcal F_k)
\quad\hbox{모든 }i\geq0\hbox{에 대하여}.
\]
이제 가정 (ii)로부터 정의에 의해 $i>0$이면 층
$\mathcal H^i(\mathcal G^\bullet)$가 0임이 바로 따른다.

따라서 모든 $i\geq0$에 대하여
\[
H^i(X,\mathcal F)=H^i(\Gamma(X,\mathcal G^\bullet))
\quad\hbox{이고}\quad
\Gamma(X,\mathcal G^\bullet)
=\varprojlim_k\Gamma(X,\mathcal C^\bullet(X,\mathcal F_k)).
\]
방금 사상들
$\Gamma(X,\mathcal C^i(X,\mathcal F_k))\to
\Gamma(X,\mathcal C^i(X,\mathcal F_h))$ $(h\leq k)$가 모두
전사임을 보았다~; 따라서 결론은 다시
(\hyperref[0.13.2.3-fr]{13.2.3})에서 따른다.
\end{proof}

\begin{env}[13.3.2]
\label{0.13.3.2-fr}
\emph{주의} (13.3.2).
\begin{enumerate}
\item[\emph{(i)}]
(\hyperref[0.13.3.1-fr]{13.3.1})의 명제는 $i>0$일 때에만 내용이 있다.
$i=0$이면 아무 가정 없이도 $h_i$가 언제나 동형이기 때문이다
(0\textsubscript{I}, 3.2.6).
\item[\emph{(ii)}]
모든 $k$, 모든 $i>0$, 모든 $U\in\mathfrak B$에 대하여
$H^i(U,\mathcal F_k)=0$이고, $U\in\mathfrak B$일 때 사상
$\Gamma(U,\mathcal F_k)\to\Gamma(U,\mathcal F_h)$가 전사이면,
(\hyperref[0.13.3.1-fr]{13.3.1})의 조건 (i)와 (ii)는 특히 성립한다.
이 경우가 (\hyperref[0.13.3.1-fr]{13.3.1})의 가장 흔한 적용 사례이다.
\end{enumerate}
\end{env}

\subsection{미타그--레플러 조건과 사영계의 연관 등급 대상}
\label{subsection:0.13.4-fr}
\label{0.13.4-fr}

\begin{env}[13.4.1]
\label{0.13.4.1-fr}
$\mathbf A=(A_k,u_{kh})_{k\in\mathbb Z}$를 아벨 범주 $C$의 사영계라
하자~; 어떤 $k_0$가 존재하여 $k<k_0$이면 $A_k=0$일 때 이 사영계가
\emph{아래로 유계}라고 하겠다. 각 $A_k$ 위에 다음 공식으로
\emph{여과} $(\mathrm F^p(A_k))_{p\in\mathbb Z}$를 정의한다.
\begin{equation}
\label{0.13.4.1.1-fr}
\tag{13.4.1.1}
\mathrm F^p(A_k)=
\begin{cases}
\operatorname{Ker}(A_k\longrightarrow A_{p-1})&\text{$p\leq k+1$일 때},\\
0&\text{$p\geq k+1$일 때}.
\end{cases}
\end{equation}
\oldpage[0\textsubscript{III}]{70}
따라서 가정에 의해 $\mathrm F^{k_0}(A_k)=A_k$이고
$\mathrm F^{k+1}(A_k)=0$이다. 다시 말해 이 여과는 \emph{유한}이다
(\hyperref[0.11.1.3-fr]{11.1.3}). 따라서 이 여과의 연관 등급 대상들은
\[
\operatorname{gr}^p(A_k)=
\operatorname{Ker}(A_k\longrightarrow A_{p-1})
/\operatorname{Ker}(A_k\longrightarrow A_p)
\]
이고, 이에 따라 $\operatorname{gr}^p(A_k)$는
$\operatorname{Ker}(A_k\to A_{p-1})$의 $A_k\to A_p$에 의한 상과
동형이다~; 사영계를 정의하는 사상들의 추이성에 의해
\begin{equation}
\label{0.13.4.1.2-fr}
\tag{13.4.1.2}
\operatorname{gr}^p(A_k)=
\operatorname{Ker}(A_p\longrightarrow A_{p-1})
\cap\operatorname{Im}(A_k\longrightarrow A_p)
\end{equation}
이다. 그런데 (\hyperref[0.13.4.1.1-fr]{13.4.1.1})에 의해
$\operatorname{Ker}(A_p\to A_{p-1})=\operatorname{gr}^p(A_p)$이므로,
또한
\begin{equation}
\label{0.13.4.1.3-fr}
\tag{13.4.1.3}
\operatorname{gr}^p(A_k)=
\operatorname{gr}^p(A_p)\cap\operatorname{Im}(A_k\longrightarrow A_p).
\end{equation}
앞의 정의들로부터 $k\leq h$일 때
\[
u_{kh}(\mathrm F^p(A_h))\subset\mathrm F^p(A_k)
\]
임도 알 수 있다. 따라서 모든 $p\in\mathbb Z$에 대하여
$\operatorname{gr}^p(u_{kh})$들은 \emph{사영계}
$(\operatorname{gr}^p(A_k))_{k\in\mathbb Z}$를 정의한다.
\end{env}

\begin{env}[13.4.2]
\label{0.13.4.2-fr}
충분히 큰 $k$에 대하여 사상 $A_{k+1}\to A_k$가 \emph{동형}이면
사영계 $\mathbf A$가 \emph{본질적으로 상수}라고 하겠다. 사상
$A_i\to A_j$ $(j\leq i)$가 \emph{에피사상}이면 사영계 $\mathbf A$가
\emph{엄밀}하다고 하겠다. $\mathbf A$가 엄밀하면
(\hyperref[0.13.4.1.3-fr]{13.4.1.3})으로부터 $p\leq k\leq h$일 때
표준 사상 $\operatorname{gr}^p(A_h)\to\operatorname{gr}^p(A_k)$가
\emph{동형}임이 따른다. 다시 말해 사영계
$(\operatorname{gr}^p(A_k))_{k\in\mathbb Z}$는 본질적으로 상수이다.
대상들 $\operatorname{gr}^p(A_p)$의 열(각 $p$에 대하여
$\varprojlim_k\operatorname{gr}^p(A_k)$와 동일시한다)을 이때
$\operatorname{gr}^\bullet(\mathbf A)$로 쓰고, 이를 엄밀 사영계
$\mathbf A=(A_k)$의 \emph{연관 등급 대상}이라 한다.

이제 (아래로 유계인) 사영계 $\mathbf A$가 (ML)을 만족한다고
가정하자. 그러면 ``보편상'' 대상들의 사영계
$\mathbf A'=(A'_k)$가 \emph{엄밀}함을 알고 있으며
(\hyperref[0.13.1.2-fr]{13.1.2}), 이 사영계는 또한 아래로 유계이다~;
$\mathbf A'$의 연관 등급 대상 $\operatorname{gr}^\bullet(\mathbf A')$를
이때에도 $\mathbf A$에 \emph{연관된} 등급 대상이라 하고
$\operatorname{gr}^\bullet(\mathbf A)$로 쓴다.
\end{env}

\begin{proposition}[13.4.3]
\label{0.13.4.3-fr}
\emph{$\mathbf A=(A_k)_{k\in\mathbb Z}$를 아벨 범주 $C$의 아래로
유계인 사영계라 하자. 다음 두 조건은 동치이다.}
\begin{enumerate}
\item[\emph{a)}] \emph{$\mathbf A$는 조건 (ML)을 만족한다.}
\item[\emph{b)}] \emph{모든 $p\in\mathbb Z$에 대하여 사영계
$(\operatorname{gr}^p(A_k))_{k\in\mathbb Z}$는 본질적으로 상수이다.}
\end{enumerate}
\emph{더 나아가 이 조건들이 성립하면 모든 $p\in\mathbb Z$에 대하여
표준 동형}
\begin{equation}
\label{0.13.4.3.1-fr}
\tag{13.4.3.1}
\operatorname{gr}^p(\mathbf A)\xrightarrow{\sim}
\varprojlim_k\operatorname{gr}^p(A_k)
\end{equation}
\emph{이 존재한다.}
\end{proposition}

\begin{proof}
(\hyperref[0.13.4.1.2-fr]{13.4.1.2})으로부터 \emph{a)}가 \emph{b)}를
함의함이 곧바로 따른다~; 같은 공식을 사영계 $\mathbf A'$에 적용하면
((\hyperref[0.13.4.2-fr]{13.4.2})의 표기) 정의에 의해 동형
(\hyperref[0.13.4.3.1-fr]{13.4.3.1})을 얻는다. $k\leq h$일 때
$A_{kh}=\operatorname{Im}(A_h\to A_k)$로 놓자~; $k\leq h\leq j$이면
$A_{kj}\subset A_{kh}\subset A_k$이다. $A_{kh}$에
$(\mathrm F^p(A_k))$가 유도하는 여과를 주자~; $\mathbf A$를 정의하는
사상들의 추이성에 의해, 이 여과가 $(\mathrm F^p(A_h))$의 몫 여과이기도
함을 곧바로 확인할 수 있다~; 따라서
\begin{equation}
\label{0.13.4.3.2-fr}
\tag{13.4.3.2}
\operatorname{gr}^p(A_{kh})=
\operatorname{Im}\bigl(\operatorname{gr}^p(A_h)\longrightarrow
\operatorname{gr}^p(A_k)\bigr).
\end{equation}
\oldpage[0\textsubscript{III}]{71}
이제 \emph{b)}가 성립한다고 가정하자~; 모든 $p\in\mathbb Z$와 모든
$k\geq p$에 대하여, (\hyperref[0.13.4.3.2-fr]{13.4.3.2})의 우변이
$h\geq\mathrm L(p,k)$일 때 일정해지는 정수 $\mathrm L(p,k)$가
존재한다~; $p<k_0$이면 $\operatorname{gr}^p(A_k)=0$이므로, 주어진
$k$에 대하여 $p$가 $k$ 이하인 정수들의 집합을 움직일 때 0이 아닌
$\mathrm L(p,k)$는 유한 개뿐이다. 이 정수들의 최댓값을
$\mathrm L(k)=m$이라 하자~; 모든 $h\geq m$에 대하여
$A_{kh}\subset A_{km}$이고, $m$의 정의에 의해 표준 단사
$A_{kh}\to A_{km}$는 \emph{동형}
$\operatorname{gr}^\bullet(A_{kh})\xrightarrow{\sim}
\operatorname{gr}^\bullet(A_{km})$를 정의한다~; 여과들이 \emph{유한}이므로
앞의 단사 자체가 전단사임을 얻는다 (Bourbaki, \emph{Alg. comm.},
chap.~III, §~2, n\textsuperscript{o}~8, th.~1). 이는 $\mathbf A$가
(ML)을 만족함을 증명한다.
\end{proof}

\begin{env}[13.4.4]
\label{0.13.4.4-fr}
$C$에서 사영극한 $A=\varprojlim_k A_k$가 존재한다고 가정하자.
(\hyperref[0.13.4.1-fr]{13.4.1})의 정의에서 $A_k$를 $A$로 바꿀 수
있으며, 이처럼 $A$ 위에 정의한 여과는 여전히
\begin{equation}
\label{0.13.4.4.1-fr}
\tag{13.4.4.1}
\operatorname{gr}^p(A)=\operatorname{gr}^p(A_p)
\cap\operatorname{Im}(A\longrightarrow A_p)
\end{equation}
을 만족한다.
\end{env}

\begin{corollary}[13.4.5]
\label{0.13.4.5-fr}
\emph{$C$가 아벨군의 범주라고 가정하자. 사영계 $\mathbf A$가 (ML)을
만족하고 $A=\varprojlim_k A_k$이면, 모든 $p\in\mathbb Z$에 대하여
표준 동형}
\begin{equation}
\label{0.13.4.5.1-fr}
\tag{13.4.5.1}
\operatorname{gr}^p(A)\xrightarrow{\sim}
\varprojlim_k\operatorname{gr}^p(A_k)
\end{equation}
\emph{이 존재한다.}
\end{corollary}
\begin{proof}
실제로 $k$가 충분히 크면
$\operatorname{Im}(A_k\to A_p)=\operatorname{Im}(A\to A_p)$이다
(Bourbaki, \emph{Top. gén.}, chap.~II,
3\textsuperscript{e} éd., §~3, n\textsuperscript{o}~5, th.~1). 따라서
결론은 (\hyperref[0.13.4.1.3-fr]{13.4.1.3})과
(\hyperref[0.13.4.4.1-fr]{13.4.4.1})에서 따른다.
\end{proof}

\subsection{여과 복합체의 스펙트럼 열들의 사영극한}
\label{subsection:0.13.5-fr}
\begin{env}[13.5.1]
\label{0.13.5.1-fr}
$C$를 아벨 범주라 하고, $X^\bullet$을 $C$의 대상들로 이루어진
복합체라 하자. $X^\bullet$에는 어떤 지표 $p_0$에 대하여
$\mathrm F^{p_0}(X^\bullet)=X^\bullet$을 만족하는 \emph{여과}
$(\mathrm F^p(X^\bullet))_{p\in\mathbb Z}$가 주어져 있다고 하자.
각 $k\in\mathbb Z$에 대하여 복합체
$X_k^\bullet=X^\bullet/\mathrm F^{k+1}(X^\bullet)$을 생각하자~; 이
복합체에는 $p\leq k$일 때의
$\mathrm F^p(X_k^\bullet)=
\mathrm F^p(X^\bullet)/\mathrm F^{k+1}(X^\bullet)$과 $p\geq k+1$일 때의
$\mathrm F^p(X_k^\bullet)=0$으로 이루어진 여과가 표준적으로 주어진다.
더 나아가 표준 사상 $X_{k+1}^\bullet\to X_k^\bullet$들이 있으며,
이들은 $\mathbf X^\bullet=(X_k^\bullet)_{k\in\mathbb Z}$를
\emph{$C$의 대상들의 여과 복합체로 이루어진 사영계}로 만든다. 이
사영계는 \emph{엄밀}하고 $k<p_0$이면 $X_k^\bullet=0$임에 유의하자.
\end{env}
\begin{env}[13.5.2]
\label{0.13.5.2-fr}
더 일반적으로, $C$의 대상들의 복합체로 이루어진 \emph{엄밀}하고
\emph{아래로 유계인} 사영계
$\mathbf X^\bullet=(X_k^\bullet)_{k\in\mathbb Z}$를 생각하자~; 각
$X_k^\bullet$ 위에는 (아래로 유계인 $C$의 복합체들의 아벨 범주에서)
(\hyperref[0.13.4.1-fr]{13.4.1})에 정의한 여과를 생각하자. 사상
$X_k^\bullet\to X_p^\bullet$ $(p\leq k)$는 유한 여과를 갖춘
\emph{여과} 복합체들의 사상이 된다. 여과 복합체의 스펙트럼 열이
갖는 함자성 (\hyperref[0.11.2.3-fr]{11.2.3})에 의해, 사영계
$\mathbf X^\bullet$을 정의하는 사상들은
$\mathrm E(\mathbf X^\bullet)=(\mathrm E(X_k^\bullet))$를
\emph{스펙트럼 열들의 사영계}로 만드는 사상들을 준다.
\end{env}
\begin{lemma}[13.5.3]
\label{0.13.5.3-fr}
\emph{여과 복합체의 사영계
$\mathbf X^\bullet=(X_k^\bullet)_{k\in\mathbb Z}$를
(\hyperref[0.13.5.2-fr]{13.5.2})에서처럼 얻었다고 가정하자. 그러면:}
\begin{enumerate}
\item[\emph{a)}]\emph{$r\geq p-p_0$이면 모든 $k\in\mathbb Z$에 대하여
$\mathrm B_r^{pq}(X_k^\bullet)=\mathrm B_\infty^{pq}(X_k^\bullet)$이다.}
\item[\emph{b)}]\emph{$k+1\leq p+r$이면
$\mathrm Z_r^{pq}(X_k^\bullet)=\mathrm Z_\infty^{pq}(X_k^\bullet)$이다.}
\item[\emph{c)}]\emph{$k+1\geq p+r$이면 모든 $h\geq k$에 대하여 사상
$\mathrm Z_r^{pq}(X_h^\bullet)\to\mathrm Z_r^{pq}(X_k^\bullet)$과
$\mathrm B_r^{pq}(X_h^\bullet)\to\mathrm B_r^{pq}(X_k^\bullet)$은
동형이다.}
\end{enumerate}
\end{lemma}
\oldpage[0\textsubscript{III}]{72}
이 세 성질은 $p-r<p_0$이면
$\mathrm F^{p-r+1}(X_k^\bullet)=X_k^\bullet$이라는 사실을 고려하면
(\hyperref[0.11.2.2-fr]{11.2.2})의 정의들에서 곧바로 따른다.

\begin{env}[13.5.4]
\label{0.13.5.4-fr}
(\hyperref[0.13.5.3-fr]{13.5.3})의 가정들이 성립한다고 하자. 그러면
고정된 $p,q,r$에 대하여($r$은 유한), 사영계들
\[
(\mathrm Z_r^{pq}(X_k^\bullet))_{k\in\mathbb Z},\quad
(\mathrm B_r^{pq}(X_k^\bullet))_{k\in\mathbb Z},\quad
(\mathrm E_r^{pq}(X_k^\bullet))_{k\in\mathbb Z}
\]
은 \emph{본질적으로 상수}이다~; 각각의 사영극한을
$\mathrm Z_r^{pq}(\mathbf X^\bullet)$,
$\mathrm B_r^{pq}(\mathbf X^\bullet)$ 및
$\mathrm E_r^{pq}(\mathbf X^\bullet)=
\mathrm Z_r^{pq}(\mathbf X^\bullet)/
\mathrm B_r^{pq}(\mathbf X^\bullet)$로 쓰겠다.
$\mathrm Z_r^{pq}(\mathbf X^\bullet)$와
$\mathrm B_r^{pq}(\mathbf X^\bullet)$는 표준적으로
$\mathrm E_1^{pq}(\mathbf X^\bullet)$의 부분대상과 동일시된다.
$d_r^{pq}$의 정의 (M, XV, 1)에 따르면 ($X_k^\bullet$들에 관한) 이
사상들 역시 본질적으로 상수이며, 따라서 사상
\begin{equation}
\label{0.13.5.4.1-fr}
\tag{13.5.4.1}
d_r^{pq}:\mathrm E_r^{pq}(\mathbf X^\bullet)
\longrightarrow\mathrm E_r^{p+r,q-r+1}(\mathbf X^\bullet)
\end{equation}
을 정의한다. 이들은 $d_r^{p+r,q-r+1}\circ d_r^{pq}=0$을 만족한다~;
더 나아가 $\operatorname{Ker}(d_r^{pq})$에서
$\mathrm Z_{r+1}^{pq}(\mathbf X^\bullet)/
\mathrm B_r^{pq}(\mathbf X^\bullet)$로 가는 표준 동형과
$\operatorname{Im}(d_r^{pq})$에서
$\mathrm B_{r+1}^{p+r,q-r+1}(\mathbf X^\bullet)/
\mathrm B_r^{p+r,q-r+1}(\mathbf X^\bullet)$로 가는 표준 동형이 있다.
\end{env}

\begin{lemma}[13.5.5]
\label{0.13.5.5-fr}
\emph{(\hyperref[0.13.5.3-fr]{13.5.3})의 가정 아래에서,
$s\geq r>p-p_0$일 때 표준 모노사상}
\begin{equation}
\label{0.13.5.5.1-fr}
\tag{13.5.5.1}
i:\mathrm E_s^{pq}(\mathbf X^\bullet)
\longrightarrow\mathrm E_r^{pq}(\mathbf X^\bullet)
\end{equation}
\emph{과 표준 동형}
\begin{equation}
\label{0.13.5.5.2-fr}
\tag{13.5.5.2}
j_r:\mathrm E_r^{pq}(\mathbf X^\bullet)
\xrightarrow{\sim}\mathrm E_\infty^{pq}(X_{p+r-1}^\bullet)
\end{equation}
\emph{이 존재하여 도식}
\begin{equation}
\label{0.13.5.5.3-fr}
\tag{13.5.5.3}
\xymatrix{
\mathrm E_s^{pq}(\mathbf X^\bullet)\ar[r]^{j_s}\ar[d]_i
&\mathrm E_\infty^{pq}(X_{p+s-1}^\bullet)\ar[d]\\
\mathrm E_r^{pq}(\mathbf X^\bullet)\ar[r]^{j_r}
&\mathrm E_\infty^{pq}(X_{p+r-1}^\bullet)
}
\end{equation}
\emph{은 가환이다} (오른쪽 수직 화살표는 사상
$X_{p+s-1}^\bullet\to X_{p+r-1}^\bullet$에서 유도된다).
\end{lemma}
\begin{proof}
$r>p-p_0$일 때
$\mathrm B_r^{pq}(X_k^\bullet)=\mathrm B_\infty^{pq}(X_k^\bullet)$이므로
(\hyperref[0.13.5.3-fr]{13.5.3}, a), $i$가 존재한다~;
$k+1\leq p+r$이면
$\mathrm Z_r^{pq}(X_k^\bullet)=\mathrm Z_\infty^{pq}(X_k^\bullet)$이므로
(\hyperref[0.13.5.3-fr]{13.5.3}, b), 특히
$\mathrm Z_r^{pq}(X_{p+r-1}^\bullet)=
\mathrm Z_\infty^{pq}(X_{p+r-1}^\bullet)$이다. 다른 한편
$\mathrm Z_r^{pq}(X_{p+r-1}^\bullet)$와
$\mathrm B_r^{pq}(X_{p+r-1}^\bullet)$는
(\hyperref[0.13.5.3-fr]{13.5.3}, c))에 의해 각각
$\mathrm Z_r^{pq}(\mathbf X^\bullet)$와
$\mathrm B_r^{pq}(\mathbf X^\bullet)$에 표준적으로 동일시된다.
이로부터 $j_r$의 존재와
(\hyperref[0.13.5.5.3-fr]{13.5.5.3})의 가환성이 따른다.
\end{proof}

\begin{corollary}[13.5.6]
\label{0.13.5.6-fr}
\emph{(\hyperref[0.13.5.3-fr]{13.5.3})의 가정 아래에서 사영극한
$\varprojlim_r\mathrm E_r^{pq}(\mathbf X^\bullet)$와
$\varprojlim_k\mathrm E_\infty^{pq}(X_k^\bullet)$ 가운데 하나가
존재하면 다른 하나도 존재하며, 표준 동형}
\begin{equation}
\label{0.13.5.6.1-fr}
\tag{13.5.6.1}
j_\infty:\varprojlim_r\mathrm E_r^{pq}(\mathbf X^\bullet)
\xrightarrow{\sim}
\varprojlim_k\mathrm E_\infty^{pq}(X_k^\bullet)
\end{equation}
\emph{이 존재한다. 더 나아가 사영계
$(\mathrm E_r^{pq}(\mathbf X^\bullet))_{r\in\mathbb Z}$가 본질적으로
상수이기 위한 필요충분조건
(\hyperref[0.13.4.2-fr]{13.4.2})은 사영계
$(\mathrm E_\infty^{pq}(X_k^\bullet))_{k\in\mathbb Z}$가
본질적으로 상수인 것이다.}
\end{corollary}

\begin{env}[13.5.7]
\label{0.13.5.7-fr}
$\mathrm E_1^{pq}(\mathbf X^\bullet)$의 부분대상 가운데 각각
$r>p-p_0$일 때의 $\mathrm B_r^{pq}(\mathbf X^\bullet)$ 및
$\inf_r\mathrm Z_r^{pq}(\mathbf X^\bullet)$와 같은 것을
($\inf_r$가 존재할 때)
$\mathrm B_\infty^{pq}(\mathbf X^\bullet)$와
$\mathrm Z_\infty^{pq}(\mathbf X^\bullet)$로 쓴다. 그러면
$\varprojlim_r\mathrm E_r^{pq}(\mathbf X^\bullet)$는 표준적으로
\oldpage[0\textsubscript{III}]{73}
$\mathrm E_\infty^{pq}(\mathbf X^\bullet)=
\mathrm Z_\infty^{pq}(\mathbf X^\bullet)/
\mathrm B_\infty^{pq}(\mathbf X^\bullet)$와 동일시된다. 대상들
$\mathrm Z_r^{pq}(\mathbf X^\bullet)$,
$\mathrm B_r^{pq}(\mathbf X^\bullet)$,
$\mathrm E_r^{pq}(\mathbf X^\bullet)$ $(1\leq r\leq+\infty)$와
$d_r^{pq}$는 (\hyperref[0.13.5.5-fr]{13.5.5})의 제한을 받는 사영계
$\mathbf X^\bullet$에 \emph{함자적으로} 의존하고,
(\hyperref[0.13.5.5-fr]{13.5.5})와
(\hyperref[0.13.5.6-fr]{13.5.6})에서 정의한 사상들도 함자적임에
유의하자.
\end{env}

\subsection{유한 여과를 갖춘 대상에 관한 함자의 스펙트럼 열}
\label{subsection:0.13.6-fr}

\begin{env}[13.6.1]
\label{0.13.6.1-fr}
$C$, $C'$을 두 아벨 범주, $\mathrm T:C\to C'$을 가법 공변 함자라
하자. $C$의 모든 대상이 어떤 단사 대상의 부분대상과 동형이라고
가정한다. 따라서 오른쪽 유도함자 $\mathrm R^p\mathrm T$ ($p\geq0$)가
존재한다.
\end{env}

\begin{lemma}[13.6.2]
\label{0.13.6.2-fr}
\emph{$A$를 유한 여과 $(\mathrm F^i(A))_{i\in\mathbb Z}$를 갖춘
$C$의 대상이라 하자. $A$의 단사 분해
$X^\bullet=(X^j)_{j\geq0}$와 그 위의 유한 여과
$(\mathrm F^i(X^\bullet))_{i\in\mathbb Z}$가 존재하여, 관계
$\mathrm F^i(A)=A$ (각각 $\mathrm F^i(A)=0$)는
$\mathrm F^i(X^\bullet)=X^\bullet$ (각각
$\mathrm F^i(X^\bullet)=0$)를 함의하고, 모든 $i\in\mathbb Z$에
대하여 $\mathrm F^i(X^\bullet)$은 $\mathrm F^i(A)$의 단사 분해이다.}
\end{lemma}
\begin{proof}
$p$를 $\mathrm F^p(A)=A$인 가장 큰 지표, $q>p$를
$\mathrm F^q(A)=0$인 가장 작은 지표라 하자. $q-p$에 관한 귀납법을
쓴다. $q-p=1$일 때 보조정리는 명백하다. 원하는 성질들을 갖는
$A/\mathrm F^{q-1}(A)$의 단사 분해 ${X''}^\bullet$을 구성했다고 하자.
완전열
\[
0\to\mathrm F^{q-1}(A)\to A\to A/\mathrm F^{q-1}(A)\to0
\]
을 생각하고, $\mathrm F^{q-1}(A)$의 단사 분해 ${X'}^\bullet$을
취한다. 이어서 앞의 완전열과 양립하는 완전열
\[
0\to {X'}^\bullet\to X^\bullet\to {X''}^\bullet\to0
\]
을 갖도록 $A$의 단사 분해 $X^\bullet$을 정한다 (M, V, 2.2).
$X^\bullet$이 요구를 만족함은 명백하다.
\end{proof}

\begin{corollary}[13.6.3]
\label{0.13.6.3-fr}
\emph{$B$를 유한 여과 $(\mathrm F^i(B))_{i\in\mathbb Z}$를 갖춘
$C$의 또 다른 대상, $s$를 정수라 하자. $u:A\to B$가 모든
$i\in\mathbb Z$에 대하여
$u(\mathrm F^i(A))\subset\mathrm F^{i+s}(B)$를 만족한다고 하자.
$Y^\bullet=(Y^j)_{j\geq0}$가
(\hyperref[0.13.6.2-fr]{13.6.2})에 열거된 성질들을 갖는 여과
$(\mathrm F^i(Y^\bullet))_{i\in\mathbb Z}$를 갖춘 $B$의 단사
분해이면, $u$와 양립하고 모든 $i\in\mathbb Z$에 대하여
$v(\mathrm F^i(X^\bullet))\subset\mathrm F^{i+s}(Y^\bullet)$를
만족하는 사상 $v:X^\bullet\to Y^\bullet$가 존재한다. 또한 그러한
두 사상 $v$, $v'$는 서로 호모토픽하다.}
\end{corollary}

이는 앞의 구성과 (M, V, 2.3)을 $q-p$에 관하여 귀납적으로 적용하면
곧바로 따른다.

\begin{env}[13.6.4]
\label{0.13.6.4-fr}
(\hyperref[0.13.6.2-fr]{13.6.2})의 가정 아래 이제 $C'$의 복합체
$\mathrm T(X^\bullet)$을 생각하자. $\mathrm F^i(X^\bullet)$이
$X^\bullet$의 직합인자이므로, 이 복합체는 명백히 복합체들
$\mathrm T(\mathrm F^i(X^\bullet))$에 의해 여과된다.
(\hyperref[0.13.6.3-fr]{13.6.3})에 따라 이 여과 복합체의
\emph{스펙트럼 열}은 동형을 제외하면 여과 대상 $A$에만 의존한다.
그 수렴대상은 코호몰로지 $\mathrm R^\bullet\mathrm T(A)$이고, 그 위의
여과는
\begin{equation}
\label{0.13.6.4.1-fr}
\tag{13.6.4.1}
\begin{split}
\mathrm F^p(\mathrm R^n\mathrm T(A))
&=\operatorname{Im}\bigl(\mathrm R^n\mathrm T(\mathrm F^p(A))
\longrightarrow\mathrm R^n\mathrm T(A)\bigr)\\
&=\operatorname{Ker}\bigl(\mathrm R^n\mathrm T(A)
\longrightarrow\mathrm R^n\mathrm T(A/\mathrm F^p(A))\bigr)
\end{split}
\end{equation}
이다 (\hyperref[0.11.2.2-fr]{11.2.2}). 그 $\mathrm E_1$항은
\begin{equation}
\label{0.13.6.4.2-fr}
\tag{13.6.4.2}
\mathrm E_1^{pq}=\mathrm R^{p+q}\mathrm T(\operatorname{gr}^p(A))
\end{equation}
로 주어진다. 여기서 통상과 같이
$\operatorname{gr}^p(A)=\mathrm F^p(A)/\mathrm F^{p+1}(A)$이다.
(\hyperref[0.11.2.2-fr]{11.2.2})에 따라 수렴대상의 여과는 유한이고,
주어진 $p$, $q$에 대하여 열
\oldpage[0\textsubscript{III}]{74}
$\mathrm B_r^{pq}(A)=\mathrm B_r^{pq}(\mathrm T(X^\bullet))$와
$\mathrm Z_r^{pq}(A)=\mathrm Z_r^{pq}(\mathrm T(X^\bullet))$는
정상적이다. 따라서 앞의 스펙트럼 열은 \emph{쌍정칙}이다
(\hyperref[0.11.1.3-fr]{11.1.3}). 이 스펙트럼 열을
$\mathrm E(A)=(\mathrm E_r^{pq}(A))$로 나타내고,
\emph{여과 대상 $A$에 관한 함자 $\mathrm T$의 스펙트럼 열}이라 한다.
\end{env}

\begin{env}[13.6.5]
\label{0.13.6.5-fr}
이제 (\hyperref[0.13.6.3-fr]{13.6.3})의 가정들이 성립한다고 하고 그
표기를 유지하자. $\mathrm F^i(X^\bullet)$ (각각
$\mathrm F^i(Y^\bullet)$)은 $X^\bullet$ (각각 $Y^\bullet$)의
직합인자이므로, 모든 $i\in\mathbb Z$에 대하여
\[
(\mathrm T v)(\mathrm T(\mathrm F^i(X^\bullet)))
\subset\mathrm T(\mathrm F^{i+s}(Y^\bullet)).
\]
그러면 (\hyperref[0.11.2.2-fr]{11.2.2})의 정의에 따라
$1\leq r\leq+\infty$일 때 $\mathrm T v$는 사상
$\mathrm B_r^{pq}(\mathrm T(X^\bullet))\to
\mathrm B_r^{p+s,q-s}(\mathrm T(Y^\bullet))$와 사상
$\mathrm Z_r^{pq}(\mathrm T(X^\bullet))\to
\mathrm Z_r^{p+s,q-s}(\mathrm T(Y^\bullet))$를 정의하며, 따라서 사상
\[
w_r:\mathrm E_r^{pq}(A)\to\mathrm E_r^{p+s,q-s}(B)
\]
을 정의한다. 마찬가지로 수렴대상 위에는
$u_n:\mathrm R^n\mathrm T(A)\to\mathrm R^n\mathrm T(B)$가 있고,
$u_n(\mathrm F^p(\mathrm R^n\mathrm T(A)))\subset
\mathrm F^{p+s}(\mathrm R^n\mathrm T(B))$이다.

$d_r^{pq}$의 정의 (M, XV, 1)에 따라 도식들
\[
\xymatrix{
\mathrm E_r^{pq}(A)\ar[r]^{d_r^{pq}}\ar[d]_{w_r}
&\mathrm E_r^{p+r,q-r+1}(A)\ar[d]^{w_r}\\
\mathrm E_r^{p+s,q-s}(B)\ar[r]_{d_r^{p+s,q-s}}
&\mathrm E_r^{p+r+s,q-r-s+1}(B)
}
\]
도 가환한다. 여기서 동형 $\alpha_r^{pq}$에 대한 비슷한 가환도식을
얻으며, 그 명시는 독자에게 맡긴다. 끝으로 (\emph{같은 곳}), 수렴대상에
대해서도 가환도식
\[
\xymatrix{
\mathrm E_\infty^{pq}(A)\ar[r]^{\beta^{pq}}\ar[d]_{w_\infty}
&\operatorname{gr}^p(\mathrm R^{p+q}\mathrm T(A))
  \ar[d]^{u_{p+q}}\\
\mathrm E_\infty^{p+s,q-s}(B)\ar[r]_{\beta^{p+s,q-s}}
&\operatorname{gr}^{p+s}(\mathrm R^{p+q}\mathrm T(B))
}
\]
을 얻는다.
\end{env}

\begin{env}[13.6.6]
\label{0.13.6.6-fr}
특히 환 $S$가 여과 $(\mathrm F^i(S))_{i\in\mathbb Z}$를 갖고,
환 준동형
\begin{equation}
\label{0.13.6.6.1-fr}
\tag{13.6.6.1}
h:S\to\operatorname{Hom}_C(A,A)
\end{equation}
이 존재하여,
\oldpage[0\textsubscript{III}]{75}
모든 $t\in\mathrm F^i(S)$와 모든 두 지표 $i$, $j$에 대하여
$h_t(\mathrm F^j(A))\subset\mathrm F^{i+j}(A)$라고 하자. 줄여서 이때
$A$는 여과환 $S$ 위의 \emph{여과 $S$-$C$-가군 구조}를 갖는다고 한다.
연관 등급 대상으로 이행하면, 모든 $t\in\mathrm F^j(S)$에 대하여
$h_t$는 $\operatorname{gr}^\bullet(A)$의 차수 $j$인 동차 등급
자기준동형 $\overline h_t$를 정의한다. 이 사상은 $t$의
$\operatorname{gr}^j(S)$에서의 류에만 의존하므로, 등급환 준동형
\[
\overline h:\operatorname{gr}^\bullet(S)\to
\operatorname{Hom}_C^g(\operatorname{gr}^\bullet(A),
\operatorname{gr}^\bullet(A))
\]
을 정의한다. 여기서 우변은 $\operatorname{gr}^\bullet(A)$의 등급
자기준동형들의 환이다. 이때 $\operatorname{gr}^\bullet(A)$가
\emph{등급 $\operatorname{gr}^\bullet(S)$-$C$-가군 구조}를 갖는다고
한다. 이제 (\hyperref[0.13.6.5-fr]{13.6.5})에 따라
$1\leq r\leq+\infty$일 때 모든
$\overline t\in\operatorname{gr}^j(S)$는 이중등급 대상들
$(\mathrm B_r^{pq}(A))_{(p,q)\in\mathbb Z\times\mathbb Z}$,
$(\mathrm Z_r^{pq}(A))_{(p,q)\in\mathbb Z\times\mathbb Z}$,
$\mathrm E_r(A)=(\mathrm E_r^{pq}(A))_{(p,q)\in\mathbb Z\times\mathbb Z}$
위에 이중차수 $(j,-j)$인 이중등급 자기준동형들을 표준적으로 정의한다.
$r$이 유한할 때 $\mathrm E_r(A)$ 위의 이 자기준동형은
$d_r^{pq}$들이 정의하는 이중등급 자기준동형과 가환한다. 이 자기준동형들은
$\operatorname{gr}^\bullet(S)$ 및 고려하는 이중등급 대상들의 덧셈에 대한
통상적인 결합법칙과 분배법칙을 만족하므로, 줄여서 이 대상들을
\emph{이중등급 $\operatorname{gr}^\bullet(S)$-$C'$-가군}이라 한다.
$\alpha_r^{pq}$들이 이 구조에 관한 동형을 정의함은 즉시 알 수 있다.
모든 정수 $n$에 대하여 $\mathrm B_r^{(n)}(A)$ (각각
$\mathrm Z_r^{(n)}(A)$, $\mathrm E_r^{(n)}(A)$)를
$p+q=n$인 $\mathrm B_r^{pq}(A)$ (각각 $\mathrm Z_r^{pq}(A)$,
$\mathrm E_r^{pq}(A)$)들로 이루어진 $\mathrm B_r^\bullet(A)$ (각각
$\mathrm Z_r^\bullet(A)$, $\mathrm E_r^\bullet(A)$)의 등급 부분대상이라
하자 ($1\leq r\leq+\infty$). 이들은 명백히
\emph{등급 $\operatorname{gr}^\bullet(S)$-$C'$-가군}이다. 끝으로 모든
$\overline t\in\operatorname{gr}^j(S)$는 모든 $n$에 대하여 등급 대상
$\operatorname{gr}^\bullet(\mathrm R^n\mathrm T(A))$ 위에 차수 $j$인
등급 자기준동형을 정의한다. 따라서 이 대상도 등급
$\operatorname{gr}^\bullet(S)$-$C'$-가군 구조를 가지며, $p+q=n$일 때
$\beta^{pq}$들은 이 구조에 관한 동형
\[
\mathrm E^{(n)}(A)\xrightarrow{\sim}
\operatorname{gr}^\bullet(\mathrm R^n\mathrm T(A))
\]
을 정의한다.

$C'$이 아벨군의 범주일 때 $S$-$C'$-가군 구조(각각 등급 또는 이중등급
$\operatorname{gr}^\bullet(S)$-$C'$-가군 구조)는 통상적인 $S$-가군
구조(각각 등급 또는 이중등급 $\operatorname{gr}^\bullet(S)$-가군
구조)에 지나지 않음에 유의하자.
\end{env}

\subsection{여러 변수의 사영극한의 유도 함자}
\label{subsection:0.13.7-fr}

\begin{env}[13.7.1]
\label{0.13.7.1-fr}
$C$, $C'$을 두 아벨 범주라 하고, $C$의 모든 대상은 어떤 단사 대상의
부분대상이라고 가정하자. $\mathrm T:C\to C'$을 가법 공변 함자라 하자.
$C$ 안의 \emph{아래로 유계인} 엄밀 사영계
$\mathbf A=(A_k)_{k\in\mathbb Z}$를 생각하자. 구체적으로 어떤 $k_0$가
존재하여 $k<k_0$일 때 $A_k=0$이라고 가정한다. 공식
(\hyperref[0.13.4.1.1-fr]{13.4.1.1})에 따라 각 $A_k$ 위에 여과
$(\mathrm F^p(A_k))_{p\in\mathbb Z}$를 이 사영계에 표준적으로
연관시킨다. 사영계가 엄밀하므로 $i\leq j\leq k+1$ 및 $h\geq k$일 때
표준 사상
\begin{equation}
\label{0.13.7.1.1-fr}
\tag{13.7.1.1}
\mathrm F^i(A_h)/\mathrm F^j(A_h)\to
\mathrm F^i(A_k)/\mathrm F^j(A_k)
\end{equation}
은 동형이다. 또한 모든 $k$에 대하여
$\mathrm F^{k_0}(A_k)=A_k$이고 $\mathrm F^{k+1}(A_k)=0$임을 상기하자.
\end{env}

\begin{env}[13.7.2]
\label{0.13.7.2-fr}
이제 각 $k$에 대하여 (\hyperref[0.13.6.2-fr]{13.6.2})의 성질들을 갖는
$A_k$의 단사 분해 $X_k^\bullet=(X_k^j)_{j\geq0}$을 구성하자. 표준
사상 $A_{k+1}\to A_k$를 사용하면
(\hyperref[0.13.6.3-fr]{13.6.3})에 따라 각 $k$에 대하여 여과들과
양립하는 복합체의 사상 $X_{k+1}^\bullet\to X_k^\bullet$을 정의할 수
있다. 이 사상들은
$\mathbf X^\bullet=(X_k^\bullet)_{k\in\mathbb Z}$를 복합체들의
사영계로 만든다.
\oldpage[0\textsubscript{III}]{76}
더 나아가 이 사영계가 엄밀하다고 가정할 수 있다. 실제로 동형
(\hyperref[0.13.7.1.1-fr]{13.7.1.1})에 따라 $A_k$는
$A_{k+1}/\mathrm F^{k+1}(A_{k+1})$와 동형이다. 따라서
$X_{k+1}^\bullet$을 구성할 때
$A_{k+1}/\mathrm F^{k+1}(A_{k+1})$의 단사 분해를
$X_k^\bullet$로 취할 수 있다. 그러면 (M, V, 2.3)에 따라 복합체의
사상 $X_{k+1}^\bullet\to X_k^\bullet$을 구성할 때 몫으로 이행하여
여과들을 보존하는 동형
\[
X_{k+1}^\bullet/\mathrm F^{k+1}(X_{k+1}^\bullet)
\xrightarrow{\sim}X_k^\bullet
\]
을 주도록 할 수 있다. 이것이
(\hyperref[0.13.5.1-fr]{13.5.1})의 조건이다.
\end{env}

\begin{env}[13.7.3]
\label{0.13.7.3-fr}
구성에 따라 복합체 $\mathrm T(X_k^\bullet)$들의 사영계
$\mathrm T(\mathbf X^\bullet)$은
(\hyperref[0.13.5.3-fr]{13.5.3})의 가정들을 만족한다. 따라서
(\hyperref[0.13.5.4-fr]{13.5.4}),
(\hyperref[0.13.5.5-fr]{13.5.5}),
(\hyperref[0.13.5.6-fr]{13.5.6})의 결과들을 스펙트럼 열
$\mathrm E(\mathrm T(X_k^\bullet))=\mathrm E(A_k)$에 적용할 수 있다.
$1\leq r\leq+\infty$일 때
$\mathrm E_r^{pq}(\mathrm T(\mathbf X^\bullet))$ 대신
$\mathrm E_r^{pq}(\mathbf A)$라 쓰고($r=+\infty$인 경우에는
(\hyperref[0.13.5.7-fr]{13.5.7}) 참조), 비슷한 표기들에도 같은 약속을
쓴다. 특히 (\hyperref[0.13.6.4.2-fr]{13.6.4.2})와 사영계
$(\operatorname{gr}^p(A_k))$가 본질적으로 상수라는 사실에 따라
\begin{equation}
\label{0.13.7.3.1-fr}
\tag{13.7.3.1}
\mathrm E_1^{pq}(\mathbf A)=
\mathrm R^{p+q}\mathrm T(\operatorname{gr}^p(\mathbf A))
\end{equation}
이다.

이 결과들과 (\hyperref[0.13.4.3-fr]{13.4.3})으로부터 다음 명제를
얻는다. 이 명제는 Shih Weishu가 다른 미출판 방법으로 처음 증명하였다.
\end{env}

\begin{proposition}[13.7.4]
\label{0.13.7.4-fr}
\emph{(Shih). --- $n$을 정수라 하자. 다음 두 조건은 동치이다.}
\begin{enumerate}
\item[\emph{a)}] \emph{$p+q=n$인 모든 쌍 $(p,q)$에 대하여 사영계
$(\mathrm E_r^{pq}(\mathbf A))_{r\geq2}$는 본질적으로 상수이다.}
\item[\emph{b)}] \emph{사영계
$\mathrm R^n\mathrm T(\mathbf A)=
(\mathrm R^n\mathrm T(A_k))_{k\in\mathbb Z}$는 (ML)을 만족한다.}
\end{enumerate}
\emph{또한 이 조건들이 성립하면 모든 $p\in\mathbb Z$에 대하여 표준
동형}
\begin{equation}
\label{0.13.7.4.1-fr}
\tag{13.7.4.1}
\operatorname{gr}^p(\mathrm R^n\mathrm T(\mathbf A))
\xrightarrow{\sim}\mathrm E_\infty^{p,n-p}(\mathbf A)
\qquad\emph{모든 $p\in\mathbb Z$에 대하여}
\end{equation}
\emph{가 성립한다.}
\end{proposition}
\begin{proof}
(\hyperref[0.13.5.6-fr]{13.5.6})에 따라 조건 \emph{a)}는 $p+q=n$일
때 사영계 $(\mathrm E_\infty^{pq}(A_k))_{k\in\mathbb Z}$가 본질적으로
상수라는 조건과 동치이다. 한편
$\operatorname{gr}^p(\mathrm R^n\mathrm T(A_k))$는 표준적으로
$\mathrm E_\infty^{p,n-p}(A_k)$와 동형이다. 따라서
(\hyperref[0.13.4.3-fr]{13.4.3})에서 \emph{a)}와 \emph{b)}의 동치가
따른다. 동형 (\hyperref[0.13.7.4.1-fr]{13.7.4.1})은 바로
(\hyperref[0.13.5.6.1-fr]{13.5.6.1})을 지금의 경우에 적용한 것이다.
\end{proof}

\begin{corollary}[13.7.5]
\label{0.13.7.5-fr}
\emph{$\mathcal F=(\mathcal F_k)_{k\in\mathbb N}$를
(\hyperref[0.13.3.1-fr]{13.3.1})의 조건 \emph{(i)}, \emph{(ii)},
\emph{(iii)}을 만족하는 아벨군층들의 사영계라 하고
$\mathcal F=\varprojlim_k\mathcal F_k$라 하자. 함자
$\mathcal G\mapsto\Gamma(X,\mathcal G)$에 대하여, $p+q=n$ 또는
$p+q=n+1$인 모든 쌍 $(p,q)$에서 사영계
$(\mathrm E_r^{pq}(\mathcal F))_{r\in\mathbb Z}$가 본질적으로
상수라고 가정하자. $H^{n+1}(X,\mathcal F)$ 위의 여과를
$\mathrm F^p(H^{n+1}(X,\mathcal F))=
\operatorname{Ker}(H^{n+1}(X,\mathcal F)\to
H^{n+1}(X,\mathcal F_{p-1}))$로 정의하면, 모든 $p\in\mathbb Z$에
대하여 표준 동형}
\begin{equation}
\label{0.13.7.5.1-fr}
\tag{13.7.5.1}
\operatorname{gr}^p(H^{n+1}(X,\mathcal F))
\xrightarrow{\sim}\mathrm E_\infty^{p,n-p+1}(\mathcal F)
\qquad\emph{모든 $p\in\mathbb Z$에 대하여}
\end{equation}
\emph{이 성립한다.}
\end{corollary}
\begin{proof}
\oldpage[0\textsubscript{III}]{77}
$C$가 $X$ 위의 아벨군층들의 범주, $C'$이 아벨군의 범주,
$\mathrm T=\Gamma$인 경우에
(\hyperref[0.13.7.4-fr]{13.7.4})를 적용하면, 모든 $p\in\mathbb Z$에
대하여 표준 동형
$\operatorname{gr}^p(\mathrm R^{n+1}\Gamma(\mathcal F))
\xrightarrow{\sim}\mathrm E_\infty^{p,n-p+1}(\mathcal F)$을 얻는다.
한편 (\hyperref[0.13.7.4-fr]{13.7.4})에 따라 사영계
$(H^n(X,\mathcal F_k))_{k\in\mathbb Z}$가 (ML)을 만족하므로,
(\hyperref[0.13.3.1-fr]{13.3.1})에서 표준 동형
\[
\label{0.13.7.5.1b-fr}
H^{n+1}(X,\mathcal F)\xrightarrow{\sim}
\varprojlim_k H^{n+1}(X,\mathcal F_k)
\]
을 얻는다.

사영계 $\mathrm R^{n+1}\Gamma(\mathcal F)$도
(\hyperref[0.13.7.4-fr]{13.7.4})에 따라 (ML)을 만족하므로 표준 동형
\[
\operatorname{gr}^p(\mathrm R^{n+1}\Gamma(\mathcal F))
\xrightarrow{\sim}
\varprojlim_k\operatorname{gr}^p(H^{n+1}(X,\mathcal F_k))
\]
이 있고 (\hyperref[0.13.4.3-fr]{13.4.3}), 표준 동형
\[
\varprojlim_k\operatorname{gr}^p(H^{n+1}(X,\mathcal F_k))
\xrightarrow{\sim}
\operatorname{gr}^p(\varprojlim_k H^{n+1}(X,\mathcal F_k))
\]
도 있다 (\hyperref[0.13.4.5-fr]{13.4.5}). 따라서 동형
(\hyperref[0.13.7.5.1-fr]{13.7.5.1})이 양변의 여과와 양립함을 보이면
충분하다. 이는 정의들과 모든 $p$에 대한 가환도식
\[
\xymatrix{
H^{n+1}(X,\mathcal F)\ar[r]^{\sim}\ar[dr]
&\varprojlim_k H^{n+1}(X,\mathcal F_k)\ar[d]\\
&H^{n+1}(X,\mathcal F_{p-1})
}
\]
에서 즉시 따른다.
\end{proof}

\begin{env}[13.7.6]
\label{0.13.7.6-fr}
$S$를 $\mathrm F^0(S)=S$인 여과 $(\mathrm F^i(S))_{i\in\mathbb Z}$를
갖는 환이라 하자. 따라서 $i<0$일 때
$\operatorname{gr}^i(S)=0$이다. 각 $A_k$가
(\hyperref[0.13.7.1-fr]{13.7.1})에서 정의한 여과에 관하여 여과
$S$-$C$-가군이고 (\hyperref[0.13.6.6-fr]{13.6.6}), $k\leq h$일 때
사상 $A_h\to A_k$는 여과 $S$-$C$-가군 구조의 사상이라고 가정하자.
줄여서 $\mathbf A$를 \emph{여과 $S$-$C$-가군들의 사영계}라 한다.
그러면 $k\leq h$, $1\leq r\leq+\infty$일 때 사상
$\mathrm B_r^{pq}(A_h)\to\mathrm B_r^{pq}(A_k)$와
$\mathrm Z_r^{pq}(A_h)\to\mathrm Z_r^{pq}(A_k)$는
(\hyperref[0.13.6.5-fr]{13.6.5})의 이중등급
$\operatorname{gr}^\bullet(S)$-$C'$-가군 구조의 사상이고, $r$이
유한할 때 족들 $(\mathrm Z_r^{pq}(\mathbf A))$,
$(\mathrm B_r^{pq}(\mathbf A))$, $(\mathrm E_r^{pq}(\mathbf A))$은
이중등급 $\operatorname{gr}^\bullet(S)$-$C'$-가군이다. 앞의 두 족은
$(\mathrm E_1^{pq}(\mathbf A))$의 부분가군이다. $p+q=n$인 항들만
취하여 얻는 앞의 대상들의 부분대상을 각각
$\mathrm Z_r^{(n)}(\mathbf A)$, $\mathrm B_r^{(n)}(\mathbf A)$,
$\mathrm E_r^{(n)}(\mathbf A)$로 나타낸다. 이들은 등급
$\operatorname{gr}^\bullet(S)$-$C'$-가군이다.

사영계 $(\mathrm E_r^{pq}(\mathbf A))_{r\in\mathbb Z}$가 본질적으로
상수이면, $(\mathrm E_\infty^{pq}(\mathbf A))$도 이중등급
$\operatorname{gr}^\bullet(S)$-$C'$-가군이고, 각
$\mathrm E_\infty^{(n)}(\mathbf A)$는 등급
$\operatorname{gr}^\bullet(S)$-$C'$-가군이다. 또한
$\beta^{p,n-p}:\mathrm E_\infty^{p,n-p}(A_k)\xrightarrow{\sim}
\operatorname{gr}^p(\mathrm R^n\mathrm T(A_k))$들은 각 $k$에 대하여
$\mathrm E_\infty^{(n)}(A_k)$에서
$\operatorname{gr}^\bullet(\mathrm R^n\mathrm T(A_k))$로 가는 등급
$\operatorname{gr}^\bullet(S)$-$C'$-가군의 동형을 이룬다. 앞의
조건들이 성립하면
\[
\beta^{p,n-p}:\mathrm E_\infty^{(n)}(\mathbf A)\xrightarrow{\sim}
\varprojlim_k\operatorname{gr}^\bullet(\mathrm R^n\mathrm T(A_k))
\]
도 이 구조에 관한 동형이다. 표준 동형
\[
\operatorname{gr}^\bullet(\mathrm R^n\mathrm T(\mathbf A))
\xrightarrow{\sim}
\varprojlim_k\operatorname{gr}^\bullet(\mathrm R^n\mathrm T(A_k))
\]
도 명백히 마찬가지이므로, 동형들
(\hyperref[0.13.7.4.1-fr]{13.7.4.1})은 등급
$\operatorname{gr}^\bullet(S)$-$C'$-가군 구조에 관한 동형을 이룬다.
\end{env}

\begin{proposition}[13.7.7]
\label{0.13.7.7-fr}
$S$를 뇌터 $\mathfrak J$-아딕 환이라 하자. $C$를 모든 대상이 어떤
단사 대상의 부분대상과 동형인 아벨 범주라 하고, $\mathrm T$를 $C$에서
아벨군의 범주로 가는 가법 공변 함자라 하자.
$\mathbf A=(A_k)_{k\in\mathbb Z}$를 $S$ 위의 $\mathfrak J$-아딕
여과에 관한 아래로 유계인 여과 $S$-$C$-가군들의 엄밀 사영계라 하자.
\oldpage[0\textsubscript{III}]{78}
주어진 정수 $n$에 대하여 다음 조건이 성립한다고 가정한다.
\begin{itemize}
\item[$(\mathrm F_n)$] (\hyperref[0.13.7.3.1-fr]{13.7.3.1})의 등급
$\operatorname{gr}^\bullet(S)$-가군
\[
\mathrm E_1^{(m)}(\mathbf A)=
\bigl(\mathrm R^m\mathrm T(\operatorname{gr}^p(\mathbf A))\bigr)_{p\in\mathbb Z}
\]
은 $m=n$과 $m=n+1$에 대하여 유한 생성이다.
\end{itemize}
이 조건들 아래 다음이 성립한다.
\begin{enumerate}
\item[\rm(i)] 사영계들
$(\mathrm R^n\mathrm T(A_k))_{k\in\mathbb Z}$와
$(\mathrm R^{n+1}\mathrm T(A_k))_{k\in\mathbb Z}$는 (ML)을 만족한다.
\item[\rm(ii)]
\[
\mathrm R'^{\,n}\mathrm T(\mathbf A)=
\varprojlim_k\mathrm R^n\mathrm T(A_k)
\]
로 놓으면 $\mathrm R'^{\,n}\mathrm T(\mathbf A)$는 유한 생성
$S$-가군이다.
\item[\rm(iii)] $\mathrm R'^{\,n}\mathrm T(\mathbf A)$ 위의 여과를
\[
\mathrm F^p(\mathrm R'^{\,n}\mathrm T(\mathbf A))=
\operatorname{Ker}\bigl(\mathrm R'^{\,n}\mathrm T(\mathbf A)
\longrightarrow\mathrm R^n\mathrm T(A_{p-1})\bigr)
\quad(p\in\mathbb Z)
\]
로 정의하면, 이 여과는 $\mathfrak J$-좋은 여과이다. 즉 모든 $p$에
대하여
$\mathfrak J\mathrm F^p(\mathrm R'^{\,n}\mathrm T(\mathbf A))
\subset\mathrm F^{p+1}(\mathrm R'^{\,n}\mathrm T(\mathbf A))$이고,
$p$가 충분히 크면 양변이 같다. 특히 이 여과가
$\mathrm R'^{\,n}\mathrm T(\mathbf A)$ 위에 정의하는 위상은
$\mathfrak J$-아딕 위상과 같다.
\item[\rm(iv)] $p+q=n$ 또는 $p+q=n+1$일 때 사영계
$(\mathrm E_r^{pq}(\mathbf A))_{r\in\mathbb Z}$는 본질적으로 상수이다.
따라서 $\mathrm E_\infty^{pq}(\mathbf A)$가 정의되고
(\hyperref[0.13.5.7-fr]{13.5.7}), 등급
$\operatorname{gr}^\bullet(S)$-가군의 표준 동형
\begin{equation}
\label{0.13.7.7.1-fr}
\tag{13.7.7.1}
\operatorname{gr}^p(\mathrm R'^{\,n}\mathrm T(\mathbf A))
\xrightarrow{\sim}\mathrm E_\infty^{p,n-p}(\mathbf A)
\quad(p\in\mathbb Z)
\end{equation}
이 성립한다.
\end{enumerate}
\end{proposition}

\begin{proof}
동형 (\hyperref[0.13.7.7.1-fr]{13.7.7.1})과 동형들
(\hyperref[0.13.7.4.1-fr]{13.7.4.1})을 고려하면, 용어를 남용하여
사영계 $\mathrm R^n\mathrm T(A_k)$의 사영극한
$\mathrm R'^{\,n}\mathrm T(\mathbf A)$을
$\mathrm R^n\mathrm T(\mathbf A)$로 나타낼 수 있다.

등급환 $\operatorname{gr}^\bullet(S)$는 뇌터이다
(Bourbaki, \emph{Alg. comm.}, chap.~III, \S~2,
n\textsuperscript{o}~9, cor.~5 du th.~2). 따라서
$\mathrm E_1^{(m)}(\mathbf A)$의 등급
$\operatorname{gr}^\bullet(S)$-부분가군들의 증가열
$\mathrm B_r^{(m)}(\mathbf A)$
(\hyperref[0.13.6.6-fr]{13.6.6})은 $m=n$과 $m=n+1$에 대하여
정상적이다. 그러므로 (\hyperref[0.11.1.10-fr]{11.1.10})의 조건 b)가
성립한다. 이에 따라 (\hyperref[0.13.7.4-fr]{13.7.4})의 조건 a)가
$n$과 $n+1$에 대하여 성립하며, 이것으로 (i)가 증명된다. 또한 동형들
(\hyperref[0.13.7.4.1-fr]{13.7.4.1})과
(\hyperref[0.13.7.6-fr]{13.7.6})의 주석에 따라
$\operatorname{gr}^\bullet(\mathrm R^n\mathrm T(\mathbf A))$는
\[
\mathrm E_\infty^{(n)}(\mathbf A)=
\mathrm Z_\infty^{(n)}(\mathbf A)/\mathrm B_\infty^{(n)}(\mathbf A)
\]
와 동형인 등급 $\operatorname{gr}^\bullet(S)$-가군이다.
$\mathrm Z_\infty^{(n)}(\mathbf A)$는
$\mathrm E_1^{(n)}(\mathbf A)$의 부분가군이므로 유한 생성이고, 따라서
$\mathrm E_\infty^{(n)}(\mathbf A)$도 유한 생성이다. 한편 여과
$(\mathrm F^p(\mathrm R'^{\,n}\mathrm T(\mathbf A)))$에 대하여
(\hyperref[0.13.4.5-fr]{13.4.5})에 따라
$\operatorname{gr}^\bullet(\mathrm R^n\mathrm T(\mathbf A))$와
$\operatorname{gr}^\bullet(\mathrm R'^{\,n}\mathrm T(\mathbf A))$는
동형인 $\operatorname{gr}^\bullet(S)$-가군이다. 이것으로 (iv)가
증명된다. 이제 (ii)와 (iii)은 앞의 결과들과 다음 보조정리에서 따른다.
\end{proof}

\begin{lemma}[13.7.7.2]
\label{0.13.7.7.2-fr}
$S$를 뇌터 $\mathfrak J$-아딕 환이라 하고, $M$을 공이산 여과
$(\mathrm F^p(M))_{p\in\mathbb Z}$를 갖는 $S$-가군이라 하자. 모든
$p$에 대하여
$\mathfrak J\mathrm F^p(M)\subset\mathrm F^{p+1}(M)$이라고 하자.
이는 $M$이 $\mathfrak J$-아딕 여과를 갖춘 여과환 $S$ 위의 여과
가군이라는 뜻이다. 더 나아가 $M$이 여과 $(\mathrm F^p(M))$이 정의하는
위상에 관하여 분리되어 있다고 가정하자. 그러면 다음 조건들은
동치이다.
\begin{enumerate}
\item[\rm a)] $M$은 유한 생성 $S$-가군이고 $(\mathrm F^p(M))$은
$\mathfrak J$-좋은 여과이다.
\item[\rm b)] $\operatorname{gr}^\bullet(M)$은 유한 생성
$\operatorname{gr}^\bullet(S)$-가군이다.
\item[\rm c)] 각 $\operatorname{gr}^p(M)$은 유한 생성 $S$-가군이고,
$p$가 충분히 크면 표준 준동형
\begin{equation}
\label{0.13.7.7.3-fr}
\tag{13.7.7.3}
\mathfrak J\otimes_S\operatorname{gr}^p(M)\longrightarrow
\operatorname{gr}^{p+1}(M)
\end{equation}
은 전사이다.
\oldpage[0\textsubscript{III}]{79}
이 준동형들은 $\mathfrak J\otimes_S\mathrm F^p(M)\to
\mathrm F^{p+1}(M)$에서 얻는다. 이때 합성 준동형
\[
\mathfrak J\otimes_S\mathrm F^{p+1}(M)\longrightarrow
\mathfrak J\otimes_S\mathrm F^p(M)\longrightarrow\mathrm F^{p+1}(M)
\]
의 상이
$\mathfrak J\mathrm F^{p+1}(M)\subset\mathrm F^{p+2}(M)$임을
고려한다.
\end{enumerate}
증명은 Bourbaki, \emph{Alg. comm.}, chap.~III, \S~3,
n\textsuperscript{o}~1, prop.~3을 보라.
\end{lemma}

\begin{env}[13.7.7.4]
\label{0.13.7.7.4-fr}
보조정리 (\hyperref[0.13.7.7.2-fr]{13.7.7.2})를 적용하려면, 고려하는
여과가 $\mathrm R'^{\,n}\mathrm T(\mathbf A)$ 위에 정의하는 위상에서
$\mathrm R'^{\,n}\mathrm T(\mathbf A)$가 분리 완비 $S$-가군임을
관찰하면 된다. 실제로 이 위상은 이산군들
$\mathrm R^n\mathrm T(A_k)$의 사영극한 위상이다. 한편 $k<k_0$일 때
$A_k=0$이면 $k<k_0$일 때 $\mathrm R^n\mathrm T(A_k)=0$도 성립하므로,
\[
\mathrm F^{k_0}(\mathrm R'^{\,n}\mathrm T(\mathbf A))=
\mathrm R'^{\,n}\mathrm T(\mathbf A).
\]
따라서 보조정리의 적용 조건들이 모두 충족된다.
\end{env}

\begin{corollary}[13.7.8]
\label{0.13.7.8-fr}
가정 $(\mathrm F_n)$이 성립하면, 모든 원소 $f\in S$에 대하여 표준 동형
\begin{equation}
\label{0.13.7.8.1-fr}
\tag{13.7.8.1}
\varprojlim_k\bigl((\mathrm R^n\mathrm T(A_k))_f\bigr)
\xrightarrow{\sim}\mathrm R^n\mathrm T(\mathbf A)\otimes_S S_{\{f\}}
\end{equation}
이 성립한다.
\end{corollary}

\begin{proof}
실제로 $\mathrm R^n\mathrm T(\mathbf A)$는 유한 생성 $S$-가군이고,
$S_{\{f\}}$는 뇌터 아딕 $S$-대수이며
(\hyperref[0.7.6.11-ko]{0\textsubscript{I}, 7.6.11}),
$S_f$의 $\mathfrak J$-준아딕 위상에 관한 분리 완비화이다
(\hyperref[0.7.6.2-ko]{0\textsubscript{I}, 7.6.2}).
(\hyperref[0.7.7.8-ko]{0\textsubscript{I}, 7.7.8})과
(\hyperref[0.7.7.1-ko]{0\textsubscript{I}, 7.7.1})에 따라
$\mathrm R^n\mathrm T(\mathbf A)\otimes_S S_{\{f\}}$는
$\mathrm R^n\mathrm T(\mathbf A)\otimes_S S_f$의
$\mathfrak J$-준아딕 위상에 관한 분리 완비화와 동형이다. 이 위상에서
0의 기본 근방계는
$(\mathfrak J^p\mathrm R^n\mathrm T(\mathbf A))\otimes_S S_f$이고,
따라서
$\mathrm F^p(\mathrm R^n\mathrm T(\mathbf A))\otimes_S S_f$도 그러한
기본계이다. 뒤의 가군은 표준 사상
$(\mathrm R^n\mathrm T(\mathbf A))_f\to
(\mathrm R^n\mathrm T(A_{p-1}))_f$의 핵이다. 그러므로
$\mathrm R^n\mathrm T(\mathbf A)\otimes_S S_f$에 연관된 분리군은
$\varprojlim_k((\mathrm R^n\mathrm T(A_k))_f)$의 어떤 부분군 $G$와
동일시된다. 그런데 사영계
$((\mathrm R^n\mathrm T(A_k))_f)$는 명백히 (ML)을 만족하고,
각 $(\mathrm R^n\mathrm T(A_k))_f$ 안에서
$(\mathrm R^n\mathrm T(\mathbf A))_f$의 상은 충분히 큰 $h\geq k$에
대한 $(\mathrm R^n\mathrm T(A_h))_f$들의 공통 상과 같다. 따라서
$G$는 $\varprojlim_k((\mathrm R^n\mathrm T(A_k))_f)$에서 조밀하다.
마지막 군은 분리 완비이므로 따름정리가 증명된다.
\end{proof}

\begin{flushright}

\textit{(계속.)}
\end{flushright}

\clearpage
\thispagestyle{empty}
\null
\clearpage
\thispagestyle{plain}
\markboth{연접층의 코호몰로지적 연구}{제III장}
\oldpage[III]{81}
\phantomsection
\label{chapter:III-fr}

\begin{center}
{\large 제III장\par}
\vspace{1.5em}
{\Large\bfseries 연접층의\par}
{\Large\bfseries 코호몰로지적 연구\par}
\vspace{1.8em}
{\bfseries 목차\par}
\end{center}

\medskip
\begin{tabular}{@{}r@{\ }p{0.82\textwidth}@{}}
\S~1. & 아핀 스킴의 코호몰로지.\\
\S~2. & 사영 사상의 코호몰로지적 연구.\\
\S~3. & 고유 사상에 대한 유한성 정리.\\
\S~4. & 고유 사상의 기본 정리~; 응용.\\
\S~5. & 연접 대수적 층의 존재 정리.\\
\S~6. & 국소 및 전역 Tor, Ext 함자~; 퀸네트 공식.\\
\S~7. & 가군의 공변 코호몰로지 함자에서의 밑변경에 대한
연구.\\
\S~8. & 사영 다발 위의 쌍대성 정리.\\
\S~9. & 상대 코호몰로지와 국소 코호몰로지~; 국소 쌍대성.\\
\S~10. & 사영 코호몰로지와 국소 코호몰로지 사이의 관계.
인자를 따른 형식적 완비화 기법.\\
\S~11. & 전역 및 국소 피카르 군\footnotemark.
\end{tabular}

\medskip
이 장에서는 연접 대수적 층의 코호몰로지에 관한 기본 정리들을 다룬다.
다만 유수 이론에서 나오는 정리들(쌍대성 정리들)은 제외하며, 이들은 뒤의
한 장에서 다룰 것이다. 앞의 정리들 가운데에는 본질적으로 여섯 개의 기본
정리가 있고, 이 장의 처음 여섯 절이 각각 이를 다룬다. 이 결과들은 뒤에서
본래 코호몰로지적인 성격을 갖지 않는 문제들에서도 핵심 도구가 될 것이며,
독자는 이미 \S~4에서 그 첫 사례들을 보게 될 것이다. \S~7은 더 기술적인
성격의 결과들을 제시하지만, 이 결과들은 응용에서 줄곧 사용된다. 마지막으로
\S\S~8부터 11까지에서는 연접층의 쌍대성과 관련되고 응용에 특히 중요한 몇몇
결과를 전개한다. 이 결과들은 일반 유수 이론보다 앞서 설명할 수 있다.

\footnotetext{제IV장은 \S\S~8부터 11까지에 의존하지 않으며, 아마도 이들보다
먼저 출판될 것이다.}

\oldpage[III]{82}
\S\S~1과 2의 내용은 J.-P. Serre에 의한 것이며, 독자는 우리가 (FAC)의
서술을 따르기만 하면 되었음을 알게 될 것이다. \S\S~8과 9 역시 (FAC)에서
착안한 것이다(다만 서로 다른 맥락 때문에 필요한 변환은 그리 자명하지 않다).
마지막으로 서론에서 말했듯이, \S~4는 Zariski의 ``정칙함수론''에서 말하는
기본적인 ``불변성 정리''를 현대 언어로 정식화한 것으로 보아야 한다.

끝으로 3.4번의 결과들(그리고 (0, 13.4부터 13.7까지)의 예비 명제들)은
제III장의 나머지 부분에서 사용되지 않으므로 첫 독해에서는 생략해도 된다.

\setcounter{section}{0}
\section{아핀 스킴의 코호몰로지}
\label{section:III.1-fr}

\subsection{외대수 복합체에 관한 복습}
\label{subsection:III.1.1-fr}

\begin{env}[1.1.1]
\label{III.1.1.1-fr}
$A$를 환이라 하고, $\mathbf f=(f_i)_{1\leq i\leq r}$를 $A$의 $r$개
원소로 이루어진 계라 하자. $\mathbf f$에 대응하는 \emph{외대수 복합체
$K_\bullet(\mathbf f)$}는 다음과 같이 정의되는 사슬 복합체이다
(G, I, 2.2). 등급 $A$-가군 $K_\bullet(\mathbf f)$는 통상적인 방식으로
등급이 매겨진 \emph{외대수 $\bigwedge(A^r)$}이고, 경계 연산자는
쌍대가군 $(A^r)^\vee$의 원소로 본 $\mathbf f$에 의한 \emph{내부곱
$i_{\mathbf f}$}이다. $i_{\mathbf f}$는 $\bigwedge(A^r)$의 차수 $-1$인
\emph{반미분}이고, $(\mathbf e_i)_{1\leq i\leq r}$가 $A^r$의 표준기저이면
$i_{\mathbf f}(\mathbf e_i)=f_i$임을 상기하자. 조건
$i_{\mathbf f}\circ i_{\mathbf f}=0$은 곧바로 확인된다.

동치인 정의는 다음과 같다. 각 $i$에 대하여 다음과 같이 정의되는 사슬
복합체 $K_\bullet(f_i)$를 생각한다.
$K_0(f_i)=K_1(f_i)=A$, $K_n(f_i)=0$ ($n\neq0,1$)로 두고, 경계 연산자는
$d_1:A\to A$가 \emph{$f_i$에 의한 곱셈}이 되도록 정의한다. 그러면
$K_\bullet(\mathbf f)$를 전체 차수를 부여한 \emph{텐서곱
$K_\bullet(f_1)\otimes K_\bullet(f_2)\otimes\cdots\otimes K_\bullet(f_r)$}
(G, I, 2.7)로 잡는다. 이 복합체와 앞에서 정의한 복합체 사이의 동형사상은
곧바로 확인된다.
\end{env}

\begin{env}[1.1.2]
\label{III.1.1.2-fr}
임의의 $A$-가군 $M$에 대하여 \emph{사슬 복합체}
\begin{equation}
\label{III.1.1.2.1-fr}
\tag{1.1.2.1}
K_\bullet(\mathbf f,M)=K_\bullet(\mathbf f)\otimes_A M
\end{equation}
와 \emph{공사슬 복합체} (G, I, 2.2)
\begin{equation}
\label{III.1.1.2.2-fr}
\tag{1.1.2.2}
K^\bullet(\mathbf f,M)=\operatorname{Hom}_A(K_\bullet(\mathbf f),M).
\end{equation}
를 정의한다. $g$가 뒤 복합체의 $k$-공사슬이고
\[
g(i_1,\ldots,i_k)=g(\mathbf e_{i_1}\wedge\cdots\wedge\mathbf e_{i_k})
\]
로 놓으면, $g$는 $[1,r]^k$에서 $M$으로 가는 \emph{교대 사상}과 동일시된다.
앞의 정의들로부터 다음을 얻는다.
\begin{equation}
\label{III.1.1.2.3-fr}
\tag{1.1.2.3}
d^kg(i_1,i_2,\ldots,i_{k+1})=
\sum_{h=1}^{k+1}(-1)^{h-1}f_{i_h}
g(i_1,\ldots,\widehat{i_h},\ldots,i_{k+1}).
\end{equation}
\end{env}

\oldpage[III]{83}
\begin{env}[1.1.3]
\label{III.1.1.3-fr}
앞의 복합체들로부터 통상적인 방식으로 \emph{$A$-호몰로지 가군과
$A$-코호몰로지 가군} (G, I, 2.2)
\begin{align}
\label{III.1.1.3.1-fr}
\tag{1.1.3.1}
H_\bullet(\mathbf f,M)&=H_\bullet(K_\bullet(\mathbf f,M)),\\
\label{III.1.1.3.2-fr}
\tag{1.1.3.2}
H^\bullet(\mathbf f,M)&=H^\bullet(K^\bullet(\mathbf f,M))
\end{align}
을 얻는다. 한편 각 사슬
$z=\sum(\mathbf e_{i_1}\wedge\cdots\wedge\mathbf e_{i_k})\otimes
z_{i_1\ldots i_k}$에 공사슬 $g_z$를 대응시켜 $A$-동형사상
$K_i(\mathbf f,M)\xrightarrow{\sim}K^{r-i}(\mathbf f,M)$을 정의한다. 여기서
$g_z(j_1,\ldots,j_{r-k})=\varepsilon z_{i_1\ldots i_k}$이고,
$(j_h)_{1\leq h\leq r-k}$는 $[1,r]$에서 엄밀히 증가하는 수열
$(i_h)_{1\leq h\leq k}$의 여수열이며,
$\varepsilon=(-1)^\nu$이다. 이때 $\nu$는 $[1,r]$의 순열
$i_1,\ldots,i_k,j_1,\ldots,j_{r-k}$의 역전수이다. $g_{dz}=d(g_z)$임을
확인할 수 있으므로 다음 동형사상을 얻는다.
\begin{equation}
\label{III.1.1.3.3-fr}
\tag{1.1.3.3}
H^i(\mathbf f,M)\xrightarrow{\sim}H_{r-i}(\mathbf f,M)
\qquad(0\leq i\leq r).
\end{equation}
이 장에서는 주로 코호몰로지 가군 $H^\bullet(\mathbf f,M)$을 다룬다.

$\mathbf f$를 고정하면 $M\mapsto H^\bullet(\mathbf f,M)$이 $A$-가군의
범주에서 등급 $A$-가군의 범주로 가며 차수 $<0$ 및 $>r$에서 영인
\emph{코호몰로지 함자} (T, II, 2.1)임은 자명하다 (G, I, 2.1).
또한
\begin{equation}
\label{III.1.1.3.4-fr}
\tag{1.1.3.4}
H^0(\mathbf f,M)=\operatorname{Hom}_A(A/(\mathbf f),M)
\end{equation}
이다. 여기서 $(\mathbf f)$는 $f_1,\ldots,f_r$로 생성되는 $A$의
아이디얼을 나타낸다. 이는 (\hyperref[III.1.1.2.3-fr]{1.1.2.3})에서
곧바로 따르며, $H^0(\mathbf f,M)$은 \emph{$(\mathbf f)$에 의해 소멸되는}
$M$의 부분가군과 동일시된다. 마찬가지로
(\hyperref[III.1.1.2.3-fr]{1.1.2.3})에 의해
\begin{equation}
\label{III.1.1.3.5-fr}
\tag{1.1.3.5}
H^r(\mathbf f,M)=M/\Bigl(\sum_{i=1}^r f_iM\Bigr)
=(A/(\mathbf f))\otimes_A M.
\end{equation}
완전성을 위해 증명을 상기할 다음의 알려진 결과를 사용할 것이다.
\end{env}

\begin{proposition}[1.1.4]
\label{III.1.1.4-fr}
$A$를 환, $\mathbf f=(f_i)_{1\leq i\leq r}$를 $A$의 원소들로 이루어진
유한족, $M$을 $A$-가군이라 하자. $1\leq i\leq r$일 때마다
$M_{i-1}=M/(f_1M+\cdots+f_{i-1}M)$에서의 곱셈 사상
$z\mapsto f_i z$가 단사이면, $i\neq r$인 모든 $i$에 대하여
$H^i(\mathbf f,M)=0$이다.
\end{proposition}

\begin{proof}
(\hyperref[III.1.1.3.3-fr]{1.1.3.3})에 의해 모든 $i>0$에 대하여
$H_i(\mathbf f,M)=0$임을 보이면 충분하다. $r$에 관하여 귀납하자.
$r=0$인 경우는 자명하다. $\mathbf f'=(f_i)_{1\leq i\leq r-1}$로 놓자.
이 족은 명제의 조건을 만족하므로 $L_\bullet=K_\bullet(\mathbf f',M)$로
놓으면 귀납가정에 의해 $i>0$에서 $H_i(L_\bullet)=0$이고,
(\hyperref[III.1.1.3.3-fr]{1.1.3.3})과
(\hyperref[III.1.1.3.5-fr]{1.1.3.5})에 의해
$H_0(L_\bullet)=M_{r-1}$이다. 줄여서
$K_\bullet=K_\bullet(f_r)=K_0\oplus K_1$로 쓰자. 여기서
$K_0=K_1=A$이고 $d_1:K_1\to K_0$는 $f_r$에 의한 곱셈이다.
정의 (\hyperref[III.1.1.1-fr]{1.1.1})에 의해
$K_\bullet(\mathbf f,M)=K_\bullet\otimes_A L_\bullet$이다. 이제 다음 보조정리를
사용한다.

\begin{lemma}[1.1.4.1]
\label{III.1.1.4.1-fr}
$K_\bullet$을 차원 $0$과 $1$ 이외에서는 영인 자유 $A$-가군들로 이루어진
사슬 복합체라 하자. $A$-가군들로 이루어진 임의의 사슬 복합체
$L_\bullet$과 임의의 첨자 $p$에 대하여 완전열
\[
0\longrightarrow H_0(K_\bullet\otimes H_p(L_\bullet))
\longrightarrow H_p(K_\bullet\otimes L_\bullet)
\longrightarrow H_1(K_\bullet\otimes H_{p-1}(L_\bullet))
\longrightarrow0
\]
이 존재한다.
\end{lemma}

\oldpage[III]{84}
이는 퀸네트 스펙트럼 열의 저차항 완전열의 특수한 경우이다
(M, XVII, 5.2 a) 및 G, I, 5.5.2). 다음과 같이 직접 증명할 수도 있다.
$K_0$과 $K_1$을 각각 차원 $\neq0$, $\neq1$에서 영인 사슬 복합체로 보자.
그러면 복합체의 완전열
\[
0\longrightarrow K_0\otimes L_\bullet
\longrightarrow K_\bullet\otimes L_\bullet
\longrightarrow K_1\otimes L_\bullet\longrightarrow0
\]
을 얻으며, 여기에 호몰로지 완전열
\[
\cdots\longrightarrow H_{p+1}(K_1\otimes L_\bullet)
\xrightarrow{\partial}H_p(K_0\otimes L_\bullet)
\longrightarrow H_p(K_\bullet\otimes L_\bullet)
\longrightarrow H_p(K_1\otimes L_\bullet)
\xrightarrow{\partial}H_{p-1}(K_0\otimes L_\bullet)
\longrightarrow\cdots.
\]
을 적용할 수 있다. 그런데 모든 $p$에 대하여
$H_p(K_0\otimes L_\bullet)=K_0\otimes H_p(L_\bullet)$이고
$H_p(K_1\otimes L_\bullet)=K_1\otimes H_{p-1}(L_\bullet)$임은 자명하다.
또한 연산자
$\partial:K_1\otimes H_p(L_\bullet)\to K_0\otimes H_p(L_\bullet)$가 바로
$d_1\otimes1$임을 곧바로 확인한다. 따라서 앞의 완전열과
$H_0(K_\bullet\otimes H_p(L_\bullet))$ 및
$H_1(K_\bullet\otimes H_{p-1}(L_\bullet))$의 정의로부터 보조정리가 따른다.

보조정리가 증명되었으므로 (\hyperref[III.1.1.4-fr]{1.1.4})의 증명의
나머지는 즉각적이다. 귀납가정과 보조정리
(\hyperref[III.1.1.4.1-fr]{1.1.4.1})에 의해 $p\geq2$에서
$H_p(K_\bullet\otimes L_\bullet)=0$이다. 더 나아가
$H_1(K_\bullet,H_0(L_\bullet))=0$임을 보이면 보조정리
(\hyperref[III.1.1.4.1-fr]{1.1.4.1})에서
$H_1(K_\bullet\otimes L_\bullet)=0$도 얻는다. 그런데 정의에 의해
$H_1(K_\bullet,H_0(L_\bullet))$은 $M_{r-1}$에서의 곱셈 사상
$z\mapsto f_rz$의 핵과 다름없고, 가정에 따라 이 핵은 영이다. 이로써
증명이 끝난다.
\end{proof}

\begin{env}[1.1.5]
\label{III.1.1.5-fr}
$\mathbf g=(g_i)_{1\leq i\leq r}$를 $A$의 $r$개 원소로 이루어진 또 다른
수열이라 하고, $\mathbf f\mathbf g=(f_ig_i)_{1\leq i\leq r}$로 놓자.
$A^r$의 $A$-선형 사상
$(x_1,\ldots,x_r)\mapsto(g_1x_1,\ldots,g_rx_r)$을 외대수
$\bigwedge(A^r)$로 표준 연장하여 복합체의 표준 준동형
\begin{equation}
\label{III.1.1.5.1-fr}
\tag{1.1.5.1}
\varphi_{\mathbf g}:K_\bullet(\mathbf f\mathbf g)\longrightarrow
K_\bullet(\mathbf f)
\end{equation}
을 정의할 수 있다. 이것이 실제로 복합체의 준동형임을 보이려면 일반적으로
$u:E\to F$가 $A$-선형 사상이고, $\mathbf x\in F^\vee$이며,
$\mathbf y={}^tu(\mathbf x)\in E^\vee$이면 다음 공식이 성립함을 확인하면
충분하다.
\begin{equation}
\label{III.1.1.5.2-fr}
\tag{1.1.5.2}
(\bigwedge u)\circ i_{\mathbf y}=i_{\mathbf x}\circ(\bigwedge u)~;
\end{equation}
실제로 양변은 $\bigwedge F$의 반미분이므로, 이들이 $F$에서 일치함만
확인하면 충분하고 이는 정의에서 곧바로 따른다.

$K_\bullet(\mathbf f)$를 $K_\bullet(f_i)$들의 텐서곱과 동일시하면
(\hyperref[III.1.1.1-fr]{1.1.1}), $\varphi_{\mathbf g}$는
$\varphi_{g_i}$들의 텐서곱이다. 여기서 $\varphi_{g_i}$는 차수 $0$에서는
항등사상이고 차수 $1$에서는 $g_i$에 의한 곱셈이다.
\end{env}

\begin{env}[1.1.6]
\label{III.1.1.6-fr}
특히 $0\leq n\leq m$인 임의의 두 정수 $m,n$에 대하여 복합체의 준동형
\begin{equation}
\label{III.1.1.6.1-fr}
\tag{1.1.6.1}
\varphi_{\mathbf f^{m-n}}:K_\bullet(\mathbf f^m)\longrightarrow
K_\bullet(\mathbf f^n)
\end{equation}
을 얻고, 따라서 준동형들
\begin{align}
\label{III.1.1.6.2-fr}
\tag{1.1.6.2}
\varphi_{\mathbf f^{m-n}} &:K^\bullet(\mathbf f^n,M)\longrightarrow
K^\bullet(\mathbf f^m,M),\\
\label{III.1.1.6.3-fr}
\tag{1.1.6.3}
\varphi_{\mathbf f^{m-n}} &:H^\bullet(\mathbf f^n,M)\longrightarrow
H^\bullet(\mathbf f^m,M)
\end{align}
을 얻는다.

\oldpage[III]{85}
뒤의 준동형들은 $p\leq n\leq m$일 때의 추이성 조건
$\varphi_{\mathbf f^{m-p}}=
\varphi_{\mathbf f^{m-n}}\circ\varphi_{\mathbf f^{n-p}}$를 자명하게 만족한다.
따라서 이들은 $A$-가군의 두 \emph{귀납계}를 정의한다. 다음과 같이 놓는다.
\begin{align}
\label{III.1.1.6.4-fr}
\tag{1.1.6.4}
C^\bullet((\mathbf f),M)&=\varinjlim_n K^\bullet(\mathbf f^n,M),\\
\label{III.1.1.6.5-fr}
\tag{1.1.6.5}
H^\bullet((\mathbf f),M)&=H^\bullet(C^\bullet((\mathbf f),M))
=\varinjlim_n H^\bullet(\mathbf f^n,M).
\end{align}
마지막 등식은 귀납극한을 취하는 것이 함자 $H^\bullet$과 교환한다는 사실에서
따른다 (G, I, 2.1). 뒤에서 (\hyperref[III.1.4.3-fr]{1.4.3})
$H^\bullet((\mathbf f),M)$이 실제로는 $A$의 \emph{아이디얼
$(\mathbf f)$}에만 의존함(더 나아가 $A$ 위의 $(\mathbf f)$-준아딕 위상에만
의존함)을 볼 것이며, 이것이 이 표기를 정당화한다.

$M\mapsto C^\bullet((\mathbf f),M)$은 완전 $A$-선형 함자이고,
$M\mapsto H^\bullet((\mathbf f),M)$은 코호몰로지 함자임이 자명하다.
\end{env}

\begin{env}[1.1.7]
\label{III.1.1.7-fr}
$\mathbf f=(f_i)\in A^r$, $\mathbf g=(g_i)\in A^r$라 하자. 외대수
$\bigwedge(A^r)$에서 벡터 $\mathbf g\in A^r$에 의한 왼쪽 곱셈을
$e_{\mathbf g}$로 나타내자. $A$-가군 $\bigwedge(A^r)$에서 다음의
\emph{호모토피 공식}이 성립함이 알려져 있다(여기서 $1$은
$\bigwedge(A^r)$의 항등 자기동형사상을 뜻한다).
\begin{equation}
\label{III.1.1.7.1-fr}
\tag{1.1.7.1}
i_{\mathbf f}e_{\mathbf g}+e_{\mathbf g}i_{\mathbf f}
=\langle\mathbf g,\mathbf f\rangle 1
\end{equation}
이 관계는 또한 \emph{복합체 $K_\bullet(\mathbf f)$ 안에서}
\begin{equation}
\label{III.1.1.7.2-fr}
\tag{1.1.7.2}
de_{\mathbf g}+e_{\mathbf g}d=\langle\mathbf g,\mathbf f\rangle 1.
\end{equation}
임을 뜻한다. 아이디얼 $(\mathbf f)$가 $A$와 같으면
$\langle\mathbf g,\mathbf f\rangle=\sum_{i=1}^r g_if_i=1$인
$\mathbf g\in A^r$가 존재한다. 따라서 (G, I, 2.4) 다음을 얻는다.
\end{env}

\begin{proposition}[1.1.8]
\label{III.1.1.8-fr}
$f_i$들이 생성하는 아이디얼 $(\mathbf f)$가 $A$와 같다고 하자. 그러면
복합체 $K_\bullet(\mathbf f)$는 호모토피적으로 자명하고, 따라서 임의의
$A$-가군 $M$에 대하여 복합체 $K_\bullet(\mathbf f,M)$과
$K^\bullet(\mathbf f,M)$도 그러하다.
\end{proposition}

\begin{corollary}[1.1.9]
\label{III.1.1.9-fr}
$(\mathbf f)=A$이면 임의의 $A$-가군 $M$에 대하여
$H^\bullet(\mathbf f,M)=0$이고 $H^\bullet((\mathbf f),M)=0$이다.
\end{corollary}

\begin{proof}
실제로 이 경우 모든 $n$에 대하여 $(\mathbf f^n)=A$이다.
\end{proof}

\begin{remark}[1.1.10]
\label{III.1.1.10-fr}
위와 같은 표기 아래 $X=\operatorname{Spec}(A)$라 하고, $Y$를 아이디얼
$(\mathbf f)$로 정의되는 $X$의 닫힌 부분스킴이라 하자. \S~9에서
$H^\bullet((\mathbf f),M)$이 $Y$의 닫힌 부분집합들의 반필터 $\Phi$에
대응하는 코호몰로지 $H_Y^\bullet(X,\widetilde M)$ (T, 3.2)과 동형임을
증명할 것이다. 또한 $X$와 $\mathcal F=\widetilde M$에 적용한 명제
(\hyperref[III.1.2.3-fr]{1.2.3})가 코호몰로지 완전열
\[
\cdots\longrightarrow H_Y^p(X,\mathcal F)\longrightarrow H^p(X,\mathcal F)
\longrightarrow H^p(X-Y,\mathcal F)\longrightarrow H_Y^{p+1}(X,\mathcal F)
\longrightarrow\cdots.
\]
의 특수한 경우임도 보일 것이다.
\end{remark}

\subsection{열린 덮개의 체흐 코호몰로지}
\label{subsection:III.1.2-fr}

\begin{notation}[1.2.1]
\label{III.1.2.1-fr}
이 번호에서는 다음 표기를 사용한다.
\begin{itemize}
\item $X$는 준스킴이다.
\item $\mathcal F$는 준연접 $\mathcal O_X$-가군이다.
\oldpage[III]{86}
\item $A=\Gamma(X,\mathcal O_X)$, $M=\Gamma(X,\mathcal F)$이다.
\item $\mathbf f=(f_i)_{1\leq i\leq r}$는 $A$의 원소들로 이루어진 유한계이다.
\item $U_i=X_{f_i}$는 $f_i(x)\neq0$인 $x\in X$들로 이루어진 열린집합
(\hyperref[0.5.5.2-ko]{0\textsubscript{I}, 5.5.2})이다.
\item $U=\bigcup_{i=1}^r U_i$이다.
\item $\mathfrak U$는 $U$의 덮개 $(U_i)_{1\leq i\leq r}$이다.
\end{itemize}
\end{notation}

\begin{env}[1.2.2]
\label{III.1.2.2-fr}
$X$가 바탕공간이 \emph{뇌터}인 준스킴이거나, 바탕공간이
\emph{준콤팩트}인 \emph{스킴}이라고 하자. 그러면
(\hyperref[I.9.3.3-ko]{I, 9.3.3}) $\Gamma(U_i,\mathcal F)=M_{f_i}$임이
알려져 있다. 다음과 같이 놓자.
\[
U_{i_0i_1\ldots i_p}=\bigcap_{k=0}^p U_{i_k}
=X_{f_{i_0}f_{i_1}\cdots f_{i_p}}
\]
(\hyperref[0.5.5.3-ko]{0\textsubscript{I}, 5.5.3})이며, 따라서 또한
\begin{equation}
\label{III.1.2.2.1-fr}
\tag{1.2.2.1}
\Gamma(U_{i_0i_1\ldots i_p},\mathcal F)=M_{f_{i_0}f_{i_1}\cdots f_{i_p}}.
\end{equation}
그런데 (\hyperref[0.1.6.1-ko]{0\textsubscript{I}, 1.6.1})
$M_{f_{i_0}f_{i_1}\cdots f_{i_p}}$는 귀납극한
$\varinjlim_n M_{i_0i_1\ldots i_p}^{(n)}$과 동일시된다. 여기서 귀납계는
$M_{i_0i_1\ldots i_p}^{(n)}=M$들로 이루어지고, $m\leq n$일 때 준동형
$\varphi_{nm}:M_{i_0i_1\ldots i_p}^{(m)}\to
M_{i_0i_1\ldots i_p}^{(n)}$은
$(f_{i_0}f_{i_1}\cdots f_{i_p})^{n-m}$에 의한 곱셈이다. $C_n^p(M)$을
모든 $n$에 대하여 $[1,r]^{p+1}$에서 $M$으로 가는 \emph{교대 사상}들의
집합이라 하자. 이 $A$-가군들도 $\varphi_{nm}$에 의해 귀납계를 이룬다.
$C^p(\mathfrak U,\mathcal F)$가 덮개 $\mathfrak U$에 관하여
$\mathcal F$를 계수로 하는 체흐의 \emph{교대} $p$-공사슬군이면
(G, II, 5.1), 앞에서 본 것으로부터
\begin{equation}
\label{III.1.2.2.2-fr}
\tag{1.2.2.2}
C^p(\mathfrak U,\mathcal F)=\varinjlim_n C_n^p(M)
\end{equation}
으로 쓸 수 있다. 그런데 (\hyperref[III.1.1.2-fr]{1.1.2})의 표기 아래
$C_n^p(M)$은 $K^{p+1}(\mathbf f^n,M)$과 동일시되고, 사상
$\varphi_{nm}$은 (\hyperref[III.1.1.6-fr]{1.1.6})에서 정의한 사상
$\varphi_{\mathbf f^{n-m}}$과 동일시된다. 따라서 모든 $p\geq0$에 대하여
$\mathcal F$에 관해 함자적인 표준 동형사상
\begin{equation}
\label{III.1.2.2.3-fr}
\tag{1.2.2.3}
C^p(\mathfrak U,\mathcal F)\xrightarrow{\sim}C^{p+1}((\mathbf f),M)
\end{equation}
을 얻는다. 더 나아가 공식 (\hyperref[III.1.1.2.3-fr]{1.1.2.3})과 덮개의
코호몰로지 정의 (G, II, 5.1)는 동형사상들
(\hyperref[III.1.2.2.3-fr]{1.2.2.3})이 공경계 연산자들과 양립함을 보인다.
\end{env}

\begin{proposition}[1.2.3]
\label{III.1.2.3-fr}
$X$가 바탕공간이 뇌터인 준스킴이거나 바탕공간이 준콤팩트인 스킴이면,
$\mathcal F$에 관해 함자적인 표준 동형사상
\begin{equation}
\label{III.1.2.3.1-fr}
\tag{1.2.3.1}
H^p(\mathfrak U,\mathcal F)\xrightarrow{\sim}H^{p+1}((\mathbf f),M)
\qquad\text{모든 }p\geq1\text{에 대하여}
\end{equation}
이 존재한다. 또한 $\mathcal F$에 관해 함자적인 완전열
\begin{equation}
\label{III.1.2.3.2-fr}
\tag{1.2.3.2}
0\longrightarrow H^0((\mathbf f),M)\longrightarrow M
\longrightarrow H^0(\mathfrak U,\mathcal F)
\longrightarrow H^1((\mathbf f),M)\longrightarrow0
\end{equation}
이 존재한다.
\end{proposition}

\oldpage[III]{87}
\begin{proof}
관계들 (\hyperref[III.1.2.3.1-fr]{1.2.3.1})은 실제로
(\hyperref[III.1.2.2-fr]{1.2.2})에서 본 것으로부터 즉시 따른다. 한편
$C^0(\mathfrak U,\mathcal F)=C^1((\mathbf f),M)$이다. 따라서
$H^0(\mathfrak U,\mathcal F)$는 $C^1((\mathbf f),M)$의 $1$-공순환들로
이루어진 부분군과 동일시된다. $M=C^0((\mathbf f),M)$이므로 완전열
(\hyperref[III.1.2.3.2-fr]{1.2.3.2})은 코호몰로지군
$H^0((\mathbf f),M)$과 $H^1((\mathbf f),M)$의 정의에서 나오는 완전열과
다름없다.
\end{proof}

\begin{corollary}[1.2.4]
\label{III.1.2.4-fr}
$X_{f_i}$들이 준콤팩트이고
$\sum_i g_i(f_i\vert U)=1\vert U$를 만족하는
$g_i\in\Gamma(U,\mathcal O_X)$들이 존재한다고 하자. 그러면 임의의
준연접 $(\mathcal O_X\vert U)$-가군 $\mathcal G$와 모든 $p>0$에 대하여
$H^p(\mathfrak U,\mathcal G)=0$이다. 더 나아가 $U=X$이면 표준 준동형
(\hyperref[III.1.2.3.2-fr]{1.2.3.2})
$\Gamma(X,\mathcal G)\to H^0(\mathfrak U,\mathcal G)$는 전단사이다.
\end{corollary}

\begin{proof}
가정에 따라 $U_i=X_{f_i}$들이 준콤팩트이므로 $U$도 그러하다. 따라서
$U=X$인 경우로 한정해도 된다. 이때 가정으로부터 모든 $p\geq0$에 대하여
$H^p((\mathbf f),M)=0$이다 (\hyperref[III.1.1.9-fr]{1.1.9}). 이제 따름정리는
(\hyperref[III.1.2.3.1-fr]{1.2.3.1})과
(\hyperref[III.1.2.3.2-fr]{1.2.3.2})에서 곧바로 따른다.
\end{proof}

$H^0(\mathfrak U,\mathcal F)=H^0(U,\mathcal F)$ (G, II, 5.2.2)이므로,
이로써 (I, 1.3.7)을 특수한 경우로 다시 증명했음을 주목하자.

\begin{remark}[1.2.5]
\label{III.1.2.5-fr}
$X$가 \emph{아핀 스킴}이라고 하자. 그러면
$U_i=X_{f_i}=D(f_i)$들과 $U_{i_0i_1\ldots i_p}$들은 아핀 열린집합이다
(다만 $U$는 반드시 아핀일 필요는 없다). 이 경우 함자
$\Gamma(X,\mathcal F)$와 $\Gamma(U_{i_0i_1\ldots i_p},\mathcal F)$는
$\mathcal F$에 관해 완전하다 (I, 1.3.11). 준연접
$\mathcal O_X$-가군들의 완전열
$0\to\mathcal F'\to\mathcal F\to\mathcal F''\to0$이 주어지면,
복합체의 열
\[
0\longrightarrow C^\bullet(\mathfrak U,\mathcal F')
\longrightarrow C^\bullet(\mathfrak U,\mathcal F)
\longrightarrow C^\bullet(\mathfrak U,\mathcal F'')
\longrightarrow0
\]
은 완전하고, 따라서 코호몰로지 완전열
\[
\cdots\longrightarrow H^p(\mathfrak U,\mathcal F')
\longrightarrow H^p(\mathfrak U,\mathcal F)
\longrightarrow H^p(\mathfrak U,\mathcal F'')
\xrightarrow{\partial}H^{p+1}(\mathfrak U,\mathcal F')
\longrightarrow\cdots.
\]
을 준다. 한편 $M'=\Gamma(X,\mathcal F')$,
$M''=\Gamma(X,\mathcal F'')$로 놓으면
$0\to M'\to M\to M''\to0$은 완전하다.
$C^\bullet((\mathbf f),M)$이 $M$에 관해 완전한 함자이므로 다음 코호몰로지
완전열도 얻는다.
\[
\cdots\longrightarrow H^p((\mathbf f),M')
\longrightarrow H^p((\mathbf f),M)
\longrightarrow H^p((\mathbf f),M'')
\xrightarrow{\partial}H^{p+1}((\mathbf f),M')
\longrightarrow\cdots.
\]
이제 도식
\[
\xymatrix{
0\ar[r]&C^\bullet(\mathfrak U,\mathcal F')\ar[r]\ar[d]
&C^\bullet(\mathfrak U,\mathcal F)\ar[r]\ar[d]
&C^\bullet(\mathfrak U,\mathcal F'')\ar[r]\ar[d]&0\\
0\ar[r]&C^\bullet((\mathbf f),M')\ar[r]
&C^\bullet((\mathbf f),M)\ar[r]
&C^\bullet((\mathbf f),M'')\ar[r]&0
}
\]
\oldpage[III]{88}
이 가환이므로, 모든 $p$에 대하여 도식들
\begin{equation}
\label{III.1.2.5.1-fr}
\tag{1.2.5.1}
\xymatrix{
H^p(\mathfrak U,\mathcal F'')\ar[r]^{\partial}\ar[d]
&H^{p+1}(\mathfrak U,\mathcal F')\ar[d]\\
H^{p+1}((\mathbf f),M'')\ar[r]^{\partial}
&H^{p+2}((\mathbf f),M')
}
\end{equation}
이 가환임을 얻는다 (G, I, 2.1.1).
\end{remark}

\subsection{아핀 스킴의 코호몰로지}
\label{subsection:III.1.3-fr}

\begin{theorem}[1.3.1]
\label{III.1.3.1-fr}
$X$를 아핀 스킴이라 하자. 임의의 준연접 $\mathcal O_X$-가군
$\mathcal F$와 모든 $p>0$에 대하여 $H^p(X,\mathcal F)=0$이다.
\end{theorem}

\begin{proof}
$\mathfrak U$를 아핀 열린집합 $X_{f_i}=D(f_i)$ ($1\leq i\leq r$)들로
이루어진 $X$의 유한 덮개라 하자. 이때 $f_i$들이
$A=\Gamma(X,\mathcal O_X)$에서 생성하는 아이디얼은 $A$와 같음이 알려져
있다. 따라서 (\hyperref[III.1.2.4-fr]{1.2.4})로부터 $p>0$일 때
$H^p(\mathfrak U,\mathcal F)=0$을 얻는다. 아핀 열린집합들로 이루어진
임의로 세분된 $X$의 유한 덮개들이 존재하므로 (I, 1.1.10), 체흐
코호몰로지의 정의 (G, II, 5.8)로부터 $p>0$일 때
$\check H^p(X,\mathcal F)=0$도 얻는다. 그런데 이는 임의의 $f\in A$에
대한 모든 준스킴 $X_f$에도 적용되므로 (I, 1.3.6), $p>0$일 때
$\check H^p(X_f,\mathcal F)=0$이다. $X_f\cap X_g=X_{fg}$이므로
(G, II, 5.9.2)에 의해 모든 $p>0$에 대하여
$H^p(X,\mathcal F)=0$도 따른다.
\end{proof}

\begin{corollary}[1.3.2]
\label{III.1.3.2-fr}
$Y$를 준스킴, $f:X\to Y$를 아핀 사상이라 하자 (II, 1.6.1).
임의의 준연접 $\mathcal O_X$-가군 $\mathcal F$와 모든 $q>0$에 대하여
$R^qf_*(\mathcal F)=0$이다.
\end{corollary}

\begin{proof}
실제로 정의에 의해 $R^qf_*(\mathcal F)$는 $Y$의 열린집합 $U$에
$H^q(f^{-1}(U),\mathcal F)$를 대응시키는 준층에 대응하는
$\mathcal O_Y$-가군이다. 그런데 아핀 열린집합들은 $Y$의 기저를 이루며,
그러한 열린집합 $U$에 대하여 $f^{-1}(U)$는 아핀이다 (II, 1.3.2).
따라서 (\hyperref[III.1.3.1-fr]{1.3.1})에 의해
$H^q(f^{-1}(U),\mathcal F)=0$이고, 이로써 따름정리가 증명된다.
\end{proof}

\begin{corollary}[1.3.3]
\label{III.1.3.3-fr}
$Y$를 준스킴, $f:X\to Y$를 아핀 사상이라 하자. 임의의 준연접
$\mathcal O_X$-가군 $\mathcal F$에 대하여 표준 준동형
$H^p(Y,f_*(\mathcal F))\to H^p(X,\mathcal F)$ (0, 12.1.3.1)는 모든
$p$에 대하여 전단사이다.
\end{corollary}

\begin{proof}
실제로 (0, 12.1.7)에 의해 합성함자 $\Gamma f_*$의 두 번째 스펙트럼
열의 가장자리 준동형
$''E_2^{p0}=H^p(Y,f_*(\mathcal F))\to H^p(X,\mathcal F)$들이 전단사임을
보이면 충분하다. 이 열의 $E_2$항은
$''E_2^{pq}=H^p(Y,R^qf_*(\mathcal F))$로 주어지므로 (G, II, 4.17.1),
(\hyperref[III.1.3.2-fr]{1.3.2})에 의해 $q>0$이면
$''E_2^{pq}=0$이다. 따라서 스펙트럼 열은 퇴화하며, 이로부터 주장이
따른다 (0, 11.1.6).
\end{proof}

\begin{corollary}[1.3.4]
\label{III.1.3.4-fr}
$f:X\to Y$를 아핀 사상, $g:Y\to Z$를 사상이라 하자.
\oldpage[III]{89}
임의의 준연접 $\mathcal O_X$-가군 $\mathcal F$에 대하여 표준 준동형
$R^pg_*(f_*(\mathcal F))\to R^p(g\circ f)_*(\mathcal F)$ (0, 12.2.5.1)는
모든 $p$에 대하여 전단사이다.
\end{corollary}

\begin{proof}
(\hyperref[III.1.3.3-fr]{1.3.3})에 의해 $Z$의 임의의 아핀 열린집합
$W$에 대하여 표준 준동형
$H^p(g^{-1}(W),f_*(\mathcal F))\to
H^p(f^{-1}(g^{-1}(W)),\mathcal F)$가 전단사임을 주목하면 충분하다.
이는 표준 준동형
$R^pg_*(f_*(\mathcal F))\to R^p(g\circ f)_*(\mathcal F)$를 정의하는
준층의 준동형이 전단사임을 증명한다 (0, 12.2.5).
\end{proof}

\subsection{임의의 준스킴의 코호몰로지에 대한 응용}
\label{subsection:III.1.4-fr}

\begin{proposition}[1.4.1]
\label{III.1.4.1-fr}
$X$를 스킴, $\mathfrak U=(U_\alpha)$를 아핀 열린집합들로 이루어진
$X$의 덮개라 하자. 임의의 준연접 $\mathcal O_X$-가군 $\mathcal F$에
대하여 ($\Gamma(X,\mathcal O_X)$ 위의) 코호몰로지 가군
$H^\bullet(X,\mathcal F)$와 $H^\bullet(\mathfrak U,\mathcal F)$는
표준적으로 동형이다.
\end{proposition}

\begin{proof}
$X$가 스킴이므로 덮개 $\mathfrak U$의 열린집합들의 임의의 유한 교집합
$V$는 아핀이다 (I, 5.5.6). 따라서
(\hyperref[III.1.3.1-fr]{1.3.1})에 의해 $q\geq1$이면
$H^q(V,\mathcal F)=0$이다. 이제 명제는 르레 정리 (G, II, 5.4.1)에서
따른다.
\end{proof}

\begin{remark}[1.4.2]
\label{III.1.4.2-fr}
$X$가 반드시 스킴이라고 가정하지 않더라도 집합 $U_\alpha$들의 유한
교집합들이 아핀이면 (\hyperref[III.1.4.1-fr]{1.4.1})의 결론이 여전히
성립함을 주목하자.
\end{remark}

\begin{corollary}[1.4.3]
\label{III.1.4.3-fr}
$X$를 바탕공간이 준콤팩트인 스킴,
$A=\Gamma(X,\mathcal O_X)$라 하고,
$\mathbf f=(f_i)_{1\leq i\leq r}$를 $X_{f_i}$들
(\hyperref[III.1.2.1-fr]{1.2.1})의 표기)이 아핀인 $A$의 원소들로
이루어진 유한 수열이라 하자. 그러면
(\hyperref[III.1.2.1-fr]{1.2.1})의 표기 아래 임의의 준연접
$\mathcal O_X$-가군 $\mathcal F$에 대하여 $\mathcal F$에 관해 함자적인
표준 동형사상
\begin{equation}
\label{III.1.4.3.1-fr}
\tag{1.4.3.1}
H^q(U,\mathcal F)\xrightarrow{\sim}H^{q+1}((\mathbf f),M)
\qquad(q\geq1\text{에 대하여})
\end{equation}
과 $\mathcal F$에 관해 함자적인 완전열
\begin{equation}
\label{III.1.4.3.2-fr}
\tag{1.4.3.2}
0\longrightarrow H^0((\mathbf f),M)\longrightarrow M
\longrightarrow H^0(U,\mathcal F)
\longrightarrow H^1((\mathbf f),M)\longrightarrow0
\end{equation}
을 얻는다.
\end{corollary}

\begin{proof}
이는 (\hyperref[III.1.4.1-fr]{1.4.1})과
(\hyperref[III.1.2.3-fr]{1.2.3})에서 곧바로 따른다.
\end{proof}

\begin{env}[1.4.4]
\label{III.1.4.4-fr}
$X$가 \emph{아핀 스킴}이면 (\hyperref[III.1.2.5-fr]{1.2.5})와
(\hyperref[III.1.4.1-fr]{1.4.1})에 의해 모든 $q\geq0$에 대하여,
준연접 $\mathcal O_X$-가군들의 완전열
$0\to\mathcal F'\to\mathcal F\to\mathcal F''\to0$에 대응하는
(\hyperref[III.1.2.5-fr]{1.2.5})의 표기 아래의 도식들
\begin{equation}
\label{III.1.4.4.1-fr}
\tag{1.4.4.1}
\xymatrix{
H^q(U,\mathcal F'')\ar[r]^{\partial}\ar[d]
&H^{q+1}(U,\mathcal F')\ar[d]\\
H^{q+1}((\mathbf f),M'')\ar[r]^{\partial}
&H^{q+2}((\mathbf f),M')
}
\end{equation}
은 (\hyperref[III.1.2.5-fr]{1.2.5})의 표기 아래에서 준연접
$\mathcal O_X$-가군들의 완전열
$0\to\mathcal F'\to\mathcal F\to\mathcal F''\to0$에 대응하며,
이 도표들은 가환이다.
\end{env}

\oldpage[III]{90}
\begin{proposition}[1.4.5]
\label{III.1.4.5-fr}
$X$를 준콤팩트 스킴, $\mathcal L$을 가역 $\mathcal O_X$-가군이라 하고,
등급환 $A_\bullet=\Gamma_\bullet(\mathcal L)$을 생각하자
(0\textsubscript{I}, 5.4.6). 그러면
\[
H^\bullet(\mathcal F,\mathcal L)
=\bigoplus_{n\in\mathbb Z}H^\bullet(X,\mathcal F\otimes\mathcal L^{\otimes n})
\]
은 등급 $A_\bullet$-가군이고, 모든 $f\in A_n$에 대하여
$(A_\bullet)_{(f)}$-가군의 표준 동형사상
\begin{equation}
\label{III.1.4.5.1-fr}
\tag{1.4.5.1}
H^\bullet(X_f,\mathcal F)
\xrightarrow{\sim}(H^\bullet(\mathcal F,\mathcal L))_{(f)}
\end{equation}
이 존재한다.
\end{proposition}

\begin{proof}
$X$가 준콤팩트 스킴이므로, 각 $i$에 대하여 제한
$\mathcal L\vert U_i$가 $\mathcal O_X\vert U_i$와 동형인 아핀
열린집합들의 한 유한 덮개 $\mathfrak U=(U_i)$를 사용하여 모든
$\mathcal O_X$-가군 $\mathcal F\otimes\mathcal L^{\otimes n}$의
코호몰로지를 계산할 수 있다 (\hyperref[III.1.4.1-fr]{1.4.1}). 이때
$U_i\cap X_f$들이 아핀 열린집합임은 즉각적이다 (I, 1.3.6). 따라서 제한
덮개 $\mathfrak U\vert X_f=(U_i\cap X_f)$를 사용하여 코호몰로지
$H^\bullet(X_f,\mathcal F\otimes\mathcal L^{\otimes n})$도 계산할 수 있다
(\hyperref[III.1.4.1-fr]{1.4.1}). 임의의 $f\in A_n$에 대하여 $f$에
의한 곱셈이 준동형
\[
C^\bullet(\mathfrak U,\mathcal F\otimes\mathcal L^{\otimes m})
\longrightarrow
C^\bullet(\mathfrak U,\mathcal F\otimes\mathcal L^{\otimes(m+n)})
\]
을 정의함은 자명하고, 이로부터 준동형
\[
H^\bullet(\mathfrak U,\mathcal F\otimes\mathcal L^{\otimes m})
\longrightarrow
H^\bullet(\mathfrak U,\mathcal F\otimes\mathcal L^{\otimes(m+n)})
\]
을 얻는다. 이로써 첫 번째 주장이 증명된다. 한편 주어진
$f\in A_n$에 대하여 (I, 9.3.2)와 (I, 1.3.9, (ii))를 함께 고려하면
$(A_\bullet)_{(f)}$-가군 복합체의 동형사상
\[
C^\bullet(\mathfrak U\vert X_f,\mathcal F)
\xrightarrow{\sim}
\left(C^\bullet\left(\mathfrak U,
\bigoplus_{n\in\mathbb Z}\mathcal F\otimes\mathcal L^{\otimes n}
\right)\right)_{(f)}
\]
을 얻는다. 이 두 복합체의 코호몰로지를 취하고, 등급
$A_\bullet$-가군의 범주에서 함자 $M\mapsto M_{(f)}$가 완전함을
상기하면 동형사상 (\hyperref[III.1.4.5.1-fr]{1.4.5.1})을 얻는다.
\end{proof}

\begin{corollary}[1.4.6]
\label{III.1.4.6-fr}
(\hyperref[III.1.4.5-fr]{1.4.5})의 가정들이 성립하고, 더 나아가
$\mathcal L=\mathcal O_X$라고 하자. $A=\Gamma(X,\mathcal O_X)$로 놓으면,
모든 $f\in A$에 대하여 $A_f$-가군의 표준 동형사상
\[
H^\bullet(X_f,\mathcal F)
\xrightarrow{\sim}(H^\bullet(X,\mathcal F))_f
\]
이 존재한다.
\end{corollary}

\begin{corollary}[1.4.7]
\label{III.1.4.7-fr}
$X$를 준콤팩트 스킴이라 하고, $f$를 $\Gamma(X,\mathcal O_X)$의 원소라
하자.
\begin{enumerate}
\item[(i)] 열린집합 $X_f$가 아핀이라고 하자. 그러면 임의의 준연접
$\mathcal O_X$-가군 $\mathcal F$, 임의의 $i>0$, 임의의
$\xi\in H^i(X,\mathcal F)$에 대하여 $f^n\xi=0$인 정수 $n>0$이 존재한다.
\item[(ii)] 역으로 $X_f$가 준콤팩트이고, $\mathcal O_X$의 임의의 준연접
아이디얼층 $\mathcal J$와 임의의 $\zeta\in H^1(X,\mathcal J)$에 대하여
$f^n\zeta=0$인 $n>0$이 존재한다고 하자. 그러면 $X_f$는 아핀이다.
\end{enumerate}
\end{corollary}

\begin{proof}
\begin{enumerate}
\item[(i)] $X_f$가 아핀이면 모든 $i>0$에 대하여
$H^i(X_f,\mathcal F)=0$이므로
(\hyperref[III.1.3.1-fr]{1.3.1}), 주장은
(\hyperref[III.1.4.6-fr]{1.4.6})에서 직접 따른다.
\item[(ii)] 세르의 판정법 (II, 5.2.1)에 따라
$\mathcal O_X\vert X_f$의 모든 준연접 아이디얼 $\mathcal K$에 대하여
$H^1(X_f,\mathcal K)=0$임을 보이면 충분하다. $X_f$는 준콤팩트 스킴
$X$ 안의 준콤팩트 열린집합이므로, $\mathcal O_X$의 준연접 아이디얼
$\mathcal J$가 존재하여 $\mathcal K=\mathcal J\vert X_f$이다
(I, 9.4.2). (\hyperref[III.1.4.6-fr]{1.4.6})에 따라
$H^1(X_f,\mathcal K)=(H^1(X,\mathcal J))_f$이고, 가정에 따라 우변은
0이다. 이로부터 결론이 따른다.
\end{enumerate}
\end{proof}

\begin{remark}[1.4.8]
\label{III.1.4.8-fr}
(\hyperref[III.1.4.7-fr]{1.4.7}, (i))가 관계 (II, 4.5.13.2)의 더
간단한 증명을 다시 줌에 유의하자.
\end{remark}

\oldpage[III]{91}
\begin{lemma}[1.4.9]
\label{III.1.4.9-fr}
$X$를 준콤팩트 스킴,
$\mathfrak U=(U_i)_{1\leq i\leq n}$를 $X$의 아핀 열린집합들에 의한
유한 덮개, $\mathcal F$를 준연접 $\mathcal O_X$-가군이라 하자.
그러면 덮개 $\mathfrak U$가 정의하는 층의 복합체
$\mathcal C^\bullet(\mathfrak U,\mathcal F)$ (G, II, 5.2)는 준연접
$\mathcal O_X$-가군이다.
\end{lemma}

\begin{proof}
정의 (G, II, 5.2)에 따라
$\mathcal C^p(\mathfrak U,\mathcal F)$는 표준 포함
$U_{i_0\ldots i_p}\to X$에 의한
$\mathcal F\vert U_{i_0\ldots i_p}$의 직상층들의 직합이다. $X$가
스킴이라는 가정에 따라 이 포함들은 아핀 사상이므로 (I, 5.5.6),
$\mathcal C^p(\mathfrak U,\mathcal F)$들은 준연접이다 (II, 1.2.6).
\end{proof}

\begin{proposition}[1.4.10]
\label{III.1.4.10-fr}
$u:X\to Y$를 분리 준콤팩트 사상이라 하자. 모든 준연접
$\mathcal O_X$-가군 $\mathcal F$에 대하여
$R^qu_*(\mathcal F)$들은 준연접 $\mathcal O_Y$-가군이다.
\end{proposition}

\begin{proof}
문제는 $Y$ 위에서 국소적이므로 $Y$가 아핀이라고 가정할 수 있다.
그러면 $X$는 유한 개의 아핀 열린집합 $U_i$ ($1\leq i\leq n$)의
합집합이다. $\mathfrak U=(U_i)$라 하자. 또한 $Y$가 스킴이므로
(I, 5.5.10)에 따라 모든 아핀 열린집합 $V\subset Y$에 대하여 표준
포함 $u^{-1}(V)\to X$는 아핀 사상이다. 따라서
(\hyperref[III.1.4.1-fr]{1.4.1})과 (G, II, 5.2)에 의해 표준 동형
\begin{equation}
\label{III.1.4.10.1-fr}
\tag{1.4.10.1}
H^\bullet(u^{-1}(V),\mathcal F)
\xrightarrow{\sim}H^\bullet(\Gamma(V,\mathcal K^\bullet))
\end{equation}
이 있다. 여기서
$\mathcal K^\bullet=u_*(\mathcal C^\bullet(\mathfrak U,\mathcal F))$로
놓았다. (\hyperref[III.1.4.9-fr]{1.4.9})와 (I, 9.2.2)에 따라
$\mathcal K^\bullet$은 준연접 $\mathcal O_Y$-가군이다. 한편
$\mathcal C^\bullet(\mathfrak U,\mathcal F)$가 층의 복합체이므로
$\mathcal K^\bullet$도 \emph{층의 복합체}이다. 따라서 코호몰로지
$\mathcal H^\bullet(\mathcal K^\bullet)$의 정의 (G, II, 4.1)에 따라
이는 준연접 $\mathcal O_Y$-가군들로 이루어진다 (I, 4.1.1). $Y$의
아핀 열린집합 $V$에 대하여 함자 $\Gamma(V,\mathcal G)$는 준연접
$\mathcal O_Y$-가군의 범주에서 $\mathcal G$에 관하여 완전하므로,
(G, II, 4.1)에 따라
\begin{equation}
\label{III.1.4.10.2-fr}
\tag{1.4.10.2}
H^\bullet(\Gamma(V,\mathcal K^\bullet))
=\Gamma(V,\mathcal H^\bullet(\mathcal K^\bullet)).
\end{equation}
끝으로 (G, II, 5.2)에서 준 표준 준동형
\[
H^\bullet(\mathfrak U,\mathcal F)\longrightarrow H^\bullet(X,\mathcal F)
\]
의 정의에 따라, $V'\subset V$가 $Y$의 또 다른 아핀 열린집합이면
도식
\[
\xymatrix{
H^\bullet(u^{-1}(V),\mathcal F)\ar[r]^{\sim}\ar[d]
&H^\bullet(\Gamma(V,\mathcal K^\bullet))\ar[d]\\
H^\bullet(u^{-1}(V'),\mathcal F)\ar[r]^{\sim}
&H^\bullet(\Gamma(V',\mathcal K^\bullet))
}
\]
은 가환한다. 따라서 앞의 결과들로부터 동형
(\hyperref[III.1.4.10.1-fr]{1.4.10.1})들이 $\mathcal O_Y$-가군의 동형
\begin{equation}
\label{III.1.4.10.3-fr}
\tag{1.4.10.3}
R^\bullet u_*(\mathcal F)
\xrightarrow{\sim}\mathcal H^\bullet(\mathcal K^\bullet)
\end{equation}
을 정의함을 얻는다.
\oldpage[III]{92}
따라서 $R^\bullet u_*(\mathcal F)$는 준연접이다.
\end{proof}

또한 (\hyperref[III.1.4.10.3-fr]{1.4.10.3}),
(\hyperref[III.1.4.10.2-fr]{1.4.10.2}),
(\hyperref[III.1.4.10.1-fr]{1.4.10.1})에서 다음을 얻는다.

\begin{corollary}[1.4.11]
\label{III.1.4.11-fr}
(\hyperref[III.1.4.10-fr]{1.4.10})의 가정 아래, $Y$의 모든 아핀
열린집합 $V$와 모든 $q\geq0$에 대하여 표준 준동형
\begin{equation}
\label{III.1.4.11.1-fr}
\tag{1.4.11.1}
H^q(u^{-1}(V),\mathcal F)\longrightarrow
\Gamma(V,R^qu_*(\mathcal F))
\end{equation}
은 동형이다.
\end{corollary}

\begin{corollary}[1.4.12]
\label{III.1.4.12-fr}
(\hyperref[III.1.4.10-fr]{1.4.10})의 가정들이 성립하고, 더 나아가
$Y$가 준콤팩트라고 하자. 그러면 정수 $r>0$이 존재하여 모든 준연접
$\mathcal O_X$-가군 $\mathcal F$와 모든 정수 $q>r$에 대하여
$R^qu_*(\mathcal F)=0$이다. $Y$가 아핀이면, $X$가 $r$개의 아핀
열린집합으로 이루어진 덮개를 갖도록 $r$을 취할 수 있다.
\end{corollary}

\begin{proof}
$Y$를 유한 개의 아핀 열린집합으로 덮을 수 있으므로
(\hyperref[III.1.4.11-fr]{1.4.11})에 의해 둘째 주장만 증명하면 된다.
$\mathfrak U$가 $X$의 $r$개의 아핀 열린집합에 의한 덮개이면
$q>r$일 때 $H^q(\mathfrak U,\mathcal F)=0$이다. 실제로
$C^q(\mathfrak U,\mathcal F)$의 코사슬은 교대적이다. 따라서 결론은
(\hyperref[III.1.4.1-fr]{1.4.1})에서 따른다.
\end{proof}

\begin{corollary}[1.4.13]
\label{III.1.4.13-fr}
(\hyperref[III.1.4.10-fr]{1.4.10})의 가정들이 성립하고, 더 나아가
$Y=\operatorname{Spec}(A)$가 아핀이라고 하자. 그러면 모든 준연접
$\mathcal O_X$-가군 $\mathcal F$와 모든 $f\in A$에 대하여 표준 동형을
제외하면
\[
\Gamma(Y_f,R^qu_*(\mathcal F))
=(\Gamma(Y,R^qu_*(\mathcal F)))_f
\]
이다.
\end{corollary}

\begin{proof}
이는 $R^qu_*(\mathcal F)$가 준연접 $\mathcal O_Y$-가군이라는 사실에서
따른다 (I, 1.3.7).
\end{proof}

\begin{proposition}[1.4.14]
\label{III.1.4.14-fr}
$f:X\to Y$를 분리 준콤팩트 사상, $g:Y\to Z$를 아핀 사상이라 하자.
모든 준연접 $\mathcal O_X$-가군 $\mathcal F$와 모든 $p$에 대하여
표준 준동형
$R^p(g\circ f)_*(\mathcal F)\to g_*(R^pf_*(\mathcal F))$
(0, 12.2.5.2)은 전단사이다.
\end{proposition}

\begin{proof}
실제로 $Z$의 모든 아핀 열린집합 $W$에 대하여 $g^{-1}(W)$는 $Y$의
아핀 열린집합이다. 그러므로 표준 준동형
\[
R^p(g\circ f)_*(\mathcal F)\longrightarrow g_*(R^pf_*(\mathcal F))
\]
을 정의하는 전층의 준동형 (0, 12.2.5)은
(\hyperref[III.1.4.11-fr]{1.4.11})에 따라 전단사이다.
\end{proof}

\begin{proposition}[1.4.15]
\label{III.1.4.15-fr}
$u:X\to Y$를 분리 준콤팩트 사상, $v:Y'\to Y$를 준스킴의 평탄 사상이라
하자 (0\textsubscript{I}, 6.7.1). $u'=u_{(Y')}$라 하여 가환도식
\begin{equation}
\label{III.1.4.15.1-fr}
\tag{1.4.15.1}
\xymatrix{
X\ar[d]_u & X'=X_{(Y')}\ar[l]_{v'}\ar[d]^{u'}\\
Y & Y'\ar[l]_v
}
\end{equation}
을 얻는다. 모든 준연접 $\mathcal O_X$-가군 $\mathcal F$와 모든
$q\geq0$에 대하여
$\mathcal F'=v'^*(\mathcal F)=\mathcal F\otimes_{\mathcal O_Y}\mathcal O_{Y'}$
로 놓으면, $R^qu_*'(\mathcal F')$은
$R^qu_*(\mathcal F)\otimes_{\mathcal O_Y}\mathcal O_{Y'}
=v^*(R^qu_*(\mathcal F))$와 표준적으로 동형이다.
\end{proposition}

\oldpage[III]{93}
\begin{proof}
표준 준동형
$\rho:\mathcal F\to v_*'(v'^*(\mathcal F))$
(0\textsubscript{I}, 4.4.3.2)는 함자성에 따라 준동형
\begin{equation}
\label{III.1.4.15.2-fr}
\tag{1.4.15.2}
R^qu_*(\mathcal F)\longrightarrow R^qu_*(v_*'(\mathcal F'))
\end{equation}
을 정의한다. 한편 $w=u\circ v'=v\circ u'$로 놓으면 표준 준동형들
(0, 12.2.5.1 및 12.2.5.2)
\begin{equation}
\label{III.1.4.15.3-fr}
\tag{1.4.15.3}
R^qu_*(v_*'(\mathcal F'))\longrightarrow R^qw_*(\mathcal F')
\longrightarrow v_*(R^qu_*'(\mathcal F'))
\end{equation}
이 있다. (\hyperref[III.1.4.15.3-fr]{1.4.15.3})과
(\hyperref[III.1.4.15.2-fr]{1.4.15.2})를 합성하여 준동형
\[
\psi:R^qu_*(\mathcal F)\longrightarrow v_*(R^qu_*'(\mathcal F'))
\]
를 얻고, 끝으로 표준 준동형
\begin{equation}
\label{III.1.4.15.4-fr}
\tag{1.4.15.4}
\psi^\sharp:v^*(R^qu_*(\mathcal F))\longrightarrow R^qu_*'(\mathcal F')
\end{equation}
을 얻는다. 이 정의에는 \emph{$v$에 관한 어떠한 가정도} 들어가지
않는다. $v$가 \emph{평탄}일 때 이것이 동형임을 증명해야 한다.
문제는 $Y$와 $Y'$ 위에서 국소적이므로
$Y=\operatorname{Spec}(A)$, $Y'=\operatorname{Spec}(B)$라고 가정할 수
있다. 다음 보조정리도 사용한다.

\begin{lemma}[1.4.15.5]
\label{III.1.4.15.5-fr}
$\varphi:A\to B$를 환 준동형, $Y=\operatorname{Spec}(A)$,
$X=\operatorname{Spec}(B)$, $f:X\to Y$를 $\varphi$에 대응하는 사상,
$M$을 $B$-가군이라 하자. $\mathcal O_X$-가군 $\widetilde M$이
$f$-평탄이기 위한 필요충분조건은 $M$이 평탄 $A$-가군인 것이다
(0\textsubscript{I}, 6.7.1). 특히 사상 $f$가 평탄이기 위한
필요충분조건은 $B$가 평탄 $A$-가군인 것이다.
\end{lemma}

실제로 이는 정의 (0\textsubscript{I}, 6.7.1)와
(0\textsubscript{I}, 6.3.3)에서 따르며, 이때 (I, 1.3.4)를 고려한다.

이제 (\hyperref[III.1.4.11.1-fr]{1.4.11.1})과 준동형들
(\hyperref[III.1.4.15.3-fr]{1.4.15.3})의 정의(0, 12.2.5 참조)에
따라 $\psi$는 합성사상
\[
H^q(X,\mathcal F)\xrightarrow{\rho_q}
H^q(X,v_*'(v'^*(\mathcal F)))\xrightarrow{\theta_q}
H^q(X',v'^*(v_*'(v'^*(\mathcal F))))\xrightarrow{\sigma_q}
H^q(X',v'^*(\mathcal F))
\]
에 대응한다. 여기서 $\rho_q$와 $\sigma_q$는 표준 준동형 $\rho$와
$\sigma:v'^*(v_*'(\mathcal G'))\to\mathcal G'$가 코호몰로지 위에
정의하는 준동형이고, $\theta_q$는
$\mathcal O_X$-가군 $v_*'(v'^*(\mathcal F))$에 관한
$\varphi$-사상 (0, 12.1.3.1)이다. $\theta_q$의 함자성에 따라 도식
\[
\xymatrix{
H^q(X,\mathcal F)\ar[rr]^{\rho_q}\ar[dd]_{\theta_q}
&&H^q(X,v_*'(v'^*(\mathcal F)))\ar[dd]^{\theta_q}\\\\
H^q(X',v'^*(\mathcal F))\ar[rr]^{v'^*(\rho_q)}
&&H^q(X',v'^*(v_*'(v'^*(\mathcal F))))
}
\]
은 가환한다.
\oldpage[III]{94}
정의 (0\textsubscript{I}, 4.4.3)에 따라 $v'^*(\rho)$는
$\sigma$의 역이므로, 앞의 합성사상은 결국 $\theta_q$이다. 따라서
$\psi^\sharp$은 연관된 $B$-준동형
\[
H^q(X,\mathcal F)\otimes_A B\longrightarrow H^q(X',\mathcal F')
\]
이다. $u$가 준콤팩트이므로 $X$는 유한 개의 아핀 열린집합 $U_i$
($1\leq i\leq r$)의 합집합이다. $\mathfrak U=(U_i)$라 하자. 한편
$v$가 아핀 사상이므로 $v'$도 아핀 사상이고 (II, 1.6.2, (iii)),
$U_i'=v'^{-1}(U_i)$들은 $X'$의 아핀 열린 덮개 $\mathfrak U'$를 이룬다.
이때 (0, 12.1.4.2)에 따라 도식
\[
\xymatrix{
H^q(\mathfrak U,\mathcal F)\ar[r]^{\theta_q}\ar[d]
&H^q(\mathfrak U',\mathcal F')\ar[d]\\
H^q(X,\mathcal F)\ar[r]^{\theta_q}&H^q(X',\mathcal F')
}
\]
은 가환하고, $X$와 $X'$이 스킴이므로 수직 화살표들은 동형이다
(\hyperref[III.1.4.1-fr]{1.4.1}). 따라서 표준 $\varphi$-사상
$\theta_q:H^q(\mathfrak U,\mathcal F)\to
H^q(\mathfrak U',\mathcal F')$에 연관된 $B$-준동형
\[
H^q(\mathfrak U,\mathcal F)\otimes_A B
\longrightarrow H^q(\mathfrak U',\mathcal F')
\]
이 동형임을 보이면 충분하다. $[1,r]$의 $p+1$개 지표로 이루어진 임의의
열 $\mathbf s=(i_k)_{0\leq k\leq p}$에 대하여
\[
U_{\mathbf s}=\bigcap_{k=0}^pU_{i_k},\qquad
U_{\mathbf s}'=\bigcap_{k=0}^pU_{i_k}'=v'^{-1}(U_{\mathbf s}),\qquad
M_{\mathbf s}=\Gamma(U_{\mathbf s},\mathcal F),\qquad
M_{\mathbf s}'=\Gamma(U_{\mathbf s}',\mathcal F')
\]
로 놓자. 표준 사상 $M_{\mathbf s}\otimes_A B\to M_{\mathbf s}'$은
동형이다 (I, 1.6.5). 따라서 표준 사상
$C^p(\mathfrak U,\mathcal F)\otimes_A B\to
C^p(\mathfrak U',\mathcal F')$도 동형이고, 이 동형 아래
$d\otimes1$은 코경계 연산자
$C^p(\mathfrak U',\mathcal F')\to C^{p+1}(\mathfrak U',\mathcal F')$와
동일시된다. $B$가 \emph{평탄} $A$-가군이므로, 코호몰로지 가군의
정의에서 표준 사상
$H^q(\mathfrak U,\mathcal F)\otimes_A B\to
H^q(\mathfrak U',\mathcal F')$이 동형임이 곧바로 따른다
(0\textsubscript{I}, 6.1.1). 이 결과는 뒤에서 일반화할 것이다
(\S~6).
\end{proof}

\begin{corollary}[1.4.16]
\label{III.1.4.16-fr}
$A$를 환, $X$를 유한형 $A$-스킴, $B$를 $A$ 위의 충실 평탄
$A$-대수라 하자. $X$가 아핀이기 위한 필요충분조건은
$X\otimes_A B$가 아핀인 것이다.
\end{corollary}

\begin{proof}
조건의 필요성은 명백하다 (I, 3.2.2). 충분성을 보이자. $X$는 $A$ 위에서
분리이고 사상 $\operatorname{Spec}(B)\to\operatorname{Spec}(A)$는
평탄하므로, (\hyperref[III.1.4.15-fr]{1.4.15})에 따라 모든 $i\geq0$과
모든 준연접 $\mathcal O_X$-가군 $\mathcal F$에 대하여
\begin{equation}
\label{III.1.4.16.1-fr}
\tag{1.4.16.1}
H^i(X\otimes_A B,\mathcal F\otimes_A B)
=H^i(X,\mathcal F)\otimes_A B
\end{equation}
은 모든 $i\geq0$과 모든 준연접 $\mathcal O_X$-가군 $\mathcal F$에
대하여 성립한다. $X\otimes_A B$가 아핀이면
(\hyperref[III.1.4.16.1-fr]{1.4.16.1})의 좌변은 $i=1$일 때 0이다.
$B$가 충실 평탄 $A$-가군이므로 $H^1(X,\mathcal F)=0$도 성립한다.
$X$가 준콤팩트 스킴이므로 세르의 판정법 (II, 5.2.1)으로 결론을 얻는다.
\end{proof}

\begin{proposition}[1.4.17]
\label{III.1.4.17-fr}
$X$를 준스킴,
$0\to\mathcal F\xrightarrow{u}\mathcal G\xrightarrow{v}\mathcal H\to0$
을 $\mathcal O_X$-가군들의 완전열이라 하자. $\mathcal F$와
$\mathcal H$가 준연접이면 $\mathcal G$도 준연접이다.
\end{proposition}

\oldpage[III]{95}
\begin{proof}
문제는 $X$ 위에서 국소적이므로 $X=\operatorname{Spec}(A)$가
아핀이라고 가정할 수 있다. 그러면 $\mathcal G$가 (I, 1.4.1)의 조건
d~1)과 d~2)($V=X$)를 만족함을 보이면 충분하다. d~2)의 확인은
즉시 된다. 실제로 $t\in\Gamma(X,\mathcal G)$의 $D(f)$에 대한 제한이
0이면, 그 상 $v(t)\in\Gamma(X,\mathcal H)$의 제한도 0이다. 따라서
어떤 $m>0$에 대하여 $f^mv(t)=v(f^mt)=0$이다 (I, 1.4.1).
$\Gamma$가 왼쪽 완전하므로 $f^mt=u(s)$인
$s\in\Gamma(X,\mathcal F)$가 존재한다. $u$가 단사이므로 $s$의
$D(f)$에 대한 제한은 0이고, 이에 따라 (I, 1.4.1) 어떤 정수 $n>0$에
대하여 $f^ns=0$이다. 끝으로
$f^{m+n}t=u(f^ns)=0$을 얻는다.

이제 d~1)을 확인하자. $t'\in\Gamma(D(f),\mathcal G)$라 하자.
$\mathcal H$가 준연접이므로 어떤 정수 $m$에 대하여
$f^mv(t')=v(f^mt')$는 어떤 단면 $z\in\Gamma(X,\mathcal H)$로
연장된다 (I, 1.4.1). 그런데 준연접 $\mathcal O_X$-가군
$\mathcal F$에 (\hyperref[III.1.3.1-fr]{1.3.1}) 또는
(I, 5.1.9.2)를 적용하면 열
$\Gamma(X,\mathcal G)\to\Gamma(X,\mathcal H)\to0$은 완전하다.
따라서 $z=v(t)$인 $t\in\Gamma(X,\mathcal G)$가 존재한다.
$t''$를 $t$의 $D(f)$에 대한 제한이라 하면
$v(f^mt'-t'')=0$이다. 그러므로
$f^mt'-t''=u(s')$인 $s'\in\Gamma(D(f),\mathcal F)$가 존재한다.
$\mathcal F$가 준연접이므로 어떤 정수 $n>0$에 대하여 $f^ns'$은
어떤 단면 $s\in\Gamma(X,\mathcal F)$로 연장된다. 관계
$f^{m+n}t'-f^nt''=u(f^ns')$에 따라 $f^{m+n}t'$는 단면
$f^nt+u(s)\in\Gamma(X,\mathcal G)$의 $D(f)$에 대한 제한이다. 이로써
증명이 끝난다.
\end{proof}

\section{사영 사상의 코호몰로지적 연구}
\label{section:III.2-fr}

\subsection{몇몇 코호몰로지 군의 명시적 계산}
\label{subsection:III.2.1-fr}

\begin{env}[2.1.1]
\label{III.2.1.1-fr}
$X$를 준스킴, $\mathcal L$을 가역 $\mathcal O_X$-가군이라 하고, 등급환
(0\textsubscript{I}, 5.4.6)
\begin{equation}
\label{III.2.1.1.1-fr}
\tag{2.1.1.1}
S=\Gamma_\bullet(X,\mathcal L)
=\bigoplus_{n\in\mathbb Z}\Gamma(X,\mathcal L^{\otimes n})
\end{equation}
을 생각하자. $(f_i)_{1\leq i\leq r}$를 $S$의 \emph{동차} 원소들의
유한족이라 하고 $f_i\in S_{d_i}$라 하자. $U_i=X_{f_i}$,
$U=\bigcup_iU_i$라 두고, $\mathfrak U$로 $U$의 덮개 $(U_i)$를
나타내자. 모든 준연접 $\mathcal O_X$-가군 $\mathcal F$에 대하여
다음과 같이 둔다.
\begin{align}
\label{III.2.1.1.2-fr}
\tag{2.1.1.2}
H^\bullet(U,\mathcal F(*))
&=\bigoplus_{n\in\mathbb Z}
H^\bullet(U,\mathcal F\otimes\mathcal L^{\otimes n}),\\
\label{III.2.1.1.3-fr}
\tag{2.1.1.3}
H^\bullet(\mathfrak U,\mathcal F(*))
&=\bigoplus_{n\in\mathbb Z}
H^\bullet(\mathfrak U,\mathcal F\otimes\mathcal L^{\otimes n}).
\end{align}
아벨군 (\hyperref[III.2.1.1.2-fr]{2.1.1.2})와
(\hyperref[III.2.1.1.3-fr]{2.1.1.3})은 다음과 같이 \emph{이중등급}을
갖는다.
\[
(H^\bullet(U,\mathcal F(*)))_{mn}
=H^m(U,\mathcal F\otimes\mathcal L^{\otimes n}).
\]
(\hyperref[III.2.1.1.3-fr]{2.1.1.3})도 마찬가지로 정의한다. 둘째
차수에 관하여 이 군들이 등급 $S$-가군임은 명백하다. 예를 들어 이는
$\mathcal F\mapsto H^m(U,\mathcal F)$와
$\mathcal F\mapsto H^m(\mathfrak U,\mathcal F)$가 함자라는 사실에서
따른다.
\end{env}

\begin{env}[2.1.2]
\label{III.2.1.2-fr}
이제 등급 $S$-가군 (0\textsubscript{I}, 5.4.6)
\begin{equation}
\label{III.2.1.2.1-fr}
\tag{2.1.2.1}
M=\Gamma_\bullet(\mathcal L,\mathcal F)
=H^0(X,\mathcal F(*))
=\bigoplus_{n\in\mathbb Z}
\Gamma(X,\mathcal F\otimes\mathcal L^{\otimes n})
\end{equation}
을 생각하자.
\oldpage[III]{96}
$X$가 바탕공간이 뇌터인 준스킴이거나 준콤팩트 스킴이면, 통상과 같이
$U_{i_0i_1\ldots i_p}=\bigcap_{k=0}^pU_{i_k}$로 둘 때 (I, 9.3.1)로부터
표준 동형을 제외하면
\[
\Gamma(U_{i_0i_1\ldots i_p},\mathcal F(*))
=H^0(U_{i_0i_1\ldots i_p},\mathcal F(*))
=M_{f_{i_0}f_{i_1}\cdots f_{i_p}}
\]
이다. 또한 (\hyperref[III.1.2.2-fr]{1.2.2})의 표기로
$M_{f_{i_0}f_{i_1}\cdots f_{i_p}}$를
$\varinjlim_n M_{i_0i_1\ldots i_p}^{(n)}$와 식별할 수 있다. 동차 원소
$z\in\varinjlim_n M_{i_0i_1\ldots i_p}^{(n)}$의 차수를 다음과 같이
정의하면 이 식별은 등급 $S$-가군의 동형사상이다. $z$가 차수 $m$인
동차 원소
$x\in M_{i_0i_1\ldots i_p}^{(n)}=M$의 표준 상이라고 하자. 이때
$z$의 차수를
$m-n(d_{i_0}+d_{i_1}+\cdots+d_{i_p})$로 둔다. 준동형
$\varphi_{kn}:M_{i_0i_1\ldots i_p}^{(n)}\to
M_{i_0i_1\ldots i_p}^{(k)}$의 정의
(\hyperref[III.1.2.2-fr]{1.2.2})를 고려하면, 이 정의가 택한 $z$의
``대표원'' $x$에 의존하지 않음을 즉시 알 수 있다.
(\hyperref[III.1.2.2-fr]{1.2.2})에서와 같이 $C_n^p(M)$으로
$[1,r]^{p+1}$에서 $M$으로 가는 교대사상들의 집합을 나타내자(모든
$n$에 대하여). 위와 같은 방식으로 $\varinjlim_n C_n^p(M)$에 등급
$S$-가군 구조를 정의한다. 다시
(\hyperref[III.1.2.2-fr]{1.2.2})에서와 같이
\begin{equation}
\label{III.2.1.2.2-fr}
\tag{2.1.2.2}
C^p(\mathfrak U,\mathcal F(*))=\varinjlim_n C_n^p(M)
\end{equation}
이고, 두 변의 동형사상은 \emph{차수를 보존한다}. 그러면
(\hyperref[III.1.2.2-fr]{1.2.2})에서와 같이
\begin{equation}
\label{III.2.1.2.3-fr}
\tag{2.1.2.3}
C^p(\mathfrak U,\mathcal F(*))
=C^{p+1}((\mathbf f),M)
=\varinjlim_n K^{p+1}(\mathbf f^n,M)
\end{equation}
이고, 이 동형사상도 차수를 보존한다. 즉 같은 동차 성분 $M_m$에 값
$g(i_0,\ldots,i_p)$들을 갖는 공사슬
$g\in K^{p+1}(\mathbf f^n,M)$의 표준 상으로 주어진
$\varinjlim_n K^{p+1}(\mathbf f^n,M)$의 원소의 차수는
$m-n(d_{i_0}+\cdots+d_{i_p})$이며, 이 원소의 대표원으로 택한 공사슬에
의존하지 않는다.
\end{env}

앞의 동형사상들이 공경계 연산자들과 양립하므로,
(\hyperref[III.1.2.2-fr]{1.2.2})에서와 같이 다음을 얻는다.

\begin{proposition}[2.1.3]
\label{III.2.1.3-fr}
$X$를 바탕공간이 뇌터인 준스킴이거나 준콤팩트 스킴이라 하자.
$\mathcal F$에 관해 함자적인 표준 동형
\begin{equation}
\label{III.2.1.3.1-fr}
\tag{2.1.3.1}
H^p(\mathfrak U,\mathcal F(*))
\xrightarrow{\sim}H^{p+1}((\mathbf f),M)
\qquad\text{모든 }p\geq1\text{에 대하여}
\end{equation}
이 존재한다. 또한 $\mathcal F$에 관해 함자적인 완전열
\begin{equation}
\label{III.2.1.3.2-fr}
\tag{2.1.3.2}
0\longrightarrow H^0((\mathbf f),M)\longrightarrow M
\longrightarrow H^0(\mathfrak U,\mathcal F(*))
\longrightarrow H^1((\mathbf f),M)\longrightarrow0
\end{equation}
이 있다. 더 나아가, 도입한 모든 준동형은 등급 $S$-가군 구조에 관해
차수 $0$이다. 여기서 $S$는 환
(\hyperref[III.2.1.1.1-fr]{2.1.1.1})이다.
\end{proposition}

\providecommand{\bbGamma}{\boldsymbol{\Gamma}}
\providecommand{\shProj}{\mathcal{P}\mathrm{roj}}

\oldpage[III]{97}
\begin{corollary}[2.1.4]
\label{III.2.1.4-fr}
$X$가 준콤팩트 스킴이고 $U_i=X_{f_i}$들이 아핀이면,
$\mathcal F$에 관해 함자적인 차수 $0$의 표준 동형
\begin{equation}
\label{III.2.1.4.1-fr}
\tag{2.1.4.1}
H^p(U,\mathcal F(*))\xrightarrow{\sim}H^{p+1}((\mathbf f),M)
\qquad(p\geq1\text{에 대하여})
\end{equation}
과 $\mathcal F$에 관해 함자적인 완전열
\begin{equation}
\label{III.2.1.4.2-fr}
\tag{2.1.4.2}
0\longrightarrow H^0((\mathbf f),M)\longrightarrow M
\longrightarrow H^0(U,\mathcal F(*))
\longrightarrow H^1((\mathbf f),M)\longrightarrow0
\end{equation}
이 존재하며, 모든 준동형은 차수 $0$이다.
\end{corollary}

\begin{proof}
실제로 (1.4.1)을 (\hyperref[III.2.1.3-fr]{2.1.3})의 결과에 적용하면
충분하다.
\end{proof}

(\hyperref[III.2.1.3-fr]{2.1.3})과 유사한 ``국소'' 명제는 다음과 같다.

\begin{proposition}[2.1.5]
\label{III.2.1.5-fr}
$S$를 양의 차수를 갖는 등급환, $f_i$ ($1\leq i\leq r$)를 차수
$d_i$인 $S_+$의 동차 원소, $M$을 등급 $S$-가군이라 하자.
$X=\operatorname{Proj}(S)$를 $S$의 동차 소 스펙트럼이라 하고
$U_i=D_+(f_i)$, $U=\bigcup_iU_i$로 두며,
\[
H^\bullet(U,\widetilde M(*))
=\bigoplus_{n\in\mathbb Z}H^\bullet(U,(M(n))^\sim)
\]
로 두자. 그러면 등급 $S$-가군 구조에 관해 차수 $0$이고 $M$에 관해
함자적인 표준 동형
\begin{equation}
\label{III.2.1.5.1-fr}
\tag{2.1.5.1}
H^p(U,\widetilde M(*))\xrightarrow{\sim}H^{p+1}((\mathbf f),M)
\qquad\text{모든 }p\geq1\text{에 대하여}
\end{equation}
들과 $M$에 관해 함자적인 완전열
\begin{equation}
\label{III.2.1.5.2-fr}
\tag{2.1.5.2}
0\longrightarrow H^0((\mathbf f),M)\longrightarrow M
\longrightarrow H^0(U,\widetilde M(*))
\longrightarrow H^1((\mathbf f),M)\longrightarrow0
\end{equation}
이 존재하며, 모든 준동형은 차수 $0$이다.
\end{proposition}

\begin{proof}
실제로 정의에 의해 (II, 2.5.2)
\[
\Gamma(U_{i_0i_1\ldots i_p},(M(n))^\sim)
=(M_{f_{i_0}\cdots f_{i_p}})_n
\]
이므로
$\Gamma(U_{i_0i_1\ldots i_p},\widetilde M(*))
=M_{f_{i_0}\cdots f_{i_p}}$이다. 그러면 $X$가 스킴임을 고려할 때,
나머지 논증은 (\hyperref[III.2.1.4-fr]{2.1.4})의 증명과 같다.
\end{proof}

\begin{env}[2.1.6]
\label{III.2.1.6-fr}
\emph{비고.} ---
(i) (\hyperref[III.2.1.5-fr]{2.1.5})의 조건 아래에서
(0\textsubscript{I}, 1.3.2)에 의해 함자
$\Gamma(U_{i_0i_1\ldots i_p},\widetilde M(*))$는 $M$에 관해 완전하다.
따라서 (1.2.5)에서와 같은 논증으로, 등급 $S$-가군들의 완전열
$0\to M'\to M\to M''\to0$(준동형들은 차수 $0$)이 주어지면 모든
$p\geq0$에 대하여 가환 도식
\begin{equation}
\label{III.2.1.6.1-fr}
\tag{2.1.6.1}
\xymatrix{
H^p(U,\widetilde {M''}(*))\ar[r]^-\partial\ar[d] &
H^{p+1}(U,\widetilde {M'}(*))\ar[d]\\
H^{p+1}((\mathbf f),M'')\ar[r]^-\partial &
H^{p+2}((\mathbf f),M')
}
\end{equation}
을 얻는다.

(ii) 명제 (\hyperref[III.2.1.5-fr]{2.1.5})는 특히 $A$가 뇌터이고
$S$가 차수 $1$의 유한 개 원소로 생성되는 $A$-대수일 때 중요하다.
\oldpage[III]{98}
실제로 이 경우 모든 준연접 $\mathcal O_X$-가군은 $\widetilde M$의
꼴이다 (II, 2.7.7).
\end{env}

\begin{env}[2.1.7]
\label{III.2.1.7-fr}
$A$가 임의의 환이고 $T_i$들이 부정원일 때,
$S=A[T_0,\ldots,T_r]$, $M=S$, $f_i=T_i$인 경우에
(\hyperref[III.2.1.5-fr]{2.1.5})를 적용하겠다. 따라서 본질적으로
$\mathbf T=(T_i)_{0\leq i\leq r}$일 때
$H^\bullet((\mathbf T),S)$를 계산하는 문제로 귀착된다.
\end{env}

\begin{lemma}[2.1.8]
\label{III.2.1.8-fr}
$S=A[T_0,\ldots,T_r]$이고 $\mathbf T=(T_i)_{0\leq i\leq r}$이면
\begin{align}
\label{III.2.1.8.1-fr}
\tag{2.1.8.1}
H^i(\mathbf T^n,S)&=0 &&\text{$i\neq r+1$이면},\\
\label{III.2.1.8.2-fr}
\tag{2.1.8.2}
H^{r+1}(\mathbf T^n,S)&=S/(\mathbf T^n)
\end{align}
이다. 따라서 $A$-가군 $H^{r+1}(\mathbf T^n,S)$는 다음 단항식들의
$(\mathbf T^n)$을 법으로 한 유들로 이루어진 $A$ 위의 기저를 갖는다.
\[
T_{\mathbf p}=T_0^{p_0}\cdots T_r^{p_r},
\qquad \mathbf p=(p_0,\ldots,p_r),\quad 0\leq p_i<n
\quad\text{모든 }i\text{에 대하여}.
\]
\end{lemma}

\begin{proof}
가정들이 자명하게 성립하는 (1.1.3.5)와 명제 (1.1.4)의 즉각적인
귀결이다.
\end{proof}

\begin{env}[2.1.9]
\label{III.2.1.9-fr}
$n$에 관해 귀납극한을 취하자. 관계
(\hyperref[III.2.1.8.1-fr]{2.1.8.1})로부터 $i\neq r+1$이면
$H^i((\mathbf T),S)=0$임을 얻는다. $i=r+1$일 때 귀납계는
$S/(\mathbf T^n)$들로 이루어지고, 준동형
\[
\varphi_{nm}:S/(\mathbf T^n)\longrightarrow S/(\mathbf T^m)
\qquad(0\leq n\leq m)
\]
은 $(T_0\cdots T_r)^{m-n}$을 곱하는 사상이다.
$n\geq\sup(p_i)_{0\leq i\leq r}$일 때,
$\xi_{\mathbf p}^{(n)}=\xi_{p_0,\ldots,p_r}^{(n)}$로
$T_0^{n-p_0}\cdots T_r^{n-p_r}\pmod{(\mathbf T^n)}$의 유를 나타내자.
그러면 $\varphi_{nm}(\xi_{\mathbf p}^{(n)})
=\xi_{\mathbf p}^{(m)}$이고, 따라서 이 원소들은 귀납극한
$H^{r+1}((\mathbf T),S)$에서 같은 표준 상
$\xi_{\mathbf p}=\xi_{p_0,\ldots,p_r}$를 갖는다.
(\hyperref[III.2.1.2-fr]{2.1.2})에서 준 차수의 정의에 따라
$\xi_{\mathbf p}$의 차수는
$-|\mathbf p|=-(p_0+p_1+\cdots+p_r)$이다. $0<p_i\leq n$이고
$0\leq i\leq r$인 $\xi_{\mathbf p}^{(n)}$들이
$S/(\mathbf T^n)$의 기저를 이룸은 명백하다. 따라서
(\hyperref[III.2.1.8-fr]{2.1.8})로부터 즉시 다음을 얻는다.
\end{env}

\begin{corollary}[2.1.10]
\label{III.2.1.10-fr}
(\hyperref[III.2.1.8-fr]{2.1.8})의 표기 아래에서
\begin{equation}
\label{III.2.1.10.1-fr}
\tag{2.1.10.1}
H^i((\mathbf T),S)=0\qquad\text{$i\neq r+1$이면}
\end{equation}
이고, $H^{r+1}((\mathbf T),S)$는 $0\leq i\leq r$에 대하여
$p_i>0$인 원소 $\xi_{p_0,\ldots,p_r}$들로 이루어진 기저를 갖는 자유
$A$-가군이다.
\end{corollary}

\begin{env}[2.1.11]
\label{III.2.1.11-fr}
\emph{비고.} --- $N$을 임의의 $A$-가군이라 하고 $M=S\otimes_A N$이라
하자. (\hyperref[III.2.1.8-fr]{2.1.8})의 논증으로부터 더 일반적으로
\begin{align}
\label{III.2.1.11.1-fr}
\tag{2.1.11.1}
H^i(\mathbf T^n,M)&=0 &&\text{$i\neq r+1$이면},\\
\label{III.2.1.11.2-fr}
\tag{2.1.11.2}
H^{r+1}(\mathbf T^n,M)&=(S/(\mathbf T^n))\otimes_A N
\end{align}
을 얻는다. 실제로 마지막 공식은 (1.1.3.5)에서 직접 따르고, 한편
\[
M/(T_0^nM+\cdots+T_{i-1}^nM)
=\bigl(S/(T_0^nS+\cdots+T_{i-1}^nS)\bigr)\otimes_A N
\]
임이 명백하다. 아이디얼 $T_0^nS+\cdots+T_{i-1}^nS$는 $A$-가군
$S$의 직합인자이므로 (1.1.4)를 $M$에 적용할 수 있고, 이로써
(\hyperref[III.2.1.11.1-fr]{2.1.11.1})을 얻는다.
\end{env}

(\hyperref[III.2.1.10-fr]{2.1.10})과
(\hyperref[III.2.1.5-fr]{2.1.5})를 결합하면 다음을 얻는다.

\begin{proposition}[2.1.12]
\label{III.2.1.12-fr}
$A$를 환, $r$을 양의 정수, $X=\mathbf P_A^r$라 하자 (II, 4.1.1).
그러면 다음이 성립한다.
\begin{enumerate}
\item[{\rm(i)}] $i\neq0,r$이면 $H^i(X,\mathcal O_X(*))=0$이다.
\item[{\rm(ii)}] 표준 준동형
$\alpha:S\to H^0(X,\mathcal O_X(*))$ (II, 2.6.2)는 전단사이다.
\oldpage[III]{99}
\item[{\rm(iii)}] $H^r(X,\mathcal O_X(*))$는 $0\leq i\leq r$에 대해
$p_i>0$인 원소 $\xi_{p_0,\ldots,p_r}$들로 이루어진 기저를 갖는 자유
$A$-가군이다. 여기서 $\xi_{p_0,\ldots,p_r}$의 차수는
$-|\mathbf p|=-(p_0+\cdots+p_r)$이고,
$T_i\xi_{p_0,\ldots,p_r}=\xi_{p_0,\ldots,p_i-1,\ldots,p_r}$이다.
\end{enumerate}
\end{proposition}

\begin{proof}
실제로 $M=S=A[T_0,\ldots,T_r]$에 적용한 완전열
(\hyperref[III.2.1.5.2-fr]{2.1.5.2})에서
(\hyperref[III.2.1.10.1-fr]{2.1.10.1})에 의해
$H^0((\mathbf T),S)=0$이고 $H^1((\mathbf T),S)=0$임을 유의하자.
또한 $X$는 $D_+(T_i)$들의 합집합이므로 (II, 2.3.14), 명제
(\hyperref[III.2.1.5-fr]{2.1.5})를 $U=X$에 적용할 수 있다. 이제 완전열
(\hyperref[III.2.1.5.2-fr]{2.1.5.2})의 사상
$S\to H^0(X,\mathcal O_X(*))$을 표준 사상 $\alpha$와 식별하면 된다.
이는 $H^0(U,\mathcal O_X(*))$와
$H^0(\mathfrak U,\mathcal O_X(*))$의 표준 식별에서 따른다.
\end{proof}

\begin{corollary}[2.1.13]
\label{III.2.1.13-fr}
$H^i(X,\mathcal O_X(n))\neq0$일 수 있는 $(i,n)$의 값은 다음뿐이다.
$i=0$, $n\geq0$인 경우와 $i=r$, $n\leq-(r+1)$인 경우.
\end{corollary}

$A\neq0$이면 (\hyperref[III.2.1.13-fr]{2.1.13})에 열거한 쌍들에 대해
실제로 $H^i(X,\mathcal O_X(n))\neq0$임을 유의하자. 이는 모든
$n\geq0$에 대해 $S_n\neq0$이므로
(\hyperref[III.2.1.12-fr]{2.1.12})에서 따른다.

이 장에서의 응용에는 주로 다음의 덜 정밀한 결과를 사용하겠다.

\begin{corollary}[2.1.14]
\label{III.2.1.14-fr}
$A$-가군 $H^i(X,\mathcal O_X(n))$은 유한 생성 자유 가군이다. $i>0$이면
$n>0$에 대하여 0이다.
\end{corollary}

\begin{proposition}[2.1.15]
\label{III.2.1.15-fr}
$Y$를 준스킴, $\mathcal E$를 계수 $r+1$인 국소 자유
$\mathcal O_Y$-가군, $X=\mathbf P(\mathcal E)$를 $\mathcal E$로 정의되는
사영 다발, $f:X\to Y$를 구조 사상이라 하자.
$R^if_*(\mathcal O_X(n))\neq0$일 수 있는 $i$, $n$의 값은 $i=0$,
$n\geq0$인 경우와 $i=r$, $n\leq-(r+1)$인 경우뿐이다. 또한 표준
준동형 (II, 3.3.2)
\[
\alpha:S_{\mathcal O_Y}(\mathcal E)
\longrightarrow\Gamma_*(\mathcal O_X)
=R^0f_*(\mathcal O_X(*))
=\bigoplus_{n\in\mathbb Z}f_*(\mathcal O_X(n))
\]
은 동형사상이다.
\end{proposition}

\begin{proof}
문제는 $Y$ 위에서 국소적이므로, $Y$가 환 $A$의 아핀이고
$\mathcal E=\widetilde E$, $E=A^{r+1}$이라 가정할 수 있다. 그러면
(1.4.11)을 고려하여 곧바로 (\hyperref[III.2.1.12-fr]{2.1.12})로
귀착된다.
\end{proof}

\begin{env}[2.1.16]
\label{III.2.1.16-fr}
\emph{비고.} --- 나중에 (\hyperref[III.2.1.15-fr]{2.1.15})의 결과를
다음 명제들을 증명함으로써 보완하겠다. 다음과 같이 놓자.
\[
\boldsymbol\omega=f^*\left(\bigwedge^{r+1}\mathcal E\right)(-r-1).
\]
이는 가역 $\mathcal O_X$-가군이다. 그러면 다음이 성립한다.
\begin{enumerate}
\item[{\rm(i)}] 표준 동형사상
\begin{equation}
\label{III.2.1.16.1-fr}
\tag{2.1.16.1}
\rho:R^rf_*(\boldsymbol\omega)\xleftarrow{\sim}\mathcal O_Y
\end{equation}
이 존재한다.
\item[{\rm(ii)}] 컵곱 (0, 12.2.2)에 의한 짝지음
\begin{equation}
\label{III.2.1.16.2-fr}
\tag{2.1.16.2}
R^rf_*(\mathcal O_X(n))\times R^0f_*(\boldsymbol\omega(-n))
\longrightarrow R^rf_*(\boldsymbol\omega)
\end{equation}
을 동형사상 $\rho^{-1}$과 합성하면,
$R^rf_*(\mathcal O_X(n))$에서 국소 자유 $\mathcal O_Y$-가군
\[
R^0f_*(\boldsymbol\omega(-n))
=\left(\bigwedge^{r+1}\mathcal E\right)
\otimes_{\mathcal O_Y}\bigl(S_{\mathcal O_Y}(\mathcal E)\bigr)_{-n}
\]
의 \emph{쌍대}로 가는 \emph{동형사상}이 정의된다.
\end{enumerate}
\end{env}

\subsection{사영 사상의 기본정리}
\label{subsection:III.2.2-fr}
\oldpage[III]{100}

\begin{theorem}[2.2.1]
\label{III.2.2.1-fr}
(세르). $Y$를 뇌터 준스킴, $f:X\to Y$를 고유 사상,
$\mathcal L$을 $f$에 대해 풍부한 가역 $\mathcal O_X$-가군이라 하자.
모든 $\mathcal O_X$-가군 $\mathcal F$에 대하여
\[
\mathcal F(n)=\mathcal F\otimes_{\mathcal O_X}\mathcal L^{\otimes n}
\qquad\text{모든 }n\in\mathbb Z\text{에 대하여}
\]
로 놓는다. 그러면 모든 연접 $\mathcal O_X$-가군 $\mathcal F$에 대해
다음이 성립한다.
\begin{enumerate}
\item[{\rm(i)}] $R^qf_*(\mathcal F)$들은 연접 $\mathcal O_Y$-가군이다.
\item[{\rm(ii)}] 어떤 정수 $N$이 존재하여 $n\geq N$이면
\[
R^qf_*(\mathcal F(n))=0\qquad\text{모든 }q>0\text{에 대하여}
\]
이다.
\item[{\rm(iii)}] 어떤 정수 $N$이 존재하여 $n\geq N$이면 표준 준동형
\[
f^*(f_*(\mathcal F(n)))\longrightarrow\mathcal F(n)
\]
이 전사이다.
\end{enumerate}
\end{theorem}

\begin{proof}
먼저 정리에서 $\mathcal L$을 $\mathcal L^{\otimes d}$ $(d>0)$로
바꾸었을 때 정리가 참이면 원래 형태로도 참임을 유의하자. 실제로
\[
\mathcal F(n)
=(\mathcal F\otimes\mathcal L^{\otimes r})
\otimes\mathcal L^{\otimes hd}
\]
로 쓸 수 있는데, 여기서 $h>0$, $0\leq r<d$이다. 가정에 따라 각
$r$에 대하여 어떤 정수 $N_r$이 존재하여 $h\geq N_r$이면
$\mathcal O_X$-가군 $\mathcal F\otimes\mathcal L^{\otimes r}$에 대해
(ii), (iii)이 성립한다. $dN_r$들 가운데 가장 큰 것을 $N$으로 택하면
$n\geq N$에 대해 (ii), (iii)이 성립한다. 따라서 $\mathcal L$이
$f$에 대해 매우 풍부하다고 가정할 수 있다 (II, 4.6.11). 그러므로
어떤 우세한 $Y$-열린 몰입 $i:X\to P$가 존재한다. 여기서
$P=\operatorname{Proj}(\mathcal S)$이고, $\mathcal S$는 양의 차수들을
갖는 등급 준연접 $\mathcal O_Y$-대수이며, $\mathcal S_1$은 유한
생성이고 $\mathcal S$를 생성한다. 또한 $\mathcal L$은
$i^*(\mathcal O_P(1))$과 동형이다 (II, 4.4.7). 그런데 $f$가 고유이므로
$i$도 고유이고 (II, 5.4.4), 따라서 $i$는 동형사상
$X\xrightarrow{\sim}P$이다. 그러므로
$X=\operatorname{Proj}(\mathcal S)$, $\mathcal L=\mathcal O_X(1)$인
경우로 한정할 수 있다. 이제 정리 (2.2.1)은 다음 명제의 귀결이다.
\end{proof}

\begin{proposition}[2.2.2]
\label{III.2.2.2-fr}
$A$를 뇌터 환, $S$를 양의 차수들을 갖는 등급 $A$-대수라 하자.
$S_1$은 $r+1$개의 생성원계를 갖는 $A$-가군이고 대수 $S$를 생성한다고
하자. $X=\operatorname{Proj}(S)$로 놓는다. 모든 연접
$\mathcal O_X$-가군 $\mathcal F$에 대해 다음이 성립한다.
\begin{enumerate}
\item[{\rm(i)}] $A$-가군 $H^q(X,\mathcal F)$들은 유한 생성이다.
\item[{\rm(ii)}] $q>r$이면 $H^q(X,\mathcal F)=0$이다.
\item[{\rm(iii)}] 어떤 정수 $N$이 존재하여 $n\geq N$이면 모든
$q>0$에 대해 $H^q(X,\mathcal F(n))=0$이다.
\item[{\rm(iv)}] 어떤 정수 $N$이 존재하여 $n\geq N$이면
$\mathcal F(n)$은 $X$ 위의 단면들로 생성된다.
\end{enumerate}
\end{proposition}

\begin{proof}
먼저 (2.2.2)가 (2.2.1)을 함의함을 보이자. 위에서 고려한 특수한 경우
$X=\operatorname{Proj}(\mathcal S)$로 귀착된 (2.2.1)에서 $Y$는
준콤팩트이므로, 각 환이 뇌터인 유한 개의 아핀 열린집합들로 덮인다.
각 열린집합 $U_\alpha$ 위에서 $\mathcal S_1$의 제한은
$U_\alpha$ 위의 유한 개 단면들로 생성된다. (2.2.2)가 증명되었다고
가정하면, (1.4.11)과 (II, 3.4.7)을 고려하여 (2.2.1)의 (ii), (iii)에
나타나는 $N$으로 $U_\alpha$들에 대응하는 정수들 가운데 가장 큰 것을
택하면 충분하다.

(2.2.2)를 증명하기 위해 $X$가 $P=\mathbf P_A^r$의 닫힌 부분스킴과
동일시됨을 유의하자 (II, 3.6.2). 또한 $j:X\to P$가 표준 포함이면
$j_*(\mathcal F)$는 연접 $\mathcal O_P$-가군이고
$j_*(\mathcal F(n))=(j_*(\mathcal F))(n)$이다 (II, 3.4.5 및 3.5.2).
(G, II, 정리 4.9.1의 따름정리)를 고려하면, 따라서
\[
X=\mathbf P_A^r\quad\text{이고}\quad S=A[T_0,\ldots,T_r]
\]
인 특수한 경우에 (2.2.2)를 증명하는 문제로 귀착된다.

\oldpage[III]{101}
$X$는 $r+1$개의 아핀 열린집합 $D_+(T_i)$로 덮이므로 (ii)는
(1.4.12)에서 따른다. 한편 (iv)는 이미 증명되었다 (II, 2.7.9).

(i)와 (iii)을 동시에 증명하겠다. 이 명제들은
$\mathcal F=\mathcal O_X(m)$일 때 참이다
(\hyperref[III.2.1.13-fr]{2.1.13}). 따라서 $\mathcal F$가
$\mathcal O_X(m_j)$ 꼴의 유한 개 $\mathcal O_X$-가군들의 직합일 때도
참이다. 한편 (ii)에 의해 $q>r$이면 (i), (iii)은 자명하게 참이다.
$q$에 대한 \emph{하강 귀납법}을 쓰겠다. $\mathcal F$는 유한 개의
층 $\mathcal O_X(m_j)$의 직합 $\mathcal E$의 몫과 동형임을 안다
(II, 2.7.10). 다시 말해 완전열
\[
0\longrightarrow\mathcal R\longrightarrow\mathcal E
\longrightarrow\mathcal F\longrightarrow0
\]
이 존재한다. 여기서 $\mathcal R$은 연접이고 (0\textsubscript{I}, 5.3.3),
$\mathcal E$는 (i), (iii)을 만족한다. $\mathcal F(n)$은
$\mathcal F$에 관해 완전한 함자이므로, 모든 $n\in\mathbb Z$에 대해
완전열
\[
0\longrightarrow\mathcal R(n)\longrightarrow\mathcal E(n)
\longrightarrow\mathcal F(n)\longrightarrow0
\]
도 존재한다. 이로부터 코호몰로지 완전열
\[
H^{q-1}(X,\mathcal E(n))\longrightarrow
H^{q-1}(X,\mathcal F(n))\longrightarrow
H^q(X,\mathcal R(n))
\]
을 얻는다. $\mathcal E(n)$은 $\mathcal O_X(n+m_j)$들의 직합이므로
(II, 2.5.14), $H^{q-1}(X,\mathcal E(n))$은 유한 생성이다. 귀납가정에
의해 $H^q(X,\mathcal R(n))$도 유한 생성이다. $A$가 뇌터이므로
모든 $n\in\mathbb Z$에 대해 $H^{q-1}(X,\mathcal F(n))$이 유한
생성이고, 특히 $n=0$일 때 그러함을 얻는다. 한편 귀납가정에 의해
어떤 정수 $N$이 존재하여 $n\geq N$이면
$H^q(X,\mathcal R(n))=0$이다. 또한 $\mathcal E$가 (iii)을 만족하므로,
$n\geq N$이면 $H^{q-1}(X,\mathcal E(n))=0$이 되도록 $N$을 택할 수도
있다. 따라서 $n\geq N$이면
$H^{q-1}(X,\mathcal F(n))=0$이고, 이로써 증명이 끝난다.
\end{proof}

\begin{corollary}[2.2.3]
\label{III.2.2.3-fr}
(\hyperref[III.2.2.1-fr]{2.2.1})의 가정 아래에서
$\mathcal F\to\mathcal G\to\mathcal H$를 연접
$\mathcal O_X$-가군들의 완전열이라 하자. 그러면 어떤 정수 $N$이
존재하여 $n\geq N$이면 열
\[
f_*(\mathcal F(n))\longrightarrow f_*(\mathcal G(n))
\longrightarrow f_*(\mathcal H(n))
\]
이 완전하다.
\end{corollary}

\begin{proof}
$\mathcal F\to\mathcal G$의 핵, 상, 여핵을 각각
$\mathcal F'$, $\mathcal G'$, $\mathcal G''$라 하자.
$\mathcal G'$은 $\mathcal G\to\mathcal H$의 핵이고
$\mathcal G''$는 그 상이다. 이 준동형의 여핵을 $\mathcal H''$라 하자.
이 모든 $\mathcal O_X$-가군은 연접이다 (0\textsubscript{I}, 5.3.4).
$\mathcal F(n)$은 $\mathcal F$에 관해 완전한 함자이므로, 충분히 큰
$n$에 대하여 다음 열들이 각각 완전함을 증명하면 충분하다.
\begin{align*}
0&\longrightarrow f_*(\mathcal F'(n))
\longrightarrow f_*(\mathcal F(n))
\longrightarrow f_*(\mathcal G'(n))\longrightarrow0,\\
0&\longrightarrow f_*(\mathcal G'(n))
\longrightarrow f_*(\mathcal G(n))
\longrightarrow f_*(\mathcal G''(n))\longrightarrow0,\\
0&\longrightarrow f_*(\mathcal G''(n))
\longrightarrow f_*(\mathcal H(n))
\longrightarrow f_*(\mathcal H''(n))\longrightarrow0.
\end{align*}
따라서 $0\to\mathcal F\to\mathcal G\to\mathcal H\to0$이
완전하다고 가정할 수 있다. 이때 코호몰로지 완전열
\[
0\longrightarrow f_*(\mathcal F(n))
\longrightarrow f_*(\mathcal G(n))
\longrightarrow f_*(\mathcal H(n))
\longrightarrow R^1f_*(\mathcal F(n))\longrightarrow\cdots
\]
을 얻고, 결론은 (\hyperref[III.2.2.1-fr]{2.2.1}, (ii))에서 따른다.
\end{proof}

\begin{corollary}[2.2.4]
\label{III.2.2.4-fr}
$Y$를 뇌터 준스킴, $f:X\to Y$를 유한형 사상,
$\mathcal L$을 $f$에 대해 풍부한 가역 $\mathcal O_X$-가군이라 하자.
모든 $\mathcal O_X$-가군 $\mathcal F$에 대해
$\mathcal F(n)=\mathcal F\otimes_{\mathcal O_X}\mathcal L^{\otimes n}$
$(n\in\mathbb Z)$으로 놓는다. $\mathcal F\to\mathcal G\to\mathcal H$가
연접 $\mathcal O_X$-가군들의 완전열이고,
$\operatorname{Im}(\mathcal F\to\mathcal G)$의 받침이 $Y$ 위에서
고유라고 하자 (II, 5.4.10). 그러면 어떤 정수 $N$이 존재하여
$n\geq N$이면 열
\[
f_*(\mathcal F(n))\longrightarrow f_*(\mathcal G(n))
\longrightarrow f_*(\mathcal H(n))
\]
이 완전하다.
\end{corollary}

\begin{proof}
(\hyperref[III.2.2.1-fr]{2.2.1})의 증명 첫머리와 같은 논증에 따르면,
$\mathcal L^{\otimes d}$ $(d>0)$에 대해 따름정리가 참이면
$\mathcal L$에 대해서도 참이다. 따라서 $\mathcal L$이 $f$에 대해
매우 풍부한 경우로 한정할 수 있다 (II, 4.6.11). 그러므로 $X$를
$Y$-스킴 $Z=\operatorname{Proj}(\mathcal S)$의 열린집합과 동일시할 수
있다. 여기서 $\mathcal S$는 양의 차수들을 갖는 등급 준연접
$\mathcal O_Y$-대수이고, $\mathcal S_1$은 유한 생성이며
$\mathcal S$를 생성한다. 또한 $i:X\to Z$가 표준 몰입이면
$\mathcal L=i^*(\mathcal O_Z(1))$이다 (II, 4.4.7).
$\mathcal G_1=\operatorname{Im}(\mathcal F\to\mathcal G)
=\operatorname{Ker}(\mathcal G\to\mathcal H)$로 놓으면 이는 연접이다.
완전열 $0\to\mathcal G_1(n)\to\mathcal G(n)\to\mathcal H(n)$이 있고
함자 $f_*$가 왼쪽 완전하므로, 모든 $n$에 대해 열
$0\to f_*(\mathcal G_1(n))\to f_*(\mathcal G(n))
\to f_*(\mathcal H(n))$은 완전하다. 따라서 충분히 큰 $n$에 대하여
$f_*(\mathcal F(n))\to f_*(\mathcal G_1(n))\to0$이 완전함을 보이면
충분하다. 다시 말해 $\mathcal H=0$이고 $\mathcal G$의 받침이
$Y$ 위에서 고유인 경우로 한정할 수 있다. 이 받침은 따라서
$Z$에서 \emph{닫혀 있다} (II, 5.4.10). 그러므로
$\mathcal G'=i_*(\mathcal G)$는 $\mathcal G'|X=\mathcal G$인 연접
$\mathcal O_Z$-가군이다. $\mathcal F'|X=\mathcal F$인 연접
$\mathcal O_Z$-가군 $\mathcal F'$가 존재함을 안다 (I, 9.4.3).
또한 $X$가 $Z$에서 열려 있고
$\operatorname{Supp}(\mathcal G)\subset X$가 $Z$에서 닫혀 있으므로,
$\operatorname{Supp}(\mathcal G)$와 만나지 않는 $Z$의 모든 열린집합
$U$에서는 $u'|U=0$으로, 모든 열린집합 $U\subset X$에서는
$u'|U=u|U$로 놓으면, 주어진 전사 준동형
$u:\mathcal F\to\mathcal G$로 제한되는 전사 층 준동형
$u':\mathcal F'\to\mathcal G'$을 정의함이 곧바로 확인된다.

이제 (2.2.2)의 증명 첫머리에서와 같이 $Y$가 아핀인 경우로 한정할
수 있다. 따라서 충분히 큰 $n$에 대해 준동형
$\Gamma(u):\Gamma(X,\mathcal F(n))\to\Gamma(X,\mathcal G(n))$이
전사임을 보이면 된다. 그런데 가환 도식
\[
\xymatrix@C=35mm@R=12mm{
\Gamma(Z,\mathcal F'(n)) \ar[r]^{\Gamma(u')} \ar[d]_v &
\Gamma(Z,\mathcal G'(n)) \ar[d]^w \\
\Gamma(X,\mathcal F(n)) \ar[r]_{\Gamma(u)} & \Gamma(X,\mathcal G(n))
}
\]
이 있다. 여기서 $v$, $w$는 제한 준동형들이다. $\mathcal G'$의 정의에
의해 $w$는 \emph{전단사}이다. 한편 (2.2.3)에 의해 충분히 큰 $n$에
대해 $\Gamma(u')$은 \emph{전사}이므로, $\Gamma(u)$도 전사이다.
\end{proof}

\begin{env}[2.2.5]
\label{III.2.2.5-fr}
\emph{비고.} ---
(i) (\hyperref[III.2.2.1-fr]{2.2.1})의 명제 (i)는 $Y$가 단지
\emph{국소 뇌터}라고 가정해도 여전히 참이다. 실제로 이 성질은
$Y$ 위에서 명백히 국소적이다. 한편
(\hyperref[III.2.2.1-fr]{2.2.1})의 가정들은 모든 열린집합
$U\subset Y$에 대해 $f$를 $f^{-1}(U)$에 제한한 것이 사영 사상
$f^{-1}(U)\to U$이고 (II, 5.5.5, (iii)),
$\mathcal L|f^{-1}(U)$가 이 사상에 대해 풍부함을 함의한다
(II, 4.6.4).

(ii) 앞에서 보았듯이, (\hyperref[III.2.2.1-fr]{2.2.1})의 명제 (iii)은
$X$가 준콤팩트 스킴이거나 그 바탕공간이 뇌터인 준스킴이고
$f:X\to Y$가 준콤팩트 사상이라고만 가정해도 여전히 성립한다
(II, 4.6.8). 그러나 $Y$가 체 $K$의 \emph{스펙트럼}이고 $f$가
\emph{준사영}이라고 가정해도
(\hyperref[III.2.2.1-fr]{2.2.1})의 명제 (ii)는 더 이상 반드시
성립하지 않음에 유의해야 한다. 예를 들어
$X'=\operatorname{Spec}(K[T_0,\ldots,T_r])$라 하고 $X$를 $X'$의 아핀
열린집합 $D(T_i)$ $(0\leq i\leq r)$들의 합집합이라 하자.
몰입 $X\to X'$가 준콤팩트이므로 구조 사상 $f:X\to Y$는 준아핀이고
(II, 5.1.10), 따라서 $\mathcal O_X$는 $f$에 대해 매우 풍부하다
(II, 5.1.6). 그런데 환 $\Gamma(X,\mathcal O_X)$는
$0\leq i\leq r$에 대한 분수환
$(K[T_0,\ldots,T_r])_{T_i}$들의 교집합과 동일시되며
(I, 8.2.1.1), 이는 $K[T_0,\ldots,T_r]$이다. 따라서 공식 (1.4.3.1)과
(1.1.3.5)로부터 모든 $n$에 대해
$H^r(X,\mathcal O_X^{\otimes n})=H^r(X,\mathcal O_X)\neq0$임을 얻는다.

(iii) $X$, $Y$, $f$, $\mathcal L$에 관한 (2.2.4)의 가정 아래에서,
$\mathcal H$가 받침이 \emph{$Y$ 위에서 고유}인 연접
$\mathcal O_X$-가군이면 (2.2.4)와 유사한 논증에 의해 어떤 정수 $N$이
존재하여 $n\geq N$이면 모든 $i>0$에 대해
$\mathrm R^if_*(\mathcal H(n))=0$임은 즉시 확인된다. 코호몰로지
완전열로부터 다음 결론을 얻는다. $u:\mathcal F\to\mathcal G$가 연접
$\mathcal O_X$-가군들의 준동형이고 $\operatorname{Ker}(u)$와
$\operatorname{Coker}(u)$의 받침들이 \emph{$Y$ 위에서 고유}이면,
어떤 $N$이 존재하여 $n\geq N$일 때 대응하는 준동형
$\mathrm R^if_*(\mathcal F(n))\to\mathrm R^if_*(\mathcal G(n))$은
모든 $i>0$에 대해 \emph{전단사}이다.
\end{env}

\subsection{등급 대수 및 가군의 층에 대한 응용}
\label{subsection:III.2.3-fr}

\begin{theorem}[2.3.1]
\label{III.2.3.1-fr}
$Y$를 뇌터 준스킴, $\mathcal S$를 양의 차수들을 갖는 준연접 유한 생성
등급 $\mathcal O_Y$-대수, $X=\operatorname{Proj}(\mathcal S)$,
$q:X\to Y$를 구조 사상, $\mathcal M$을 조건 \textbf{(TF)}를 만족하는
준연접 등급 $\mathcal S$-가군이라 하자. 그러면 어떤 정수 $N$이
존재하여 $n\geq N$이면 표준 준동형 (II, 8.14.5.1)
\[
\alpha_n:\mathcal M_n\longrightarrow
q_*(\shProj_0(\mathcal M(n)))
=q_*((\shProj(\mathcal M))_n)
\]
\oldpage[III]{103}
이 전단사이다. 다시 말해 표준 준동형
\[
\alpha:\mathcal M\longrightarrow\bbGamma_*(\shProj(\mathcal M))
\]
은 \textbf{(TN)}-동형사상이다.
\end{theorem}

\begin{proof}
$\mathcal M$이 유한 생성 $\mathcal S$-가군인 경우로 한정할 수 있다
(II, 3.4.2). $Y$가 준콤팩트이므로, $\mathcal S^{(d)}$가 준연접
$\mathcal O_Y$-가군 $\mathcal S_d$로 생성되도록 하는 어떤 정수
$d>0$이 존재한다. $\mathcal S_d$는 유한 생성이고 (II, 3.1.10),
$Y$가 뇌터이므로 연접이다. 이제 $\mathcal M$은
$0\leq k<d$에 대한 $\mathcal M^{(d,k)}$들의 직합이며, 문제는
$Y$ 위에서 국소적이므로 (II, 2.1.6, (iii))에 따라 각
$\mathcal M^{(d,k)}$는 유한 생성 준연접 등급
$\mathcal S^{(d)}$-가군이다. 표준 준동형들
\[
\alpha:\mathcal M^{(d,k)}
\longrightarrow\bbGamma_*((\shProj(\mathcal M))^{(d,k)})
\]
이 각각 \textbf{(TN)}-동형사상임을 증명하면 충분함은 명백하다.
(II, 8.14.13), 특히 도식 (8.14.13.4)을 고려하면,
$\mathcal S$가 $\mathcal S_1$로 생성되고 $\mathcal S_1$이 연접
$\mathcal O_Y$-가군인 경우에 정리를 증명하는 문제로 귀착됨을 알 수
있다.

$Y$가 뇌터이므로, (\hyperref[III.2.2.2-fr]{2.2.2})의 증명 첫머리와
같은 논증으로 $Y=\operatorname{Spec}(A)$,
$\mathcal S=\widetilde S$, $\mathcal M=\widetilde M$인 경우로 한정할 수
있다. 여기서 $A$는 뇌터 환, $S_1$은 유한 생성 $A$-가군이고 $M$은
유한 생성 등급 $S$-가군이다. 이때 $M=S$인 경우에 정리를 증명하면
충분함을 보이자. 실제로 일반적인 경우에는 완전열
\[
L'\longrightarrow L\longrightarrow M\longrightarrow0
\]
이 존재하며, 여기서 $L$, $L'$은 $S(m)$ 꼴의 등급가군들의 직합이다.
$M=S$일 때 결과가 참이면 $M=S(m)$일 때도 참이고, 따라서 $L$과
$L'$에 대해서도 참이다. 이제 가환 도식
\[
\xymatrix{
\widetilde {L'_n}\ar[r]\ar[d]_{\alpha_n} &
\widetilde {L_n}\ar[r]\ar[d]^{\alpha_n} &
\widetilde {M_n}\ar[r]\ar[d]^{\alpha_n} & 0\\
q_*(\widetilde {L'}(n))\ar[r] &
q_*(\widetilde L(n))\ar[r] &
q_*(\widetilde M(n))\ar[r] & 0 .
}
\]
을 생각하자. 충분히 큰 $n$에 대해 둘째 줄은
(\hyperref[III.2.2.3-fr]{2.2.3})에 의해 완전하다. 첫째 줄도 그러하며
왼쪽의 두 수직 화살표가 동형사상이므로 셋째 화살표도 동형사상이다.

이제 $M=S$일 때 정리를 증명하기 위해 먼저
$S=A[T_0,\ldots,T_r]$($T_i$들은 부정원)라고 가정하자. 이 경우 우리
명제는 (\hyperref[III.2.1.11-fr]{2.1.11}, (ii))에 지나지 않는다.
일반적인 경우 $S$는 등급 아이디얼에 의한 환
$S'=A[T_0,\ldots,T_r]$의 몫과 동일시되며, 따라서 $X$는
$X'=\mathbf P_A^r$의 닫힌 부분스킴과 동일시된다 (II, 2.9.2).
$j$가 표준 포함 $X\to X'$이면, $j_*(\widetilde S(n))$는 $S$를 등급
$S'$-가군으로 보았을 때의 $\mathcal O_{X'}$-가군
$(\shProj(\widetilde S))(n)$에 지나지 않는다. 이는 (II, 2.8.7)에서
곧바로 따른다. $j_*(\widetilde S(n))$는 조건 \textbf{(TF)}를 만족하는
$\mathcal O_{X'}$-가군이므로, 앞에서 증명한 바에 따라 충분히 큰
$n$에 대해 표준 준동형
$\alpha_n:S_n\to\Gamma(X',j_*(\widetilde S(n)))$은 전단사이다.
그런데
$\Gamma(X',j_*(\widetilde S(n)))=\Gamma(X,\widetilde S(n))$이므로
증명이 끝난다.
\end{proof}

\oldpage[III]{104}
\begin{corollary}[2.3.2]
\label{III.2.3.2-fr}
(\hyperref[III.2.3.1-fr]{2.3.1})의 가정 아래에서
\[
\mathcal S_X=\bigoplus_{n\in\mathbb Z}\mathcal O_X(n)
\]
이라 하고, $\mathcal F$를 유한 생성 준연접 등급
$\mathcal S_X$-가군이라 하자. 그러면 $\bbGamma_*(\mathcal F)$는 조건
\textbf{(TF)}를 만족한다.
\end{corollary}

\begin{proof}
(\hyperref[III.2.3.1-fr]{2.3.1})의 증명에서
$X$가 $\operatorname{Proj}(\mathcal S^{(d)})$와 동형임을 보았다
(II, 3.1.8). 따라서 $X$는 $Y$ 위에서 유한형이다 (II, 3.4.1).
그러면 (II, 8.14.9)에 의해 $\mathcal F$는
$\shProj(\mathcal M)$ 꼴의 등급 $\mathcal S_X$-가군과 동형이다.
여기서 $\mathcal M$은 유한 생성 준연접 등급 $\mathcal S$-가군이다.
(\hyperref[III.2.3.1-fr]{2.3.1})에 의해
$\bbGamma_*(\mathcal F)$는 $\mathcal M$과 \textbf{(TN)}-동형이고,
따라서 \textbf{(TF)}를 만족한다.
\end{proof}

\begin{env}[2.3.3]
\label{III.2.3.3-fr}
\emph{보주.} --- $Y$를 뇌터 준스킴, $\mathcal S$를
(\hyperref[III.2.3.1-fr]{2.3.1})의 조건들을 만족하는 등급
$\mathcal O_Y$-대수, $X=\operatorname{Proj}(\mathcal S)$라 하자.
조건 \textbf{(TF)}를 만족하는 준연접 등급 $\mathcal S$-가군들의
아벨 범주를 $K_{\mathcal S}$라 하고, 그 가운데 조건 \textbf{(TN)}을
만족하는 $\mathcal S$-가군들로 이루어진 부분범주를
$K'_{\mathcal S}$라 하자. 마지막으로, 유한 생성 준연접 등급
$\mathcal S_X$-가군 $\mathcal F$들의 범주를 $K_X$라 하자. 이는
$\mathcal S_X$가 주기적이므로 (II, 8.14.4 및 8.14.12),
$\mathcal F_i$들이 연접 $\mathcal O_X$-가군이라고 말하는 것과 같다.
그러면 함자
\[
\mathcal M\longmapsto\shProj(\mathcal M)
\quad\text{$K_{\mathcal S}$에서},
\qquad
\mathcal F\longmapsto\bbGamma_*(\mathcal F)
\quad\text{$K_X$에서}
\]
는 (II, 8.14.8 및 8.14.10)과
(\hyperref[III.2.3.2-fr]{2.3.2})에 의해 몫범주
$K_{\mathcal S}/K'_{\mathcal S}$ (T, I, 1.11)와 범주 $K_X$ 사이의
동치 (T, I, 1.2)를 정의한다. $\mathcal S$가 $\mathcal S_1$로
생성되면 $K_X$를 연접 $\mathcal O_X$-가군들의 범주로 바꿀 수 있다
(II, 8.14.12).
\end{env}

\begin{proposition}[2.3.4]
\label{III.2.3.4-fr}
$Y$를 뇌터 준스킴이라 하자.
\begin{enumerate}
\item[{\rm(i)}] $\mathcal S$를 양의 차수들을 갖는 준연접 유한 생성
등급 $\mathcal O_Y$-대수라 하자. $X=\operatorname{Proj}(\mathcal S)$,
\[
\mathcal S_X=\shProj(\mathcal S)
=\bigoplus_{n\in\mathbb Z}\mathcal O_X(n)
\]
으로 놓는다. 그러면 $\mathcal S_X$는 동차 성분
$(\mathcal S_X)_n=\mathcal O_X(n)$들이 연접 $\mathcal O_X$-가군인
주기적 등급 $\mathcal O_X$-대수이다 (II, 8.14.12). 또한 $d>0$이
$\mathcal S_X$의 주기이면 $(\mathcal S_X)_d=\mathcal O_X(d)$는
$Y$-풍부한 가역 $\mathcal O_X$-가군이다. 더 나아가 표준 준동형
\[
\alpha:\mathcal S\longrightarrow\bbGamma_*(\mathcal S_X)
\]
은 \textbf{(TN)}-동형사상이다.
\item[{\rm(ii)}] 역으로, $q:X\to Y$를 사영 사상이라 하고,
$\mathcal S'$을 동차 성분 $\mathcal S'_n$ $(n\in\mathbb Z)$들이 연접
$\mathcal O_X$-가군인 등급 $\mathcal O_X$-대수라 하자. 또한
$\mathcal S'_d$가 $q$에 대해 풍부한 가역 $\mathcal O_X$-가군이
되도록 하는 주기 $d>0$을 갖는다고 하자. 그러면
\[
\mathcal S=\bigoplus_{n\geq0}q_*(\mathcal S'_n)
\]
은 양의 차수들을 갖는 준연접 유한 생성 등급 $\mathcal O_Y$-대수이고,
$r^*(\shProj(\mathcal S))$가 등급 $\mathcal O_X$-대수로서
$\mathcal S'$과 동형이 되도록 하는 $Y$-동형사상
\[
r:X\xrightarrow{\sim}\operatorname{Proj}(\mathcal S)
\]
이 존재한다.
\end{enumerate}
\end{proposition}

\begin{proof}
(i) 모든 명제는 사실상 이미 증명되었고, 마지막 것은
(\hyperref[III.2.3.2-fr]{2.3.2})의 특수한 경우에 지나지 않는다.
$\mathcal S_X$가 주기적이라는 사실은 (II, 8.14.14)에서 보았고,
$\mathcal O_X(d)$가 가역이고 $Y$-풍부하도록 하는 주기 $d>0$이
존재한다는 사실은 (II, 4.6.18)에 지나지 않는다. 마지막으로
$0\leq k<d$에 대해 $(\mathcal S_X)^{(d,k)}$는 유한 생성
$(\mathcal S_X)^{(d)}$-가군이고 (II, 8.14.14), 따라서
(II, 2.1.6, (ii))에 의해 각 $(\mathcal S_X)_n$은 유한 생성 준연접
$\mathcal O_X$-가군이다. 문제는 국소적이고 $\mathcal O_X$가
연접이므로 $(\mathcal S_X)_n$들도 연접이다.

(ii) 주기 $d$를 그 배수로 바꾸어 $\mathcal L=\mathcal S'_d$가
$q$에 대해 \emph{매우 풍부한} 가역 $\mathcal O_X$-가군이라고 가정할
수 있다 (II, 4.6.11). 또한 가정에 의해
\[
\mathcal S'^{(d)}=\bigoplus_{n\in\mathbb Z}\mathcal L^{\otimes n}
\quad\text{이고, 따라서}\quad
\mathcal S^{(d)}=\bigoplus_{n\geq0}q_*(\mathcal L^{\otimes n}).
\]
(II, 3.1.8 및 3.2.9)에 의해
$X'=\operatorname{Proj}(\mathcal S)$에서
$X''=\operatorname{Proj}(\mathcal S^{(d)})$로 가는 $Y$-동형사상 $s$가
존재하며
\oldpage[III]{105}
\[
s^*(\mathcal O_{X''}(n))=\mathcal O_{X'}(nd)
\]
이다. 따라서 다음을 증명하면 $Y$-동형사상
$X\xrightarrow{\sim}X'$의 존재가 확립된다.

\begin{env}[2.3.4.1]
\label{III.2.3.4.1-fr}
$Y$를 뇌터 준스킴, $q:X\to Y$를 사영 사상,
$\mathcal L$을 $q$에 대해 매우 풍부한 가역 $\mathcal O_X$-가군이라
하자. 그러면
$\mathcal S=\bigoplus_{n\geq0}q_*(\mathcal L^{\otimes n})$은
$\mathcal S_n=\mathcal S_1^n$이 충분히 큰 $n$에 대해 성립하는 유한
생성 준연접 등급 $\mathcal O_Y$-대수이고, 어떤 $Y$-동형사상
$r:X\xrightarrow{\sim}P=\operatorname{Proj}(\mathcal S)$이 존재하여
$\mathcal L=r^*(\mathcal O_P(1))$이다.
\end{env}

$q$가 사영 사상이므로 (II, 5.4.4 및 4.4.7)에 의해 어떤 $Y$-동형사상
$r':X\xrightarrow{\sim}P'=\operatorname{Proj}(\mathcal T)$가 존재한다.
여기서 $\mathcal T$는 준연접 $\mathcal O_Y$-대수이고,
$\mathcal T_1$은 유한 생성이며 $\mathcal T$를 생성하고,
$\mathcal L=r'^*(\mathcal O_{P'}(1))$이다. $q':P'\to Y$가 구조
사상이면
\[
\mathcal S=\bigoplus_{n\geq0}q'_*(\mathcal O_{P'}(n))
\]
이다. (\hyperref[III.2.3.1-fr]{2.3.1})에 의해 충분히 큰 $n$에 대해
표준 준동형
$\alpha_n:\mathcal T_n\to\mathcal S_n=q'_*(\mathcal O_{P'}(n))$은
전단사이다. $\mathcal T_n=\mathcal T_1^n$이므로, 더구나 충분히 큰
$n$부터 $\mathcal S_n=\mathcal S_1^n$이다. 또한 등급
$\mathcal O_Y$-대수들의 표준 준동형
$\alpha:\mathcal T\to\mathcal S$은 \textbf{(TN)}-동형사상이다.
따라서
\phantomsection\label{III.2.3.5.1-fr}
\[
\Phi=\operatorname{Proj}(\alpha):
\operatorname{Proj}(\mathcal S)\longrightarrow
\operatorname{Proj}(\mathcal T)
\]
는 동형사상이다 (II, 3.6.1). 더 나아가 (II, 3.5.2)와, 스칼라 제한에
의해 $\mathcal S(n)$으로부터 얻은 등급 $\mathcal T$-가군들이
$\mathcal T(n)$과 \textbf{(TN)}-동형이라는 사실에 의해
\[
\Phi_*(\mathcal O_P(n))=\mathcal O_{P'}(n)
\]
이다 (II, 3.4.2). (\hyperref[III.2.3.4.1-fr]{2.3.4.1})의 증명을
끝내려면 $\mathcal S$가 유한 생성 $\mathcal O_Y$-대수임을 보이면
된다. 그런데 (\hyperref[III.2.2.1-fr]{2.2.1})에 의해
$\mathcal S_n=q'_*(\mathcal O_{P'}(n))$은 연접
$\mathcal O_Y$-가군이고, $n\geq n_0$일 때
$\mathcal S_n=\mathcal S_1^n$이면 $\mathcal S$는 연접인
$\bigoplus_{i\leq n_0}\mathcal S_i$로 생성된다. 이로부터 결론이
따른다 (I, 9.6.2).

(\hyperref[III.2.3.4-fr]{2.3.4})의 증명으로 돌아가 그 표기를 다시
사용하자. 모든 $n\in\mathbb Z$에 대해
$r''_*(\mathcal L^{\otimes n})=\mathcal O_{X''}(n)$이 되도록 하는
$Y$-동형사상 $r'':X\xrightarrow{\sim}X''$의 존재를 증명하였다.
$X''\to Y$의 구조 사상을 $q''$라 하자. 이제 $\mathcal S'$은 등급
$\mathcal S'^{(d)}$-가군 $\mathcal S'^{(d,k)}$들의 직합이다. 각
$\mathcal S'^{(d,k)}$는 $\mathcal S'$의 주기성과 $\mathcal S'_n$들이
유한 생성 $\mathcal O_X$-가군이라는 가정에 의해 유한 생성 준연접
등급 $\mathcal S'^{(d)}$-가군이다 (II, 8.14.12). 다음과 같이 놓자.
\[
\mathcal F^{(k)}=r''_*(\mathcal S'^{(d,k)}).
\]
그러면 $\mathcal F^{(k)}$들은 유한 생성 준연접 등급
$\mathcal S_{X''}$-가군이다. 따라서 (II, 8.14.8)에 의해 표준 준동형
\[
\beta:\shProj(\bbGamma_*(\mathcal F^{(k)}))
\longrightarrow\mathcal F^{(k)}
\]
은 $\mathcal S_{X''}$-가군의 동형사상이다. 한편
$q''_*((\mathcal F^{(k)})_n)=q_*((\mathcal S'^{(d,k)})_n)$이고,
$n\geq0$이면 마지막 $\mathcal O_Y$-가군은 정의에 의해
$(\mathcal S^{(d,k)})_n$과 같다. 다시 말해 표준 포함
$\mathcal S^{(d,k)}\to\bbGamma_*(\mathcal F^{(k)})$은
\textbf{(TN)}-동형사상이다. 그러므로
\[
\shProj(\mathcal S^{(d,k)})
=\shProj(\bbGamma_*(\mathcal F^{(k)})),
\]
따라서
$r''{}^*(\shProj(\mathcal S^{(d,k)}))\simeq\mathcal S'^{(d,k)}$이다.
마지막으로 표준 동형을 제외하면
\[
\shProj(\mathcal S^{(d,k)})
=s_*((\shProj(\mathcal S))^{(d,k)})
\]
임을 유의하면 (II, 8.14.13.1), $r^*(\shProj(\mathcal S))$와
$\mathcal S'$의 동형이 증명된다.

끝으로 (\hyperref[III.2.3.2-fr]{2.3.2})에 의해 각
$\bbGamma_*(\mathcal F^{(k)})$는 조건 \textbf{(TF)}를 만족하므로,
각 $\mathcal S^{(d,k)}$도 그러하다. 또한 $\mathcal S'_n$들이
연접이므로 (\hyperref[III.2.2.1-fr]{2.2.1})에 의해
$\mathcal S_n=q_*(\mathcal S'_n)$들은 연접이다. 이로부터 곧바로
$\mathcal S^{(d,k)}$들이 유한 생성 $\mathcal S^{(d)}$-가군임을
얻는다. (\hyperref[III.2.3.4.1-fr]{2.3.4.1})에 의해
$\mathcal S^{(d)}$가 유한 생성 $\mathcal O_Y$-대수이므로
$\mathcal S$도 그러하다.
\end{proof}

\oldpage[III]{106}
\begin{proposition}[2.3.5]
\label{III.2.3.5-fr}
$Y$를 뇌터 정역 준스킴, $X$를 정역 준스킴,
$f:X\to Y$를 사영 쌍유리 사상이라 하자. 그러면 어떤 연접 분수
아이디얼 $\mathcal J\subset\mathcal R(Y)$가 존재하여 (II, 8.1.2),
$X$는 $\mathcal J$를 따라 블로업하여 얻은 준스킴과 $Y$-동형이다
(II, 8.1.3). 또한 $Y$의 어떤 열린집합 $U$가 존재하여, $f$를
$f^{-1}(U)$에 제한한 것은 $f^{-1}(U)$에서 $U$로 가는 동형사상이고
(I, 6.5.5 참조), $\mathcal J|U$는 가역이다.
\end{proposition}

\begin{proof}
$f$에 대해 매우 풍부한 어떤 가역 $\mathcal O_X$-가군 $\mathcal L$이
존재하므로 (II, 4.4.2 및 5.3.2),
(\hyperref[III.2.3.4.1-fr]{2.3.4.1})을 적용할 수 있다. 따라서 $X$는
\[
\operatorname{Proj}(\mathcal S),\qquad
\mathcal S=\bigoplus_{n\geq0}f_*(\mathcal L^{\otimes n})
\]
과 동일시된다. 또한 $f_*(\mathcal L^{\otimes n})$들은 꼬임 없는
$\mathcal O_Y$-가군임을 안다 (I, 7.4.5). 따라서
$\mathcal O_Y$-가군 $\mathcal S$도 꼬임이 없고, 표준적으로
$\mathcal S\otimes_{\mathcal O_Y}\mathcal R(Y)$의 부분
$\mathcal O_Y$-가군과 동일시된다 (I, 7.4.1). 마지막 것은 그 제한을
어떤 공집합이 아닌 열린집합에서 알면 결정되는 \emph{단순한} 층이다
(I, 7.3.6). 예를 들어 $\mathcal L|f^{-1}(U')$가
$\mathcal O_X|f^{-1}(U')$와 동형이 되도록 하는 공집합이 아닌
$U'\subset U$에서 그러하다. 가정에 의해 이때
$f_*(\mathcal L^{\otimes n})|U'$들은 $\mathcal O_Y|U'$와 동형이므로,
$\mathcal S\otimes_{\mathcal O_Y}\mathcal R(Y)$는 부정원 $T$에 대한
$\mathcal R(Y)[T]$와 동형인 $\mathcal R(Y)$-가군이다. 또한
$\mathcal S$는 $f_*(\mathcal L)$의 표준 상이
$\mathcal S\otimes_{\mathcal O_Y}\mathcal R(Y)$ 안에서 생성하는
부분 $\mathcal O_Y$-대수와 \textbf{(TN)}-동형이다
(\hyperref[III.2.3.4.1-fr]{2.3.4.1}). 마지막 층을
$\mathcal R(Y)[T]$와 동일시하면 $f_*(\mathcal L)$의 상은
$\mathcal J T$와 동일시된다. 여기서 $\mathcal J$는
$\mathcal R(Y)$의 연접 부분 $\mathcal O_Y$-가군이고
(\hyperref[III.2.2.1-fr]{2.2.1}), 그 $U'$ 위의 제한은
$\mathcal O_Y|U'$와 동형이며, 따라서 $\mathcal J|U$가 가역이 되도록
한다. 이제 $\mathcal S$가 $\bigoplus_{n\geq0}\mathcal J^n$과
\textbf{(TN)}-동형임을 알 수 있고, 이로써 증명이 끝난다.
\end{proof}

\begin{corollary}[2.3.6]
\label{III.2.3.6-fr}
(\hyperref[III.2.3.5-fr]{2.3.5})의 가정에 더하여,
$\mathcal R(Y)$의 모든 영이 아닌 연접 부분 $\mathcal O_Y$-가군
$\mathcal J$에 대하여 어떤 가역 $\mathcal O_Y$-가군 $\mathcal L$이
존재하여
\[
\Gamma(Y,\mathcal L\otimes_{\mathcal O_Y}
\mathcal{H}om(\mathcal J,\mathcal O_Y))\neq0
\]
이라고 가정하자. 그러면 (\hyperref[III.2.3.5-fr]{2.3.5})의 명제에서
$\mathcal J$가 $\mathcal O_Y$의 아이디얼이라고 가정할 수 있다.
어떤 풍부한 $\mathcal O_Y$-가군이 존재하면 이 추가 조건은 항상
만족된다.
\end{corollary}

\begin{proof}
실제로 (0\textsubscript{I}, 5.4.2)에 의해
\[
\mathcal L\otimes\mathcal{H}om(\mathcal J,\mathcal O_Y)
=\mathcal{H}om(\mathcal L^{-1},\mathcal{H}om(\mathcal J,\mathcal O_Y))
=\mathcal{H}om(\mathcal J\otimes\mathcal L^{-1},\mathcal O_Y)
\]
이다. 따라서 가정은 $\mathcal J\otimes\mathcal L^{-1}$에서
$\mathcal O_Y$로 가는 영이 아닌 준동형 $u$가 존재한다는 뜻이다.
모든 $y\in Y$에 대하여
$(\mathcal J\otimes\mathcal L^{-1})_y$는 $\mathcal O_y$의 분수체
$\mathcal R(Y)_y$의 부분 $\mathcal O_y$-가군과 동일시되므로
(I, 7.1.5), $u_y$는 반드시 단사이다. 따라서 $u$는
$\mathcal J\otimes\mathcal L^{-1}$에서 $\mathcal O_Y$의 어떤 아이디얼
$\mathcal J'$로 가는 동형사상이다. 그런데
\[
\operatorname{Proj}\left(\bigoplus_{n\geq0}\mathcal J^n\right)
\quad\text{과}\quad
\operatorname{Proj}\left(
\bigoplus_{n\geq0}(\mathcal J\otimes\mathcal L^{-1})^n\right)
\]
은 $Y$-동형이므로 (II, 3.1.8), 이는 따름정리의 첫째 명제를
증명한다. 둘째 명제를 증명하기 위해
$\mathcal F=\mathcal{H}om_{\mathcal O_Y}(\mathcal J,\mathcal O_Y)$로
놓자. $\mathcal J|U$가 가역이 되도록 하는 열린집합 $U$가 존재하므로
$\mathcal F$는 연접이고 영이 아니다. $\mathcal L$이 풍부한
$\mathcal O_Y$-가군이면, 어떤 정수 $n$이 존재하여
$\mathcal F(n)=\mathcal F\otimes\mathcal L^{\otimes n}$은 $Y$ 위의
단면들로 생성된다 (II, 4.5.5). 더구나
$\Gamma(Y,\mathcal F(n))\neq0$이고, 이로부터 결론이 따른다.
\end{proof}

\begin{corollary}[2.3.7]
\label{III.2.3.7-fr}
$X$, $Y$를 체 $k$ 위의 두 사영 정역 스킴이라 하고,
$f:X\to Y$를 쌍유리 $k$-사상이라 하자. 그러면 $X$는 $Y$의 어떤
닫힌 부분스킴 $Y'$($Y'$은 반드시 축소일 필요가 없다)을 따라
블로업하여 얻은 $Y$-스킴과 $k$-동형이다.
\end{corollary}

\oldpage[III]{107}
\begin{proof}
실제로 $f$는 사영이고 (II, 5.5.5, (v)), $Y$가 $k$ 위에서
사영이므로 (\hyperref[III.2.3.6-fr]{2.3.6})의 추가 조건이 만족된다.
이제 따름정리 (\hyperref[III.2.3.6-fr]{2.3.6})의 연접 아이디얼
$\mathcal J$로 정의되는 $Y$의 닫힌 부분스킴 $Y'$을 택하면 충분하다.
\end{proof}

\begin{env}[2.3.8]
\label{III.2.3.8-fr}
\emph{비고.} 제IV장에서 약수의 개념을 연구할 때 다음을 보겠다.
(\hyperref[III.2.3.5-fr]{2.3.5})의 명제에서 환
$\mathcal O_y$ $(y\in Y)$들이 유일인수분해환이라고 가정하면
($Y$가 비특이이면 예를 들어 그러하다), $X$는 그 바탕공간이
$Y-U$에 포함되는 $Y$의 어떤 닫힌 부분준스킴 $Y'$을 따라
$Y$를 블로업하여 얻을 수 있다.
\end{env}

\subsection{기본정리의 일반화}
\label{subsection:III.2.4-fr}

\begin{theorem}[2.4.1]
\label{III.2.4.1-fr}
$Y$를 뇌터 준스킴, $\mathcal S$를 유한 생성 준연접
$\mathcal O_Y$-대수라 하자. $f:X\to Y$를 사영 사상,
$\mathcal S'=f^*(\mathcal S)$, $\mathcal M$을 유한 생성 준연접
$\mathcal S'$-가군이라 하자. 그러면 다음이 성립한다.
\begin{enumerate}
\item[(i)] 모든 $p\in\mathbb Z$에 대해 $R^pf_*(\mathcal M)$은 유한
생성 $\mathcal S$-가군이다.
\item[(ii)] 더 나아가 $\mathcal L$을 $f$에 대해 풍부한 가역
$\mathcal O_X$-가군이라 하고, 모든 $n\in\mathbb Z$에 대해
$\mathcal M(n)=\mathcal M\otimes\mathcal L^{\otimes n}$으로 놓자.
그러면 어떤 정수 $N$이 존재하여 $n\geq N$이면 모든 $p>0$에 대해
\[
R^pf_*(\mathcal M(n))=0
\tag{2.4.1.1}\label{III.2.4.1.1-fr}
\]
이고, 표준 준동형
\[
f^*(f_*(\mathcal M(n)))\longrightarrow\mathcal M(n)
\]
(0\textsubscript{I}, 4.4.3)은 전사이다.
\end{enumerate}
\end{theorem}

\begin{proof}
$Y'=\operatorname{Spec}(\mathcal S)$,
$X'=\operatorname{Spec}(\mathcal S')$로 놓자. 그러면
$X'=X\times_Y Y'$이다 (II, 1.5.5). 정의에 의해 아핀인 구조 사상들을
$g:Y'\to Y$, $g':X'\to X$라 하고,
$f'=f_{(Y')}:X'\to Y'$라 하자. 그러면 가환 도식
\[
\xymatrix{
X\ar[d]_f & X'\ar[l]_{g'}\ar[d]^{f'}\ar[dl]^{h}\\
Y & Y'\ar[l]^{g}
}
\]
이 있고, $f'$은 사영 사상이다 (II, 5.5.5, (iii)).
$h=f\circ g'=g\circ f'$로 놓자.

(i) $X'$을 $X$ 위의 아핀 스킴으로 보았을 때 준연접
$\mathcal S'$-가군 $\mathcal M$에 대응하는 $\mathcal O_{X'}$-가군을
$\widetilde{\mathcal M}$이라 하자 (II, 1.4.3). 그러면
$\mathcal M=g'_*(\widetilde{\mathcal M})$이다. $\mathcal M$이 유한
생성 $\mathcal S'$-가군이므로 $\widetilde{\mathcal M}$은 유한 생성
$\mathcal O_{X'}$-가군이다 (II, 1.4.5). 또한 $g$와 $f'$이 유한형이므로
$h$도 유한형이고 (II, 1.3.7 및 I, 6.3.4, (ii)), $X'$은 뇌터이다
(I, 6.3.7). 따라서 $\widetilde{\mathcal M}$은 연접이다.
$g'$이 아핀이므로 표준 준동형
\[
R^pf_*(\mathcal M)\longrightarrow R^ph_*(\widetilde{\mathcal M})
\]
은 전단사이다 (\hyperref[III.1.3.4-fr]{1.3.4}). 더 나아가 이 준동형은
$\mathcal S$-가군의 준동형이다. 실제로 $\mathcal M$의
$\mathcal S'$-가군 구조를 정의하는 표준 준동형
\[
g'_*(\mathcal O_{X'})\otimes_{\mathcal O_X}
g'_*(\widetilde{\mathcal M})\longrightarrow
g'_*(\widetilde{\mathcal M})
\tag{2.4.1.2}\label{III.2.4.1.2-fr}
\]
에서($\mathcal S'=g'_*(\mathcal O_{X'})$임을 상기하라), 표준적으로
준동형
\[
f_*(g'_*(\mathcal O_{X'}))\otimes
R^pf_*(g'_*(\widetilde{\mathcal M}))\longrightarrow
R^pf_*(g'_*(\widetilde{\mathcal M}))
\]
을 얻는다.

\oldpage[III]{108}
(0, 12.2.2). 또한 (\hyperref[III.2.4.1.2-fr]{2.4.1.2}) 자체가
$\widetilde{\mathcal M}$의 $\mathcal O_{X'}$-가군 구조를 정의하는
준동형
\[
\mathcal O_{X'}\otimes\widetilde{\mathcal M}
\longrightarrow\widetilde{\mathcal M}
\]
으로부터 (0\textsubscript{I}, 4.2.2.1)을 적용하여 얻어지므로, 도식
\[
\xymatrix{
f_*(g'_*(\mathcal O_{X'}))\otimes
R^pf_*(g'_*(\widetilde{\mathcal M}))
\ar[r]\ar[d] &
R^pf_*(g'_*(\widetilde{\mathcal M}))\ar[d]\\
h_*(\mathcal O_{X'})\otimes R^ph_*(\widetilde{\mathcal M})
\ar[r] & R^ph_*(\widetilde{\mathcal M})
}
\]
은 가환한다 (0, 12.2.6). 수평 화살표들을 표준 준동형
\[
\mathcal S\longrightarrow f_*(f^*(\mathcal S))=f_*(\mathcal S')
=f_*(g'_*(\mathcal O_{X'}))=h_*(\mathcal O_{X'})
\]
에서 오는 준동형과 합성하면 위의 주장을 얻는다. 한편 $g$가 아핀이고
$f'$이 분리이고 준콤팩트이므로, 표준 준동형
\[
R^ph_*(\widetilde{\mathcal M})\longrightarrow
g_*(R^pf'_*(\widetilde{\mathcal M}))
\]
은 전단사이다 (\hyperref[III.1.4.14-fr]{1.4.14}). 위와 같은 방식으로
이번에는 (0, 12.2.6.2)의 가환성을 사용하면, 이것이
$\mathcal S$-가군의 동형사상임을 증명할 수 있다. 그런데 $f'$은
사영이고 $\widetilde{\mathcal M}$은 연접이므로
(\hyperref[III.2.2.1-fr]{2.2.1})에 의해
$R^pf'_*(\widetilde{\mathcal M})$은 연접 $\mathcal O_{Y'}$-가군이다.
따라서 $g_*(R^pf'_*(\widetilde{\mathcal M}))$은 유한 생성
$\mathcal S$-가군이다 (II, 1.4.5).

(ii) $\mathcal L'=g'^*(\mathcal L)$로 놓자. 이는 가역
$\mathcal O_{X'}$-가군이다. 모든 $n\in\mathbb Z$에 대해 표준 동형을
제외하면
\[
g'_*(\widetilde{\mathcal M}\otimes\mathcal L'^{\otimes n})
=g'_*(\widetilde{\mathcal M})\otimes\mathcal L^{\otimes n}
=\mathcal M(n)
\]
이다 (0\textsubscript{I}, 5.4.10). $\widetilde{\mathcal M}$에 대해
(i)에서 한 논증을
$\widetilde{\mathcal M}\otimes\mathcal L'^{\otimes n}$에 적용하면
\[
R^pf_*(g'_*(\widetilde{\mathcal M}\otimes\mathcal L'^{\otimes n}))
\quad\text{은}\quad
g_*(R^pf'_*(\widetilde{\mathcal M}\otimes\mathcal L'^{\otimes n}))
\]
과 동형이다. $\mathcal L'$은 $f'$에 대해 풍부하므로
(II, 4.6.13, (iii)), (\hyperref[III.2.2.1-fr]{2.2.1})에 의해 어떤
정수 $N$이 존재하여 모든 $p>0$, $n\geq N$에 대해
\[
R^pf'_*(\widetilde{\mathcal M}\otimes\mathcal L'^{\otimes n})=0
\]
이다. 이는 (\hyperref[III.2.4.1.1-fr]{2.4.1.1})을 증명한다.
다시 (\hyperref[III.2.2.1-fr]{2.2.1})에 의해, $n\geq N$이면 표준
준동형
\[
f'^*(f'_*(\widetilde{\mathcal M}\otimes\mathcal L'^{\otimes n}))
\longrightarrow\widetilde{\mathcal M}\otimes\mathcal L'^{\otimes n}
\]
이 전사가 되도록 $N$을 택할 수 있다. $g'_*$은 완전한 함자이므로
(II, 1.4.4), 대응하는 준동형
\[
g'_*(f'^*(f'_*(\widetilde{\mathcal M}\otimes
\mathcal L'^{\otimes n})))
\longrightarrow
g'_*(\widetilde{\mathcal M}\otimes\mathcal L'^{\otimes n})
=\mathcal M(n)
\]
은 전사이다. 한편 (II, 1.5.2)에 의해
\[
g'_*(f'^*(f'_*(\widetilde{\mathcal M}\otimes
\mathcal L'^{\otimes n})))
=f^*(g_*(f'_*(\widetilde{\mathcal M}\otimes
\mathcal L'^{\otimes n})))
\]
이고 $g_*\circ f'_* = f_*\circ g'_*$이므로, 결국
\[
g'_*(f'^*(f'_*(\widetilde{\mathcal M}\otimes\mathcal L'^{\otimes n})))
=f^*(f_*(g'_*(\widetilde{\mathcal M}\otimes
\mathcal L'^{\otimes n})))=f^*(f_*(\mathcal M(n)))
\]
이다. 이로써 증명이 끝난다.
\end{proof}

\begin{env}[2.4.2]
\label{III.2.4.2-fr}
특히 $\mathcal S$가 양의 차수들을 갖는 등급 $\mathcal O_Y$-대수이고
$\mathcal M=\bigoplus_{k\in\mathbb Z}\mathcal M_k$가 등급
$\mathcal S'$-가군일 때 (\hyperref[III.2.4.1-fr]{2.4.1})을 적용하게
된다.

\oldpage[III]{109}
$\mathcal S$, $\mathcal M$에 대해 같은 유한성 가정을 두면,
(\hyperref[III.2.4.1-fr]{2.4.1})로부터 모든 $p$에 대해
\[
\bigoplus_{k\in\mathbb Z}R^pf_*(\mathcal M_k)
\]
가 유한 생성 $\mathcal S$-가군임을 얻는다. 또한
(\hyperref[III.2.4.1-fr]{2.4.1}, (ii))의 추가 가정 아래에서는 어떤
$N$이 존재하여 $n\geq N$이면 모든 $p>0$, $k\in\mathbb Z$에 대해
$R^pf_*(\mathcal M_k(n))=0$이고, 표준 준동형
$f^*(f_*(\mathcal M_k(n)))\to\mathcal M_k(n)$은 모든
$k\in\mathbb Z$에 대해 전사이다.
\end{env}

\subsection{오일러--푸앵카레 특성과 힐베르트 다항식}
\label{subsection:III.2.5-fr}

\begin{env}[2.5.1]
\label{III.2.5.1-fr}
$A$를 아르틴 환, $X$를 $Y=\operatorname{Spec}(A)$ 위의 사영
$A$-스킴이라 하자. 모든 연접 $\mathcal O_X$-가군 $\mathcal F$에 대해
$H^i(X,\mathcal F)$ $(i\geq0)$들은 유한 생성 $A$-가군이다
(\hyperref[III.2.2.1-fr]{2.2.1}). $A$가 아르틴이므로 여기서는 유한
길이를 갖는다. 또한 (\hyperref[III.2.2.1-fr]{2.2.1})에 의해
$H^i(X,\mathcal F)$는 유한 개의 $i\geq0$을 제외하면 0이다. 따라서
정수
\[
\chi_A(\mathcal F)=\sum_{i=0}^{\infty}(-1)^i
\operatorname{long}(H^i(X,\mathcal F))
\tag{2.5.1.1}\label{III.2.5.1.1-fr}
\]
는 모든 연접 $\mathcal O_X$-가군 $\mathcal F$에 대해 정의된다.
$A$가 아르틴 국소환이면 $\chi_A(\mathcal F)$를 $\mathcal F$의
($A$에 대한) 오일러--푸앵카레 특성이라 한다.
$\mathcal F=\mathcal O_X$일 때 $\chi_A(\mathcal O_X)$를 $X$의
($A$에 대한) 산술 종수라고 한다.
\end{env}

\begin{proposition}[2.5.2]
\label{III.2.5.2-fr}
\[
0\longrightarrow\mathcal F'\longrightarrow\mathcal F
\longrightarrow\mathcal F''\longrightarrow0
\]
을 연접 $\mathcal O_X$-가군들의 완전열이라 하자. 그러면
\[
\chi_A(\mathcal F)=\chi_A(\mathcal F')+\chi_A(\mathcal F'').
\tag{2.5.2.1}\label{III.2.5.2.1-fr}
\]
\end{proposition}

\begin{proof}
$\mathcal F$, $\mathcal F'$, $\mathcal F''$의 코호몰로지 가군들은
유한 개를 제외하면 0이므로, 어떤 정수 $r>0$이 존재하여 코호몰로지
완전열은
\[
0\to H^0(X,\mathcal F')\to H^0(X,\mathcal F)
\to H^0(X,\mathcal F'')\to H^1(X,\mathcal F')\to\cdots
\]
\[
\cdots\to H^r(X,\mathcal F')\to H^r(X,\mathcal F)
\to H^r(X,\mathcal F'')\to0
\]
으로 쓰인다. 양 끝이 0인 유한 길이 $A$-가군들의 완전열에서는
길이들의 교대합이 0임을 안다 (0, 11.10.1). 이 결과를 적용하면 공식
(\hyperref[III.2.5.2.1-fr]{2.5.2.1})을 즉시 얻는다.
\end{proof}

$X$가 준콤팩트 $A$-스킴이고 모든 연접 $\mathcal O_X$-가군
$\mathcal F$에 대해 $A$-가군 $H^i(X,\mathcal F)$들이 유한 생성임을
아는 모든 경우에도 (\hyperref[III.2.5.2-fr]{2.5.2})의 결과가 적용됨을
유의하자 (\hyperref[III.1.4.12-fr]{1.4.12}).

\begin{theorem}[2.5.3]
\label{III.2.5.3-fr}
$A$를 아르틴 국소환, $X$를 $Y=\operatorname{Spec}(A)$ 위의 사영 스킴,
$\mathcal L$을 $Y$에 대해 매우 풍부한 가역 $\mathcal O_X$-가군,
$\mathcal F$를 연접 $\mathcal O_X$-가군이라 하자. 모든
$n\in\mathbb Z$에 대해
$\mathcal F(n)=\mathcal F\otimes_{\mathcal O_X}\mathcal L^{\otimes n}$
으로 놓는다.
\begin{enumerate}
\item[(i)] 모든 $n\in\mathbb Z$에 대해
$\chi_A(\mathcal F(n))=P(n)$이 되도록 하는 유일한 다항식
$P\in\mathbb Q[T]$가 존재한다. $P$를 $\mathcal F$의 $A$에 대한
힐베르트 다항식이라 한다.
\item[(ii)] 충분히 큰 $n$에 대해
$\chi_A(\mathcal F(n))=\operatorname{long}_A
\Gamma(X,\mathcal F(n))$이다.
\item[(iii)] $\chi_A(\mathcal F(n))$의 최고차항의 계수는
$\geq0$이다.
\end{enumerate}
\end{theorem}

제IV장의 차원 개념을 다루는 절에서는
$\chi_A(\mathcal F(n))$의 차수가 $\mathcal F$의 받침의 차원과
같음도 증명하겠다.

\begin{proof}
충분히 큰 $n$이면 모든 $i>0$에 대해
$H^i(X,\mathcal F(n))=0$이므로
(\hyperref[III.2.2.1-fr]{2.2.1}),

\oldpage[III]{110}
충분히 큰 $n$에 대해
\[
\chi_A(\mathcal F(n))=\operatorname{long}H^0(X,\mathcal F(n))
=\operatorname{long}\Gamma(X,\mathcal F(n))
\]
이다. 이로써 (ii)가 따르고, 충분히 큰 $n$에 대해
$\chi_A(\mathcal F(n))\geq0$이다. 따라서 (iii)은 (i)에서 따른다.
또한 (i)의 유일성 명제는 명백하므로 다항식 $P$의 존재만 증명하면
된다.

먼저 $\mathfrak m$이 $A$의 극대 아이디얼일 때
$\mathfrak m\mathcal F=0$이라고 가정할 수 있음을 보이자. 실제로
$\mathfrak m^s=0$인 어떤 정수 $s>0$이 존재하므로 $\mathcal F(n)$은
유한 여과
\[
\mathcal F(n)\supset\mathfrak m\mathcal F(n)\supset\cdots
\supset\mathfrak m^{s-1}\mathcal F(n)\supset0
\]
를 갖는다. (\hyperref[III.2.5.2.1-fr]{2.5.2.1})과 귀납법으로
\[
\chi_A(\mathcal F(n))=\sum_{k=1}^{s}
\chi_A(\mathfrak m^{k-1}\mathcal F(n)/\mathfrak m^k\mathcal F(n))
\]
을 얻는다. 그런데
$\mathfrak m^{k-1}\mathcal F(n)/\mathfrak m^k\mathcal F(n)
=\mathcal F'_k(n)$이고
$\mathcal F'_k=\mathfrak m^{k-1}\mathcal F/\mathfrak m^k\mathcal F$이므로,
이는 우리의 주장을 증명한다.

따라서 $\mathfrak m\mathcal F=0$이라고 가정하자. $X'$을 구조 사상
$X\to\operatorname{Spec}(A)$에 의한 $\operatorname{Spec}(A)$의 유일한
닫힌점의 역상인 $X$의 닫힌 부분스킴이라 하고, $j:X'\to X$를 표준
포함이라 하자. 그러면 $\mathcal F=j_*(\mathcal F')$이며,
$\mathcal F'$은 연접 $\mathcal O_{X'}$-가군이다. 또한
$K=A/\mathfrak m$일 때 $X'$은 $\operatorname{Spec}(K)$ 위의 사영
스킴이다. $\mathcal L'=j^*(\mathcal L)$로 놓으면 $\mathcal L'$은
$\operatorname{Spec}(K)$에 대해 매우 풍부하고 (II, 4.4.10),
\[
\mathcal F(n)=j_*(\mathcal F'(n)),\qquad
\mathcal F'(n)=\mathcal F'\otimes_{\mathcal O_{X'}}
\mathcal L'^{\otimes n}
\]
이다 (0\textsubscript{I}, 5.4.10). 따라서
$\chi_A(\mathcal F(n))=\chi_K(\mathcal F'(n))$이고
(G, II, 4.9.1), $A$가 체인 경우로 귀착된다.

이제 어떤 적당한 $r$에 대해 $X$를 $P=\mathbf P_A^r$의 닫힌
부분스킴으로 볼 수 있다 (II, 5.5.4, (ii)). $i:X\to P$가 표준
포함이면 위와 같이
$\chi_A(\mathcal F(n))=\chi_A(i_*(\mathcal F)(n))$임을 알 수 있다.
따라서 $A$가 체이고
$X=\mathbf P_A^r=\operatorname{Proj}(S)$,
$S=A[T_0,\ldots,T_r]$인 경우로 한정할 수 있다.

이제 $\mathcal F=\widetilde M$이며, $M$은 유한 생성 등급
$S$-가군이다 (II, 2.7.8). 힐베르트 시지지 정리 (M, VIII, 6.5)에
의해 $M$은 유한 생성 자유 등급 $S$-가군들에 의한 유한 분해
\[
0\to L_q\to L_{q-1}\to\cdots\to L_1\to M\to0
\]
를 갖는다. $M\mapsto\widetilde M$은 $M$에 관해 완전한 함자이므로
(II, 2.5.4), 완전열
\[
0\to\widetilde L_q\to\widetilde L_{q-1}\to\cdots
\to\widetilde L_1\to\widetilde M\to0
\]
도 존재하고, 따라서 모든 $n\in\mathbb Z$에 대해 열
\[
0\to\widetilde L_q(n)\to\widetilde L_{q-1}(n)\to\cdots
\to\widetilde L_1(n)\to\widetilde M(n)\to0
\]
은 완전하다. 명제 (2.5.2)를 $q$에 대한 귀납법으로 적용하면
\[
\chi_A(\widetilde M(n))=\sum_{j=1}^{q}(-1)^{j+1}
\chi_A(\widetilde L_j(n))
\]
이다. 그러므로 (i)를 증명하기 위해 $M$이 유한 생성 자유 등급
가군인 경우, 나아가 어떤 $h\in\mathbb Z$에 대해 $M=S(h)$인 경우로
귀착된다. 이때
$\widetilde M(n)=(M(n))^\sim=(S(n+h))^\sim$이므로 (II, 2.5.15),
결국 정리는 다음 보조정리에서 따른다.
\end{proof}

\oldpage[III]{111}
\begin{lemma}[2.5.3.1]
\label{III.2.5.3.1-fr}
$A$를 체, $r>0$을 정수, $X=\mathbf P_A^r$라 하자. 그러면 모든
$n\in\mathbb Z$에 대해
\[
\chi_A(\mathcal O_X(n))=\binom{n+r}{r}
\]
이다.
\end{lemma}

\begin{proof}
실제로 $n\geq0$이면
$\chi_A(\mathcal O_X(n))=\operatorname{long}H^0(X,\mathcal O_X(n))$이고,
이는 총차수가 $n$인 $T_i$들의 단항식의 개수, 즉
$\binom{n+r}{r}$이다 (\hyperref[III.2.1.12-fr]{2.1.12}).
$n\leq-r-1$이면 마찬가지로
\[
\chi_A(\mathcal O_X(n))=(-1)^r
\operatorname{long}H^r(X,\mathcal O_X(n)).
\]
$n=-r-h$라 하면 $A$ 위에서 $H^r(X,\mathcal O_X(n))$의 차원은
$p_i>0$, $\sum_{i=0}^{r}p_i=r+h$인 정수열
$(p_i)_{0\leq i\leq r}$의 개수이다
(\hyperref[III.2.1.12-fr]{2.1.12}). 이는 $q_i\geq0$,
$\sum_{i=0}^{r}q_i=h-1$인 정수열 $(q_i)_{0\leq i\leq r}$의 개수와
같으므로
\[
\binom{h+r-1}{r}=(-1)^r\binom{n+r}{r}
\]
이다. 마지막으로 $-r\leq n<0$이면 $\binom{n+r}{r}=0$이고,
한편 모든 $i\geq0$에 대해 $H^i(X,\mathcal O_X(n))=0$이다
(\hyperref[III.2.1.12-fr]{2.1.12}). 이로써 보조정리가 증명된다.
\end{proof}

\begin{corollary}[2.5.4]
\label{III.2.5.4-fr}
$A$를 아르틴 국소환, $S$를 $S_1$로 생성되는 유한 생성 등급
$A$-대수, $M$을 유한 생성 등급 $S$-가군,
$X=\operatorname{Proj}(S)$라 하자. 충분히 큰 $n$에 대해
\[
\chi_A(\widetilde M(n))=\operatorname{long}M_n
\]
이다.
\end{corollary}

\begin{proof}
충분히 큰 $n$에 대해 $M_n$과 $\Gamma(X,\widetilde M(n))$이
동형이라는 사실에서 따른다 (\hyperref[III.2.3.1-fr]{2.3.1}).
\end{proof}

\subsection{응용: 풍부성 판정법}
\label{subsection:III.2.6-fr}

\begin{proposition}[2.6.1]
\label{III.2.6.1-fr}
$Y$를 뇌터 준스킴, $f:X\to Y$를 고유 사상,
$\mathcal L$을 가역 $\mathcal O_X$-가군이라 하자. 다음 조건들은
동치이다.
\begin{enumerate}
\item[a)] $\mathcal L$은 $f$에 대해 풍부하다.
\item[b)] 모든 연접 $\mathcal O_X$-가군 $\mathcal F$에 대해 어떤
정수 $N$이 존재하여 $n\geq N$이면 모든 $q>0$에 대해
\[
R^qf_*(\mathcal F\otimes\mathcal L^{\otimes n})=0
\]
이다.
\item[c)] $\mathcal O_X$의 모든 연접 아이디얼 $\mathcal J$에 대해
어떤 정수 $N$이 존재하여 $n\geq N$이면
\[
R^1f_*(\mathcal J\otimes\mathcal L^{\otimes n})=0
\]
이다.
\end{enumerate}
\end{proposition}

\begin{proof}
a)가 b)를 함의함은 이미 보았다
(\hyperref[III.2.2.1-fr]{2.2.1}, (ii)). b)가 c)를 함의함은 자명하므로,
c)가 a)를 함의함을 증명하면 된다. $Y$가 아핀인 경우로 한정하여
(II, 4.6.4), 이 경우 $\mathcal L$이 풍부함을 증명할 수 있다.
$h$가 $X$ 위의 $\mathcal L^{\otimes n}$ $(n>0)$의 모든 단면을
움직일 때 아핀인 $X_h$들이 $X$를 덮음을 보이면 충분하다
(II, 4.5.2). 이를 위해 $X$의 모든 닫힌점 $x$와 $x$의 모든 아핀
열린 근방 $U$에 대해, 어떤 $n$과
$h\in\Gamma(X,\mathcal L^{\otimes n})$이 존재하여
$x\in X_h\subset U$임을 보이자. 그러면 $X_h$는 반드시 아핀이고
(I, 1.3.6), 이러한 $X_h$들의 합집합은 $X$의 모든 닫힌점을 포함하는
열린집합이다. $X$가 뇌터이므로 이는 $X$ 자체이다
(I, 6.3.7 및 0\textsubscript{I}, 2.1.3).

$X-U$ (각각 $(X-U)\cup\{x\}$)를 바탕공간으로 하는 $X$의 축소 닫힌
부분준스킴을 정의하는 $\mathcal O_X$의 준연접 아이디얼층을
$\mathcal J$ (각각 $\mathcal J'$)라 하자 (I, 5.2.1).
$\mathcal J$, $\mathcal J'$은 연접이고 (I, 6.1.1),
$\mathcal J'\subset\mathcal J$이며,
$\mathcal J''=\mathcal J/\mathcal J'$은 받침이 $\{x\}$이고
$\mathcal J''_x=k(x)$인 연접 $\mathcal O_X$-가군이다
(0\textsubscript{I}, 5.3.3). $\mathcal L$이 국소 자유이므로 모든 $n$에
대해 열
\[
0\to\mathcal J'\otimes\mathcal L^{\otimes n}
\to\mathcal J\otimes\mathcal L^{\otimes n}
\to\mathcal J''\otimes\mathcal L^{\otimes n}\to0
\]
은 완전하다. 가정에 의해 충분히 큰 어떤 $n$에 대해
$H^1(X,\mathcal J'\otimes\mathcal L^{\otimes n})=0$이다. 따라서
코호몰로지 완전열로부터 준동형

\oldpage[III]{112}
\[
\Gamma(X,\mathcal J\otimes\mathcal L^{\otimes n})
\longrightarrow
\Gamma(X,\mathcal J''\otimes\mathcal L^{\otimes n})
\]
이 전사임을 얻는다. $g(x)\neq0$인
$\mathcal J''\otimes\mathcal L^{\otimes n}$의 $X$ 위의 단면 $g$는
따라서 어떤 단면
\[
h\in\Gamma(X,\mathcal J\otimes\mathcal L^{\otimes n})
\subset\Gamma(X,\mathcal L^{\otimes n})
\]
의 상이다. 여기서 포함은 (0\textsubscript{I}, 5.4.1)에 의해
$\mathcal J\otimes\mathcal L^{\otimes n}$이
$\mathcal L^{\otimes n}$의 부분 $\mathcal O_X$-가군이기 때문이다.
정의에 의해 $h(x)\neq0$이고 $z\notin U$이면 $h(z)=0$이다. 이로써
증명이 끝난다.
\end{proof}

\begin{proposition}[2.6.2]
\label{III.2.6.2-fr}
$Y$를 뇌터 준스킴, $f:X\to Y$를 유한형 사상,
$g:X'\to X$를 유한 전사 사상, $\mathcal L$을 가역
$\mathcal O_X$-가군이라 하고 $\mathcal L'=g^*(\mathcal L)$로 놓자.
다음 조건을 가정한다. $Y$ 위에서 고유인 어떤 부분집합 $Z\subset X$가
존재하여 (II, 5.4.10), 모든 $x\in X-Z$에 대해 $X$가 $x$에서
정규이거나 $(g_*(\mathcal O_{X'}))_x$가 자유 $\mathcal O_x$-가군이다.
이 조건 아래에서 $\mathcal L$이 $f$에 대해 풍부일 필요충분조건은
$\mathcal L'$이 $f\circ g$에 대해 풍부인 것이다.
\end{proposition}

\begin{proof}
\begin{env}[2.6.2.1]
\label{III.2.6.2.1-fr}
$g$가 아핀이므로 조건이 필요함은 (II, 5.1.12)에서 따른다. 충분함을
보이기 위해 $Y$가 아핀이라고 가정할 수 있다 (II, 4.6.4). 더 나아가
$X$가 축소인 경우로 한정할 수 있음을 보이자. 실제로
$j:X_{\mathrm{red}}\to X$를 표준 포함이라 하고,
$X_1=X_{\mathrm{red}}$, $X'_1=X'\times_X X_1$으로 놓자. 그러면 가환
도식
\[
\xymatrix{
X'\ar[d]_g & X'_1\ar[l]_{j'}\ar[d]^{g_1\, .}\\
X & X_1\ar[l]^{j}
}
\tag{2.6.2.2}\label{III.2.6.2.2-fr}
\]
을 얻는다. 사상 $f\circ j$는 유한형이고 (I, 6.3.4), $g_1$은 유한
사상이다 (II, 6.1.5, (iii)). $\mathcal L'$이 $f\circ g$에 대해
풍부이면 $j'$이 닫힌 몰입이므로
$j'^*(\mathcal L')$은 $f\circ g\circ j'$에 대해 풍부하다
(II, 5.1.12 및 I, 4.3.2). $Z_1=j^{-1}(Z)$로 놓으면 $Z_1$은 $Y$ 위에서
고유이다 (II, 5.4.10). 한편 $X$가 점 $x$에서 정규이면
$X_{\mathrm{red}}$도 명백히 그러하다. 마지막으로
$(g_*(\mathcal O_{X'}))_x$가 자유 $\mathcal O_x$-가군이면
(II, 1.5.2)에서 곧바로
$((g_1)_*(\mathcal O_{X'_1}))_x$가 자유 $\mathcal O_{X_1,x}$-가군임을
얻는다. 끝으로 $X$가 뇌터이므로 (I, 6.3.7), $j^*(\mathcal L)$이
풍부이면 $\mathcal L$도 풍부하다 (II, 4.5.14). 또한
$j'^*(\mathcal L')=g_1^*(j^*(\mathcal L))$이므로 앞서 말한 귀착이
끝난다. 이제부터 $Y$가 아핀이고 $X$가 축소라고 가정한다.
\end{env}

이제 (II, 6.6.11)의 가정들이 만족되므로, 어떤 축소 $Y$-준스킴
$X_2$와 유한 쌍유리 $Y$-사상 $h:X_2\to X$가 존재하여,
$h$를 $h^{-1}(X-Z)$에 제한한 것은 $X-Z$ 위의 동형사상이고
$h^*(\mathcal L)$은 풍부하다. $X'$을 $X_2$로 바꾸면, $g$가 방금
$h$에 대해 열거한 성질들을 더 만족한다고 가정하여 명제를 증명하는
문제로 귀착된다. 계속해서 $Z$를 바탕공간으로 $Z$를 가지며 $Y$ 위에서
고유인 $X$의 부분준스킴도 $Z$로 나타내겠다 (II, 5.4.10).

\begin{env}[2.6.2.3]
\label{III.2.6.2.3-fr}
이제 $X_1$을 $X$의 닫힌 부분준스킴, $j:X_1\to X$를 표준 포함,
$X'_1=g^{-1}(X_1)=X'\times_X X_1$을 그 역상,
$j':X'_1\to X'$을 표준 포함이라 하자. 그러면 가환 도식
(\hyperref[III.2.6.2.2-fr]{2.6.2.2})가 있다.
$\mathcal L_1=j^*(\mathcal L)$,
$\mathcal L'_1=j'^*(\mathcal L')=g_1^*(\mathcal L_1)$로 놓으면,
$\mathcal L'_1$은 $f\circ g\circ j'$에 대해 풍부하다
(II, 5.1.12). $Z_1=j^{-1}(Z)$로 놓으면 $X_1$의 닫힌 부분준스킴
$Z_1$은 $Y$ 위에서 고유이다 (II, 5.4.2, (ii)). 다시 말해 $X_1$,
$\mathcal L_1$, $g_1$, $Z_1$에 대해
(\hyperref[III.2.6.2-fr]{2.6.2})의 가정들이 만족된다.
\end{env}

\oldpage[III]{113}
이 사실을 이용하여, $g$를 $g^{-1}(X-Z)$에 제한한 것이 $X-Z$ 위의
동형사상인 경우에 뇌터 귀납법 (0\textsubscript{I}, 2.2.2)으로
(\hyperref[III.2.6.2-fr]{2.6.2})를 증명할 수 있다. 앞서
(\hyperref[III.2.6.2.1-fr]{2.6.2.1})에서 본 대로 이것이면 충분하다.
즉 바탕공간이 $X$와 다른 $X$의 모든 닫힌 부분준스킴 $X_1$에 대해
(\hyperref[III.2.6.2-fr]{2.6.2})의 결론이 $\mathcal L_1$에 대해
참이라고 가정하고, 그 결론이 $\mathcal L$에 대해서도 참임을 보이면
된다.

\begin{env}[2.6.2.4]
\label{III.2.6.2.4-fr}
$\mathcal A=\mathcal O_X$, $\mathcal B=g_*(\mathcal O_{X'})$로 놓자.
그러면 $\mathcal B$는 $\mathcal R(X)$의 부분 $\mathcal A$-대수이고,
연접 $\mathcal A$-가군이다. 또한
$\mathcal B|(X-Z)=\mathcal A|(X-Z)$이다. $\mathcal K$를
$\mathcal A$ 위의 $\mathcal B$의 \emph{도체}, 즉
$\mathcal B\mathcal K\subset\mathcal A$를 만족하는
$\mathcal A$의 가장 큰 준연접 부분 $\mathcal A$-가군이라 하자. 이는
$\mathcal A$-가군 $\mathcal B/\mathcal A$의 소멸자와도 같다
(0\textsubscript{I}, 5.3.7). 따라서
$\mathcal B\mathcal K=\mathcal K$이다. $g$가 어떤 근방 $W_x$의 역상
$g^{-1}(W_x)$에서 $W_x$로 가는 동형사상이 되는 모든 점 $x$에서는
$\mathcal K_x=\mathcal A_x$임이 명백하다. 특히 이는 $X-Z$의 모든
점과 $X$의 각 기약 성분의 일반점의 어떤 근방에서 성립한다.
$\mathcal K$가 정의하는 $X$의 닫힌 부분준스킴
$Z_1=\operatorname{Spec}(\mathcal A/\mathcal K)$를 생각하자.
바탕공간 $Z_1$이 $Z$에서 닫혀 있으므로 $Z_1$도 $Y$ 위에서 고유이다
(II, 5.4.10). 또한 $\mathcal K$의 정의로부터
$\mathcal B|(X-Z_1)=\mathcal A|(X-Z_1)$임을 얻는다. 따라서 언제나
$Z=\operatorname{Spec}(\mathcal A/\mathcal K)$인 경우로 귀착할 수
있다. 위에서 $X-Z_1$이 $X$의 공집합이 아닌 열린집합임을 보았으므로,
항상 바탕공간 $Z$가 $X$와 다르다고 가정할 수 있다.
\end{env}

\begin{env}[2.6.2.5]
\label{III.2.6.2.5-fr}
$g$가 아핀이므로 $X'=\operatorname{Spec}(\mathcal B)$로 보고,
$\mathcal K'=\widetilde{\mathcal K}$라 하자. 이는
$g_*(\mathcal K')=\mathcal K$인 $\mathcal O_{X'}$의 연접 아이디얼이다
(II, 1.4.1). $X'$의 닫힌 부분준스킴
$Z'=g^{-1}(Z)=Z\times_X X'$은 $\mathcal K'$로 정의되고
$\operatorname{Spec}(\mathcal B/\mathcal K)$와 같다 (II, 1.4.10).
$h:Z'\to Z$가 유한 사상이므로 (II, 6.1.5, (iii)), $Z'$은 $Y$ 위에서
고유이다 (II, 6.1.11 및 5.4.2, (ii)).

이제 모든 $x\in X$와 모든 열린 근방 $U$에 대해, 어떤 $n>0$과
$X$ 위의 $\mathcal L^{\otimes n}$의 단면 $s$가 존재하여
$x\in X_s\subset U$임을 증명해야 한다 (II, 4.5.2). 두 경우로
나눈다.

$1^\circ$ $x\in X-Z$인 경우. 이때 $U\subset X-Z$라고 가정해도 되고,
따라서 열린집합 $U'=g^{-1}(U)$은 $Z'$과 만나지 않는다.
$\mathcal L'$이 풍부하므로 어떤 $n>0$과 $X'$ 위의
$\mathcal L'^{\otimes n}$의 단면 $s'$이 존재하여
$x'=g^{-1}(x)\in X'_{s'}\subset g^{-1}(U)$이다 (II, 4.5.2).
또한 $\mathcal K'\otimes\mathcal L'^{\otimes n}$이 $X'$ 위의
단면들로 생성된다고 가정할 수 있다 (II, 4.5.5). 따라서
$\mathcal K'_{x'}=\mathcal O_{x'}$이므로 이 단면들 가운데
$s''(x')\neq0$인 $s''$이 존재한다. 이를 $s'$과 곱하면($n$을 $2n$으로
바꾸는 것과 같다), $x'\in X'_{s''}\subset g^{-1}(U)$라고도 가정할 수
있다. 이제 (0\textsubscript{I}, 5.4.10)에 의해 표준 동형사상
\[
\Gamma(X,\mathcal K\otimes\mathcal L^{\otimes n})
\xrightarrow{\sim}
\Gamma(X',\mathcal K'\otimes\mathcal L'^{\otimes n})
\]
이 있다. 이 동형사상에서 $s''$에 대응하는
$\mathcal K\otimes\mathcal L^{\otimes n}$의 단면 $s$는 명백히 원하는
성질들을 갖는다.

$2^\circ$ $x\in Z$인 경우. 바탕공간이 $X-U$인 $X$의 축소 닫힌
부분준스킴을 정의하는 $\mathcal O_X$의 연접 아이디얼을 $\mathcal J$라
하자. $\mathcal B$ 안의 두 연접 아이디얼
$\mathcal J\mathcal B$와
$\mathcal J\mathcal K=\mathcal J(\mathcal K\mathcal B)
=\mathcal K(\mathcal J\mathcal B)$를 생각하면 포함 도식

\oldpage[III]{114}
\[
\xymatrix@R=1.5em{
\mathcal J\mathcal B\ar@{^{(}->}[r] & \mathcal B\\
\mathcal J\ar@{^{(}->}[r]\ar@{^{(}->}[u] &
\mathcal A\ar@{^{(}->}[u]\\
\mathcal J\mathcal K\mathcal B=\mathcal J\mathcal K
\ar@{^{(}->}[r]\ar@{^{(}->}[u] &
\mathcal K\ar@{^{(}->}[u]
}
\tag{2.6.2.6}\label{III.2.6.2.6-fr}
\]
을 얻는다. $\mathcal J'=(\mathcal J\mathcal B)^\sim$를
$\mathcal O_{X'}$의 연접 아이디얼이라 하자. 그러면
$\mathcal J\mathcal B=g_*(\mathcal J')$이고,
$\mathcal J'\mathcal K'=(\mathcal J\mathcal K\mathcal B)^\sim$이므로
\[
\mathcal J'/\mathcal J'\mathcal K'
=(\mathcal J\mathcal B/\mathcal J\mathcal K\mathcal B)^\sim
\]
이다 (II, 1.4.4). $Z$와 만나지 않는 모든 열린집합 $V$에서
$\mathcal J|V=\mathcal J\mathcal K|V$이므로,
$\mathcal J'/\mathcal J'\mathcal K'$의 받침은 $Z'$에 포함된다.
$Z'$이 $Y$ 위에서 고유이므로 (\hyperref[III.2.2.4-fr]{2.2.4})를
적용할 수 있다. 따라서 충분히 큰 $n$에 대해 표준 사상
\[
\Gamma(X',\mathcal J'\otimes\mathcal L'^{\otimes n})
\longrightarrow
\Gamma(X',(\mathcal J'/\mathcal J'\mathcal K')
\otimes\mathcal L'^{\otimes n})
\]
은 전사이다. (0\textsubscript{I}, 5.4.10)에 의해, 이로부터 표준 사상
\[
\Gamma(X,\mathcal J\mathcal B\otimes\mathcal L^{\otimes n})
\longrightarrow
\Gamma(X,(\mathcal J\mathcal B/\mathcal J\mathcal K\mathcal B)
\otimes\mathcal L^{\otimes n})
\]
이 전사임을 얻는다.

$i:Z\to X$, $i':Z'\to X'$을 표준 포함들이라 하자. 그러면 가환 도식
\[
\xymatrix{
X'\ar[d]_g & Z'\ar[l]_{i'}\ar[d]^{h}\\
X & Z\ar[l]^{i}
}
\]
이 있다. $\mathcal M=i^*(\mathcal L)$,
$\mathcal M'=i'^*(\mathcal L')$로 놓자. $\mathcal L'$이 풍부하므로
$\mathcal M'$도 풍부하고 (II, 5.1.12), 한편
$\mathcal M'=h^*(\mathcal M)$이다. 따라서 뇌터 귀납가정에 의해
($Z\neq X$이므로) $\mathcal M$은 풍부하다. 그러므로 충분히 큰
$n$에 대해
$i^*(\mathcal J/\mathcal J\mathcal K)\otimes\mathcal M^{\otimes n}$은
$Z$ 위의 단면들로 생성된다 (II, 4.5.5). 또한
$\mathcal J/\mathcal J\mathcal K
=i_*(i^*(\mathcal J/\mathcal J\mathcal K))$이므로,
(0\textsubscript{I}, 5.4.10)에 의해 충분히 큰 어떤 $n$에 대해
$X$ 위의
$(\mathcal J/\mathcal J\mathcal K)\otimes\mathcal L^{\otimes n}$의
단면 $s$가 존재하여 $s(x)\neq0$이다. 실제로 $\mathcal J$의 정의에
의해 $\mathcal J_x=\mathcal A_x$이고, 가정에 의해
$\mathcal K_x\neq\mathcal A_x$이다. 도식
(\hyperref[III.2.6.2.6-fr]{2.6.2.6})에 따르면 $s$는 또한 $X$ 위의
\[
(\mathcal J\mathcal B/\mathcal J\mathcal K\mathcal B)
\otimes\mathcal L^{\otimes n}
\]
의 단면이다. 따라서 $s$는 $X$ 위의
$(\mathcal J\mathcal B)\otimes\mathcal L^{\otimes n}$의 어떤 단면
$t$의 표준 상이다. 그런데 정의에 의해
$(\mathcal J\mathcal K\mathcal B)\otimes\mathcal L^{\otimes n}$을
법으로 한 $t$의 표준 상 $s$는
\[
(\mathcal J\otimes\mathcal L^{\otimes n})/
(\mathcal J\mathcal K\mathcal B\otimes\mathcal L^{\otimes n})
\]
에 속한다. 따라서 (\hyperref[III.2.6.2.6-fr]{2.6.2.6})에 의해
$t$는 $X$ 위의 $\mathcal J\otimes\mathcal L^{\otimes n}$의 단면이고,
더구나 $\mathcal L^{\otimes n}$의 단면이다. 위에서 $t(x)\neq0$임을
보았으므로 $x\in X_t$이다. 또한 $\mathcal J$의 정의에 의해
$\mathcal O_X/\mathcal J$의 받침인 $X-U$에서 $t(y)=0$이다. 따라서
$X_t\subset U$이고, 이로써 증명이 끝난다.
\end{env}
\end{proof}

\oldpage[III]{115}
\begin{env}[2.6.3]
\label{III.2.6.3-fr}
\emph{비고.} $X$가 $Y$ 위에서 고유이면, 슈발레 정리 (II, 6.7.1)에서와
같이 (\hyperref[III.2.6.1-fr]{2.6.1})과 보조정리 (II, 6.7.1.1)을
사용하여 (\hyperref[III.2.6.2-fr]{2.6.2})를 더 간단히 증명할 수 있다.
\end{env}

\section{고유 사상의 유한성 정리}
\label{section:III.3-fr}

\subsection{데비사주 보조정리}
\label{subsection:III.3.1-fr}

\begin{definition}[3.1.1]
\label{III.3.1.1-fr}
$K$를 아벨 범주라 하자. $K$의 대상들의 집합의 부분집합 $K'$이
$0\in K'$을 만족하고, $K$ 안의 모든 완전열
\[
0\longrightarrow A'\longrightarrow A\longrightarrow A''\longrightarrow0
\]
에서 $A$, $A'$, $A''$ 가운데 둘이 $K'$에 속하면 나머지 하나도
$K'$에 속할 때, $K'$을 \emph{완전}하다고 한다.
\end{definition}

\begin{theorem}[3.1.2]
\label{III.3.1.2-fr}
$X$를 뇌터 준스킴이라 하고, $K$로 연접 $\mathcal O_X$-가군들의 아벨
범주를 나타내자. $K'$을 $K$의 완전한 부분집합, $X'$을 $X$의
바탕공간의 닫힌 부분집합이라 하자. 일반점이 $y$인 $X'$의 모든 기약
닫힌 부분집합 $Y$에 대하여, 받침이 $Y$이고 $\mathcal G_y$가 차원
$1$인 $k(y)$-벡터공간인 어떤 $\mathcal O_X$-가군
$\mathcal G\in K'$이 존재한다고 가정하자. 그러면 받침이 $X'$에
포함되는 모든 연접 $\mathcal O_X$-가군은 $K'$에 속한다. 특히
$X'=X$이면 $K'=K$이다.
\end{theorem}

\begin{proof}
$X'$의 닫힌 부분집합 $Y$에 대한 다음 성질을 $P(Y)$라 하자. 받침이
$Y$에 포함되는 모든 연접 $\mathcal O_X$-가군이 $K'$에 속한다.
뇌터 귀납법의 원리 (0\textsubscript{I}, 2.2.2)에 의해 다음을
증명하면 충분하다. $Y$가 $X'$의 닫힌 부분집합이고, $Y$와 다른
$Y$의 모든 닫힌 부분집합 $Y'$에 대해 $P(Y')$이 참이면 $P(Y)$도
참이다.

따라서 받침이 $Y$에 포함되는 $\mathcal F\in K$를 택하고
$\mathcal F\in K'$임을 증명하자. 계속해서 $Y$로 바탕공간이 $Y$인
$X$의 축소 닫힌 부분준스킴을 나타내자 (I, 5.2.1). 이는
$\mathcal O_X$의 어떤 연접 아이디얼 $\mathcal J$로 정의된다.
$\mathcal J^n\mathcal F=0$이 되도록 하는 어떤 정수 $n>0$이 존재한다
(I, 9.3.4). 따라서 $1\leq k\leq n$에 대해 연접
$\mathcal O_X$-가군들의 완전열
\[
0\longrightarrow
\mathcal J^{k-1}\mathcal F/\mathcal J^k\mathcal F
\longrightarrow \mathcal F/\mathcal J^k\mathcal F
\longrightarrow \mathcal F/\mathcal J^{k-1}\mathcal F
\longrightarrow0
\]
이 있다 (0\textsubscript{I}, 5.3.6 및 5.3.3). $K'$이 완전하므로
$k$에 대한 귀납법에 의해 각
$\mathcal F_k=\mathcal J^{k-1}\mathcal F/\mathcal J^k\mathcal F$가
$K'$에 속함을 보이면 충분하다. 따라서 추가로
$\mathcal J\mathcal F=0$이라고 가정하여 $\mathcal F\in K'$임을
증명하는 문제로 귀착된다. 이는 $j:Y\to X$가 표준 포함일 때
$\mathcal F=j_*(j^*(\mathcal F))$라고 말하는 것과 같다. 이제 두
경우로 나눈다.
\begin{enumerate}
\item[a)] $Y$가 \emph{가약}인 경우. $Y=Y'\cup Y''$라 하자. 여기서
$Y'$, $Y''$은 $Y$와 다른 $Y$의 닫힌 부분집합들이다. 계속해서
$Y'$, $Y''$로 각각 바탕공간이 $Y'$, $Y''$인 $X$의 축소 닫힌
부분준스킴들도 나타내고, 이를 각각 정의하는 $\mathcal O_X$의 연접
아이디얼들을 $\mathcal J'$, $\mathcal J''$라 하자. 다음과 같이 놓자.
\[
\mathcal F'=\mathcal F\otimes_{\mathcal O_X}
(\mathcal O_X/\mathcal J'),\qquad
\mathcal F''=\mathcal F\otimes_{\mathcal O_X}
(\mathcal O_X/\mathcal J'').
\]
표준 준동형 $\mathcal F\to\mathcal F'$,
$\mathcal F\to\mathcal F''$은 준동형
$u:\mathcal F\to\mathcal F'\oplus\mathcal F''$를 정의한다.
$z\notin Y'\cap Y''$이면 준동형
$u_z:\mathcal F_z\to\mathcal F'_z\oplus\mathcal F''_z$가
\emph{전단사}임을 보이자. 실제로
$\mathcal J'\cap\mathcal J''=\mathcal J$이다. 문제는 국소적이고

\oldpage[III]{116}
이 등식은 (I, 5.2.1 및 1.1.5)에서 따르기 때문이다. $z\notin Y''$이면
$\mathcal J'_z=\mathcal J_z$이고, 따라서
$\mathcal F'_z=\mathcal F_z$, $\mathcal F''_z=0$이다. 이는 이 경우의
주장을 증명하며, $z\notin Y'$인 경우도 마찬가지이다. 그러므로 $K$에
속하는 $u$의 핵과 여핵은 (0\textsubscript{I}, 5.3.4) 받침이
$Y'\cap Y''$에 포함되어 귀납가정에 의해 $K'$에 속한다. 같은 이유로
$\mathcal F'$, $\mathcal F''$도 $K'$에 속한다. $K'$이 완전하므로
$\mathcal F'\oplus\mathcal F''$도 $K'$에 속한다. 이제 두 완전열
\[
0\longrightarrow\operatorname{Im}u\longrightarrow
\mathcal F'\oplus\mathcal F''\longrightarrow
\operatorname{Coker}u\longrightarrow0,
\]
\[
0\longrightarrow\operatorname{Ker}u\longrightarrow\mathcal F
\longrightarrow\operatorname{Im}u\longrightarrow0
\]
과 $K'$의 완전성에서 결론이 따른다.

\item[b)] $Y$가 기약인 경우. 따라서 $X$의 부분준스킴 $Y$는
\emph{정역}이다. 그 일반점을 $y$라 하면
$(\mathcal O_Y)_y=k(y)$이다. $j^*(\mathcal F)$가 연접
$\mathcal O_Y$-가군이므로
$\mathcal F_y=(j^*(\mathcal F))_y$는 어떤 유한 차원 $m$의
$k(y)$-벡터공간이다. 가정에 의해 어떤 연접 $\mathcal O_X$-가군
$\mathcal G$가 존재하여($\mathcal G$의 받침은 반드시 $Y$이다),
$\mathcal G_y$는 차원 $1$인 $k(y)$-벡터공간이다. 따라서
$k(y)$-동형사상
$(\mathcal G_y)^m\xrightarrow{\sim}\mathcal F_y$가 존재하며, 이는
$\mathcal O_Y$-동형사상이기도 하다. $\mathcal G^m$과 $\mathcal F$가
연접이므로, $X$에서 $y$의 어떤 열린 근방 $W$와 동형사상
$\mathcal G^m|W\xrightarrow{\sim}\mathcal F|W$가 존재한다
(0\textsubscript{I}, 5.2.7). 이 동형사상의 그래프를 $\mathcal H$라
하자. 이는 $(\mathcal G^m\oplus\mathcal F)|W$의 연접 부분
$(\mathcal O_X|W)$-가군이고, 표준적으로 $\mathcal G^m|W$ 및
$\mathcal F|W$와 동형이다. 따라서 $\mathcal G^m\oplus\mathcal F$의
어떤 연접 부분 $\mathcal O_X$-가군 $\mathcal H_0$가 존재하여,
$W$ 위에서는 $\mathcal H$를 유도하고 $X-Y$ 위에서는 0을 유도한다.
이는 $\mathcal G^m$과 $\mathcal F$의 받침이 $Y$이기 때문이다
(I, 9.4.7).

$\mathcal G^m\oplus\mathcal F$의 표준 사영들을 제한하여 얻는
$v:\mathcal H_0\to\mathcal G^m$과
$w:\mathcal H_0\to\mathcal F$는 연접 $\mathcal O_X$-가군들의
준동형이고, $W$ 및 $X-Y$에서는 동형사상이다. 다시 말해 $v$, $w$의
핵과 여핵들의 받침은 $Y$와 다른 닫힌집합 $Y-(Y\cap W)$에 포함된다.
따라서 이들은 $K'$에 속한다. 한편 $\mathcal G\in K'$이고 $K'$이
완전하므로 $\mathcal G^m\in K'$이다. $K'$의 완전성에 의해 차례로
$\mathcal H_0\in K'$, 이어서 $\mathcal F\in K'$임을 얻는다. 이로써
증명이 끝난다.
\end{enumerate}
\end{proof}

\begin{corollary}[3.1.3]
\label{III.3.1.3-fr}
$K$의 완전한 부분집합 $K'$이 다음 성질도 만족한다고 가정하자.
$\mathcal M\in K'$인 연접 $\mathcal O_X$-가군의 모든 연접 직합인자는
다시 $K'$에 속한다. 이 조건 아래에서
(\hyperref[III.3.1.2-fr]{3.1.2})의 조건
``$\mathcal G_y$는 차원 $1$인 $k(y)$-벡터공간이다''를
$\mathcal G_y\neq0$으로 바꾸어도 같은 결론이 성립한다. 마지막 조건은
$\operatorname{Supp}(\mathcal G)=Y$와 동치이다.
\end{corollary}

\begin{proof}
실제로 (\hyperref[III.3.1.2-fr]{3.1.2})의 논증은 경우 b)에서만
바꾸면 된다. 이번에는 $\mathcal G_y$가 어떤 차원 $q>0$의
$k(y)$-벡터공간이므로
$(\mathcal G_y)^m\xrightarrow{\sim}(\mathcal F_y)^q$인
$\mathcal O_Y$-동형사상이 존재한다. 그러면
(\hyperref[III.3.1.2-fr]{3.1.2})의 논증 끝부분은
$\mathcal F^q\in K'$임을 증명하고, $K'$에 대한 추가 가정으로부터
$\mathcal F\in K'$이 따른다.
\end{proof}

\subsection{유한성 정리: 통상적 스킴의 경우}
\label{subsection:III.3.2-fr}

\begin{theorem}[3.2.1]
\label{III.3.2.1-fr}
$Y$를 국소 뇌터 준스킴, $f:X\to Y$를 고유 사상이라 하자. 모든 연접
$\mathcal O_X$-가군 $\mathcal F$에 대해 $q\geq0$인
$\mathcal O_Y$-가군 $R^qf_*(\mathcal F)$들은 연접이다.
\end{theorem}

\begin{proof}
문제는 $Y$ 위에서 국소적이므로 $Y$가 뇌터라고 가정할 수 있고,
따라서 $X$도 뇌터이다 (I, 6.3.7). 정리
(\hyperref[III.3.2.1-fr]{3.2.1})의 결론이 성립하는 연접
$\mathcal O_X$-가군 $\mathcal F$들은 연접 $\mathcal O_X$-가군들의
범주 $K$의 완전한 부분집합 $K'$을 이룬다.

\oldpage[III]{117}
실제로
$0\to\mathcal F'\to\mathcal F\to\mathcal F''\to0$을 연접
$\mathcal O_X$-가군들의 완전열이라 하자. 예를 들어
$\mathcal F'$, $\mathcal F''$이 $K'$에 속한다고 가정하면 코호몰로지
완전열
\[
R^{q-1}f_*(\mathcal F'')\xrightarrow{\partial}R^qf_*(\mathcal F')
\longrightarrow R^qf_*(\mathcal F)\longrightarrow R^qf_*(\mathcal F'')
\xrightarrow{\partial}R^{q+1}f_*(\mathcal F')
\]
에서 가정에 의해 바깥 네 항이 연접이다. 따라서 가운데 항
$R^qf_*(\mathcal F)$도 연접이다
(0\textsubscript{I}, 5.3.4 및 5.3.3). 같은 방식으로
$\mathcal F$, $\mathcal F'$ (각각 $\mathcal F$, $\mathcal F''$)이
$K'$에 속하면 $\mathcal F''$ (각각 $\mathcal F'$)도 $K'$에 속함을
보인다. 더 나아가 $\mathcal F\in K'$인 $\mathcal O_X$-가군의 모든
연접 \emph{직합인자} $\mathcal F'$도 $K'$에 속한다. 실제로 이때
$R^qf_*(\mathcal F')$은 $R^qf_*(\mathcal F)$의 직합인자이고
(G, II, 4.4.4), 따라서 유한 생성이다. 또한
(\hyperref[III.1.4.10-fr]{1.4.10})에 의해 준연접이고 $Y$가
뇌터이므로 연접이다.

(\hyperref[III.3.1.3-fr]{3.1.3})에 의해 다음을 증명하는 문제로
귀착된다. $X$가 일반점 $x$를 갖는 \emph{기약} 준스킴이면,
$\mathcal F_x\neq0$이고 $K'$에 속하는 어떤 연접
$\mathcal O_X$-가군 $\mathcal F$가 존재한다. 실제로 이 사실이
증명되면 $X$의 모든 기약 닫힌 부분준스킴 $Y$에 적용할 수 있다.
$j:Y\to X$가 표준 포함이면 $f\circ j$는 고유이고 (II, 5.4.2),
$\mathcal G$가 받침이 $Y$인 연접 $\mathcal O_Y$-가군일 때
$j_*(\mathcal G)$는 연접 $\mathcal O_X$-가군이며
\[
R^q(f\circ j)_*(\mathcal G)=R^qf_*(j_*(\mathcal G))
\]
이다 (G, II, 4.9.1). 따라서
(\hyperref[III.3.1.3-fr]{3.1.3})을 적용할 조건을 만족한다.

이제 차우 보조정리 (II, 5.6.2)에 의해 어떤 기약 준스킴 $X'$과
\emph{사영 전사} 사상 $g:X'\to X$가 존재하여
$f\circ g:X'\to Y$는 \emph{사영}이다. $g$에 대해 풍부한 어떤
$\mathcal O_{X'}$-가군 $\mathcal L$이 존재한다 (II, 5.3.1).
$g:X'\to X$와 $\mathcal L$에 사영 사상의 기본정리
(\hyperref[III.2.2.1-fr]{2.2.1})을 적용하자. 그러면 어떤 정수 $n$이
존재하여
\[
\mathcal F=g_*(\mathcal O_{X'}(n))
\]
은 연접 $\mathcal O_X$-가군이고 모든 $q>0$에 대해
$R^qg_*(\mathcal O_{X'}(n))=0$이다. 또한 충분히 큰 $n$에 대해
\[
g^*(g_*(\mathcal O_{X'}(n)))\longrightarrow\mathcal O_{X'}(n)
\]
이 전사이므로 (\hyperref[III.2.2.1-fr]{2.2.1}), $X$의 일반점 $x$에서
$\mathcal F_x\neq0$이라고 가정할 수 있다 (II, 3.4.7). 한편
$f\circ g$가 사영이고 $Y$가 뇌터이므로
$R^p(f\circ g)_*(\mathcal O_{X'}(n))$들은 \emph{연접}이다
(\hyperref[III.2.2.1-fr]{2.2.1}). 이제
$R^*(f\circ g)_*(\mathcal O_{X'}(n))$은 $E_2$항이
\[
E_2^{pq}=R^pf_*(R^qg_*(\mathcal O_{X'}(n)))
\]
인 르레 스펙트럼 열의 수렴 대상이다. 앞의 결과에 의해 이 스펙트럼
열은 퇴화하고, 따라서 (0, 11.1.6)에 의해
$E_2^{p0}=R^pf_*(\mathcal F)$는
$R^p(f\circ g)_*(\mathcal O_{X'}(n))$과 동형이다. 이로써 증명이
끝난다.
\end{proof}

\begin{corollary}[3.2.2]
\label{III.3.2.2-fr}
$Y$를 국소 뇌터 준스킴이라 하자. 모든 고유 사상 $f:X\to Y$에
대하여 연접 $\mathcal O_X$-가군의 $f$에 의한 직접상은 연접
$\mathcal O_Y$-가군이다.
\end{corollary}

\begin{corollary}[3.2.3]
\label{III.3.2.3-fr}
$A$를 뇌터 환, $X$를 $A$ 위의 고유 스킴이라 하자. 모든 연접
$\mathcal O_X$-가군 $\mathcal F$에 대해 $H^p(X,\mathcal F)$는 유한
생성 $A$-가군이다. 또한 어떤 정수 $r>0$이 존재하여 모든 연접
$\mathcal O_X$-가군 $\mathcal F$와 모든 $p>r$에 대해
$H^p(X,\mathcal F)=0$이다.
\end{corollary}

\begin{proof}
둘째 명제는 이미 증명하였다
(\hyperref[III.1.4.12-fr]{1.4.12}). 첫째 명제는
(\hyperref[III.1.4.11-fr]{1.4.11})을 고려하면 유한성 정리
(\hyperref[III.3.2.1-fr]{3.2.1})에서 따른다.
\end{proof}

특히 $X$가 체 $k$ 위의 \emph{고유 대수적 스킴}이면, 모든 연접
$\mathcal O_X$-가군 $\mathcal F$에 대해 $H^p(X,\mathcal F)$들은
\emph{유한 차원} $k$-벡터공간이다.

\begin{corollary}[3.2.4]
\label{III.3.2.4-fr}
$Y$를 국소 뇌터 준스킴, $f:X\to Y$를 유한형 사상이라 하자. 받침이
$Y$ 위에서 고유인 모든 연접 $\mathcal O_X$-가군 $\mathcal F$에 대해
(II, 5.4.10), $\mathcal O_Y$-가군 $R^qf_*(\mathcal F)$들은 연접이다.
\end{corollary}

\oldpage[III]{118}
\begin{proof}
문제는 $Y$ 위에서 국소적이므로 $Y$가 뇌터라고 가정할 수 있고,
따라서 $X$도 뇌터이다 (I, 6.3.7). 가정에 의해 바탕공간이
$\operatorname{Supp}(\mathcal F)$인 $X$의 모든 닫힌 부분준스킴
$Z$는 $Y$ 위에서 고유이다. 다시 말해 $j:Z\to X$가 표준 포함이면
$f\circ j:Z\to Y$는 고유이다. 한편
$\mathcal F=j_*(\mathcal G)$가 되도록 $Z$를 택할 수 있으며, 여기서
$\mathcal G=j^*(\mathcal F)$는 연접 $\mathcal O_Z$-가군이다
(I, 9.3.5). (\hyperref[III.1.3.4-fr]{1.3.4})에 의해
$R^qf_*(\mathcal F)=R^q(f\circ j)_*(\mathcal G)$이므로, 결론은
(\hyperref[III.3.2.1-fr]{3.2.1})에서 즉시 따른다.
\end{proof}

\subsection{유한성 정리의 일반화 (통상적 스킴)}
\label{subsection:III.3.3-fr}

\begin{proposition}[3.3.1]
\label{III.3.3.1-fr}
$Y$를 뇌터 준스킴, $\mathcal S$를 양의 차수들을 갖는 준연접 유한
생성 등급 $\mathcal O_Y$-대수,
$Y'=\operatorname{Proj}(\mathcal S)$라 하고, $g:Y'\to Y$를 구조
사상이라 하자. $f:X\to Y$를 고유 사상,
$\mathcal S'=f^*(\mathcal S)$,
$\mathcal M=\bigoplus_{k\in\mathbb Z}\mathcal M_k$를 유한 생성 준연접
등급 $\mathcal S'$-가군이라 하자. 그러면 모든 $p$에 대해
\[
R^pf_*(\mathcal M)=\bigoplus_{k\in\mathbb Z}R^pf_*(\mathcal M_k)
\]
는 유한 생성 등급 $\mathcal S$-가군이다. 더 나아가 $\mathcal S$가
$\mathcal S_1$로 생성된다고 가정하자. 그러면 각 $p\in\mathbb Z$에
대해 어떤 정수 $k_p$가 존재하여, 모든 $k\geq k_p$, $r\geq0$에 대해
\[
R^pf_*(\mathcal M_{k+r})=\mathcal S_rR^pf_*(\mathcal M_k)
\tag{3.3.1.1}\label{III.3.3.1.1-fr}
\]
이다.
\end{proposition}

\begin{proof}
첫째 명제는 (\hyperref[III.2.4.1-fr]{2.4.1}, (i))의 명제에서 단지
``사영 사상''을 ``고유 사상''으로 바꾼 것이다. 그런데
(\hyperref[III.2.4.1-fr]{2.4.1}, (i))의 증명에서 $f$에 대한 가정은
그 증명의 표기로 $R^pf'_*(\widetilde{\mathcal M})$이 연접
$\mathcal O_{Y'}$-가군임을 보이는 데에만 사용되었다.
(\hyperref[III.3.3.1-fr]{3.3.1})의 가정 아래에서는 $f'$이 고유이므로
(II, 5.4.2, (iii)), 유한성 정리
(\hyperref[III.3.2.1-fr]{3.2.1})에 의해
(\hyperref[III.2.4.1-fr]{2.4.1}, (i))의 증명을 그대로 반복할 수 있다.

둘째 명제에 대해서는, $Y$의 어떤 유한 아핀 열린 덮개 $(U_i)$가
존재하여 모든 $k\geq k_{p,i}$에 대해
(\hyperref[III.3.3.1.1-fr]{3.3.1.1})의 두 변을 $U_i$에 제한한 것들이
같음을 유의하면 충분하다 (II, 2.1.6, (ii)). $k_{p,i}$들 가운데 가장
큰 것을 $k_p$로 택하면 된다.
\end{proof}

\begin{corollary}[3.3.2]
\label{III.3.3.2-fr}
$A$를 뇌터 환, $\mathfrak m$을 $A$의 아이디얼, $X$를 고유
$A$-스킴, $\mathcal F$를 연접 $\mathcal O_X$-가군이라 하자. 그러면
모든 $p\geq0$에 대해 직합
\[
\bigoplus_{k\geq0}H^p(X,\mathfrak m^k\mathcal F)
\]
은 환 $S=\bigoplus_{k\geq0}\mathfrak m^k$ 위의 유한 생성 가군이다.
특히 어떤 정수 $k_p\geq0$이 존재하여 모든 $k\geq k_p$, $r\geq0$에
대해
\[
H^p(X,\mathfrak m^{k+r}\mathcal F)
=\mathfrak m^rH^p(X,\mathfrak m^k\mathcal F)
\tag{3.3.2.1}\label{III.3.3.2.1-fr}
\]
이다.
\end{corollary}

\begin{proof}
$Y=\operatorname{Spec}(A)$, $\mathcal S=\widetilde S$,
$\mathcal M_k=\mathfrak m^k\mathcal F$로 놓고
(\hyperref[III.1.4.11-fr]{1.4.11})을 고려하여
(\hyperref[III.3.3.1-fr]{3.3.1})을 적용하면 된다.
\end{proof}

$\bigoplus_{k\geq0}H^p(X,\mathfrak m^k\mathcal F)$의 $S$-가군 구조는
다음과 같이 얻어짐을 상기하자. 모든 $a\in\mathfrak m^r$에 대해
곱셈
$\mathfrak m^k\mathcal F\to\mathfrak m^{k+r}\mathcal F$에서
코호몰로지로 넘어가 얻는 사상
\[
H^p(X,\mathfrak m^k\mathcal F)
\longrightarrow H^p(X,\mathfrak m^{k+r}\mathcal F)
\]
을 생각한다 (\hyperref[III.2.4.1-fr]{2.4.1}).

\oldpage[III]{119}
\subsection{유한성 정리: 형식적 스킴의 경우}
\label{subsection:III.3.4-fr}

이 절의 결과들 가운데 정의
(\hyperref[III.3.4.1-fr]{3.4.1})를 제외한 것들은 이 장의 나머지
부분에서 사용하지 않는다.

\begin{env}[3.4.1]
\label{III.3.4.1-fr}
$\mathfrak X$와 $\mathfrak S$를 국소 뇌터 형식적 준스킴들
(I, 10.4.2), $f:\mathfrak X\to\mathfrak S$를 형식적 준스킴의 사상이라
하자. 다음 조건들을 만족할 때 $f$를 \emph{고유 사상}이라고 한다.
\begin{enumerate}
\item[$1^\circ$] $f$는 유한형 사상이다 (I, 10.13.3).
\item[$2^\circ$] $\mathcal K$가 $\mathfrak S$의 정의 아이디얼이고
\[
\mathcal J=f^*(\mathcal K)\mathcal O_{\mathfrak X},\qquad
X_0=(\mathfrak X,\mathcal O_{\mathfrak X}/\mathcal J),\qquad
S_0=(\mathfrak S,\mathcal O_{\mathfrak S}/\mathcal K)
\]
로 놓으면, $f$에서 유도되는 사상 $f_0:X_0\to S_0$ (I, 10.5.6)는
고유이다.
\end{enumerate}

이 정의가 고려한 $\mathfrak S$의 정의 아이디얼 $\mathcal K$에 의존하지
않음은 즉시 알 수 있다. 실제로 $\mathcal K'\subset\mathcal K$인 둘째
정의 아이디얼 $\mathcal K'$에 대하여
$\mathcal J'=f^*(\mathcal K')\mathcal O_{\mathfrak X}$,
$X'_0=(\mathfrak X,\mathcal O_{\mathfrak X}/\mathcal J')$,
$S'_0=(\mathfrak S,\mathcal O_{\mathfrak S}/\mathcal K')$로 놓으면,
$f$에서 유도되는 사상 $f'_0:X'_0\to S'_0$에 대한 도표
\[
\xymatrix{
X_0\ar[r]^{f_0}\ar[d]_i & S_0\ar[d]^j\\
X'_0\ar[r]_{f'_0} & S'_0
}
\]
는 가환이고 $i$와 $j$는 전사인 몰입 사상이다. 따라서 (II, 5.4.5)에
의해 $f_0$가 고유라는 것과 $f'_0$가 고유라는 것은 동치이다.

모든 $n\geq0$에 대하여
$X_n=(\mathfrak X,\mathcal O_{\mathfrak X}/\mathcal J^{n+1})$,
$S_n=(\mathfrak S,\mathcal O_{\mathfrak S}/\mathcal K^{n+1})$로 놓자.
$f$에서 유도되는 사상 $f_n:X_n\to S_n$ (I, 10.5.6)는 $n=0$일 때
고유이면 모든 $n$에 대해 고유이다 (II, 5.4.6).

$g:Y\to Z$가 국소 뇌터인 통상적 준스킴들의 고유 사상이고, $Z'$가
$Z$의 닫힌 부분집합, $Y'$가 $g(Y')\subset Z'$를 만족하는 $Y$의 닫힌
부분집합이면, $g$를 완비화들로 연장한
$\widehat g:Y_{/Y'}\to Z_{/Z'}$ (I, 10.9.1)는 정의와 (II, 5.4.5)에
의해 형식적 준스킴들의 고유 사상이다.

$\mathfrak X$와 $\mathfrak S$를 국소 뇌터 형식적 준스킴들,
$f:\mathfrak X\to\mathfrak S$를 유한형 사상이라 하자 (I, 10.13.3).
표기는 위와 같이 하고, $\mathfrak X$의 바탕공간의 부분집합 $Z$를
$X_0$의 부분집합으로 보았을 때 $S_0$ 위에서 고유이면
(II, 5.4.10), $Z$가 \emph{$\mathfrak S$ 위에서 고유하다} (또는
$f$에 대해 고유하다)고 한다. 통상적 준스킴의 고유한 부분집합에 관하여
(II, 5.4.10)에서 서술한 모든 성질은 정의에서 즉시 알 수 있듯이 형식적
준스킴의 고유한 부분집합에도 그대로 성립한다.
\end{env}

\begin{theorem}[3.4.2]
\label{III.3.4.2-fr}
$\mathfrak X$, $\mathfrak Y$를 국소 뇌터 형식적 준스킴들,
$f:\mathfrak X\to\mathfrak Y$를 고유 사상이라 하자. 모든 연접
$\mathcal O_{\mathfrak X}$-가군 $\mathcal F$에 대하여
$\mathcal O_{\mathfrak Y}$-가군 $R^qf_*(\mathcal F)$들은 모든
$q\geq0$에 대해 연접이다.
\end{theorem}

$\mathcal J$를 $\mathfrak Y$의 정의 아이디얼,
$\mathcal K=f^*(\mathcal J)\mathcal O_{\mathfrak X}$라 하고,
$\mathcal O_{\mathfrak X}$-가군
\[
\mathcal F_k=\mathcal F\otimes_{\mathcal O_{\mathfrak Y}}
(\mathcal O_{\mathfrak Y}/\mathcal J^{k+1})
=\mathcal F/\mathcal K^{k+1}\mathcal F\qquad(k\geq0)
\tag{3.4.2.1}\label{III.3.4.2.1-fr}
\]
들을 생각하자. 이들은 명백히 위상 $\mathcal O_{\mathfrak X}$-가군의
사영계를 이루고,

\oldpage[III]{120}
\[
\mathcal F=\varprojlim_k\mathcal F_k
\]
이다 (I, 10.11.3). 한편
(\hyperref[III.3.4.2-fr]{3.4.2})에 의해 각
$R^qf_*(\mathcal F)$는 연접이므로 자연스럽게 위상
$\mathcal O_{\mathfrak Y}$-가군의 구조를 가지며 (I, 10.11.6),
$R^qf_*(\mathcal F_k)$도 마찬가지이다. 표준 준동형사상
$\mathcal F\to\mathcal F_k=\mathcal F/\mathcal K^{k+1}\mathcal F$에는
표준 준동형사상
\[
R^qf_*(\mathcal F)\longrightarrow R^qf_*(\mathcal F_k)
\]
들이 대응한다. 이들은 앞의 위상 $\mathcal O_{\mathfrak Y}$-가군 구조에
관하여 반드시 연속이고 (I, 10.11.6), 사영계를 이룬다. 따라서 극한을
취하면 함자적인 표준 준동형사상
\[
R^qf_*(\mathcal F)\longrightarrow\varprojlim_kR^qf_*(\mathcal F_k)
\tag{3.4.2.2}\label{III.3.4.2.2-fr}
\]
을 얻으며, 이는 위상 $\mathcal O_{\mathfrak Y}$-가군의 연속
준동형사상이다. 우리는 (\hyperref[III.3.4.2-fr]{3.4.2})와 함께 다음을
증명할 것이다.

\begin{corollary}[3.4.3]
\label{III.3.4.3-fr}
각 준동형사상 (\hyperref[III.3.4.2.2-fr]{3.4.2.2})는 위상
동형사상이다. 더 나아가 $\mathfrak Y$가 뇌터이면 사영계
$(R^qf_*(\mathcal F/\mathcal K^{k+1}\mathcal F))_{k\geq0}$는
조건 (ML)을 만족한다 (0, 13.1.1).
\end{corollary}

먼저 $\mathfrak Y$가 뇌터 아핀 형식적 스킴인 경우 (I, 10.4.1)에
(\hyperref[III.3.4.2-fr]{3.4.2})와
(\hyperref[III.3.4.3-fr]{3.4.3})를 증명한다.

\begin{corollary}[3.4.4]
\label{III.3.4.4-fr}
(\hyperref[III.3.4.2-fr]{3.4.2})의 가정들에 더하여
$\mathfrak Y=\operatorname{Spf}(A)$이고 $A$가 뇌터 아딕 환이라고
가정하자. $\mathfrak J$를 $A$의 정의 아이디얼이라 하고 $k\geq0$에
대해 $\mathcal F_k=\mathcal F/\mathfrak J^{k+1}\mathcal F$로 놓자.
그러면 $H^n(\mathfrak X,\mathcal F)$들은 유한 생성 $A$-가군이고,
사영계 $(H^n(\mathfrak X,\mathcal F_k))_{k\geq0}$는 모든 $n$에 대해
조건 (ML)을 만족한다. 또한
\[
N_{n,k}=\operatorname{Ker}\bigl(H^n(\mathfrak X,\mathcal F)
\longrightarrow H^n(\mathfrak X,\mathcal F_k)\bigr)
\tag{3.4.4.1}\label{III.3.4.4.1-fr}
\]
로 놓으면, 이는 코호몰로지 완전열에 의해
$\operatorname{Im}(H^n(\mathfrak X,\mathfrak J^{k+1}\mathcal F)
\to H^n(\mathfrak X,\mathcal F))$와도 같다. $N_{n,k}$들은
$H^n(\mathfrak X,\mathcal F)$ 위에 $\mathfrak J$-좋은 여과를 정의한다
(0, 13.7.7). 끝으로 모든 $n$에 대해 표준 준동형사상
\[
H^n(\mathfrak X,\mathcal F)
\longrightarrow\varprojlim_kH^n(\mathfrak X,\mathcal F_k)
\tag{3.4.4.2}\label{III.3.4.4.2-fr}
\]
은 위상 동형사상이다. 여기서 왼쪽에는 $\mathfrak J$-아딕 위상을,
$H^n(\mathfrak X,\mathcal F_k)$들에는 이산 위상을 준다.
\end{corollary}

다음과 같이 놓자.
\[
S=\operatorname{gr}(A)=\bigoplus_{k\geq0}
\mathfrak J^k/\mathfrak J^{k+1},\qquad
\mathcal M=\operatorname{gr}(\mathcal F)=\bigoplus_{k\geq0}
\mathfrak J^k\mathcal F/\mathfrak J^{k+1}\mathcal F.
\tag{3.4.4.3}\label{III.3.4.4.3-fr}
\]
$\mathfrak J^\Delta$가 $\mathfrak Y$의 정의 아이디얼임은 알려져 있다
(I, 10.3.1). 이제
\[
\mathcal K=f^*(\mathfrak J^\Delta)\mathcal O_{\mathfrak X},\qquad
X_0=(\mathfrak X,\mathcal O_{\mathfrak X}/\mathcal K),\qquad
Y_0=(\mathfrak Y,\mathcal O_{\mathfrak Y}/\mathfrak J^\Delta)
=\operatorname{Spec}(A_0)
\]
라 하자. 여기서 $A_0=A/\mathfrak J$이다. 가군
$\mathcal M_k=\mathfrak J^k\mathcal F/\mathfrak J^{k+1}\mathcal F$들이
연접 $\mathcal O_{X_0}$-가군임은 명백하다 (I, 10.11.3). 한편 준연접
등급 $\mathcal O_{X_0}$-대수
\[
\mathcal S=\mathcal O_{X_0}\otimes_{A_0}S
\tag{3.4.4.4}\label{III.3.4.4.4-fr}
\]
를 생각하자. 곱셈은 등급 대수의 전사 준동형사상
\[
\mathcal S\longrightarrow
\operatorname{gr}(\mathcal O_{\mathfrak X})
=\bigoplus_{k\geq0}\mathcal K^k/\mathcal K^{k+1}
\]
을 정의하며, 이를 통해 $\mathcal M$은 등급 $\mathcal S$-가군이 된다.

$\mathcal F$가 유한 생성 $\mathcal O_{\mathfrak X}$-가군이라는 가정에서
먼저 $\mathcal M$이 유한 생성 등급 $\mathcal S$-가군임이 따른다.

\oldpage[III]{121}
실제로 이 문제는 $\mathfrak X$ 위에서 국소적이므로, 이를 다룰 때
$\mathfrak X=\operatorname{Spf}(B)$라고 가정할 수 있다. 여기서 $B$는
뇌터 아딕 환이고 $\mathcal F=N^\Delta$이며, $N$은 유한 생성
$B$-가군이다 (I, 10.10.5). $B_0=B/\mathfrak J B$,
$N_0=N/\mathfrak JN$으로 놓자. 그러면
$X_0=\operatorname{Spec}(B_0)$이고, 준연접 $\mathcal O_{X_0}$-가군
$\mathcal S$와 $\mathcal M$은 각각
\[
S'=\bigoplus_{k\geq0}
((\mathfrak J^k/\mathfrak J^{k+1})\otimes_{A_0}B_0),\qquad
M'=\bigoplus_{k\geq0}\mathfrak J^kN/\mathfrak J^{k+1}N
\]
에 대응한다. 모든 $k\geq0$에 대해 곱셈은 전사 사상
\[
(\mathfrak J^k/\mathfrak J^{k+1})\otimes_{A_0}N_0
\longrightarrow\mathfrak J^kN/\mathfrak J^{k+1}N,
\qquad \overline a\otimes\overline n\longmapsto\overline{an}
\]
을 유도한다. 이 사상들은 등급 $S'$-가군의 전사 준동형사상
$S'\otimes_{B_0}N_0\to M'$을 정의한다. $N_0$는 유한 생성
$B_0$-가군이므로 $M'$는 유한 생성 $S'$-가군이며, 여기서 주장이 따른다
(I, 1.3.13).

사상 $f_0:X_0\to Y_0$는 가정에 의해 \emph{고유}이므로,
$\mathcal S$, $\mathcal M$, $f_0$에
(\hyperref[III.3.3.1-fr]{3.3.1})을 적용할 수 있다.
(\hyperref[III.1.4.11-fr]{1.4.11})도 고려하면 \emph{모든 $n\geq0$}에
대해 $\bigoplus_{k\geq0}H^n(X_0,\mathcal M_k)$가 유한 생성 등급
$S$-가군임을 얻는다. 이는 아벨 군의 층들로 이루어진 엄밀 사영계
$(\mathcal F/\mathfrak J^k\mathcal F)_{k\geq0}$에 각자의 자연스러운
‘여과된 $A$-가군’ 구조를 줄 때, (0, 13.7.7)의 조건
$(\mathrm F_n)$이 \emph{모든 $n\geq0$}에 대해 성립함을 증명한다.
따라서 (0, 13.7.7)을 적용하면 다음을 얻는다.
\begin{enumerate}
\item[$1^\circ$] 사영계
$(H^n(\mathfrak X,\mathcal F_k))_{k\geq0}$는 조건 (ML)을 만족한다.
\item[$2^\circ$]
$H'^{,n}=\varprojlim_kH^n(\mathfrak X,\mathcal F_k)$로 놓으면,
$H'^{,n}$은 유한 생성 $A$-가군이다.
\item[$3^\circ$] 표준 준동형사상
$H'^{,n}\to H^n(\mathfrak X,\mathcal F_k)$들의 핵들이 $H'^{,n}$ 위에
정의하는 여과는 $\mathfrak J$-좋은 여과이다.
\end{enumerate}

한편 $X_k=(\mathfrak X,\mathcal O_{\mathfrak X}/\mathcal K^{k+1})$로
놓으면, $\mathcal F_k$는 연접 $\mathcal O_{X_k}$-가군이다
(I, 10.11.3). $U$가 $X_0$의 아핀 열린집합이면 $U$는 각 $X_k$에서도
아핀 열린집합이므로 (I, 5.1.9), 모든 $n>0$과 모든 $k$에 대해
$H^n(U,\mathcal F_k)=0$이고
(\hyperref[III.1.3.1-fr]{1.3.1}), $h\leq k$일 때
$H^0(U,\mathcal F_k)\to H^0(U,\mathcal F_h)$는 전사이다 (I, 1.3.9).
그러므로 (0, 13.3.2)의 조건들이 성립하며, (0, 13.3.1)의 사상에 의해
$H'^{,n}$은
\[
H^n(\mathfrak X,\varprojlim_k\mathcal F_k)
=H^n(\mathfrak X,\mathcal F)
\]
와 표준적으로 동일시된다. 이것으로
(\hyperref[III.3.4.4-fr]{3.4.4})의 증명이 끝난다.

\begin{env}[3.4.5]
\label{III.3.4.5-fr}
이제 (\hyperref[III.3.4.2-fr]{3.4.2})와
(\hyperref[III.3.4.3-fr]{3.4.3})의 증명으로 돌아가자. 먼저
(\hyperref[III.3.4.4-fr]{3.4.4})에서 다룬
$\mathfrak Y=\operatorname{Spf}(A)$인 경우에 이 명제들을 증명한다.
이를 위해 모든 $g\in A$에 대하여 $\mathfrak Y$의 열린집합
$\mathfrak Y_g=\mathfrak D(g)$ 위에 유도되는 뇌터 아핀 형식적 스킴과
$f^{-1}(\mathfrak Y_g)$ 위에 유도되는 형식적 준스킴에
(\hyperref[III.3.4.4-fr]{3.4.4})를 적용하자. 여기서
$\mathfrak Y_g=\operatorname{Spf}(A_{\{g\}})$이다. $\mathfrak Y_g$는
준스킴
\[
Y_k=(\mathfrak Y,\mathcal O_{\mathfrak Y}/
(\mathfrak J^\Delta)^{k+1})
\]
에서도 아핀 열린집합이고, $\mathcal F_k$는 연접
$\mathcal O_{X_k}$-가군이므로
(\hyperref[III.1.4.11-fr]{1.4.11})에 의해 모든 $k\geq0$에 대해
\[
H^n(f^{-1}(\mathfrak Y_g),\mathcal F_k)
=\Gamma(\mathfrak Y_g,R^nf_*(\mathcal F_k))
\]
이다. 표준 준동형사상
\[
H^n(f^{-1}(\mathfrak Y_g),\mathcal F)
\longrightarrow\varprojlim_k\Gamma(\mathfrak Y_g,R^nf_*(\mathcal F_k))
\]
은 동형사상이다. 그런데 (0\textsubscript{I}, 3.2.6)에 의해
\[
\varprojlim_k\Gamma(\mathfrak Y_g,R^nf_*(\mathcal F_k))
=\Gamma(\mathfrak Y_g,\varprojlim_kR^nf_*(\mathcal F_k))
\]
이다.

\oldpage[III]{122}
또한 층 $R^nf_*(\mathcal F)$는 $\mathfrak Y_g$들 위의 준층
\[
\mathfrak Y_g\longmapsto
H^n(f^{-1}(\mathfrak Y_g),\mathcal F)
\]
에 대응하므로 (0\textsubscript{I}, 3.2.1), 준동형사상
(\hyperref[III.3.4.2.2-fr]{3.4.2.2})가 전단사임을 보였다. 이제
$R^nf_*(\mathcal F)$가 연접 $\mathcal O_{\mathfrak Y}$-가군임을,
더 정확히는
\[
R^nf_*(\mathcal F)=(H^n(\mathfrak X,\mathcal F))^\Delta
\tag{3.4.5.1}\label{III.3.4.5.1-fr}
\]
임을 증명하자. 앞의 표기를 쓰면 $\mathcal F_k$가 연접
$\mathcal O_{X_k}$-가군이므로
(\hyperref[III.1.4.13-fr]{1.4.13})
\[
\Gamma(\mathfrak Y_g,R^nf_*(\mathcal F_k))
=(\Gamma(\mathfrak Y,R^nf_*(\mathcal F_k)))_g
=(H^n(\mathfrak X,\mathcal F_k))_g.
\]
한편 $H^n(\mathfrak X,\mathcal F_k)$들은 (ML)을 만족하는 사영계를
이루고, 그 사영극한 $H^n(\mathfrak X,\mathcal F)$는 유한 생성
$A$-가군이다. 따라서 (0, 13.7.8)과 $A$, $A_{\{g\}}$에 적용한
(I, 10.10.8)에 의해
\[
\varprojlim_k(H^n(\mathfrak X,\mathcal F_k))_g
=H^n(\mathfrak X,\mathcal F)\otimes_AA_{\{g\}}
=\Gamma(\mathfrak Y_g,(H^n(\mathfrak X,\mathcal F))^\Delta)
\]
이다. 또한
\[
\Gamma(\mathfrak Y_g,R^nf_*(\mathcal F))
=\varprojlim_k\Gamma(\mathfrak Y_g,R^nf_*(\mathcal F_k))
\]
이므로, 이것이 (\hyperref[III.3.4.5.1-fr]{3.4.5.1})을 증명한다.
이제 (\hyperref[III.3.4.2.2-fr]{3.4.2.2})는 연접
$\mathcal O_{\mathfrak Y}$-가군의 동형사상이므로 반드시 \emph{위상}
동형사상이다 (I, 10.11.6). 끝으로
$R^nf_*(\mathcal F_k)=(H^n(\mathfrak X,\mathcal F_k))^\Delta$라는
관계에서 사영계 $(R^nf_*(\mathcal F_k))_{k\geq0}$가 (ML)을 만족함이
따른다 (I, 10.10.2).

(\hyperref[III.3.4.2-fr]{3.4.2})와
(\hyperref[III.3.4.3-fr]{3.4.3})를 $\mathfrak Y$가 뇌터 아핀
형식적 준스킴인 경우에 증명하고 나면,
(\hyperref[III.3.4.2-fr]{3.4.2})와
(\hyperref[III.3.4.3-fr]{3.4.3})의 첫째 명제는 $\mathfrak Y$ 위에서
국소적이므로 일반적인 경우로 넘어가는 것은 즉시 이루어진다.
(\hyperref[III.3.4.3-fr]{3.4.3})의 둘째 명제에 대해서는
$\mathfrak Y$가 뇌터이므로 이를 유한 개의 뇌터 아핀 열린집합
$U_i$들로 덮고, 사영계 $(R^qf_*(\mathcal F_k))$를 각 $U_i$에 제한한
것들이 (ML)을 만족함을 유의하면 충분하다.
\end{env}

더 나아가 위의 증명 과정에서 다음도 증명하였다.

\begin{corollary}[3.4.6]
\label{III.3.4.6-fr}
(\hyperref[III.3.4.4-fr]{3.4.4})의 가정 아래에서 표준 준동형사상
\[
H^q(\mathfrak X,\mathcal F)
\longrightarrow\Gamma(\mathfrak Y,R^qf_*(\mathcal F))
\tag{3.4.6.1}\label{III.3.4.6.1-fr}
\]
은 전단사이다.
\end{corollary}

\section{고유 사상의 기본 정리와 그 응용}
\label{section:III.4-fr}

\subsection{기본 정리}
\label{subsection:III.4.1-fr}

\begin{env}[4.1.1]
\label{III.4.1.1-fr}
$X$, $Y$를 통상적인 뇌터 준스킴들, $f:X\to Y$를 고유 사상,
$Y'$를 $Y$의 닫힌 부분집합, $X'=f^{-1}(Y')$를 그 역상이라 하자.
$X$와 $Y$를 각각 $X'$과 $Y'$을 따라 완비화하여 얻는 형식적 준스킴
$X_{/X'}$, $Y_{/Y'}$ (I, 10.8.5)를 각각 $\widehat X$,
$\widehat Y$로 나타내고, $f$를 이 완비화들로 연장한 형식적 준스킴의
사상 $\widehat f:\widehat X\to\widehat Y$ (I, 10.9.1)를
$\widehat f$로 나타내겠다. 모든 연접 $\mathcal O_X$-가군
$\mathcal F$에 대하여,

\oldpage[III]{123}

$X'$을 따른 그 완비화 $\mathcal F_{/X'}$ (I, 10.8.4)를
$\widehat{\mathcal F}$로 나타내겠다. 이는 연접
$\mathcal O_{\widehat X}$-가군이다 (I, 10.8.8).
\end{env}

\begin{env}[4.1.2]
\label{III.4.1.2-fr}
$\operatorname{Supp}(\mathcal O_Y/\mathcal J)=Y'$인
$\mathcal O_Y$의 연접 아이디얼 $\mathcal J$를 택하자 (I, 5.2.1).
(I, 4.4.5)에 의해
\[
\mathcal K=f^*(\mathcal J)\mathcal O_X
\]
는
\[
\operatorname{Supp}(\mathcal O_X/\mathcal K)=X'
\]
인 $\mathcal O_X$의 연접 아이디얼이다. 모든 $k\geq0$에 대하여 연접
$\mathcal O_X$-가군
\[
\mathcal F_k
=\mathcal F\otimes_{\mathcal O_Y}(\mathcal O_Y/\mathcal J^{k+1})
=\mathcal F/\mathcal K^{k+1}\mathcal F
\]
를 생각하자. 모든 $n$에 대해 $\mathcal O_Y$-가군
$R^nf_*(\mathcal F)$와 $R^nf_*(\mathcal F_k)$는 연접이다
(\hyperref[III.3.2.1-fr]{3.2.1}). 모든 $k\geq0$과 모든 $n$에 대해
표준 준동형사상 $\mathcal F\to\mathcal F_k$는 함자성에 의해
준동형사상
\[
R^nf_*(\mathcal F)\longrightarrow R^nf_*(\mathcal F_k)
\tag{4.1.2.1}\label{III.4.1.2.1-fr}
\]
을 정의한다. 또한 $\mathcal F_k$는
$\mathcal O_X/\mathcal K^{k+1}$-가군이므로
$R^nf_*(\mathcal F_k)$는
$\mathcal O_Y/\mathcal J^{k+1}$-가군이다 (0, 12.2.1). 따라서
(\hyperref[III.4.1.2.1-fr]{4.1.2.1})에서 준동형사상
\[
R^nf_*(\mathcal F)\otimes_{\mathcal O_Y}
(\mathcal O_Y/\mathcal J^{k+1})
\longrightarrow R^nf_*(\mathcal F_k)
\tag{4.1.2.2}\label{III.4.1.2.2-fr}
\]
을 얻는다. (\hyperref[III.4.1.2.2-fr]{4.1.2.2})의 양변은 각각
사영계를 이루며, 왼쪽의 사영극한은 다름 아닌
$(R^nf_*(\mathcal F))_{/Y'}$의 완비화이다. 이를
$\widehat{R^nf_*(\mathcal F)}$로 나타내겠다. 또한
(\hyperref[III.4.1.2.2-fr]{4.1.2.2})의 준동형사상들이 사영계를
이룸은 즉시 알 수 있으므로, 극한을 취하여 표준 준동형사상
\[
\varphi_n:\widehat{R^nf_*(\mathcal F)}
\longrightarrow\varprojlim_kR^nf_*(\mathcal F_k)
\tag{4.1.2.3}\label{III.4.1.2.3-fr}
\]
을 얻는다. 한편 (\hyperref[III.4.1.2.2-fr]{4.1.2.2})는
$(\mathcal O_Y/\mathcal J^{k+1})$-가군의 준동형사상이므로
(I, 10.8.3), 의사이산 위상
$\mathcal O_{\widehat Y}$-가군의 연속 준동형사상으로 볼 수 있다
(0\textsubscript{I}, 3.8.1). 따라서 $\varphi_n$은 위상
$\mathcal O_{\widehat Y}$-가군의 연속 준동형사상이다.
\end{env}

\begin{env}[4.1.3]
\label{III.4.1.3-fr}
$i:\widehat X\to X$를 (I, 10.8.7)에서 정의한 환 달린 공간의 표준
사상이라 하자. 그러면 가환 도표
\[
\xymatrix{
X_k\ar[r]^{h_k}\ar[dr]_{i_k}&\widehat X\ar[d]^i\\
&X
}
\tag{4.1.3.1}\label{III.4.1.3.1-fr}
\]
를 얻는다. 여기서 $X_k$는 아이디얼 $\mathcal K^{k+1}$로 정의되는
$X$의 닫힌 부분준스킴이고, $i_k$는 표준 단사이며, $h_k$는 바탕공간의
항등사상과 표준 준동형사상
$\mathcal O_{\widehat X}\to\mathcal O_X/\mathcal K^{k+1}$에 대응하는
환 달린 공간의 사상이다 (I, 10.5.2). 또한 표준 동형사상을 제외하면
$\widehat{\mathcal F}=i^*(\mathcal F)$이다 (I, 10.8.8). 다음 동형사상이
알려져 있다.
\[
H^n(X_k,i_k^*(\mathcal F_k))=H^n(X,\mathcal F_k)
\tag{4.1.3.2}\label{III.4.1.3.2-fr}
\]

\oldpage[III]{124}

실제로 $\mathcal F_k=(i_k)_*(i_k^*(\mathcal F_k))$이기 때문이다
(G, II, 4.9.1). 그러므로 표준 준동형사상
\[
H^n(\widehat X,\widehat{\mathcal F})
\longrightarrow H^n(X_k,h_k^*(\widehat{\mathcal F}))
\qquad\text{(0, 12.1.3.5)}
\]
은
\[
H^n(\widehat X,\widehat{\mathcal F})\longrightarrow H^n(X,\mathcal F_k)
\tag{4.1.3.3}\label{III.4.1.3.3-fr}
\]
로도 쓸 수 있다. 이 준동형사상들은 명백히 사영계를 이루므로, 극한을
취하여 표준 준동형사상
\[
\psi_{n,X}:H^n(\widehat X,\widehat{\mathcal F})
\longrightarrow\varprojlim_kH^n(X,\mathcal F_k)
\tag{4.1.3.4}\label{III.4.1.3.4-fr}
\]
을 얻는다. $X$를 $Y$의 아핀 열린집합 $V$에 대한
$f^{-1}(V)$ 꼴의 열린집합으로 바꾸고
(\hyperref[III.1.4.11-fr]{1.4.11})을 고려하면 준동형사상
\[
\psi_{n,V}:H^n(\widehat X\cap f^{-1}(V),\widehat{\mathcal F})
\longrightarrow\varprojlim_k\Gamma(V,R^nf_*(\mathcal F_k))
\tag{4.1.3.5}\label{III.4.1.3.5-fr}
\]
을 얻는다. 이 준동형사상들은 $V$를 더 작은 아핀 열린집합으로 제한하는
것과 명백히 가환하므로, 결국 층의 표준 준동형사상
\[
\psi_n:R^n\widehat f_*(\widehat{\mathcal F})
\longrightarrow\varprojlim_kR^nf_*(\mathcal F_k)
\tag{4.1.3.6}\label{III.4.1.3.6-fr}
\]
을 정의한다.
\end{env}

\begin{env}[4.1.4]
\label{III.4.1.4-fr}
끝으로 $j:\widehat Y\to Y$를 환 달린 공간의 표준 사상이라 하자
(I, 10.8.7). $R^nf_*(\mathcal F)$는 연접 $\mathcal O_Y$-가군이므로
(\hyperref[III.3.2.1-fr]{3.2.1}), 표준 동형사상을 제외하면
\[
j^*(R^nf_*(\mathcal F))=\widehat{R^nf_*(\mathcal F)}
\]
이다 (I, 10.8.8). 따라서 환 달린 공간에 대해 일반적으로 정의되는
표준 준동형사상
\[
\rho_n:\widehat{R^nf_*(\mathcal F)}=j^*(R^nf_*(\mathcal F))
\longrightarrow R^n\widehat f_*(i^*(\mathcal F))
=R^n\widehat f_*(\widehat{\mathcal F})
\tag{4.1.4.1}\label{III.4.1.4.1-fr}
\]
을 얻는다
((\hyperref[III.1.4.15-fr]{1.4.15})의 증명을 보라). 도표
\[
\xymatrix{
\widehat{R^nf_*(\mathcal F)}\ar[rr]^{\rho_n}\ar[dr]_{\varphi_n}
&&R^n\widehat f_*(\widehat{\mathcal F})\ar[dl]^{\psi_n}\\
&\displaystyle\varprojlim_kR^nf_*(\mathcal F_k)&
}
\tag{4.1.4.2}\label{III.4.1.4.2-fr}
\]
가 가환임을 보이자. 대응하는 준층 준동형사상의 도표가 가환임을 보이면
충분하므로 $Y$가 아핀인 경우로 한정할 수 있다. 그러면 문제는 도표
\[
\xymatrix{
\widehat{H^n(X,\mathcal F)}\ar[rr]^{\rho_n}\ar[dr]_{\varphi_n}
&&H^n(\widehat X,\widehat{\mathcal F})\ar[dl]^{\psi_{n,X}}\\
&\displaystyle\varprojlim_kH^n(X,\mathcal F_k)&
}
\tag{4.1.4.3}\label{III.4.1.4.3-fr}
\]
가 가환임을 보이는 것으로 귀착된다. 그런데
(\hyperref[III.4.1.3.1-fr]{4.1.3.1})의 가환성과
(\hyperref[III.4.1.3-fr]{4.1.3})에서 본 코호몰로지 군 사이의 관계들은
즉시 가환 도표

\oldpage[III]{125}

\[
\xymatrix{
H^n(X,\mathcal F)\ar[r]\ar[dr]&
H^n(\widehat X,\widehat{\mathcal F})
=H^n(\widehat X,i^*(\mathcal F))\ar[d]\\
&H^n(X,\mathcal F_k)=H^n(X_k,i_k^*(\mathcal F))
}
\]
를 주며, 여기서 (\hyperref[III.4.1.4.3-fr]{4.1.4.3})의 가환성이 즉시
따른다.
\end{env}

\begin{theorem}[4.1.5]
\label{III.4.1.5-fr}
$f:X\to Y$를 뇌터 준스킴들의 고유 사상, $Y'$를 $Y$의 닫힌 부분집합,
$X'=f^{-1}(Y')$라 하자. 그러면 모든 연접 $\mathcal O_X$-가군
$\mathcal F$에 대해 $R^n\widehat f_*(\widehat{\mathcal F})$는 연접
$\mathcal O_{\widehat Y}$-가군이고, 도표
(\hyperref[III.4.1.4.2-fr]{4.1.4.2})의 준동형사상 $\varphi_n$,
$\psi_n$, $\rho_n$은 위상 동형사상이다.
\end{theorem}

\begin{proof}
$\varphi_n$과 $\psi_n$이 동형사상임을 보이면 충분하다. 실제로
$R^nf_*(\mathcal F)$는 연접이므로
(\hyperref[III.3.2.1-fr]{3.2.1}),
$\widehat{R^nf_*(\mathcal F)}$도 연접이고 (I, 10.8.8),
$\varphi_n$, $\psi_n$, $\rho_n$의 양방향 연속성은 자동으로 따른다
(I, 10.11.6).
\end{proof}

\begin{env}[4.1.6]
\label{III.4.1.6-fr}
\emph{주의.}
\begin{enumerate}
\item[(i)]
$\widehat{\mathcal F}_k=
\widehat{\mathcal F}/\widehat{\mathcal K}^{k+1}\widehat{\mathcal F}$로
놓으면 $\mathcal F_k=i_*(\widehat{\mathcal F}_k)$임은 즉시 알 수 있다.
표준 준동형사상 (\hyperref[III.4.1.3.6-fr]{4.1.3.6})은
(\hyperref[III.3.4.2.2-fr]{3.4.2.2})에서 이미 정의한 준동형사상
\[
R^n\widehat f_*(\widehat{\mathcal F})
\longrightarrow\varprojlim_kR^n\widehat f_*(\widehat{\mathcal F}_k)
\tag{4.1.6.1}\label{III.4.1.6.1-fr}
\]
과 다르지 않다. 따라서 $\psi_n$이 동형사상이라는 사실은
(\hyperref[III.3.4.3-fr]{3.4.3})의 특수한 경우이다. 그러나 아래에서는
일반 정리 (\hyperref[III.3.4.3-fr]{3.4.3})가 의존하는 스펙트럼 열들의
사영극한에 관한 섬세한 논의 (0, 13.7)를 피하는 직접 증명을 제시하겠다.

\item[(ii)]
$\psi_n$들이 동형사상이라는 사실을 고려하면, $\varphi_n$들이
동형사상이라는 것과
$\rho_n=\psi_n^{-1}\circ\varphi_n$들이 동형사상이라는 것은 동치이다.
정리 (\hyperref[III.4.1.5-fr]{4.1.5})는 특히 $R^nf_*$의 형성이
완비화와 가환함을 표현하며, 이를 ‘대수적’ 이론과 ‘형식적’ 이론의
제1 비교 정리라고 부를 수 있다.
\end{enumerate}
이제 먼저 (\hyperref[III.4.1.5-fr]{4.1.5})의 아핀 형태를 증명하겠다.
\end{env}

\begin{corollary}[4.1.7]
\label{III.4.1.7-fr}
(\hyperref[III.4.1.5-fr]{4.1.5})의 가정들에 더하여
$Y=\operatorname{Spec}(A)$이고 $A$가 뇌터이며,
$\mathcal J=\widetilde{\mathfrak J}$이고 $\mathfrak J$가 $A$의
아이디얼이라고 가정하자. 따라서
$\mathcal F_k=\mathcal F/\mathfrak J^{k+1}\mathcal F$이다. 표준
준동형사상
\[
\varphi_n:\widehat{H^n(X,\mathcal F)}
\longrightarrow\varprojlim_kH^n(X,\mathcal F_k)
\tag{4.1.7.1}\label{III.4.1.7.1-fr}
\]
은 동형사상이다. 여기서 왼쪽은 $\mathfrak J$-준아딕 위상에 관한
$H^n(X,\mathcal F)$의 분리 완비화이다. 사영계
$(H^n(X,\mathcal F_k))_{k\geq0}$는 모든 $n$에 대해 조건 (ML)을
만족하고, 표준 준동형사상
\[
\psi_n:H^n(\widehat X,\widehat{\mathcal F})
\longrightarrow\varprojlim_kH^n(X,\mathcal F_k)
\tag{4.1.7.2}\label{III.4.1.7.2-fr}
\]

\oldpage[III]{126}

은 동형사상이다. 끝으로 표준 준동형사상
$H^n(X,\mathcal F)\to H^n(X,\mathcal F_k)$들의 핵들이
$H^n(X,\mathcal F)$ 위에 정의하는 여과는 $\mathfrak J$-좋은 여과이고
(0, 13.7.7), $\varphi_n$은 위상 동형사상이다.
\footnote{뒤따르는 증명은 최초의 증명보다 간단하며,
$H^n(X,\mathcal F)$의 여과에 관한 보충과 함께 J.-P. Serre가
알려 주었다.}
\end{corollary}

\begin{proof}
이 증명에서 정수 $n\geq0$을 고정하고, 표기를 간단히 하기 위해
\[
\mathrm H=H^n(X,\mathcal F),\qquad
\mathrm H_k=H^n(X,\mathcal F_k)
\tag{4.1.7.3}\label{III.4.1.7.3-fr}
\]
및
\[
\mathrm R_k=\operatorname{Ker}(\mathrm H\longrightarrow\mathrm H_k),
\qquad\text{$\mathrm H$의 부분 $A$-가군}
\tag{4.1.7.4}\label{III.4.1.7.4-fr}
\]
로 놓겠다. 코호몰로지 완전열
\[
H^n(X,\mathfrak J^{k+1}\mathcal F)\longrightarrow H^n(X,\mathcal F)
\longrightarrow H^n(X,\mathcal F_k)\xrightarrow{\partial}
H^{n+1}(X,\mathfrak J^{k+1}\mathcal F)
\longrightarrow H^{n+1}(X,\mathcal F)
\]
에서 또한
\[
\mathrm R_k=\operatorname{Im}\bigl(
H^n(X,\mathfrak J^{k+1}\mathcal F)\longrightarrow H^n(X,\mathcal F)
\bigr)
\]
임을 알 수 있다. 다음과 같이 놓자.
\[
\begin{aligned}
\mathrm Q_k
&=\operatorname{Ker}\bigl(H^{n+1}(X,\mathfrak J^{k+1}\mathcal F)
\longrightarrow H^{n+1}(X,\mathcal F)\bigr)\\
&=\operatorname{Im}\bigl(H^n(X,\mathcal F_k)
\longrightarrow H^{n+1}(X,\mathfrak J^{k+1}\mathcal F)\bigr).
\end{aligned}
\tag{4.1.7.5}\label{III.4.1.7.5-fr}
\]
따라서 완전열
\[
0\longrightarrow\mathrm R_k\longrightarrow\mathrm H
\longrightarrow\mathrm H_k\longrightarrow\mathrm Q_k\longrightarrow0
\tag{4.1.7.6}\label{III.4.1.7.6-fr}
\]
을 얻는다.

\begin{env}[4.1.7.7]
\label{III.4.1.7.7-fr}
$x$를 $\mathfrak J^m$의 원소라 하자 ($m\geq0$).
$\mathfrak J^k\mathcal F$에서 $x$를 곱하는 것은 준동형사상
$\mathfrak J^k\mathcal F\to\mathfrak J^{k+m}\mathcal F$이며, 따라서
준동형사상
\[
\mu_{x,m}:H^n(X,\mathfrak J^k\mathcal F)
\longrightarrow H^n(X,\mathfrak J^{k+m}\mathcal F)
\tag{4.1.7.8}\label{III.4.1.7.8-fr}
\]
을 준다. $S=\bigoplus_{k\geq0}\mathfrak J^k$라는 등급 $A$-대수를
생각하면, 곱셈 $\mu_{x,m}$들은
\[
\mathrm E=\bigoplus_{k\geq0}H^n(X,\mathfrak J^k\mathcal F)
\]
에 등급환 $S$ 위의 유한 생성 등급 가군 구조를 정의한다
(\hyperref[III.3.3.2-fr]{3.3.2}). 이 등급환은 뇌터이다
(II, 2.1.5).
\end{env}

\begin{lemma}[4.1.7.9]
\label{III.4.1.7.9-fr}
$\mathrm H$의 부분가군 $\mathrm R_k$들은 $\mathrm H$ 위에
$\mathfrak J$-좋은 여과를 정의한다.
\end{lemma}

\begin{proof}
먼저
\[
\mathfrak J^m\mathrm R_k\subset\mathrm R_{k+m}
\tag{4.1.7.10}\label{III.4.1.7.10-fr}
\]
임을 보이자. 여기서 $x\in\mathfrak J^m$을
$\mathrm H=H^n(X,\mathcal F)$에 곱하는 사상은 $\mu_{x,0}$이다. 모든
$x\in\mathfrak J^m$에 대해 도표
\[
\xymatrix{
\mathfrak J^{k+1}\mathcal F\ar[r]^x\ar[d]&
\mathfrak J^{k+m+1}\mathcal F\ar[d]\\
\mathcal F\ar[r]_x&\mathcal F
}
\tag{4.1.7.11}\label{III.4.1.7.11-fr}
\]

\oldpage[III]{127}

는 가환이다. 여기서 수평 화살표들은 $x$의 곱셈이고 수직 화살표들은
표준 단사이다. 따라서 대응하는 도표
\[
\xymatrix{
H^n(X,\mathfrak J^{k+1}\mathcal F)\ar[r]^{\mu_{x,m}}\ar[d]&
H^n(X,\mathfrak J^{k+m+1}\mathcal F)\ar[d]\\
H^n(X,\mathcal F)\ar[r]_{\mu_{x,0}}&H^n(X,\mathcal F)
}
\]
도 가환한다. $\mathrm R_k$를
$H^n(X,\mathfrak J^{k+1}\mathcal F)\to H^n(X,\mathcal F)$의 상으로
해석하면 이것이 (\hyperref[III.4.1.7.10-fr]{4.1.7.10})을 증명한다.
또한 등급 $S$-가군
$\mathrm R=\bigoplus_{k\geq0}\mathrm R_k$가 $\mathrm E$의 부분
$S$-가군
\[
\mathrm M=\bigoplus_{k\geq0}
H^n(X,\mathfrak J^{k+1}\mathcal F)
\]
의 몫임을 보여 준다. 그러므로 위의 주의에 의해 $\mathrm R$은 유한
생성 $S$-가군이다. 이는
(\hyperref[III.4.1.7.9-fr]{4.1.7.9})의 주장과 동치이다
(Bourbaki, \emph{Alg. comm.}, chap. III, §~3,
n\textsuperscript{o}~1, th.~1).
\end{proof}

\begin{env}[4.1.7.12]
\label{III.4.1.7.12-fr}
이제 (4.7.1.8)과 같이 정의되는 등급 $S$-가군
\[
\mathrm N=\bigoplus_{k\geq0}
H^{n+1}(X,\mathfrak J^{k+1}\mathcal F)
\]
을 생각하자. 이것도
(\hyperref[III.3.3.2-fr]{3.3.2})에 의해 유한 생성 $S$-가군이다.
(\hyperref[III.4.1.7.5-fr]{4.1.7.5})에 의해 모든 $k$에 대해
$\mathrm Q_k\subset\mathrm N_k$이고,
(\hyperref[III.4.1.7.11-fr]{4.1.7.11})에서 $n$을 $n+1$로 바꾼 도표는
\[
S_m\mathrm Q_k=\mathfrak J^m\mathrm Q_k\subset\mathrm Q_{k+m}
\]
임을 보여 준다. 다시 말해 $\mathrm Q$는 $\mathrm N$의 등급
부분 $S$-가군이고, 따라서 유한 생성이다.
\end{env}

\begin{env}[4.1.7.13]
\label{III.4.1.7.13-fr}
$S_m\to S_0$로도 쓸 수 있는 표준 단사
$\alpha_m:\mathfrak J^m\to A$를 생각하자.
$\mathfrak J^{k+1}\mathcal F_k=0$이므로 $A$-가군
$H^n(X,\mathcal F_k)$는 $\mathfrak J^{k+1}$에 의해 소멸한다.
$\mathrm Q_k$는 $A$-준동형사상
$H^n(X,\mathcal F_k)\to
H^{n+1}(X,\mathfrak J^{k+1}\mathcal F)$의 상이므로,
$A$-가군 $\mathrm Q_k$도 $\mathfrak J^{k+1}$에 의해 소멸한다.
이는 $S$-가군 $\mathrm Q$에서
\[
\alpha_{k+1}(S_{k+1})\mathrm Q_k=0
\tag{4.1.7.14}\label{III.4.1.7.14-fr}
\]
임을 뜻한다. $\mathrm Q$는 유한 생성 $S$-가군이므로 어떤 정수
$k_0$와 정수 $h$가 존재하여 $k\geq k_0$일 때
$\mathrm Q_{k+h}=S_h\mathrm Q_k$이다 (II, 2.1.6, (ii)).
이 관계와 (\hyperref[III.4.1.7.14-fr]{4.1.7.14})에서 어떤 정수
$r>0$이 존재하여
\[
\alpha_r(S_r)\mathrm Q=0
\tag{4.1.7.15}\label{III.4.1.7.15-fr}
\]
임이 따른다.
\end{env}

\begin{env}[4.1.7.16]
\label{III.4.1.7.16-fr}
표준 단사 $\mathfrak J^{k+m}\mathcal F\to\mathfrak J^k\mathcal F$는
코호몰로지에서 $A$-준동형사상
\[
\nu_m:H^{n+1}(X,\mathfrak J^{k+m}\mathcal F)
\longrightarrow H^{n+1}(X,\mathfrak J^k\mathcal F)
\tag{4.1.7.17}\label{III.4.1.7.17-fr}
\]
을 준다. 모든 $x\in\mathfrak J^m$에 대해 명백한 분해
\[
\xymatrix@C=4.5em{
H^{n+1}(X,\mathfrak J^k\mathcal F)\ar[r]^{\mu_{x,m}}&
H^{n+1}(X,\mathfrak J^{k+m}\mathcal F)\ar[r]^{\nu_m}&
H^{n+1}(X,\mathfrak J^k\mathcal F)
}
\tag{4.1.7.18}\label{III.4.1.7.18-fr}
\]
가 있으며, 그 합성은 $\mu_{x,0}$이다.
\end{env}

\oldpage[III]{128}

따라서 $H^{n+1}(X,\mathfrak J^k\mathcal F)$의 모든 부분
$A$-가군 $P$에 대하여 $S$-가군 $\mathrm N$ 안에서
\[
\nu_m(S_mP)=\alpha_m(S_m)P
\tag{4.1.7.19}\label{III.4.1.7.19-fr}
\]
이다.

\begin{lemma}[4.1.7.20]
\label{III.4.1.7.20-fr}
어떤 정수 $m>0$이 존재하여 모든 $k\geq k_0$에 대해
$\nu_m(\mathrm Q_{k+m})=0$이다.
\end{lemma}

\begin{proof}
$h$의 배수이면서 $m\geq r$인 $m$을 택하자.
$k\geq k_0$일 때 $\mathrm Q_{k+m}=S_m\mathrm Q_k$이므로
(\hyperref[III.4.1.7.19-fr]{4.1.7.19})와
(\hyperref[III.4.1.7.15-fr]{4.1.7.15})에 의해
\[
\nu_m(\mathrm Q_{k+m})=\alpha_m(S_m)\mathrm Q_k
\subset\alpha_r(S_r)\mathrm Q_k=0.
\]
\end{proof}

\begin{env}[4.1.7.21]
\label{III.4.1.7.21-fr}
가환 도표
\[
\xymatrix@C=2.5em{
H^n(X,\mathcal F)\ar[r]&H^n(X,\mathcal F_k)\ar[r]^-\partial&
H^{n+1}(X,\mathfrak J^{k+1}\mathcal F)\ar[r]&H^{n+1}(X,\mathcal F)\\
H^n(X,\mathcal F)\ar[r]\ar[u]&H^n(X,\mathcal F_{k+m})\ar[r]^-\partial\ar[u]&
H^{n+1}(X,\mathfrak J^{k+m+1}\mathcal F)\ar[r]\ar[u]&
H^{n+1}(X,\mathcal F)\ar[u]
}
\]
는 수직 화살표들이 표준 사상인 가환 도표
\[
\xymatrix{
0\ar[r]&\mathfrak J^{k+1}\mathcal F\ar[r]&\mathcal F\ar[r]&
\mathcal F_k\ar[r]&0\\
0\ar[r]&\mathfrak J^{k+m+1}\mathcal F\ar[r]\ar[u]&
\mathcal F\ar[r]\ar[u]&\mathcal F_{k+m}\ar[r]\ar[u]&0
}
\]
에서 유도된다. 따라서 행들이 완전인 가환 도표
\[
\xymatrix@C=2.2em{
0\ar[r]&\mathrm R_k\ar[r]&\mathrm H\ar[r]&\mathrm H_k\ar[r]&
\mathrm Q_k\ar[r]&0\\
0\ar[r]&\mathrm R_{k+m}\ar[r]\ar[u]&\mathrm H\ar[r]\ar@{=}[u]&
\mathrm H_{k+m}\ar[r]\ar[u]&\mathrm Q_{k+m}\ar[r]\ar[u]^{\nu_m}&0
}
\]
를 얻는다. $k\geq k_0$일 때 마지막 수직 화살표는 영이므로
(\hyperref[III.4.1.7.20-fr]{4.1.7.20}), $\mathrm H_{k+m}$의
$\mathrm H_k$ 안의 상은
\[
\operatorname{Ker}(\mathrm H_k\to\mathrm Q_k)
=\operatorname{Im}(\mathrm H\to\mathrm H_k)
\]
에 포함된다. 한편 도표의 가환성에 의해 이 상은
$\operatorname{Im}(\mathrm H\to\mathrm H_k)$를 포함하므로 둘은 같다.
따라서 모든 $k'\geq k+m$에 대해 $\mathrm H_{k'}$의 $\mathrm H_k$
안의 상도 이와 같으며, 이것이 사영계
$(\mathrm H_k)_{k\geq0}$에 대한 조건 (ML)을 증명한다. 또한 $X$의 모든
아핀 열린집합 $U$에 대해 $i>0$이면
$H^i(U,\mathcal F_k)=0$이고
(\hyperref[III.1.3.1-fr]{1.3.1}), $m>0$이면
$H^0(U,\mathcal F_{k+m})\to H^0(U,\mathcal F_k)$가 전사이다
(I, 1.3.9). 따라서 (0, 13.3.1)을 적용할 수 있고,

\oldpage[III]{129}

표준 준동형사상
\[
H^n(\widehat X,\widehat{\mathcal F})
\longrightarrow\varprojlim_kH^n(X,\mathcal F_k)
\]
은 모든 $n\geq0$에 대해 전단사이다.

사영계 $(\mathrm H/\mathrm R_k)_{k\geq0}$는 엄밀하므로 완전열
\[
0\longrightarrow\mathrm H/\mathrm R_k\longrightarrow\mathrm H_k
\longrightarrow\mathrm Q_k\longrightarrow0
\tag{4.1.7.22}\label{III.4.1.7.22-fr}
\]
에서 사영극한을 취할 수 있다 (0, 13.2.2).
$\nu_m(\mathrm Q_{k+m})=0$이므로
$\varprojlim_k\mathrm Q_k=0$이며, 따라서 위상 동형사상
\[
\varprojlim_k(\mathrm H/\mathrm R_k)\xrightarrow{\sim}
\varprojlim_k\mathrm H_k
\]
을 얻는다. 그런데 $\mathrm H$의 여과 $(\mathrm R_k)$는
$\mathfrak J$-좋은 여과이므로 $\mathrm H$ 위에
$\mathfrak J$-준아딕 위상을 정의한다. 따라서
$\varprojlim_k(\mathrm H/\mathrm R_k)$는 이 위상에 관한
$\mathrm H$의 분리 완비화이다. 이것으로
(\hyperref[III.4.1.7-fr]{4.1.7})의 증명이 끝난다.
\end{env}
\end{proof}

\begin{env}[4.1.8]
\label{III.4.1.8-fr}
끝으로 (\hyperref[III.4.1.5-fr]{4.1.5})를 증명하자. $Y$의 모든 아핀
열린집합 $V$에 대하여
$\Gamma(V,\widehat{R^nf_*(\mathcal F)})$는
$\Gamma(V,R^nf_*(\mathcal F))$의 $\mathfrak J$-준아딕 위상에 관한
분리 완비화이다. 여기서
$\mathcal J|V=\widetilde{\mathfrak J}$로 쓴다. 실제로
$R^nf_*(\mathcal F)$는 연접 $\mathcal O_Y$-가군이다 (I, 10.8.4).
또한
\[
\Gamma\bigl(V,\varprojlim_kR^nf_*(\mathcal F_k)\bigr)
=\varprojlim_k\Gamma(V,R^nf_*(\mathcal F_k))
\]
이다 (0\textsubscript{I}, 3.2.6). 따라서 $\varphi_n$이 위상
동형사상이라는 사실은 (\hyperref[III.4.1.7-fr]{4.1.7})과
(\hyperref[III.1.4.11-fr]{1.4.11})에서 따른다. 한편 역시
(\hyperref[III.1.4.11-fr]{1.4.11})에 의해
(\hyperref[III.4.1.7-fr]{4.1.7})에서
(\hyperref[III.4.1.3.3-fr]{4.1.3.3})의 준동형사상 $\psi_{n,V}$가
동형사상임이 따른다. 그러므로
$R^n\widehat f_*(\widehat{\mathcal F})$의 정의에 의해 $\psi_n$도
동형사상이다.
\end{env}

\begin{corollary}[4.1.9]
\label{III.4.1.9-fr}
(\hyperref[III.4.1.4-fr]{4.1.4})의 가정 아래에서 $Y$의 모든 아핀
열린집합 $V$에 대해 표준 준동형사상
\[
H^n(\widehat X\cap f^{-1}(V),\widehat{\mathcal F})
\longrightarrow\Gamma(\widehat Y\cap V,
R^n\widehat f_*(\widehat{\mathcal F}))
\]
은 전단사이다.
\end{corollary}

\begin{env}[4.1.10]
\label{III.4.1.10-fr}
\emph{주의.}
$f:X\to Y$를 통상적인 뇌터 준스킴들의 유한형 사상이라 하고,
$\mathcal F$를 그 받침이 $Y$ 위에서 고유인 연접
$\mathcal O_X$-가군이라 하자 (II, 5.4.10). 그러면 모든 $n\geq0$에
대해 $R^nf_*(\mathcal F)$가 연접 $\mathcal O_Y$-가군임은 알려져 있다
(\hyperref[III.3.2.4-fr]{3.2.4}). 또한 항상
$\mathcal F=u_*(\mathcal G)$라고 가정할 수 있다. 여기서 $Z$는
바탕공간이 $\operatorname{Supp}(\mathcal F)$인 $X$의 적당한 닫힌
부분준스킴이고, $u:Z\to X$는 표준 단사이며,
$\mathcal G=u^*(\mathcal F)$는 연접 $\mathcal O_Z$-가군이다
(I, 9.3.5). 다음과 같이 놓자.
\[
\mathcal G_k=\mathcal G\otimes_{\mathcal O_Y}
(\mathcal O_Y/\mathcal J^{k+1}).
\]
그러면 $\mathcal G_k=u^*(\mathcal F_k)$이고,
\[
R^nf_*(\mathcal F_k)=R^n(f\circ u)_*(\mathcal G_k),\qquad
R^nf_*(\mathcal F)=R^n(f\circ u)_*(\mathcal G)
\]
이다 (\hyperref[III.1.3.4-fr]{1.3.4}). 끝으로 (I, 10.9.5)를 고려하면
\[
R^n\widehat f_*(\widehat{\mathcal F})
=R^n\widehat{(f\circ u)}_*(\widehat{\mathcal G}).
\]
따라서 $\mathcal G$와 고유 사상 $f\circ u$에
(\hyperref[III.4.1.5-fr]{4.1.5})를 적용할 수 있고, 이 가정 아래에서도
(\hyperref[III.4.1.5-fr]{4.1.5})의 결과들이 $\mathcal F$와 $f$에 대해
성립한다.
\end{env}

\subsection{특수한 경우들과 변형들}
\label{subsection:III.4.2-fr}

비교 정리 (\hyperref[III.4.1.5-fr]{4.1.5})의 가장 유용한 형태는
다음과 같다.

\begin{proposition}[4.2.1]
\label{III.4.2.1-fr}
$Y$를 국소 뇌터 준스킴, $f:X\to Y$를 고유 사상, $\mathcal F$를 연접
$\mathcal O_X$-가군이라 하자. 그러면 모든 $y\in Y$와 모든 $p$에 대해
$(R^pf_*(\mathcal F))_y$는

\oldpage[III]{130}

유한 생성 $\mathcal O_y$-가군이고, 따라서 $\mathfrak m_y$-준아딕
위상에 관하여 분리이다. 또한 표준 위상 동형사상
\[
\widehat{(R^pf_*(\mathcal F))_y}
\xrightarrow{\sim}
\varprojlim_kH^p\bigl(f^{-1}(y),
\mathcal F\otimes_{\mathcal O_Y}(\mathcal O_y/\mathfrak m_y^k)\bigr)
\tag{4.2.1.1}\label{III.4.2.1.1-fr}
\]
이 있다. 여기서 왼쪽은 $(R^pf_*(\mathcal F))_y$의
$\mathfrak m_y$-준아딕 위상에 관한 완비화이고, 오른쪽에서는 모든
$k\geq0$에 대해 $f^{-1}(y)$를 준스킴
\[
X\times_Y\operatorname{Spec}(\mathcal O_y/\mathfrak m_y^k)
\]
의 바탕공간으로 본다 (I, 3.6.1).
\end{proposition}

\begin{proof}
$\mathcal O_y$는 뇌터 국소환이고
$(R^pf_*(\mathcal F))_y$는 유한 생성 $\mathcal O_y$-가군이므로
(\hyperref[III.3.2.1-fr]{3.2.1}), 그 위의
$\mathfrak m_y$-준아딕 위상은 분리이다 (0\textsubscript{I}, 7.3.5).
$Y$가 뇌터이고 $y$가 닫힌점인 경우에는 나머지 주장들이
(\hyperref[III.4.1.7-fr]{4.1.7})에서 따른다. 실제로 $Y$를 $y$의
아핀 근방으로 바꾸고 $Y'=\{y\}$로 택하면 된다
(G, II, 4.9.1). 일반적인 경우에는
\[
Y_1=\operatorname{Spec}(\mathcal O_y),\qquad
X_1=X\times_YY_1,\qquad
\mathcal F_1=\mathcal F\otimes_{\mathcal O_Y}\mathcal O_{Y_1}
\]
로 놓고 $f_1=f\times1_{Y_1}:X_1\to Y_1$라 하자. $Y_1$은 뇌터이고,
$f_1$은 고유이며 (II, 5.4.2, (iii)), $\mathcal F_1$은 연접이다
(0\textsubscript{I}, 5.3.11). $Y_1$의 유일한 닫힌점을 $y_1$이라 하자.
명제는 $f_1$, $\mathcal F_1$, $y_1$에 대해 성립한다. 또한
\[
\mathcal O_{y_1}=\mathcal O_y,\qquad
f_1^{-1}(y_1)=f^{-1}(y)
\]
이고 (I, 3.6.5), 준스킴
\[
X\times_Y\operatorname{Spec}(\mathcal O_y/\mathfrak m_y^k)
\quad\text{및}\quad
X_1\times_{Y_1}\operatorname{Spec}(\mathcal O_y/\mathfrak m_y^k)
\]
은 표준적으로 동일시된다 (I, 3.3.9). 더 나아가
\[
\mathcal F_1\otimes_{\mathcal O_{Y_1}}
(\mathcal O_y/\mathfrak m_y^k)
\]
는
$\mathcal F\otimes_{\mathcal O_Y}(\mathcal O_y/\mathfrak m_y^k)$와
동일시된다 (I, 9.1.6). 이제
\[
R^pf_{1*}(\mathcal F_1)\simeq
R^pf_*(\mathcal F)\otimes_{\mathcal O_Y}\mathcal O_{Y_1}
\]
임만 보이면 된다. 이는 국소 사상
$\operatorname{Spec}(\mathcal O_y)\to Y$가 평탄하므로
(0\textsubscript{I}, 6.7.1 및 I, 2.4.2),
(\hyperref[III.1.4.15-fr]{1.4.15})에서 따른다.
\end{proof}

다음 따름정리는 차원 이론의 용어 (제IV장)를 사용하며 제IV장 전에는
적용하지 않는다.

\begin{corollary}[4.2.2]
\label{III.4.2.2-fr}
$Y$를 국소 뇌터 준스킴, $f:X\to Y$를 고유 사상, $y$를 $Y$의 점,
$r$을 $f^{-1}(y)$의 차원이라 하자. 그러면 모든 연접
$\mathcal O_X$-가군 $\mathcal F$에 대해, 모든 $p>r$에서 층
$R^pf_*(\mathcal F)$는 $y$의 어떤 근방에서 영이다.
\end{corollary}

\begin{proof}
실제로 모든 $k$에 대해
\[
H^p\bigl(f^{-1}(y),
\mathcal F\otimes_{\mathcal O_Y}(\mathcal O_y/\mathfrak m_y^k)\bigr)=0
\]
이다 (G, II, 4.15.2). 따라서
(\hyperref[III.4.2.1-fr]{4.2.1})에 의해
$(R^pf_*(\mathcal F))_y$의 $\mathfrak m_y$-준아딕 위상에 관한 분리
완비화는 영이다. 이 위상은 분리이므로
$(R^pf_*(\mathcal F))_y=0$이다. $R^pf_*(\mathcal F)$가 연접이므로
(0\textsubscript{I}, 5.2.2), 결론이 따른다.
\end{proof}

\begin{env}[4.2.3]
\label{III.4.2.3-fr}
결과 (\hyperref[III.4.2.1-fr]{4.2.1})은 주로 $p=0$일 때 사용된다.
따라서 다음 따름정리를 얻는다.
\end{env}

\begin{corollary}[4.2.4]
\label{III.4.2.4-fr}
(\hyperref[III.4.2.1-fr]{4.2.1})의 가정 아래에서 표준 위상 동형사상
\[
\widehat{(f_*(\mathcal F))_y}
\xrightarrow{\sim}
\varprojlim_k\Gamma\bigl(f^{-1}(y),
\mathcal F_y/\mathfrak m_y^k\mathcal F_y\bigr)
\]
이 있다.
\end{corollary}

\subsection{자리스키의 연결성 정리}
\label{subsection:III.4.3-fr}
\label{III.4.3-fr}

이 번호와 다음 번호의 결과들은 자리스키의 잘 알려진 정리들을
일반화하며, 모두 (\hyperref[III.4.2.4-fr]{4.2.4})에서 이끌어낼 수 있다.
이들은 다음 정리의 귀결이다.

\begin{theorem}[4.3.1]
\label{III.4.3.1-fr}
\textnormal{(연결성 정리).}
$Y$를 국소 뇌터 준스킴, $f:X\to Y$를 고유 사상이라 하자. 그러면
\[
\mathcal A(X)=f_*(\mathcal O_X)
\]
는 연접 $\mathcal O_Y$-대수이다.

\oldpage[III]{131}

$\mathcal A(Y')=\mathcal A(X)$이고 $Y$ 위에서 유한인 $Y$-스킴
$Y'$를 택하자. 이는 $Y$-동형사상을 제외하면 유일하게 결정된다
(II, 1.3.1 및 6.1.3). 항등 동형사상
$e:\mathcal A(Y')\to\mathcal A(X)$에서 얻는 $Y$-사상
$f'=\mathcal A(e):X\to Y'$ (II, 1.2.7)는 고유이고,
$f'_*(\mathcal O_X)$는 $\mathcal O_{Y'}$와 동형이며, 사상 $f'$의
모든 올 $f'^{-1}(y')$는 연결이고 공집합이 아니다.
\end{theorem}

\begin{proof}
$g:Y'\to Y$를 구조 사상이라 하자. $f'$의 정의에 들어가는 준동형사상
\[
\theta:\mathcal O_{Y'}\longrightarrow f'_*(\mathcal O_X)
\]
이 전단사임을 보이자. $Y'$는 $Y$ 위에서 아핀이므로,
\[
g_*(\theta):g_*(\mathcal O_{Y'})\longrightarrow
g_*(f'_*(\mathcal O_X))=f_*(\mathcal O_X)
\]
이 항등사상임을 보이면 충분하다 (II, 1.4.2). 그런데 정의에 의해
$g_*(\mathcal O_{Y'})=\mathcal A(Y')$이고
$f_*(\mathcal O_X)=\mathcal A(X)$이므로 그러하다.
$\mathcal A(X)$가 연접이라는 사실은 유한성 정리
(\hyperref[III.3.2.1-fr]{3.2.1})의 특수한 경우이다. $f$는 고유이고
$g$는 분리이므로 $f'$은 고유이다 (II, 5.4.3, (i)). 따라서
(\hyperref[III.4.3.1-fr]{4.3.1})의 증명을 끝내려면 다음 따름정리를
증명하면 충분하다.
\end{proof}

\begin{corollary}[4.3.2]
\label{III.4.3.2-fr}
(\hyperref[III.4.3.1-fr]{4.3.1})의 가정에 더하여
$f_*(\mathcal O_X)\simeq\mathcal O_Y$라고 가정하자. 그러면 모든
$y\in Y$에 대해 $f$의 올 $f^{-1}(y)$는 연결이고 공집합이 아니다.
\end{corollary}

\begin{proof}
$f_*(\mathcal O_X)\simeq\mathcal O_Y$라는 가정에서 이미 $f$가 우세임이
따른다. $f$는 닫힌 사상이므로 전사이다.
(\hyperref[III.4.2.1-fr]{4.2.1})에서와 같이 $y$가 $Y$의 닫힌점인
경우로 귀착시킬 수 있다. $f^{-1}(y)$는 뇌터 공간이므로 유한 개의 연결
성분을 가지며, $f^{-1}(y)$를 따른 완비화 $\widehat X$의 바탕공간이다.
이 연결 성분들을 $Z_i$ ($1\leq i\leq n$)라 하자. 그러면
$\Gamma(\widehat X,\mathcal O_{\widehat X})$는 환들
$\Gamma(Z_i,\mathcal O_{\widehat X})$의 직접곱이고, 각 환은 영환이
아니다. 실제로 $\widehat X$의 모든 점에서 단위원 단면은 영과 다르다.
한편 $\mathcal F=\mathcal O_X$에
(\hyperref[III.4.1.5-fr]{4.1.5})를 적용하면, $f^{-1}(y)$를 따른 그
완비화가 $\mathcal O_{\widehat X}$이므로
$\Gamma(\widehat X,\mathcal O_{\widehat X})$는 국소환
$\mathcal O_y$의 $\mathfrak m_y$-아딕 분리 완비화
$\widehat{\mathcal O}_y$와 동형이다. 따라서 이는 국소환이며, 여러
영이 아닌 환들의 직접곱일 수 없다. 그렇지 않으면 서로 다른 극대
아이디얼들을 여럿 가질 것이기 때문이다. 그러므로 $n=1$이고,
따름정리가 증명되었다.
\end{proof}

\begin{corollary}[4.3.3]
\label{III.4.3.3-fr}
(\hyperref[III.4.3.1-fr]{4.3.1})의 가정 아래에서 모든 $y\in Y$에 대해
올 $f^{-1}(y)$의 연결 성분들의 집합은 구조 사상 $g:Y'\to Y$의 올
$g^{-1}(y)$의 유한한 점 집합과 전단사로 대응한다. 다시 말해 이는
$(f_*(\mathcal O_X))_y$의 극대 아이디얼들의 집합이다.
\end{corollary}

\begin{proof}
$Y'$는 $Y$ 위에서 유한이므로 $g^{-1}(y)$는 유한 이산 공간이다
(II, 6.1.7). 또한 $f^{-1}(y)=f'^{-1}(g^{-1}(y))$이므로, 이 관찰과
(\hyperref[III.4.3.1-fr]{4.3.1})에서 따름정리가 따른다.
\end{proof}

이로써 (\hyperref[III.4.3.1-fr]{4.3.1})에서 정의한 $Y$-준스킴
$Y'$에 대한 주목할 만한 해석을 얻었다. 고유 사상 $f$의 인수분해
$f=g\circ f'$은 해석공간의 정칙 사상에 대해 K.~Stein이 얻은
인수분해와 유사하다. 앞으로 이를 $f$의 \emph{스테인 인자분해}라
부르겠다.

\begin{remark}[4.3.4]
\label{III.4.3.4-fr}
$k$를 체 $k(y)$의 확대라 하자. 준스킴
\[
f^{-1}(y)\otimes_{k(y)}k=X\times_Y\operatorname{Spec}(k)
\]
가 연결이면, 그 사영 사상에 의한 상인 $f^{-1}(y)$도 연결이다
(I, 3.4.7). 준스킴의 사상 $f:X\to Y$와 점 $y\in Y$에 대해, 모든
$k(y)$의 확대 $k$에 대하여 준스킴
\[
f^{-1}(y)\otimes_{k(y)}k=X\times_Y\operatorname{Spec}(k)
\]
가 연결일 때 올 $f^{-1}(y)$가 \emph{기하적으로 연결}이라고 한다.

\oldpage[III]{132}

(\hyperref[III.4.3.2-fr]{4.3.2})의 가정 아래에서는 결론을 강화할 수
있다. 실제로 올 $f^{-1}(y)$들은 \emph{기하적으로 연결}이다. 이를
보이기 위해, 모든 $k(y)$의 확대 $k$에 대하여 어떤 뇌터 국소환 $A$와
국소 준동형사상 $\varphi:\mathcal O_y\to A$가 존재하여 $A$가 평탄
$\mathcal O_y$-가군이고 $A$의 잉여체가 $k$와 $k(y)$-동형임을
유의하자 (0, 10.3.1). $Y_1=\operatorname{Spec}(A)$라 하고,
$h:Y_1\to Y$를 $\varphi$에 대응하며 $Y_1$의 유일한 닫힌점 $y_1$을
$y$로 보내는 국소 사상이라 하자 (I, 2.4.1). 또한
\[
X_1=X\times_YY_1,\qquad f_1=f\times1_{Y_1}
\]
로 놓자. $f_1$은 고유이고 (II, 5.4.2, (iii)),
$f_1^{-1}(y_1)$은 $X\times_Y\operatorname{Spec}(k)$와 동형인
$k(y_1)$-준스킴이다. 따라서
(\hyperref[III.4.3.2-fr]{4.3.2})를 $f_1$에 적용하기 위해
$f_{1*}(\mathcal O_{X_1})=\mathcal O_{Y_1}$임만 보이면 된다.
$h$는 평탄 사상이다 (I, 2.4.2 및
\hyperref[III.1.4.15-fr]{1.4.15.5}). 그러므로 $q=0$에 대해
(\hyperref[III.1.4.15-fr]{1.4.15})를 적용하면
\[
f_{1*}(\mathcal O_{X_1})
=h^*(f_*(\mathcal O_X))
=h^*(\mathcal O_Y)
=\mathcal O_{Y_1}
\]
이다.

일반적인 경우 (\hyperref[III.4.3.1-fr]{4.3.1})에도 같은 논증은 그
정리의 표기로
\[
f_{1*}(\mathcal O_{X_1})=h^*(g_*(\mathcal O_{Y'}))
\]
임을 보여 준다. $f_1$의 스테인 인자분해
$f_1=g_1\circ f'_1$에서
\[
g_1=g\times1_{Y_1}
\]
이고 (II, 1.5.2), 대응하는 유한 $Y_1$-스킴은
$Y'_1=Y'\times_YY_1$이다. 올의 추이성을 고려하면 (I, 3.6.4),
(\hyperref[III.4.3.3-fr]{4.3.3})에 의해 $f_1^{-1}(y_1)$의 연결
성분 수는
\[
g_1^{-1}(y_1)=g^{-1}(y)\otimes_{k(y)}k
\]
의 원소 수와 같다. $k$를 $k(y)$의 대수적으로 닫힌 확대체로 택하면,
이 수는 선택한 대수적으로 닫힌 확대체에 의존하지 않으며
$g^{-1}(y)$의 \emph{기하적 점의 수} (I, 6.4.7), 즉
\emph{분리차수들의 합}
\[
\sum_i[k(y'_i):k(y)]_s
\]
과 같다. 여기서 $y'_i$는 유한집합 $g^{-1}(y)$를 달린다. 이 수를
$f^{-1}(y)$의 \emph{기하적 연결 성분의 수}라고도 한다. 체
$k(y'_i)$들은 반국소환 $(f_*(\mathcal O_X))_y$의 잉여체들과
다르지 않다.
\end{remark}

\begin{proposition}[4.3.5]
\label{III.4.3.5-fr}
$X$, $Y$를 국소 뇌터 정역 준스킴들, $f:X\to Y$를 고유 우세
사상이라 하자. 모든 $y\in Y$에 대해 $f^{-1}(y)$의 연결 성분 수는
유리함수체 $\mathrm R(X)$ 안에서 $\mathcal O_y$의 정수적 폐포
$\mathcal O'_y$의 극대 아이디얼 수 이하이다.
\end{proposition}

\begin{proof}
실제로 $Y$의 모든 열린집합 $U$에 대해
\[
\Gamma(U,f_*(\mathcal O_X))
=\Gamma(f^{-1}(U),\mathcal O_X)
\]
는 $x\in f^{-1}(U)$인 국소환 $\mathcal O_x$들의 교집합이다
(I, 8.2.1.1). 따라서 $(f_*(\mathcal O_X))_y$는
$\mathcal O_y$를 포함하는 $\mathrm R(X)$의 부분환이다. 또한
$f_*(\mathcal O_X)$는 연접 $\mathcal O_Y$-가군이므로
$(f_*(\mathcal O_X))_y$는 유한 생성 $\mathcal O_y$-가군이고,
따라서 $\mathcal O'_y$에 포함된다. 이러한 환 $A$의 모든 극대
아이디얼이 $A$와 $\mathcal O'_y$의 어떤 극대 아이디얼의 교집합임은
알려져 있다 ([13], vol. I, p. 257 및 259). 여기서 명제가 따른다.
\end{proof}

\begin{definition}[4.3.6]
\label{III.4.3.6-fr}
정역 국소환의 정수적 폐포가 국소환이면 그 국소환을
\emph{단일분기(unibranche)}라고 한다. 정역 준스킴 $Y$의 점 $y$에서
국소환 $\mathcal O_y$가 단일분기이면 $y$가 단일분기라고 한다.
특히 $Y$가 $y$에서 정규이면 그러하다.
\end{definition}

$A$를 정역 국소환, $K$를 그 분수체라 하자. $A$가 단일분기이기 위한
필요충분조건은 $A$를 포함하고 유한 $A$-대수인 $K$의 모든 부분환
$A_1$이 국소환인 것이다. 실제로 $A'$을 $A$의 정수적 폐포라 하자.
첫째 Cohen--Seidenberg 정리에 의해 (Bourbaki,
\emph{Alg. comm.}, chap. V, \S~2, n\textsuperscript{o}~1, th.~1),
$A_1$의 모든 극대 아이디얼은 $A'$의 어떤 극대 아이디얼의 자취이다.
따라서 $A'$이 국소이면 $A_1$도 국소이다. 역으로 $A'$은 그 안의 유한
부분 $A$-대수 $A_\alpha$들로 이루어진 증가 여과족의 귀납극한이다.

\oldpage[III]{133}

각 $A_\alpha$가 국소환이면, 위와 같은 논증에 의해
$A_\alpha\subset A_\beta$일 때 $A_\alpha$의 극대 아이디얼은
$A_\beta$의 극대 아이디얼이 $A_\alpha$ 위에 남기는 자취이다. 따라서
$A'$은 국소환이다 (0, 10.3.1.3).

뇌터 국소환 $A$의 완비화가 정역이면, 즉 $A$가
\emph{해석적으로 정역}이면, $A$는 단일분기임도 유의하자. 실제로
$\mathfrak m$을 $A$의 극대 아이디얼, $K$를 그 분수체,
$K'$을 $\widehat A$의 분수체라 하면
\[
K'=K\otimes_A\widehat A.
\]
$A'_F$를 $K$ 안의 유한 부분 $A$-대수라 하자. $\widehat A$와
$A'_F$로 생성되는 $K'$의 부분환 $B_F$는
$A'_F\otimes_A\widehat A$와 동형이다. 이는 유한 생성
$\widehat A$-가군이며, $A'_F$의 $\mathfrak m$-아딕 위상에 관한
완비화이다 (0\textsubscript{I}, 7.3.3 및 7.3.6).
$A'_F$는 반국소환이고 (Bourbaki, \emph{Alg. comm.}, chap. IV,
\S~2, n\textsuperscript{o}~5, prop. 9의 cor. 3), 그 완비화가
정역이므로 $A'_F$는 오직 하나의 극대 아이디얼 $\mathfrak m'_F$만
가질 수 있으며 $\mathfrak m'_F\cap A=\mathfrak m$이다. 여기서
주장이 따른다.

\begin{corollary}[4.3.7]
\label{III.4.3.7-fr}
(\hyperref[III.4.3.5-fr]{4.3.5})의 가정 아래에서
$\mathrm R(X)$ 안의 $\mathrm R(Y)$의 대수적 폐포의 분리차수가
$n$이고 $y\in Y$가 단일분기라고 가정하자. 그러면 올 $f^{-1}(y)$는
연결 성분을 많아야 $n$개 가진다. 특히 $\mathrm R(X)$ 안의
$\mathrm R(Y)$의 대수적 폐포가 $\mathrm R(Y)$ 위에서 순수 비분리이면
$f^{-1}(y)$는 연결이다.
\end{corollary}

\begin{proof}
$\mathcal O'_y$를 $\mathcal O_y$의 정수적 폐포라 하자.
$\mathrm R(X)$ 안의 $\mathcal O_y$의 정수적 폐포
$\mathcal O''_y$는 $\mathcal O'_y$의 정수적 폐포이기도 하다.
$\mathcal O'_y$가 국소환이면 $\mathcal O''_y$는 반국소환이고 그 극대
아이디얼 수는 $n$ 이하임이 알려져 있다
([13], vol. I, p. 289, th. 22).

이 따름정리는 자리스키가 대수적 스킴에 대한 자신의
‘연결성 정리’를 서술한 형태와 본질적으로 같다.
\end{proof}

\begin{remark}[4.3.8]
\label{III.4.3.8-fr}
(\hyperref[III.4.3.7-fr]{4.3.7})의 가정에 더하여 $Y$가 $y$에서
정규라고 가정하면, 올 $f^{-1}(y)$는 \emph{기하적으로 연결}이다.
실제로 (\hyperref[III.4.3.4-fr]{4.3.4})의 표기로 $g^{-1}(y)$는 한
점 $y'$로 이루어지고 $k(y')$는 $k(y)$ 위에서 순수 비분리이다.
\end{remark}

\begin{definition}[4.3.9]
\label{III.4.3.9-fr}
$Y$를 국소 뇌터 준스킴이라 하자. 모든 기약 국소 뇌터 준스킴
$Y'$와 모든 우세 사상 $g:Y'\to Y$에 대해
$X'=X\times_YY'$의 모든 기약 성분이 $Y'$ 위에서 우세이면, 유한형
사상 $f:X\to Y$를 \emph{보편적으로 열린 사상}이라고 한다.
\end{definition}

$Y$가 기약인 경우에는 다음 조건과 동치이다. $Y$, $Y'$의 일반점을
각각 $\eta$, $\eta'$라 하고 $g(\eta')=\eta$라 하며
$f'=f_{(Y')}$로 놓으면, $X'$의 모든 기약 성분이
$f'^{-1}(\eta')$와 만난다 (0\textsubscript{I}, 2.1.8). 따라서
$Y$의 모든 열린집합 $U$에 대해 $f$의 제한
$f^{-1}(U)\to U$도 보편적으로 열린 사상이다.

\begin{corollary}[4.3.10]
\label{III.4.3.10-fr}
$X$, $Y$를 국소 뇌터 정역 준스킴들이라 하고, $f:X\to Y$를 고유하고
우세이며 보편적으로 열린 사상이라 하자. $\mathrm R(X)$ 안의
$\mathrm R(Y)$의 대수적 폐포가 $\mathrm R(Y)$ 위에서 순수
비분리이면, 모든 올 $f^{-1}(y)$ ($y\in Y$)는 기하적으로 연결이다.
\end{corollary}

\begin{proof}
$Y=\operatorname{Spec}(B)$이고 $B$가 뇌터 정역인 경우로 한정할 수
있다. (II, 7.1.7)에 의해 $\mathcal O_y$를 우세 지배하고 분수체가
$\mathrm R(Y)$인 정수적으로 닫힌 뇌터 국소환 $A$가 존재한다.
$Y'=\operatorname{Spec}(A)$라 하고, 표준 단사 $B\to A$에 대응하는
쌍유리 사상 $h:Y'\to Y$를 생각하자. 이는 특히 우세이다. $Y'$의
유일한 닫힌점을 $y'$라 하면 $h(y')=y$이다. 다음과 같이 놓자.
\[
X'=X\times_YY',\qquad f'=f\times1_{Y'}.
\]
$Y$, $Y'$, $X$의 일반점을 각각 $\eta$, $\eta'$, $\xi$라 하자.
그러면 $f(\xi)=\eta$이고 $h^{-1}(\eta)=\{\eta'\}$이다. 또한
$k(\eta)=k(\eta')=\mathrm R(Y)$이므로
$f'^{-1}(\eta')\simeq f^{-1}(\eta)$이다 (I, 3.6.4).
$\xi$는 $f^{-1}(\eta)$의 일반점이므로
(0\textsubscript{I}, 2.1.8), $f'^{-1}(\eta')$는 일반점을 하나만
가진다. 그런데 가정에 의해 $X'$의 모든 기약 성분의 일반점은
$f'^{-1}(\eta')$ 안에 있다. 따라서 $X'$은 기약이고 그 일반점
$\xi'$은 $f'^{-1}(\eta')$의 일반점이며
$k(\xi')=k(\xi)$이다.

\oldpage[III]{134}

$X''=X'_{\mathrm{red}}$, $f''=f'_{\mathrm{red}}$로 놓자. 그러면
$X''$은 뇌터 정역 준스킴이고, $f''$은 고유이며 (II, 5.4.6),
$f'^{-1}(y')$와 $f''^{-1}(y')$의 바탕공간은 같다. 또한
\[
\mathrm R(X'')=k(\xi')=\mathrm R(X)
\]
이므로 $f''$은 (\hyperref[III.4.3.8-fr]{4.3.8})의 가정을 만족하고
$f''^{-1}(y')$는 기하적으로 연결이다. 이제 $k$를 $k(y)$의 임의의
확대라 하자. $k(y')$와 $k$를 모두 부분확대로 갖는 $k(y)$의 확대
$k_1$이 존재한다 (Bourbaki, \emph{Alg.}, chap. V, \S~4, prop. 2).
가정에 의해
\[
f''^{-1}(y')\times_{Y'}\operatorname{Spec}(k_1)
\]
은 연결이고, 이는
$f'^{-1}(y')\times_{Y'}\operatorname{Spec}(k_1)$과 같은 축약
준스킴을 가진다 (I, 5.1.8). 따라서 후자도 연결이다. 후자는
$f^{-1}(y)\times_Y\operatorname{Spec}(k_1)$과 동형이므로
(I, 3.6.4), 이 준스킴도 연결이다. 그러므로
(\hyperref[III.4.3.4-fr]{4.3.4}) 첫머리의 주의에 의해
\emph{더욱이} $f^{-1}(y)\times_Y\operatorname{Spec}(k)$도 연결이다.
이것으로 증명이 끝난다.
\end{proof}

\begin{env}[4.3.11]
\label{III.4.3.11-fr}
\emph{주의} (4.3.11). ---
\begin{enumerate}
\item[(i)]
앞의 논증은 본질적으로 자리스키 [20]의 것이다. 다만 고전 대수기하의
국소환에서는 $\mathcal O_y$의 정수적 폐포가 뇌터환이므로, 자리스키는
$A$로 그 정수적 폐포를 택할 수 있었다. 한편 자리스키는 $Y$가 체
$k$ 위의 사영공간 $\mathbf P_k^r$의 Chow 다양체이고, $X$가
$\mathbf P_k^r$와 $Y$ 사이의 Chow 대응을 정의하는
$\mathbf P_k^r\times_kY$의 닫힌 부분집합이면 사영
$X\to Y$가 보편적으로 열린 사상임을 증명하였다
(\emph{loc. cit.}, p. 82의 lemma). 어떤 응용들에서 쓰인
‘Chow 좌표’의 형식적 성질은 이것뿐인 것으로 보인다. 따라서 이러한
상황에서는 ‘사영공간 안의 순환의 특수화’라는 언어 대신 ‘고유 사상의
올’이라는 언어를 쓰는 것이 유익하다. 필요하면 그 사상이 보편적으로
열려 있다고 가정하거나, 국소 정칙성에 관한 다른 유사한 제한을
부과한다.

\item[(ii)]
제IV장에서 보편적으로 열린 사상 $f:X\to Y$를 다음과 같이 정의할
수 있음도 보겠다. 이 성질이 그 용어를 정당화한다. 모든 사상
$Z\to Y$에 대해 사상 $f_{(Z)}:X_{(Z)}\to Z$가 열린 사상이다. 또한
$f$가 (\hyperref[III.4.3.10-fr]{4.3.10})의 가정들을 만족하고
$y,y'$가 $Y$의 두 점이며 $y$가 $y'$의 특수화이면, $f^{-1}(y)$의
기하적 연결 성분 수는 $f^{-1}(y')$의 연결 성분 수 이하임을 보일 수
있다.
\end{enumerate}
\end{env}

\begin{corollary}[4.3.12]
\label{III.4.3.12-fr}
(\hyperref[III.4.3.5-fr]{4.3.5})의 가정 아래에서
$\mathrm R(Y)$가 $\mathrm R(X)$ 안에서 대수적으로 닫혀 있고 $y$가
$Y$의 정규점이라고 가정하자. 그러면 $f^{-1}(y)$는 기하적으로
연결이고, $Y$ 안의 $y$의 어떤 열린 근방 $U$가 존재하여
\[
f_*(\mathcal O_X|f^{-1}(U))\simeq\mathcal O_Y|U
\]
이다. 특히 $Y$가 정규이고 $\mathrm R(Y)$가 $\mathrm R(X)$ 안에서
대수적으로 닫혀 있으면
$f_*(\mathcal O_X)\simeq\mathcal O_Y$이다.
\end{corollary}

\begin{proof}
$f^{-1}(y)$에 관한 첫째 주장은
(\hyperref[III.4.3.8-fr]{4.3.8})의 특수한 경우이다. 따라서
$X\xrightarrow{f'}Y'\xrightarrow{g}Y$가 $f$의 스테인 인자분해이면
(\hyperref[III.4.3.3-fr]{4.3.3}), $g^{-1}(y)$는 한 점 $y'$로
이루어진다. 또한

\oldpage[III]{135}

\[
\mathcal O_y\subset\mathcal O_{y'}=(f_*(\mathcal O_X))_y
\subset\mathrm R(X).
\]
$\mathcal O_{y'}$는 $\mathcal O_y$ 위에서 유한이므로
$\mathrm R(Y)$ 위에서 대수적이고, 가정에 의해 $\mathrm R(Y)$에
포함된다. $y$는 정규이므로 반드시
$\mathcal O_{y'}=\mathcal O_y$이다. 따라서 $g$는 $y'$에서 국소
동형사상이고 (I, 6.5.4), 이것으로 따름정리의 첫째 부분이 증명된다.
둘째 부분은 첫째 부분에서 따른다. 실제로 추가 가정 아래에서 $g$는
전단사이고 $Y'$의 모든 점의 근방에서 국소 동형사상이므로
동형사상이다.
\end{proof}

(\hyperref[III.4.3.7-fr]{4.3.7})을 스킴의 틀에서 증명했기 때문에
다음과 같은 응용이 가능하다.

\begin{proposition}[4.3.13]
\label{III.4.3.13-fr}
$A$를 단일분기 뇌터 국소환, $\mathfrak a$를 $A$의 정의 아이디얼,
$A_0=A/\mathfrak a$,
\[
S=\operatorname{gr}_{\mathfrak a}(A)
\]
를 $\mathfrak a$-준아딕 여과에 관해 $A$에 대응하는 등급환이라 하자.
$S$는 $S_1$로 생성되는 등급 $A_0$-대수이고, $S_1$은 유한 생성
$A_0$-가군이다. 그러면 $\operatorname{Proj}(S)$는 연결
$A_0$-스킴이다.
\end{proposition}

\begin{proof}
$\mathfrak m$을 $A$의 극대 아이디얼이라 하자.
$Y=\operatorname{Spec}(A)$는 정역 스킴이고 $\mathfrak m$에 대응하는
점 $y$는 그 유일한 닫힌점이다. 가정에 의해 어떤 정수 $p$에 대해
\[
\mathfrak m^p\subset\mathfrak a\subset\mathfrak m
\]
이므로 $V(\mathfrak a)=\{\mathfrak m\}$이다.
$S'=\bigoplus_{n\geq0}\mathfrak a^n$이라 하고,
$X=\operatorname{Proj}(S')$라 하자. 이는 아이디얼 $\mathfrak a$를
부풀려 얻는 $Y$-스킴이다. $X$는 정역이고 구조 사상 $f:X\to Y$는
쌍유리이며 (II, 8.1.4) 명백히 사영적이다. 따라서
(\hyperref[III.4.3.7-fr]{4.3.7})을 적용할 수 있고
$f^{-1}(y)$는 연결이다. 그런데 $f^{-1}(y)$는
\[
\operatorname{Proj}(S'\otimes_AA_0)
\]
의 바탕공간이며 (I, 3.6.1 및 II, 2.8.10),
정의에 의해 $S'\otimes_AA_0=S$이다. 따라서 명제가 증명되었다.
\end{proof}

\subsection{자리스키의 “주정리”}
\label{subsection:III.4.4-fr}

\begin{proposition}[4.4.1]
\label{III.4.4.1-fr}
$Y$를 국소 뇌터 준스킴, $f:X\to Y$를 고유 사상이라 하자.
$X'$을 자신의 올 $f^{-1}(f(x))$ 안에서 고립된 점 $x\in X$들의
집합이라 하자. 그러면 $X'$은 $X$의 열린집합이다. 또한
$f=g\circ f'$이 $f$의 스테인 인자분해이면
(\hyperref[III.4.3.3-fr]{4.3.3}), $f'$의 $X'$에 대한 제한은
$X'$에서 $Y'$의 어떤 열린집합 $U$ 위에 유도된 부분준스킴으로 가는
동형사상이고 $X'=f'^{-1}(U)$이다.
\end{proposition}

\begin{proof}
$g^{-1}(f(x))$는 유한 이산 공간이므로
(\hyperref[III.4.3.3-fr]{4.3.3} 및 II, 6.1.7), $x$가
$f^{-1}(f(x))$ 안에서 고립되기 위한 필요충분조건은
$f'^{-1}(f'(x))$ 안에서 고립되는 것이다. 따라서 $f'=f$, 즉
$f_*(\mathcal O_X)=\mathcal O_Y$인 경우로 한정할 수 있다.
$x\in X'$이면 연결공간 $f^{-1}(f(x))$
(\hyperref[III.4.3.2-fr]{4.3.2})은 반드시 한 점 $x$로 이루어진다.
$f$는 닫힌 사상이므로, $X$ 안의 $x$의 모든 열린 근방 $V$에 대해
$f(X-V)$는 $Y$의 닫힌집합이다. 이는 $y=f(x)$를 포함하지 않는다.
실제로 $f^{-1}(y)=\{x\}$이다. $U$를 $Y$에서 $f(X-V)$의 여집합이라
하면 $f^{-1}(U)\subset V$이다. 따라서 $y$의 열린 근방 기본계의
$f$에 의한 역상들은 $x$의 열린 근방 기본계를 이룬다.

$f_*(\mathcal O_X)=\mathcal O_Y$라는 가정과 층의 직상 정의
(0\textsubscript{I}, 3.4.1 및 4.2.1)에 의해, $f=(\psi,\theta)$로
쓰면 준동형사상
\[
\theta_y^\#:\mathcal O_y\longrightarrow\mathcal O_x
\]
는 동형사상이다. 따라서 $x$의 열린 근방 $V$와 $y$의 열린 근방
$U$가 존재하여 $f$의 $V$에 대한 제한이 $V$에서 $U$로 가는
동형사상이다 (I, 6.5.4). 더 나아가 앞의 관찰에 의해

\oldpage[III]{136}

$f^{-1}(U)=V$라고 가정할 수 있다. 그러면 정의에서 즉시
$V\subset X'$임이 따르며, 증명이 끝난다.
\end{proof}

다음 명제는 대수적 스킴의 경우에 Chevalley가 증명하였다.

\begin{proposition}[4.4.2]
\label{III.4.4.2-fr}
$Y$를 국소 뇌터 준스킴, $f:X\to Y$를 사상이라 하자. 다음 조건들은
서로 동치이다.
\begin{enumerate}
\item[a)] $f$는 유한 사상이다.
\item[b)] $f$는 아핀이고 고유이다.
\item[c)] $f$는 고유이고, 모든 $y\in Y$에 대해 $f^{-1}(y)$는
유한집합이다.
\end{enumerate}
\end{proposition}

\begin{proof}
a)가 b)를 함의함은 알려져 있다 (II, 6.1.2 및 6.1.11). $f$가
고유이고 아핀이면
\[
f^{-1}(y)\longrightarrow\operatorname{Spec}(k(y))
\]
도 그러하다 (II, 1.6.2, (iii) 및 5.4.2, (iii)). $f^{-1}(y)$의
구조층에 유한성 정리 (\hyperref[III.3.2.1-fr]{3.2.1})를 적용하면
$f^{-1}(y)=\operatorname{Spec}(A)$이고 $A$는 유한 $k(y)$-대수이다.
따라서 $f^{-1}(y)$는 유한집합이다 (II, 6.1.7). 이로써 b)가 c)를
함의한다.

끝으로 $f^{-1}(y)$는 $k(y)$ 위의 대수적 준스킴이므로, 집합
$f^{-1}(y)$가 유한이라는 가정에서 공간 $f^{-1}(y)$가 이산임이
따른다 (I, 6.4.4). (\hyperref[III.4.4.1-fr]{4.4.1})의 표기로
$X'=X$이고 $f':X\to Y'$은 동형사상이다. $g$는 유한 사상이므로
c)가 a)를 함의한다.
\end{proof}

\begin{theorem}[4.4.3]
\label{III.4.4.3-fr}
\textnormal{(자리스키의 “주정리”).}
$Y$를 뇌터 준스킴, $f:X\to Y$를 준사영 사상, $X'$을 자신의 올
$f^{-1}(f(x))$ 안에서 고립된 점 $x\in X$들의 집합이라 하자. 그러면
$X'$은 $X$의 열린부분집합이고, 유도 부분준스킴 $X'$은 $Y$ 위에서
유한인 어떤 $Y$-준스킴 $Y'$의 열린부분집합 위에 유도된 준스킴과
동형이다.
\end{theorem}

\begin{proof}
가정에 의해 어떤 사영적 $Y$-준스킴 $Z$가 존재하여 $X$는 $Z$의
어떤 열린집합 위에 유도된 부분준스킴과 $Y$-동형이다
(II, 5.3.2 및 5.5.1). 따라서 $f$가 사영 사상인 경우로 귀착된다.
이때 $f$는 고유이고 (II, 5.5.3), 정리는
(\hyperref[III.4.4.1-fr]{4.4.1})에서 즉시 따른다.
\end{proof}

\begin{remark}[4.4.4]
\label{III.4.4.4-fr}
$X$가 축약이면 (각각 $X$가 기약이고 $X'$이 공집합이 아니면),
(\hyperref[III.4.4.3-fr]{4.4.3})의 명제에서 $Y'$이 축약이라고
(각각 기약이라고) 가정할 수 있다. 실제로 $Y'$을 언제나 $Y'$ 안의
$X'$의 폐포 부분준스킴 $\overline{X'}$로 바꿀 수 있다
(I, 9.5.11 및 II, 6.1.5, (i), (ii)). $X'$이 축약이면
$\overline{X'}$도 축약임이 알려져 있다 (I, 9.5.9, (i)). 또한
$X'$이 공집합이 아니고 $X$가 기약이면 $X'$도 기약이며,
$\overline{X'}$도 기약이다.
\end{remark}

\begin{corollary}[4.4.5]
\label{III.4.4.5-fr}
$Y$를 국소 뇌터 스킴, $f:X\to Y$를 유한형 사상, $x$를 자신의 올
$f^{-1}(f(x))$ 안에서 고립된 $X$의 점이라 하자. 그러면 $X$ 안의
$x$의 어떤 열린 근방은 $Y$ 위에서 유한인 어떤 $Y$-준스킴의
열린부분집합과 동형이다.
\end{corollary}

\begin{proof}
$y=f(x)$라 하고, $U$를 $Y$ 안의 $y$의 아핀 열린 근방,
$V$를 $f^{-1}(U)$에 포함되는 $X$ 안의 $x$의 아핀 열린 근방이라
하자. $Y$는 분리이므로 단사 $U\to Y$는 아핀 사상이다 (II, 1.6.3).
$V$도 $U$ 위에서 아핀이므로 (\emph{ibid.}), $f$의 $V$에 대한
제한은 아핀 사상 $V\to Y$이다 (II, 1.6.2, (ii)). 이 제한은
유한형이므로 \emph{더욱이} 준사영 사상이다
(I, 6.3.5 및 II, 5.3.4, (i)). 이제 이 제한에
(\hyperref[III.4.4.3-fr]{4.4.3})을 적용하면 된다.
\end{proof}

따름정리 (\hyperref[III.4.4.5-fr]{4.4.5})를 가환대수의 언어로
서술하면 다음과 같다.

\oldpage[III]{137}

\begin{corollary}[4.4.6]
\label{III.4.4.6-fr}
$A$를 뇌터환, $B$를 유한 생성 $A$-대수, $\mathfrak q$를 $B$의
소 아이디얼, $\mathfrak p$를 그 $A$ 안의 역상이라 하자.
$\mathfrak q$가 역상이 $\mathfrak p$인 $B$의 소 아이디얼들의
집합에서 동시에 극대이고 극소라고 가정하자. 그러면 어떤
$g\in B-\mathfrak q$, 유한 $A$-대수 $A'$, 원소 $f'\in A'$이
존재하여 $A$-대수 $B_g$와 $A'_{f'}$이 동형이다.
\end{corollary}

\begin{proof}
$Y=\operatorname{Spec}(A)$와 $X=\operatorname{Spec}(B)$에
(\hyperref[III.4.4.5-fr]{4.4.5})를 적용하면 충분하다.
$\mathfrak q$에 대한 가정은 그것이 자신의 올 안에서 고립됨을 정확히
뜻한다 (I, 1.1.7).
\end{proof}

여기서 겉보기에는 덜 일반적인 다음 결과를 얻는다.

\begin{corollary}[4.4.7]
\label{III.4.4.7-fr}
$A$를 뇌터 국소환, $B$를 유한 생성 $A$-대수, $\mathfrak n$을
$A$ 안의 역상이 극대 아이디얼 $\mathfrak m$인 $B$의 소
아이디얼이라 하자. $\mathfrak n$이 $B$에서 극대이고, 역상이
$\mathfrak m$인 $B$의 소 아이디얼들의 집합에서 극소라고 가정하자.
이는 $B\mathfrak m$이 $\mathfrak n$에 대해 일차라는 뜻이기도 하다.
그러면 유한 $A$-대수 $A'$과 $A'$의 극대 아이디얼
$\mathfrak m'$이 존재하여, $\mathfrak m$은 $\mathfrak m'$의
$A$ 안의 역상이고 $B_{\mathfrak n}$은 $A$-대수
$A'_{\mathfrak m'}$과 동형이다.
\end{corollary}

다음은 (\hyperref[III.4.4.7-fr]{4.4.7})의 특수한 경우이며,
때때로 “주정리”라고도 부른다.

\begin{corollary}[4.4.8]
\label{III.4.4.8-fr}
(\hyperref[III.4.4.7-fr]{4.4.7})의 조건 아래에서 $A$, $B$가
정역이고 같은 분수체 $K$를 가진다고 더 가정하자. $A$가 정수적으로
닫혀 있으면 $B=A$이다.
\end{corollary}

\begin{proof}
주의 (\hyperref[III.4.4.4-fr]{4.4.4})에 의해
(\hyperref[III.4.4.7-fr]{4.4.7})을 적용할 때 $A'$이 정역이고
분수체가 $K$라고 가정할 수 있다. $A$에 대한 가정에서
$A'=A$가 따르므로 $B_{\mathfrak n}=A$이다. 이제
\[
A\subset B\subset B_{\mathfrak n}
\]
이므로 $A=B$이다.
\end{proof}

명제 (\hyperref[III.4.4.8-fr]{4.4.8})은 자리스키가 자신의
“주정리”에 부여한 형태이다. 여기서는 임의의 뇌터 정역 국소환으로
확장하였다.

앞의 따름정리들은 전역적 결과인
(\hyperref[III.4.4.3-fr]{4.4.3})의 국소적 변형들이었다. 다음은
또 다른 전역적 귀결이다.

\begin{corollary}[4.4.9]
\label{III.4.4.9-fr}
$X$, $Y$를 국소 뇌터 정역 준스킴들, $f:X\to Y$를 분리이고 국소
유한형인 쌍유리 사상이라 하자. 더 나아가 $Y$가 정규이고 모든
$y\in Y$에 대해 올 $f^{-1}(y)$가 유한이라고 가정하자. 그러면
$f$는 열린 몰입 사상이다. 더 나아가 $f$가 닫힌 사상이면
(특히 $f$가 고유이면) $f$는 동형사상이다.
\end{corollary}

\begin{proof}
$x\in X$라 하고 $y=f(x)$로 놓자. $f^{-1}(y)$는 $k(y)$ 위에서
국소 유한형인 스킴이므로, 유한이라는 가정에서 이산임이 따른다
(I, 6.4.4). 또한 $\mathcal O_y$는 정수적으로 닫혀 있고
$\mathcal O_x$, $\mathcal O_y$는 같은 분수체를 가진다 (I, 7.1.5).
따라서 (\hyperref[III.4.4.8-fr]{4.4.8})을 적용할 수 있다.
$f=(\psi,\theta)$로 쓰면
\[
\theta_y^\#:\mathcal O_y\longrightarrow\mathcal O_x
\]
는 전단사이다. 그러므로 $f$는 국소 동형사상이다 (I, 6.5.4).
그런데 $f$는 분리이고 $X$는 정역이므로 $f$는 열린 몰입 사상이다
(I, 8.2.8). 마지막 주장은 $f$가 우세라는 사실에서 따른다.
\end{proof}

\begin{proposition}[4.4.10]
\label{III.4.4.10-fr}
$Y$를 국소 뇌터 준스킴, $f:X\to Y$를 국소 유한형 사상이라 하자.
자신의 올 $f^{-1}(f(x))$ 안에서 고립된 점 $x\in X$들의 집합 $X'$은
$X$에서 열린집합이다.
\end{proposition}

\begin{proof}
문제는 $X$와 $Y$ 위에서 국소적이므로 $X$, $Y$가 뇌터 아핀이고
$f$가 유한형이라고 가정할 수 있다. 이때 $f$는 유한형 아핀
사상이므로 준사영이다 (II, 5.3.4, (i)). 이제
(\hyperref[III.4.4.3-fr]{4.4.3})을 적용하면 된다.
\end{proof}

\oldpage[III]{138}

\begin{corollary}[4.4.11]
\label{III.4.4.11-fr}
$Y$를 국소 뇌터 준스킴, $f:X\to Y$를 고유 사상이라 하자.
$f^{-1}(y)$가 이산인 점 $y\in Y$들의 집합 $U$는 $Y$에서 열려 있고,
$f$의 제한 $f^{-1}(U)\to U$는 유한 사상이다. 특히 고유 준유한
사상 $X\to Y$는 유한 사상이다.
\end{corollary}

\begin{proof}
$Y$에서 $U$의 여집합은 $X-X'$의 $f$에 의한 상이다.
(\hyperref[III.4.4.10-fr]{4.4.10})에 의해 $X-X'$은 $X$에서
닫혀 있고 $f$는 닫힌 사상이므로 $U$는 열려 있다. 또한
(II, 6.2.2)에 의해 모든 $y\in U$에 대해 $f^{-1}(y)$는 유한하다.
$f$의 제한 $f^{-1}(U)\to U$는 고유이므로 (II, 5.4.1),
(\hyperref[III.4.4.2-fr]{4.4.2})에 의해 유한이다.
\end{proof}

\begin{env}[4.4.12]
\label{III.4.4.12-fr}
\emph{주의} (4.4.12). ---
\begin{enumerate}
\item[(i)]
(II, 6.2.7)에서 예고했듯이, 제V장에서 $Y$가 국소 뇌터이면 모든
준유한 분리 사상 $f:X\to Y$가 준아핀이고 따라서 준사영임을 보일
것이다. 그러면 주정리 (\hyperref[III.4.4.3-fr]{4.4.3})에서 $f$가
분리이고 유한형이라고만 가정해도 결론이 성립한다. 실제로
(\hyperref[III.4.4.10-fr]{4.4.10})에 의해 $X'$은 $X$에서 열려 있다.
$X$가 국소 뇌터이므로 $f$의 $X'$에 대한 제한은 여전히 유한형이고
(I, 6.3.5), $X'$의 정의에 의해 준유한이며 명백히 분리이다. 따라서
이 제한에 (\hyperref[III.4.4.3-fr]{4.4.3})을 적용하면 결론이 따른다.

\item[(ii)]
제IV장에서 차원 이론을 이용하여
(\hyperref[III.4.4.10-fr]{4.4.10})의 더 초등적인 증명을 제시하겠다.
\end{enumerate}
\end{env}

\subsection{준동형사상 가군의 완비화}
\label{subsection:III.4.5-fr}

\begin{proposition}[4.5.1]
\label{III.4.5.1-fr}
$A$를 뇌터환, $\mathfrak J$를 $A$의 아이디얼, $X$를 유한형
$A$-준스킴이라 하고, $\mathcal F$, $\mathcal G$를 그 받침들의
교집합이 $Y=\operatorname{Spec}(A)$ 위에서 고유인 두 연접
$\mathcal O_X$-가군이라 하자 (II, 5.4.10). 그러면 모든 정수
$n\geq0$에 대해
$\operatorname{Ext}^n_{\mathcal O_X}(X;\mathcal F,\mathcal G)$는 유한 생성
$A$-가군이고, 그 $\mathfrak J$-준아딕 위상에 대한 분리 완비화는
(\hyperref[III.4.1.7-fr]{4.1.7})의 표기 아래에서
$\operatorname{Ext}^n_{\mathcal O_{\widehat X}}(\widehat X;
\widehat{\mathcal F},\widehat{\mathcal G})$와 표준적으로 동일시된다.
\end{proposition}

\begin{proof}
(T, 4.2)에 의해, 수렴 대상이
$\operatorname{Ext}_{\mathcal O_X}(X;\mathcal F,\mathcal G)$이고
$\mathrm E_2$ 항들이
\[
\mathrm E_2^{pq}=\mathrm H^p\bigl(X,
\mathcal{E}\!xt^q_{\mathcal O_X}(\mathcal F,\mathcal G)\bigr)
\]
로 주어지는 쌍정칙 스펙트럼 열
$\mathrm E(\mathcal F,\mathcal G)$가 존재한다. 층
$\mathcal{E}\!xt^q_{\mathcal O_X}(\mathcal F,\mathcal G)$는 연접
$\mathcal O_X$-가군이고 (0, 12.3.3), 그 받침은 $\mathcal F$와
$\mathcal G$의 받침들의 교집합에 포함되므로 (T, 4.2.2), $Y$ 위에서
고유이다 (II, 5.4.10). 따라서
(\hyperref[III.3.2.4-fr]{3.2.4})에 의해 $\mathrm E_2^{pq}$들은
유한 생성 $A$-가군이고, 따라서 (0, 11.1.8)에 의해 스펙트럼 열의
모든 항 $\mathrm E_r^{pq}$와 그 수렴 대상도 그러하다.

한편 $i:\widehat X\to X$가 표준 사상이면 $\widehat{\mathcal F}$와
$\widehat{\mathcal G}$는 각각 $i^*(\mathcal F)$와
$i^*(\mathcal G)$로 표준적으로 동일시되고, $i$는 평탄하다
(I, 10.8.8 및 10.8.9). 그러면 (0, 12.3.4)에 의해 모든 $q\geq0$에
대해 표준 $i$-사상
\[
u_q:\mathcal{E}\!xt^q_{\mathcal O_X}(\mathcal F,\mathcal G)
\longrightarrow
\mathcal{E}\!xt^q_{\mathcal O_{\widehat X}}
(\widehat{\mathcal F},\widehat{\mathcal G})
\]
이 존재하고, 이에 대응하는 $\mathcal O_{\widehat X}$-준동형사상
\[
u_q^\#:i^*\bigl(\mathcal{E}\!xt^q_{\mathcal O_X}
(\mathcal F,\mathcal G)\bigr)
\longrightarrow
\mathcal{E}\!xt^q_{\mathcal O_{\widehat X}}
(\widehat{\mathcal F},\widehat{\mathcal G})
\]
은 동형사상이다 (0, 12.3.5). 다시 말해 (I, 10.8.8),
$\mathcal{E}\!xt^q_{\mathcal O_{\widehat X}}
(\widehat{\mathcal F},\widehat{\mathcal G})$는
(\hyperref[III.4.1.7-fr]{4.1.7})의 표기 아래에서 완비화
$(\mathcal{E}\!xt^q_{\mathcal O_X}(\mathcal F,\mathcal G))^\wedge$와
표준적으로 동일시된다. 이제 비교 정리
(\hyperref[III.4.1.10-fr]{4.1.10})에 의해 모든 $p\geq0$에 대해
\[
\mathrm H^p\bigl(\widehat X,
\mathcal{E}\!xt^q_{\mathcal O_{\widehat X}}
(\widehat{\mathcal F},\widehat{\mathcal G})\bigr)
\]
는 $\mathrm H^p\bigl(X,\mathcal{E}\!xt^q_{\mathcal O_X}
(\mathcal F,\mathcal G)\bigr)$의 $\mathfrak J$-준아딕 위상에 대한
분리 완비화와 표준적으로 동일시된다.

\oldpage[III]{139}

$\mathrm E(\widehat{\mathcal F},\widehat{\mathcal G})$를 (T, 4.2)에서
정의된 $\widehat{\mathcal F}$와 $\widehat{\mathcal G}$에 대한 쌍정칙
스펙트럼 열이라 하자. $\widehat A$가 $A$의 $\mathfrak J$-준아딕
위상에 대한 분리 완비화이면, 위에서 표준 동형사상 아래
\[
\mathrm E_2^{pq}(\widehat{\mathcal F},\widehat{\mathcal G})
=\mathrm E_2^{pq}(\mathcal F,\mathcal G)\otimes_A\widehat A
\qquad (0\textsubscript{I}, 7.3.3)
\]
를 얻는다.

이제 평탄 사상 $i$는 스펙트럼 열들의 표준 준동형사상
\[
\varphi:\mathrm E(\mathcal F,\mathcal G)\longrightarrow
\mathrm E(\widehat{\mathcal F},\widehat{\mathcal G})
=\mathrm E(i^*(\mathcal F),i^*(\mathcal G))
\]
을 정한다. 그 $\mathrm E_2$ 항들에서의 사상은
\[
\varphi_2^{pq}:\mathrm H^p\bigl(X,
\mathcal{E}\!xt^q_{\mathcal O_X}(\mathcal F,\mathcal G)\bigr)
\longrightarrow
\mathrm H^p\bigl(\widehat X,
\mathcal{E}\!xt^q_{\mathcal O_{\widehat X}}
(\widehat{\mathcal F},\widehat{\mathcal G})\bigr)
\]
이고, 수렴 대상에서의 사상은
\[
\varphi^n:\operatorname{Ext}^n_{\mathcal O_X}
(X;\mathcal F,\mathcal G)\longrightarrow
\operatorname{Ext}^n_{\mathcal O_{\widehat X}}
(\widehat X;\widehat{\mathcal F},\widehat{\mathcal G})
\]
이며, 이들은 함수성에 의해 각각 $u_q$와 $u_0$에서 얻어진다
(0, 12.3.4). $\widehat A$로 텐서하면 $\varphi_r^{pq}$와
$\varphi^n$에서 $\widehat A$-가군의 준동형사상
\[
\psi_r^{pq}:\mathrm E_r^{pq}(\mathcal F,\mathcal G)
\otimes_A\widehat A\longrightarrow
\mathrm E_r^{pq}(\widehat{\mathcal F},\widehat{\mathcal G})
\]
과
\[
\psi^n:\operatorname{Ext}^n_{\mathcal O_X}(X;\mathcal F,\mathcal G)
\otimes_A\widehat A\longrightarrow
\operatorname{Ext}^n_{\mathcal O_{\widehat X}}
(\widehat X;\widehat{\mathcal F},\widehat{\mathcal G})
\]
을 얻는다. $\widehat A$는 평탄 $A$-가군이므로
(0\textsubscript{I}, 7.3.3), $\widehat A$-가군
$\mathrm E_r^{pq}(\mathcal F,\mathcal G)\otimes_A\widehat A$들은
$\operatorname{Ext}^n_{\mathcal O_X}(X;\mathcal F,\mathcal G)
\otimes_A\widehat A$를 수렴 대상으로 하는 쌍정칙 스펙트럼 열을
이루고, $\psi_r^{pq}$와 $\psi^n$은 스펙트럼 열들의 사상을 이룬다.
$\psi_2^{pq}$들이 동형사상이므로 $\psi^n$들도 동형사상이다
(0, 11.1.5).
\end{proof}

\begin{corollary}[4.5.2]
\label{III.4.5.2-fr}
(\hyperref[III.4.5.1-fr]{4.5.1})의 가정 아래에서 $A$가 뇌터
$\mathfrak J$-아딕 환이라고 더 가정하자. 그러면 모든 정수
$n\geq0$에 대해
$\operatorname{Ext}^n_{\mathcal O_X}(X;\mathcal F,\mathcal G)$는
$\operatorname{Ext}^n_{\mathcal O_{\widehat X}}
(\widehat X;\widehat{\mathcal F},\widehat{\mathcal G})$와 표준적으로
동일시된다.
\end{corollary}

\begin{proof}
$\operatorname{Ext}^n_{\mathcal O_X}(X;\mathcal F,\mathcal G)$는 유한 생성
$A$-가군이므로 $\mathfrak J$-준아딕 위상에 대해 분리이고 완비임을
관찰하면 충분하다 (0\textsubscript{I}, 7.3.6).
\end{proof}

(\hyperref[III.4.5.1-fr]{4.5.1})에서 $n=0$인 특수한 경우는 다음과
같이 서술된다.

\begin{corollary}[4.5.3]
\label{III.4.5.3-fr}
(\hyperref[III.4.5.1-fr]{4.5.1})의 가정 아래에서, 임의의 준동형사상
$u:\mathcal F\to\mathcal G$에 대해 $\widehat u:\widehat{\mathcal F}
\to\widehat{\mathcal G}$를 그 완비화된 준동형사상이라 하자
(I, 10.8.4). 그러면 표준 동형사상
\[
\label{III.4.5.3.1-fr}
\left(\operatorname{Hom}_{\mathcal O_X}(\mathcal F,\mathcal G)\right)^\wedge
\xrightarrow{\sim}
\operatorname{Hom}_{\mathcal O_{\widehat X}}
(\widehat{\mathcal F},\widehat{\mathcal G})
\tag{4.5.3.1}
\]
이 존재한다. 여기서 왼쪽 항은 $A$-가군
$\operatorname{Hom}_{\mathcal O_X}(\mathcal F,\mathcal G)$의
$\mathfrak J$-준아딕 위상에 대한 분리 완비화이고, 이 동형사상은
준동형사상 $u\mapsto\widehat u$에서 분리 완비화로 이행하여 얻어진다.
\end{corollary}

\subsection{형식 사상과 보통 사상의 관계}

\begin{proposition}[4.6.1]
\label{III.4.6.1-fr}
$Y$를 국소 뇌터 준스킴, $f:X\to Y$를 고유 사상,
$\mathcal F$를 연접이고 $f$-평탄인 $\mathcal O_X$-가군, $y$를
$Y$의 점이라 하자. 어떤 정수 $n$에 대해
\[
H^n\bigl(f^{-1}(y),\mathcal F\otimes_{\mathcal O_Y}k(y)\bigr)=0
\]
이라고 가정하자. 그러면 $Y$ 안의 $y$의 어떤 근방 $U$가 존재하여
$R^nf_*(\mathcal F)|U=0$이고, 모든 정수 $p\geq0$에 대해
(\hyperref[III.4.2.1.1-fr]{4.2.1.1})의 표준 준동형사상
\[
\bigl(R^{n-1}f_*(\mathcal F)\bigr)_y
\longrightarrow
H^{n-1}\bigl(f^{-1}(y),
\mathcal F\otimes_{\mathcal O_Y}
(\mathcal O_y/\mathfrak m_y^{p+1})\bigr)
\]
은 전사이다.
\end{proposition}

\begin{proof}
$R^nf_*(\mathcal F)$는 연접 $\mathcal O_Y$-가군이므로
(\hyperref[III.3.2.1-fr]{3.2.1}), 명제의 첫째 주장을 보이려면
$(R^nf_*(\mathcal F))_y=0$임을 보이면 충분하다
(0\textsubscript{I}, 5.2.2). (\hyperref[III.4.2.1-fr]{4.2.1})에
의해, 이를 위해서는 모든 $p$에 대해
\[
H^n\bigl(f^{-1}(y),\mathcal F\otimes_{\mathcal O_Y}
(\mathcal O_y/\mathfrak m_y^{p+1})\bigr)=0
\]
임을 보이면 충분하다. $p=0$일 때는 가정에 의해 참이므로 $p$에
관한 귀납법으로 보이겠다. 다음과 같이 놓자.
\[
X_p=X\times_Y\operatorname{Spec}
(\mathcal O_y/\mathfrak m_y^{p+1}).
\]
그러면 $X_{p-1}$은 $X_p$와 같은 바탕공간을 갖는 닫힌 부분준스킴이다
(I, 3.6.1). 따라서 귀납 가정
\[
H^n\bigl(X_{p-1},\mathcal F\otimes_{\mathcal O_Y}
(\mathcal O_y/\mathfrak m_y^p)\bigr)=0
\]
에서
\[
H^n\bigl(X_p,\mathcal F\otimes_{\mathcal O_Y}
(\mathcal O_y/\mathfrak m_y^p)\bigr)=0
\]
이 따른다. 한편 $\mathcal O_{X_p}$-가군의 완전열
\[
0\longrightarrow
\mathfrak m_y^p\mathcal F/\mathfrak m_y^{p+1}\mathcal F
\longrightarrow
\mathcal F/\mathfrak m_y^{p+1}\mathcal F
\longrightarrow
\mathcal F/\mathfrak m_y^p\mathcal F
\longrightarrow0
\]
에서 얻는 코호몰로지 완전열은
\[
H^n\bigl(X_p,
\mathfrak m_y^p\mathcal F/\mathfrak m_y^{p+1}\mathcal F\bigr)
\longrightarrow
H^n\bigl(X_p,\mathcal F/\mathfrak m_y^{p+1}\mathcal F\bigr)
\longrightarrow
H^n\bigl(X_p,\mathcal F/\mathfrak m_y^p\mathcal F\bigr)
\]
이 완전임을 준다. 따라서
\[
\label{III.4.6.1.1-fr}
H^n\bigl(X_p,
\mathfrak m_y^p\mathcal F/\mathfrak m_y^{p+1}\mathcal F\bigr)=0
\tag{4.6.1.1}
\]
임을 보이면 충분하다. 실제로 그러면
$H^n(X_p,\mathcal F/\mathfrak m_y^{p+1}\mathcal F)$는
$H^n(X_p,\mathcal F/\mathfrak m_y^p\mathcal F)$의 부분가군이고,
귀납 가정에 의해 $0$이다.

\oldpage[III]{140}

이제 올
$Z=f^{-1}(y)=X\times_Y\operatorname{Spec}(k(y))$는 $X_p$의 닫힌
부분준스킴이고,
$\mathfrak m_y^p\mathcal F/\mathfrak m_y^{p+1}\mathcal F$는
$\mathfrak m_y$에 의해 소멸된다. 따라서 이를 $\mathcal O_Z$-가군으로
볼 수 있으며
\[
H^n\bigl(Z,\mathfrak m_y^p\mathcal F/
\mathfrak m_y^{p+1}\mathcal F\bigr)
=H^n\bigl(X_p,\mathfrak m_y^p\mathcal F/
\mathfrak m_y^{p+1}\mathcal F\bigr).
\]
표준 $\mathcal O_Z$-준동형사상
\[
\label{III.4.6.1.2-fr}
(\mathcal F/\mathfrak m_y\mathcal F)
\otimes_{k(y)}(\mathfrak m_y^p/\mathfrak m_y^{p+1})
\longrightarrow
\mathfrak m_y^p\mathcal F/\mathfrak m_y^{p+1}\mathcal F
\tag{4.6.1.2}
\]
이 전단사임을 보이자. 이를 보이면
$\mathfrak m_y^p/\mathfrak m_y^{p+1}$가 자유 $k(y)$-가군이므로
\[
\begin{split}
H^n\bigl(Z,\mathfrak m_y^p\mathcal F/
\mathfrak m_y^{p+1}\mathcal F\bigr)
&=H^n\bigl(Z,\mathcal F/\mathfrak m_y\mathcal F\bigr)
\otimes_{k(y)}(\mathfrak m_y^p/\mathfrak m_y^{p+1})\\
&=0
\end{split}
\]
을 얻는다 (0, 12.2.3). 마지막 등식은 가정에 의해
$H^n(Z,\mathcal F/\mathfrak m_y\mathcal F)=0$이기 때문이고, 이로써
(\hyperref[III.4.6.1.1-fr]{4.6.1.1})을 얻는다. 이제 첫째 주장을
마치려면 (\hyperref[III.4.6.1.2-fr]{4.6.1.2})가 전단사임만 보이면
된다. 이 문제는 $X$ 위에서 점별 문제이고, 가정에 의해 모든
$x\in f^{-1}(y)$에 대해 $\mathcal F_x$는 평탄
$\mathcal O_y$-가군이다. 또한
$\mathcal F_x/\mathfrak m_y\mathcal F_x$는 체
$k(y)=\mathcal O_y/\mathfrak m_y$ 위의 평탄 가군이므로
(0, 10.2.1, c))를 적용하면 된다.

(\hyperref[III.4.6.1-fr]{4.6.1})의 둘째 주장을 보이기 위해서는
(\hyperref[III.4.2.1-fr]{4.2.1})에서와 같이 곧바로 $Y$가 아핀이고
$y$가 닫힌 경우로 귀착된다. 관계
(\hyperref[III.4.6.1.1-fr]{4.6.1.1})은 같은 논증으로 모든 $k>0$과
모든 $p\geq0$에 대해
\[
\label{III.4.6.1.3-fr}
H^n\bigl(X_p,
\mathfrak m_y^k\mathcal F/
\mathfrak m_y^{k+p+1}\mathcal F\bigr)=0
\tag{4.6.1.3}
\]
을 얻는다.

\oldpage[III]{141}

그러므로 (\hyperref[III.4.2.1-fr]{4.2.1})에 의해
\[
\label{III.4.6.1.4-fr}
\bigl(R^nf_*(\mathfrak m_y^k\mathcal F)\bigr)_y=0
\tag{4.6.1.4}
\]
도 성립한다. 이제 코호몰로지 완전열에서
\[
\bigl(R^{n-1}f_*(\mathcal F)\bigr)_y
\longrightarrow
\bigl(R^{n-1}f_*(\mathcal F/\mathfrak m_y^p\mathcal F)\bigr)_y
\longrightarrow
\bigl(R^nf_*(\mathfrak m_y^p\mathcal F)\bigr)_y=0
\]
의 완전성을 얻는다. $y$가 닫혀 있고 $Y$가 아핀이므로
(\hyperref[III.1.4.11-fr]{1.4.11})
\[
R^{n-1}f_*(\mathcal F/\mathfrak m_y^p\mathcal F)
=\bigl(H^{n-1}(X,\mathcal F/\mathfrak m_y^p\mathcal F)\bigr)^\sim
=\bigl(H^{n-1}(f^{-1}(y),
\mathcal F/\mathfrak m_y^p\mathcal F)\bigr)^\sim
\]
이다 (G, II, 4.9.1). 그런데
$H^{n-1}(f^{-1}(y),\mathcal F/\mathfrak m_y^p\mathcal F)$는
$(\mathcal O_y/\mathfrak m_y^p)$-가군이므로
\[
\bigl(R^{n-1}f_*(\mathcal F/\mathfrak m_y^p\mathcal F)\bigr)_y
=H^{n-1}(f^{-1}(y),\mathcal F/\mathfrak m_y^p\mathcal F).
\]
이로써 (\hyperref[III.4.6.1-fr]{4.6.1})의 증명이 끝난다.
\end{proof}

\begin{corollary}[4.6.2]
\label{III.4.6.2-fr}
$Y$를 국소 뇌터 준스킴, $f:X\to Y$를 고유이고 평탄인 사상,
$\mathcal F$, $\mathcal G$를 두 국소 자유 $\mathcal O_X$-가군,
$y$를 $Y$의 점이라 하자. 다음과 같이 놓자.
\[
X_y=f^{-1}(y)=X\otimes_Yk(y),\qquad
\mathcal F_y=\mathcal F\otimes_{\mathcal O_Y}k(y),\qquad
\mathcal G_y=\mathcal G\otimes_{\mathcal O_Y}k(y).
\]
또한
\[
\label{III.4.6.2.1-fr}
H^1\bigl(X_y,
\mathcal{H}\!om_{\mathcal O_{X_y}}(\mathcal F_y,\mathcal G_y)\bigr)=0
\tag{4.6.2.1}
\]
이라고 가정하자. 그러면 모든 준동형사상
$u_0:\mathcal F_y\to\mathcal G_y$에 대해 $y$의 열린 근방 $U$와
준동형사상
\[
u:\mathcal F|f^{-1}(U)\longrightarrow\mathcal G|f^{-1}(U)
\]
가 존재하여 $u_0=u\otimes1$이다.
\end{corollary}

\begin{proof}
실제로 이 가정에 의해 $n=1$, $p=0$으로
(\hyperref[III.4.6.1-fr]{4.6.1})을 연접 $\mathcal O_X$-가군
\[
\mathcal H=\mathcal{H}\!om_{\mathcal O_X}(\mathcal F,\mathcal G)
\]
에 적용할 수 있다. $\mathcal H$는 국소 자유이므로 특히
$f$-평탄이고, 이때 $\mathcal O_{X_y}$-가군
$\mathcal H\otimes_{\mathcal O_Y}k(y)$는
$\mathcal{H}\!om_{\mathcal O_{X_y}}(\mathcal F_y,\mathcal G_y)$와
동일시된다 (0\textsubscript{I}, 6.2.2). $Y=\operatorname{Spec}(A)$가
아핀이라고 가정할 수 있고, 그러면
(\hyperref[III.1.4.11-fr]{1.4.11})
\[
R^0f_*(\mathcal H)=
\bigl(\operatorname{Hom}_{\mathcal O_X}(\mathcal F,\mathcal G)\bigr)^\sim,
\]
따라서
\[
\bigl(R^0f_*(\mathcal H)\bigr)_y
=\operatorname{Hom}_{\mathcal O_X}(\mathcal F,\mathcal G)
\otimes_A\mathcal O_y.
\]
(\hyperref[III.4.6.1-fr]{4.6.1})에 의해 표준 준동형사상
\[
\operatorname{Hom}_{\mathcal O_X}(\mathcal F,\mathcal G)
\otimes_A\mathcal O_y
\longrightarrow
\operatorname{Hom}_{\mathcal O_{X_y}}(\mathcal F_y,\mathcal G_y)
\]
은 전사이다. 또한
$\operatorname{Hom}_{\mathcal O_X}(\mathcal F,\mathcal G)
\otimes_A\mathcal O_y$의 모든 원소는 $s\notin\mathfrak m_y$인
$s\in A$와 어떤 $u$를 써서 $u\otimes(1/s)$의 꼴로 나타낼 수 있다.
이로써 따름정리가 성립한다.
\end{proof}

이 따름정리는 다음 결과로 보충된다.

\begin{corollary}[4.6.3]
\label{III.4.6.3-fr}
(\hyperref[III.4.6.2-fr]{4.6.2})의 가정 아래에서 $u_0$이 단사이면
(각각 전사이면, 전단사이면) $u$도 그러하다고 가정할 수 있다.
\end{corollary}

\begin{proof}
$U=Y$인 경우로 한정할 수 있다. $u_0$이 단사이면 (각각 전사이면)
모든 $x\in f^{-1}(y)$에 대해 $\operatorname{Ker}u_x=0$임을
(각각 $\operatorname{Coker}u_x=0$임을) 보이면 충분하다. 실제로
$\operatorname{Ker}u$와 $\operatorname{Coker}u$는 연접
$\mathcal O_X$-가군이므로 (0\textsubscript{I}, 5.3.4),
$f^{-1}(y)$의 어떤 근방 $V$가 존재하여
$\operatorname{Ker}u$의 (각각 $\operatorname{Coker}u$의) $V$에
대한 제한이 $0$이다 (0\textsubscript{I}, 5.2.2). $f$는 닫힌
사상이므로 $y$의 어떤 근방 $U'\subset U$가 존재하여
$f^{-1}(U')\subset V$이고, 이로써
(\hyperref[III.4.6.3-fr]{4.6.3})이 증명된다.

가정에 의해
\[
u_x\otimes1:
\mathcal F_x\otimes_{\mathcal O_y}k(y)
\longrightarrow
\mathcal G_x\otimes_{\mathcal O_y}k(y)
\]
는 단사이다 (각각 전사이다). $\mathcal F_x$와 $\mathcal G_x$는
유한 생성 자유 $\mathcal O_x$-가군이고, $\mathcal O_x$는 평탄
$\mathcal O_y$-가군이다. $u_x\otimes1$이 단사라고 가정하면
(0, 10.2.4)에 의해 $u_x$도 단사이다. $u_x\otimes1$이 전사라고
가정하면 몫을 취하여 얻는 준동형사상
\[
\mathcal F_x/\mathfrak m_x\mathcal F_x
\longrightarrow
\mathcal G_x/\mathfrak m_x\mathcal G_x
\]
은 더욱이 전사이다. $\mathcal G_x$는 유한 생성
$\mathcal O_x$-가군이고 $\mathcal O_x$는 극대 아이디얼이
$\mathfrak m_x$인 국소환이므로, 나카야마 보조정리에서 결론이
따른다 (Bourbaki, \emph{Alg.}, chap. VIII, \S\,6, no~3,
cor. 4 de la prop. 6).
\end{proof}

특히 (\hyperref[III.4.6.3-fr]{4.6.3})에서 다음이 따른다.

\begin{corollary}[4.6.4]
\label{III.4.6.4-fr}
$Y$를 국소 뇌터 준스킴, $f:X\to Y$를 고유이고 평탄인 사상,
$y$를 $Y$의 점이라 하고 $X_y=X\otimes_Yk(y)$로 놓자.
$\mathcal E_0$을
\[
\label{III.4.6.4.1-fr}
H^1\bigl(X_y,
\mathcal{H}\!om_{\mathcal O_{X_y}}(\mathcal E_0,\mathcal E_0)\bigr)=0
\tag{4.6.4.1}
\]
을 만족하는 국소 자유 $\mathcal O_{X_y}$-가군이라 하자.
$\mathcal F$, $\mathcal G$를 두 국소 자유 $\mathcal O_X$-가군이라
하되, (\hyperref[III.4.6.2-fr]{4.6.2})의 표기에서
$\mathcal F_y$와 $\mathcal G_y$가 $\mathcal E_0$과 동형이라고
하자. 그러면 $y$의 어떤 열린 근방 $U$가 존재하여
$\mathcal F|f^{-1}(U)$와 $\mathcal G|f^{-1}(U)$는 동형이다.
\end{corollary}

더 특별히 다음을 얻는다.

\begin{corollary}[4.6.5]
\label{III.4.6.5-fr}
$f$, $X$, $Y$에 대해 (\hyperref[III.4.6.4-fr]{4.6.4})의 가정이
성립하고 $H^1(X_y,\mathcal O_{X_y})=0$이라고 하자.
$\mathcal F$와 $\mathcal G$가 두 가역 $\mathcal O_X$-가군이고
$\mathcal F_y$와 $\mathcal G_y$가 동형이면, $y$의 어떤 열린 근방
$U$가 존재하여 $\mathcal F|f^{-1}(U)$와
$\mathcal G|f^{-1}(U)$는 동형이다.
\end{corollary}

\begin{proof}
(\hyperref[III.4.6.4-fr]{4.6.4})을
$\mathcal F^{-1}\otimes\mathcal G$와 $\mathcal O_X$에 적용하면
충분하다.
\end{proof}

\begin{remarks}[4.6.6]
\label{III.4.6.6-fr}
\begin{enumerate}
\item[(i)]
(\hyperref[III.4.6.5-fr]{4.6.5})를 이용하여 제V장에서
(II, 4.2.7)에 예고한 사영 다발 위의 가역층의 분류를 확립하겠다.
\item[(ii)]
(\hyperref[III.4.6.1-fr]{4.6.1})의 결과는 \S\,7에서 더 일반적인
명제들의 귀결로 다시 나타날 것이다.
\end{enumerate}
\end{remarks}

\begin{proposition}[4.6.7]
\label{III.4.6.7-fr}
$Z$를 국소 뇌터 준스킴, $X$, $Y$를 두 $Z$-준스킴이라 하고, 구조
사상 $g:X\to Z$, $h:Y\to Z$가 고유라고 하자. $f:X\to Y$를
$Z$-사상, $z$를 $Z$의 점이라 하고
\[
f_z=f\times_Z1:
X\otimes_Zk(z)\longrightarrow Y\otimes_Zk(z)
\]
로 놓자.
\begin{enumerate}
\item[(i)]
$f_z$가 유한 사상이면 (각각 닫힌 몰입 사상이면), $z$의 어떤 열린
근방 $U$가 존재하여 $f$의 제한
$g^{-1}(U)\to h^{-1}(U)$는 유한 사상이다 (각각 닫힌 몰입
사상이다).
\item[(ii)]
더 나아가 $g$가 평탄 사상이라고 가정하자. $f_z$가 동형사상이면
$z$의 어떤 열린 근방 $U$가 존재하여 $f$의 제한
$g^{-1}(U)\to h^{-1}(U)$는 동형사상이다.
\end{enumerate}
\end{proposition}

\begin{proof}
두 경우 모두 모든 $y\in h^{-1}(z)$에 대해 $y$의 어떤 근방
$V_y$가 존재하여 $f$의 제한 $f^{-1}(V_y)\to V_y$가 유한
사상임을 (각각 닫힌 몰입 사상임을, 동형사상임을) 보이면 충분하다.
실제로 $V$를 $V_y$들의 합집합이라 하면 $f$의 제한
$f^{-1}(V)\to V$도 유한 사상이다 (각각 닫힌 몰입 사상이다,
동형사상이다) (II, 6.1.1 및 I, 4.2.4). $h$가 닫힌 사상이므로
$z$의 어떤 열린 근방 $U$가 존재하여 $h^{-1}(U)\subset V$이고,
명제가 증명된다.

\noindent\textit{(i)}
먼저 $f$가 고유 사상임을 주의하자 (II, 5.4.3). $f_z$가 유한이라고
가정하면, 모든 $y\in h^{-1}(z)$에 대해
$f^{-1}(V_y)\to V_y$가 유한이 되는 근방 $V_y$의 존재는
(\hyperref[III.4.4.11-fr]{4.4.11})에서 따른다.

\oldpage[III]{143}

$f_z$가 닫힌 몰입 사상인 경우를 다루기 위해서는 이미 $f$가 유한
사상이라고 가정할 수 있다. 따라서
$X=\operatorname{Spec}(\mathcal B)$이고, 여기서 $\mathcal B$는 연접
$\mathcal O_Y$-대수이며, $f$는 표준 준동형사상
$u:\mathcal O_Y\to\mathcal B$에 대응한다 (II, 1.2.7).
모든 $y\in h^{-1}(z)$에 대해
$u_y:\mathcal O_y\to\mathcal B_y$가 전사임을 보이면, 층
$\operatorname{Coker}u$가 연접이므로 (0\textsubscript{I}, 5.3.4
및 5.2.2) $y$의 어떤 근방 $V_y$에서 $u|V_y$가 전사임이 따를
것이다. 한편 유한 사상 $f_z$는 준동형사상
\[
v=u\otimes1:
\mathcal O_Y\otimes_{\mathcal O_Z}k(z)
\longrightarrow
\mathcal B\otimes_{\mathcal O_Z}k(z)
\]
에 대응하고, $f_z$가 닫힌 몰입 사상이라는 가정에서
\[
u_y\otimes1:
\mathcal O_y\otimes_{\mathcal O_z}k(z)
\longrightarrow
\mathcal B_y\otimes_{\mathcal O_z}k(z)
\]
가 전사임이 따른다. $\mathcal B_y$는 유한 생성
$\mathcal O_y$-가군이고 $\mathcal O_y$는 뇌터 국소환이므로,
(\hyperref[III.4.6.3-fr]{4.6.3})에서와 같이 나카야마 보조정리에서
결론이 따른다.

\noindent\textit{(ii)}
위와 같은 논증에 의해, 이제는 $u_y\otimes1$이 전단사임을 알고
$u_y$가 전단사임을 보이면 충분하다. 이는 다음 보조정리에서 따른다.
\end{proof}

\begin{lemma}[4.6.7.1]
\label{III.4.6.7.1-fr}
$A$, $B$를 두 뇌터 국소환, $\rho:A\to B$를 국소 준동형사상,
$u:N\to M$을 $B$-가군의 준동형사상이라 하자. $M$은 평탄
$A$-가군이고 $N$은 유한 생성 $B$-가군이며
\[
u\otimes1:N\otimes_Ak\longrightarrow M\otimes_Ak
\]
가 단사라고 가정하자. 여기서 $k$는 $A$의 잉여체이다. 그러면
$N$은 평탄 $A$-가군이고 $u$는 단사이다.
\end{lemma}

\begin{proof}
첫째 주장을 보이려면 유한 생성 $A$-가군의 임의의 쌍 $P$, $Q$와
임의의 단사 $A$-준동형사상 $v:P\to Q$에 대해
\[
1_N\otimes v:N\otimes_AP\longrightarrow N\otimes_AQ
\]
가 단사임을 보이면 된다. 그런데 가환 그림
\[
\xymatrix{
N\otimes_AP\ar[r]^{1_N\otimes v}\ar[d]_{u\otimes1_P}
  &N\otimes_AQ\ar[d]^{u\otimes1_Q}\\
M\otimes_AP\ar[r]_{1_M\otimes v}
  &M\otimes_AQ
}
\]
이 있고, 가정에 의해 $1_M\otimes v$가 단사이므로
$u\otimes1_P$가 단사임을 보이면 충분하다.

$\mathfrak m$을 $A$의 극대 아이디얼이라 하자.
$A$-가군 $N\otimes_AP$ 위의 $\mathfrak m$-아딕 여과는 이를
$B$-가군으로 보았을 때의 $\mathfrak mB$-아딕 여과이기도 하다.
$B$는 뇌터환이고 $\mathfrak mB$는 $B$의 야코브슨 근기에 포함되며,
$N$이 유한 생성 $B$-가군이고 $P$가 유한 생성 $A$-가군이므로
$N\otimes_AP$는 유한 생성 $B$-가군이다. 따라서 이 여과가 정하는
위상은 분리이다 (0\textsubscript{I}, 7.3.5). 그러므로
\[
\operatorname{gr}(u\otimes1_P):
\operatorname{gr}_\bullet(N\otimes_AP)
\longrightarrow
\operatorname{gr}_\bullet(M\otimes_AP)
\]
가 단사임을 보이면 충분하다. 여기서 등급 가군들은
$\mathfrak m$-아딕 여과에 관한 것이다
(Bourbaki, \emph{Alg. comm.}, chap. III, \S\,2, no~8,
cor. 1 du th. 1).

$M$은 평탄 $A$-가군이므로 준동형사상
\[
M\otimes_A(\mathfrak m^nP)
\longrightarrow
\mathfrak m^n(M\otimes_AP)
\]
들은 전단사이다. 따라서 표준 준동형사상
\[
\varphi_M:
\operatorname{gr}_0(M)\otimes_A\operatorname{gr}_\bullet(P)
\longrightarrow
\operatorname{gr}_\bullet(M\otimes_AP)
\]
도 전단사이다.

\oldpage[III]{144}

이제 가환 그림
\[
\xymatrix{
\operatorname{gr}_0(N)\otimes_A\operatorname{gr}_\bullet(P)
  \ar[r]^{\operatorname{gr}_0(u)\otimes1}\ar[d]_{\varphi_N}
  &\operatorname{gr}_0(M)\otimes_A\operatorname{gr}_\bullet(P)
  \ar[d]^{\varphi_M}\\
\operatorname{gr}_\bullet(N\otimes_AP)
  \ar[r]_{\operatorname{gr}(u\otimes1)}
  &\operatorname{gr}_\bullet(M\otimes_AP)
}
\]
이 있다. 여기서 $\varphi_M$은 전단사이고 $\varphi_N$은 전사이다.
또한 가정에 의해 $\operatorname{gr}_0(u)$가 단사이고
\[
\begin{split}
\operatorname{gr}_0(N)\otimes_A\operatorname{gr}_\bullet(P)
&=\operatorname{gr}_0(N)\otimes_k\operatorname{gr}_\bullet(P),\\
\operatorname{gr}_0(M)\otimes_A\operatorname{gr}_\bullet(P)
&=\operatorname{gr}_0(M)\otimes_k\operatorname{gr}_\bullet(P)
\end{split}
\]
이므로 $\operatorname{gr}_0(u)\otimes1$도 단사이다. 따라서
$\operatorname{gr}(u\otimes1)$은 단사이고, 이로써 첫째 주장이
증명된다. 앞의 논증에서 $P=A$로 놓으면 둘째 주장이 따른다.
\end{proof}

\begin{proposition}[4.6.8]
\label{III.4.6.8-fr}
$Z$를 국소 뇌터 준스킴, $X$, $Y$를 두 $Z$-준스킴이라 하고, 구조
사상 $g:X\to Z$, $h:Y\to Z$가 고유라고 하자. $Z'$을 $Z$의 닫힌
부분집합, $X'=g^{-1}(Z')$, $Y'=h^{-1}(Z')$을 그 역상들이라 하고,
\[
\widehat X=X_{/X'},\qquad
\widehat Y=Y_{/Y'},\qquad
\widehat Z=Z_{/Z'}
\]
을 이 닫힌 부분집합들을 따른 $X$, $Y$, $Z$의 형식적 완비화라
하자. $f:X\to Y$를 $Z$-사상,
$\widehat f:\widehat X\to\widehat Y$를 그 완비화로의 연장이라
하자. $\widehat f$가 동형사상이기 위한 (각각 닫힌 몰입 사상이기
위한) 필요충분조건은, $Z'$의 어떤 열린 근방 $U$가 존재하여 $f$의
제한 $g^{-1}(U)\to h^{-1}(U)$가 동형사상인 것이다
(각각 닫힌 몰입 사상인 것이다).
\end{proposition}

\begin{proof}
조건의 충분성은 즉시 따른다 (I, 10.14.7). 필요성을 보이기
위해서는 (\hyperref[III.4.6.7-fr]{4.6.7})에서와 같은 논증에 의해
모든 $y\in Y'$에 대해 $y$의 어떤 열린 근방 $V_y$가 존재하여
$f$의 제한 $f^{-1}(V_y)\to V_y$가 동형사상임을 (각각 닫힌 몰입
사상임을) 보이면 충분하다. 따라서 $Y=Z=\operatorname{Spec}(A)$이고
$A$가 뇌터환인 경우로 귀착된다.

가정에 의해 (I, 10.9.1 및 10.14.2) 모든 $y\in Y'$에 대해 올
$f^{-1}(y)$는 한 점 이하이다. 따라서 $f$가 고유이므로
(II, 5.4.3), $y$의 어떤 열린 근방 $U$가 존재하여 $f$의 제한
$f^{-1}(U)\to U$가 유한 사상이다
(\hyperref[III.4.4.11-fr]{4.4.11}). 그러므로 $f$가 유한
사상이라고 가정할 수 있고, 따라서 $X=\operatorname{Spec}(B)$이며
$B$는 유한 $A$-대수이다. $Y'=V(\mathfrak J)$이면
\[
\widehat Y=\operatorname{Spf}(\widehat A),\qquad
\widehat X=\operatorname{Spf}(\widehat B),
\]
여기서 $\widehat A$는 $A$의 $\mathfrak J$-준아딕 위상에 대한
분리 완비화이고, $\widehat B$는 $B$의 $\mathfrak JB$-준아딕
위상에 대한 분리 완비화, 또는 이와 같은 것으로서 $A$-가군 $B$의
$\mathfrak J$-준아딕 위상에 대한 분리 완비화이다. 또한
$\widehat f$는 표준 환 준동형사상 $\varphi:A\to B$의 연속적 연장
\[
\widehat\varphi:\widehat A\longrightarrow\widehat B
\]
에 대응하는 아핀 형식 스킴의 사상이다. 가정은
$\widehat\varphi$가 전단사라는 것이다 (각각 전사라는 것이다)
(I, 10.14.2). 그런데 $\widehat\varphi$는 $\varphi$를
$A$-가군의 준동형사상으로 보았을 때의 연속적 연장이기도 하다.
따라서 (I, 10.8.14)에 의해 $Y'$의 어떤 열린 근방 $U$가 존재하여
$\mathcal O_Y$-가군의 준동형사상
$\widetilde\varphi:\widetilde A\to\widetilde B$의 $U$에 대한
제한은 전단사이다 (각각 전사이다). 이로써 증명이 끝난다.
\end{proof}

\oldpage[III]{145}

\subsection{풍부성 판정법}

\begin{theorem}[4.7.1]
\label{III.4.7.1-fr}
$Y$를 국소 뇌터 준스킴, $f:X\to Y$를 고유 사상,
$\mathcal L$을 가역 $\mathcal O_X$-가군, $y$를 $Y$의 점이라 하자.
$X_y=X\otimes_Yk(y)=f^{-1}(y)$라 하고 $g$를 $X_y$에서 $X$로 가는
사영이라 하자. $\mathcal L_y=g^*(\mathcal L)
=\mathcal L\otimes_{\mathcal O_Y}k(y)$가 $X_y$ 위에서 풍부하면,
$Y$ 안의 $y$의 어떤 열린 근방 $U$가 존재하여
$\mathcal L|f^{-1}(U)$는 $f$의 $f^{-1}(U)$에 대한 제한에 관해
풍부하다.
\end{theorem}

\begin{proof}
\noindent\textit{I)}
다음과 같이 놓자.
\[
Y'=\operatorname{Spec}(\mathcal O_y),\qquad
X'=X\times_YY',\qquad
\mathcal L'=\mathcal L\otimes_{\mathcal O_Y}\mathcal O_y.
\]
먼저 $\mathcal L'$이 $f'=f_{(Y')}$에 관해 풍부함을 보이겠다.
가환 그림
\[
\xymatrix{
X\ar[d]_f & X'\ar[l]\ar[d]_{f'} & X_y\ar[l]\ar[d]_{f_y}\\
Y & Y'\ar[l] & \operatorname{Spec}(k(y))\ar[l]
}
\]
이 있다. $f'$은 고유이고 (II, 5.4.2, (iii)) $\mathcal O_y$는
뇌터환이므로, $Y=Y'=\operatorname{Spec}(\mathcal O_y)$이고
따라서 $X=X'$인 경우로 한정할 수 있다. 이때 $\mathcal L_y$가
$f_y$에 관해 풍부하다고 가정하고 $\mathcal L$이 $f$에 관해
풍부함을 보이면 된다 (II, 4.6.6).

판정법 (\hyperref[III.2.6.1-fr]{2.6.1, c)})을 적용하겠다. 실제로
모든 연접 $\mathcal O_X$-가군 $\mathcal F$에 대해 어떤 정수 $N$이
존재하여
\[
H^1(X,\mathcal F(n))=0\quad\text{$n\geq N$인 모든 $n$에 대해},
\qquad
\mathcal F(n)=\mathcal F\otimes\mathcal L^{\otimes n}
\]
임을 보이겠다. $y$는 $\mathcal O_y$의 극대 아이디얼
$\mathfrak m$에 대응하는 $Y$의 닫힌점이다. 따라서 $X_y$는 연접
아이디얼
\[
\mathcal J=f^*(\widetilde{\mathfrak m})\mathcal O_X
=\mathfrak m\mathcal O_X
\]
에 의해 정의되는 $X$의 닫힌 부분준스킴이고 (I, 4.4.5), $g$는
표준 단사이다. 이제 등급 $k(y)$-대수
\[
S=\operatorname{gr}(\mathcal O_y)
=\bigoplus_{j\geq0}\mathfrak m^j/\mathfrak m^{j+1}
\]
를 생각하자. $\mathcal O_y$가 뇌터환이므로 $S$는 유한형이다.
그러므로 $\mathcal O_X$-대수
$\mathcal S=f^*(\widetilde S)$는 준연접이고 유한형이다.
또한 이는 분명히 $\mathcal J$에 의해 소멸된다. 따라서
$\mathcal S_y=g^*(\mathcal S)$로 놓으면 $\mathcal S_y$는 준연접
유한형 $\mathcal O_{X_y}$-대수이고
$\mathcal S=g_*(\mathcal S_y)$이다.

한편
\[
\mathcal M_j=\mathfrak m^j\mathcal F/
\mathfrak m^{j+1}\mathcal F,\qquad
\mathcal M=\bigoplus_{j\geq0}\mathcal M_j
=\operatorname{gr}(\mathcal F)
\]
로 놓자. $\mathcal F$가 연접이므로 $\mathcal M$은 준연접 유한형
$\mathcal S$-가군이고 (0, 10.1.1), 역시 $\mathcal J$에 의해
소멸된다. 따라서
\[
\mathcal M'_j=g^*(\mathcal M_j),\qquad
\mathcal M'=g^*(\mathcal M)
=\bigoplus_{j\geq0}\mathcal M'_j
\]
로 놓으면 $\mathcal M'$은 준연접 유한형 등급
$\mathcal S_y$-가군이고 $\mathcal M=g_*(\mathcal M')$이다.
더 나아가
$\mathcal M'_j(n)=\mathcal M'_j\otimes\mathcal L_y^{\otimes n}$으로
놓으면 $\mathcal M'_j(n)=g^*(\mathcal M_j(n))$이다.

$f_y$는 고유이고 (II, 5.4.2, (iii)) $\mathcal L_y$는
풍부하므로 $f_y$는 사영 사상이다 (II, 5.5.4 및 4.6.11).
따라서 $\operatorname{Spec}(k(y))$, $f_y$, $\mathcal S_y$,
$\mathcal L_y$, $\mathcal M'$에 정리
(\hyperref[III.2.4.1-fr]{2.4.1, (ii)})를 적용할 수 있다.
어떤 정수 $N$이 존재하여 $n\geq N$이면 모든 $q>0$과 모든 $j$에
대해
\[
H^q(X_y,\mathcal M'_j(n))=0
\]
이다. 따라서 모든 $q>0$과 모든 $j$에 대해
$H^q(X,\mathcal M_j(n))=0$이기도 하다 (G, II, 4.9.1).
이제
\[
\mathcal F(n)_j
=\mathcal F(n)/\mathfrak m^{j+1}\mathcal F(n)
\]
로 놓자. 그러면 $j\geq1$일 때
$\mathcal F(n)_{j-1}=\mathcal F(n)_j/\mathcal M_j(n)$이고
$\mathcal F(n)_0=\mathcal M_0(n)$이다.
$H^1(X,\mathcal F(n)_0)=0$이고, 코호몰로지 완전열에 의해 모든
$j\geq1$에 대해
\[
H^1(X,\mathcal F(n)_j)=H^1(X,\mathcal F(n)_{j-1}).
\]
따라서 모든 $j\geq0$에 대해
$H^1(X,\mathcal F(n)_j)=0$이다. 이제
(\hyperref[III.4.2.1-fr]{4.2.1})에 의해
$H^1(X,\mathcal F(n))=0$이고, 이로써 첫째 주장이 증명된다.

\oldpage[III]{146}

\noindent\textit{II)}
증명의 처음 표기로 돌아가자. 언제나 $Y=\operatorname{Spec}(A)$가
아핀이라고 가정할 수 있다. $f'$은 유한형이고 $\mathcal L'$은
$f'$에 관해 풍부하므로, 어떤 정수 $m>0$이 존재하여
$\mathcal L'^{\otimes m}$은 $f'$에 관해 매우 풍부하다
(II, 4.6.11). 필요하면 $\mathcal L$을
$\mathcal L^{\otimes m}$으로 바꾸어 $\mathcal L'$이 $f'$에 관해
매우 풍부한 경우로 한정할 수 있다. 이때
$\mathcal L|f^{-1}(U)$가 $f$에 관해 매우 풍부함을 보이면 된다.

$f'$은 고유이므로, 적당한 정수 $r>0$과 $Y'$-닫힌 몰입 사상
\[
j:X'\longrightarrow P=\mathbf P_{Y'}^r
\]
이 존재하여 $\mathcal L'$은 $j^*(\mathcal O_P(1))$과 동형이다
(II, 5.5.4, (ii)). 이 몰입 사상은 전사
$\mathcal O_{X'}$-준동형사상
\[
u:\mathcal O_{X'}^{r+1}\longrightarrow\mathcal L'
\]
에 표준적으로 대응한다 (II, 4.2.3). 후자는 $X'$ 위의
$\mathcal L'$의 $r+1$개 절단 $s'_i$ $(0\leq i\leq r)$로서 이
$\mathcal O_{X'}$-가군을 생성하는 것들의 데이터와 대응한다
(0\textsubscript{I}, 5.1.1). 정의에 의해 이 절단들은 $Y'$ 위의
$f'_*(\mathcal L')$의 절단들이기도 하다.
\[
f'_*(\mathcal L')
=f_*(\mathcal L)\otimes_{\mathcal O_Y}\mathcal O_{Y'}
\]
이다 (0\textsubscript{I}, 5.4.10). $Y$는 아핀이고
$\mathcal O_y$는 $A$의 소아이디얼 $\mathfrak j_y$에서의
국소환이므로
\[
s'_i=s''_i/t_i,
\]
여기서 $s''_i$들은 $Y$ 위의 $f_*(\mathcal L)$의 절단들이고
$t_i$들은 $\mathfrak j_y$에 속하지 않는 $A$의 원소들이다.
따라서 $Y$ 안의 $y$의 어떤 아핀 열린 근방 $V$와
$f_*(\mathcal L)|V$의 절단 $s_i$들이 존재하여 $s'_i=s_i/1$이다.
여기서 공간 $Y'$은 $V$에 포함됨을 상기한다 (I, 2.4.2).

$s_i$들은 $f^{-1}(V)$ 위의 $\mathcal L$의 절단들이므로
준동형사상
\[
v:\bigl(\mathcal O_X|f^{-1}(V)\bigr)^{r+1}
\longrightarrow\mathcal L|f^{-1}(V)
\]
을 정의한다. 가정에 의해 이는 $f^{-1}(y)$의 모든 점에서 전사이다.
$\operatorname{Coker}(v)$는 연접이므로
(0\textsubscript{I}, 5.3.4), 그 받침은 닫혀 있다
(0\textsubscript{I}, 5.2.2). 따라서 $f^{-1}(y)$의 어떤 열린 근방
$W\subset f^{-1}(V)$가 존재하여 $v$의 $W$에 대한 제한은 전사이다.
$f$는 닫힌 사상이므로 $W=f^{-1}(U)$라고 가정할 수 있다. 여기서
$U$는 $y$의 열린 근방이다. 이제 (II, 4.2.3)에서 결론이 따른다.
\end{proof}

\subsection{형식적 준스킴의 유한 사상}

\begin{proposition}[4.8.1]
\label{III.4.8.1-fr}
$\mathfrak Y$를 국소 뇌터 형식적 준스킴, $\mathcal K$를
$\mathfrak Y$의 정의 아이디얼, $f:\mathfrak X\to\mathfrak Y$를
형식적 준스킴의 사상이라 하자. 다음 조건들은 서로 동치이다.
\begin{enumerate}
\item[a)]
$\mathfrak X$는 국소 뇌터이고 $f$는 아딕 사상이며 (I, 10.12.1),
$\mathcal J=f^*(\mathcal K)\mathcal O_{\mathfrak X}$로 놓을 때
$f$에서 유도되는 사상
\[
f_0:(\mathfrak X,\mathcal O_{\mathfrak X}/\mathcal J)
\longrightarrow
(\mathfrak Y,\mathcal O_{\mathfrak Y}/\mathcal K)
\]
은 유한이다.
\item[b)]
$\mathfrak X$는 국소 뇌터이고, $X_0\to Y_0$가 유한인 어떤 아딕
$(Y_n)$-귀납계 $(X_n)$의 귀납극한이다.
\item[c)]
$\mathfrak Y$의 모든 점은 어떤 뇌터 아핀 형식적 열린 근방 $V$를
가지며, $f^{-1}(V)$는 아핀 형식적 열린집합이고
\[
\Gamma(f^{-1}(V),\mathcal O_{\mathfrak X})
\]
는 $\Gamma(V,\mathcal O_{\mathfrak Y})$ 위의 유한 생성 가군이다.
\end{enumerate}
\end{proposition}

\begin{proof}
(I, 10.12.3)에 의해 a)가 b)를 함의함은 즉시 따른다. b)가 c)를
함의함을 보이기 위해
\[
\mathfrak Y=\operatorname{Spf}(B),\qquad
\mathcal K=\mathfrak R^\Delta
\]
라고 가정할 수 있다. 여기서 $B$는 뇌터 아딕 환이고
$\mathfrak R$은 $B$의 정의 아이디얼이다. 가정에 의해 $X_0$은
아핀 스킴이고 그 환 $A_0$은 유한 생성
$B/\mathfrak R$-가군이다 (II, 6.1.3). (I, 5.1.9)에 의해 각
$X_n$은 아핀 스킴이다. 그 환을 $A_n$이라 하면, 가정 b)에 의해
$m\leq n$일 때
\[
A_m\simeq A_n/\mathfrak R^{m+1}A_n.
\]
따라서 $\mathfrak X$는 $\operatorname{Spf}(A)$와 동형이고, 여기서
$A=\varprojlim A_n$이다. 이제 (0\textsubscript{I}, 7.2.9)에서
결론이 따른다. 끝으로 c)가 a)를 함의함을 보이기 위해 다시

\oldpage[III]{147}

$\mathfrak Y=\operatorname{Spf}(B)$,
$\mathfrak X=\operatorname{Spf}(A)$이고 $A$가 유한
$B$-대수인 경우로 한정할 수 있다. 이때
$A/\mathfrak RA$는 유한 $B/\mathfrak R$-대수이므로
(I, 10.10.9)에 의해 a)의 조건들이 성립한다.
\end{proof}

\begin{definition}[4.8.2]
\label{III.4.8.2-fr}
(\hyperref[III.4.8.1-fr]{4.8.1})의 서로 동치인 성질 a), b), c)가
성립할 때 사상 $f$를 \emph{유한}이라 하고, 또는
$\mathfrak X$를 \emph{유한 $\mathfrak Y$-형식적 준스킴} 혹은
\emph{$\mathfrak Y$ 위에서 유한인 형식적 준스킴}이라 한다.
\end{definition}

\begin{proposition}[4.8.3]
\label{III.4.8.3-fr}
\begin{enumerate}
\item[(i)]
국소 뇌터 형식적 준스킴의 닫힌 몰입 사상은 유한 사상이다.
\item[(ii)]
국소 뇌터 형식적 준스킴의 두 유한 사상의 합성은 유한 사상이다.
\item[(iii)]
$\mathfrak X$, $\mathfrak Y$, $\mathfrak S$를 세 국소 뇌터
형식적 준스킴, $f:\mathfrak X\to\mathfrak S$를 유한 사상,
$g:\mathfrak Y\to\mathfrak S$를 사상이라 하자. 그러면 사상
\[
\mathfrak X\times_{\mathfrak S}\mathfrak Y\longrightarrow\mathfrak Y
\]
은 유한이다.
\item[(iv)]
$\mathfrak S$를 국소 뇌터 형식적 준스킴,
$\mathfrak X'$, $\mathfrak Y'$을 두 국소 뇌터 형식적 준스킴이라
하되
$\mathfrak X'\times_{\mathfrak S}\mathfrak Y'$이 국소 뇌터라고
하자. $\mathfrak X$, $\mathfrak Y$가 국소 뇌터
$\mathfrak S$-형식적 준스킴이고
$f:\mathfrak X\to\mathfrak X'$,
$g:\mathfrak Y\to\mathfrak Y'$이 두 유한 $\mathfrak S$-사상이면
$f\times_{\mathfrak S}g$는 유한 사상이다.
\item[(v)]
$f:\mathfrak X\to\mathfrak Y$,
$g:\mathfrak Y\to\mathfrak Z$를 국소 뇌터 형식적 준스킴의 두
사상이라 하고, $g$가 유한형이고 분리라고 하자. $g\circ f$가
유한 사상이면 $f$도 유한 사상이다.
\end{enumerate}
\end{proposition}

\begin{proof}
(i)은 자명하다. 나머지 주장들은
(\hyperref[III.4.8.1-fr]{4.8.1})의 판정법 a)를 이용하면 보통
준스킴의 사상에 대한 대응하는 명제들 (II, 6.1.5)로 곧바로
귀착된다. 세부사항은 (I, 10.13.5)를 본떠 독자에게 맡긴다.
\end{proof}

\begin{corollary}[4.8.4]
\label{III.4.8.4-fr}
(I, 10.9.9)의 가정 아래에서 $f$가 유한 사상이면, 그 완비화들로의
연장 $\widehat f$도 유한 사상이다.
\end{corollary}

\begin{corollary}[4.8.5]
\label{III.4.8.5-fr}
$\mathfrak X$가 $\mathfrak Y$ 위에서 유한인 형식적 준스킴이고
$f:\mathfrak X\to\mathfrak Y$가 구조 사상이면, 모든 열린집합
$U\subset\mathfrak Y$에 대해 $f^{-1}(U)$는 $U$ 위에서 유한이다.
\end{corollary}

\begin{proposition}[4.8.6]
\label{III.4.8.6-fr}
$f:\mathfrak X\to\mathfrak Y$가 국소 뇌터 형식적 준스킴의 유한
사상이면 $f_*(\mathcal O_{\mathfrak X})$는 연접
$\mathcal O_{\mathfrak Y}$-대수이다.
\end{proposition}

\begin{proof}
$f$를 사상 $f_n:X_n\to Y_n$들의 귀납계 $(f_n)$의 귀납극한으로
볼 수 있다. $f_n$들이 유한 사상이고
$f_*(\mathcal O_{\mathfrak X})$가
$(f_n)_*(\mathcal O_{X_n})$들의 사영극한과 동형임을 보이자.
그러면 (I, 10.10.5)에 의해 주장이 성립한다.
$\mathfrak Y=\operatorname{Spf}(B)$,
$\mathfrak X=\operatorname{Spf}(A)$인 경우로 한정하고 다음을
관찰하면 충분하다. $\mathfrak R$이 $B$의 정의 아이디얼이고
$A$가 유한 생성 $B$-가군이면
$A/\mathfrak R^{n+1}A$는
$B/\mathfrak R^{n+1}B$ 위의 유한 생성 가군이고, $A$는
$A/\mathfrak R^{n+1}A$들의 사영극한이다.
\end{proof}

역으로 다음이 성립한다.

\begin{proposition}[4.8.7]
\label{III.4.8.7-fr}
$\mathfrak Y$를 국소 뇌터 형식적 준스킴,
$\mathcal A$를 연접 $\mathcal O_{\mathfrak Y}$-대수라 하자.
$\mathfrak Y$-동형사상을 제외하고 유일한, $\mathfrak Y$ 위에서
유한인 형식적 준스킴 $\mathfrak X$가 존재하여
\[
f_*(\mathcal O_{\mathfrak X})=\mathcal A
\]
이다. 여기서 $f:\mathfrak X\to\mathfrak Y$는 구조 사상이다.
\end{proposition}

\begin{proof}
$\mathcal K$를 $\mathfrak Y$의 정의 아이디얼이라 하고
\[
Y_n=(\mathfrak Y,\mathcal O_{\mathfrak Y}/\mathcal K^{n+1}),
\qquad
\mathcal A_n=\mathcal A/\mathcal K^{n+1}\mathcal A
\]
로 놓자. $\mathcal A_n$은 분명히 유한
$\mathcal O_{Y_n}$-대수이고, 따라서 유한 $Y_n$-준스킴

\oldpage[III]{148}

$X_n=\operatorname{Spec}(\mathcal A_n)$을 정의한다 (II, 6.1.3).
$m\leq n$일 때 표준 전사 준동형사상
$h_{mn}:\mathcal A_n\to\mathcal A_m$은 사상
$u_{mn}:X_m\to X_n$을 정의하며, 이 사상은 구조 사상을 $f_n$이라
할 때 그림
\[
\xymatrix{
X_n\ar[d]_{f_n} & X_m\ar[l]_{u_{mn}}\ar[d]^{f_m}\\
Y_n & Y_m\ar[l]
}
\]
을 가환으로 만들고 $X_m$을
$X_n\times_{Y_n}Y_m$과 동일시한다 (II, 1.4.6).
귀납계 $(X_n)$의 귀납극한인 형식적 준스킴 $\mathfrak X$는 국소
뇌터이고, 귀납계 $(f_n)$의 귀납극한인 구조 사상
$f:\mathfrak X\to\mathfrak Y$는 유한이다
(\hyperref[III.4.8.1-fr]{4.8.1} 및 II, 10.12.3.1).
또한 (\hyperref[III.4.8.6-fr]{4.8.6})의 증명에서
$f_*(\mathcal O_{\mathfrak X})$가 $\mathcal A_n$들의 사영극한임을
보았으므로 이는 $\mathcal A$와 같다 (I, 10.10.6).
유일성 주장은 다음의 더 일반적인 결과의 귀결이다.
\end{proof}

\begin{proposition}[4.8.8]
\label{III.4.8.8-fr}
$\mathfrak Y$를 국소 뇌터 형식적 준스킴,
$\mathfrak X$, $\mathfrak X'$을 $\mathfrak Y$ 위에서 유한인 두
$\mathfrak Y$-형식적 준스킴이라 하고,
$f:\mathfrak X\to\mathfrak Y$,
$f':\mathfrak X'\to\mathfrak Y$을 구조 사상들이라 하자. 다음 두
집합 사이에는 표준 전단사 대응이 존재한다.
\[
\operatorname{Hom}_{\mathfrak Y}(\mathfrak X,\mathfrak X')
\quad\hbox{와}\quad
\operatorname{Hom}_{\mathcal O_{\mathfrak Y}}
\bigl(f'_*(\mathcal O_{\mathfrak X'}),
f_*(\mathcal O_{\mathfrak X})\bigr)
\footnotemark.
\]
\footnotetext{마지막 표현은
$\mathcal O_{\mathfrak Y}$-대수 준동형사상
$f'_*(\mathcal O_{\mathfrak X'})
\to f_*(\mathcal O_{\mathfrak X})$들의 집합을 뜻한다.}
\end{proposition}

\begin{proof}
이 대응 $h\mapsto\mathcal A(h)$의 정의는 (II, 1.1.2)에서와 같다.
그것이 전단사임을 보이기 위해
$\mathfrak Y=\operatorname{Spf}(B)$가 뇌터 아핀 형식적 스킴인
경우로 곧바로 귀착된다. 이때
\[
\mathfrak X=\operatorname{Spf}(A),\qquad
\mathfrak X'=\operatorname{Spf}(A'),
\]
여기서 $A$, $A'$은 두 유한 $B$-대수이고
\[
f_*(\mathcal O_{\mathfrak X})=A^\Delta,\qquad
f'_*(\mathcal O_{\mathfrak X'})=A'^\Delta.
\]
한편 $\mathfrak Y$-사상
$\mathfrak X\to\mathfrak X'$과 대수 준동형사상인
$B$-준동형사상 $A'\to A$ 사이에는 전단사 대응이 있다. 이
준동형사상들은 반드시 연속이다 (I, 10.2.2). 다른 한편
$B$-가군 준동형사상 $A'\to A$와
$\mathcal O_{\mathfrak Y}$-가군 준동형사상
$A'^\Delta\to A^\Delta$ 사이에도 전단사 대응이 있다
(I, 10.10.2.3). 이 두 대응에서 결론이 따른다.
\end{proof}

\begin{corollary}[4.8.9]
\label{III.4.8.9-fr}
(\hyperref[III.4.8.8-fr]{4.8.8})에서 정의된 표준 전단사 대응에서
닫힌 몰입 사상 $\mathfrak X\to\mathfrak X'$은 전사
$\mathcal O_{\mathfrak Y}$-대수 준동형사상
\[
f'_*(\mathcal O_{\mathfrak X'})
\longrightarrow f_*(\mathcal O_{\mathfrak X})
\]
과 대응한다.
\end{corollary}

\begin{proof}
문제는 다시 $\mathfrak Y$ 위에서 국소적이므로 국소 뇌터 형식적
준스킴의 닫힌 몰입 사상의 정의로 귀착된다 (I, 10.14.2).
\end{proof}

\begin{corollary}[4.8.10]
\label{III.4.8.10-fr}
(\hyperref[III.4.8.1-fr]{4.8.1})의 표기와 가정 아래에서, 아딕 사상
$f$가 닫힌 몰입 사상이기 위한 필요충분조건은 $f_0$가 보통
준스킴의 닫힌 몰입 사상인 것이다.
\end{corollary}

\begin{proof}
이는 (\hyperref[III.4.8.9-fr]{4.8.9})와 연접
$\mathcal O_{\mathfrak Y}$-가군의 준동형사상에 대한 전사 조건
(I, 10.11.5)에서 곧바로 따른다.
\end{proof}

\begin{proposition}[4.8.11]
\label{III.4.8.11-fr}
국소 뇌터 형식적 준스킴의 사상
$f:\mathfrak X\to\mathfrak Y$가 유한이기 위한 필요충분조건은

\oldpage[III]{149}

$f$가 고유이고 모든 $y\in\mathfrak Y$에 대해 올
$f^{-1}(y)$가 유한인 것이다.
\end{proposition}

\begin{proof}
정의 (3.4.1 및 4.8.2)에 의해 문제는
(\hyperref[III.4.8.1-fr]{4.8.1})의 표기에서 $f_0$에 대한 같은
명제로 곧바로 귀착되며, 이는 바로
(\hyperref[III.4.4.2-fr]{4.4.2})이다.
\end{proof}

\section{연접 대수적 층의 존재 정리}
\label{section:III.5-fr}

\subsection{정리의 진술}
\label{subsection:III.5.1-fr}

\begin{env}[5.1.1]
\label{III.5.1.1-fr}
$A$를 뇌터 아딕 환, $\mathfrak J$를 $A$의 정의 아이디얼이라 하자.
그러면 $A$는 $\mathfrak J$-아딕 위상에 대해 분리이고 완비이다.
$Y=\operatorname{Spec}(A)$이면 아핀 형식적 스킴
$\operatorname{Spf}(A)$는 닫힌부분집합
$Y'=\mathrm V(\mathfrak J)$을 따른 $Y$의 완비화
$\widehat Y$와 동일시된다 (I, 10.10.1).
$X$를 유한형 $Y$-준스킴, $f:X\to Y$를 구조 사상이라 하자.
$X'=f^{-1}(Y')$을 따른 $X$의 완비화, 또는 같은 것으로서
$\widehat Y$-형식적 준스킴 $X\times_Y\widehat Y$를
$\widehat X$로 나타내겠다. $f$의 완비화들로의 연장을
$\widehat f:\widehat X\to\widehat Y$로 나타내겠다. 끝으로 모든
연접 $\mathcal O_X$-가군 $\mathcal F$에 대해 그 완비화
$\mathcal F_{/X'}$를 $\widehat{\mathcal F}$로 나타내겠다. 이는
연접 $\mathcal O_{\widehat X}$-가군이다.
\end{env}

\begin{proposition}[5.1.2]
\label{III.5.1.2-fr}
(\hyperref[III.5.1.1-fr]{5.1.1})의 가정과 표기 아래에서,
$\mathcal F$를 그 받침이 $Y$ 위에서 고유인 연접
$\mathcal O_X$-가군이라 하자 (II, 5.4.10). 그러면 표준
준동형사상들 (4.1.4)
\[
\rho_i:\mathrm H^i(X,\mathcal F)\longrightarrow
\mathrm H^i(\widehat X,\widehat{\mathcal F})
\]
은 동형사상이다.
\end{proposition}

\begin{proof}
$\mathrm H^i(X,\mathcal F)$는 유한 생성 $A$-가군이고 (3.2.4),
따라서 자신의 $\mathfrak J$-준아딕 위상에 대한 분리 완비화와
같다 (0\textsubscript{I}, 7.3.6). 그러므로 이 명제는
(\hyperref[III.4.1.10-fr]{4.1.10})의 특수한 경우일 뿐이다.
\end{proof}

연접 $\mathcal O_X$-가군의 모든 완전열
\[
0\longrightarrow\mathcal F'\longrightarrow\mathcal F
\longrightarrow\mathcal F''\longrightarrow0
\]
에 대해 표준 동형사상 $\rho_i$들이 연결 준동형사상들과 가환함을
상기하자 (0, 12.1.6 및 I, 10.8.9).

\begin{corollary}[5.1.3]
\label{III.5.1.3-fr}
$\mathcal F$, $\mathcal G$를 두 연접 $\mathcal O_X$-가군이라 하되,
그 받침들의 교집합이 $Y$ 위에서 고유라고 하자. 그러면 모든
준동형사상 $u:\mathcal F\to\mathcal G$에 그 완비화
$\widehat u:\widehat{\mathcal F}\to\widehat{\mathcal G}$를
대응시키는 표준 준동형사상
\begin{equation}
\operatorname{Hom}_{\mathcal O_X}(\mathcal F,\mathcal G)
\longrightarrow
\operatorname{Hom}_{\mathcal O_{\widehat X}}
(\widehat{\mathcal F},\widehat{\mathcal G})
\tag{5.1.3.1}
\label{III.5.1.3.1.eq-fr}
\end{equation}
은 동형사상이다. 더 나아가 $f$가 닫힌 사상이면
$\widehat u$가 단사이기 위한 (각각 전사이기 위한) 필요충분조건은
$u$가 단사인 것이다 (각각 전사인 것이다).
\end{corollary}

\begin{proof}
첫째 주장은 (\ref{III.5.1.3.1.eq-fr})의 왼쪽 항이 유한 생성
$A$-가군이고 따라서 자신의 분리 완비화와 같다는 사실에 의한
(\hyperref[III.4.5.3-fr]{4.5.3})의 특수한 경우이다.
둘째 주장을 보이기 위해 (I, 10.8.14)에 의해
$\widehat u$가 단사이기 위한 (각각 전사이기 위한) 필요충분조건은
$X'$의 어떤 근방에서 $u$가 단사인 것이다 (각각 전사인 것이다).
따라서 결론은 다음 보조정리에서 따른다.
\end{proof}

\oldpage[III]{150}

\begin{lemma}[5.1.3.1]
\label{III.5.1.3.1-fr}
(\hyperref[III.5.1.1-fr]{5.1.1})의 가정 아래에서 $f$가 닫힌
사상이라고 더 가정하면, $X$ 안의 $X'$의 모든 근방은 $X$와 같다.
\end{lemma}

\begin{proof}
먼저 $f(X)=Y$인 경우로 귀착될 수 있다. 실제로 가정에 의해
$f(X)$는 $Y$의 닫힌부분집합 $Z$이다. 더 나아가 $f$를
$f_{\mathrm{red}}$로 바꾸어 (I, 6.3.4) $X$와 $Y$가 축약이라고
가정할 수 있다. 그러면 $Y$를 바탕공간이 $Z$인 $Y$의 축약 닫힌
부분준스킴으로 바꿀 수 있다 (I, 5.2.2). $A$의 모든 아이디얼은
닫혀 있으므로 $A$의 모든 몫환은 아딕이고 뇌터이기 때문이다.

이제 $f(X')=Y'$이다. $V$가 $X$ 안의 $X'$의 열린 근방이면
$f(X-V)$는 가정에 의해 $Y$에서 닫혀 있고 $Y'$과 만나지 않는다.
그런데 $\mathfrak J$가 $A$의 야코브슨 근기에 포함되므로
(I, 1.1.15 및 0\textsubscript{I}, 7.1.10), 이는 $X-V$가
공집합이지 않는 한 불가능하다. 따라서 $V=X$이다.
\end{proof}

받침이 $Y$ 위에서 고유인 연접 $\mathcal O_X$-가군들로 한정하면,
(\hyperref[III.5.1.3-fr]{5.1.3})은 범주의 언어로 다음과 같이
말할 수 있다. 함자
\[
\mathcal F\longmapsto\widehat{\mathcal F}
\]
는 그러한 $\mathcal O_X$-가군들의 범주에서 연접
$\mathcal O_{\widehat X}$-가군들의 범주로 가는 충만하고 충실한
함자이고, 따라서 전자의 범주와 후자의 어떤 충만한 부분범주
사이의 동치를 준다 (0, 8.1.6). 존재 정리는 $X$가 $Y$ 위에서
고유일 때 이 부분범주가 실제로 모든 연접
$\mathcal O_{\widehat X}$-가군들의 범주임을 보일 것이다. 정확히는
다음과 같다.

\begin{theorem}[5.1.4]
\label{III.5.1.4-fr}
$A$를 뇌터 아딕 환, $Y=\operatorname{Spec}(A)$,
$\mathfrak J$를 $A$의 정의 아이디얼,
$Y'=\mathrm V(\mathfrak J)$이라 하자.
$f:X\to Y$를 분리이고 유한형인 사상,
$X'=f^{-1}(Y')$이라 하자. $Y'$과 $X'$을 따른 $Y$와 $X$의
완비화들을
\[
\widehat Y=Y_{/Y'}=\operatorname{Spf}(A),\qquad
\widehat X=X_{/X'}
\]
로 나타내고, $f$의 완비화들로의 연장을
$\widehat f:\widehat X\to\widehat Y$라 하자. 그러면 함자
\[
\mathcal F\longmapsto
\mathcal F_{/X'}=\widehat{\mathcal F}
\]
는 받침이 $\operatorname{Spec}(A)$ 위에서 고유인 연접
$\mathcal O_X$-가군들의 범주와, 받침이
$\operatorname{Spf}(A)$ 위에서 고유인 연접
$\mathcal O_{\widehat X}$-가군들의 범주 사이의 동치이다.
\end{theorem}

다시 말해, (\hyperref[III.5.1.3-fr]{5.1.3})을 고려하면 다음과
같다.

\begin{corollary}[5.1.5]
\label{III.5.1.5-fr}
$\mathcal O_{\widehat X}$-가군이 받침이
$\operatorname{Spec}(A)$ 위에서 고유인 어떤 연접
$\mathcal O_X$-가군의 완비화와 동형이기 위한 필요충분조건은,
그 가군이 연접이고 그 받침이 $\operatorname{Spf}(A)$ 위에서
고유인 것이다.
\end{corollary}

가장 중요한 경우는 다음이다.

\begin{corollary}[5.1.6]
\label{III.5.1.6-fr}
$X$가 $Y=\operatorname{Spec}(A)$ 위에서 고유라고 가정하자.
그러면 함자
$\mathcal F\mapsto\widehat{\mathcal F}$는 연접
$\mathcal O_X$-가군들의 범주와 연접
$\mathcal O_{\widehat X}$-가군들의 범주 사이의 동치이다.
\end{corollary}

\begin{env}[5.1.7]
\label{III.5.1.7-fr}
\emph{보충.}
형식적 준스킴 위의 연접층의 특성화 (I, 10.11.3)를 고려하면,
(\hyperref[III.5.1.1-fr]{5.1.1})의 조건 아래에서 받침이
$\operatorname{Spec}(A)$ 위에서 고유인 연접
$\mathcal O_X$-가군을 주는 것은 다음 데이터를 주는 것과 동치임을
알 수 있다. 다음과 같이 놓자.
\[
Y_n=\operatorname{Spec}(A/\mathfrak J^{n+1}),\qquad
X_n=X\times_Y Y_n.
\]
각 $\mathcal F_n$이 연접 $\mathcal O_{X_n}$-가군인 사영계
$(\mathcal F_n)$을 주되, $m\leq n$일 때
\[
\mathcal F_m
=\mathcal F_n\otimes_{\mathcal O_{Y_n}}\mathcal O_{Y_m}
\quad\text{(또는 같은 것으로서 }\mathcal F_m
=\mathcal F_n/\mathfrak J^{m+1}\mathcal F_n\text{)}
\]
이고 $\mathcal F_0$의 받침이 $Y_0$ 위에서 고유인 $X_0$의
부분집합이어야 한다. (I, 10.11.4)를 이용하면 연접
$\mathcal O_X$-가군의 준동형사상들도 마찬가지로 연접
$\mathcal O_{X_n}$-가군들의 사영계 사이의 준동형사상으로
해석된다.

\oldpage[III]{151}

알려진 모든 응용의 경우 $A$는 실제로 뇌터 아딕 국소환이다.
따라서 $Y_n$들은 아르틴 국소환의 스펙트럼이고, 이 절과 앞 절들의
결과들은 뇌터 아딕 국소환 위의 대수기하학을 상당한 정도로 아르틴
국소환 위의 대수기하학으로 환원한다.
\end{env}

\begin{corollary}[5.1.8]
\label{III.5.1.8-fr}
(\hyperref[III.5.1.4-fr]{5.1.4})의 조건 아래에서 대응
\[
Z\longmapsto\widehat Z=Z_{/(Z\cap X')}
\]
은 $Y$ 위에서 고유인 $X$의 닫힌 부분준스킴들의 집합에서
$\widehat Y$ 위에서 고유인 $\widehat X$의 닫힌 형식적
부분준스킴들의 집합으로 가는 전단사이다.
\end{corollary}

\begin{proof}
실제로 $\widehat X$의 닫힌 형식적 부분준스킴은
\[
(T,(\mathcal O_{\widehat X}/\mathcal A)|T)
\]
의 꼴이다. 여기서 $\mathcal A$는 $\mathcal O_{\widehat X}$의 연접
아이디얼이다 (I, 10.14.2). $T$가 $\widehat Y$ 위에서 고유이면
(\hyperref[III.5.1.4-fr]{5.1.4})에 의해
$\mathcal O_{\widehat X}/\mathcal A$는
$\widehat{\mathcal F}$ 꼴의 $\mathcal O_{\widehat X}$-가군과
동형이다. 여기서 $\mathcal F$는 받침이 $Y$ 위에서 고유인 연접
$\mathcal O_X$-가군이다. 또한
(\hyperref[III.5.1.3-fr]{5.1.3})에 의해 표준 준동형사상
\[
\mathcal O_{\widehat X}\longrightarrow
\mathcal O_{\widehat X}/\mathcal A
\]
은 $\widehat u$의 꼴이고, 여기서
$u:\mathcal O_X\to\mathcal F$는 전사
$\mathcal O_X$-가군 준동형사상이다. 따라서
$\mathcal F=\mathcal O_X/\mathcal N$의 꼴이고, 여기서
$\mathcal N$은 $\mathcal O_X$의 연접 아이디얼이다. 또한
$\mathcal A=\widehat{\mathcal N}$이다 (I, 10.8.8).
그러므로 결론은 (I, 10.14.7)에서 따른다.
\end{proof}

\subsection{존재 정리의 증명: 사영인 경우와 준사영인 경우}
\label{subsection:III.5.2-fr}

\begin{env}[5.2.1]
\label{III.5.2.1-fr}
(\hyperref[III.5.1.4-fr]{5.1.4})의 조건 아래에서, 연접
$\mathcal O_{\widehat X}$-가군이 받침이 $Y$ 위에서 고유인 어떤
연접 $\mathcal O_X$-가군 $\mathcal F$의 완비화
$\widehat{\mathcal F}$와 동형일 때 잠정적으로 그 가군을
\emph{대수화 가능}이라고 하겠다.
\end{env}

\begin{lemma}[5.2.2]
\label{III.5.2.2-fr}
$\mathcal F'$, $\mathcal G'$을 두 대수화 가능한
$\mathcal O_{\widehat X}$-가군이라 하자. 모든 준동형사상
$u:\mathcal F'\to\mathcal G'$에 대해
$\operatorname{Ker}(u)$, $\operatorname{Im}(u)$,
$\operatorname{Coker}(u)$는 대수화 가능하다.
\end{lemma}

\begin{proof}
실제로
\[
\mathcal F'=\widehat{\mathcal F},\qquad
\mathcal G'=\widehat{\mathcal G}
\]
이고 $\mathcal F$, $\mathcal G$가 받침이 $Y$ 위에서 고유인 연접
$\mathcal O_X$-가군들이면, (\hyperref[III.5.1.3-fr]{5.1.3})에
의해 어떤 준동형사상 $v:\mathcal F\to\mathcal G$에 대해
$u=\widehat v$이다. 함자
$\mathcal F\mapsto\widehat{\mathcal F}$의 완전성에 의해
$\operatorname{Ker}(\widehat v)$는
$\widehat{\operatorname{Ker}(v)}$와 동형이다. 또한
$\operatorname{Ker}(v)$의 받침은 $\mathcal F$의 받침에
포함되므로 $\operatorname{Ker}(u)$는 대수화 가능하다.
$\operatorname{Im}(u)$와 $\operatorname{Coker}(u)$에도 같은
논증을 적용한다.
\end{proof}

\begin{proposition}[5.2.3]
\label{III.5.2.3-fr}
$A$를 뇌터 아딕 환, $\mathfrak J$를 $A$의 정의 아이디얼,
$\mathfrak Y=\operatorname{Spf}(A)$,
$f:\mathfrak X\to\mathfrak Y$를 형식적 준스킴의 고유 사상이라
하자. 다음과 같이 놓자.
\[
Y_k=\operatorname{Spec}(A/\mathfrak J^{k+1}),\qquad
X_k=\mathfrak X\times_{\mathfrak Y}Y_k,
\]
그리고 모든 $\mathcal O_{\mathfrak X}$-가군 $\mathcal F$에 대해
\[
\mathcal F_k
=\mathcal F\otimes_{\mathcal O_{\mathfrak X}}\mathcal O_{X_k}
=\mathcal F/\mathfrak J^{k+1}\mathcal F.
\]
$\mathcal L$을 가역 $\mathcal O_{\mathfrak X}$-가군이라 하고
$\mathcal L_0=\mathcal L/\mathfrak J\mathcal L$이 풍부한
$\mathcal O_{X_0}$-가군이라고 가정하자. 모든
$\mathcal O_{\mathfrak X}$-가군 $\mathcal F$와 모든 정수 $n$에
대해
\[
\mathcal F(n)=\mathcal F\otimes\mathcal L^{\otimes n}
\]
으로 놓자. 그러면 모든 연접 $\mathcal O_{\mathfrak X}$-가군
$\mathcal F$에 대해 어떤 정수 $n_0$가 존재하여, 모든
$n\geq n_0$에 대해 다음 성질들이 성립한다.
\begin{enumerate}
\item[(i)]
모든 $k\geq0$에 대해 표준 준동형사상
\[
\mathrm H^0(\mathfrak X,\mathcal F(n))
\longrightarrow
\mathrm H^0(\mathfrak X,\mathcal F_k(n))
\]
은 전사이다.
\item[(ii)]
모든 $q>0$에 대해
$\mathrm H^q(\mathfrak X,\mathcal F(n))=0$이다.
\end{enumerate}
\end{proposition}

\begin{proof}
$\mathfrak X$와 $X_0$의 바탕공간들은 같음이 알려져 있다. 층
\[
\mathcal M_k
=\mathfrak J^k\mathcal F/\mathfrak J^{k+1}\mathcal F
\]
는 $\mathfrak J$에 의해 소멸되므로 연접
$\mathcal O_{X_0}$-가군으로 볼 수 있다
(0\textsubscript{I}, 5.3.10). 또한
\[
\mathcal M_k(n)
=\mathcal M_k\otimes_{\mathcal O_{X_0}}\mathcal L_0^{\otimes n}
\]
으로 놓으면 즉시
\[
\mathcal M_k(n)
=\mathfrak J^k\mathcal F(n)/
\mathfrak J^{k+1}\mathcal F(n)
\]
임을 알 수 있다. $\mathcal L_0$은
$f_0:X_0\to Y_0$에 관해 풍부하고

\oldpage[III]{152}

$f_0$는 고유이므로 $f_0$는 사영 사상이다 (II, 5.5.4).
$A$의 $\mathfrak J$-아딕 여과에 결부된 등급 $A$-대수
\[
S=\bigoplus_{k\geq0}
\mathfrak J^k/\mathfrak J^{k+1}
\]
를 생각하자. $A$가 뇌터환이므로 $S$는 유한형이다.
\[
\mathcal S'=f_0^*(\widetilde S),\qquad
\mathcal M=\bigoplus_{k\geq0}\mathcal M_k
\]
로 놓으면 $\mathcal S'$은 준연접 유한형
$\mathcal O_{X_0}$-대수이고, $\mathcal M$은 준연접 유한형 등급
$\mathcal S'$-가군이다. 실제로 $\mathcal F_0$가 연접이고
$\mathcal S'$-가군 $\mathcal M$을 생성한다.

따라서 정리 (\hyperref[III.2.4.1-fr]{2.4.1, (ii)})를 적용할 수
있다. 어떤 $n_0$가 존재하여 $n\geq n_0$이면 모든 $k$와 모든
$q>0$에 대해
\begin{equation}
\mathrm H^q(X_0,\mathcal M_k(n))=0.
\tag{5.2.3.1}
\label{III.5.2.3.1-fr}
\end{equation}
$\mathcal M_k(n)$을 이번에는 $\mathcal O_{\mathfrak X}$-가군으로
보아도 $q>0$, $n\geq n_0$일 때
\[
\mathrm H^q(\mathfrak X,\mathcal M_k(n))=0
\]
이다. 완전열
\[
0\longrightarrow
\frac{\mathfrak J^k\mathcal F(n)}
{\mathfrak J^{k+1}\mathcal F(n)}
\longrightarrow
\frac{\mathfrak J^h\mathcal F(n)}
{\mathfrak J^{k+1}\mathcal F(n)}
\longrightarrow
\frac{\mathfrak J^h\mathcal F(n)}
{\mathfrak J^k\mathcal F(n)}
\longrightarrow0
\]
에 코호몰로지 완전열을 적용하자. 먼저 $0\leq h<k$,
$n\geq n_0$, $q>0$일 때 $k-h$에 관한 귀납법으로
\begin{equation}
\mathrm H^q\left(\mathfrak X,
\frac{\mathfrak J^h\mathcal F(n)}
{\mathfrak J^k\mathcal F(n)}\right)=0
\tag{5.2.3.2}
\label{III.5.2.3.2-fr}
\end{equation}
을 얻는다. 특히 $h=0$이면
\begin{equation}
\mathrm H^q(\mathfrak X,\mathcal F_k(n))=0
\qquad
(n\geq n_0,\ k\geq0,\ q>0).
\tag{5.2.3.3}
\label{III.5.2.3.3-fr}
\end{equation}
$h=0$일 때 코호몰로지 완전열의 다른 부분은 완전열
\begin{equation}
\mathrm H^0(\mathfrak X,\mathcal F_{k+1}(n))
\longrightarrow
\mathrm H^0(\mathfrak X,\mathcal F_k(n))
\longrightarrow
\mathrm H^1\left(\mathfrak X,
\frac{\mathfrak J^k\mathcal F(n)}
{\mathfrak J^{k+1}\mathcal F(n)}\right)=0
\tag{5.2.3.4}
\label{III.5.2.3.4-fr}
\end{equation}
을 준다. 따라서 $h\leq k$일 때 표준 사상
\begin{equation}
\mathrm H^0(\mathfrak X,\mathcal F_k(n))
\longrightarrow
\mathrm H^0(\mathfrak X,\mathcal F_h(n))
\tag{5.2.3.5}
\label{III.5.2.3.5-fr}
\end{equation}
은 전사이다.

그러므로 모든 $q$에 대해 사영계
\[
\bigl(\mathrm H^q(\mathfrak X,\mathcal F_k(n))\bigr)_{k\geq0}
\]
는 $n\geq n_0$일 때 조건 (ML)을 만족한다. 한편
$\mathfrak X$의 모든 아핀 형식적 열린집합 $U$는 각 $X_k$에서도
아핀 열린집합이므로 (I, 10.5.2), 모든 $q>0$에 대해
$\mathrm H^q(U,\mathcal F_k(n))=0$이고 (I, 1.3.1),
$h\leq k$일 때
\[
\mathrm H^0(U,\mathcal F_k(n))
\longrightarrow\mathrm H^0(U,\mathcal F_h(n))
\]
은 전사이다 (I, 1.3.9). 따라서 (0, 13.3.1)의 적용 조건들이
충족되며, $n\geq n_0$일 때 다음을 얻는다.
\begin{enumerate}
\item[$1^\circ$]
모든 $q>0$에 대해
\[
\mathrm H^q(\mathfrak X,\mathcal F(n))
\longrightarrow
\varprojlim_k\mathrm H^q(\mathfrak X,\mathcal F_k(n))
\]
은 전단사이다. 따라서 (\ref{III.5.2.3.3-fr})에 의해
$\mathrm H^q(\mathfrak X,\mathcal F(n))=0$이다.
\item[$2^\circ$]
준동형사상
\[
\mathrm H^0(\mathfrak X,\mathcal F(n))
\longrightarrow
\varprojlim_k\mathrm H^0(\mathfrak X,\mathcal F_k(n))
\]
은 전단사이다. 또한 준동형사상
(\ref{III.5.2.3.5-fr})들이 전사이므로 각 준동형사상
\[
\varprojlim_k\mathrm H^0(\mathfrak X,\mathcal F_k(n))
\longrightarrow
\mathrm H^0(\mathfrak X,\mathcal F_h(n))
\]
도 전사이다.
\end{enumerate}
이로써 증명이 끝난다.
\end{proof}

\begin{corollary}[5.2.4]
\label{III.5.2.4-fr}
(\hyperref[III.5.2.3-fr]{5.2.3})의 가정 아래에서 모든 연접
$\mathcal O_{\mathfrak X}$-가군 $\mathcal F$에 대해 어떤 정수
$N$이 존재하여, $n\geq N$이면 $\mathcal F(n)$은

\oldpage[III]{153}

$\mathfrak X$ 위의 자신의 절단들로 생성된다. 다시 말해
$\mathcal F$는
\[
(\mathcal O_{\mathfrak X}(-n))^k
\]
꼴의 $\mathcal O_{\mathfrak X}$-가군의 몫과 동형이다.
\end{corollary}

\begin{proof}
$X_0$가 뇌터이므로, $\mathcal L_0$에 대한 가정과 (II, 4.5.5)에
의해 어떤 $n_0$가 존재하여 $n\geq n_0$이면
$\mathcal F_0(n)$은 $\mathfrak X$ 위의 자신의 절단들로 생성된다.
또한 $n_0$를 충분히 크게 잡아 $n\geq n_0$일 때 준동형사상
\[
\Gamma(\mathfrak X,\mathcal F(n))
\longrightarrow
\Gamma(\mathfrak X,\mathcal F_0(n))
\]
이 전사가 되게 할 수 있다
(\hyperref[III.5.2.3-fr]{5.2.3}). 따라서 유한 개의 절단
$s_i\in\Gamma(\mathfrak X,\mathcal F(n))$가 존재하여, 그
$\Gamma(\mathfrak X,\mathcal F_0(n))$ 안의 상들이
$\mathcal F_0(n)$을 생성한다 (0\textsubscript{I}, 5.2.3).
$\mathfrak J$는 $\mathfrak X$의 모든 점에서 국소환의 극대
아이디얼에 포함되므로, 이 국소환들에 나카야마 보조정리를 적용하면
$s_i$들이 $\mathcal F(n)$을 생성함을 얻는다
(0\textsubscript{I}, 5.1.1).
\end{proof}

\begin{env}[5.2.5]
\label{III.5.2.5-fr}
\emph{존재 정리의 증명: 사영인 경우.}
(\hyperref[III.5.1.4-fr]{5.1.4})의 표기를 사용하고 $f$가
사영이라고 가정하자. 그러면 풍부한 $\mathcal O_X$-가군
$\mathcal L$이 존재한다 (I, 5.5.4). 정의에 의해
\[
X_n=\widehat X\times_{\widehat Y}Y_n
\]
은 $X$의 닫힌 부분준스킴
$X\times_Y Y_n=f^{-1}(Y_n)$과 같다.
$\mathcal L'=\mathcal L\otimes_{\mathcal O_X}
\mathcal O_{\widehat X}$이 $\mathcal L$의 완비화이면
$\mathcal L'_0=\mathcal L/\mathfrak J\mathcal L$이며, 이를
$\mathcal O_{X_0}$-가군으로 본다. $\mathcal L'_0$은 풍부함이
알려져 있다 (II, 4.6.13, (i bis)).

따라서 $\mathcal L'$과 임의의 연접
$\mathcal O_{\widehat X}$-가군 $\mathcal F$에 따름정리
(\hyperref[III.5.2.4-fr]{5.2.4})를 적용할 수 있다. 적당한 정수
$n>0$, $k>0$에 대해 $\mathcal F$는
\[
\mathcal G=(\mathcal L'^{\otimes(-n)})^k
\]
의 몫과 동형이다. 그런데 $\mathcal G$는 분명히
$(\mathcal L^{\otimes(-n)})^k$의 완비화이므로 (I, 10.8.10)
대수화 가능하다. 이제 표준 전사 준동형사상
$u:\mathcal G\to\mathcal F$를 생각하고
\[
\mathcal H=\operatorname{Ker}(u)
\]
로 놓자. 이는 연접 $\mathcal O_{\widehat X}$-가군이다
(0\textsubscript{I}, 5.3.4). 같은 방식으로 어떤 대수화 가능한
$\mathcal O_{\widehat X}$-가군 $\mathcal K$와 준동형사상
$v:\mathcal K\to\mathcal G$가 존재하여
$\mathcal H=\operatorname{Im}(v)$임을 알 수 있다. 그러면
$\mathcal F=\operatorname{Coker}(v)$이고,
(\hyperref[III.5.2.2-fr]{5.2.2})에 의해 $\mathcal F$는 대수화
가능하다.
\end{env}

\begin{env}[5.2.6]
\label{III.5.2.6-fr}
\emph{존재 정리의 증명: 준사영인 경우.}
계속해서 (\hyperref[III.5.1.4-fr]{5.1.4})의 표기를 사용하고,
이제 $f$가 준사영이라고 가정하자. 어떤 사영 사상
$g:Z\to Y$가 존재하여 $X$는 $Z$의 어떤 열린집합 위에 유도된
$Y$-준스킴과 동일시된다 (II, 5.3.2).
$Z'=g^{-1}(Y')$로 놓으면 $X'=X\cap Z'$이다. 따라서 완비화
$\widehat X=X_{/X'}$는 완비화
$\widehat Z=Z_{/Z'}$의 열린집합 $X\cap Z'$ 위에 유도된 형식적
준스킴과 동일시된다 (I, 10.8.5).

$\mathcal F'$을 연접 $\mathcal O_{\widehat X}$-가군이라 하고 그
받침 $T'$이 $\widehat Y$ 위에서 고유라고 하자. 이는 정의에 의해,
바탕공간이 $T'\subset X'$이고 $f$의 제한 $T'\to Y'$이 고유인
$X'$의 닫힌 부분준스킴이 존재한다는 뜻이다. 따라서 $T'$은
$Y$ 위에서 고유이고, 그러므로 $Z'$에서 닫혀 있다
(II, 5.4.10).

이로부터 $\mathcal F'$은 다음 연접
$\mathcal O_{\widehat Z}$-가군 $\mathcal G'$이
$\widehat X$ 위에 유도하는 층임이 따른다. 즉
$\widehat Z$의 열린집합 $\widehat X$ 위에 정의된
$\mathcal F'$과 열린집합 $\widehat Z-T'$ 위의 영층을 붙여
$\mathcal G'$을 얻는다. 이 두 층은 교집합
$\widehat X-T'$에서 일치한다. $\mathcal G'$이 연접임은
명백하다. (\hyperref[III.5.2.5-fr]{5.2.5})에 의해 어떤 연접
$\mathcal O_Z$-가군 $\mathcal G$가 존재하여
\[
\mathcal G'=\widehat{\mathcal G}.
\]
$T$를 $\mathcal G$의 받침이라 하면
$T'=T\cap Z'$이다 (I, 10.8.12). $h$를 바탕공간이 $T$인 $Z$의
축약 닫힌 부분준스킴에 대한 $g$의 제한이라 하면
\[
T'=h^{-1}(Y')=T\cap g^{-1}(Y').
\]
따라서 $X\cap T$는 $T'$을 포함하는 $T$의 열린집합이다.

\oldpage[III]{154}

$h$는 고유이고 (II, 5.4.2), 따라서 닫힌 사상이므로
(\hyperref[III.5.1.3.1-fr]{5.1.3.1})에 의해
$X\cap T=T$이다. 다시 말해 $T\subset X$이다. 또한 $T$는
$Z$에서 닫혀 있으므로 $Y$ 위에서 고유이다.
$\mathcal F$를 $\mathcal G$가 $X$ 위에 유도하는 층이라 하자.
그 완비화 $\widehat{\mathcal F}$은
$\widehat{\mathcal G}$가 $\widehat X$ 위에 유도하는 층이므로
(I, 10.8.4) $\mathcal F'$과 같다. 이로써 증명이 끝난다.
\end{env}

\subsection{존재 정리의 증명: 일반적인 경우}
\label{subsection:III.5.3-fr}

\begin{lemma}[5.3.1]
\label{III.5.3.1-fr}
(\hyperref[III.5.1.4-fr]{5.1.4})의 조건 아래에서,
\[
0\to\mathcal H\to\mathcal F\to\mathcal G\to0
\]
이 연접 $\mathcal O_{\widehat X}$-가군들의 완전열이고
$\mathcal G$와 $\mathcal H$가 대수화 가능하면, $\mathcal F$도
대수화 가능하다.
\end{lemma}

\begin{proof}
실제로 $\mathcal G=\widehat{\mathcal B}$,
$\mathcal H=\widehat{\mathcal C}$라고 하자. 여기서 $\mathcal B$와
$\mathcal C$는 받침이 $Y$ 위에서 고유인 연접
$\mathcal O_X$-가군들이다. $\mathcal F$는 표준적으로 $A$-가군
\[
\operatorname{Ext}^1_{\mathcal O_{\widehat X}}
(\widehat X;\widehat{\mathcal B},\widehat{\mathcal C})
\]
의 한 원소를 정한다 ($0$, 12.3.2). 가정에 의해 이 $A$-가군은
표준적으로
\[
\operatorname{Ext}^1_{\mathcal O_X}(X;\mathcal B,\mathcal C)
\]
와 동형이다 (4.5.2). 따라서 연접 $\mathcal O_X$-가군들의 완전열
\[
0\to\mathcal C\to\mathcal A\to\mathcal B\to0
\]
이 존재하여, $\mathcal A$에 대응하는
$\operatorname{Ext}^1_{\mathcal O_X}(X;\mathcal B,\mathcal C)$의
원소의 표준상은 $\mathcal F$에 대응하는
$\operatorname{Ext}^1_{\mathcal O_{\widehat X}}
(\widehat X;\widehat{\mathcal B},\widehat{\mathcal C})$의 원소이다.
그런데 정의에 의해, 그리고 (I, 10.8.8, (iii))를 고려하면, 이는
$\mathcal F$가 $\widehat{\mathcal A}$와 동형이라는 뜻이다. 또한
$\operatorname{Supp}(\mathcal A)$는
$\operatorname{Supp}(\mathcal B)$와
$\operatorname{Supp}(\mathcal C)$의 합집합에 포함되므로 $Y$ 위에서
고유이다. 이로써 보조정리가 증명된다.
\end{proof}

\begin{corollary}[5.3.2]
\label{III.5.3.2-fr}
(\hyperref[III.5.1.1-fr]{5.1.1})의 조건 아래에서,
$u:\mathcal F\to\mathcal G$를 연접
$\mathcal O_{\widehat X}$-가군들의 준동형사상이라 하자.
$\mathcal G$, $\operatorname{Ker}(u)$ 및
$\operatorname{Coker}(u)$가 대수화 가능하면 $\mathcal F$도
대수화 가능하다.
\end{corollary}

\begin{proof}
준동형사상 $\mathcal G\to\operatorname{Coker}(u)$에 보조정리
(\hyperref[III.5.2.2-fr]{5.2.2})를 적용하면
$\operatorname{Im}(u)$가 대수화 가능함을 알 수 있다. 이제 완전열
\[
0\to\operatorname{Ker}(u)\to\mathcal F
\to\operatorname{Im}(u)\to0
\]
에 보조정리 (\hyperref[III.5.3.1-fr]{5.3.1})를 적용하면 충분하다.
\end{proof}

\begin{lemma}[5.3.3]
\label{III.5.3.3-fr}
(\hyperref[III.5.1.1-fr]{5.1.1})의 조건 아래에서,
$h:Z\to Y$를 유한형 사상, $\widehat Z$를
$Z'=h^{-1}(Y')$을 따른 $Z$의 완비화, $g:Z\to X$를 고유
$Y$-사상, $\widehat g:\widehat Z\to\widehat X$를 완비화들로
연장한 사상이라 하자. 모든 대수화 가능한
$\mathcal O_{\widehat Z}$-가군 $\mathcal F'$에 대해
$\widehat g_*(\mathcal F')$은 대수화 가능한
$\mathcal O_{\widehat X}$-가군이다.
\end{lemma}

\begin{proof}
실제로 $\mathcal F'=\widehat{\mathcal F}$인 연접
$\mathcal O_Z$-가군 $\mathcal F$가 있으면, 제1 비교 정리
(4.1.5)에 의해 $\widehat g_*(\widehat{\mathcal F})$는
$g_*(\mathcal F)$의 완비화와 동형이다.
\end{proof}

\begin{lemma}[5.3.4]
\label{III.5.3.4-fr}
$X$를 (통상적인) 뇌터 스킴, $X'$을 $X$의 닫힌 부분,
$f:Z\to X$를 고유 사상, $Z'=f^{-1}(X')$라 하고,
\[
\widehat X=X_{/X'},\qquad \widehat Z=Z_{/Z'}
\]
라 하자. 또한 $\widehat f:\widehat Z\to\widehat X$를 $f$의
완비화들로의 연장이라 하자. $\mathcal M$을 다음 성질을 갖는
$\mathcal O_X$의 연접 아이디얼이라 하자. 즉
\[
U=X-\operatorname{Supp}(\mathcal O_X/\mathcal M)
\]
이면 $f$의 제한 $f^{-1}(U)\to U$는 동형사상이다. 그러면 모든
연접 $\mathcal O_{\widehat X}$-가군 $\mathcal F$에 대해 어떤 정수
$n>0$가 존재하여, 표준 준동형사상
\[
\rho_{\mathcal F}:\mathcal F\longrightarrow
\widehat f_*\bigl(\widehat f^{,*}(\mathcal F)\bigr)
\]
($0_{\mathrm I}$, 4.4.3)의 핵과 여핵은
$\widehat{\mathcal M}^{,n}$에 의해 소멸된다.
\end{lemma}

\begin{proof}
$B$가 뇌터 환인 $X=\operatorname{Spec}(B)$의 경우로 한정할 수
있다. 이때 어떤 $B$의 아이디얼 $\mathfrak K$에 대해
$X'=\mathrm V(\mathfrak K)$이다. 먼저 $B$가 뇌터 아딕 환이고
$\mathfrak K$가 $B$의 정의 아이디얼인 경우로 환원할 수 있음을
보이자. 실제로 $B_1$을 $\mathfrak K$-선아딕 위상에 대한 $B$의
분리 완비환이라 하자. $\mathfrak K_1=\mathfrak K B_1$이면
$B_1$은 $\mathfrak K_1$을 정의 아이디얼로 갖는

\oldpage[III]{155}

뇌터 아딕 환이다. $X_1=\operatorname{Spec}(B_1)$로 놓고,
$h:X_1\to X$를 표준 준동형사상 $B\to B_1$에 대응하는 사상이라
하자. $X'_1=h^{-1}(X')$이면
$X'_1=\mathrm V(\mathfrak K_1)$이다. 마지막으로
\[
Z_1=Z\times_X X_1=Z_{(X_1)},\qquad
f_1=f_{(X_1)}:Z_1\to X_1
\]
로 놓자. $f_1$은 고유 사상이다 (II, 5.4.2). $\widehat X_1$을
$X'_1$을 따른 $X_1$의 완비화,
\[
\widehat Z_1=Z_1\times_{X_1}\widehat X_1
\]
을 $Z'_1=f_1^{-1}(X'_1)$을 따른 $Z_1$의 완비화,
$\widehat f_1$을 $f_1$의 완비화들로의 연장이라 하자.

$h$의 완비화들로의 연장
$\widehat h:\widehat X_1\to\widehat X$는 $B_1$의 항등사상에
대응하는 동형사상임이 즉시 따른다 (I, 10.9.1). 따라서 이에
대응하는 준동형사상 $\widehat Z_1\to\widehat Z$도
동형사상이고, 이 동형사상들은 $\widehat f_1$과
$\widehat f$를 동일시한다. 한편
$\mathcal M_1=h^*(\mathcal M)$은 $\mathcal O_{X_1}$의 연접
아이디얼이고
\[
\operatorname{Supp}(\mathcal O_{X_1}/\mathcal M_1)
=h^{-1}(\operatorname{Supp}(\mathcal O_X/\mathcal M))
\]
이다 (I, 9.1.13). 따라서
\[
U_1=X_1-\operatorname{Supp}(\mathcal O_{X_1}/\mathcal M_1)
\]
로 놓으면 $U_1=h^{-1}(U)$이고, 곧바로 $f_1$의 제한
$f_1^{-1}(U_1)\to U_1$이 동형사상임을 알 수 있다
(I, 3.2.7). 또한 완비화 $\widehat{\mathcal M}$과
$\widehat{\mathcal M}_1$은 $\widehat h$에 의해 동일시된다
(I, 10.9.5). 그러므로 $X_1$, $X'_1$, $f_1$, $\mathcal M_1$은
(\hyperref[III.5.3.4-fr]{5.3.4})의 모든 가정을 만족한다. 이제부터
$B$가 뇌터 아딕 환이고 $\mathfrak K$가 $B$의 정의 아이디얼이라고
가정할 수 있다.

이 경우 $\widehat X=\operatorname{Spf}(B)$이고,
$N$이 유한형 $B$-가군인 $\mathcal F=N^\Delta$로 쓸 수 있다.
따라서 $\mathcal G$를 연접 $\mathcal O_X$-가군
$\widetilde N$이라 하면
$\mathcal F=\widehat{\mathcal G}$이다 (I, 10.10.5). 그러므로
\[
\widehat f^{,*}(\mathcal F)=\widehat{f^*(\mathcal G)}
\]
이다 (I, 10.9.5). 또한 제1 비교 정리 (4.1.5)에 의해
\[
\widehat f_*\bigl(\widehat{f^*(\mathcal G)}\bigr)
\simeq\widehat{f_*\bigl(f^*(\mathcal G)\bigr)},
\]
이고, (\hyperref[III.5.1.3-fr]{5.1.3})에 의해 표준 준동형사상
$\rho_{\mathcal F}$는 바로 $\widehat{\rho}_{\mathcal G}$이다.
그런데
\[
\rho_{\mathcal G}:\mathcal G\longrightarrow
f_*\bigl(f^*(\mathcal G)\bigr)
\]
의 핵 $\mathcal P$와 여핵 $\mathcal R$은 연접
$\mathcal O_X$-가군들이고, 가정에 의해 이들의 $U$에 대한 제한은
명백히 영이다. 따라서 어떤 정수 $n>0$가 존재하여
\[
\mathcal M^n\mathcal P=\mathcal M^n\mathcal R=0
\]
이다 (I, 9.3.4). 이로부터
\[
\widehat{\mathcal M}^{,n}\widehat{\mathcal P}=
\widehat{\mathcal M}^{,n}\widehat{\mathcal R}=0
\]
이 따른다 (I, 10.8.8 및 10.8.10).
\end{proof}

\begin{env}[5.3.5]
\label{III.5.3.5-fr}
\emph{존재 정리의 증명의 끝.}
(\hyperref[III.5.1.4-fr]{5.1.4})의 가정 아래에서 뇌터 귀납법의
원리 ($0_{\mathrm I}$, 2.2.2)를 사용하자. 즉 바탕공간이 $X$와
다른 $X$의 모든 닫힌 부분준스킴 $T$에 대해 정리가 참이라고
가정한다. 여기서 $\widehat T$는 물론 $T\cap X'$을 따른 $T$의
완비화이다. $X$는 공집합이 아니라고 가정할 수 있다. $f$가
분리이고 유한형이므로 초우 보조정리 (II, 5.6.1)를 적용할 수
있다. 따라서 $Y$-스킴 $Z$와 $Y$-사상 $g:Z\to X$가 존재하여,
구조 사상 $h:Z\to Y$는 준사영이고 $g$는 사영이며 전사이다.
더 나아가 $X$의 공집합이 아닌 열린집합 $U$가 존재하여 제한
$g^{-1}(U)\to U$는 동형사상이다.

$\mathcal M$을 바탕공간이 $X-U$인 닫힌 부분준스킴을 정의하는
$\mathcal O_X$의 연접 아이디얼이라 하고 (I, 5.2.2),
$\mathcal F$를 받침 $E$가 $Y$ 위에서 고유인 연접
$\mathcal O_{\widehat X}$-가군이라 하자. $\widehat Z$를
$h^{-1}(Y')$을 따른 $Z$의 완비화,
$\widehat g:\widehat Z\to\widehat X$를 $g$의 완비화들로의
연장이라 하자. 그러면 $\widehat g^{,*}(\mathcal F)$는 받침이
$g^{-1}(E)$에 포함되는 연접 $\mathcal O_{\widehat Z}$-가군이다.
$g$는 사영이고 따라서 고유이므로, 이 받침은 $Y$ 위에서 고유이다
(II, 5.4.6). $h$가 준사영이므로
$\widehat g^{,*}(\mathcal F)$는
(\hyperref[III.5.2.6-fr]{5.2.6})에 의해 대수화 가능하다. 따라서

\oldpage[III]{156}

$g$가 고유임과 (\hyperref[III.5.3.3-fr]{5.3.3})에 의해
\[
\widehat g_*\bigl(\widehat g^{,*}(\mathcal F)\bigr)
\]
은 대수화 가능한 $\mathcal O_{\widehat X}$-가군이다. 이제
$\mathcal F$와 $g$에 (\hyperref[III.5.3.4-fr]{5.3.4})를 적용할
수 있다. 준동형사상
\[
\rho_{\mathcal F}:\mathcal F\to
\widehat g_*\bigl(\widehat g^{,*}(\mathcal F)\bigr)
\]
의 핵 $\mathcal P$와 여핵 $\mathcal R$은 어떤 거듭제곱
$\widehat{\mathcal M}^{,n}$에 의해 소멸된다. $T$를
$\mathcal M^n$으로 정의되는 $X$의 닫힌 부분준스킴이라 하자.
그 바탕공간은 $X-U$이다. 또한 $j:T\to X$를 표준 단사라 하자.
그러면 완비화들로의 연장 $\widehat j:\widehat T\to\widehat X$는
표준 단사이다 (I, 10.14.7). 따라서
\[
\mathcal P=\widehat j_*\bigl(\widehat j^{,*}(\mathcal P)\bigr),
\qquad
\mathcal R=\widehat j_*\bigl(\widehat j^{,*}(\mathcal R)\bigr).
\]
$U$는 공집합이 아니므로 귀납가정에 의해
$\widehat j^{,*}(\mathcal P)$와
$\widehat j^{,*}(\mathcal R)$은 대수화 가능한
$\mathcal O_{\widehat T}$-가군들이다.
(\hyperref[III.5.3.3-fr]{5.3.3})에 의해 $\mathcal P$와
$\mathcal R$도 대수화 가능하다. 이제
(\hyperref[III.5.3.2-fr]{5.3.2})를 적용하면 마침내
$\mathcal F$가 대수화 가능함이 증명된다.

\hfill C.Q.F.D.
\end{env}

\subsection{응용: 통상적 스킴의 사상과 형식적 스킴의 사상의 비교.
대수화 가능한 형식적 스킴}
\label{subsection:III.5.4-fr}

\begin{theorem}[5.4.1]
\label{III.5.4.1-fr}
$A$를 뇌터 아딕 환, $\mathfrak J$를 $A$의 정의 아이디얼이라 하고
\[
S=\operatorname{Spec}(A),\qquad S'=\mathrm V(\mathfrak J)
\]
로 놓자. $u:X\to S$를 고유 사상, $v:Y\to S$를 분리 유한형
사상이라 하자. $\widehat S$, $\widehat X$, $\widehat Y$를 각각
$S'$, $u^{-1}(S')$, $v^{-1}(S')$을 따른 $S$, $X$, $Y$의
완비화라 하자. 모든 $S$-사상 $f:X\to Y$에 대해
$\widehat f:\widehat X\to\widehat Y$가 $f$의 완비화들로의
연장이면, 사상 $f\mapsto\widehat f$는 전단사
\[
\operatorname{Hom}_S(X,Y)\xrightarrow{\ \sim\ }
\operatorname{Hom}_{\widehat S}(\widehat X,\widehat Y)
\]
이다.
\end{theorem}

\begin{proof}
먼저 $f\mapsto\widehat f$가 단사임을 보이자. 실제로 두
$S$-사상 $f,g:X\to Y$가 $\widehat f=\widehat g$를 만족한다고
하자. 그러면 $f$와 $g$가 일치하는 $X'=u^{-1}(S')$의 어떤 열린
근방 $V$가 존재한다 (I, 10.9.4). 그런데 $u$는 닫힌 사상이므로
(\hyperref[III.5.1.3.1-fr]{5.1.3.1})에 의해 $V=X$이고, 따라서
$f=g$이다.

이제 $f\mapsto\widehat f$가 전사임을 증명하자. 이를 위해
$h:\widehat X\to\widehat Y$를 $\widehat S$-사상이라 하자.
$Z=X\times_S Y$로 놓고 $p:Z\to X$, $q:Z\to Y$를 표준 사영이라
하자. $Z$는 $S$ 위에서 유한형이므로 (I, 6.3.4) 뇌터이다.
$\widehat Z$를 $Z'=p^{-1}(u^{-1}(S'))$을 따른 $Z$의 완비화라
하자. $\widehat Z$는 표준적으로
$\widehat X\times_{\widehat S}\widehat Y$와 동일시되고,
$\widehat Z\to\widehat X$, $\widehat Z\to\widehat Y$의 사영들은
$\widehat p$, $\widehat q$와 동일시된다 (I, 10.9.7).
$Y$가 $S$ 위에서 분리이므로 $\widehat Y$는 $\widehat S$ 위에서
분리이다 (I, 10.15.7). 따라서 그래프 사상
\[
\Gamma_h=(1_{\widehat X},h):\widehat X\to\widehat Z
\]
은 닫힌 몰입이다 (I, 10.15.4). 이 몰입에 대응하는
$\widehat Z$의 닫힌 형식적 부분준스킴을 $\mathfrak T$라 하고,
$j:\mathfrak T\to\widehat Z$를 표준 단사라 하자. 그러면
\[
\Gamma_h=j\circ w,
\]
여기서 $w:\widehat X\to\mathfrak T$는 동형사상이고
(I, 10.14.3), 그 역동형사상은 $\widehat p\circ j$이다. 또한
$\widehat X$가 $\widehat S$ 위에서 고유이므로 $\mathfrak T$도
분명히 $\widehat S$ 위에서 고유이다. 따라서
(\hyperref[III.5.1.8-fr]{5.1.8})에 의해 $Z$의 닫힌 부분준스킴
$T$가 존재하여
\[
\mathfrak T=\widehat T=T_{/(T\cap Z')},
\]
이고 $j=\widehat i$이다. 여기서 $i:T\to Z$는 표준 단사이다
(I, 10.14.7). 사상 $p\circ i:T\to X$는 동형사상이다. 실제로
가정에 의해 그 완비화
\[
\widehat{p\circ i}=\widehat p\circ\widehat i
\]
가 동형사상이므로, 위에서와 같이 $S$가 $S'$의 $S$ 안에서의
유일한 근방임을 주의하고

\oldpage[III]{157}

(4.6.8)을 적용하면 된다. $g:X\to T$를 $p\circ i$의
역동형사상이라 하고 $f=q\circ i\circ g$로 놓자. 이는
$X\to Y$인 사상이며, 정의에 의해 그 그래프는
$\Gamma_f=i\circ g$이다. $\widehat g$는
$\widehat{p\circ i}=w$의 역동형사상이므로
\[
\widehat{\Gamma_f}=\widehat i\circ\widehat g
=j\circ w=\Gamma_h.
\]
한편 $\widehat{\Gamma_f}=\Gamma_{\widehat f}$임을 알고 있으므로
(I, 10.9.8), 마침내 $h=\widehat f$이다. 이로써 증명이 끝난다.
\end{proof}

따라서 범주의 언어로 말하면, 모든 뇌터 아딕 환 $A$에 대해
함자 $X\mapsto\widehat X$는 $\operatorname{Spec}(A)$ 위의 고유
스킴들의 범주에서 $\operatorname{Spf}(A)$ 위의 고유 형식적
스킴들의 범주로 가는 충만하고 충실한 함자이다 ($0$, 8.1.6).
그러므로 이 함자는 전자의 범주와 후자의 어떤 부분범주 사이의
동치를 정한다. 이 부분범주의 대상들을 
\emph{대수화 가능한 형식적 스킴}이라 하겠다. 그러한 형식적 스킴
$\mathfrak X$에 대해, $\mathfrak X$가 $\widehat X$와 동형이 되게
하는 $\operatorname{Spec}(A)$ 위의 고유한 통상적 스킴 $X$가
유일한 동형사상까지 존재한다.

\begin{env}[5.4.2]
\label{III.5.4.2-fr}
\emph{보충.}
(\hyperref[III.5.4.1-fr]{5.4.1})의 표기를 사용하여
\[
S_n=\operatorname{Spec}(A/\mathfrak J^{n+1}),\qquad
X_n=X\times_S S_n,\qquad Y_n=Y\times_S S_n
\]
로 놓자. (\hyperref[III.5.4.1-fr]{5.4.1})과 (I, 10.12.3)에
의하면, $S$-사상 $f:X\to Y$를 주는 것은 $S_n$-사상
$f_n:X_n\to Y_n$들로 이루어진 아딕 $(S_n)$-귀납계를 주는 것과
동치이다 (I, 10.12.2).
\end{env}

\begin{remark}[5.4.3]
\label{III.5.4.3-fr}
존재 정리 (\hyperref[III.5.1.6-fr]{5.1.6})가 암시하는 듯한 것과
달리, $\operatorname{Spf}(A)$ 위에서 고유이지만 대수화 가능하지
않은 형식적 스킴들이 존재한다. 이는 복소 대수다양체에서 나오지
않는 콤팩트 해석공간이 존재하는 것과 마찬가지이다. 뒤에서 이러한
스킴들을 ``모듈라이 이론''에서 만나게 될 것이다. 이 이론은
(기초체가 $\mathbb C$일 때) 완비 대수다양체의 복소 구조의 무한소
변형을 바로 다루며, 그러한 변형이 대수적이지 않은 해석다양체를
낳을 수 있음이 알려져 있다.
\end{remark}

\begin{proposition}[5.4.4]
\label{III.5.4.4-fr}
$A$를 뇌터 아딕 환, $\mathfrak S=\operatorname{Spf}(A)$라 하고,
$g:\mathfrak X\to\mathfrak S$와
$h:\mathfrak Y\to\mathfrak S$를 형식적 스킴들의 두 고유 사상,
$f:\mathfrak X\to\mathfrak Y$를 $\mathfrak S$-사상이라 하자.
$f$가 유한이고 $\mathfrak Y$가 대수화 가능하면
$\mathfrak X$도 대수화 가능하다.
\end{proposition}

\begin{proof}
$g$와 $h$에 대한 가정만으로도 이미 $f$는 고유임을 주의하자
(3.4.1). 또한 $f$가 유한이기 위해서는 모든
$y\in\mathfrak Y$에 대해 올 $f^{-1}(y)$이 유한이면 충분하다
(\hyperref[III.4.8.11-fr]{4.8.11}). 가정에 의해
\[
\mathcal B=f_*(\mathcal O_{\mathfrak X})
\]
는 연접 $\mathcal O_{\mathfrak Y}$-대수이다
(\hyperref[III.4.8.6-fr]{4.8.6}). 이제
$\mathfrak Y=\widehat Y$이고 $h=\widehat w$이며
$w:Y\to\operatorname{Spec}(A)$가 통상적 스킴들의 고유 사상이라
하면, 존재 정리에 의해
\[
\mathcal B=\widehat{\mathcal C}
\]
를 만족하는 연접 $\mathcal O_Y$-대수 $\mathcal C$가 존재한다.
$X=\operatorname{Spec}(\mathcal C)$라 하고 $u:X\to Y$를 구조
사상이라 하자. 그러면 $\mathcal B$로부터 $\mathfrak X$를
정의하는 방식 (\hyperref[III.4.8.7-fr]{4.8.7})에서 곧바로
$\mathfrak X$가 표준적으로 $\widehat X$와 동형이고 $f$가
$\widehat u$와 동일시됨을 알 수 있다. 이를 확인하려면 $Y$가
아핀인 경우를 생각하면 충분하다.

(\hyperref[III.5.1.8-fr]{5.1.8})은
(\hyperref[III.5.4.4-fr]{5.4.4})의 특수한 경우임을 주의하자.
\end{proof}

\begin{theorem}[5.4.5]
\label{III.5.4.5-fr}
$A$를 뇌터 아딕 환, $\mathfrak J$를 $A$의 정의 아이디얼이라 하고,
\[
S=\operatorname{Spec}(A),\qquad
\mathfrak S=\widehat S=\operatorname{Spf}(A)
\]
로 놓자. $f:\mathfrak X\to\mathfrak S$를 형식적 스킴들의 고유
사상이라 하자. 다음과 같이 놓는다.
\[
S_k=\operatorname{Spec}(A/\mathfrak J^{k+1}),\qquad
X_k=\mathfrak X\times_{\mathfrak S}S_k,
\]
그리고 모든 $\mathcal O_{\mathfrak X}$-가군 $\mathcal F$에 대해
\[
\mathcal F_k
=\mathcal F\otimes_{\mathcal O_{\mathfrak X}}\mathcal O_{X_k}
=\mathcal F/\mathfrak J^{k+1}\mathcal F.
\]

\oldpage[III]{158}

$\mathcal L$을 가역 $\mathcal O_{\mathfrak X}$-가군이라 하고,
$\mathcal L_0=\mathcal L/\mathfrak J\mathcal L$이 풍부한
$\mathcal O_{X_0}$-가군이라고 가정하자. 그러면 $\mathfrak X$는
대수화 가능하다. 또한 $X$가 $\mathfrak X\simeq\widehat X$를
만족하는 고유 $S$-스킴이면, $\mathcal L\simeq\widehat{\mathcal M}$을
만족하는 풍부한 $\mathcal O_X$-가군 $\mathcal M$이 존재한다.
따라서 $X$는 $S$ 위에서 사영이다.
\end{theorem}

\begin{proof}
(\hyperref[III.5.2.3-fr]{5.2.3})을
$\mathcal F=\mathcal O_{\mathfrak X}$에 적용하자. 그러면 어떤 정수
$n_0$가 존재하여 모든 $n\geq n_0$에 대해 표준 준동형사상
\[
\Gamma(\mathfrak X,\mathcal L^{\otimes n})
\longrightarrow\Gamma(X_0,\mathcal L_0^{\otimes n})
\]
은 전사이다. $\mathcal L_0^{\otimes n}$이 $S_0$에 대해 매우
풍부하도록 $n\geq n_0$를 충분히 크게 택할 수 있다
(II, 4.5.10). 사상 $f_0:X_0\to S_0$가 고유이므로
$\Gamma(X_0,\mathcal L_0^{\otimes n})$은 유한형 $A$-가군이다
(3.2.1). 따라서
$\Gamma(\mathfrak X,\mathcal L^{\otimes n})$의 어떤 유한형
부분 $A$-가군 $E$가 존재하여, 그
$\Gamma(X_0,\mathcal L_0^{\otimes n})$ 안에서의 상은 이 가군
전체이다.

이제 모든 $k\geq0$에 대해 표준 단사
$E\to\Gamma(\mathfrak X,\mathcal L^{\otimes n})$에서 유도되는
준동형사상
\[
u_k:E/\mathfrak J^{k+1}E\longrightarrow
\Gamma(X_k,\mathcal L_k^{\otimes n})
\]
을 생각하자. $(f_k)_*(\mathcal L_k^{\otimes n})$은 준연접이고
\[
\Gamma(S_k,(f_k)_*(\mathcal L_k^{\otimes n}))
=\Gamma(X_k,\mathcal L_k^{\otimes n})
\]
이므로, $u_k$는 준동형사상
\[
\widetilde u_k:(E/\mathfrak J^{k+1}E)^\sim
\longrightarrow(f_k)_*(\mathcal L_k^{\otimes n})
\]
를 정하고, 따라서 준동형사상
\[
\widetilde u_k^\sharp:
f_k^*((E/\mathfrak J^{k+1}E)^\sim)
\longrightarrow\mathcal L_k^{\otimes n}
\]
도 정한다. 한편
\[
\mathcal G_k=f_k^*((E/\mathfrak J^{k+1}E)^\sim)
\]
로 놓으면 $\mathcal G_0=\mathcal G_k/\mathfrak J\mathcal G_k$이다
(I, 9.1.5). 따라서
\[
\widetilde u_0^\sharp:
\mathcal G_k/\mathfrak J\mathcal G_k
\longrightarrow
\mathcal L_k^{\otimes n}/\mathfrak J\mathcal L_k^{\otimes n}
\]
은 $\widetilde u_k^\sharp$에서 몫으로 넘어가서 얻어진다. 그런데
$E$의 정의에 의해 $\widetilde u_0^\sharp$는 바로 표준 준동형사상
\[
\sigma:f_0^*((f_0)_*(\mathcal L_0^{\otimes n}))
\longrightarrow\mathcal L_0^{\otimes n}
\]
이다. $\mathcal L_0^{\otimes n}$이 매우 풍부하다는 가정에 의해
$\widetilde u_0^\sharp$는 전사이다 (II, 4.4.3). 그러므로 나카야마
보조정리에 의해 모든 $\widetilde u_k^\sharp$도 전사이다.

따라서 각각의 $\widetilde u_k^\sharp$는 $S_k$-사상
\[
g_k:X_k\longrightarrow P_k
=\mathbf P(E/\mathfrak J^{k+1}E)
\]
를 정한다 (II, 4.2.2). (II, 4.1.3)에 의해 $h\leq k$이면
$P_h=P_k\times_{S_k}S_h$이다. 그러므로
$\widetilde u_k^\sharp$들 사이의 관계와 (II, 4.2.10)에 의해
$(g_k)$는 아딕 $(S_k)$-귀납계이다 (I, 10.12.2). 따라서 $g_k$들은
형식적 스킴들의 $\mathfrak S$-사상
\[
g:\mathfrak X\to\mathfrak P
\]
를 정한다. 여기서 $\mathfrak P$는 계 $(P_k)$의 귀납극한이며,
동일하게 $P=\mathbf P(E)$의 완비화 $\widehat P$이다.
$\mathcal L_0^{\otimes n}$이 매우 풍부하다는 가정에 의해 $g_0$는
닫힌 몰입이다 (II, 4.4.3). 따라서 $g$는 형식적 스킴들의 닫힌
몰입이고 (4.8.10), 그러므로 $\mathfrak X$는 대수화 가능하다
(5.1.8).

이제 존재 정리 (5.1.6)에 의해 $\mathcal L$은 어떤 가역
$\mathcal O_X$-가군 $\mathcal M$의 완비화
$\widehat{\mathcal M}$과 동형이다. 또한
$\mathcal L^{\otimes n}$은 $\mathcal M^{\otimes n}$의 완비화이다
(I, 10.8.10). 준동형사상 $\widetilde u_k^\sharp$들은 잘 정의된
준동형사상
\[
v:f^*(\widetilde E)\longrightarrow\mathcal M^{\otimes n}
\]
을 정한다 (5.1.7). $\widetilde u_k^\sharp$가 전사이므로
$\widehat v$도 전사이고 (I, 10.11.5), 따라서 $v$도 전사이다
(5.1.3). $v$가 정하는 사상 $r:X\to P$의 (II, 4.2.2)
완비화들로의 연장은 $g$이다. $g$가 닫힌 몰입이므로 (5.1.8)과
(5.4.1)에 의해 $r$도 닫힌 몰입이다. 따라서
$\mathcal M^{\otimes n}$은 매우 풍부하고 (II, 4.4.6),
$\mathcal M$은 풍부하다 (II, 4.5.10).
\end{proof}

\begin{remark}[5.4.6]
\label{III.5.4.6-fr}
$A$를 뇌터 아딕 환, $\mathfrak S=\operatorname{Spf}(A)$라 하고,
$f:\mathfrak X\to\mathfrak S$를 형식적 스킴들의 고유 사상이라
하자. 모든 아핀 형식적 열린집합 $U$에 대해
$\Gamma(U,\mathcal N)$이
$\Gamma(U,\mathcal O_{\mathfrak X})$의 멱영근기가 되게 하는
$\mathcal O_{\mathfrak X}$의 연접 아이디얼을 $\mathcal N$이라
하자. 이 아이디얼의 존재는 (I, 10.10.2)와 모든 환 준동형사상
$B\to C$가 $B$의 멱영근기를 $C$의 멱영근기 안으로 보낸다는
사실에서 쉽게 따른다. $\mathfrak X'$을 $\mathcal N$으로 정의되는
$\mathfrak X$의 닫힌 형식적 부분스킴이라 하자 (I, 10.14.2).
$\mathfrak X'$이 대수화 가능할 때 $\mathfrak X$ 자체도 대수화
가능한지 아는 것은 흥미로운 문제이다. 구조층
$\mathcal O_{\mathfrak X}$을 제곱이 영인 아이디얼로 확장한 것들,
즉 $\mathcal J$가 $\mathcal A$의 제곱이 영인 아이디얼이고
\[
\mathcal O_{\mathfrak X}\simeq\mathcal A/\mathcal J
\]
인 $\mathcal O_{\mathfrak X}$-대수 $\mathcal A$들을, 예를 들어

\oldpage[III]{159}

코호몰로지적 성격의 불변량으로 분류할 수 있다면 이 문제의 해답에
도달할 수 있을 것이다. 여기서 $\mathfrak X$는 통상적 준스킴일
수도 있고 형식적 준스킴일 수도 있다.
\end{remark}

\subsection{어떤 스킴들의 분해}
\label{subsection:III.5.5-fr}

\begin{proposition}[5.5.1]
\label{III.5.5.1-fr}
$A$를 뇌터 아딕 환, $\mathfrak J$를 $A$의 정의 아이디얼이라 하고
$Y=\operatorname{Spec}(A)$로 놓자. $f:X\to Y$를 분리 유한형
사상이라 하고
\[
Y_0=\operatorname{Spec}(A/\mathfrak J),\qquad
X_0=X\times_Y Y_0=f^{-1}(Y_0)
\]
로 놓자. $Z_0$를 $Y_0$ 위에서 고유인 $X_0$의 열린 부분이라
하자. 그러면 $X$ 안에 $Y$ 위에서 고유이고
$Z\cap X_0=Z_0$를 만족하는 열린닫힌 부분 $Z$가 존재한다.
\end{proposition}

\begin{proof}
가정에 의해 $T\cap X_0=Z_0$인 $X$의 열린집합 $T$가 존재한다.
$\widehat T$를 열린집합 $T$ 위에 $X$가 유도하는 스킴의 $Z_0$을
따른 완비화라 하자. $\mathcal O_{\widehat T}$의 받침은 $Y_0$
위에서 고유인 $Z_0$이므로, $\widehat T$는
$\widehat Y=\operatorname{Spf}(A)$ 위에서 고유이다 (3.4.1).
(5.1.8)에 의해 $Y$ 위에서 고유인 $T$의 닫힌 부분스킴 $Z$가
존재하여, 표준 단사 $i:Z\to T$의 완비화
\[
\widehat i:\widehat Z\to\widehat T
\]
는 동형사상이다. 여기서 $\widehat Z$는 $Z_0$을 따른 $Z$의
완비화이다. 따라서 (4.6.8)에 의해 $Z_0$의 $T$ 안에서의 어떤
열린 근방 $V$가 존재하여 $i$의 제한
$i^{-1}(V)\to V$는 동형사상이다. 그런데 $i^{-1}(V)$는 $Z$ 안에서
$Z_0$의 근방이므로 반드시 $Z$와 같다 (5.1.3.1). 따라서 $Z$는
$T$에서 열려 있고, 그러므로 $X$에서도 열려 있다. 이로써 증명이
끝난다.
\end{proof}

\begin{corollary}[5.5.2]
\label{III.5.5.2-fr}
$X_0$가 $Y_0$ 위에서 고유이면 $X$는 서로소인 두 열린 부분
$Z$와 $Z'$의 합집합이다. 여기서 $Z$는 $Y$ 위에서 고유이고
$X_0$를 포함한다. 또한 $Y$ 위에서 고유인 $X$의 모든 닫힌 부분
$P$는 $Z$에 포함된다.
\end{corollary}

\begin{proof}
마지막 주장은 다음 사실에서 따른다. $P\cap Z'$은 $P$에서 닫혀
있으므로 $Y$ 위에서 고유이다. 만일 $P\cap Z'$이 공집합이 아니면
$f(P\cap Z')$은 $Y$의 공집합이 아닌 닫힌집합이므로 $Y_0$와
만나야 한다 (5.1.3.1). 이는 $Z$의 정의에 모순된다.
\end{proof}

\begin{flushright}
\emph{(계속.)}
\end{flushright}

% 인쇄본 504쪽은 백지이다.
\clearpage
\thispagestyle{empty}
\null
\clearpage
\thispagestyle{plain}

\oldpage[III]{161}

\section*{참고문헌 (계속)}

\begin{flushleft}
[27] P. Gabriel, \emph{Des catégories abéliennes}, thèse, Paris (1961).

\medskip
[28] J. E. Roos, \emph{Sur les foncteurs dérivés de $\varprojlim$.
Applications}, C. R. Acad. Sci. Paris, t. CCLII, 1961, p. 3702--3704.

\medskip
[29] H. Cartan, \emph{Séminaire de l'École Normale Supérieure},
13\textsuperscript{e} année (1960--1961), exposé n\textsuperscript{o} 11.
\end{flushleft}

\clearpage
\oldpage[III]{162}
\section*{기호 색인}

\begin{flushleft}
$\mathbf{Ens}$ : 0, 8.1.1.

$h_X(Y)$, $h_X(u)$ ($X$, $Y$는 범주 $C$의 대상, $u$는 $C$의
사상) : 0, 8.1.1.

$h_w(Y)$ ($Y$는 $C$의 대상, $w$는 $C$의 사상) : 0, 8.1.2.

$Z_k(E_r^{pq})$, $B_k(E_r^{pq})$ ($k\geq r$) : 0, 11.1.1.

$Z_\infty(E_2^{pq})$, $B_\infty(E_2^{pq})$ : 0, 11.1.1.

$K^\bullet$, $K_\bullet$ (복합체) : 0, 11.2.1.

$K^{\bullet\bullet}$, $K_{\bullet\bullet}$ (이중복합체) : 0, 11.2.1.

$B^i(K^\bullet)$, $Z^i(K^\bullet)$, $H^i(K^\bullet)$,
$H^\bullet(K^\bullet)$ : 0, 11.2.1.

$B_i(K_\bullet)$, $Z_i(K_\bullet)$, $H_i(K_\bullet)$,
$H_\bullet(K_\bullet)$ : 0, 11.2.1.

$T(K^\bullet)$, $T(K_\bullet)$ ($T$는 함자) : 0, 11.2.1.

$E(K^\bullet)$ ($K^\bullet$은 여과 복합체) : 0, 11.2.2.

$K^{i,\bullet}$, $K^{\bullet,j}$, $Z_{\mathrm{II}}^p(K^{i,\bullet})$,
$B_{\mathrm{II}}^p(K^{i,\bullet})$, $H_{\mathrm{II}}^p(K^{i,\bullet})$,
$Z_{\mathrm I}^p(K^{\bullet,j})$, $B_{\mathrm I}^p(K^{\bullet,j})$,
$H_{\mathrm I}^p(K^{\bullet,j})$ : 0, 11.3.1.

$H_{\mathrm{II}}^p(K^{\bullet\bullet})$,
$Z_{\mathrm I}^q(H_{\mathrm{II}}^p(K^{\bullet\bullet}))$,
$B_{\mathrm I}^q(H_{\mathrm{II}}^p(K^{\bullet\bullet}))$,
$H_{\mathrm I}^q(H_{\mathrm{II}}^p(K^{\bullet\bullet}))$,
$H^n(K^{\bullet\bullet})$ : 0, 11.3.1.

$F_{\mathrm I}^p(K^{\bullet\bullet})$,
$F_{\mathrm{II}}^p(K^{\bullet\bullet})$, ${}'E(K^{\bullet\bullet})$,
${}''E(K^{\bullet\bullet})$ : 0, 11.3.2.

$K_{i,\bullet}$, $K_{\bullet,j}$, $H_p^{\mathrm{II}}(K_{i,\bullet})$,
$H_p^{\mathrm I}(K_{\bullet,j})$, $H_p^{\mathrm{II}}(K_{\bullet\bullet})$,
$H_p^{\mathrm I}(K_{\bullet\bullet})$,
$H_q^{\mathrm I}(H_p^{\mathrm{II}}(K_{\bullet\bullet}))$,
$H_q^{\mathrm{II}}(H_p^{\mathrm I}(K_{\bullet\bullet}))$,
$H_n(K_{\bullet\bullet})$ : 0, 11.3.5.

$\mathrm R^\bullet T(K^\bullet)$ ($T$는 공변 가법 함자) : 0, 11.4.3.

$\mathrm R^\bullet T(K^\bullet,K'^{\bullet})$ ($T$는 공변 가법
쌍함자) : 0, 11.4.6.

$\mathrm L_\bullet T(K_\bullet)$ ($T$는 공변 가법 함자) : 0, 11.6.2.

$\mathrm L_\bullet T(K_\bullet,K'_\bullet)$ ($T$는 공변 가법
쌍함자) : 0, 11.6.6.

$\mathrm L_\bullet T(K_{\bullet\bullet})$ ($T$는 공변 가법 함자) :
0, 11.7.2.

$|\sigma|$, $\Sigma(A)$, $\mathrm C_\bullet(A)$, $\mathrm L_\bullet(A)$
($A$는 유한집합, $\sigma$는 $A$의 단체) : 0, 11.8.1.

$\mathrm C^\bullet(A,B;\mathscr S)$, $\mathrm L^\bullet(A,B;\mathscr S)$
($A$, $B$는 유한집합, $\mathscr S$는 계수계) : 0, 11.8.4.

$\mathrm P^\bullet(A,B,\mathscr S)$ ($A$, $B$는 유한집합,
$\mathscr S$는 계수계) : 0, 11.8.6.

$\mathrm C^\bullet(A,B;\mathscr S^\bullet)$,
$\mathrm L^\bullet(A,B;\mathscr S^\bullet)$ ($A$, $B$는 유한집합,
$\mathscr S^\bullet$은 계수계들의 복합체) : 0, 11.8.8.

$\mathrm P^\bullet(A,B;\mathscr S^\bullet)$ ($A$, $B$는 유한집합,
$\mathscr S^\bullet$은 계수계들의 복합체) : 0, 11.8.10.

$\chi(M^\bullet)$ ($M^\bullet$은 유한 길이 가군들로 이루어진 유한
복합체) : 0, 11.10.1.

$\check C^\bullet(\mathfrak U,\mathcal F)$ ($\mathcal F$는 $X$ 위의
아벨군의 층, $\mathfrak U$는 $X$의 열린 덮개) : 0, 12.1.2.

$\mathrm R^pf_*(\mathcal F)$ ($f:X\to Y$는 환 달린 공간들의 사상,
$\mathcal F$는 $\mathcal O_X$-가군의 층) : 0, 12.2.1.

$\mathcal H^p(f,\mathcal K^\bullet)$,
$\mathcal H_f^p(\mathcal K^\bullet)$, ${}'E(f,\mathcal K^\bullet)$,
${}''E(f,\mathcal K^\bullet)$ ($f:X\to Y$는 환 달린 공간들의 사상,
$\mathcal K^\bullet$은 $\mathcal O_X$-가군들의 복합체) : 0, 12.4.1.

$\mathbf H^\bullet(X,\mathcal K^\bullet)$,
$\mathbf H^\bullet(V,\mathcal K^\bullet)$ ($X$는 환 달린 공간,
$\mathcal K^\bullet$은 $\mathcal O_X$-가군들의 복합체, $V$는
$X$의 열린집합) : 0, 12.4.2.

$\operatorname{gr}^\bullet(A)$ ($A$는 조건 (ML)을 만족하는 사영계) :
0, 13.4.2.

$E(A)$ ($A$는 여과 대상) : 0, 13.6.4.

$K_\bullet(\mathbf f)$, $i_{\mathbf f}$ ($\mathbf f$는 환의 원소들로
이루어진 유한족) : III, 1.1.1.

$K^\bullet(\mathbf f,M)$, $K_\bullet(\mathbf f,M)$ ($\mathbf f$는 환
$A$의 원소들로 이루어진 유한족, $M$은 $A$-가군) : III, 1.1.2.

$H^\bullet(\mathbf f,M)$, $H_\bullet(\mathbf f,M)$ ($\mathbf f$는 환
$A$의 원소들로 이루어진 유한족, $M$은 $A$-가군) : III, 1.1.3.

$C^\bullet((\mathbf f),M)$, $H^\bullet((\mathbf f),M)$ ($\mathbf f$는
환 $A$의 원소들로 이루어진 유한족, $M$은 $A$-가군) : III, 1.1.6.

$H^\bullet(U,\mathcal F^{(*)})$,
$H^\bullet(\mathfrak U,\mathcal F^{(*)})$ ($\mathcal F$는 준연접
$\mathcal O_X$-가군, $U$는 $X$의 열린집합, $\mathfrak U$는
$X$의 열린 덮개) : III, 2.1.1.
\end{flushleft}

\clearpage
\oldpage[III]{163}

\section*{용어 색인}

\begin{flushleft}
스펙트럼 열의 수렴 대상 : 0, 11.1.1.

대수화 가능한 ($\mathcal O_X$-가군) : III, 5.2.1.

대수화 가능한 (형식적 스킴) : III, 5.4.2.

해석적으로 정역인 (환) : III, 4.3.6.

준콤팩트 사상 : 0, 9.1.1.

분해의 증대 : 0, 11.4.1.

복합체의 오일러--푸앵카레 특성 : 0, 11.10.1.

연접 $\mathcal O_X$-가군의 오일러--푸앵카레 특성
($X$는 $\operatorname{Spec}(A)$ 위에서 사영, $A$는 아르틴 환) :
III, 2.5.1.

쌍교대 공사슬 : 0, 11.8.4.

이중복합체가 정하는 복합체 : 0, 11.3.1.

외대수의 복합체 : III, 1.1.1.

조건 (ML) : 0, 13.1.2.

구성가능인 (부분, 집합) : 0, 9.1.2.

구성가능인 (함수) : 0, 9.3.1.

컵곱 : 0, 12.1.2.

범주에서의 대수적 구조에 관한 디-준동형 : 0, 8.2.1.

아벨 범주에서 완전한 (부분집합) : III, 3.1.1.

본질적으로 상수인 (사영계) : 0, 13.4.2.

스테인 인자분해 : III, 4.3.3.

공이산, 공분리, 이산, 소진적, 유한, 분리인 여과 : 0, 11.1.3.

범주의 끝 (대상) : 0, 8.1.10.

형식적 준스킴의 유한 (사상) : III, 4.8.2.

$\mathfrak Y$ 위에서 유한인 (형식적 준스킴),
$\mathfrak Y$-유한인 (형식적 준스킴) : III, 4.8.2.

표준 공변 함자 $C\to\operatorname{Hom}(C^o,\mathbf{Ens})$ : 0, 8.1.2.

산술 종수 : III, 2.5.1.

기하적으로 연결인 (올) : III, 4.3.4.

범주에서의 대수적 구조에 관한 준동형 : 0, 8.2.1.

복합체 $K^\bullet$에 관한 함자의 초코호몰로지 : 0, 11.4.3.

두 복합체 $K^\bullet$, $K'^{\bullet}$에 관한 쌍함자의
초코호몰로지 : 0, 11.4.6.

복합체 $K_\bullet$에 관한 함자의 초호몰로지 : 0, 11.6.2.

두 복합체 $K_\bullet$, $K'_\bullet$에 관한 쌍함자의 초호몰로지 :
0, 11.6.6.

이중복합체에 관한 함자의 초호몰로지 : 0, 11.7.2.

국소 구성가능인 (집합) : 0, 9.1.11.

범주의 대상 위의 외부 (내부) 합성법칙 : 0, 8.2.1.

\oldpage[III]{164}

여과 $S$-$C$-가군, 등급 $\operatorname{gr}^\bullet(S)$-$C$-가군,
이중등급 $\operatorname{gr}^{\bullet\bullet}(S)$-$C$-가군 : 0, 13.6.6.

스펙트럼 열의 사상 : 0, 11.1.2.

올의 기하적 연결 성분 수 : III, 4.3.4.

$C$의 군 대상, $C$-군, $C$-환, $C$-가군 : 0, 8.2.3.

경계 대상, 공경계 대상, 공순환 대상, 순환 대상 : 0, 11.2.1.

보편상 대상 : 0, 13.1.1.

조건 (ML)을 만족하는 사영계의 연관 등급 대상 : 0, 13.4.2.

아래로 (위로) 유계인 등급 대상 : 0, 11.2.1.

형식적 준스킴의 ($\mathfrak Y$ 위에서) 고유인 부분 : III, 3.4.1.

충만한 (부분범주) : 0, 8.1.5.

충만하고 충실한 (함자) : 0, 8.1.5.

힐베르트 다항식 : III, 2.5.3.

형식적 준스킴의 고유 (사상) : III, 3.4.1.

표현 가능한 (함자) : 0, 8.1.8.

코호몰로지, 오른쪽, 왼쪽, 호몰로지, 단사, 자유, 평탄, 사영 분해 :
0, 11.4.1.

오른쪽 카르탕--아일렌베르크 분해, 단사 카르탕--아일렌베르크 분해 :
0, 11.4.2.

왼쪽 카르탕--아일렌베르크 분해, 사영 카르탕--아일렌베르크 분해 :
0, 11.6.1.

역콤팩트인 (집합) : 0, 9.1.1.

엄밀한 (사영계) : 0, 13.4.2.

아벨 범주에서의 스펙트럼 열 : 0, 11.1.1.

약수렴, 정칙, 공정칙, 쌍정칙 스펙트럼 열 : 0, 11.1.3.

퇴화하는 스펙트럼 열 : 0, 11.1.6.

여과 대상에 관한 함자의 스펙트럼 열 : 0, 13.6.4.

이중복합체의 스펙트럼 열 : 0, 11.3.2 및 11.3.5.

복합체에 관한 함자의 초코호몰로지 스펙트럼 열 : 0, 11.4.3.

복합체에 관한 함자의 초호몰로지 스펙트럼 열 : 0, 11.6.2.

이중복합체에 관한 함자의 초호몰로지 스펙트럼 열 : 0, 11.7.2.

계수계 : 0, 11.8.4.

단일분기인(unibranche) (환, 점) : III, 4.3.6.

보편적으로 열린 (사상) : III, 4.3.9.
\end{flushleft}

\clearpage
\oldpage[III]{165}

\section*{목차}

\begin{flushleft}
\textsc{제0장.} --- \textbf{예비사항 (계속)} \dotfill 5

\medskip
\S~8. 표현 가능 함자 \dotfill 5

8.1. 표현 가능 함자 \dotfill 5

8.2. 범주에서의 대수적 구조 \dotfill 9

\medskip
\S~9. 구성가능 집합 \dotfill 12

9.1. 구성가능 집합 \dotfill 12

9.2. 뇌터 공간에서의 구성가능 집합 \dotfill 14

9.3. 구성가능 함수 \dotfill 16

\medskip
\S~10. 평탄 가군에 관한 보충 \dotfill 17

10.1. 평탄 가군과 자유 가군 사이의 관계 \dotfill 17

10.2. 평탄성의 국소 판정법 \dotfill 18

10.3. 국소환의 평탄 확대의 존재 \dotfill 20

\medskip
\S~11. 호몰로지 대수학에 관한 보충 \dotfill 23

11.1. 스펙트럼 열에 관한 복습 \dotfill 23

11.2. 여과 복합체의 스펙트럼 열 \dotfill 27

11.3. 이중복합체의 스펙트럼 열 \dotfill 29

11.4. 복합체 $K^\bullet$에 관한 함자의 초코호몰로지
\dotfill 32

11.5. 초코호몰로지에서 귀납극한으로의 이행 \dotfill 35

11.6. 복합체 $K_\bullet$에 관한 함자의 초호몰로지
\dotfill 39

11.7. 이중복합체 $K_{\bullet\bullet}$에 관한 함자의 초호몰로지
\dotfill 41

11.8. 단체 복합체의 코호몰로지에 관한 보충 \dotfill 43

11.9. 유한 생성 복합체에 관한 보조정리 \dotfill 46

11.10. 유한 길이 가군 복합체의 오일러--푸앵카레 특성
\dotfill 48

\medskip
\S~12. 층의 코호몰로지에 관한 보충 \dotfill 49

12.1. 환 달린 공간 위의 가군층의 코호몰로지 \dotfill 49

12.2. 고차 직상 \dotfill 57

12.3. 층의 Ext 함자에 관한 보충 \dotfill 60

12.4. 직상 함자의 초코호몰로지 \dotfill 62

\oldpage[III]{166}

\S~13. 호몰로지 대수학에서의 사영극한 \dotfill 64

13.1. 미타그--레플러 조건 \dotfill 64

13.2. 아벨군에 대한 미타그--레플러 조건 \dotfill 65

13.3. 응용: 층들의 사영극한의 코호몰로지 \dotfill 68

13.4. 미타그--레플러 조건과 사영계의 연관 등급 대상 \dotfill 69

13.5. 여과 복합체의 스펙트럼 열들의 사영극한 \dotfill 71

13.6. 유한 여과를 갖춘 대상에 관한 함자의 스펙트럼 열 \dotfill 73

13.7. 여러 변수의 사영극한의 유도 함자 \dotfill 75

\medskip
\textsc{제III장.} --- \textbf{연접층의 코호몰로지적 연구} \dotfill 81

\medskip
\S~1. 아핀 스킴의 코호몰로지 \dotfill 82

1.1. 외대수 복합체에 관한 복습 \dotfill 82

1.2. 열린 덮개의 체흐 코호몰로지 \dotfill 85

1.3. 아핀 스킴의 코호몰로지 \dotfill 88

1.4. 임의의 준스킴의 코호몰로지에 대한 응용 \dotfill 89

\medskip
\S~2. 사영 사상의 코호몰로지적 연구 \dotfill 95

2.1. 몇몇 코호몰로지 군의 명시적 계산 \dotfill 95

2.2. 사영 사상의 기본정리 \dotfill 100

2.3. 등급 대수 및 가군의 층에 대한 응용 \dotfill 102

2.4. 기본정리의 일반화 \dotfill 107

2.5. 오일러--푸앵카레 특성과 힐베르트 다항식 \dotfill 109

2.6. 응용: 풍부성 판정법 \dotfill 111

\medskip
\S~3. 고유 사상의 유한성 정리 \dotfill 115

3.1. 데비사주 보조정리 \dotfill 115

3.2. 유한성 정리: 통상적 스킴의 경우 \dotfill 116

3.3. 유한성 정리의 일반화 (통상적 스킴) \dotfill 118

3.4. 유한성 정리: 형식적 스킴의 경우 \dotfill 119

\medskip
\S~4. 고유 사상의 기본 정리와 그 응용 \dotfill 122

4.1. 기본 정리 \dotfill 122

4.2. 특수한 경우들과 변형들 \dotfill 129

4.3. 자리스키의 연결성 정리 \dotfill 130

4.4. 자리스키의 ``주정리'' \dotfill 135

4.5. 준동형사상 가군의 완비화 \dotfill 138

4.6. 형식 사상과 보통 사상의 관계 \dotfill 139

4.7. 풍부성 판정법 \dotfill 145

4.8. 형식적 준스킴의 유한 사상 \dotfill 146

\oldpage[III]{167}

\S~5. 연접 대수적 층의 존재 정리 \dotfill 149

5.1. 정리의 진술 \dotfill 149

5.2. 존재 정리의 증명: 사영인 경우와 준사영인 경우
\dotfill 151

5.3. 존재 정리의 증명: 일반적인 경우 \dotfill 154

5.4. 응용: 통상적 스킴의 사상과 형식적 스킴의 사상의 비교.
대수화 가능한 형식적 스킴 \dotfill 156

5.5. 어떤 스킴들의 분해 \dotfill 159

\medskip
\textsc{참고문헌 (계속)} \dotfill 161

\textsc{기호 색인} \dotfill 162

\textsc{용어 색인} \dotfill 163

\begin{flushright}
\emph{1961년 6월 28일 접수.}
\end{flushright}
\end{flushleft}
